% Preamble for the swarm-transcribed Weber volumes.
% Deliberately close to weber_v1_ge.tex so the verified p1-148 region and the swarm
% region can live in one document.
\documentclass[11pt]{article}
\usepackage[ngerman]{babel}
\usepackage[a4paper,margin=0.9in]{geometry}
\usepackage[T1]{fontenc}
\usepackage[utf8]{inputenc}
\usepackage{lmodern}
\usepackage[leqno]{amsmath}
\usepackage{amssymb,mathtools,mathrsfs}
\usepackage{microtype}
\usepackage{fancyhdr}
\usepackage{array,longtable,enumitem,multicol}
\usepackage{tikz}
\setlength{\parindent}{1.5em}
\setlength{\parskip}{0.18em}
\emergencystretch=6em
\setlength{\headheight}{14pt}
\pagestyle{fancy}
\fancyhf{}
\cfoot{\thepage}

% --- Weber conventions -------------------------------------------------------
% Sperrung (letter-spacing) is Weber's emphasis. There is no bold in his body text.
\providecommand{\gesperrt}[1]{\textls[130]{#1}}
% Printed-page marker: a no-op in the PDF. The machine-readable anchor is the
% "% !PRINTED n" comment in the source, which the collector and Codex both read.
\providecommand{\weberpage}[1]{}
% Displayed equation numbered in the LEFT margin, as Weber prints them.
\providecommand{\weq}[2]{\begin{equation}\tag{#1}#2\end{equation}}
% Figure placeholder. A later pass replaces it with TikZ or a cleaned scan crop.
% A marked gap is honest; a fabricated diagram is not.
\providecommand{\weberfigure}[2]{%
  \begin{center}\fbox{\parbox{0.78\linewidth}{\centering\small [Figur, S.~#1: #2]}}\end{center}}

% Text symbols agents reach for that are not always defined. Providing them in the preamble
% fixes every occurrence at once and leaves the generated fragments untouched, which is the
% right trade whenever the fix is "define the macro" rather than "change the text".
\usepackage{textcomp}
% amsmath caps matrix environments at 10 columns. Weber prints permutation matrices far wider
% than that -- Bd III has an 11-column pmatrix of the w_nu -- and the failure is the opaque
% "Extra alignment tab has been changed to \cr".
\setcounter{MaxMatrixCols}{30}
\providecommand{\textquotedblbase}{\glqq}
\providecommand{\textquotedblleft}{``}
\providecommand{\textquotedblright}{''}
\providecommand{\textquotesinglbase}{\glq}

% --- Greek capitals are SLANTED in this fount -------------------------------------------
% RULING 2026-08-07, after the same observation was reported from five separate pages
% (printed 271, 275, 285, 333 and the nu/v cousin on 172).
%
% Verified at 8x zoom on printed 333 (b1 idx 358): every capital Delta on that page --
% "f(x) - f(alpha) = Delta", "Delta_alpha = f(alpha+delta) - f(alpha)" -- is the INCLINED
% form, which LaTeX calls \varDelta. \Delta, LaTeX's default, is upright and is not what
% Vieweg set. The same inclination is visible on capital Phi (printed 171, 282).
%
% This is a property of the TYPEFACE, not of any occurrence -- the same kind of fact as the
% varkappa ruling -- so it holds wherever the letter appears, and the right fix is one line
% here rather than several thousand edits across the fragments. \let leaves every fragment
% untouched, applies to the whole corpus at once, and is reversible by deleting these lines.
%
% Delta and Phi are the two verified by eye. The rest follow from its being a single fount;
% if any capital is later found upright on the page, delete its line and the fragments still
% say exactly what they said before.
\let\Gamma\varGamma
\let\Delta\varDelta
\let\Theta\varTheta
\let\Lambda\varLambda
\let\Xi\varXi
\let\Pi\varPi
\let\Sigma\varSigma
\let\Upsilon\varUpsilon
\let\Phi\varPhi
\let\Psi\varPsi
\let\Omega\varOmega
% Lower-case rho is the OPEN-TAILED form in this fount. Verified at 10x on printed 419
% ("wenn also überhaupt ein \varrho vorhanden ist"): the tail curves down and left rather than
% descending straight, which is \varrho, not \rho. Same kind of fact as the capitals above, and
% the same one-line treatment. (kappa is already handled: the fragments write \varkappa
% explicitly, on the ruling recorded in resolve_class.py.)
\let\rho\varrho

% Common shorthands inherited from the volume I file, so fragments that use them compile.
\providecommand{\eps}{\varepsilon}
\providecommand{\ii}{\mathrm{i}}
\providecommand{\tg}{\operatorname{tg}}
\providecommand{\ctg}{\operatorname{cotg}}
\providecommand{\cotg}{\operatorname{cotg}}
\providecommand{\sn}{\operatorname{sn}}
\providecommand{\cn}{\operatorname{cn}}
\providecommand{\dn}{\operatorname{dn}}
\providecommand{\am}{\operatorname{am}}
\providecommand{\dd}{\,d}
\providecommand{\tht}{\vartheta}
\providecommand{\vark}{\varkappa}
\providecommand{\Norm}{\operatorname{N}}
\providecommand{\ord}{\operatorname{ord}}
\providecommand{\ind}{\operatorname{ind}}
\providecommand{\abs}[1]{\left|#1\right|}
\providecommand{\sect}[2]{\section*{\S.\,#1. #2}}
\providecommand{\frp}{\mathfrak p}
\providecommand{\frP}{\mathfrak P}
\providecommand{\frG}{\mathfrak G}
\providecommand{\frA}{\mathfrak A}
\providecommand{\frB}{\mathfrak B}

\begin{document}
\section*{Gordan--Kerschensteiner, Vorlesungen über Invariantentheorie --- Zweiter Band: Binäre Formen}
\noindent\textit{Vollständiger erster deutscher Transkriptionsdurchgang der erfassten Blätter; ausdrückliche Unsicherheiten bleiben; keine kanonische Neukollation.}\par\bigskip

% ===== pdf idx 15, printed 1.0 =====
% !PRINTED 1
\weberpage{1}
\begin{center}
\textbf{Erster Theil.}
\end{center}

\begin{center}
\textbf{Processse für invariante Bildungen.}
\end{center}

\begin{center}
\textbf{§ 1. Lineare Transformation und Faltungsprocess.}
\end{center}

1. \textit{Einführung der Symbolik.} Jeder homogene algebraische Ausdruck $n^{\text{ten}}$ Grades zwischen zwei Variabeln $x_1$ und $x_2$ wird als \glqq binäre Form\grqq{} bezeichnet. Eine solche binäre Form ist allgemein dargestellt durch:
\weq{1}{f(x) = ax_1{}^n + bx_1{}^{n-1}x_2 + cx_1{}^{n-2}x_2{}^2 + \cdots + qx_2{}^n\,.}
Viele Untersuchungen lassen es bequemer erscheinen, an Stelle der einfachen Coefficienten $a$, $b$, $c \ldots$ etc. Grössen mit Binomialcoefficienten multiplicirt einzuführen, also $f(x)$ in der Form zu schreiben:
\weq{2}{f(x) = a_0x_1{}^n + \binom{n}{1}\overline{a_1}x_1{}^{n-1}x_2 + \binom{n}{2}\overline{a_2}x_1{}^{n-2}x_2{}^2 + \cdots \overline{a_n}x_2{}^n\,.}
Man kann noch einen Schritt weiter gehen. Der algebraische Ausdruck rechts erinnert in mehr als einer Beziehung an die binomische Entwicklung
\weq{3}{(a_1x_1+a_2x_2)^n = a_1{}^nx_1{}^n + \binom{n}{1}a_1{}^{n-1}a_2x_1{}^{n-1}x_2 + \binom{n}{2}a_1{}^{n-2}a_2{}^2x_1{}^{n-2}x_2{}^2 + \ldots a_2{}^nx_2{}^n\,,}
welche mit der Darstellung (2) die Homogeneität in $x_1$ und $x_2$, die Zahl und Dimension der Glieder und die Binomialcoefficienten gemein hat. Fügen wir hinzu, dass die $(n+1)$ Constanten der Form $f(x)$ durch die $(n+1)$ Producte $a_1{}^n$, $a_1{}^{n-1}a_2$, $a_1{}^{n-2}a_2{}^2\ldots a_2{}^n$ charakterisirt sein mögen, so ist die binäre Form $f(x)$ auch vollständig definirt durch den symbolischen Ausdruck:
\[
f(x) = (\dot a_1x_1 + a_2x_2)^n = a_x{}^n\,,
\]
vermöge welchen
\[
\begin{array}{lcccl}
a_1{}^n & \text{als} & \text{Symbol} & \text{für} & \overline{a_0} \\
a_1{}^{n-1}a_2 & \text{\textquotedbl} & \text{\textquotedbl} & \text{\textquotedbl} & \overline{a_1} \\
a_1{}^{n-2}a_2{}^2 & \text{\textquotedbl} & \text{\textquotedbl} & \text{\textquotedbl} & \overline{a_2} \\
\multicolumn{5}{c}{\cdot\;\cdot\;\cdot\;\cdot\;\cdot\;\cdot\;\cdot\;\cdot\;\cdot\;\cdot} \\
a_2{}^n & \text{als} & \text{Symbol} & \text{für} & \overline{a_n}
\end{array}
\]
zu gelten hat.


% ===== pdf idx 16, printed 2.0 =====
% !PRINTED 2
\weberpage{2}
2. \textit{Mehrere Symbole für ein- und dieselbe Form $f$.} Diese eine Symbolreihe
\[
a_1{}^{n-k}a_2{}^k = \overline{a_k}
\]
genügt indess nicht, um jedes beliebige Aggregat aus den Coefficienten von $f$ in der Weise symbolisch darzustellen, dass man auch rückwärts aus dem symbolischen Ausdruck den wirklichen eindeutig zu bilden vermöchte. So wäre beispielsweise das Product $\overline{a_2}\cdot\overline{a_3}$ durch
\[
\overline{a_2}\cdot\overline{a_3} = a_1{}^{n-2}a_2{}^2\cdot a_1{}^{n-3}a_2{}^3 = a_1{}^{2n-5}a_2{}^5
\]
symbolisch darzustellen. Aus dem Ausdrucke $a_1{}^{2n-5}a_2{}^5$ lässt sich aber umgekehrt das Product der wirklichen Coefficienten nicht mehr eindeutig bestimmen. Derselbe kann auch noch als Symbol für $\overline{a_0}\cdot\overline{a_5}$ und $\overline{a_1}\cdot\overline{a_4}$ gelten. Diese Undeutlichkeit lässt sich aber vermeiden, wenn man neben der einen Symbolreihe $a_1{}^{n-k}a_2{}^k$ noch so viel weitere Symbolreihen $b_1{}^{n-k}b_2{}^k$, $c_1{}^{n-k}c_2{}^k$, \ldots etc. einführt, als das betreffende Glied des Coefficientenaggregates unsymbolische Factoren $\overline{a_k}$ besitzt. In diesem Falle sind die symbolischen Ausdrücke
\[
\overline{a_2}\cdot\overline{a_3} = a_1{}^{n-2}a_2{}^2\cdot b_1{}^{n-3}b_3 = c_1{}^{n-2}c_2\cdot d_1{}^{n-3}d_2{}^3 = \text{etc.}
\]
\[
\overline{a_4}{}^3 = a_1{}^{n-4}a_2{}^4\cdot b_1{}^{n-4}b_2{}^4\ c_1{}^{n-4}c_2{}^4 = \text{etc.}\,, % ?? uncertain: printed spacing between b_2^4 and c_1^{n-4} shows no visible multiplication dot/comma; reproduced as printed
\]
hin und zurück nur auf eine Weise zu deuten, in welcher Vertauschung auch die symbolischen Factoren rechts auftreten. Die Form $f$ selbst ist alsdann durch jedes der symbolischen Producte
\[
f = a_x{}^n = b_x{}^n = c_x{}^n = d_x{}^n = \text{etc.}
\]
in gleich berechtigter Weise repräsentirt.

Diese Symbolik, von der wir nunmehr in den meisten Untersuchungen Gebrauch machen werden, wurde zuerst von \gesperrt{Aronhold} (Crelle's Journ. Bd. 39 und 55) eingeführt, und von \gesperrt{Clebsch} (Crelle's Journ. Bd. 59) systematisch ausgebildet. Etwas früher als diese Beiden hatte \gesperrt{Cayley} eine andere symbolische Methode entwickelt, die im 30. Band von Crelle's Journ. 1846 niedergelegt ist. Beide Methoden erweisen sich in gleicher Weise fruchtbringend für die Untersuchungen der Invariantentheorie.

\textit{Beispiel.} Die quadratische binäre Form
\[
f(x) = \overline{a_0}x_1{}^2 + 2\overline{a_1}x_1x_2 + \overline{a_2}x_2{}^2
\]
ist symbolisch dargestellt durch
\[
f = a_1{}^2x_1{}^2 + 2a_1a_2x_1x_2 + a_2{}^2x_2{}^2 = a_x{}^2 = b_x{}^2 = c_x{}^2 = \text{etc.}
\]
Ihre Discriminante ist bekanntlich (vgl. auch Bd. I Nr. 182)
\[
\Delta_f = \overline{a_0}\,\overline{a_2} - \overline{a_1}{}^2\,.
\]


% ===== pdf idx 17, printed 3.0 =====
% !PRINTED 3
\weberpage{3}
Dieses Coefficientenaggregat geht unter Einführung der symbolischen Coefficienten über in
\[
\Delta_f = a_1{}^2b_2{}^2 - a_1a_2b_1b_2\,,
\]
oder auch in
\[
\Delta_f = b_1{}^2a_2{}^2 - b_1b_2a_1a_2\,.
\]
Durch Addition dieser beiden Ausdrücke erhält man
\[
2\Delta_f = (a_1b_2 - b_1a_2)^2 = (ab)^2 \quad \text{(vgl. auch Bd. I Nr. 18)}
\]
als einfache symbolische Darstellung der Discriminante von $f$.

3. \textit{Lineare Transformation.} Wir wollen nun zunächst untersuchen, in welcher Weise eine lineare Transformation den symbolischen Ausdruck einer binären Form ändert, und zu dem Zwecke uns auf ein einfaches Beispiel beschränken. Die quadratische Form
\weq{1}{f(x) = a_x{}^2 = \overline{a_0}x_1{}^2 + 2\overline{a_1}x_1x_2 + \overline{a_2}x_2{}^2}
unterwerfen wir der linearen Transformation:
\weq{I}{\left.\begin{aligned}
x_1 &= \xi_1y_1 + \eta_1y_2 \\
x_2 &= \xi_2y_1 + \eta_2y_2\,.
\end{aligned}\right\}}
Die Determinante
\[
\Delta = \begin{vmatrix} \xi_1 & \eta_1 \\ \xi_2 & \eta_2 \end{vmatrix} = (\xi\eta)
\]
der Substitutionscoefficienten nennen wir den \glqq Modul\grqq{} der Transformation. Das Resultat der Substitution (I) in (1) liefert:
\[
\begin{aligned}
\varphi(y) &= \overline{a_0}(\xi_1y_1+\eta_1y_2)^2 + 2a_1(\xi_1y_1+\eta_1y_2)(\xi_2y_1+\eta_2y_2) % [sic] a_1 printed here without overline, unlike the neighbouring terms
\\
&\quad + \overline{a_2}(\xi_2y_1+\eta_2y_2)^2 \\
&= y_1{}^2(\overline{a_0}\xi_1{}^2 + 2\overline{a_1}\xi_1\xi_2 + \overline{a_2}\xi_2{}^2) \\
&\quad + 2y_1y_2(\overline{a_0}\xi_1\eta_1 + \overline{a_1}(\xi_2\eta_1+\xi_1\eta_2) + \overline{a_2}\xi_2\eta_2) \\
&\quad + y_2{}^2(\overline{a_0}\eta_1{}^2 + 2\overline{a_1}\eta_1\eta_2 + \overline{a_2}\eta_2{}^2)\,.
\end{aligned}
\]
Setzt man hierin:
\weq{II}{\left.\begin{aligned}
\overline{a_0}\xi_1{}^2 + 2\overline{a_1}\xi_1\xi_2 + \overline{a_2}\xi_2{}^2 &= \overline{A_0} \\
\overline{a_0}\xi_1y_1 + \overline{a_1}(\xi_2\eta_1+\xi_1\eta_2) + \overline{a_2}\xi_2\eta_2 &= \overline{A_1} % ?? uncertain: printed as xi_1 y_1 (not eta_1); reproduced as read, possibly a printed slip for xi_1 eta_1
\\
\overline{a_0}\eta_1{}^2 + 2\overline{a_1}\eta_1\eta_2 + \overline{a_2}\eta_2{}^2 &= \overline{A_2}
\end{aligned}\right\}}
so wird durch die Transformation (I) aus $f(x)$
\weq{2}{\varphi(y) = \overline{A_0}y_1{}^2 + 2\overline{A_1}y_1y_2 + \overline{A_2}y_2{}^2\,.}
Führt man nun einestheils in $f(x)$ für die unsymbolischen Coefficienten $\overline{a_0}$, $\overline{a_1}$, $\overline{a_2}$ die Symbole $a_1{}^2$, $a_1a_2$, $a_2{}^2$ und ebenso für die Coefficienten $\overline{A_0}$, $\overline{A_1}$, $\overline{A_2}$ der transformirten Form $\varphi(y)$ die Symbole $A_1{}^2$, $A_1A_2$, $A_2{}^2$ ein, so erhalten die Relationen (II) die Form:
%CONT>>


% ===== pdf idx 18, printed 4.0 =====
% !PRINTED 4
\weberpage{4}
%<<CONT
\[
\begin{aligned}
(a_1\xi_1+a_2\xi_2)^2 &= a_\xi{}^2 = A_1{}^2 = \overline{A_0} \\
(a_1\xi_1+a_2\xi_2)(a_1\eta_1+a_2\eta_2) &= a_\xi a_\eta = A_1A_2 = \overline{A_1} \\
(a_1\eta_1+a_2\eta_2)^2 &= a_\eta{}^2 = A_2{}^2 = \overline{A_2}\,.
\end{aligned}
\]
In derselben Weise ergeben sich für die Coefficienten $\overline{A_k}$ der transformirten $\varphi = A_y{}^n$ einer Form $f = a_x{}^n$ die Relationen:
\[
\overline{A_k} = A_1{}^{n-k}A_2{}^k = (a_1\xi_1+a_2\xi_2)^{n-k}(a_1\eta_1+a_2\eta_2)^k = a_\xi{}^{n-k}a_\eta{}^k\,.
\]
Man erkennt daraus folgendes Gesetz:

\glqq Bezeichnet man die ursprüngliche Form mit $f = a_x{}^n$, die transformirte mit $\varphi = A_y{}^n$, so gehen die symbolischen Coefficienten $A_1$, $A_2$ aus den symbolischen $a_1$, $a_2$ durch die Transformation:
\weq{III}{\left.\begin{aligned}
A_1 &= \xi_1a_1+\xi_2a_2 = a_\xi \\
A_2 &= \eta_1a_1+\eta_2a_2 = a_\eta
\end{aligned}\right\}}
hervor, eine Transformation, die den gleichen Modul $\Delta$ wie (I) besitzt, und sich von der ursprünglichen (I) nur dadurch unterscheidet, dass in den Coefficienten derselben die Horizontal- mit den Verticalreihen vertauscht sind.\grqq{}

4. \textit{Cogredienz und Contragredienz.} Die Grössen $a_i$ einestheils und die Grössen $x_i$ anderntheils treten durch die Transformation in eine gewisse Beziehung, die man mit Contragredienz bezeichnet hat. Während nämlich die Grössen $x_i$ durch die Substitution (I) in die neuen Variabeln $y_i$ übergehen, transformiren sich die Grössen $a_i$ durch die Substitution (III) in die neuen Grössen $A_i$. Man nennt die Substitutionen (I) und (III) \gesperrt{contragrediente Substitutionen}, und ebendeshalb die Grössen $a_i$ und $x_i$ \gesperrt{contragrediente Grössen}.

Grössen, die denselben Substitutionen unterworfen sind, nennt man \gesperrt{cogrediente Grössen}. So werden wir im nächsten § Formen mit zwei Reihen Variabler $x$ und $y$ kennen lernen, welche wir als cogredient bezeichnen müssen.

Bei den binären Formen sind übrigens cogrediente und contragrediente Grössen nicht wesentlich verschieden. Denn lösen wir Substitutionen (III) nach $a_1$ und $a_2$ auf, so finden wir
\[
\begin{aligned}
(-a_2)\cdot\Delta &= \xi_1(-A_2) + \eta_1A_1 \\
a_1\cdot\Delta &= \xi_2(-A_2) + \eta_2A_1\,.
\end{aligned}
\]
Für $-a_2$ und $a_1$ ist also die Transformation genau die nämliche, wie für $x_1$ und $x_2$; oder mit anderen Worten: $x_1$ und $x_2$ sind mit $-a_2$ und $a_1$ cogredient. — Wir werden später von dieser Eigenschaft öfter Gebrauch machen, wenn wir in symbolischen Producten gemäss dieser Cogredienz $x_1$ durch $-a_2$, und $x_2$ durch $a_1$ ersetzen.


% ===== pdf idx 19, printed 5.0 =====
% !PRINTED 5
\weberpage{5}
5. \textit{Definition der Invarianten.} Betrachtet man nun zwei Reihen
symbolischer Co\"efficienten $A$ und $B$, nämlich
\[
\begin{array}{ll}
A_1 = a_1\xi_1 + a_2\xi_2 & B_1 = b_1\xi_1+b_2\xi_2\\
A_2 = a_1\eta_1+a_2\eta_2, & B_2 = b_1\eta_1+b_2\eta_2
\end{array}
\]
und bildet aus ihnen die Determinante
\[
\begin{vmatrix} A_1A_2\\ B_1B_2 \end{vmatrix} = (AB),
\]
so ist nach dem Multiplicationsgesetze für Determinanten (Bd.\,I Nr.\,70)
\weq{1}{(AB) = \Delta\,.\,(ab),}
wobei $\Delta$ der Modul der Transformation (I) Nr.\,3 ist.

Liegt daher ein beliebiges Product solcher Determinanten $(AB)$
vor, etwa $(AB)^\mu(AC)^\nu(BC)^\varrho$, so besteht immer die Beziehung
\weq{2}{(AB)^\mu(AC)^\nu(BC)^\varrho = \Delta^{\mu+\nu+\varrho}\cdot(ab)^\mu(ac)^\nu(bc)^\varrho\,.}
Dieselbe lehrt:

\glqq Irgend ein Product von Determinanten $(ab)\ldots$ aus symbolischen
Co\"efficienten $a, b, c\ldots$ ändert sich nicht [bis auf den
Factor $\Delta^\lambda$ des Substitutionsmoduls], sobald die Form
\[
f = a_x^n = b_x^n = c_x^n = \text{etc.}
\]
linear transformirt wird. Denn das von den Coefficienten der
transformirten Form
\[
\varphi = A_y^n = B_y^n = C_y^n = \text{etc.}
\]
gebildete gleiche Product unterscheidet sich seinem Werthe nach
von dem ersteren nur um eine bekannte Potenz des Moduls, wie die
Beziehung (1) lehrt.\grqq{}

Man sieht, solche symbolische Determinantenproducte sind linearen
Transformationen gegenüber unveränderliche algebraische Ausdrücke.
Nun stellen dieselben aber stets Aggregate der wirklichen
Co\"efficienten von $f$ dar, wenn nur jedes Symbol $a, b, c$ in
einer Potenz auftritt, welche gleich ist dem Grade der Form $f$
in $x$. Es existiren also für jede Form $f$ gewisse weitere
algebraische Formen ihrer Co\"efficienten, die bei linearer
Transformation unverändert bleiben.

Wir nennen dieselben \glqq Invarianten\grqq{} von $f$ und jedes symbolische
Product
\weq{3}{i = (ab)^\mu(ac)^\nu(bc)^\varrho\ldots}
stellt eine solche Invariante dar, vorausgesetzt, dass die Summe
aller Exponenten von $a$ sowohl, als von $b, c$ etc. gleich $n$
ist, wenn $n$ den Grad von $f$ in $x$ angiebt.


% ===== pdf idx 20, printed 6.0 =====
% !PRINTED 6
\weberpage{6}
Die Gleichung:
\weq{4}{(AB)^\mu(AC)^\nu(BC)^\varrho\ldots = \Delta^{\mu+\nu+\varrho}\cdot\cdot(ab)^\mu(ac)^\nu(bc)^\varrho\ldots}
ist also Definitionsgleichung für Invarianten. Sie sagt:

\glqq Eine Invariante ist eine solche Verbindung der Co\"efficienten von
$f$, welche die Eigenschaft hat, bis auf einen bekannten Factor
$\Delta^\lambda$ in ihrem Werthe unverändert zu bleiben, sei es,
dass man sie aus den Co\"efficienten der ursprünglichen, sei es,
dass man sie aus denen der transformirten Form bildet.\grqq{}

\textit{Anmerkung.} Wir haben in Nr.\,3 gesehen, dass durch lineare
Transformation das Symbol $a_1$ übergeht in $a_\xi$ und das
Symbol $a_2$ in $a_\eta$. Ersetzen wir also in dem
nichtsymbolischen Ausdruck der Invariante $i$ den Co\"efficienten
$\overline{a_0}=a_1{}^n$ durch $a_\xi{}^n$,
$\overline{a_1}=a_1{}^{n-1}a_2$ durch $a_\xi{}^{n-1}a_\eta$ etc.,
so erhalten wir die entsprechende Invariante $I$ der
transformirten Form. Da nun aber nach Definitionsgleichung (4)
die Beziehung
\[
I = \Delta^\lambda\,.\,i
\]
bestehen muss, so folgt daraus durch dieselben Schlüsse, die wir
schon in Bd.\,I Nr.\,141 bei Untersuchung der Resultante zweier
Formen gezogen haben, dass jede Invariante \gesperrt{homogen und isobar
vom Gewichte $\lambda$ ist.}

6. \textit{Definition der Covarianten.} Das symbolische Product $i$
(Nr.\,5, (3)) behält aber diese Invarianteneigenschaft noch bei,
wenn wir es mit symbolischen Factoren $a_x, b_x, c_x\ldots$
multipliciren. Denn ein solcher linearer Factor geht durch die
lineare Transformation nur wieder in einen linearen Factor über.
Multiplicirt man nämlich die erste der Relationen (III):
\[
A_1 = \xi_1 a_1 + \xi_2 a_2
\]
\[
A_2 = \eta_1 a_1 + \eta_2 a_2
\]
mit $y_1$, die zweite mit $y_2$ und addirt, so erhält man:
\[
A_1y_1 + A_2y_2 = a_1(\xi_1y_1+\eta_1y_2) + a_2(\xi_2y_1+\eta_2y_2)
\]
oder
\[
A_y = a_x\,.
\]

Wir nennen ein solches symbolisches Product, das ausser den
Determinanten $(ab), (bc), \ldots$ etc.\ldots, — die wir in
Zukunft \gesperrt{Factoren zweiter Art} oder \gesperrt{Klammerfactoren} nennen
wollen, — auch noch Factoren $a_x, b_x, c_x, \ldots$ etc., d.\,i.
\gesperrt{Factoren erster Art} enthält, eine \gesperrt{Covariante.}

Die Definitionsgleichung für Covarianten ist sonach:
\weq{5}{(AB)^\mu(AC)^\nu(BC)^\varrho\ldots A_y^\sigma B_y^\tau C_y^\varkappa\ldots = \Delta^{\mu+\nu+\varrho\cdot\cdot}\cdot(ab)^\mu(ac)^\nu(bc)^\varrho\ldots a_x^\sigma b_x^\tau c_x^\varkappa\ldots}


% ===== pdf idx 21, printed 7.0 =====
% !PRINTED 7
\weberpage{7}
Sie sagt: \glqq Eine Covariante ist eine solche Verbindung der
Co\"efficienten und Variabeln von $f$, welche die Eigenschaft hat,
bei linearer Transformation von $f$ sich nur um einen Factor
$\Delta^\lambda$ zu ändern.\grqq{}

Auch hier muss — soll das symbolische Product eine wirkliche
Covariante der Form $n^{\text{ten}}$ Grades $f=a_x^n$ sein —
jedes Symbol in einer Potenz auftreten, die gleich $n$ ist,
d.\,h. es muss sein:
\weq{6}{\left.\begin{aligned}
\mu+\nu+\sigma\ldots &= n\\
\mu+\varrho+\tau\ldots &= n\\
\nu+\varrho+\varkappa\ldots &= n
\end{aligned}\right\}}

7. \textit{Beispiele von simultanen Invarianten.} Ein symbolisches
Product wie (4) oder (5) kann auch eine Co\"efficientenverbindung
verschiedener Formen
\[
f = a_x^n, \quad \varphi = b_x^m, \quad \psi = c_x^r\ldots
\]
darstellen, und wir nennen es alsdann \glqq simultane Invariante,
bezw. Covariante\grqq{} dieser Formen. In diesem Falle muss
natürlich sein:
\weq{7}{\left.\begin{aligned}
\mu+\nu+\sigma\ldots &= n\\
\mu+\varrho+\tau\ldots &= m\\
\nu+\varrho+\varkappa\ldots &= r
\end{aligned}\right\}}
Eine solche simultane Invariante ist z.\,B. die Resultante zweier
Formen. (Vgl. Bd.\,I Nr.\,147.) So kann auch der Klammerfactor
$(a\alpha)$ als simultane Invariante der beiden linearen Formen
\[
f = a_1x_1 + a_2x_2
\]
\[
\varphi = \alpha_1x_1 + \alpha_2x_2
\]
aufgefasst werden.

Ebenso ist die Determinante der Co\"efficienten dreier
quadratischer Formen
\[
f = \overline{a_0}x_1{}^2+2\overline{a_1}x_1x_2+\overline{a_2}x_2{}^2
\]
\[
\varphi = \overline{b_0}x_1{}^2+2\overline{b_1}x_1x_2+\overline{b_2}x_2{}^2
\]
\[
\psi = \overline{c_0}x_1{}^2+2\overline{c_1}x_1x_2+\overline{c_2}x_2{}^2
\]
eine simultane Invariante. Denn führt man in diese Determinante
\[
R = \begin{vmatrix}
\overline{a_0}, & \overline{a_1}, & \overline{a_2}\\
\overline{b_0}, & \overline{b_1}, & \overline{b_2}\\
\overline{c_0}, & \overline{c_1}, & \overline{c_2}
\end{vmatrix}
\]
die symbolischen Co\"efficienten ein, so wird sie
\[
R = \begin{vmatrix}
a_1{}^2, & a_1a_2, & a_2{}^2\\
b_1{}^2, & b_1b_2, & b_2{}^2\\
c_1{}^2, & c_1c_2, & c_2{}^2
\end{vmatrix}\,.
\]


% ===== pdf idx 22, printed 8.0 =====
% !PRINTED 8
\weberpage{8}
Diese Determinante hat den Werth (vgl. Bd.\,I Nr.\,40, Anmerk.)
\[
R = (ab)(bc)(ac)\,.
\]
Die Function $R$ lässt sich also durch ein symbolisches Product
darstellen und ist folglich Invariante.

\textit{Beispiele von In- und Covarianten einzelner Formen.} Die
cubische Form
\[
f(x) = a_x{}^3 = b_x{}^3 = c_x{}^3 = d_x{}^3 = \ldots =
\overline{a_0}x_1{}^3+3\overline{a_1}x_1{}^2x_2+3\overline{a_2}x_1x_2{}^2+\overline{a_3}x_2{}^3
\]
hat eine Covariante $(ab)^2a_xb_x$. Ihre nichtsymbolische Form
erhält man, indem man die symbolischen Co\"efficienten in dem
Producte
\[
\Pi = (a_1b_2-b_1a_2)^2(a_1x_1+a_2x_2)(b_1x_1+b_2x_2)
\]
\[
= (a_1{}^2b_2{}^2-2a_1a_2b_1b_2+b_1{}^2a_2{}^2)(a_1b_1x_1{}^2+(a_1b_2+b_1a_2)x_1x_2+a_2b_2x_2) % [sic]
\]
\[
= (a_1{}^3b_1b_2{}^2-2a_1{}^2a_2b_1{}^2b_2+a_1a_2{}^2b_1{}^3)x_1{}^2
\]
\[
\quad + (a_1{}^3b_2{}^3-2a_1{}^2a^2b_1b_2{}^2+a_1a_2{}^2b_1{}^2b_2+\cdots)x_1x_2 % [sic]
\]
\[
\quad + (a_2a_1{}^2b_2{}^3-2a_1a_2{}^2b_1b_2{}^2+a_1{}^2a_2b_1{}^2b_2)x_2{}^2
\]
durch die nichtsymbolischen Co\"efficienten ersetzt.

Man erhält:
\[
\tfrac{1}{2}\Pi = (a_0a_2-a_1{}^2)x_1{}^2+(a_0a_3-a_1a_2)x_1x_2+(a_1a_3-a_2{}^2)x_2{}^2\,.
\]
Für die biquadratische Form
\[
f = a_x{}^4 = b_x{}^4 = \text{etc.} =
a_0x^4+4a_1x_1{}^3x_2+6a_2x_1{}^2x_2{}^2+4a_3x_1x_2{}^3+a_4 % [sic]
\]
existirt eine Invariante $(ab)^4$. Ihre nichtsymbolische Gestalt
geht aus der binomischen Entwicklung
\[
\Pi_1 = (a_1b_2-b_1a_2)^4 = a_1{}^4b_2{}^4-4a_1{}^3a_2b_1b_2{}^3
\]
\[
\quad + 6a_1{}^2a_2{}^2b_1{}^2b_2{}^2-4a_1a_2{}^3b_1{}^3b_2+a_2b_1{}^4
\]
hervor, wenn man wieder die symbolischen Co\"efficienten durch
die wirklichen ersetzt. Man erhält:
\[
\Pi = a_0a_4-4a_1a_3+6a_2{}^2-4a_3a_1+a_4a_0,
\]
oder
\[
\Pi = (ab)^4 = 2(a_0a_4-4a_1a_3+3a_2{}^2)\,.
\]

8. \textit{Eigenschaften der In- und Covarianten, die aus dem
symbolischen Producte direct erkennbar sind.} Aus den bisherigen
Entwicklungen ist klar, dass alle symbolischen Producte, wenn sie
nur die Symbole in der geeigneten Anzahl enthalten, Invarianten
oder Covarianten von binären Formen repräsentiren. Daraus geht
hervor, dass diese In- und Covarianten, soweit sie sich durch
solche Producte darstellen lassen, \gesperrt{ganze und rationale Functionen
der Co\"efficienten der ursprünglichen Formen sind.} Sie sind
überdies \gesperrt{homogen in den Variabeln}
%CONT>>


% ===== pdf idx 23, printed 9.0 =====
% !PRINTED 9
\weberpage{9}
wie in den Co\"efficienten; den Grad in den Variabeln erkennt man aus
der Anzahl der Factoren erster Art, den Grad in den Co\"efficienten
aus der Anzahl der verschiedenen Symbole, welche in dem Producte
auftreten. Da die Factoren erster Art bei Transformation in sich
selbst übergehen, so ist der Exponent des Moduls $\Delta$, den die
Covariante nach der Transformation als Factor erhält, gleich der Zahl
der Klammerfactoren überhaupt.

Nun lehrt die Definitionsgleichung für Invarianten $I = \Delta^\varrho \cdot i$
auch, dass die Invariante isobar ist vom Gewichte $\varrho$, wie wir in
Nr. 5 gesehen haben. Da aber der Modulexponent gleich der Zahl der
Klammerfactoren ist, so giebt sonach diese Zahl stets auch das Gewicht
der Invariante an, während das Gewicht der Covariante sich
zusammensetzt aus der Zahl $\varrho$ der Klammerfactoren plus dem Grade
$m$ der Covariante in den Variabeln $x$. Ist $\varrho$ eine gerade Zahl, so
nennen wir die betreffende Form auch eine \gesperrt{Form geraden
Charakters} (\textit{forme droite}); so die in Nr. 7 berechnete
Invariante $(ab)^4$ und die Covariante $(ab)^2 a_x b_x$. Ist dagegen
$\varrho$ ungerade, so heissen wir die Form eine \gesperrt{Form ungeraden
Charakters} (\textit{forme gauche}) oder auch schiefe Invariante,
resp. Covariante; so die in Nr. 7 berechnete simultane Invariante
$(ab)(bc)(ac)$. Wir werden später sehen: Die Quadrate von schiefen
Formen lassen sich durch Formen geraden Charakters darstellen.

Ein symbolisches Product $\Pi$ kann indessen auch, wenn die
symbolischen Co\"efficienten durch die wirklichen ersetzt werden, auf
einen Ausdruck führen, der identisch verschwindet. Dies tritt unter
anderm immer dann ein, sobald es bei Vertauschung zweier
gleichwerthiger Symbole nur sein Zeichen ändert. Denn durch
Vertauschung zweier Symbole, welche dieselbe Form darstellen, kann
sich die wirkliche Covariante $C$ nicht ändern, da es ja gleichgiltig
ist, ob ich beispielsweise $\overline{a_3}\cdot\overline{a_5}$ durch
$a_1^{n-3}a_2^3 \cdot b_1^{n-5}b_2^5$ oder durch
$b_1^{n-3}b_2^3 \cdot a_1^{n-5}a_2^5$ repräsentire. Nun ist also vor
der Vertauschung
\[
C = \Pi,
\]
nach derselben
\[
C = -\Pi,
\]
demnach:
\[
2C = 0, \quad \text{also} \quad C = 0.
\]

So ist z.\,B. für die quadratische Form $f = a_x^2 = b_x^2$ die
Covariante $(ab) a_x b_x$ identisch Null; so verschwindet für die
cubische Form $\varphi = a_x^3 = b_x^3$ die Invariante $(ab)^3$. Der
Umstand, dass die Co- und Invarianten ganze rationale und homogene
Functionen in den Co\"efficienten einer Form $f = a_x^n$ sind, lehrt
uns auch, dass sie symmetrische Functionen der Wurzeln $a_i$ von
$f = 0$ sein müssen. Bezeichnen wir mit $\alpha_{ix}$ die
%CONT>>


% ===== pdf idx 24, printed 10.0 =====
% !PRINTED 10
\weberpage{10}
%<<CONT
$n$ linearen Factoren von $f = 0$, so können wir auch sagen: Jede
Covariante oder Invariante von $f$ ist eine \gesperrt{simultane
Covariante}*) oder Invariante der $n$ linearen Formen $\alpha_{ix}$,
welche symmetrisch ist in den Symbolen $\alpha_i$, und umgekehrt:
Jede simultane und in $\alpha_i$ symmetrische Covariante der $n$
linearen Formen $\alpha_{ix}$ ist eine Covariante von
\[
f = \alpha_{1x} \,.\, \alpha_{2x} \,.\, \alpha_{3x} \dots \alpha_{nx} = a_x^n .
\]

Es tritt nun die Frage auf: Wie ist eine in symbolischen
Co\"efficienten von $f = a_x^n = b_x^n = $ etc. gegebene Covariante
darstellbar durch die Wurzeln von $f = 0$, und umgekehrt: wie kann man
eine simultane symmetrische Covariante der $\alpha_{ix}$ zurückführen
auf ein symbolisches Product in den Co\"efficienten von $f$. Wenn es
nun auch nicht unser Zweck ist mit Wurzeln zu rechnen, so werden wir
doch an geeigneter Stelle (vergl. Nr. 39 und Nr. 94) wenigstens eine
allgemeine Methode zur Beantwortung beider Fragen geben.

Den Beweis, dass jede Covariante oder Invariante sich überhaupt durch
ein symbolisches Product, oder durch eine Summe solcher, darstellen
lässt, werden wir später kennen lernen. (\S\,9). Ich werde indess bis
dahin von diesem Satze auch keinen Gebrauch machen.

9. \textit{Definition des Faltungsprocesses.} Aus einem symbolischen
Producte
\[
P = (ab)^\mu (ac)^\nu (bc)^\varrho \, \dots\dots \, a_x^\sigma b_x^\tau c_x^\varkappa
\]
lassen sich eine Reihe anderer ableiten durch eine sehr einfache
Operation. Man greift zwei Factoren erster Art mit verschiedenen
Symbolen heraus, etwa $a_x$ und $b_x$ und setzt an ihre Stelle den
Klammerfactor $(ab)$. Dieses Verfahren nennt \gesperrt{Gordan}
\gesperrt{falten}, und den Process selbst \glqq Faltungsprocess\grqq{}.
Man kann auch umgekehrt das Product \gesperrt{entfalten}, indem man
einen Klammerfactor $(ab)$ in $P$ durch $a_x b_x$ ersetzt. Der
Faltungsprocess ist ein rein symbolischer Process. Es kann ihm keine
Bedeutung unterlegt werden, die durch eine wirklich rechnerische
Operation ihren Ausdruck findet. Die Faltung kann gleichzeitig mehr
als einmal vorgenommen werden, nämlich gerade so oft, als sich
verschiedene Factorenpaare erster Art im Producte gleichzeitig
gruppiren lassen.

Ich will hier durch Faltung einige In- und Covarianten der cubischen
Form herleiten. Es sei
\[
f = a_x^3 = b_x^3 = c_x^3 = d_x^3 = \text{etc.}
\]

\medskip
\noindent *) Vgl. auch Nr.\,39 und Nr.\,94.


% ===== pdf idx 25, printed 11.0 =====
% !PRINTED 11
\weberpage{11}
1) Durch zweimalige Faltung des Productes
\[
\Pi = a_x^3 \,.\, b_x^3
\]
erhält man die Covariante
\[
\Delta = (ab)^2 a_x b_x \,.
\]
Die Formen $(ab) a_x^2 b_x^2$, $(ab)^3$, welche durch einmalige bezw.
dreimalige Faltung daraus erhalten werden, sind identisch null nach
den Bemerkungen am Schlusse von Nr. 8.

2) Aus dem Producte
\[
\Delta^2 = (ab)^2 a_x b_x \,.\, (cd)^2 c_x d_x
\]
erhält man durch zweimalige Faltung entweder:
\[
R = (ab)^2 (cd)^2 (ac) (bd)
\]
oder:
\[
R_1 = (ab)^2 (cd)^2 (bc) (ad) .
\]
Beide Producte stellen offenbar dieselbe Invariante dar; denn das
zweite geht aus dem ersten durch die berechtigte Vertauschung von $b$
mit $a$ hervor.

3) Aus dem Producte
\[
\Delta \,.\, f = (ab)^2 a_x b_x \,.\, c_x^3
\]
erhält man durch einmalige Faltung entweder
\[
Q = (ab)^2 (ac) b_x c_x^2
\]
oder:
\[
Q_1 = (ab)^2 (bc) a_x c_x^2 ,
\]
und auch hier ist wieder $Q_1 = Q$, wie man erkennt, wenn $a$ mit $b$
vertauscht wird.

10. \textit{Zusammenhang der durch Faltung entstehenden Producte.}
Für alle späteren Untersuchungen, insbesondere für die
Systembildung, ist der Faltungsprocess von einschneidender Bedeutung
geworden. (Vergl. auch \gesperrt{Camille Jordan}, Liouville Journ.
sér. 3, Bd. 2 und 5.) Alle durch ihn aus einem Producte $A$ — wir
wollen es \gesperrt{Stammproduct} nennen — hervorgehenden neuen
Producte haben symbolisch wenigstens einen engen Zusammenhang. Sie
haben von vornherein einen gemeinsamen symbolischen Factor, nämlich
die Klammerfactoren, welche bereits das Stammproduct besass und auf
welche ja der Faltungsprocess keinen Einfluss hat. Diese
Klammerfactoren werden unter Umständen von Wichtigkeit, insofern sie
eine ganze Gruppe von symbolischen Producten charakterisiren können,
die aus gewissen Betrachtungen ausgeschlossen werden dürfen. (Vergl.
auch \glqq Reducenten\grqq{} Nr. 123.)


% ===== pdf idx 26, printed 12.0 =====
% !PRINTED 12
\weberpage{12}
Eine bestimmte Zahl von Faltungen kann natürlich, wie schon an den
letzten Beispielen Nr. 9 ersichtlich war, an einem Producte auf
verschiedene Arten vorgenommen werden. Sind insbesondere $A_1$ und
$A_2$ zwei solche durch $\mu$ Faltungen aus $A$ entstandene Producte,
dass $A_1$ durch einmalige in zwei Symbolfactoren vorgenommene
Vertauschung gleichberechtigter Symbole $a$ und $b$, oder $a$ und $c$,
$b$ und $c$ etc. in $A_2$ übergeht, so heissen $A_1$ und $A_2$
\gesperrt{benachbarte symbolische Producte}. So sind:
\[
\begin{aligned}
A_1 &= (ac)\, b_x \,.\, Q \\
A_2 &= (ab)\, c_x \,.\, Q
\end{aligned}
\]
einestheils, und
\[
\begin{aligned}
A_1 &= (ad)\, (bc) \,.\, Q \\
A_2 &= (cd)\, (ba) \,.\, Q
\end{aligned}
\]
anderntheils \gesperrt{Nachbarproducte}, wobei mit $Q$ die den
Producten $A_1$ und $A_2$ gemeinschaftlichen Factoren bezeichnet sein
mögen. Es ist klar, dass alle durch $\mu$-malige Faltung aus $A$
entstehenden symbolischen Producte sich stets so ordnen lassen, dass
immer das vorausgehende mit dem nachfolgenden benachbart ist. Hat man
eine Reihe symbolischer Producte
\[
A_1, \ A_2, \ A_3 \dots A_m ,
\]
welche alle homogen in den Co\"efficienten der nämlichen Formen $f$,
$\varphi$, $\psi \dots$, also auch deren Symbolen $a$, $b$, $c \dots$,
$\alpha$, $\beta$, $\gamma \dots$ und homogen in den Variabeln $x_1$,
$x_2$ sind, so kann man sie stets aus ein und demselben Stammproduct
durch $\mu$-malige Faltung entstanden denken, und demnach ist es immer
möglich, wenn sie nicht benachbart sein sollten, zwischen je zwei
$A_i$ und $A_k$ eine Reihe von Producten
\[
M_1 M_2 \dots M_r; \ N_1, N_2 \dots N_s; \ \dots.
\]
einzuschalten, so dass $A_1$ mit $A_m$ durch eine Kette von
Nachbargliedern zusammenhängt. Dies wird bei manchen grundlegenden
Untersuchungen von Wichtigkeit und ich werde daher die ebengemachten
Bemerkungen am Schlusse dieses \S{} durch ein Beispiel erläutern.

11. \textit{Der Identitäts- und Productsatz.} Sehr häufig erwächst die
Aufgabe, symbolische Producte, die durch Faltung aus einem
Stammproduct entstehen, in andere umzuformen. Hierzu wird unter andern
Hilfsmitteln in sehr fruchtbringender Weise die Identität
\weq{I}{a_x(bc) + b_x(ca) + c_x(ab) = 0\text{*)}}

\medskip
\noindent *) Die Richtigkeit dieser Identität haben wir in Bd.\,I
Nr.\,34 gezeigt; man kann sie übrigens leicht verificiren, wenn man in
der Determinante
\[
\begin{vmatrix} a_1 & b_1 & c_1 \\ a_2 & b_2 & c_2 \\ a_x & b_x & c_x \end{vmatrix} ,
\]
% ?? uncertain: bottom-row entries blurred in scan; read as a_x, b_x, c_x
%CONT>>


% ===== pdf idx 27, printed 13.0 =====
% !PRINTED 13
\weberpage{13}
%<<CONT
verwendet, welche wir auch, indem wir $x_1$ durch $d_2$ und $x_2$ durch
$-d_1$ ersetzen, in der Form schreiben können

\weq{II}{(ad)(bc) + (bd)(ca) + (cd)(ab) = 0,}

oder auch, wenn wir in (I)% ?? uncertain: footnote anchor asterisk not visible on this crop; footnote text below is attached here as best guess
\footnote{die mit $x_1$ multiplicirt gedachte erste und die mit $x_2$ multiplicirt gedachte zweite Zeile von der dritten subtrahirt.}
für $c_1$ schreiben $y_2$, für $c_2$ dagegen $-y_1$,
in der Form

\weq{III}{a_xb_y - b_xa_y = (ab)(xy).}

Mit Hilfe dieser dritten Form des Identitätssatzes können wir z. B.
direct zeigen, dass sich ein symbolisches Product, welches nicht identisch
verschwindet, und dessen zumeist auftretender Klammerfactor mit
gleichberechtigten Symbolen in ungerader Potenz vorhanden ist, immer
in andere verwandeln lässt, die diesen Klammerfactor in gerader Potenz
besitzen. Denn nehmen wir an, die im Producte $P$ zu höchst vorkommende
Potenz eines Klammerfactors sei $(ab)^{2\mu+1}$, dann hat $P$ die Form

\[ P = (ab)^{2\mu+1}\cdot Q \]

und $Q$ enthält im Allgemeinen die Symbole $a$ und $b$ --- mindestens
einmal, wenn $P\gtrless 0$ --- noch in beliebigen andern Verbindungen,
die wir, wenn $\xi, \eta, \zeta, \tau \ldots$ beliebige Grössen, mit
$a_\xi, b_\eta, a_\zeta, b_\tau \ldots$ etc. bezeichnen können. Es ist dann

\[ P = (ab)^{2\mu+1} a_\xi b_\eta a_\zeta b_\tau \ldots \times \overline{Q}, \]

wo nun $\overline{Q}$ weder ein Symbol $a$ noch ein Symbol $b$ enthält.
Für jedes Factorenpaar $a_\xi b_\eta, a_\zeta b_\tau \ldots$ können wir
schreiben:

\[ \tfrac{1}{2}\{(a_\xi b_\eta + b_\xi a_\eta) + (a_\xi b_\eta - b_\xi a_\eta)\}, \ldots \]

oder nach (III):

\[ \tfrac{1}{2} (a_\xi b_\eta + b_\xi a_\eta) + \tfrac{1}{2} (ab)(\xi\eta), \text{ etc.,} \]

so dass, wenn $\varrho$ solche Factorenpaare vorhanden sind, $P$ übergeht in

\[
P = \frac{(ab)^{2\mu+1}}{2^\varrho}\{a_\xi b_\eta + b_\xi a_\eta\}\cdot
\{a_\zeta b_\tau + b_\zeta a_\tau\}\ldots \times \overline{Q}
+ \frac{(ab)^{2\mu+1}}{2^\varrho}\sum_i \overline{Q}_i.
\]

Der erste Term rechts ist Null, da er durch Vertauschung von $a$ mit
$b$ nur sein Zeichen ändert. Jedes Glied $\overline{Q}_i$ in der Summe des zweiten
%CONT>>


% ===== pdf idx 28, printed 14.0 =====
% !PRINTED 14
\weberpage{14}
%<<CONT
Termes enthält aber mindestens \gesperrt{Einen} Factor $(ab)$, und es wird
sonach in der That:

\[ P = \frac{1}{2^\varrho}\cdot (ab)^{2\mu+2}\sum \overline{\overline{Q}}_i. \]

Wir werden jede der drei obigen Relationen (I), (II), (III) stets
kurz als \glqq Identitätssatz\grqq{} bezeichnen. Aus der ersten von ihnen
leiten wir eine weitere ebenfalls häufig zu benutzende Identität ab,
indem wir $c_x (ab)$ auf die rechte Seite schaffen und dann quadriren.
Es kommt nach einer einfachen Umstellung:

\weq{IV}{a_xb_x(ac)(bc) = \tfrac{1}{2}\{(bc)^2a_x^2 + (ac)^2b_x^2 - (ab)^2c_x^2\}}

oder für $x_1 = -d_2$, $x_2 = d_1$

\weq{V}{(ad)(bd)(ac)(bc) = \tfrac{1}{2}\{(bc)^2(ad)^2 + (ac)^2(bd)^2 - (ab)^2(cd)^2\}.}

Diese beiden Identitäten belegen wir mit dem Namen \glqq Productsatz\grqq{}.

12. \textit{Ein Beispiel symbolischer Rechnung.} Ich will die allgemeinen
Betrachtungen dieses § über Symbolik schliessen, indem ich an einem
Beispiel einen Theil der bisher entwickelten Begriffe zusammenfassend
erläutere. Wir stellen uns nämlich die Aufgabe, irgend eine Relation
zwischen symbolischen Producten $A_i$, also eine ganze rationale Function

\[ F(A_i) = 0, \]

welche homogen in den Coëfficienten und Variabeln ist, und deren
einzelne Glieder $G_i$ weder identisch verschwinden noch sich gegenseitig
aufheben, so umzuformen, dass sie die Gestalt

\weq{VI}{F(A_i) = \sum_1^\lambda \{(bc)a_x + (ca)b_x + (ab)c_x\}\varphi_\lambda = 0}

annimmt, aus welcher das identische Verschwinden der Function $F(A_i)$
direct ersichtlich ist. Die Möglichkeit dieser Umformung liegt in der
Beziehung, welche vermöge des Identitätssatzes zwischen zwei Nachbargliedern
der Function $F(A_i)$ besteht\footnote{Einen strengen Beweis der Relation (VI) wollen wir später bringen. (\S.\,10.)}. Sind nämlich $G_\lambda$ und $G_\mu$
zwei Nachbarglieder

\[
\begin{aligned}
G_\lambda &{}= (ac)b_x\cdot Q &&\text{oder} &&{}= (ad)(bc)Q\\
G_\mu &{}= (ab)c_x\cdot Q &&\text{oder} &&{}= (cd)(ba)Q,
\end{aligned}
\]

so ist ihre Differenz unter Anwendung des Identitätssatzes durch

\weq{1}{
\begin{aligned}
G_\mu - G_\lambda &{}= \{(ab)c_x + (bc)a_x + (ca)b_x\}Q - a_x(bc)\cdot Q\\
\text{resp.}\quad G_\mu - G_\lambda &{}= \{(ad)(bc) + (bd)(ca) + (cd)(ab)\}Q - (ca)(bd)Q
\end{aligned}
}
%CONT>>


% ===== pdf idx 29, printed 15.0 =====
% !PRINTED 15
\weberpage{15}
%<<CONT
darstellbar, wobei $Q$ immer die gemeinschaftlichen Symbolfactoren von
$G_\mu$ und $G_\lambda$ bedeutet.

Sind aber $G_\varkappa$ und $G_\varrho$ keine benachbarten Glieder, so können wir
gleichwohl noch ihre Differenz in der Form darstellen

\weq{2}{
\begin{aligned}
G_\varkappa - G_\varrho &{}= \sum_1^\lambda \{(bc)a_x + b_x(ca) + (ab)c_x\}Q_\lambda + \sum_1^\lambda r_x(pq)\cdot Q_\lambda\\
&{}+ \sum_1^\lambda \{(ad)(bc)+(bd)(ca)+(cd)(ab)\}Q'_\lambda+\sum_1^\lambda (rs)(pq)\cdot Q'_\lambda
\end{aligned}
}

wobei $p, q, r, s$ für irgend welche Combination der Symbole $a, b, c, d$
stehen. Denn wegen der vorausgesetzten Homogeneität in Coëfficienten
und Variabeln, können wir zwischen $G_\varkappa$ und $G_\varrho$ eine Reihe von
Gliedern

\[ M_1, \ M_2, \ M_3 \ldots M_s \]

einschalten, so dass in der Reihe

\[ G_\varkappa \ M_1 \ M_2 \ M_3 \ldots G_\varrho \]

jedes vorausgehende mit dem nachfolgenden Gliede benachbart ist.
Dann ist aber

\[ G_\varkappa - G_\varrho = (G_\varkappa - M_1) + (M_1 - M_2) + (M_2 - M_3) + \ldots (M_s - G_\varrho) \]

und für jede Differenz rechts lässt sich nun aus Formel (1) ein Werth
einführen, wodurch wir zu Formel (2) gelangen.

In welcher Weise nun diese beiden Identitäten (1) und (2) verwendet
werden, um irgend eine Relation $F(A_i) = 0$ in die Form

\[ \sum \{(bc)a_x + (ca)b_x + (ab)c_x\}\varphi_\lambda = 0 \]

zu bringen, will ich an folgendem Beispiel zeigen. In der allgemeinen
Darstellung § 10 werde ich darauf zurückkommen.

Es seien $f = a_x^3 = b_x^3$ und $\varphi = \alpha_x^3 = \beta_x^3$ zwei cubische Formen,
dann besteht zwischen den beiden simultanen Covarianten

\[
\begin{aligned}
G_1 &{}= (\alpha\beta)^2 (ab)(a\alpha)(a\beta)b_x^2\\
G_2 &{}= (ab)^2 (\alpha\beta)(\alpha a)(\alpha b)\beta_x^2
\end{aligned}
\]

die Relation

\[ F(G_i) = G_1 + G_2 = 0. \quad \text{(Vergl. auch Nr. 312, 5.)} \]

Es handelt sich nun darum, $F(G_i)$ durch eine Summe darzustellen, deren
einzelne Summanden Differenzen von Nachbargliedern sind. Zu dem
Zwecke genügt ist hier in $G_1$: $a$ mit $b$, in $G_2$: $\alpha$ mit $\beta$ zu vertauschen.
Die Summe der so entstehenden zwei Glieder, die wir mit
%CONT>>


% ===== pdf idx 30, printed 16.0 =====
% !PRINTED 16
\weberpage{16}
%<<CONT
\[
\begin{aligned}
-G_1' &{}= (\alpha\beta)^2 (ba)(b\alpha)(b\beta)a_x^2\\
-G_2' &{}= (ab)^2 (\beta\alpha)(\beta a)(\beta b)a_x^2
\end{aligned}
\]

bezeichnen wollen, ist dann wiederum gleich $F(G_i)$ und demnach:

\weq{3}{2F(G_i) = G_1 - G_1' + G_2 - G_2'.}

Nun sind aber anderntheils die symbolischen Producte

\[
\begin{aligned}
M_1 &{}= (\alpha\beta)^2 (ab)(a\alpha)(b\beta)a_xb_x\\
M_2 &{}= (ab)^2 (\alpha\beta)(\alpha a)(\beta b)\alpha_x\beta_x
\end{aligned}
\]

zu $G_1$ und $+G_1'$, resp. $G_2$ und $+G_2'$ benachbart; wenn wir daher
Gleichung (3) in die Form

\weq{4}{2F(G_i) = (G_1 - M_1) + (M_1 - G_1') + (G_2 - M_2) + (M_2 - G_2')}

bringen, so sind die einzelnen Differenzen rechts Differenzen von
Nachbargliedern.

Unter Benutzung der Gleichung (1) geht demnach die Relation (4)
über in:

\[
\begin{aligned}
2F(G_i) &{}= \{(b\beta)a_x + (\beta a)b_x + (ab)\beta_x\}(\alpha\beta)^2(ab)(a\alpha)b_x\\
&{}\quad - (\alpha\beta)^2(ab)^2(a\alpha)b_x\beta_x\\
&{}\quad + \{(b\alpha)a_x + (\alpha a)b_x + (ab)\alpha_x\}(\alpha\beta)^2(ab)(b\beta)a_x\\
&{}\quad - (\alpha\beta)^2(ab)^2(b\beta)a_x\alpha_x\\
&{}\quad - \{(b\beta)\alpha_x + (\beta\alpha)b_x + (\alpha b)\beta_x\}(ab)^2(\alpha\beta)(\alpha a)\beta_x
% ?? uncertain: the beta_x in this line prints with a heavier/bolder glyph than elsewhere on the page; reading as ordinary beta_x
\\
&{}\quad + (\alpha\beta)^2(ab)^2(a\alpha)b_x\beta_x\\
&{}\quad - \{(a\beta)\alpha_x + (\beta\alpha)a_x + (\alpha a)\beta_x\}(ab)^2(\alpha\beta)(\beta b)\alpha_x\\
&{}\quad + (\alpha\beta)^2(ab)^2(b\beta)a_x\alpha_x.
\end{aligned}
\]

Die letzten Terme der vier Zeilen heben sich gegenseitig auf; die
ersten Terme aber sind paarweise gleich, indem der erste Term der
zweiten Zeile durch Vertauschung der gleichberechtigten Symbole $\alpha$
und $\beta$ und der erste Term der dritten aus dem ersten der vierten
Zeile durch Vertauschung von $a$ mit $b$ hervorgeht. Also ist in der
That

\[
\begin{aligned}
F(G_i) &{}= \{(b\beta)a_x + (\beta a)b_x + (ab)\beta_x\}(\alpha\beta)^2(ab)(a\alpha)b_x\\
&{}\quad - \{(b\beta)\alpha_x + (\beta\alpha)b_x + (\alpha b)\beta_x\}(ab)^2(\alpha\beta)(\alpha a)\beta_x = 0
\end{aligned}
\]

oder

\weq{I}{F(G_i) = \sum \{(rs)p_x + (sp)r_x + (pr)s_x\}\varphi_\lambda = 0.}

In ähnlicher Weise kann man jede Relation in die Form (I) bringen.
(Siehe Evectantenprocess).


% ===== pdf idx 31, printed 17.0 =====
% !PRINTED 17
\weberpage{17}

\begin{center}
\S.\,2. \textbf{Process der Polarenbildung.}
\end{center}

13. \textit{Definition der Polare einfacher Formen.} An der binären Form $f=a_x^n$ können wir folgende symbolische Operation vornehmen. Wir sondern $\mu$ Factoren $a_x$ von $a_x^n$ ab und ersetzen sie durch gleich viel Factoren $a_y$. Die so entstehende Form
\weq{1}{f_{y^\mu} = a_x^{n-\mu}a_y^\mu}
nennen wir $\mu^{\text{te}}$ Polare der Form $a_x^n$, und den eben geschilderten einfachen Process: Polarenprocess. Man erkennt, dass jede Form $n^{\text{ten}}$ Grades $n+1$ Polaren besitzt, indem wir der Zahl $\mu$ alle Werthe von 0 bis $n$ beilegen. Eine höhere Polare von $f$ als die $n^{\text{te}}$ ist stets identisch null. Die $(n-\mu)^{\text{te}}$ Polare ist:
\weq{2}{f_{y^{n-\mu}} = a_x^\mu a_y^{n-\mu}\,.}
Man sieht, dass beide Polaren (1) und (2) durch Vertauschung von $x$ mit $y$ aus einander hervorgehen. Diese Thatsache, dass alle Polaren, deren Grad $>\dfrac{n}{2}$ ist, aus den vorhergehenden durch Vertauschung von $x$ mit $y$ erhalten werden, wird später bei Berechnung von Polaren complicirter Formen von Wichtigkeit.

Ich will diesen sehr einfachen Process im Interesse späterer Betrachtungen etwas umständlicher formuliren. Um $f_{y^\mu}$ zu bilden, kann man sich nämlich die $n$ Factoren $a_x$ von $f$ momentan durch $n$ lineare Factoren
\[ r_{1x},\ r_{2x},\ r_{3x}\ldots r_{nx} \]
ersetzt denken. Das entstehende Product wird nun $\mu$ mal polarisirt, indem man in allen Combinationen je $\mu$ Factoren $r_{kx}$ durch die $\mu!$ Permutationen gleichviel entsprechender Factoren $r_{ky}$ ersetzt. Dadurch entstehen $\binom{n}{\mu}\mu! = n(n-1)(n-2)\ldots(n-\mu+1)$ Glieder, nämlich gerade so viel, als die Anzahl aller Variationen von $n$ Elementen zur $k^{\text{ten}}$ Klasse beträgt. Bildet man deren Summe, dividirt dieselbe durch diese Anzahl, und ersetzt hierauf wieder alle Symbole $r_k$ durch $a$, so entsteht gerade wieder das Product $a_x^{n-\mu}a_y^\mu$. Es ist also:
\[
\begin{aligned}
f_{y^{n-\mu}} &= a_x^{n-\mu}a_y^\mu\\
&= \frac{1}{\binom{n}{\mu}\mu!}\sum r_{i_1y}r_{i_2y}\ldots r_{i_\mu y}\cdot\frac{f}{r_{i_1x}r_{i_2x}\ldots r_{i_\mu x}}\,.
\end{aligned}
\]
% ?? uncertain: the exact placement/form of the coefficient before the sum on the second line is partly obscured by an ink mark in the scan; read here as $1/\binom{n}{\mu}\mu!$, consistent with the surrounding text ("dividirt dieselbe durch diese Anzahl").

Der Polarenprocess ist nun aber nicht, wie der Faltungsprocess rein symbolischer Natur, vielmehr lässt er sich in mehrfacher Weise
%CONT>>


% ===== pdf idx 32, printed 18.0 =====
% !PRINTED 18
\weberpage{18}

%<<CONT
durch rechnerische Operationen ersetzen, die wir im Folgenden kennen lernen wollen.

14. \textit{Differentiationsmethode.} Bilden wir von $f=a_x^n$ die beiden ersten Differentialquotienten nach $x_1$ und $x_2$, so erhalten wir:
\weq{1}{\frac{\partial f}{\partial x_1} = na_1a_x^{n-1} = nf_1}
\weq{2}{\frac{\partial f}{\partial x_2} = na_2a_x^{n-1} = nf_2\,.}
Multiplicirt man die erste der beiden Gleichungen mit $y_1$, die zweite mit $y_2$ und addirt die entstehenden Producte, so kommt:
\[ y_1f_1+y_2f_2 = (a_1y_1+a_2y_2)a_x^{n-1} = a_ya_x^{n-1} = f_y\,. \]
Die erste Polare von $f=a_x^n$ ist also auch durch die Gleichung definirt:
\weq{3}{f_y = y_1f_1+y_2f_2 = \frac{1}{n}\left(y_1\frac{\partial f}{\partial x_1}+y_2\frac{\partial f}{\partial x_2}\right).}
Bilden wir ebenso die zweiten Differentialquotienten
\weq{4}{\frac{\partial^2 f}{\partial x_1^2} = n(n-1)a_1^2a_x^{n-2} = n(n-1)f_{11}}
\weq{5}{\frac{\partial^2 f}{\partial x_1\partial x_2} = n(n-1)a_1a_2a_x^{n-2} = n(n-1)f_{12}}
\weq{6}{\frac{\partial^2 f}{\partial x_2^2} = n(n-1)a_2^2a_x^{n-2} = n(n-1)f_{22}\,,}
multipliciren sodann die erste Gleichung mit $y_1^2$, die zweite mit $2y_1y_2$, die dritte mit $y_2^2$ und addiren, so kommt:
\[
\begin{aligned}
y_1^2f_{11}+2y_1y_2f_{12}+y_2^2f_{22} &= (a_1^2y_1^2+2a_1a_2y_1y_2+a_2^2y_2^2)a_x^{n-2}\\
&= a_y^2a_x^{n-2} = f_{y^2}\,.
\end{aligned}
\]
Die zweite Polare ist also durch die Gleichung definirt:
\[ f_{y^2} = y_1^2f_{11}+2y_1y_2f_{12}+y_2^2f_{22} \]
oder
\weq{7}{f_{y^2} = \frac{1}{n(n-1)}\left\{y_1^2\frac{\partial^2 f}{\partial x_1^2}+2y_1y_2\frac{\partial^2 f}{\partial x_1\partial x_2}+y_2^2\frac{\partial^2 f}{\partial x_2^2}\right\}.}
Fährt man auf diese Weise fort, so kommt man der Reihe nach zu allen Polaren. Die rechten Seiten dieser Gleichungen (3) und (7) sind aber nichts anderes als die totalen Differentiale von $f=a_x^n$, wobei nur die Increments $dx_1$, $dx_2$ und durch $y_1$ resp. $y_2$ ersetzt sind. Daher der Satz:

\glqq Die $\mu^{\text{te}}$ Polare $a_x^{n-\mu}a_y^\mu$ wird erhalten, indem man $a_x^n$ $\mu$ mal total differentiirt, die Increments durch $y_1$ bezw. $y_2$ ersetzt und endlich mit
%CONT>>


% ===== pdf idx 33, printed 19.0 =====
% !PRINTED 19
\weberpage{19}

%<<CONT
$\dfrac{1}{n(n-1)\ldots(n-\mu+1)}$ multiplicirt\grqq{}; also in leicht verständlicher Form:
\weq{8}{f_{y^k} = \frac{(m-k)!}{k!}\left\{\frac{\partial f}{\partial x_1}y_1+\frac{\partial f}{\partial x_2}y_2\right\}^k.} % [sic] numerator printed with $m$ rather than $n$
Durch diese Methode bestätigt man leicht folgende Sätze:

1) Die Polare einer Summe von Formen ist gleich der Summe der Polaren dieser Formen, also
\[ \left\{\sum f^{(i)}\right\}_{y^k} = \sum f^{(i)}{}_{y^k}\,. \]

2) Die Polare eines Productes von zwei Factoren, deren einer eine Constante $c$ ist, ist gleich dem Producte dieser Constanten in die Polare des andern Factors, also
\[ \{c\cdot f\}_{y^k} = c\cdot f_{y^k}\,. \]

3) Die $\nu^{\text{te}}$ Polare einer Form $n^{\text{ten}}$ Grades ist identisch Null, wenn
%CONT>>

$\nu>n$.

Es ist gleichgiltig, ob man das $\nu^{\text{te}}$ totale Differential nach $x$ von der $\mu^{\text{ten}}$ Polare, oder das $\mu+\nu^{\text{te}}$ totale Differential der nullten Polare, also der Form $f$ selbst bildet, ferner ob man die Originalform $k$ mal total nach $x$ differentiirt und die Increments durch $y$ ersetzt oder die $(\varrho+k)^{\text{te}}$ Polare von $f$ $\varrho$ mal nach $y$ differentiirt und an Stelle der Increments wieder die Variable $x$ einführt. Man bestätigt dies leicht durch Rechnung und überzeugt sich so von der Richtigkeit des Satzes:

4) \gesperrt{Polaren von Polaren sind selbst wieder Polaren der Originalform.}

15. \textit{Beispiel.} Wir bilden die erste Polare der quadratischen Form
\[ f = a_x^2 = \overline{a_0}x_1^2+2\overline{a_1}x_1x_2+\overline{a_2}x_2^2 \]
zuerst nach der ursprünglichen symbolischen Methode, sodann nach der Differentiationsmethode und erhalten:
\[
\begin{aligned}
f_y = a_xa_y &= (a_1x_1+a_2x_2)(a_1y_1+a_2y_2)\\
&= (a_1^2x_1+a_1a_2x_2)y_1+(a_1a_2x_1+a_2^2x_2)y_2\\
&= (\overline{a_0}x_1+\overline{a_1}x_2)y_1+(\overline{a_1}x_1+\overline{a_2}x_2)y_2\,.
\end{aligned}
\]
Ebenso findet man
\[
f_y = \frac{1}{2}\left(y_1\frac{\partial f}{\partial x_2}+y_2\frac{\partial f}{\partial x_2}\right) = \frac{1}{2}\left\{y_1(2\overline{a_0}x_1+2\overline{a_1}x_2)+y_2(2\overline{a_1}x_1+2\overline{a_2}x_2)\right\}
\]
% [sic] both partial derivatives in this line are printed with $\partial x_2$ in the denominator; expected $\partial x_1$ for the first term
\[ = (\overline{a_0}x_1+\overline{a_1}x_2)y_1+(\overline{a_1}x_1+\overline{a_2}x_2)y_2\,. \]

16. \textit{Methode der binomischen Entwicklung.} Die zweite rechnerische Operation, welche auf Polaren führt, stützt sich auf den binomischen Lehrsatz. Ersetzt man nämlich in dem Symbol $a_x$ die Variable
\[
\begin{aligned}
x_1 &\ \text{durch}\ x_1+\lambda y_1\\
x_2 &\ \text{durch}\ x_2+\lambda y_2\,,
\end{aligned}
\]
%CONT>>


% ===== pdf idx 34, printed 20.0 =====
% !PRINTED 20
\weberpage{20}

%<<CONT
so geht dieser lineare Factor über in: $a_x+\lambda a_y$ und sonach $a_x^n$ in:
\[ (a_x+\lambda a_y)^n = a_x^n+\binom{n}{1}\lambda a_x^{n-1}a_y+\binom{n}{2}\lambda^2a_x^{n-2}a_y^2+\cdots+a_y^n = f(x+\lambda y)\,. \]
Hieraus erkennt man:

\glqq Die $\mu^{\text{te}}$ Polare von $a_x^n$ wird erhalten, wenn man in der binomischen Entwicklung von $(a_x+\lambda a_y)^n$ den Coefficienten von $\lambda^\mu$ nimmt, und denselben durch $\binom{n}{\mu}$ dividirt.\grqq{}

\textit{Beispiel.} Um die zweite Polare der Form $f=a_x^3$ zu erhalten, berechnen wir demnach in:
\[
\begin{aligned}
f(x+\lambda y) &= \overline{a_0}(x_1+\lambda y_1)^3+3\overline{a_1}(x_1+\lambda y_1)^2(x_2+\lambda y_2)\\
&\quad+3\overline{a_2}(x_1+\lambda y_1)(x_2+\lambda y_2)^2+\overline{a_3}(x_2+\lambda y_2)^3
\end{aligned}
\]
das in $\lambda^2$ multiplicirte Glied. Wir erhalten
\[
\begin{aligned}
\binom{3}{2}\lambda^2f_{y^2} &= 3\lambda^2\overline{a_0}x_1y_1^2+3\overline{a_1}\lambda^2y_1^2x_2+6\overline{a_1}\lambda^2x_1y_1y_2+6\overline{a_2}\lambda^2x_2y_2y_1\\
&\quad+3\overline{a_2}\lambda^2x_1y_2^2+3\overline{a_3}x_2y_2^2\lambda^2\\
&= 3\lambda^2\left\{(\overline{a_0}x_1+\overline{a_1}x_2)y_1^2+2y_1y_2(\overline{a_1}x_1+\overline{a_2}x_2)+y_2^2(\overline{a_2}x_1+\overline{a_3}x_2)\right\}\cdot
\end{aligned}
\]
Dividiren wir diese Gleichung mit $\binom{3}{2}\lambda^2$, so folgt:
\[ f_{y^2} = (\overline{a_0}x_1+\overline{a_1}x_2)y_1^2+2y_1y_2(\overline{a_1}x_1+\overline{a_2}x_2)+y_2^2(\overline{a_2}x_1+\overline{a_3}x_2)\,. \]

17. \textit{Methode der Transformation.} Wir sahen bereits in Nr. 3, dass durch die Transformation
\weq{I}{x_1 = \xi_1y_1+\eta_1y_2,\quad x_2 = \xi_2y_1+\eta_2y_2}
mit dem Modul $\Delta=(\xi\eta)$ die Form $f=a_x^n$ übergeht in die Form
\[ F = A_y^n = (a_\xi y_1+a_\eta y_2)^n\,, \]
dass also die Beziehungen bestehen:
\[ \overline{A}_k = A_1^{n-k}A_2^k = a_\xi^{n-k}a_\eta^k\,. \]
Die $\overline{A}_k$ sind die wirklichen Coefficienten der transformirten Form; $a_\xi^{n-k}a_\eta^k$ ist aber eine Polare von $f$, nur geschrieben in den Transformationscoefficienten. Daher der Satz:

\glqq Die Coefficienten von $y$ der transformirten Form sind die Polaren der ursprünglichen Form, nur geschrieben in den Substitutionscoefficienten $\xi$ und $\eta$.\grqq{}

Wir können diesen Satz auch noch in anderer Weise formuliren. Multipliciren wir nämlich die Transformationsgleichungen (I) mit $a_1$ resp. $a_2$ und addiren die Resultate, so kommt:
\weq{1}{a_x = a_\xi y_1+a_\eta y_2\,.}
%CONT>>


% ===== pdf idx 35, printed 21.0 =====
% !PRINTED 21
\weberpage{21}
Anderntheils ist nach dem Identitätssatze:
\weq{2}{a_x = a_\xi\,\frac{(x\eta)}{(\xi\eta)} + a_\eta\,\frac{(\xi x)}{(\xi\eta)}.}
Durch Comparation dieser beiden Gleichungen erhält man
\weq{3}{y_1 = \frac{(x\eta)}{(\xi\eta)}, \quad y_2 = \frac{(\xi x)}{(\xi\eta)}.}
Erhebt man aber die Identität (2) auf die $n^{\text{te}}$ Potenz, so kommt:
\[
a_x^n \cdot (\xi\eta)^n = \{a_\xi(x\eta) + a_\eta(\xi x)\}^n
\]
und diese Gleichung lehrt in Verbindung mit (3):
\glqq Multiplicirt man die Form $f = a_x^n$ mit der $n^{\text{ten}}$ Potenz des Transformationsmoduls, so sind die Coefficienten von $(x\eta)^{n-k}(\xi x)^k = (\xi\eta)^n y_1^{n-k} y_2^k$, dividirt durch den betreffenden Binomialcoefficienten, die $k^{\text{ten}}$ Polaren der Originalform.\grqq{}

\textit{18. Erweiterung des Polarenbegriffes.} Im weiteren Sinne nennt man auch Formen, die mehr als zwei Reihen Veränderlicher enthalten, also Formen wie
\[
a_x^{n-k-\lambda}a_y^k a_z^\lambda, \quad a_x^{n-k-\lambda-\mu}a_y^k a_z^\lambda a_u^\mu, \text{ etc.,}
\]
Polaren der Form $f$. Ihr Bildungsprocess ist folgendermassen definirt. Man sondere vom Producte $a_x^n$ zunächst $\varkappa+\lambda+\mu$ Factoren ab, ersetze $\varkappa$ derselben durch $a_y$, $\lambda$ Factoren durch $a_z$, $\mu$ Factoren durch $a_u$; oder: Man nehme zuerst $\varkappa+\lambda+\mu$ Factoren $a_x$ hinweg und ersetze sie durch Factoren $a_y$; von diesen trenne man wieder $\lambda+\mu$ Factoren ab und ersetze sie durch Grössen $a_z$, hierauf bringe man an Stelle von $\mu$ Factoren $a_z$ wiederum $\mu$ Factoren $a_u$.

Man kann auch die so definirten Polaren durch rechnerische Processe erhalten. So würde die Differentiationsmethode folgende Operationen verlangen: Die Form $f$ differentiire man $k+\lambda+\mu$ mal total nach $x$, ersetze die Incremente durch $y$.
% ?? uncertain: "Increment"/"Incremente" — plural form as printed on this page could not be fully confirmed at the line break
Hierauf differentiire man $\lambda+\mu$ mal nach $y$, ersetze die neuen Incremente durch $z$, endlich differentiire man $\mu$ mal total nach $z$, und ersetze die Incremente durch $u$.

\textit{19. Differentialgleichung der Polaren.} Die Polaren sind also Formen mit mehreren Reihen cogredienter Variabler. Doch lässt sich aus ihrer Bildung erkennen, dass sie sehr specielle Formen dieser Art sind. In der That genügen sie auch einer bestimmten Differentialgleichung, die sich leicht aufstellen lässt. Bezeichnet man die Polare mit $U$, so dass also
\[
U = a_x^{n-\lambda}a_y^\lambda\,,
\]
dann ist:
\[
\frac{\partial U}{\partial x_1} = (n-\lambda)\,a_x^{n-\lambda-1}a_y^\lambda a_1
\]
\[
\frac{\partial U}{\partial x_2} = (n-\lambda)\,a_x^{n-\lambda-1}a_y^\lambda a_2\,,
\]
%CONT>>


% ===== pdf idx 36, printed 22.0 =====
% !PRINTED 22
\weberpage{22}
%<<CONT
und daher:
\[
\left.\begin{aligned}
\frac{\partial^2 U}{\partial x_1\partial y_2} &= \lambda(n-\lambda)\,a_x^{n-\lambda-1}a_y^{\lambda-1}a_1a_2 \\
\frac{\partial^2 U}{\partial x_2\partial y_1} &= \lambda(n-\lambda)\,a_x^{n-\lambda-1}a_y^{\lambda-1}a_2a_1
\end{aligned}\right\}
\]
also:
\weq{I}{\frac{\partial^2 U}{\partial x_1\partial y_2} - \frac{\partial^2 U}{\partial x_2\partial y_1} = 0.}
Das ist die partielle Differentialgleichung, welcher jede Polare genügt; wir werden später umgekehrt sehen, dass jede rationale ganze Function von $x$ und $y$, welche diese Gleichung befriedigt, in der That eine Polare ist.

Besitzt die Polare mehr als zwei Reihen Veränderlicher, dann treten zu dieser Gleichung (I) noch weitere analog gebildete, so dass bei $n$ Reihen von Variabeln die Anzahl derselben auf $N$ anwächst, wo $N$ die Anzahl der Combinationen ohne Wiederholung von $n$ Elementen zu je zweien bedeutet. Dieselben sind natürlich nicht durchwegs selbstständige Gleichungen.

\textit{20. Der Omegaprocess.} Durch die Differentialgleichung (I) ist gleichzeitig ein Process definirt, der wie jener der Faltung und Polarenbildung ein Invariantenprocess ist. Wir schreiben denselben, wenn wir ihn auf eine Form $f$, die in $x$ vom Grade $m$, in $y$ vom Grade $n$ ist, angewendet wissen wollen, gewöhnlich in der Gestalt:
\[
\Omega(f) = \frac{1}{m\cdot n}\left\{\frac{\partial^2 f}{\partial x_1\partial y_2} - \frac{\partial^2 f}{\partial x_2\partial y_1}\right\}
\]
und nennen ihn demgemäss Omegaprocess. Man kann ihn durch
\[
\Omega(f) = \begin{vmatrix} x_1 & x_2 \\ y_1 & y_2 \end{vmatrix}
\]
repräsentiren.

Dass durch ihn die Invarianteneigenschaft eines symbolischen Productes nicht aufgehoben wird, erkennt man daraus, weil durch diese Operation aus demselben wieder nur symbolische Producte entstehen.

Wenden wir nämlich den Omegaprocess zunächst auf das symbolische Product
\[
P = a_x^m b_y^n
\]
an; wir erhalten:
\[
\frac{\partial P}{\partial x_1\partial y_2} = n\cdot m\cdot a_1\,a_x^{m-1}b_2b_y^{n-1}
\]
\[
\frac{\partial P}{\partial x_2\partial y_1} = n\cdot m\cdot a_2\,a_x^{m-1}b_1b_y^{n-1}\,.
\]


% ===== pdf idx 37, printed 23.0 =====
% !PRINTED 23
\weberpage{23}
Durch Subtraction folgt aus beiden Gleichungen:
\[
\frac{\partial P}{\partial x_1\partial y_2} - \frac{\partial P}{\partial x_2\partial y_1} = n\cdot m\cdot a_x^{m-1}b_y^{n-1}\{a_1b_2 - b_1a_2\} = n\cdot m\cdot(ab)\,a_x^{m-1}b_y^{n-1}\,,
\]
also:
\weq{I}{\Omega(P) = (ab)\,a_x^{m-1}b_y^{n-1}\,.}
Unterwerfen wir diese Relation abermals demselben Process
\[
\Omega^2(P) = \Omega(\Omega(P)) = \frac{1}{m-1\cdot n-1}\left\{\frac{\partial\Omega(P)}{\partial x_1\partial y_2} - \frac{\partial\Omega(P)}{\partial y_1\partial x_2}\right\}\,,
\]
so erhalten wir:
\[
\Omega^2(P) = (ab)^2 a_x^{m-2}b_y^{n-2}\,, \text{ u.\,s.\,w.}
\]
Aus symbolischen Producten entstehen also durch diese Operation immer wieder symbolische Producte, und man erkennt aus dem Beispiel die enge Beziehung, die zwischen Omega- und Faltungsprocess besteht. Alle symbolischen Producte, welche aus einem Stammproducte $P_x\cdot Q_x$ dadurch abgeleitet werden, dass man immer je einen linearen Factor von $P_x$ mit einem von $Q_x$ faltet, entstehen auch aus $P_x\cdot Q_y$, oder $Q_x\cdot P_y$ durch den Omegaprocess.

Seine Invarianteneigenschaft wird übrigens auch noch durch eine gewisse Relation bestätigt, die ihn mit Polarenbildung in Verbindung bringt.

Polarisirt man nämlich das Product
\[
P = a_x^m b_y^n
\]
nach $x$ in Bezug auf $y$, so erhält man
\[
P_y = a_x^{m-1}a_y b_y^n\,.
\]
Bildet man aber hiervon wieder die erste Polare nach $y$ in Bezug auf $x$, so kommt unter Anwendung der Differentiationsmethode:
\[
(n+1)\sum\frac{\partial P_y}{\partial y_i}x_i = a_x^m b_y^n + n\,a_x^{m-1}a_y b_x b_y^{n-1}\,.
\]
Hievon auf beiden Seiten
\[
(n+1)P = a_x^m b_y^n + n\,a_x^m b_y^n
\]
subtrahirt, liefert die Relation:
\[
(n+1)\left\{P - \sum\frac{\partial P_y}{\partial y_i}x_i\right\} = n(ab)(xy)\,a_x^{m-1}b_y^{n-1} = n\,(xy)\cdot\Omega(P)\,,
\]
oder:
\weq{II}{P = (P_y)_x + \frac{n}{n+1}(xy)\,\Omega(P).}
Clebsch hat diese Relation benutzt, um die Form $P$ in eine Reihe zu entwickeln, die nach Potenzen von $(xy)$ fortschreitet, und deren Coefficienten Polaren gewisser Formen sind (Clebsch, \textit{Theorie der binären Formen}, § 7. — Vgl. auch Nr. 25 und 28 dieses §).


% ===== pdf idx 38, printed 24.0 =====
% !PRINTED 24
\weberpage{24}
\textit{21. Polaren eines Productes mehrerer Formen.} In den meisten Fällen hat man nicht die Polare einer einzigen Form zu bilden, sondern eines Productes von Formen. Um dieselbe aufzustellen, benutze ich zunächst die am Schlusse von Nr. 13 entwickelte Methode. Es sei das symbolische Product, dessen Polare zu berechnen ist:
\[
\Pi = (ab)(ac)(bc)\dots a_x^\varrho b_x^\sigma c_x^\tau \dots\,.
\]
Die Klammerfactoren bleiben als Constante $C$ von dem Polarenprocess unberührt. Die Factoren erster Art, — es seien $\lambda$ an der Zahl — ersetzen wir wieder durch $\lambda$ lineare Factoren $r_{ix}$ und bilden sonach die Polare von
\[
F = C\,r_{1x}r_{2x}r_{3x}\dots r_{\lambda x}\,.
\]
Dann erhalte ich hieraus die $k^{\text{te}}$ Polare, indem ich in allen möglichen Combinationen ohne Wiederholung von $\lambda$ Elementen je $k$ Elemente $r_{ix}$ herausgreife und jedesmal durch die $k!$ Permutationen der entsprechenden Elemente $r_{iy}$ ersetze. Die Summe aller dieser Glieder, dividirt durch ihre Anzahl $\binom{\lambda}{k}k!$ liefert die Polare von $F$. Wir können sie also schreiben:
\weq{I}{F_{y^k} = \frac{C}{k!\binom{\lambda}{k}}\sum_{i_1i_2\dots i_k} r_{i_1y}r_{i_2y}\dots r_{i_ky}\,\frac{F}{r_{i_1x}\dots r_{i_kx}}\,.}
Ersetzen wir hierin wieder die Symbole
\[
r_{1y}, r_{2y}\dots r_{\varrho y} \text{ durch } a_x
\]
\[
r_{\varrho+1,\,y}\dots r_{\sigma y} \text{ durch } b_x
\]
\[
r_{\sigma+1,\,y},\, r_{\sigma+2,\,y}\dots r_{\tau y} \text{ durch } c_x \text{ etc.,}
\]
so stellt das Resultat die $k^{\text{te}}$ Polare von $\Pi$ dar.

\textit{22. Darstellung dieses Polarenprocesses durch binomische Entwicklung.} Wir können nun wieder diese symbolische Methode durch Rechnung ersetzen und wählen zu dem Zwecke die binomische Entwicklung. Ich will mit einem Producte von zwei Formen beginnen. Es sei
\[
F = a_x^m\cdot b_x^n = p_x^{m+n}\,.
\]
Die Polare von $F$
\[
F_{y^k} = p_x^{m+n-k}p_y^k
\]
entsteht alsdann, wenn wir in $F$ die Variable $x_i$ durch $x_i+\lambda y_i$ ersetzen und den Coefficienten von $\binom{m+n}{k}\lambda^k$ in der Entwicklung von $p(x+\lambda y)^{m+n}$ nehmen.

Diesen Coefficienten erhalten wir auf folgende Weise. Es ist:
%CONT>>


% ===== pdf idx 39, printed 25.0 =====
% !PRINTED 25
\weberpage{25}
\[
\begin{aligned}
(a_x+\lambda a_y)^m &= \binom{m}{0}a_x^m+\binom{m}{1}\lambda a_x^{m-1}a_y+\binom{m}{2}\lambda^2 a_x^{m-2}a_y^2+\binom{m}{3}\lambda^3 a_x^{m-3}a_y^3+\cdots\\
(b_x+\lambda b_y)^n &= \binom{n}{0}b_x^n+\binom{n}{1}\lambda b_x^{n-1}b_y+\binom{n}{2}\lambda^2 b_x^{n-2}b_y^2+\binom{n}{3}\lambda^3 b_x^{n-3}b_y^3+\cdots
\end{aligned}
\]
Durch Multiplication der beiden Entwicklungen bekommt man als Co\"efficienten von $\lambda^k$:
\[
\begin{aligned}
\binom{m+n}{k}F_{y^k} &= \binom{m+n}{k}p_x^{m+n-k}p_y^k\\
&= \binom{m}{0}\binom{n}{k}a_x^m b_x^{n-k}b_y^k+\binom{m}{1}\binom{n}{k-1}a_x^{m-1}a_y b_x^{n-k+1}b_y^{k-1}\\
&\quad+\binom{m}{2}\binom{n}{k-2}a_x^{m-2}a_y^2 b_x^{n-k+2}b_y^{k-2}+\cdots\\
&\quad\cdots+\binom{m}{k}\binom{n}{0}a_x^{m-k}a_y^k b_x^m\,.
\end{aligned}
\]
Die rechte Seite dieser Gleichung, dividirt durch $\binom{m+n}{k}$, stellt also die $k^{\text{te}}$ Polare von $a_x^m\cdot b_x^n$ dar.

23. \textit{Relation zwischen den Co\"efficienten der Polare.} In die zuletzt erhaltene Entwicklung führen wir folgende Abkürzungen ein. Wir setzen:
\[
c_0=\frac{\binom{m}{0}\binom{n}{k}}{\binom{m+n}{k}},\quad c_1=\frac{\binom{m}{1}\binom{n}{k-1}}{\binom{m+n}{k}},\quad\cdots,\quad c_\varrho=\frac{\binom{m}{\varrho}\binom{n}{k-\varrho}}{\binom{m+n}{k}},
\]
ferner:
\[
G_0=a_x^m b_x^{n-k}b_y^k,\ G_1=a_x^{m-1}a_y b_x^{n-k+1}b_y^{k-1},\ \cdots,\ G_\varrho=a_x^{m-\varrho}a_y^\varrho b_x^{n-k+\varrho}b_y^{k-\varrho},
\]
dann ist:
\[
F_{y^k}=c_0G_0+c_1G_1+c_2G_2+\cdots c_\varrho G_\varrho\cdots c_kG_k
\]
\weq{I}{F_{y^k}=\sum_0^k c_\varrho G_\varrho\,.}
Die Grössen $G_\varrho$ bezeichnen wir kurz als \glqq Glieder\grqq{} der Polare, das Glied $G_0$ führt den Namen \glqq Anfangsglied\grqq{}. Jedes Glied hat, wie man unmittelbar ersieht, die Eigenschaft, direct in $F$ überzugehen, sobald man $y$ durch $x$ ersetzt. Da anderntheils $F_{y^k}$ selbst in $F$ übergeht durch diese Substitution, so folgt aus (I) die Identität:
\[
F=c_0F+c_1F+c_2F+\cdots c_kF
\]
oder nach Division mit $F$
\weq{II}{1=c_0+c_1+c_2+\cdots c_k\,,}
d.\,h. \glqq die Summe der Co\"efficienten in der Polare $F_{yk}$ ist immer gleich Eins.\grqq{}

24. \textit{Benachbarte Glieder.} Jedes Glied $G_\varrho$ kann man sich auf folgende Art symbolisch entstanden denken: Man schreibt $m$ Factoren $a_x$ und $n$ Factoren $b_x$ an, nimmt beliebig $\varrho$ ($\varrho\leqq k$) Factoren hinweg und
%CONT>>


% ===== pdf idx 40, printed 26.0 =====
% !PRINTED 26
\weberpage{26}
%<<CONT
ersetzt sie entsprechend durch Factoren $a_y$ resp. $b_y$. Dabei giebt der Zähler des zugehörigen Co\"efficienten $c_\varrho$ immer an, wie oft durch Herbeiführung immer anderer und anderer Factoren diese Operation vorgenommen werden kann, nämlich $\binom{m}{\varrho}\binom{n}{k-\varrho}$ mal. Zwei aufeinanderfolgende Glieder $G_\varrho$ und $G_{\varrho+1}$
\[
\begin{aligned}
G_\varrho &= a_x^{m-\varrho}a_y^\varrho b_x^{n-k+\varrho}b_y^{k-\varrho}\\
G_{\varrho+1} &= a_x^{m-\varrho-1}a_y^{\varrho+1}b_x^{n-k+\varrho+1}b_y^{k-\varrho-1}
\end{aligned}
\]
unterscheiden sich nur dadurch, dass an Stelle des Factorenpaares $a_xb_y$, das Product $a_yb_x$ getreten ist, d.\,h. ein vorhergehendes Glied hat einen Factor $a_x$ und einen Factor $b_y$ mehr, als ein nachfolgendes. Zwei Glieder, die sich nur um zwei solchen linearen Factoren unterscheiden, nennen wir \glqq\gesperrt{benachbarte Glieder}\grqq{}. (Vgl. auch Nr. 10.) Ihre Differenz liefert:
\[
G_\varrho-G_{\varrho+1}=a_x^{m-\varrho-1}a_y^\varrho b_x^{n-k+\varrho}b_y^{k-\varrho-1}(a_xb_y-b_xa_y)\,.
\]
Nun ist aber nach dem Identitätssatze Nr.\ 11, (III):
\[
(a_xb_y-b_xa_y)=(ab)(xy)\,.
\]
Daher wird:
\[
G_\varrho-G_{\varrho+1}=(ab)(xy)\cdot a_x^{m-\varrho-1}a_y^\varrho\cdot b_x^{n-k+\varrho}b_y^{k-\varrho-1},\ \text{d.\,h.}
\]
\glqq Die Differenz je zweier benachbarter Glieder einer Polare enthält immer den Factor $(ab)(xy)$.\grqq{}

25. \textit{Verallgemeinerung dieses Satzes.} Dieser Satz gilt aber auch für die Differenz zweier beliebiger Glieder
\[
G_\varrho-G_\sigma\,.
\]
Denn es ist:
\[
G_\varrho-G_\sigma=G_\varrho-G_{\varrho+1}+G_{\varrho+1}-G_{\varrho+2}+G_{\varrho+2}-\cdots+G_{\sigma-1}-G_\sigma\,,
\]
wo nun rechts immer je zwei Glieder benachbart sind. Er gilt endlich aber auch für die Differenz
\[
F_{y^k}-G_\varrho\,.
\]
Denn multipliciren wir die Gleichung (II) Nr.\ 23 mit $G_\varrho$ und subtrahiren sie von Gleichung (I), dann erhalten wir
\[
F_{y^k}-G_\varrho=c_0(G_0-G_\varrho)+c_1(G_1-G_\varrho)+\cdots+c_k(G_k-G_\varrho)\,.
\]
Daher der Satz:

\glqq Die Differenz zwischen Polare und einem ihrer Glieder enthält immer den Factor $(ab)(xy)$.\grqq{}

Wir wollen diesen fundamentalen Satz noch etwas weiter verfolgen. Der andere Factor, der mit $(ab)(xy)$ zusammen die Differenz $F_{y^k}-G_\varrho$ darstellt, besteht aus einer Summe von Gliedern $H_\sigma$, die in
%CONT>>


% ===== pdf idx 41, printed 27.0 =====
% !PRINTED 27
\weberpage{27}
%<<CONT
$x$ wie in $y$ um je einen Grad niedriger sind. Multipliciren wir mit dem einen Factor $(ab)$ in diese Summe hinein, so lässt sich jedes Glied
\[
(ab)\,H_\sigma
\]
als Glied einer gewissen Polare auffassen, die um einen Grad niedriger ist als die ursprüngliche Polare. Ueberdies ist diese $(k-1)^{\text{te}}$ Polare nicht wie $F_{y^k}$ von dem Producte
\[
f\cdot\varphi=a_x^m b_x^n
\]
der beiden ursprünglichen Formen gebildet, sondern — eben wegen des Klammerfactors $(ab)$ — von einer Form, die aus ihm durch einmalige Faltung entstanden ist. Bezeichnen wir diese Form mit
\[
(f,\varphi)^1\,,
\]
ihre Polare mit:
\[
(f,\varphi)_{y^{k-1}}
\]
und ein Glied derselben mit
\[
G_{(f,\varphi)_{y^{k-1}}}\,,
\]
so kann man den eben erhaltenen Satz in die algebraische Form bringen
\weq{I}{G_\varrho=F_{y^k}+(xy)\sum c_iG_{(f,\varphi)_{y^{k-1}}}\,.}
Die Iteration liegt nahe. Ich ersetze die einzelnen Glieder der Summe in (I) gerade nach diesem Gesetze wieder durch die ihnen zugehörigen $(k-1)^{\text{ten}}$ Polaren und eine Summe anderer Glieder $K$, die neuerdings den Factor $(ab)(xy)$ erhalten, also $(k-2)^{\text{ten}}$ Polaren angehören von Formen, die durch zweimalige Faltung aus $f\cdot\varphi$ entstanden sind. Der Ausdruck (I) wird dann:
\[
G_\varrho=F_{y^k}+(xy)\sum c_i(f,\varphi)_{y^{k-1}}+(xy)^2\sum c_i^{(1)}K_{(f,\varphi)^2_{y^{k-2}}}
\]
% ?? uncertain: the exponent/subscript on K in the last term is partly obscured in the scan (looks like "k-1\cdot 2"); transcribed here as $(f,\varphi)^2_{y^{k-2}}$ by analogy with the corresponding term on the next line of print.
Wiederholt man diesen Process so lange, bis man schliesslich auf nullte Polaren kommt, von Formen, die aus $f\cdot\varphi$ durch $k$malige Faltung hervorgingen, so gelangt man endlich zu dem Resultat:
\weq{II}{G_\varrho=F_{y^k}+(xy)\sum c_i(f,\varphi)_{y^{k-1}}+(xy)^2\sum c_i^{(1)}(f,\varphi)^2_{y^{k-2}}+\cdots+(xy)^k\sum c_i^{k-1}(f,\varphi)^k_{y^0}\,,\ \text{d.\,h.}}
\glqq Jedes Glied einer Polare lässt sich entwickeln in eine Summe von Polaren, multiplicirt mit Potenzen von $(xy)$. Das erste Glied dieser Summe ist die ursprüngliche Polare selbst und hat den Co\"efficienten Eins.\grqq{}

Für den hier betrachteten einfachen Fall werden wir später ($\S.\,8$) die Reihenentwicklung mit Angabe der Co\"efficienten geben.

26. \textit{Beispiel.} Ich will die bisher gewonnenen Resultate an einem Beispiele verificiren. Es sei gegeben das Product
\[
F=f\cdot\varphi=a_x^3\cdot b_x^3=r_{1x}r_{2x}r_{3x}\cdot r_{4x}r_{5x}r_{6x}\,.
\]


% ===== pdf idx 42, printed 28.0 =====
% !PRINTED 28
\weberpage{28}
Die zweite Polare dieses Products ist nach der ursprünglichen Definition (Nr.\ 21) dargestellt durch:
\[
F_{y^2}=\frac{1}{2!\binom{6}{2}}\sum r_{i_1y}r_{i_2y}\cdot\frac{F}{r_{i_1x}r_{i_2x}}
\]
\[
\begin{aligned}
=\frac{2!\,1}{2!\,15}\Big\{
&r_{1y}r_{2y}\,r_{3x}r_{4x}r_{5x}r_{6x}+r_{1y}r_{3y}\,r_{2x}r_{4x}r_{5x}r_{6x}+r_{1y}r_{4y}\,r_{2x}r_{3x}r_{5x}r_{6x}\\
&+r_{1y}r_{5y}\,r_{2x}r_{3x}r_{4x}r_{6x}+r_{1y}r_{6y}\,r_{2x}r_{3x}r_{4x}r_{5x}+r_{2y}r_{3y}\,r_{1x}r_{4x}r_{5x}r_{6x}\\
&+r_{2y}r_{4y}\,r_{1x}r_{3x}r_{5x}r_{6x}+r_{2y}r_{5y}\,r_{1x}r_{3x}r_{4x}r_{6x}+r_{2y}r_{6y}\,r_{1x}r_{3x}r_{4x}r_{5x}\\
&+r_{3y}r_{4y}\,r_{1x}r_{2x}r_{5x}r_{6x}+r_{3y}r_{5y}\,r_{1x}r_{2x}r_{4x}r_{6x}+r_{3y}r_{6y}\,r_{1x}r_{2x}r_{4x}r_{5x}\\
&+r_{4y}r_{5y}\,r_{1x}r_{2x}r_{3x}r_{6x}+r_{4y}r_{6y}\,r_{1x}r_{2x}r_{3x}r_{5x}+r_{5y}r_{6y}\,r_{1x}r_{2x}r_{3x}r_{4x}
\Big\}\,.
\end{aligned}
\]
Ersetzen wir hierin $r_1r_2r_3$ durch $a$, $r_4r_5r_6$ durch $b$, so erhalten wir:
\[
\begin{aligned}
F_{y^2}=\frac{1}{15}\Big\{
&a_y^2a_xb_x^3+a_y^2a_xb_x^3+a_yb_ya_x^2b_x^2+a_yb_ya_x^2b_x^2+a_y^2b_ya_x^2b_x\\
&+a_y^2a_xb_x^3+a_yb_ya_x^2b_x^2+a_yb_ya_x^2b_x^2+a_yb_ya_x^2b_x^2+a_yb_ya_x^2b_x\\
&+a_yb_ya_x^2b_x^2+a_yb_ya_x^2b_x^2+b_y^2a_x^3b_x+b_y^2a_x^3b_x+b_y^2a_x^3b_x
\Big\}\,.
\end{aligned}
\]
Fassen wir alle gleichen Glieder zusammen, so kommt:
\[
F_{y^2}=\frac{1}{5}\Big\{a_y^2a_xb_x^3+3a_yb_ya_x^2b_x^2+b_y^2a_x^3b_x\Big\}\,.
\]
In dem nämlichen Resultate gelangen wir nach Nr.\ 22 durch die binomische Entwicklung. Darnach ist:
\[
F_{y^2}=\frac{1}{\binom{6}{2}}\Big\{\binom{3}{2}a_x^3b_xb_y^2+\binom{3}{1}\binom{3}{1}a_x^2a_yb_x^2b_y+\binom{3}{2}a_xa_y^2b_x^3\Big\}
\]
\[
=\frac{1}{5}\Big\{a_x^3b_xb_y^2+3a_x^2a_yb_x^2b_y+a_xa_y^2b_x^3\Big\}\,.
\]
Bilden wir nun die Differenz der Polare und des zweiten Gliedes. Man erhält:
\[
\begin{aligned}
F_{y^2}-\frac{5}{5}a_x^2a_yb_x^2b_y &=\frac{1}{5}(a_x^3b_xb_y^2-a_x^2a_yb_x^2b_y)+\frac{1}{5}(a_xa_y^2b_x^3-a_x^2a_yb_x^2b_y)\\
&=\frac{1}{5}(ab)a_x^2b_xb_y(xy)+\frac{1}{5}(ab)a_xa_yb_x^2(xy)
\end{aligned}
\]
oder
\[
F_{y^2}-a_x^2a_yb_x^2b_y=\frac{(xy)}{5}\Big\{(ab)a_x^2b_xb_y+(ab)a_xa_yb_x^2\Big\}\,.
\]
Jeder der beiden Terme im zweiten Factor rechts stellt ein Glied der ersten Polare einer Form dar, die aus $a_x^3b_x^3$ durch Faltung hervorgeht, nämlich der Form:
\[
\Phi=(ab)a_x^2b_x^2\,.
\]
Die erste Polare dieser Form ist aber gerade
\[
\Phi_y=\frac{1}{2}\Big\{(ab)a_x^2b_xb_y+(ab)a_xa_yb_x^2\Big\}\,.
\]


% ===== pdf idx 43, printed 29.0 =====
% !PRINTED 29
\weberpage{29}
Folglich ist
\[
a_x^2 b_x^2 a_y b_y = F_{y^2} - \frac{2}{5} (xy)\, \Phi_y .
\]
27. \textit{Beispiel.} Berechnen wir ebenso die vierte Polare von
\[
P = a_x^6 b_x^4 .
\]
Sie ist, wenn wir uns wieder lineare Factoren $r_{ix}$ eingeführt denken, dargestellt durch
\[
\binom{10}{4} P_{y^4} = \sum r_{1y} r_{2y} r_{3y} r_{4y} \cdot \frac{P}{r_{1x} r_{2x} r_{3x} r_{4x}} .
\]
Die 6 Factoren $a_x$ lassen sich $\binom{6}{4}$ mal zu je 4 Factoren $r_{ix}$ gruppiren.
Also ist das Anfangsglied sammt Coefficienten
\[
c_0 G_0 = \binom{6}{4} \binom{4}{0} a_y^4 a_x^2 b_x^4 .
\]
Sie lassen sich ferner $\binom{6}{3}$ mal zu je 3 Factoren $r_{ix}$ verbinden. Jede dieser Complexionen kann mit je einem Factor $b_x$ multiplicirt auftreten; folglich ist
\[
c_1 G_1 = \binom{6}{3} \binom{4}{1} a_y^3 b_y a_x^3 b_x^3 .
\]
Auf gleiche Weise findet man die drei anderen Glieder und erhält so
\[
\binom{10}{4} P_{y^4} = \binom{6}{4} \binom{4}{0} a_y^4 a_x^2 b_x^4 + \binom{6}{3} \binom{4}{1} a_y^3 b_y a_x^3 b_x^3 + \binom{6}{2} \binom{4}{2} a_y^2 b_y^2 a_x^4 b_x^2 + \binom{6}{1} \binom{4}{3} a_y b_y^3 a_x^5 b_x + \binom{6}{0} \binom{4}{4} b_x^4 a_x^6
\]
oder:
\[
P_{y^4} = \frac{1}{210} \left\{ 15 a_y^4 a_x^2 b_x^4 + 80 a_y^3 b_y a_x^3 b_x^3 + 90 a_y^2 b_y^2 a_x^4 b_x^2 + 24 a_y b_y^3 a_x^5 b_x + b_y^4 a_x^6 \right\} .
\]
Zunächst erkennt man: die Summe aller Coefficienten ist Eins. Subtrahiren wir nun von $P_{y^4}$ das Glied
\[
G_2 = \frac{1}{210} \left\{ 15 a_y^2 b_y^2 a_x^4 b_x^2 + 80 a_y^2 b_y^2 a_x^4 b_x^2 + 90 a_y^2 b_y^2 a_x^4 b_x^2 + 24 a_y^2 b_y^2 a_x^4 b_x^2 + a_y^2 b_y^2 a_x^4 b_x^2 \right\} ,
\]
so erhalten wir:
\[
P_{y^4} - G_2 = \frac{1}{210} \left\{ 15 (G_0 - G_2) + 80 (G_1 - G_2) + 24 (G_3 - G_2) + G_4 - G_2 \right\} .
\]
Nun ist
\[
G_1 - G_2 = (ab)(yx)\, a_y^2 b_y a_x^3 b_x^2
\]
\[
G_3 - G_2 = (ab)(xy)\, a_y b_y^2 a_x^4 b_x
\]
und folglich:
\[
G_0 - G_2 = G_0 - G_1 + G_1 - G_2 = (ab)(yx)\, a_y^3 a_x^2 b_x^3 + (ab)(yx)\, a_y^2 b_y a_x^3 b_x^2
\]
\[
G_4 - G_2 = G_4 - G_3 + G_3 - G_2 = (ab)(xy)\, b_y^3 a_x^5 b_x + (ab)(xy)\, a_y b_y^2 a_x^2 b_x .
\]
%CONT>>


% ===== pdf idx 44, printed 30.0 =====
% !PRINTED 30
\weberpage{30}
%<<CONT
Folglich ist:
\[
P_{y^4} - G_2 = \frac{(ab)(xy)}{210} \left\{ -15 \left( a_y^2 b_y a_x^3 b_x^2 + a_y^3 a_x^2 b_x \right) - 80 a_y^2 b_y a_x^3 b_x^2 + 24 a_y b_y^2 a_x^4 b_x + b_y^3 a_x^5 + a_y b_y^2 a_x^4 b_x \right\}
\]
\[
= \frac{(ab)(xy)}{210} \left\{ -15 a_y^3 a_x^2 b_x^3 - 95 a_y^2 b_y a_x^3 b_x^2 + 25 a_y b_y^2 a_x^4 b_x + b_y^3 a_x^5 \right\} .
\]
Demnach wird:
\[
G_2 = P_{y^4} + (xy) \left\{ \frac{1}{14} (ab)\, a_y^3 a_x^2 b_x^3 + \frac{19}{42} (ab)\, a_y^2 b_y a_x^3 b_x^2 - \frac{5}{42} (ab)\, a_y b_y^2 a_x^4 b_x - \frac{1}{210} (ab)\, b_y^3 a_x^5 \right\} .
\]
Die Glieder, welche sich rechts innerhalb der Klammer befinden, sind durchwegs Glieder der dritten Polare von
\[
Q = (ab)\, a_x^5 b_x^3 .
\]
Man könnte sie also wieder ersetzen durch $Q_{y^3}$ plus einer Summe von Gliedern zweiter Polaren u.\,s.\,w.

28. \textit{Polaren eines beliebigen Productes von Formen.} Wir haben bereits in Nr. 21 den Process definirt, durch welchen die Polare eines Productes von mehr als zwei Formen erhalten wird. Es erübrigt noch, rechnerische Operationen anzugeben, durch welche dieser symbolische Process ersetzt werden kann. Am einfachsten bedienen wir uns wieder der binomischen Entwicklung. Sei nunmehr $F$ das Product von $\tau$ Formen
\[
f = a_x^m, \quad \varphi = b_x^n, \quad \psi = c_x^p, \quad \ldots \ \text{etc.}
\]
Dann ist die $k^{\text{te}}$ Polare von $F$ durch folgende Grösse definirt. Wir ersetzen in $F$ $x_i$ durch $x_i + \lambda y_i$, bilden die binomischen Entwicklungen der Formen
\[
(a_x + \lambda a_y)^m, \quad (b_x + \lambda b_y)^n, \quad (c_x + \lambda c_y)^n, \quad \ldots \ \text{etc.}
\]
und multipliciren die so erhaltenen Summen. Der Coefficient von $\lambda^k$ dividirt durch
\[
\binom{m+n+p+\cdots s}{k} ,
\]
liefert die $k^{\text{te}}$ Polare von $F$. Die Coefficienten dieser Polare haben demnach die Form
\[
\frac{\binom{m}{i_1}\binom{n}{i_2}\binom{p}{i_3}\cdots\binom{s}{i_k}}{\binom{m+n+p+\cdots s}{k}} % ?? uncertain: subscript under lower denominator binomial shows a smudged mark that may not simply be k
\]
und ihre Summe ist wiederum gleich Eins, wie aus der ursprünglichen Definition Nr. 21 unmittelbar ersichtlich ist.

29. \textit{Beispiele.} Wir berechnen die dritte Polare des Productes:
\[
F = a_x^3 b_x^2 c_x .
\]


% ===== pdf idx 45, printed 31.0 =====
% !PRINTED 31
\weberpage{31}
Es ist:
\[
(a_x + \lambda a_y)^3 = \binom{3}{0} a_x^3 + \binom{3}{1} \lambda a_x^2 a_y + \binom{3}{2} \lambda^2 a_x a_y^2 + \binom{3}{3} \lambda^3 a_y^3
\]
\[
(b_x + \lambda b_y)^2 = \binom{2}{0} b_x^2 + \binom{2}{1} \lambda b_x^2 b_y + \binom{2}{2} \lambda^2 b_y^2
\]
\[
c_x + \lambda c_y = \binom{1}{0} c_x + \binom{1}{1} \lambda c_y .
\]
Die dritte Polare ist somit:
\[
F_{y^3} = \frac{1}{\binom{6}{3}} \Big\{ \binom{1}{0}\binom{2}{0}\binom{3}{3}\, a_y^3 b_x^2 c_x + \binom{1}{0}\binom{2}{1}\binom{3}{2}\, a_x^2 a_x b_y b_x c_x % [sic]
\]
\[
+ \binom{1}{0}\binom{2}{2}\binom{1}{1}\, a_y b_y^2 a_x^2 c_x + \binom{1}{1}\binom{2}{1}\binom{1}{1}\, a_y b_y c_y a_x^2 b_x % [sic]
\]
\[
+ \binom{1}{1}\binom{2}{2}\binom{3}{0}\, a_x^3 b_y^2 c_x \Big\} . % [sic] upper binomial indices printed inconsistently across the terms of this display
\]
Ich möchte hierbei noch aufmerksam machen auf die Zahl der Glieder dieser Polare. Sie ist nämlich gleich der Zahl der Variationen mit Wiederholungen von 4 Elementen $(0, 1, 2, 3)$ zur Summe 3. Dies lässt sich an diesem einfachen Beispiel direct ersehen; denn in den Coefficienten der einzelnen Glieder:
\[
\binom{1}{0}\binom{2}{0}\binom{3}{3}, \quad \binom{1}{0}\binom{2}{1}\binom{3}{2}, \quad \binom{1}{0}\binom{2}{2}\binom{3}{1}, \quad \binom{1}{1}\binom{2}{1}\binom{3}{1}, \quad \binom{1}{1}\binom{2}{2}\binom{3}{0}
\]
stellen gerade die untern Zahlenreihen die erwähnten Complexionen dar. Dieses Gesetz, das aus dieser Art der Polarenberechnung unmittelbar hervorgeht, gilt auch allgemein.

Ueber die Anzahl dieser Variationen hat \gesperrt{Cayley} im Jahre 1856, Phil. Trans. Bd. 146, folgendes Gesetz aufgestellt:
\glqq Eine Zahl $p$ kann als Summe von $\varrho$ Zahlen der Reihe
$0, 1, 2, \ldots$ so oft dargestellt werden, als der Coefficient $x^p z^{\varrho}$
in der Entwicklung nach steigenden Potenzen von $z$ der
Function
\[
\Phi = \frac{1}{(1-z)(1-xz)(1-x^2 z) \ldots (1-x^{\varrho} z)}
\]
beträgt. (Vgl.\ auch \gesperrt{Faà di Bruno}, Theorie der binären Formen, übersetzt von \gesperrt{Walter}, Seite 124.)

30. \textit{Benachbarte Glieder.} Auch hier, bei Polaren eines Productes von Formen, sind wiederum zwei Glieder benachbart, wenn sie bis auf ein Factorenpaar $q_x \cdot r_y$ übereinstimmen (wobei $r$ und $q$ für irgend eines der Symbole $a, b, c, d, \ldots$ etc. stehen). Die Differenz zweier benachbarter Glieder besitzt wieder das Factorenpaar $(qr)(xy)$ zu Factoren. Bildet man alsdann die Differenz zweier beliebiger Glieder, oder die Differenz von Polare und einem ihrer Glieder, so hat diese Differenz zunächst den einen Factor $(xy)$. Der zweite Factor gliedert
%CONT>>


% ===== pdf idx 46, printed 32.0 =====
% !PRINTED 32
\weberpage{32}
%<<CONT
sich in eine Reihe von Summen, welche je mit einem der Klammerfactoren $(qr)$ multiplicirt auftreten. Jede solche Theilsumme ist somit ein Aggregat von Gliedern, die $(k-1)^{\text{ten}}$ Polaren angehören von Formen, welche aus $F$ durch Faltung entstanden sind. Es gilt somit noch immer der Satz:

\glqq Jedes Glied einer Polare von $F$ lässt sich darstellen durch die Polare selbst plus einer Summe von Gliedern, welche selbst wieder niedrigerern Polaren von Formen angehören, die aus $F$ durch Faltung entstanden sind.\grqq{} % [sic] \textquotedbl{}niedrigerern\textquotedbl{}

Wendet man auf diese Glieder niedrigerer Polaren den gleichen Satz an und wiederholt ihn so lange, bis alle Glieder durch Polaren ersetzt sind, so erkennt man die Richtigkeit des Satzes:

\glqq Jedes Glied einer Polare und damit jedes symbolische Product mit zwei Reihen Veränderlicher $x$ und $y$ kann in eine Summe von Polaren entwickelt werden, die nach Potenzen von $(xy)$ fortschreitet.\grqq{}

31. \textit{Polaren mit mehr als zwei Reihen cogredienter Variabler.} Wir haben schon in Nr. 18 den Polarenbegriff für eine Form dahin erweitert, dass wir als Polare auch das symbolische Product definirten:
\[
f_{y^k z^{\lambda}} = a_x^{n-k-\lambda} a_y^k a_z^{\lambda} .
\]
Aehnlich können wir allgemein ein beliebiges symbolisches Product nach mehreren Reihen Variabler $y, z, u, v \ldots$ etc. polarisiren. Wir ersetzen in dem symbolischen Product
\[
F = f \cdot \varphi \cdot \psi \ldots ,
\]
welches vom Grade $m + n + p + \cdots = s$ in $x$ sein möge, die Variabeln $x_i$ durch $x_i + \varrho_1 y_i + \varrho_2 z_i + \cdots$, entwickeln die einzelnen Factoren $f, \varphi, \psi \ldots$ nach dem polynomischen Lehrsatz und multipliciren die so erhaltenen Summen. Irgend eine Polare, etwa die Polare $F_{y^k z^{\lambda} u^{\nu}}$ ist alsdann dargestellt durch den Coefficienten von $\varrho_1^k \varrho_2^{\lambda} \varrho_3^{\nu}$, dividirt durch den Polynomialcoefficienten:
\[
\frac{s!}{k! \, \lambda! \, \nu! \, (s-k-\lambda-\nu)!} .
\]
Wir nennen sie wiederum eine \gesperrt{gemischte Polare} und alle Sätze, welche über die einfachen Polaren und deren Glieder aufgestellt wurden, lassen sich auf diese Polaren und deren Glieder übertragen. Insbesondere gilt der Satz:

\glqq Jedes symbolische Product mit beliebig vielen Reihen cogredienter Veränderlicher $x, y, z, u \ldots$ kann in eine Summe von Polaren entwickelt werden, multiplicirt mit Potenzen von $(xy), (yz), (xz) \ldots$ etc.


% ===== pdf idx 47, printed 33.0 =====
% !PRINTED 33
\weberpage{33}

\begin{center}
\S.\,3.

\textbf{Ueberschiebungsprocess.}
\end{center}

32. \textit{Definition.} Die einfachsten symbolischen Producte $u$ sind
diejenigen, welche nur zwei Symbole enthalten:
\[
u = (ab)^k a_x^{m-k} b_x^{n-k}.
\]
Man nennt sie \glqq Ueberschiebungen\grqq{}, und der Process, durch den sie
entstehen, ist der Faltungsprocess: Das Product
\[
P = a_x^m \cdot b_x^n
\]
der beiden Formen $f = a_x^m$, $\varphi = b_x^n$ wird $k$ mal gefaltet. Wir
werden eine solche $k^{\text{te}}$ Ueberschiebung in der Folge auch sehr
häufig durch
\[
U = (f, \varphi)^k
\]
darstellen. Ist $n < m$, so kann $k$ alle Werthe von 0 bis $n$ annehmen.
Insbesondere ist:
\[
(f, \varphi)^0 = f \cdot \varphi = a_x^m \cdot b_x^n
\]
\[
(f, \varphi)^n = (ab)^n a_x^{m-n}.
\]
Jede höhere Ueberschiebung als die $n^{\text{te}}$ verschwindet alsdann
identisch.

Die Ueberschiebung ist eine auch unsymbolisch darstellbare Form, eine
simultane Covariante von $f$ und $\varphi$, und der rein symbolische
Faltungsprocess liefert sonach hier ein Resultat, das sonst durch rein
rechnerische Operationen erhalten wird. (Siehe Nr. 35.) Aus der Definition
der Ueberschiebung als Resultat der Faltung eines Productes $f \cdot \varphi$
erkennt man unmittelbar die Gültigkeit folgender Relationen:

\weq{$\alpha$}{(c \cdot f, \varphi)^k = c \cdot (f, \varphi)^k, \quad \text{d. h.:}}

\begin{quote}
\glqq Die Ueberschiebung eines Productes $c \cdot f$, wo $c$ eine
Constante, über eine Form $\varphi$ ist gleich dem Producte der
Ueberschiebung $(f, \varphi)^k$ in die Constante $c$.\grqq{}
\end{quote}

\weq{$\beta$}{(f + c_1 \cdot \psi,\ \varphi + c_2 \cdot \chi)^k = (f, \varphi)^k + c_1 \cdot (\psi, \varphi)^k + c_2 (f, \chi)^k + c_1 \cdot c_2 \cdot (\psi, \chi)^k, \quad \text{d. h.:}}

\begin{quote}
\glqq Die Ueberschiebung von linearen Systemen einzelner Formen ist gleich
einem linearen Systeme der Ueberschiebungen dieser Formen.\grqq{}
\end{quote}

\weq{$\gamma$}{(\varphi, f)^k = (ba)^k b_x^{n-k} a_x^{m-k} = (-1)^k \cdot (f, \varphi)^k, \quad \text{d. h.:}}

\begin{quote}
\glqq Die $k^{\text{te}}$ Ueberschiebung ändert, wenn man $f$ mit $\varphi$
vertauscht, ihr Zeichen, je nachdem $k$ gerade oder ungerade.\grqq{}
\end{quote}

Im Interesse späterer Untersuchungen will ich die Faltungen, durch welche
eine Ueberschiebung entsteht, noch auf einem umständlicheren Wege
ausführen. Wir denken uns nämlich einen Augenblick
%CONT>>


% ===== pdf idx 48, printed 34.0 =====
% !PRINTED 34
\weberpage{34}

%<<CONT
die $m$ symbolischen Factoren von $f$, und die $n$ Factoren von $\varphi$
verschieden, etwa:
\[
f = r_{1x} r_{2x} r_{3x} \cdots r_{mx} = a_x^m
\]
\[
\varphi = s_{1x} s_{2x} s_{3x} \cdots s_{nx} = b_x^n.
\]
Die Ueberschiebung $(f, \varphi)^k$ entsteht alsdann, wenn wir das Product
\[
P = f \cdot \varphi = r_{1x} r_{2x} \cdots r_{mx} \times s_{1x} s_{2x} \cdots s_{nx}
\]
$k$ mal falten und nachträglich $r_k$ durch $a$, $s_k$ durch $b$ wieder
ersetzen. Um aber hierbei keine der einzelnen Factoren auszuzeichnen, werden
wir diese $k$ Faltungen auf alle möglichen Arten vornehmen, die Summe aller
so entstehenden symbolischen Producte bilden, und, damit wir nach
Einführung der alten Symbole $a$ und $b$ wieder auf das einzige Product
$(ab)^k a_x^{m-k} b_x^{n-k}$ kommen, diese Summe mit der Anzahl aller
Glieder dividiren. Diese Anzahl ist im Allgemeinen gleich dem Producte
\[
\binom{m}{k} k! \times \binom{n}{k} k! = m(m-1) \cdots (m-k+1) \cdot n(n-1) \cdots (n-k+1),
\]
da jede Variation ohne Wiederholung von $m$ Elementen zur $k^{\text{ten}}$
Klasse mit jeder Variation ohne Wiederholung von $n$ Elementen zu einer
Faltung Veranlassung giebt, wenn wir völlig symmetrisch verfahren wollen.
Man hat demnach:
\[
(f, \varphi)^k = \frac{1}{m(m-1) \cdots (m-k+1)} \cdot \frac{1}{n(n-1) \cdots (n-k+1)}
\cdot \sum (r_{i_1} s_{l_1}) \cdots (r_{i_k} s_{l_k}) \cdot \frac{f \cdot \varphi}{r_{i_1 x} r_{i_2 x} \cdots s_{l_1 x} \cdots s_{l_k x}},
\]
und dieser Ausdruck rechts reducirt sich auf
\[
(f, \varphi)^k = (ab)^k a_x^{m-k} b_x^{n-k},
\]
wenn wir $r_k$ durch $a$, und $s_k$ durch $b$ ersetzen.

33. \textit{Der Omegaprocess und die Ueberschiebung.} Das gleiche
Resultat, das wir hier durch $k$ malige Faltung der beiden Formen $f$ und
$\varphi$ erhielten, entsteht auch durch $k$ malige Anwendung des
Omegaprocesses, wenn wir vorher in einer der beiden Formen $x$ mit $y$
vertauschen, und nach ausgeführter Operation wieder $y$ durch $x$
ersetzen. Man erkennt dies aus den Entwicklungen in Nr. 20, nach welchen
der Omegaprocess angewendet auf das Product
\[
P = a_x^m b_y^n
\]
die Resultate lieferte:
\[
\Omega(P) = (ab)\, a_x^{m-1} b_y^{n-1}
\]
\[
\Omega^2(P) = (ab)^2\, a_x^{m-2} b_y^{n-2} \quad \text{etc.}
\]


% ===== pdf idx 49, printed 35.0 =====
% !PRINTED 35
\weberpage{35}

Wir haben somit auch als weitere Definitionsgleichung:

\weq{III}{(f, \varphi)^k = \{\Omega^k (f_x \cdot \varphi_y)\}_{y=x}.}

34. \textit{Ueberschiebung der Polaren.} Nach der ursprünglichen Definition
wurde die Ueberschiebung in der Weise gebildet, dass man je zwei Factoren
$a_x b_x$ durch den Klammerfactor $(ab)$ ersetzte.

Wir können in diesen Process noch eine scheinbar überflüssige
Zwischenoperation einführen, die wir aber sofort als einen für die
Bildung der Ueberschiebung fundamentalen Process erkennen werden.
Ersetzen wir nämlich in dem Producte
\[
P = a_x^m \cdot b_x^n
\]
jedes der $k$ Factorenpaare $a_x b_x$ zuerst durch $k$ Factorenpaare
$a_y b_y$, und falten wir nun erst $k$ mal nach $y$, so entsteht wiederum
\[
U = (ab)^k a_x^{m-k} b_x^{n-k}.
\]
Die Zwischenoperation war aber nichts anderes als der Polarenprocess,
durch welchen aus $P$ die Form
\[
f_{y^k} \cdot \varphi_{y^k} = a_x^{m-k} a_y^k \cdot b_x^{n-k} b_y^k
\]
geworden war.

Die Ueberschiebung entstand hieraus durch Faltung der Polaren nach $y$.
Daraus erkennen wir den Satz, der namentlich für complicirtere Formen $f$
und $\varphi$ von Wichtigkeit wird:

\begin{quote}
\glqq \gesperrt{Die $k^{\text{te}}$ Ueberschiebung der $k^{\text{ten}}$
Polaren ist gleich der $k^{\text{ten}}$ Ueberschiebung der Formen,}\grqq{}
oder:
\end{quote}

\weq{IV}{U = (f, \varphi)^k = (f_{y^k}, \varphi_{y^k})^k.}

Bei dieser Ueberschiebung der Polaren werden hier die Variabeln $x$ als
constant angesehen und die Ueberschiebung geht so weit als die Variabeln
$y$ gehen, d. h. die Operation liefert eine simultane Invariante der
Polaren in Bezug auf die $y$.

35. \textit{Unsymbolische Darstellung der Ueberschiebung zweier einfacher
Formen $f$ und $\varphi$.} Bezeichnen wir nun die $k^{\text{te}}$ Polare
von $f$ mit
\[
r_y^k = a_x^{m-k} a_y^k,
\]
und ebenso die $k^{\text{te}}$ Polare von $\varphi$ mit
\[
s_y^k = b_x^{n-k} b_y^k,
\]
dann ist also die $k^{\text{te}}$ Ueberschiebung dieser beiden Formen $f$
und $\varphi$ bezeichnet durch:
\[
(rs)^k = (f_{y^k}, \varphi_{y^k})^k = (f, \varphi)^k = a_x^{m-k} b_x^{n-k} (ab)^k.
\]


% ===== pdf idx 50, printed 36.0 =====
% !PRINTED 36
\weberpage{36}

Es ist aber $r_y^k = r$ als Function von $y$ betrachtet unsymbolisch
gegeben durch:
\[
r = \bar r_0 y_1^k + \binom{k}{1} \bar r_1 y_1^{k-1} y_2 + \binom{k}{2} \bar r_2 y_1^{k-2} y_2^2 + \cdots + \bar r_k y_2^k
\]
und ebenso
\[
s = \bar s_k y_2^k + \binom{k}{1} \bar s_{k-1} y_2^{k-1} y_1 + \binom{k}{2} \bar s_{k-2} y_2^{k-2} y_1^2 + \cdots + \bar s_0 y_1^k.
\]
Hierbei sind die unsymbolischen Coefficienten $r_i$, $\bar s_i$ im
Allgemeinen Functionen von $x$. Stellt man nun anderntheils $r$ und $s$
symbolisch dar, setzt also
\[
r = (r_1 y_1 + r_2 y_2)^k
\]
\[
s = (s_1 y_1 + s_2 y_2)^k;
\]
dann ist in symbolischer Bezeichnung die Invariante

\weq{I}{(rs)^k = r_1^k s_2^k - \binom{k}{1} r_1^{k-1} r_2 s_2^{k-1} s_1 + \binom{k}{2} r_1^{k-2} r_2^2 s_2^{k-2} s_1^2 + \cdots \pm r_2^k s_1^k.}

Ersetzt man aber hierin die symbolischen Coefficienten
\[
r_1^k, \quad r_1^{k-1} \cdot r_2, \quad r_2^{k-2} \cdot r_2^2, \quad \ldots % [sic]
\]
\[
s_1^k, \quad s_1^{k-1} \cdot s_2, \quad s_2^{k-2} \cdot s_2^2, \quad \ldots
\]
durch die wirklichen
\[
\bar r_0, \quad \bar r_1, \quad \bar r_2, \quad \ldots
\]
\[
\bar s_0, \quad \bar s_1, \quad \bar s_2, \quad \ldots,
\]
so erhält man als unsymbolische Darstellung der Ueberschiebung
$(f, \varphi)^k$ die Invariante der Polaren:

\weq{II}{(rs)^k = \bar r_0 \bar s_k - \binom{k}{1} \bar r_1 \bar s_{k-1} + \binom{k}{2} \bar r_2 \bar s_{k-2} - + \cdots \pm \bar r_k \bar s_0.}

Dieser Ausdruck ist eine simultane Covariante von $f$ und $\varphi$. Er
stellt direct die bilineare Invariante der beiden Formen dar, sobald
$m = n = k$, d. h. der Grad der beiden Formen mit dem Grade der
Ueberschiebung übereinstimmt.

Nun sind die Coefficienten der Polare nichts anderes als die Grössen
\[
\frac{(m-k)!}{m!} \frac{\partial^k f}{\partial x_1^{k-\varrho} \cdot \partial x_2^\varrho} \quad \text{(vergl. Nr. 14, (8)).}
\]
Ersetzen wir daher dieselben in (II) durch diese Differentialquotienten,
so erhalten wir als Ausdruck der $k^{\text{ten}}$ Ueberschiebung
\[
(f, \varphi)^k = (rs)^k = \frac{(m-k)!}{m!} \frac{(n-k)!}{n!} \left\{ \frac{\partial^k f}{\partial x_1^k} \cdot \frac{\partial^k \varphi}{\partial x_2^k} \right.
\]
\[
\left. - \binom{k}{1} \frac{\partial^k f}{\partial x_1^{k-1} \partial x_2} \cdot \frac{\partial^k \varphi}{\partial x_2^{k-1} \partial x_1} + \cdots \pm \frac{\partial^k f}{\partial x_2^k} \frac{\partial^k \varphi}{\partial x_1^k} \right\},
\]
%CONT>>


% ===== pdf idx 51, printed 37.0 =====
% !PRINTED 37
\weberpage{37}
% running head: "Processe für invariante Bildungen.  37"
%<<CONT
oder in leicht verständlicher kürzerer Schreibweise:
\weq{III}{(f, \varphi)^k = \frac{(m-k)!}{m!}\,\frac{(n-k)!}{n!}\left\{\frac{\partial f}{\partial x_1}\frac{\partial \varphi}{\partial x_2} - \frac{\partial f}{\partial x_2}\cdot\frac{\partial \varphi}{\partial x_1}\right\}^k \cdot{}^{*)}}
Vergleicht man diese Darstellung mit der Darstellung der $k^{\text{ten}}$ Polare von $f=a_x^m$, Nr.\,14 (8)
\[
f_{\cdot y^k} = \frac{(m-k)!}{m!}\left\{\frac{\partial f}{\partial x_1}y_1+\frac{\partial f}{\partial x_2}y_2\right\}^k,
\]
so erkennt man, dass sich der Ueberschiebungsprocess auch folgendermassen formuliren lässt:

\glqq Die $k^{\text{te}}$ Ueberschiebung von $f$ über $\varphi$ wird erhalten, wenn
man in der $k^{\text{ten}}$ Polare von $f$ $y_1$ durch $b_2$ und $y_2$ durch $-b_1$
ersetzt, und die entsprechende $(n-k)^{\text{te}}$ Potenz des Symbols
$b_x$ von $\varphi$ als Factor hinzufügt.\grqq{}

% running head at start of mid band: "36. Beispiele von simultanen Co- und Invarianten."
36. \textit{Beispiele von simultanen Co- und Invarianten.} Ich will nun
eine Reihe von Ueberschiebungen wirklich bilden. Wie schon erwähnt,
erhält man die simultane bilineare Invariante zweier Formen gleichen
Grades unmittelbar durch die Relation (II).

Ist z.\,B. $f=a_x^2=\bar a_0x_1^2+2\bar a_1x_1x_2+\bar a_2x_2^2$
\[
\varphi=b_x^2=\bar b_0x_1^2+2\bar b_1x_1x_2+\bar b_2x_2^2,
\]
so ist
\weq{1}{(f,\varphi)^2=(a_y^2,b_y^2)^2=(ab)^2=\bar a_0\bar b_2-2\bar a_1\bar b_1+\bar a_2\bar b_0.}
Ebenso ist, wenn $f=a_x^3$ und $\varphi=b_x^3$:
\weq{2}{(ab)^3=\bar a_0\bar b_3-3\bar a_1\bar b_2+3\bar a_2\bar b_1-\bar a_3\bar b_0.}
Berechnen wir auch die zweite Ueberschiebung der beiden Formen $f$ und
$\varphi$. Die zweite Polare von $f$ ist unsymbolisch
\[
f_{y^2}=a_y^2a_x=(\bar a_0x_1+\bar a_1x_2)y_1^2+2(\bar a_1x_1+\bar a_2x_2)y_1y_2+(\bar a_2x_1+\bar a_3x_3)y_2^2. % [sic] printed x_3, pattern elsewhere on page reads x_2
\]
Ersetzt man hierin
\[
\begin{aligned}
y_1 &\ \text{durch}\ \ b_2\\
y_2 &\ \text{durch}\ -b_1
\end{aligned}
\]
und multiplicirt mit $b_x$, so erhält man
\[
(ab)^2a_xb_x=(\bar a_0x_1+\bar a_1x_2)b_2^2b_x-2(\bar a_1x_1+\bar a_2x_2)b_1b_2b_x+(\bar a_2x_1+\bar a_3x_3)b_1^2b_x, % [sic] printed x_3, cf. above
\]
oder, wenn man $b_x$ durch $b_1x_1+b_2x_2$ ersetzt:
\[
\begin{aligned}
(ab)^2a_xb_x={}&(\bar a_0x_1+\bar a_1x_2)(b_1b_2^2x_1+b_2^3x_2)-2(\bar a_1x_1+\bar a_2x_2)(b_1^2b_2x_1+b_1b_2^2x_2)\\
&+(\bar a_2x_1+\bar a_3x_2)(b_1^3x_1+b_1^2b_2x_2).
\end{aligned}
\]
Führt man hier für die symbolischen Co\"efficienten $b_1^3, b_1^2, b_2$ etc.\ die
wirklichen ein und ordnet nach den Potenzen der Variabeln $x$, so kommt:

\noindent\rule{4cm}{0.4pt}

{\footnotesize *) Die Gleichung (III) lehrt, dass $(f, \varphi)^k = 0$ eine lineare
Differentialgleichung zur Bestimmung der Form $\varphi$ ist, sobald $f$ als bekannt
angenommen wird.}
%CONT>>


% ===== pdf idx 52, printed 38.0 =====
% !PRINTED 38
\weberpage{38}
% running head: "38   Erster Theil."
%<<CONT
\weq{3}{
\begin{aligned}
(ab)^2a_xb_x={}&(\bar a_0\bar b_2-2\bar a_1\bar b_1+\bar a_2\bar b_0)x_1^2 - (\bar a_1\bar b_2+\bar a_2\bar b_1-\bar a_0\bar b_3+\bar a_3\bar b_0)x_1x_2\\
&+(\bar a_1\bar b_3-2\bar a_2\bar b_2+\bar a_3\bar b_1)x_2^2.
\end{aligned}
}

37. \textit{Beispiele von Co- und Invarianten einer Form.} Bei der
Darstellung des Ueberschiebungsprocesses waren $f$ und $\varphi$
beliebige Formen. Wir können nun als Form $\varphi$ die Form $f$
selbst wieder wählen und erhalten auf diese Weise Co- und Invarianten
einer einzigen Form, durch Ueberschiebung derselben über sich selbst.

So liefert das Beispiel (2) in Nr.\,36 für die cubische Form
\[
f=a_x^3=b_x^3
\]
\[
(ab)^3=\bar a_0\bar a_3-3\bar a_1\bar a_2+3\bar a_2\bar a_1-\bar a_3a_0=0,
\]
d.\,h.\ die dritte Ueberschiebung von $f$ über sich selbst verschwindet,
was a priori aus dem symbolischen Product $(ab)^3$ hervorgeht. (Vergleiche
Nr.\,8.)

Ebenso wird nach Beispiel (3) Nr.\,36
\weq{4}{(ab)^2a_xb_x=2\left[(\bar a_0\bar a_2-\bar a_1^2)x_1^2+(\bar a_0\bar a_3-\bar a_1\bar a_2)x_1x_2+(\bar a_1\bar a_3-\bar a_2^2)x_2^2\right].}

Bezeichnet man diese Covariante von $f$ mit $\Delta_x^2$, also ihre
Co\"efficienten mit $\bar\Delta_0, \bar\Delta_1, \bar\Delta_2$, so kann man von
ihr selbst einmal die erste Ueberschiebung über $f$, dann die zweite über
sich selbst auf folgende Weise bilden. Es ist:
\[
(f, \Delta)=(a\Delta)a_x^2\Delta_x=\frac{1}{2}\cdot\frac{1}{3}\left\{\frac{\partial f}{\partial x_1}\frac{\partial \Delta}{\partial x_2} - \frac{\partial f}{\partial x_2}\frac{\partial \Delta}{\partial x_1}\right\}.
\]

Man erhält somit
\[
\begin{aligned}
(a\Delta)a_x^2\Delta_x={}&(\bar a_0x_1^2+2\bar a_1x_1x_2+\bar a_2x_2^2)\Delta_2\Delta_x-(\bar a_1x_1^2+2\bar a_2x_1x_2+\bar a_3x_2^2)\Delta_1\Delta_x\\
={}&(\bar a_0x_1^2+2\bar a_1x_1x_2+\bar a_2x_2^2)(\Delta_1x_1+\Delta_2x_2)\\
&-(\bar a_1x_1^2+2\bar a_2x_1x_2+\bar a_3x_2^2)(\bar\Delta_0x_1+\bar\Delta_1x_2).
\end{aligned}
\]

Ersetzt man hierin die Co\"efficienten von $\Delta$ durch ihre Werthe,
also:
\[
\begin{aligned}
\bar\Delta_0 &\ \text{durch}\ \ 2(\bar a_0\bar a_2-\bar a_1^2)\\
\bar\Delta_1 &\ \text{durch}\ \ (\bar a_0\bar a_3-\bar a_1\bar a_2)\\
\bar\Delta_2 &\ \text{durch}\ \ 2(\bar a_1\bar a_3-a_2^2),
\end{aligned}
\]
und ordnet nach fallenden Potenzen von $x_1$, so erhält man:
\weq{5}{
\begin{aligned}
(a\Delta)a_x^2\Delta_x={}&(\bar a_0^2\bar a_3-3\bar a_0\bar a_1\bar a_2+2\bar a_1^3)x_1^3+3(\bar a_0\bar a_1\bar a_3-2\bar a_0\bar a_2^2+\bar a_1^2\bar a_2)x_1^2x_2\\
&-3(\bar a_0\bar a_2\bar a_3-2\bar a_1^2\bar a_3+\bar a_1\bar a_2^2)x_1x_2^2-(\bar a_0\bar a_3^2-3\bar a_1\bar a_2\bar a_3+2\bar a_2^3)x_2^3.
\end{aligned}
}

Endlich liefert die zweite Ueberschiebung der quadratischen Form
$\Delta$ über sich selbst
%CONT>>


% ===== pdf idx 53, printed 39.0 =====
% !PRINTED 39
\weberpage{39}
% running head: "Processe für invariante Bildungen.  39"
%<<CONT
\[
(\Delta, \Delta)^2=(\Delta \Delta^{(1)})^2=(\Delta_1^2\Delta_2^{(1)2}-2\Delta_1\Delta_2^{(1)}\Delta_2\Delta_1^{(1)}+\Delta_2^2\Delta_1^{(1)2})=2(\bar\Delta_0\bar\Delta_2-\bar\Delta_1^2)
\]
\[
=2\left\{4(\bar a_0\bar a_2-\bar a_1^2)(\bar a_1\bar a_3-\bar a_2^2)-(\bar a_0\bar a_3-\bar a_1\bar a_2)^2\right\},
\]
oder
\weq{6}{(\Delta, \Delta)^2=2\left\{6a_0a_1a_2a_3-4a_2^3a_0-4a_1^3a_3+3a_1^2a_2^2-a_0^2a_3^2\right\}.}

Wir haben somit durch einfache Ueberschiebungsprocesse eine
Invariante $4^{\text{ten}}$ Grades in den Co\"efficienten, eine
Covariante $3^{\text{ten}}$ Grades in Co\"efficienten und Variabeln,
und ferner eine Covariante $2^{\text{ten}}$ Grades der cubischen Form
$f$ kennen gelernt. Wir werden später sehen, dass diese drei Formen
überhaupt die einzigen \gesperrt{selbstständigen} Covarianten von $f$
sind, d.\,h.\ dass sich alle übrigen rational und ganz durch $f$ und
diese drei ausdrücken lassen.

38. \textit{Ueberschiebungen von Producten von Formen.} Wir gehen nun
zum complicirteren Falle des Ueberschiebungsprocesses über, in welchem
jede der beiden Formen, die zur Ueberschiebung gelangen sollen, ein
Product einfacher Formen ist. Die Ueberschiebung ist dann nicht mehr
wie im vorigen Falle durch ein \gesperrt{einziges} symbolisches
Product dargestellt, sondern durch eine ganze Reihe derselben, die wir
nunmehr ermitteln wollen.

Sei $F$ das Product der Formen
\[
f_1=a_x^m,\quad f_2=b_x^n,\quad\cdots f_k=k_x^p;
\]
ferner $\Phi$ das Product der Formen
\[
\varphi_1=\alpha_x^\mu,\quad \varphi_2=\beta_x^\nu,\quad\cdots \varphi_\varkappa=\varkappa_x^\pi.
\]

Für die linearen symbolischen Factoren der Formen $f_i$ führen wir der
Reihe nach die Factoren $r_{ix}$ ein; ihre Anzahl $s$ ist also gleich
der Summe der Grade aller Formen, also
\[
s=m+n+\cdots+p.
\]
Die linearen Factoren der Formen $\varphi$ ersetzen wir ebenso durch
die Factoren $\varrho_{ix}$; die Anzahl $\sigma$ derselben ist dann
\[
\sigma=\mu+\nu+\cdots+\pi.
\]
Die $k^{\text{te}}$ Ueberschiebung der beiden Formen $F$ und $\Phi$
wird dann erhalten, wenn wir das Product
\[
\Pi=r_{1x}r_{2x}\cdots r_{sx}\cdot\varrho_{1x}\varrho_{2x}\cdots\varrho_{\sigma x}
\]
auf alle möglichen Arten in der Weise $k$ mal falten, dass immer
andere Variationen von je $k$ Factoren $r_{ix}$ mit immer neuen
Variationen $k$ Factoren $\varrho_{\lambda x}$ zu $k$ Klammerfactoren
vereinigt werden. Die Summe aller so entstehenden symbolischen
Producte dividirt durch ihre Anzahl gibt die Ueberschiebung von $F$
über $\Phi$. Diese Anzahl ist daher
%CONT>>


% ===== pdf idx 54, printed 40.0 =====
% !PRINTED 40
\weberpage{40}
% running head: "40   Erster Theil."
%<<CONT
gleich dem Producte der Anzahl der Variationen von $s$ Elementen zur
$k^{\text{ten}}$ Klasse in die Anzahl der Variationen von $\sigma$
Elementen zur $k^{\text{ten}}$ Klasse. Es ist also die Ueberschiebung
definirt durch die Gleichung:
\weq{I}{
\begin{aligned}
(F, \Phi)^k={}&\frac{(s-k)!}{s!}\cdot\frac{(\sigma-k)!^{*)}}{\sigma!}\\
&\cdot\sum(r_{i_1}\varrho_{\lambda_1})(r_{i_2}\varrho_{\lambda_2})\cdots(r_{i_k}\varrho_{\lambda_k})r_{i_{k+1}x}\cdots r_{i_sx}\cdot\varrho_{\lambda_{k+1}x}\cdots\varrho_{\lambda_\sigma x}.
\end{aligned}
}

Ersetzt man in dieser Summe
\[
\begin{aligned}
r_{1x},\ r_{2x},\ \cdots r_{mx} &\ \text{durch}\ \ a_x\\
r_{m+1,x},\ \cdots r_{m+n,x} &\ \text{durch}\ \ b_x\\
&\cdot\ \cdot\ \cdot\ \cdot\ \cdot\ \cdot\ \cdot\ \cdot\ \cdot\ \cdot\\
\varrho_{1x},\ \varrho_{2x},\ \cdots \varrho_{\mu x} &\ \text{durch}\ \ \alpha_x\\
\varrho_{\mu+1,x},\ \cdots \varrho_{\mu+\nu,x} &\ \text{durch}\ \ \beta_x,\ \text{etc.}
\end{aligned}
\]
und fasst alle gleichen Glieder zusammen, so erhält man die
Ueberschiebung ausgedrückt durch die ursprünglichen Symbole von $F$
und $\Phi$. Die Summe aller Co\"efficienten dieser Summe ist gleich
Eins, wie aus der Definition der Ueberschiebung direct hervorgeht.
(Vergl.\ auch Nr.\,32.)

\textit{Anmerkung.} Es mag noch bemerkt werden, dass bei wirklicher
Berechnung solcher Ueberschiebungen nicht jede Variation der Elemente
$r_k$ mit jeder Variation der Elemente $\varrho_k$ verbunden zu werden
braucht. Es genügt vielmehr, jede \gesperrt{Combination} von $k$
Factoren $r_{kx}$ mit jeder Variation von $k$ Factoren $\varrho_{kx}$
zu falten, da bereits hierdurch alle möglichen Arten von wirklich
verschiedenen Verbindungen geliefert werden. Die Zahl der Glieder ist
alsdann gleich der oben erwähnten dividirt durch $k!$. Wenn ich in der
oben gegebenen Entwicklung mehr Faltungen als nöthig vornahm, so
geschah das einmal um die Symmetrie nicht zu verletzen, und
andernteils um die Uebereinstimmung des numerischen Co\"efficienten
\[
\frac{(s-k)!}{s!}\cdot\frac{(\sigma-k)!}{\sigma!}
\]
mit dem Nr.\,35 (III) gegebenen hervorzuheben.

39. \textit{Beispiel 1.} Ich will diese Darstellung des
Ueberschiebungsprocesses an einem einfachen Beispiele deduciren, das
bereits in Nr.\,37 behandelt wurde. Wir hatten dort die zweite
Ueberschiebung der quadratischen Covariante einer cubischen Form über
sich selbst berechnet, und zu dem Zwecke für dieselbe das neue Symbol
$\Delta$ eingeführt. Bilden wir nun direct die zweite Ueberschiebung
dieser Form
\[
\Delta=(ab)^2a_xb_x=(cd)^2c_xd_x=\text{etc.}
\]
über sich selbst, so kommt nach den vorhergehenden Entwicklungen
\[
\left((ab)^2a_xb_x,\ (cd)^2c_xd_x\right)^2=(ab)^2(cd)^2(a_xb_x,\ c_xd_x)^2
\]
\[
=\frac{1}{2}\left\{(ab)^2(cd)^2(ac)(bd)+(ab)^2(cd)^2(ad)(bc)\right\}.
\]


% ===== pdf idx 55, printed 41.0 =====
% !PRINTED 41
\weberpage{41}
Beide Glieder sind einander gleich; denn sie gehen in einander über
durch Vertauschung der gleichwerthigen Symbole $a$ und $b$, oder $c$
und $d$. Ihre Summe kann also durch das doppelte erste Glied ersetzt
werden und somit ist diese Invariante $(\Delta,\ \Delta)^2$ auch dargestellt durch
\[
(\Delta,\ \Delta)^2 = (ab)^2(cd)^2(ac)(bd) = (ab)^2(cd)^2(ad)(bc).
\]
\textit{Beispiel 2.} Wir stellen uns ferner die Aufgabe, eine durch ein
symbolisches Product dargestellte Covariante oder Invariante einer
Form $f$ in den Wurzeldifferenzen dieser Form zu berechnen. Die
Methode will ich an einem Beispiele entwickeln. Es sei gegeben die
Form $f = a_x^4$; sie besitzt eine Invariante $i = (ab)^4$. Um sie in den
Wurzeln $\alpha_1$, $\alpha_2$, $\alpha_3$, $\alpha_4$ von $f$ darzustellen, bilden wir zunächst diese
Invariante durch vierte Ueberschiebung von $a_x^4$ mit:
\[
f = \alpha_{1x}\ \alpha_{2x}\ \alpha_{3x}\ \alpha_{4x}
\]
und erhalten nach Nr. 38 (I):
\weq{1}{i = (ab)^4 = (a\alpha_1)(a\alpha_2)(a\alpha_3)(a\alpha_4);}
auf diese Weise ist bereits ein Symbol $b$ durch die Wurzeln $\alpha_1\alpha_2\alpha_3\alpha_4$
ersetzt. Es handelt sich also noch darum, auch für das Symbol $a$ die
Wurzeln $\alpha_i$ einzuführen. Wir ersetzen zu dem Zwecke in (1) $a$ durch
$z$ und überschieben die Polare $b_z^4 = \alpha_{1z}\alpha_{2z}\alpha_{3z}\alpha_{4z}$ viermal über die
Polare $a_z^4 = \alpha_{1z}\alpha_{2z}\alpha_{3z}\alpha_{4z}$. So erhalten wir nach Nr. 38 (I) und An\-merkung:
\[
(a_z^4,\ b_z^4)^4 = (ab)^4 = (\alpha_{1z}\cdot\alpha_{2z}\cdot\alpha_{3z}\cdot\alpha_{4z},\ \alpha_{1z}\alpha_{2z}\alpha_{3z}\alpha_{4z})^4
\]
\[
= \frac{1}{4!}\sum (\alpha_{i_1}\alpha_1)(\alpha_{i_2}\alpha_2)(\alpha_{i_3}\alpha_3)(\alpha_{i_4}\alpha_4),
\]
oder mit Vernachlässigung jener Glieder der Summe, welche wegen
Factoren $(\alpha_{i_1}\alpha_{i_1}) = 0$ identisch verschwinden:
\[
(ab)^4 = i = \frac{1}{4!}\left\{(\alpha_1\alpha_2)(\alpha_3\alpha_4) + (\alpha_1\alpha_3)(\alpha_2\alpha_4) + (\alpha_1\alpha_4)(\alpha_2\alpha_3)\right\}^2.
\]
Ganz in derselben Weise verfährt man auch allgemein. Man ersetzt
zunächst ein Symbol der Covariante durch $y$, und bildet die
$n^{\text{te}}$ Ueberschiebung dieser Polare mit der Polare $f_{y^n} = \alpha_{1y}\alpha_{2y}\cdots\alpha_{ny}$.
So erhält man ein Aggregat von Gliedern, welche ein Symbol
weniger enthalten. In jedem dieser Glieder ersetzt man wiederum ein
Symbol durch $z$, und überschiebt dasselbe abermals $n$ mal über
$f_{z^n} = \alpha_{1z}\alpha_{2z}\cdots\alpha_{nz}$ etc. Dadurch ist die Zahl der Symbole um 2
reducirt, und so setzt man den Process fort.] % [sic] stray closing bracket as printed; no opening bracket visible on this page

40. \textit{Darstellung der Ueberschiebung durch Ueberschiebung von
Polaren. Erster Fall.} Bei Berechnung einer Ueberschiebung, welche
an zwei Producten von Formen ausgeführt werden soll, kann man
sich nun
%CONT>>


% ===== pdf idx 56, printed 42.0 =====
% !PRINTED 42
\weberpage{42}
%<<CONT
aber auch wieder des bereits in Nr. 34 bewiesenen Satzes bedienen.
Die Sätze über Eigenschaften der Polaren lassen sich alsdann auf die
Ueberschiebung übertragen. Ich will mich indess bei den folgenden
Entwicklungen auf den Fall beschränken, in dem $F$ und $\Phi$ nur als
Producte von je zwei Formen vorausgesetzt werden, da die Verallge\-
meinerung nur eine öftere Anwendung der Operationen des speciellen
Falles erheischt, und keine neuen Ueberlegungen nöthig macht. Ueber\-
dies sei zunächst vorausgesetzt, dass der Grad der beiden Producte
in $x$ gleich dem Grade $k$ der Ueberschiebung sei, dass also, wenn
\[
F = f\cdot\varphi = a_x^m\cdot b_x^n
\]
\[
\Phi = \psi\cdot\chi = \alpha_x^\mu\cdot\beta_x^\nu
\]
für den Exponenten der Ueberschiebung
\[
U = (F,\ \Phi)^k
\]
die Beziehung bestehe
\[
k = m+n = \mu+\nu.
\]
Unter dieser Voraussetzung besteht nun aber die $k^{\text{te}}$ Polare von
$F$ aus dem einzigen Gliede
\[
F_{y^k} = a_y^m\cdot b_y^n
\]
und ebenso
\[
\Phi_{y^k} = \alpha_y^\mu\cdot\beta_y^\nu.
\]
Faltet man das Product $F_{y^k}\cdot \Phi_{y^k}$ auf alle möglichen Arten, etwa
indem man die linearen Factoren der einen Polare einen Moment durch
$r_{ix}$, die der andern durch $\varrho_{kx}$ ersetzt und nun die einzige Combination
der $r_{ix}$ mit allen Permutationen der $\varrho_{ix}$ durch Faltung verbindet, so
bildet die Summe aller dieser Glieder, dividirt durch ihre Anzahl, die
$k^{\text{te}}$ Ueberschiebung von $F$ über $\Phi$. Ersetzt man wieder $r_{ix}$ und $\varrho_{ix}$
in geeigneter Weise durch die ursprünglichen Factoren $a_x$ $b_x$ $\alpha_x$ $\beta_x$,
fasst hierauf alle gleichen Glieder zu einem einzigen Gliede $L_s$ zu\-
sammen, so ist die Ueberschiebung dargestellt durch
\[
(F,\ \Phi)^k = \sum q_s L_s,
\]
wo $L_s$ von der Form $(a\alpha)^{\nu_1}(b\beta)^{\nu_2}(a\beta)^{\nu_3}(b\beta)^{\nu_4}$, $\nu_1+\nu_2+\nu_3+\nu_4=k$,
und $q_s$ constante Coefficienten sind, deren Summe gleich Eins ist.

41. \textit{Zweiter Fall.} Ist nun der Grad des Productes $F$, und der
des Productes $\Phi$ grösser als der der Ueberschiebung, oder höchstens
eines der beiden Producte von gleichem Grade wie die Ueberschiebung,
so ist die $k^{\text{te}}$ Polare von $F$ (vergl. Nr. 22)
\[
F_{y^k} = \binom{m}{0}\binom{n}{k} a_x^m b_x^{n-k} b_y^k + \binom{m}{1}\binom{n}{k-1} a_x^{m-1} a_y b_x^{n-k+1} b_y^{k-1} + \cdots = \sum c_i G_i
\]
und ebenso
%CONT>>


% ===== pdf idx 57, printed 43.0 =====
% !PRINTED 43
\weberpage{43}
%<<CONT
\[
\Phi_{y^k} = \binom{\mu}{0}\binom{\nu}{k} \alpha_x^\mu \beta_x^{\nu-k} \beta_y^k + \binom{\mu}{1}\binom{\nu}{k-1} \alpha_x^{\mu-1} \alpha_y \beta_x^{\nu-k+1} \beta_y^{k-1} + \cdots = \sum p_\varrho H_\varrho.
\]
Die Ueberschiebung erhält man, indem man das Product
\[
F_{y^k}\cdot \Phi_{y^k} = \sum c_i G_i \cdot \sum p_\varrho H_\varrho
\]
auf alle möglichen Arten nach den in $y$ linearen Factoren faltet. Sie
ist also dargestellt durch die Summe
\weq{I}{(F,\ \Phi)^k = \sum\sum c_i\cdot p_\varrho (G_i,\ H_\varrho)^k,}
wo nun das Glied $(G_i,\ H_\varrho)^k$ entsteht, indem man auf irgend zwei
Glieder der $k^{\text{ten}}$ Polaren $k$ mal den Faltungsprocess anwendet. Die Ueber\-
schiebung der beiden Formen $F$ und $\Phi$ ist also zurückgeführt auf
eine Summe von Ueberschiebungen solcher Formen $G_i$ und $H_\varrho$, deren
Grad in $y$ genau mit dem Grade der Ueberschiebung übereinstimmt.
Jede solche Ueberschiebung ist nach Nr. 40 dargestellt durch
\[
(G_i,\ H_\varrho)^k = \sum q_s L_s.
\]
Trägt man also dieselben in (I) ein, so erhält man:
\[
(F,\ \Phi)^k = \sum\sum c_i p_\varrho \sum q_s L_s = \sum\sum\sum c_i p_\varrho q_s L_s
\]
\[
= \sum \lambda_{i\varrho s}(a\alpha)^{\nu_1}(b\beta)^{\nu_2}(a\beta)^{\nu_3}(b\beta)^{\nu_4}\, a_x^{m-(\nu_1+\nu_3)}\, b_x^{n-(\nu_2+\nu_4)}\, \alpha_x^{\mu-(\nu_1+\nu_4)}\, \beta_x^{\nu-(\nu_2+\nu_3)}.
\]
Hierbei ist die Summe der Coefficienten $\lambda_{i\varrho s}$ gleich Eins. Denn
es ist
\[
\sum c_i = 1, \qquad \sum p_\varrho = 1 \quad \text{nach Nr. 23}
\]
und ebenso
\[
\sum q_s = 1 \quad \text{nach Nr. 40,}
\]
folglich auch
\[
\sum c_i \cdot \sum p_\varrho \cdot \sum q_s = \sum \lambda_{i\varrho s} = 1.
\]

42. \textit{Benachbarte Glieder einer Ueberschiebung.} Wie bei den Po\-
laren so spricht man auch hier von benachbarten Gliedern. Doch ist
die Mannigfaltigkeit der Nachbarschaft nun eine grössere geworden.
So können zwei Glieder benachbart sein, wenn in dem einen an
Stelle von
\weq{1}{a_x (b\alpha)\ \text{sich befindet}\ b_x (a\alpha),}
oder an Stelle von
\weq{2}{\alpha_x (a\beta)\ \text{sich befindet}\ \beta_x (a\alpha),}
während alle übrigen Factoren übereinstimmen. Im ersten Falle hat
die Differenz der betreffenden Glieder den Factor
%CONT>>


% ===== pdf idx 58, printed 44.0 =====
% !PRINTED 44
\weberpage{44}
%<<CONT
\[
a_x(b\alpha) - b_x(a\alpha) = \mathfrak{a}_x(ab), % ?? uncertain: unusual letterform before \textquotedbl{}_x(ab)\textquotedbl{} in the scan, possibly a print defect rather than intentional Fraktur
\]
im zweiten Falle den Factor
\[
\alpha_x(a\beta) - \beta_x(a\alpha) = a_x(\alpha\beta),
\]
wie aus dem Identitätssatz Nr. 11 hervorgeht.
Die Differenz zweier in dieser Weise benachbarter Glieder hat
demnach immer einen der beiden Factoren
\[
(ab)\quad\text{oder}\quad(\alpha\beta).
\]
Die Nachbarschaft zweier Glieder kann aber auch lediglich durch
Klammerfactoren bedingt sein. Es kann die Stelle von
\weq{3}{(a\alpha)(b\beta)\ \text{das Product}\ (a\beta)(b\alpha)}
einnehmen; dann hat nach dem Identitätssatze die Differenz dieser
beiden Glieder
\[
G_1 - G_2 = \bar{G}\left\{(a\alpha)(b\beta) - (a\beta)(b\alpha)\right\}
\]
\[
= \bar{G}(ab)(\alpha\beta)
\]
den Factor $(ab)(\alpha\beta)$.

\glqq Diese Differenzen haben also in allen drei Fällen mindestens
einen Klammerfactor, der durch Faltung zwischen Formen $F$
allein, oder $\Phi$ allein entsteht, und nicht durch gegenseitige
Faltung von $F$ mit $\Phi$.\grqq{}
[So ist z.\,B. in der schon mehrfach berechneten Ueberschiebung
\[
U = \left((ab)^2 a_x b_x,\ (cd)^2 c_x d_x\right)^2
\]
\[
= \frac{(ab)^2(cd)^2}{2}\left\{(ac)(bd) + (ad)(bc)\right\},
\]
die Differenz der beiden Glieder dargestellt durch
\[
G_1 - G_2 = (ab)^2(cd)^2\left\{(ac)(bd) - (ad)(bc)\right\}
\]
\[
= (ab)^3(cd)^3.
\]
Die rechte Seite verschwindet aus bekannten Gründen, wenn $a$,
$b$, $c$, $d$ Symbole der cubischen Form $f = a_x^3$ sind. Dann ist also
\[
G_1 = G_2,
\]
folglich
\[
U = \frac{1}{2}\cdot 2 G_1 = (ab)^2(cd)^2(ac)(bd),
\]
was wir schon auf anderem Wege gezeigt haben. (Vergl. Nr. 39.)
Die eben besprochene Eigenschaft der Differenz benachbarter
Glieder kann man auch für die Differenz zweier beliebiger Glieder
$L_i$ und $L_p$ nachweisen. Denn es ist
\[
L_i - L_p = (L_i - L_{i+1}) + (L_{i+1} - L_{i+2}) + \cdots + (L_{p-1} - L_p).
\]


% ===== pdf idx 59, printed 45.0 =====
% !PRINTED 45
\weberpage{45}
Jede Differenz rechts besitzt den Factor $(ab)$ oder $(\alpha\beta)$, also ist
die Differenz links durch ein Aggregat von Gliedern darstellbar, welche
entweder den Factor $(ab)$ oder den Factor $(\alpha\beta)$ oder beide besitzen.

43. \textit{Anwendungen.} Man kann von diesen Sätzen interessante Anwendungen
machen, die für uns in der Folge von Wichtigkeit sein werden. Denken wir uns
eine Covariante $A_i$ der Form $f_n=a_x^n$, so können wir uns genau dieselbe
Covariante $B_i$ der Form $\varphi=\alpha_x^m$ bilden, wenn $m\geqq n$, indem
wir an der Potenz $\varphi^\varrho$ die nämlichen Faltungen vornehmen, die wir
an $f^\varrho$ machten, um die Form $A_i$ zu erhalten. Nennen wir der
Einfachheit halber $A_i$ das Modell, $B_i$ das Bild. Ist insbesondere
\[
\varphi=(ab)^{2\mu}a_x^{n-2\mu}b_x^{n-2\mu}=(aa_1)^{2\mu}a_x^{n-\mu}a_{1x}^{n-\mu}, % ?? uncertain: faint diacritic marks over several exponents in this line; read by degree-balance as n-2\mu (left) and n-\mu (right)
\]
also eine gerade Ueberschiebung der Form $f$ über sich selbst, so können wir
mit Hilfe der Eigenschaft benachbarter Glieder für Bilder $B_i$, welche aus
dieser Form nach den Modellen $A_i$ gestaltet werden, eine Reihe von Sätzen
ableiten.

In diesem Falle nämlich kann man aus der Potenz
\[
\varphi^\varrho=(aa_1)^{2\mu}a_x^{n-\mu}a_{1x}^{n-\mu}\cdot(bb_1)^{2\mu}b_x^{n-\mu}b_{1x}^{n-\mu}\cdots(rr_1)^{2\mu}r_x^{n-\mu}r_{1x}^{n-\mu} % ?? uncertain: same faint diacritics over the exponents as above
\]
eine Menge Bilder $B_i$ nach dem Modelle $A_i$ herstellen, je nachdem man die
durch $A_i$ dictirten Faltungen der Potenz
\[
f^\varrho=a_x^n\cdot b_x^n\cdot c_x^n\cdots r_x^n
\]
in $\varphi^\varrho$ ausführt, d.\,h. je nachdem man in $\varphi^\varrho$ die
entsprechenden Faltungen an $a_x^{n-\mu}b_x^{n-\mu}$, oder an
$a_x^{n-\mu}b_{1x}^{n-\mu}$, oder an $a_{1x}^{n-\mu}b_x^{n-\mu}$, oder
endlich an $a_{1x}^{n-\mu}b_{1x}^{n-\mu}$ vornimmt.

Wir nennen dann wieder zwei Bilder $B_i$ und $B_i'$ benachbarte Bilder, wenn
sie dadurch aus einander hervorgehen, dass man in zwei symbolischen Factoren
derselben $a$ mit $a_1$ vertauscht. Ihre Differenz hat dann immer um einen
Factor $(aa_1)=(bb_1)=(cc_1)$ mehr als das einzelne Bild.

Ist z.\,B.
\[
\begin{aligned}
B_i&=(ab)(a_1b_1)G_i\\
B_i'&=(a_1b)(ab_1)G_i,
\end{aligned}
\]
dann ist:
\[
\begin{aligned}
B_i-B_i'&=G_i\{(ab)(a_1b_1)-(a_1b)(ab_1)\}\\
&=G_i(aa_1)(bb_1);
\end{aligned}
\]
oder ist:
\[
\begin{aligned}
B_i&=a_x(a_1b)G_i\\
B_i'&=a_{1x}(ab)G_i,
\end{aligned}
\]
so ist
\[
B_i-B_i'=G_i\{a_x(a_1b)-a_{1x}(ab)\}=G_i(aa_1)b_x.
\]


% ===== pdf idx 60, printed 46.0 =====
% !PRINTED 46
\weberpage{46}
Genau wie vorhin erkennt man daher, dass auch die Differenz zwei beliebiger
Bilder $B_i$ und $\overline{B}_i$ sich durch ein Aggregat von Gliedern
darstellen lässt, deren jedes einen Factor $(aa_1)$ mehr besitzt als das
einzelne Bild. Da nun ein Bild mindestens $(aa_1)^{2\mu}$ bereits zum Factor
hat, so erkennt man die Richtigkeit der Behauptung:

\glqq Die Differenz zweier Bilder $B_i$ ist gleich einem Aggregat von
Gliedern, welche den Klammerfactor $(ab)^{2\mu+1}$ oder wegen Nr.\,12
$(ab)^{2\mu+2}$ besitzen.\grqq{}

Man kann auch sagen:

\glqq Die Differenz zweier Bilder ist null, bis auf Glieder mit dem Factor
$(ab)^{2\mu+2}$.\grqq{}

Wir werden einen derartigen aus Klammerfactoren zusammengesetzten
charakteristischen Factor Modul nennen, und demnach diesen Satz in die
algebraische Form kleiden:
\weq{1}{\overline{B}_i-B_i\equiv 0\ \mathrm{mod}\ (ab)^{2\mu+2}.}
Hat man nun drei symbolische Producte $A$, so dass
\[
\begin{aligned}
A_1&=a_x(bc)\psi\\
A_2&=b_x(ca)\psi\\
A_3&=c_x(ab)\psi,
\end{aligned}
\]
dann verschwindet wegen des Identitätssatzes ihre Summe identisch. Das
Gleiche gilt für die erste Gattung der oben erwähnten vier Bilder,
\[
\begin{aligned}
B_1&=a_x(bc)\psi'\\
B_2&=b_x(ca)\psi'\\
B_3&=c_x(ab)\psi',
\end{aligned}
\]
die ja neben andern Klammerfactoren genau dieselben enthalten wie die
Modelle $A_i$.

Nimmt man demnach drei beliebige Bilder $\overline{B}_1\ \overline{B}_2\
\overline{B}_3$, so ist wegen Gleichung (1)
\[
\sum_1^3\overline{B}_i-\sum_1^3 B_i\equiv 0\ \mathrm{mod}\ (ab)^{2\mu+2},
\]
weil aber $\sum B_i=0$, so erhält man
\weq{2}{\sum_1^3\overline{B}_i\equiv 0\ \mathrm{mod}\ (ab)^{2\mu+2}.}
Dieses Resultat lässt uns aber einen sehr wichtigen Satz beweisen:

\glqq Ist $F(A_i)$ eine ganze Function von symbolischen Producten $A_i$ der
Art, dass weder ihre Glieder einzeln null sind, noch
%CONT>>


% ===== pdf idx 61, printed 47.0 =====
% !PRINTED 47
\weberpage{47}
%<<CONT
sich gegenseitig aufheben; so ist, wenn diese Function $F$ gleichwohl
identisch verschwindet, die entsprechende Function von Bildern, also von
symbolischen Producten $B_i$ der Form $\varphi=(f,f)^{2\mu}$
\[
F(\overline{B}_i)\equiv 0\ \mathrm{mod}\ (ab)^{2\mu+2}.
\]
\grqq{} Denn man kann setzen (vergl.\ \S\,1 Nr.\,12, und \S\,10 Nr.\,117)
\[
F(A_i)=0=\sum\{(bc)a_x+(ca)b_x+(ab)c_x\}\varphi
\]
und demnach für die genau entsprechenden Bilder (erster Gattung)
\[
F(B_i)=\sum\{(bc)a_x+(ca)b_x+(ab)c_x\}\varphi'=0.
\]
Sind alsdann $\overline{B}_i$ beliebige Bilder, so ist
\[
F(\overline{B}_i)-F(B_i)\equiv 0\ \mathrm{mod}\ (ab)^{2\mu+2};
\]
also weil der Subtrahend links verschwindet
\weq{3}{F(\overline{B}_i)\equiv 0\ \mathrm{mod}\ (ab)^{2\mu+2}.}
Wir werden später von diesem Satze eine hervorragende Anwendung machen.

44. \textit{Differenz der Ueberschiebung selbst mit einem ihrer Glieder.}
Die Eigenschaft, die wir für die Differenz zweier beliebiger Glieder einer
Ueberschiebung nachgewiesen haben, besitzt auch die Differenz zwischen einem
Gliede und der ganzen Ueberschiebung. Denn da in der Ueberschiebung die
Summe aller Coefficienten Eins ist, so folgt durch Subtraction der beiden
Gleichungen
\[
U=\sum\lambda_{i\varrho s}L_s
\]
\[
L_q=L_q\sum\lambda_{i\varrho s}
\]
die Gleichung
\[
L_q=U+\sum_{i,\varrho,s}\lambda_{i\varrho s}(L_q-L_s)
\]
oder, wenn wir links in den Differenzen die Factoren $(ab)$ resp.
$(\alpha\beta)$ absondern:
\weq{I}{L_q=U+(ab)^{\varepsilon_1}(\alpha\beta)^{\varepsilon_2}\sum\lambda_{i\varrho s}P_s,\quad (\varepsilon_1\ \mathrm{u.}\ \varepsilon_2=0\ \mathrm{oder}\ 1)\ \varepsilon_1\geqq\varepsilon_2=0.}
Nun sind die symbolischen Producte $(ab)^{\varepsilon_1}(\alpha\beta)^{\varepsilon_2}P_s$
Glieder, welche irgend einen Klammerfactor, wie ihn die Ueberschiebung $U$
involviren würde, weniger besitzen, dafür aber mindestens einen mehr, welcher
dieser Ueberschiebung $U$ fremd ist, insofern er nur durch Faltung der Formen
von $F$ unter sich, oder der Formen $\Phi$ unter sich entstehen kann. Daraus
lernen wir zweierlei: Erstens, diese Glieder in den
%CONT>>


% ===== pdf idx 62, printed 48.0 =====
% !PRINTED 48
\weberpage{48}
%<<CONT
Summen rechts gehören Ueberschiebungen an, die einen um mindestens Eine
Einheit niedrigeren Grad $k-1$ haben als die Ueberschiebung $U$; zweitens,
diese Ueberschiebungen beziehen sich auf Formen niedrigeren Grades, welche,
sei es aus $F$ oder sei es aus $\Phi$, durch Faltung entstanden sind.

Wir gewinnen also zunächst den Satz:

\glqq Jedes Glied einer Ueberschiebung lässt sich darstellen durch die
Ueberschiebung selbst plus einer Summe von Gliedern, welche niedrigeren
Ueberschiebungen von Formen niedrigeren Grades angehören.\grqq{}

Ersetzen wir nun vermöge dieses Satzes abermals die Glieder
\[
(ab)^{\varepsilon_1}(\alpha\beta)^{\varepsilon_2}P_i
\]
durch ihre Ueberschiebungen und fahren in diesem Processe so lange fort, bis
wir zu nullten Ueberschiebungen von Formen gelangen, die aus $F$ oder $\Phi$
durch vielfache Faltung entstanden sind, so werden wir endlich zu dem Satze
geführt:

\glqq Jedes Glied einer Ueberschiebung lässt sich darstellen durch eine
Summe von Ueberschiebungen,\grqq{} oder
\weq{II}{L=U+\sum c_i^{(1)}U_i^{(1)}+\sum c_i^{(2)}U_i^{(2)}+\cdots+\sum c_i^{(k)}U_i^{(k)}.}
Dabei sind die Grössen $c_i^{(\varrho)}$ Constante, die Grössen
$U_i^{(\varrho)}$ der Reihe nach Ueberschiebungen $(k-1)^{\mathrm{ten}}$,
$(k-2)^{\mathrm{ten}}$ bis nullten Grades von Formen, die aus $F$ oder
$\Phi$ durch einmalige, zweimalige resp. $k$ malige Faltung entstanden sind.

Diese beiden Sätze (I) und (II) entsprechen den analogen Sätzen (I) und (II)
Nr.\,25 über Polaren.

45. \textit{Jedes symbolische Product lässt sich auf eine Summe einfacher
Ueberschiebungen zweier Formen reduciren.} Der letzte Satz kann noch
erweitert werden, indem wir ihn auf ein beliebiges symbolisches Product,
ohne Rücksicht darauf, ob ihm die Bedeutung einer wirklichen binären Form
zukommt oder nicht, ausdehnen. Das Princip, das uns die Darstellung eines
solchen symbolischen Productes durch eine Summe von einfachen
Ueberschiebungen ermöglicht, liegt in der wiederholten Einführung neuer
Symbole, die so lange fortgesetzt wird, bis in jedem Gliede der
entstehenden Summe nur mehr zwei Symbole übrig sind. (Vgl.\ auch Camille
Jordan, Liouville Journal, sér.\,3, Bd.\,2 und 5.) Ich will der
Deutlichkeit halber ein Beispiel voraussenden.

Es sei gegeben das symbolische Product:
\[
P=(ab)^2(ac)^2(bc)^2 a_x b_x c_x.
\]


% ===== pdf idx 63, printed 49.0 =====
% !PRINTED 49
\weberpage{49}
Halten wir in demselben die beiden Symbole $a$ und $b$ fest, ersetzen dagegen das Symbol $c$ durch eine Variable $y$, so wird aus $P$, abgesehen vom Factor $(xy)$, der Ausdruck:
\[
G_\varrho = (ab)^2\,a_y^2\,b_y^2\,a_x\,b_x.
\]
Derselbe stellt nichts anderes dar, als ein Glied der vierten Polare einer Form:
\[
\alpha_x^6 = (ab)^2\,a_x^3\,b_x^3.
\]
Dieses Glied lässt sich in die Polare $(\alpha_x^6)_{y^4}$ plus einer Summe von Polaren niedrigerer Formen entwickeln, die aus $\alpha_x^6$ durch Faltung entstehen. Das Resultat der Entwicklung können wir direct aus dem Beispiel Nr. 26 ableiten, indem wir in der Gleichung
\[
a_x^2\,b_x^2\,a_y\,b_y = F_{y^2} - \frac{2}{5}\,(xy)\,\Phi_y
\]
$x$ mit $y$ vertauschen und mit $(ab)^2$ multipliciren; dies liefert
\[
(ab)^2\,a_y^2\,b_y^2\,a_x\,b_x = (\alpha_x^6)_{y^4} + \frac{2}{5}\,(xy)\,(\beta_x^4)_{y^3},
\]
wobei mit $\beta_x^4$ die Form $(ab)^3\,a_x^2\,b_x^2$ bezeichnet ist. Ersetzen wir hierin $y_1$ wieder durch $c_2$ und $y_2$ durch $-c_1$ und multipliciren mit dem aus dem unterdrückten Factor $(xy)$ durch diese Substitution wieder hergestellten Factor $c_x$, so erhalten wir:
\[
P = (ab)^2\,(ac)^2\,(bc)^2\,a_x\,b_x\,c_x = (\alpha_x^6,\, c_x^5)^4 + \frac{2}{5}\,(\beta_x^4,\, c_x^5)^3.
\]
Das symbolische Product $P$ ist somit durch eine vierte und eine dritte Ueberschiebung dargestellt.

46. \textit{Allgemeine Darstellung.} Um also ein symbolisches Product (mit beliebig vielen Reihen von Symbolen)
\[
\begin{array}{c}
a,\ b,\ c \dots \text{etc.} \\
\alpha,\ \beta,\ \gamma \dots \text{etc.} \\
\mathfrak{a},\ \mathfrak{b},\ \mathfrak{c} \dots \text{etc.} \\
\cdot\ \ \cdot\ \ \cdot\ \ \cdot\ \ \cdot\ \ \cdot\ \ \cdot
\end{array}
\]
durch eine Summe von Ueberschiebungen darzustellen, verfahren wir folgendermassen. Wir halten eine Reihe von Symbolen fest, etwa die Symbole $a, b, c \dots$ etc., ersetzen alle übrigen durch Variable $y, z, u, v \dots$ etc. und sondern einen Augenblick die dadurch aus Factoren erster Art entstehenden identischen Covarianten $(xy), (xz), (xu), (xv) \dots$ ab. Der so entstehende Ausdruck ist dann nichts anderes, als ein Polarenglied einer Form, die wir mit $f$ bezeichnen wollen. Dieses Glied lässt sich aber nach Nr. 24 und 31 in eine Reihe von Polaren entwickeln, welche Formen angehören mögen, die wir mit
\[
\varphi = A_x^u, \quad \psi = B_x^v, \quad \chi = C_x^\varrho \dots \text{etc.}
\]
%CONT>>


% ===== pdf idx 64, printed 50.0 =====
% !PRINTED 50
\weberpage{50}
%<<CONT
bezeichnen. Ersetzt man nun in dieser Reihe die Formen $\varphi, \psi, \chi \dots$ etc. durch das ihnen zugeordnete Symbol $A, B, C$ etc., substituirt sodann für die Variabeln $y, z, u \dots$ wieder die ursprünglichen Symbole
\[
\begin{array}{c}
\alpha,\ \beta,\ \gamma,\ \dots \text{etc.} \\
\mathfrak{a},\ \mathfrak{b},\ \mathfrak{c},\ \dots \text{etc.}
\end{array}
\]
und multiplicirt nun mit den vorhin abgesonderten Factoren $(xy), (xz) \dots$ d. i. mit:
\[
\begin{array}{c}
\alpha_x,\ \beta_x,\ \gamma_x \dots \\
\mathfrak{a}_x,\ \mathfrak{b}_x,\ \mathfrak{c}_x \dots,
\end{array}
\]
so entsteht auf der linken Seite das ursprüngliche Product $P$. Rechts dagegen erhalten wir eine Summe von symbolischen Producten, deren jedes statt der ganzen Symbolreihe
\[
a,\ b,\ c,\ \dots
\]
nur je ein einziges enthält, sei es das Symbol $A$ der Form $\varphi$, oder das Symbol $B$ der Form $\psi, \dots$ etc. Ein solches Glied enthält also weniger Symbole wie vorher, und wir können nun in jedem derselben, genau auf dieselbe Weise wie in $P$ selbst, die Zahl der Symbole reduciren.

Auf diese Weise gelangen wir schliesslich zu einer Summe von Gliedern, deren jedes nur mehr zwei Symbole enthält. Solche symbolische Producte haben wir aber als einfache Ueberschiebungen definirt, und wir sind somit berechtigt den Satz aufzustellen: \glqq\gesperrt{Jedes symbolische Product lässt sich durch eine Summe von einfachen Ueberschiebungen darstellen.}\grqq{}

\S.\,4. \textbf{Die einfachsten Ueberschiebungen.}

47. \textit{Die Functionaldeterminante.} Unter den Ueberschiebungen haben insbesondere zwei eine hervorragende Bedeutung. Es sind gerade die beiden einfachsten Ueberschiebungen zweier Formen, nämlich die erste und die zweite. Die erste ist bekannt unter dem Namen \glqq Functionaldeterminante\grqq{} oder auch \glqq Jacobi'sche Determinante\grqq{} (vergl. Bd. I \S.\,9). Sie ist, wenn
\[
f = a_x^m, \quad \varphi = b_x^n
\]
die beiden Formen sind, dargestellt durch
\[
(f, \varphi) = a_x^{m-1}\,b_x^{n-1}\,(ab),
\]
und unterscheidet sich von der im ersten Bande \S.\,9 definirten Function nur um den Factor $m \cdot n$. Man erkennt dies unmittelbar aus der unsymbolischen Schreibweise dieser Form (vgl. Nr. 35, III)
%CONT>>


% ===== pdf idx 65, printed 51.0 =====
% !PRINTED 51
\weberpage{51}
%<<CONT
\weq{I}{(f, \varphi) = \begin{vmatrix} f_1, & f_2 \\ \varphi_1, & \varphi_2 \end{vmatrix},}
wobei
\[
f_i = \frac{1}{m}\,\frac{\partial f}{\partial x_i}, \quad \varphi_i = \frac{1}{n}\cdot\frac{\partial \varphi}{\partial x_i}.
\]
Diese Determinante lässt sich übrigens, was für die folgenden Betrachtungen bequemer ist, auch als eine dreigliedrige darstellen. Das symbolische Product
\[
(ab)\,a_x\,b_x = (a_1b_2 - b_1a_2)(a_1x_1 + a_2x_2)(b_1x_1 + b_2x_2)
\]
hat nämlich, wie die Determinantentheorie lehrt (vergl. Bd. I Nr. 40, Anmerkung), als Determinante geschrieben die Matrix
\[
\begin{vmatrix} a_1^2, & a_1a_2, & a_2^2 \\ b_1^2, & b_1b_2, & b_2^2 \\ x_2^2, & -x_1x_2, & x_2^2 \end{vmatrix} = (ab)\,a_x\,b_x.
\]
Multiplicirt man hier die erste Zeile mit $a_x^{m-2}$, die zweite mit $b_x^{n-2}$ und berücksichtigt, dass
\[
\begin{array}{ll}
a_1^2\,a_x^{m-2} = f_{11}, & b_1^2\,b_x^{n-2} = \varphi_{11} \\
a_1a_2\,a_x^{m-2} = f_{12}, & b_1b_2\,b_x^{n-2} = \varphi_{12} \\
a_2^2\,a_x^{m-2} = f_{22}, & b_2^2\,b_x^{n-2} = \varphi_{22},
\end{array}
\]
so erhält man für die Functionaldeterminante die dreigliedrige Determinante:
\weq{II}{(f, \varphi) = (ab)\,a_x^{m-1}\,b_x^{n-1} = \begin{vmatrix} f_{11}, & f_{12}, & f_{22} \\ \varphi_{11}, & \varphi_{12}, & \varphi_{22} \\ x_2^2, & -x_1x_2, & x_1^2 \end{vmatrix}.}

48. \textit{Die Functionaldeterminante als Form ungeraden Charakters.} Da die Functionaldeterminante nur einen einzigen Klammerfactor besitzt, so ist sie eine Form ungeraden Charakters (\textit{forme gauche}). Wir können zeigen, dass ihr Quadrat sich als eine ganze Function von Formen geraden Charakters darstellen lassen muss, eine Eigenschaft, die sie mit allen schiefen Co- und Invarianten theilt. (Vergl. Nr. 8.)

In der That, vertauschen wir in (II) Nr. 47 die erste und letzte Colonne und multipliciren die zweite Colonne mit $-2$, so erhalten wir
\[
+2\,(f, \varphi) = \begin{vmatrix} f_{22}, & -2f_{12}, & f_{11} \\ \varphi_{22}, & -2\varphi_{12}, & \varphi_{11} \\ x_2^2, & +2x_1x_2, & x_1^2 \end{vmatrix}.
\]


% ===== pdf idx 66, printed 52.0 =====
% !PRINTED 52
\weberpage{52}
Bilden wir nun das Product dieser Determinante mit der ursprünglichen, so kommt:
\[
\begin{aligned}
2\,(f, \varphi)\cdot(f, \varphi) &= \begin{vmatrix} f_{11}, & f_{12}, & f_{22} \\ \varphi_{11}, & \varphi_{12}, & \varphi_{22} \\ x_1^2, & -x_1x_2, & x_2^2 \end{vmatrix} \cdot \begin{vmatrix} f_{22}, & -2f_{12}, & f_{11} \\ \varphi_{22}, & -2\varphi_{12}, & \varphi_{11} \\ x_2^2, & 2x_1x_2, & x_1^2 \end{vmatrix} \\
&= \begin{vmatrix} 2(f_{11}f_{22} - f_{12}^2), & f_{22}\varphi_{11} - 2f_{12}\varphi_{12} + f_{11}\varphi_{22}, & f \\ f_{22}\varphi_{11} - 2f_{12}\varphi_{12} + f_{11}\varphi_{22}, & 2(\varphi_{22}\varphi_{11} - \varphi_{12}^2), & \varphi \\ f, & \varphi, & 0 \end{vmatrix}.
\end{aligned}
\]
Hierin ist nach Nr. 45, III:
\[
\begin{array}{l}
f_{11}\varphi_{22} - 2f_{12}\varphi_{12} + f_{22}\varphi_{11} = (f, \varphi)^2 \\
2(f_{11}f_{22} - f_{12}^2) = (f, f)^2 \\
2(\varphi_{11}\varphi_{22} - \varphi_{12}^2) = (\varphi, \varphi)^2.
\end{array}
\]
Demnach erhalten wir als Ausdruck für das Quadrat der Functionaldeterminante $J = (f, \varphi)$
\[
J^2 = \frac{1}{2} \begin{vmatrix} (f, f)^2, & (f, \varphi)^2, & f \\ (f, \varphi)^2, & (\varphi, \varphi)^2, & \varphi \\ f, & \varphi, & 0 \end{vmatrix}
\]
oder
\weq{I}{J^2 = -\frac{1}{2}\left\{(f, f)^2\cdot\varphi^2 - 2(f, \varphi)^2 f\cdot\varphi + (\varphi, \varphi)^2\cdot f^2\right\}.}

Die Formen $f$ und $\varphi$ sowohl, als ihre zweiten Ueberschiebungen, sind aber Formen geraden Charakters; denn $f$ und $\varphi$ haben keine und ihre zweiten Ueberschiebungen haben zwei Klammerfactoren, gehen also bei Transformation bis auf eine gerade Potenz des Moduls $\Delta$ in sich selbst über.

Genau in derselben Weise kann man auch das Product zweier verschiedener Functionaldeterminanten darstellen. Man erhält
\[
\begin{aligned}
2\,(f, \varphi)\cdot(\chi, \psi) &= \begin{vmatrix} f_{11}, & f_{12}, & f_{22} \\ \varphi_{11}, & \varphi_{12}, & \varphi_{22} \\ x_1^2, & -x_1x_2, & x_2^2 \end{vmatrix} \cdot \begin{vmatrix} \chi_{22}, & -2\chi_{12}, & \chi_{11} \\ \psi_{22}, & -2\psi_{12}, & \psi_{11} \\ x_2^2, & +2x_1x_2, & x_1^2 \end{vmatrix} \\
&= \begin{vmatrix} (f, \chi)^2, & (\varphi, \chi)^2, & \chi \\ (f, \psi)^2, & (\varphi, \psi)^2, & \psi \\ f, & \varphi, & 0 \end{vmatrix}
\end{aligned}
\]
oder
\weq{II}{(f, \varphi)\cdot(\chi, \psi) = -\frac{1}{2}\left\{(f, \chi)^2\varphi\cdot\psi - (f, \psi)^2\varphi\cdot\chi - (\varphi, \chi)^2 f\cdot\psi + (\varphi, \psi)^2 f\cdot\chi\right\}.}


% ===== pdf idx 67, printed 53.0 =====
% !PRINTED 53
\weberpage{53}

49. \textit{Die Functionaldeterminante der zweiten Polaren dreier Formen.} Bildet man von den drei Formen
\[
f = a_x^m, \quad \varphi = b_x^n, \quad \psi = c_x^p
\]
die zweiten Polaren
\[
\begin{aligned}
f_{y^2} &= f_{11}y_1{}^2 + 2f_{12}y_1y_2 + f_{22}y_2{}^2 \\
\varphi_{y^2} &= \varphi_{11}y_1{}^2 + 2\varphi_{12}y_1y_2 + \varphi_{22}y_2{}^2 \\
\psi_{y^2} &= \psi_{11}y_1{}^2 + 2\psi_{12}y_1y_2 + \psi_{22}y_2{}^2,
\end{aligned}
\]
so können wir im Sylvester'schen Sinne (vergl. Bd. I. Nr. 143) die Resultante der drei in $y$ quadratischen Formen als Functionaldeterminante betrachten und erhalten, abgesehen von einem Factor 2, als Ausdruck derselben die Determinante
\weq{I}{
\begin{aligned}
U &= \begin{vmatrix} f_{11}, & f_{12}, & f_{22} \\ \varphi_{11}, & \varphi_{12}, & \varphi_{22} \\ \psi_{11}, & \psi_{12}, & \psi_{22} \end{vmatrix}
= a_x^{m-2}b_x^{n-2}c_x^{p-2} \begin{vmatrix} a_1{}^2, & a_1a_2, & a_2{}^2 \\ b_1{}^2, & b_1b_2, & b_2{}^2 \\ c_1{}^2, & c_1c_2, & c_2{}^2 \end{vmatrix} \\
&= (ab)(ac)(bc)\,a_x^{m-2}b_x^{n-2}c_x^{p-2}.
\end{aligned}
}
Dieselbe Determinante geht auch aus $J = (f, \varphi)$ Nr. 47 (II) hervor, wenn wir $x_1$ durch $-\psi_2$ und $x_2$ durch $\psi_1$ ersetzen.

Auch diese Covariante $U$ ist eine Form ungeraden Charakters und wir können sie ganz in derselben Weise wie die Functionaldeterminante $J$ durch gerade Formen darstellen. Multipliciren wir sie nämlich mit
\[
2\,U = \begin{vmatrix} f_{22}, & -2f_{12}, & f_{11} \\ \varphi_{22}, & -2\varphi_{12}, & \varphi_{11} \\ \psi_{22}, & -2\psi_{12}, & \psi_{11} \end{vmatrix},
\]
so erhalten wir:
% ?? uncertain: isolated mark/dot in the left margin before this display, possibly a printer's ornament or dust mark
\weq{II}{
2\,U^2 = \begin{vmatrix} (f,\,f)^2, & (f,\,\varphi)^2, & (f,\,\psi)^2 \\ (f,\,\varphi)^2, & (\varphi,\,\varphi)^2, & (\varphi,\,\psi)^2 \\ (f,\,\psi)^2, & (\varphi,\,\psi)^2, & (\psi,\,\psi)^2 \end{vmatrix}.
}
Die Determinante rechts ist eine Covariante, die, wie wir später sehen werden, unter die Gattung der Combinanten zu rechnen ist. Ihre Elemente sind die Coefficienten von $\lambda$ in der zweiten Ueberschiebung
\[
(F,\,F)^2
\]
der Form
\[
F = \lambda_1 f + \lambda_2 \varphi + \lambda_3 \psi,
\]
über sich selbst. Denn bilden wir die zweite Polare
\[
F_{y^2} = \lambda_1 f_{y^2} + \lambda_2 \varphi_{y^2} + \lambda_3 \psi_{y^2},
\]
%CONT>>


% ===== pdf idx 68, printed 54.0 =====
% !PRINTED 54
\weberpage{54}

%<<CONT
so erhält man daraus die Ueberschiebung, wenn man in dem Producte
\[
\{\lambda_1 f_{y^2} + \lambda_2 \varphi_{y^2} + \lambda_3 \psi_{y^2}\}\{\lambda_1 f_{y^2} + \lambda_2 \varphi_{y^2} + \lambda_3 \psi_{y^2}\}
\]
alle nach der Multiplication entstehenden Glieder nach $y$ zweimal faltet. Dies ist, da wir hier einfache Formen zweiten Grades in $y$ haben, immer nur auf eine einzige Art möglich, so dass wir erhalten:
\[
(F,\,F)^2 = \lambda_1{}^2(f,\,f)^2 + \lambda_2{}^2(\varphi,\,\varphi)^2 + \cdots + 2\lambda_1\lambda_2(f,\,\varphi)^2 + \cdots
\]
50. \textit{Relation zwischen drei und vier binären Formen.} Das in diesem Paragraph nun schon zum wiederholten Male benutzte Princip, durch Determinantenmultiplication Relationen zu erhalten, wollen wir noch benutzen, um Relationen zwischen beliebigen Formen herzustellen.

Wir erhalten eine Relation zwischen den drei beliebigen Formen $f, \varphi$ und $\psi$, wenn wir die beiden Determinanten vom Werthe null
\weq{I}{
D = \begin{vmatrix} f_{11}, & f_{12}, & f_{22}, & 0 \\ \varphi_{11}, & \varphi_{12}, & \varphi_{22}, & 0 \\ \psi_{11}, & \psi_{12}, & \psi_{22}, & 0 \\ x_1{}^2, & -x_1x_2, & x_2{}^2, & 0 \end{vmatrix}
}
und
\weq{II}{
2\,D = \begin{vmatrix} f_{22}, & -2f_{12}, & f_{11}, & 0 \\ \varphi_{22}, & -2\varphi_{12}, & \varphi_{11}, & 0 \\ \psi_{22}, & -2\psi_{12}, & \psi_{11}, & 0 \\ x_2{}^2, & +2x_1x_2, & x_1{}^2, & 0 \end{vmatrix}
}
mit einander multipliciren. Das Resultat der Multiplication ist:
% ?? uncertain: isolated mark/dot in the left margin near this display, possibly a printer's ornament or dust mark
\weq{III}{
0 = \begin{vmatrix} (f,\,f)^2, & (f,\,\varphi)^2, & (f,\,\psi)^2, & f \\ (f,\,\varphi)^2, & (\varphi,\,\varphi)^2, & (\varphi,\,\psi)^2, & \varphi \\ (f,\,\psi)^2, & (\varphi,\,\psi)^2, & (\psi,\,\psi)^2, & \psi \\ f, & \varphi, & \psi, & 0 \end{vmatrix} = \varDelta.
}
Nach den Elementen der letzten Zeile und Reihe ausgewerthet, liefert sie, wenn man die Unterdeterminanten zweiter Ordnung mit $A_{\substack{4k\\4i}}$ bezeichnet (vergl. Bd. I. Nr. 64 und Beispiel 65)
\[
\begin{aligned}
\varDelta &= -\sum_{4i} a_{4i}\,a_{k4}\,A_{\substack{4k\\4i}} \\
&= -\left\{A_{\substack{41\\41}}f^2 + A_{\substack{42\\42}}\varphi^2 + A_{\substack{43\\43}}\psi^2 + 2A_{\substack{41\\42}}f\cdot\varphi + 2A_{\substack{41\\43}}f\cdot\psi + 2A_{\substack{42\\43}}\varphi\cdot\psi\right\},
\end{aligned}
\]
oder in leicht verständlicher Abkürzung:
\[
\varDelta = -(A_1f + A_2\varphi + A_3\psi)^2 = 0.
\]


% ===== pdf idx 69, printed 55.0 =====
% !PRINTED 55
\weberpage{55}

Vertauscht man in den beiden Determinanten (I) und (II) wieder $x_1$ mit $\chi_2$ und $x_2$ mit $-\chi_1$, so erhält man durch Multiplication derselben die Relation:
\[
\begin{vmatrix}
(f,\,f)^2, & (f,\,\varphi)^2, & (f,\,\psi)^2, & (f,\,\chi)^2 \\
(f,\,\varphi)^2, & (\varphi,\,\varphi)^2, & (\varphi,\,\psi)^2, & (\varphi,\,\chi)^2 \\
(f,\,\psi)^2, & (\psi,\,\varphi)^2, & (\psi,\,\psi)^2, & (\psi,\,\chi)^2 \\
(f,\,\chi)^2, & (\chi,\,\varphi)^2, & (\chi,\,\psi)^2, & (\chi,\,\chi)^2
\end{vmatrix} = 0,
\]
und diese Determinante giebt eine Relation zwischen Covarianten von vier beliebigen Formen. Wir werden später, insbesondere bei der Theorie der quadratischen Formen, von den bisher entwickelten Relationen vielfach Gebrauch machen.

51. \textit{Die Functionaldeterminante $(f, \varphi)$ und die erste Polare $(f\cdot\varphi)_y$.} Die Functionaldeterminante giebt auch noch in anderer Richtung zu Relationen Veranlassung, die besonders für die symbolische Rechnung von Nutzen sind. Hierher gehört ihre Beziehung zur Polare $(f\cdot\varphi)_y$. Diese Polare ist nämlich, wenn $f = a_x^m = b_x^m\cdots$ und $\varphi = \alpha_x^n = \beta_x^n\cdots$, dargestellt durch:
\weq{1}{
(f\cdot\varphi)_y = \frac{m}{m+n}\cdot a_x^{m-1}a_y\alpha_x^n + \frac{n}{m+n}a_x^m\alpha_x^{n-1}\alpha_y = c_0G_0 + c_1G_1 \text{ (vgl. Nr. 22).}
}
Die Differenz der beiden Glieder ist:
\[
G_1 - G_2 = a_x^{m-1}\alpha_x^{n-1}(a_y\alpha_x - \alpha_ya_x) = a_x^{m-1}\alpha_x^{n-1}(a\alpha)(yx),
\]
oder
\weq{2}{
(yx)\cdot(f,\varphi) = a_x^{m-1}a_y\alpha_x^n - a_x^m\alpha_x^{n-1}\alpha_y.
}
Eliminiren wir aus beiden Gleichungen (1) und (2) einmal das erste, sodann das zweite Glied rechts, so erhalten wir:
\weq{I}{
\begin{aligned}
(f\cdot\varphi)_y - \frac{m}{m+n}(yx)\cdot(f,\varphi) &= a_x^m\alpha_x^{n-1}\alpha_y = f\cdot\varphi_y \\
(f\cdot\varphi)_y + \frac{n}{m+n}(yx)\cdot(f,\varphi) &= a_x^{m-1}a_y\alpha_x^n = f_y\cdot\varphi
\end{aligned}
}
Diese beiden Gleichungen, von denen die zweite auch durch Vertauschung von $f$ mit $\varphi$ aus der ersten erhalten werden kann, setzen die Polare $(f\cdot\varphi)_y$ mit der Functionaldeterminante $(f,\varphi)$ und den Polaren $f_y$ und $\varphi_y$ in Beziehung.

Dabei haben wir nicht vorausgesetzt, dass $f$ und $\varphi$, unsymbolisch genommen, wirkliche binäre Formen sind. Die Relationen (I) gelten also auch, wenn wir in ihnen $f$ durch $a_x^{m-1}$ und $\varphi$ durch $\alpha_x^{n-1}$ ersetzen und nachträglich mit $(a\alpha)$ multipliciren. Dann geht aber
\[
(f\cdot\varphi)_y \text{ über in } (a\alpha)\cdot(a_x^{m-1}\cdot\alpha_x^{n-1})_y = (f,\varphi)_y,
\]
%CONT>>


% ===== pdf idx 70, printed 56.0 =====
% !PRINTED 56
\weberpage{56}

%<<CONT
ferner:
\[
(f,\varphi) \text{ über in } (a\alpha)(a_x^{m-1},\,\alpha_x^{n-1}) = (a\alpha)^2a_x^{m-2}\alpha_x^{n-2} = (f,\varphi)^2.
\]
Demnach können wir aus den Relationen (I) die beiden neuen gewinnen:
\weq{II}{
\begin{aligned}
(f,\varphi)_y - \frac{m-1}{m+n-2}(yx)\cdot(f,\varphi)^2 &= (a\alpha)\,a_x^{m-1}\alpha_y\alpha_x^{n-2} \\
(f,\varphi)_y + \frac{n-1}{m+n-2}(yx)\cdot(f,\varphi)^2 &= (a\alpha)\,a_x^{m-2}a_y\alpha_x^{n-1}
\end{aligned}
}
In ihnen sind die Ausdrücke rechts die Glieder der Polare $(f,\varphi)_y$ und die Relationen (I) und (II) sind somit nichts anderes, als die in Nr. 25 gegebene Darstellung eines Gliedes durch seine Polare und die Polaren von niedrigeren Formen.

52. \textit{Die Functionaldeterminante von $(f,\varphi)$ und $f$ oder $\varphi$.} Eine besondere Rolle spielen bei späteren Untersuchungen die Functionaldeterminanten von Functionaldeterminanten, also die erste Ueberschiebung von $(f,\varphi)$ über eine dritte Form $\psi$, welche wir symbolisch bezeichnen mit
\[
((f,\varphi),\psi) = U.
\]
Diese Formen haben nämlich die Eigenschaft, dass sie sich immer durch ein Aggregat niedrigerer Formen darstellen lassen. Nehmen wir zunächst an, $\psi$ sei eine der beiden Formen $f$ oder $\varphi$ selbst. Die erste Ueberschiebung $U = ((f,\varphi),f)$ entsteht nämlich, wenn wir in der Polare $(f,\varphi)_y$ die Variabeln $y_1$ durch $b_2$, und $y_2$ durch $-b_1$ ersetzen (vergl. Nr. 35) und mit $b_x^{m-1}$ multipliciren. Aus der zweiten der Relationen (II) erhalten wir demnach:
\[
((f,\varphi),f) = (a\alpha)(ab)a_x^{m-2}b_x^{m-1}\alpha_x^{n-1} - \frac{n-1}{m+n-2}(f,\varphi)^2\cdot b_x\cdot b_x^{m-1}.
\]
Hierin ist das zweite Glied rechts bereits das Product von zwei einfacheren Formen $(f,\varphi)^2$ und $b_x^m = f$. Das erste Glied $G_1$ lässt sich mit Hilfe des Identitätssatzes in ein ähnliches Product verwandeln. Vertauscht man nämlich in ihm die gleichwerthigen Symbole $a$ und $b$, so wird aus dem Gliede
\[
G_1 = (a\alpha)(ab)a_x^{m-2}b_x^{m-1}\alpha_x^{n-1}
\]
wiederum
\[
G_1 = -(b\alpha)(ab)b_x^{m-2}a_x^{m-1}\alpha_x^{n-1}.
\]
Durch Addition erhält man:
\[
2G_1 = (ab)a_x^{m-2}b_x^{m-2}\alpha_x^{n-1}\{(a\alpha)b_x - (b\alpha)a_x\},
\]
oder, weil gemäss dem Identitätssatze die Differenz in der Klammer gleich $(ab)\alpha_x$ ist,
\[
G_1 = \frac{1}{2}(ab)^2a_x^{m-2}b_x^{m-2}\alpha_x^n = \frac{1}{2}(f,f)^2\cdot\varphi.
\]


% ===== pdf idx 71, printed 57.0 =====
% !PRINTED 57
\weberpage{57}
Die Functionaldeterminante $((f,\ \varphi),\ f)$ ist somit dargestellt durch:

\weq{I}{((f,\ \varphi),\ f) = \frac{1}{2}(f,f)^2\cdot\varphi - \frac{n-1}{m+n-2}\cdot(f,\ \varphi)^2\cdot f.}

Dabei ist natürlich vorausgesetzt, dass beide Formen $f$ und $\varphi$ von höherem als dem ersten Grade in $x$ sind.

Vertauschen wir in der Relation (I) $f$ mit $\varphi$, so erhalten wir auch die andere erste Ueberschiebung $((\varphi,\ f),\ \varphi) = -((f,\ \varphi),\ \varphi)$ dargestellt durch die Gleichung:

\weq{II}{((f,\ \varphi),\ \varphi) = -\frac{1}{2}(\varphi,\ \varphi)^2\cdot f + \frac{m-1}{m+n-2}(f,\ \varphi)^2\cdot\varphi\,.}

53. \textit{Die Functionaldeterminante von $(f,\ \varphi)$ und $\psi$.} Ist nun $\psi = r_x^p$ eine von $f$ und $\varphi$ verschiedene Form (sie kann auch selbst wieder eine Functionaldeterminante sein), so können wir die erste Ueberschiebung

\[
U = ((f,\ \varphi),\ \psi)
\]

auf folgende Weise durch Formen niederern Grades ausdrücken. Wir ersetzen zunächst wieder in der zweiten der beiden Relationen (II) Nr. 51 $y$ durch $r$, d. h. $y_1$ durch $r_2$, $y_2$ durch $-r_1$ und erhalten nach Multiplication mit $r_x^{n-1}$

\weq{1}{((f,\ \varphi),\ \psi) = (a\alpha)(ar)a_x^{m-2}\alpha_x^{n-1}r_x^{p-1} - \frac{n-1}{m+n-2}(f,\ \varphi)^2\cdot\psi.}

Zur Umformung des ersten Gliedes benutzen wir den Productsatz

\[
(a\alpha)(ar)\alpha_x r_x = \frac{1}{2}\left\{(a\alpha)^2 r_x^2 + (ar)_x^2\alpha_x^2 - (\alpha r)^2 a_x^2\right\}\cdot
\]

Substituiren wir diesen Werth von $(\alpha a)(ar)\alpha_x r_x$ in (1), so kommt:

\[
((f,\ \varphi),\ \psi) = \frac{1}{2}\left\{(a\alpha)^2 a_x^{m-2}\alpha_x^{n-2}\cdot\psi + (ar)^2 a_x^{m-2}r_x^{p-2}\cdot\varphi - (ar)^2\alpha_x^{m-2}r_x^{p-2}\cdot f\right\}
\]
\[
- \frac{n-1}{m+n-2}(f,\ \varphi)^2\cdot\psi
\]
\[
= \frac{1}{2}(f,\ \varphi)^2\cdot\psi + \frac{1}{2}(f,\ \psi)^2\cdot\varphi - \frac{1}{2}(\varphi,\ \psi)^2\cdot f - \frac{n-1}{m+n-2}(f,\ \varphi)^2\cdot\psi,
\]

oder endlich:

\weq{III}{((f,\ \varphi),\ \psi) = \frac{m-n}{2(m+n-2)}(f,\ \varphi)^2\cdot\psi + \frac{1}{2}\left\{(f,\ \psi)^2\cdot\varphi - (\varphi,\ \psi)^2\cdot f\right\}.}

Wir sind somit berechtigt, den Satz aufzustellen:

\glqq Die Functionaldeterminante einer Form $\psi$ mit einer andern Functionaldeterminante ist, wenn alle Formen vom höheren als vom ersten Grade in $x$ sind, immer durch Producte niedrigerer Formen ausdrückbar.\grqq{}

Aus der Relation (III) gehen selbstverständlich die beiden früheren Relationen (I) oder (II) hervor, wenn wir $\psi$ durch $f$ oder $\varphi$ ersetzen.
%CONT>>


% ===== pdf idx 72, printed 58.0 =====
% !PRINTED 58
\weberpage{58}
54. \textit{Die zweite Ueberschiebung einer Form $f$ über sich selbst.} Als zweites Beispiel einer einfachen Ueberschiebung wollen wir noch die Covariante $(ab)^2 a_x^{m-2}b_x^{m-2}$ betrachten. Diese Form wurde zuerst von Hesse als eine Covariante einer Form $f = a_x^m = b_x^m$ entdeckt (Crelle Journal Bd. 28) und führt nach ihm den Namen Hesse'sche Determinante oder Hesse'sche Form. Sie ist dargestellt symbolisch durch

\[
(f,\ f)^2 = (ab)^2 a_x^{m-2}b_x^{m-2} = H,
\]

unsymbolisch durch:

\[
H = \frac{2}{m^2(m-1)^2}
\begin{vmatrix}
\dfrac{\partial^2 f}{\partial x_1{}^2}, & \dfrac{\partial^2 f}{\partial x_1\,\partial x_2} \\[2ex]
\dfrac{\partial^2 f}{\partial x_1\,\partial x_2}, & \dfrac{\partial^2 f}{\partial x_2{}^2}
\end{vmatrix}\,.
\]

Man kann sie auch auffassen als Functionaldeterminante der beiden ersten Differentialquotienten $\dfrac{\partial f}{\partial x_1}$ und $\dfrac{\partial f}{\partial x_2}$. Wir haben diese Covariante bereits in mehreren Beispielen kennen gelernt. So ist die Discriminante der quadratischen Form $f = a_x^2 = b_x^2$ nichts anderes als die Hesse'sche Form

\[
H = (ab)^2 = 2(\bar a_0\bar a_2 - \bar a_1{}^2),
\]

ebenso war die in Nr. 37 berechnete Covariante zweiten Grades einer cubischen Form die Hesse'sche Covariante

\[
H = (ab)^2 a_x b_x = 2\left\{(\bar a_0\bar a_2 - \bar a_1{}^2)x_1{}^2 + (\bar a_0\bar a_3 - a_1\bar a_2)x_1 x_2 + \cdots\right\}.
\]

55. \textit{Lehrsatz von Hesse.} Von dieser Covariante hatte Hesse den allgemeinen Satz aufgestellt und bewiesen:

\glqq Ist die Determinante

\[
H =
\begin{vmatrix}
\dfrac{\partial^2 f}{\partial x_1{}^2}, & \dfrac{\partial^2 f}{\partial x_1\,\partial x_2}, & \dfrac{\partial^2 f}{\partial x_1\,\partial x_3} & \cdots & \dfrac{\partial^2 f}{\partial x_1\,\partial x_n} \\
\dfrac{\partial^2 f}{\partial x_1\,\partial x_2}, & \dfrac{\partial^2 f}{\partial x_2{}^2}, & \dfrac{\partial^2 f}{\partial x_2\,\partial x_3} & \cdots & \cdot \\
\cdot & \cdot & \cdot & \cdots & \cdot \\
\dfrac{\partial^2 f}{\partial x_1\,\partial x_n}, & \dfrac{\partial^2 f}{\partial x_2\,\partial x_n}, & \cdots & \cdots & \dfrac{\partial^2 f}{\partial x_n{}^2}
\end{vmatrix}
\]

der zweiten Differentialquotienten einer Function $f$ mit $n$ Variabeln identisch Null, so reducirt sich die Dimension der Form $f$ um einen Grad.\grqq{}

Im Jahre 1876 zeigten nun Gordan und Nöther, dass der Satz in dieser Allgemeinheit unrichtig ist.

Durch das Verschwinden der Hesse'schen Form reduciren sich zwar binäre Formen auf reine Potenzen eines linearen Ausdruckes, Formen mit drei Variabeln, also ternäre auf binäre, quaternäre auf ternäre; dagegen giebt es bereits quinäre Formen, deren $H$ zwar
%CONT>>


% ===== pdf idx 73, printed 59.0 =====
% !PRINTED 59
\weberpage{59}
%<<CONT
identisch verschwindet, die sich aber keineswegs auf quaternäre reduciren lassen.

Wir geben hier den Beweis des von Hesse aufgestellten Satzes für binäre Formen, indem wir ihn folgendermassen formuliren.

\glqq Ist die Covariante $(ab)^2 a_x^{m-2}b_x^{m-2}$ der Form $f(x_1 x_2)$ identisch null, so ist $f(x_1 x_2)$ eine reine Potenz: $(\bar a_1 x_1 + \bar a_2 x_2)^n$.\grqq{}

Zum Beweise bringen wir zunächst die Voraussetzung

\weq{1}{H = (ab)^2 a_x^{m-2}b_x^{m-2} = 0}

durch eine andere Relation zum Ausdruck. Wir multipliciren Gleichung (1) mit $\dfrac{(xy)^2}{2}$, und können dann schreiben:

\[
\frac{(xy)^2}{2}(ab)^2 a_x^{m-2}b_x^{m-2} = \frac{1}{2}(a_x b_y - b_x a_y)^2 a_x^{m-2}b_x^{m-2}
\]
\[
= \frac{1}{2}(a_x{}^2 b_x{}^2 - 2a_x b_y b_x a_y + b_x{}^2 a_y{}^2)a_x^{m-2}b_x^{m-2},
\]

oder, weil $a_x^m\cdot b_y^2 b_x^{n-2} = b_x^m a_y^2 a_x^{m-2} = f\cdot f_{y^2}$,

\weq{2}{\frac{(xy)^2}{2}H = f\cdot f_{y^2} - f_y{}^2 = 0.}

Zu dieser Voraussetzung (2) tritt noch eine zweite, dass nämlich jede Form $f(x)$ mindestens einen linearen Factor $\alpha_x$ besitzt (vgl. Fundamentalsatz der Algebra Bd. I \S.\,12), oder, um ganz allgemein zu sein, dass $f$ von der Form ist:

\weq{3}{f = \alpha_x^{\varrho}\cdot\psi_x^{\sigma} = \alpha_x^{\varrho}\cdot\psi\,,^{*)}\quad(\sigma + \varrho = m).}

Hier ist $\alpha_x$ wirklicher vielfacher Factor von $f$, $\varrho$ sein Exponent und $\psi$ das Product aller anderen Factoren. Unter diesen beiden Voraussetzungen lautet dann die Behauptung:

\[
\varrho = m\,.
\]

Wir bilden nun die erste und zweite Polare von $f = \alpha_x^{\varrho}\cdot\psi$, und erhalten nach Nr. 22:

\weq{4}{m\cdot f_y = \varrho\cdot\alpha_x^{\varrho-1}\alpha_y\cdot\psi + \sigma\cdot\alpha_x^{\varrho}\cdot\psi_x^{\sigma-1}\psi_y}

\weq{5}{m(m-1)f_{y^2} = \varrho(\varrho-1)\alpha_x^{\varrho-2}\alpha_y{}^2\cdot\psi + 2\varrho\cdot\sigma\alpha_x^{\varrho-1}\alpha_y\psi_x^{\sigma-1}\psi_y + \sigma(\sigma-1)\alpha_x^{\varrho}\cdot\psi_x^{\sigma-2}\psi_y\,.}

Quadriren wir also den Ausdruck (4) der ersten Polare und subtra-
%CONT>>

\medskip
\noindent{}$^{*)}$ Man darf annehmen $\varrho > 1$. Für $\varrho = 1$ wird der Beweis genau ebenso geführt.


% ===== pdf idx 74, printed 60.0 =====
% !PRINTED 60
\weberpage{60}
%<<CONT
hiren das $(m-1)$fache dieses Quadrates von dem $m$fachen Producte der Form $f$ in die Form (5), so erhalten wir:

\[
m^2(m-1)\{f\cdot f_{y^2} - f_y{}^2\} = \{m\varrho(\varrho-1) - (m-1)\varrho^2\}\alpha_x^{2\varrho-2}\alpha_y{}^2\cdot\psi^2
\]
\[
+ \{2m\varrho\sigma - 2\varrho\sigma(m-1)\}\alpha_x^{2\varrho-1}\alpha_y\cdot\psi_y\cdot\psi
\]
\[
+ \{m\sigma(\sigma-1)\psi\cdot\psi_{y^2} - (m-1)\sigma^2\psi_y{}^2\}\alpha_x^{2\varrho} = 0.
\]

Da nun in dieser Identität die späteren Glieder rechts den Factor $\alpha_x$ in einer höheren Potenz enthalten als das erste, so muss dessen Zahlencoëfficient für sich verschwinden, d. h. es muss sein:

\[
m\varrho(\varrho-1) = (m-1)\varrho^2
\]

oder

\[
m = \varrho\,,
\]

was zu beweisen war.

\begin{center}
\S.\,5. \textbf{Der Aronhold'sche Process.}
\end{center}

56. \textit{Aufgabe.} Viele Fragen der Algebra und Geometrie erheischen die Lösung folgender Aufgabe: Gegeben ist eine Form $f$ und eine ihrer Invarianten (Covarianten) $i$; gegeben ist ferner eine Form $\varphi$ von gleichem Grade in $x$ wie die Form $f$. Dann besitzt das lineare System $f + \lambda\varphi$ eine entsprechende simultane Invariante $J$, welche aus $i$ entsteht, wenn man darin die Coefficienten $\bar a_i$ von $f$ durch die Coefficienten $\bar a_i + \lambda\bar\alpha_i$ von $f + \lambda\varphi$ ersetzt$^{*)}$. Diese Invariante $J$ lässt sich nach Potenzen von $\lambda$ anordnen, und die Coefficienten der Potenzen von $\lambda$ werden alsdann selbst wieder simultane Invarianten von $f$ und $\varphi$ sein. Die Aufgabe nun, die wir uns stellen, ist: womöglich eine independente Darstellung der Coefficienten von $\lambda^k$ zu geben. Es wird sich zeigen, dass sich, so lange $f$ und $\varphi$ von einander unabhängig sind, immer ein Process angeben lässt, eng verwandt mit dem Polarenprocess, vermöge welchem jeder dieser Coefficienten unmittelbar aus der Invariante $i$ hergestellt werden kann; im anderen Falle erhält man für dieselben Recursionsformeln.

Um die Begriffe zu fixiren, nehmen wir an, es sei:

\[
f = a_x^n = b_x^n = \cdots \text{ etc.}
\]
\[
= \bar a_0 x_1^n + \binom{n}{1}\bar a_1 x_1^{n-1}x_2 + \binom{n}{2}\bar a_2 x_1^{n-2}x_2{}^2 + \cdots,
\]
\[
\varphi = \alpha_x^n = \beta_x^n = \cdots \text{ etc.}
\]
\[
= \bar\alpha_0 x^n + \binom{n}{1}\bar\alpha_1 x_1^{n-1}x_2 + \binom{n}{2}\bar\alpha_2 x_1^{n-2}x_2 + \cdots, % [sic]
\]

\medskip
\noindent{}$^{*)}$ Ist $\varphi$ eine Form höheren als $n^{\text{ten}}$ Grades, so bildet man die geeignete Polare und kann sich dann die gleiche Frage für die Form $f + \lambda\varphi_{yk}$ vorlegen.


% ===== pdf idx 75, printed 61.0 =====
% !PRINTED 61
\weberpage{61}
%<<CONT
und die Invariante $\nu$ Grades in den Coefficienten von $f$
\[
i = \psi (\bar a_0,\ \bar a_1,\ \bar a_2,\ \ldots \bar a_n).
\]
Dann ist:
\[
f + \lambda\varphi = (\bar a_0 + \lambda \bar\alpha_0)x_1^n + \binom{n}{1}(\bar a_1+\lambda\bar\alpha_1)x_1^{n-1}x_2 + \cdots
\]
und demnach
\[
J = \psi (\bar a_0+\lambda\bar\alpha_0,\ \bar a_1+\lambda\bar\alpha_1,\ \ldots \bar a_n+\lambda\bar\alpha_n).
\]

Denken wir uns diese Invariante $J$ nach steigenden Potenzen von
$\lambda$ geordnet, und die Coefficienten der ersten bis $\nu^{\text{ten}}$ Potenz der Reihe
nach mit $\delta i,\ \delta^2 i \ldots \delta^\nu i$, so stellt sich $J$ in der Form
dar
\[
J = \lambda^0 i + \lambda^1 \delta i + \lambda^2 \delta^2 i + \lambda^3 \delta^3 i + \cdots \lambda^\nu \delta^\nu i. \tag{I}
\]
Dass der Coefficient von $\lambda^0$ die Invariante $i$ ist, geht daraus hervor,
dass für $\lambda = 0$ die Form $f+\lambda\varphi$ in $f$, also $J$ in $i$ übergehen muss.

57. \textit{Definition des Aronhold'schen Processes.} Erster Fall: $\varphi$ ist
von $f$ unabhängig. Um die Werthe der Grössen $\delta^k i$ zu ermitteln,
bedienen wir uns einen Augenblick einer neuen Art der symbolischen
Darstellung für die Invariante $i$. Da nämlich $i$ als Invariante eine
homogene Function $\nu^{\text{ten}}$ Grades in den Coefficienten $a_k$ ist, so kann
man sie symbolisch darstellen durch
\[
i = (\bar a_0 p_0 + \bar a_1 p_1 + \bar a_2 p_2 + \cdots \bar a_n p_n)^\nu = p_a^\nu.
\]
Hiebei sind die Coefficienten $p_i$ nur als Symbole numerischer
Grössen zu denken, d.\,h. der Zahlencoefficient --- Null nicht ausge\-schlossen --- irgend des Gliedes $\bar a_{i_0}^{\mu_0}\bar a_{i_1}^{\mu_1}\ldots\bar a_{i_\nu}^{\mu_\nu}$ der Invariante ist repräsentirt durch
\[
\frac{\nu!}{\mu_0!\,\mu_1!\ldots}p_{i_0}^{\mu_0}p_{i_1}^{\mu_1}\cdots p_{i_\nu}^{\mu_\nu},\ \text{wo}\ \mu_0+\mu_1+\cdots\mu_\nu = \nu.
\]
Im allgemeinen, wenn $i$ Covariante, können diese Grössen $p_i$ noch
Functionen von $x$ sein.

Da nun $J$ aus $i$ dadurch entstanden ist, dass an Stelle von $\bar a_i$ die
Grösse $\bar a_i + \lambda\bar\alpha_i$ trat, so ist das analoge Symbol für $J$
\[
J = (\bar a_0 p_0+\bar a_1 p_1+\cdots+\lambda\bar\alpha_0 p_0+\lambda\bar\alpha_1 p_1+\cdots)^\nu = (p_a+\lambda p_\alpha)^\nu.
\]
Diesen Ausdruck kann man aber nach dem binomischen Lehrsatz ent\-wickeln, wodurch man erhält:
\[
J = p_a^\nu + \binom{\nu}{1}p_a^{\nu-1}p_\alpha\lambda + \binom{\nu}{2}p_a^{\nu-2}p_\alpha^2\lambda^2 + \cdots \lambda^\nu p_\alpha^\nu. \tag{II}
\]
Vergleicht man diese Entwicklung mit der Entwicklung (I) Nr. 56,
so erhält man allgemein:
\[
\delta^k i = \binom{\nu}{k}p_a^{\nu-k}p_\alpha^k. \tag{I}
\]


% ===== pdf idx 76, printed 62.0 =====
% !PRINTED 62
\weberpage{62}
Der Ausdruck rechts ist aber bis auf einen numerischen Coefficienten
nichts anderes als die $k^{\text{te}}$ Polare von $p_a^\nu$
\[
p_a^{\nu-k}p_\alpha^k = (p_a^\nu)_{\alpha^k} = \frac{1}{\nu(\nu-1)\cdots(\nu-k+1)}\left(\frac{\partial i}{\partial \bar a_0}\bar a_0+\frac{\partial i}{\partial \bar a_1}\bar a_1+\cdots\frac{\partial i}{\partial \bar a_n}\bar a_n\right)^k.
\]
Daher ist
\[
\delta^k i = \binom{\nu}{k}\ \frac{1}{\nu(\nu-1)\cdots(\nu-k+1)}\sum\frac{\partial^k i}{\partial\bar a_{i_1}\cdots\partial\bar a_{i_k}}\ \bar\alpha_{i_1}\bar\alpha_{i_2}\ldots\bar\alpha_{i_k}
\]
oder
\[
\delta^k i = \frac{1}{k!}\sum\frac{\partial^k i}{\partial\bar a_{i_1}\partial\bar a_{i_2}\cdots\partial\bar a_{i_k}}\ \bar\alpha_{i_1}\bar\alpha_{i_2}\cdots\bar\alpha_{i_k}. \tag{II}
\]
Wir erhalten demnach die Lösung:

\glqq Die Grösse $\delta^k i$ ist bis auf einen numerischen Coefficienten
gleich der $k^{\text{ten}}$ Polare von $i$ nach den Coefficienten $\bar a_i$, wobei
die Incremente jedesmal durch $\bar\alpha_i$ ersetzt werden.\grqq{}

Nach den Eigenschaften der Polaren entsteht also $\delta^2 i$ ebenso aus
$\delta i$ wie $\delta i$ aus $i$, oder allgemein: Man erhält $\delta^k i$, wenn man $\delta^{k-1}i$ nach
jedem Coefficienten differentiirt, die Incremente durch $\bar\alpha_i$ ersetzt und
alle so erhaltenen Ausdrücke summirt. Der Process, welcher durch
\[
\delta i = \frac{\partial i}{\partial \bar a_0}\bar\alpha_0 + \frac{\partial i}{\partial \bar a_1}\bar\alpha_1 + \cdots + \frac{\partial i}{\partial \bar a_n}\bar\alpha_n
\]
ausgedrückt ist, wird somit, immer vorausgesetzt, dass die Coefficienten $\bar\alpha_i$
unabhängig sind von den Coefficienten $\bar a_i$, ein Iterationsprocess. Wir
nennen ihn den Aronhold'schen Process und bezeichnen die Aus\-führung desselben an einer In- oder Covariante mit dem Worte \glqq del\-tairen\grqq{}. Die Formen $f$ und $\varphi$ sind hiebei als lineare Functionen der
Variabeln $a_i$, $b_i$ aufgefasst, während die Potenzen $x^m$, $x^{m-1}$, $\ldots x^0$
als Constante betrachtet werden. Beide Formen haben in diesem
Sinne die nämlichen Coefficienten, aber verschiedene Variable und $i$
ist dann Covariante $\nu^{\text{ten}}$ Grades der Form $f$.

58. \textit{Beispiel.} 1) Die quadratische Form
\[
f = \bar a_0 x_1^2 + 2\bar a_1 x_1 x_2 + \bar a_2 x_2^2
\]
besitzt die Invariante
\[
i = \bar a_0 \bar a_2 - \bar a_1^2 = (f,\,f)^2.
\]
Ist nun $\varphi$ eine zweite quadratische Form
\[
\varphi = \bar\alpha_0 x_1^2 + 2\bar\alpha_1 x_1 x_2 + \bar\alpha_2 x_2^2,
\]
so ist die der Form $i$ entsprechende Invariante von $f+\lambda\varphi$
\[
\begin{aligned}
J &= (\bar a_0+\lambda\bar\alpha_0)(\bar a_2+\lambda\bar\alpha_2)-(\bar a_1+\lambda\bar\alpha_1)^2\\
&= (\bar a_0\bar a_2-\bar a_1^2) + \lambda(\bar a_0\bar a_2-2\bar a_1\bar\alpha_1+a_2\alpha_0) + \lambda^2(\bar\alpha_0\bar\alpha_2-\bar\alpha_1^2) % [sic] third term prints \textquotedbl{}a_2 alpha_0\textquotedbl{} without overlines, unlike the parallel terms
\end{aligned}
\]


% ===== pdf idx 77, printed 63.0 =====
% !PRINTED 63
\weberpage{63}
Das gleiche Resultat erhalten wir, wenn wir die erste und zweite
Polare von $i$ bilden und die Werte von
\[
\begin{aligned}
2p_a\,p_\alpha &= \delta i\\
p_\alpha^2 &= \delta^2 i.
\end{aligned}
\]
in die Entwicklung
\[
J = i + \lambda\,\delta i + \lambda^2\,\delta^2 i
\]
eintragen. Denn es ist:
\[
p_a\,p_\alpha = \frac{1}{2}\left(\frac{\partial i}{\partial \bar a_0}\bar\alpha_0+\frac{\partial i}{\partial \bar a_1}\bar\alpha_1+\frac{\partial i}{\partial \bar a_2}\bar\alpha_2\right) = \frac{1}{2}(\bar a_2\bar\alpha_0-2\bar a_1\bar\alpha_1+\bar\alpha_2\bar a_0)
\]
\[
= \frac{1}{2}(f,\,\varphi)^2 \ \text{(vergl. Nr. 36, (1))}
\]
und ebenso (vgl. Nr. 14, 8)
\[
p_\alpha^2 = \frac{1}{2}\left(\frac{\partial i}{\partial \bar a_0}\bar\alpha_0+\frac{\partial i}{\partial \bar a_1}\bar\alpha_1+\frac{\partial i}{\partial \bar a_2}\bar\alpha_2\right)^2 = (\bar\alpha_0\bar\alpha_2-\bar\alpha_1^2) = (\varphi,\,\varphi)^2.
\]
Also ist
\[
J = i + \lambda(f,\,\varphi)^2 + \lambda^2(\varphi,\,\varphi)^2.
\]
59. \textit{Beispiel} 2. Die biquadratische Form $f=a_x^4$ besitzt eine In\-variante dritten Grades in den Coefficienten, welche entsteht, wenn
man die Hesse'sche Form $(f,\,f)^2=(ab)^2a_x^2b_x^2$ viermal über $f=c_x^4$
selbst schiebt, nämlich die Invariante
\[
j = (ab)^2(ac)^2(bc)^2 = ((f,\,f)^2,\,f)^4.
\]
Ersetzt man hier die symbolischen Coefficienten durch die unsymbo\-lischen, so findet man (vgl. auch Nr. 155 d. Bd.), vom Factor 6 abge\-sehen,
\[
j = \begin{vmatrix} \bar a_0 & \bar a_1 & \bar a_2\\ \bar a_1 & \bar a_2 & \bar a_3\\ \bar a_2 & \bar a_3 & \bar a_4 \end{vmatrix}
\]
oder
\[
\begin{aligned}
j &= \bar a_0\bar a_2\bar a_4 + 2\bar a_1\bar a_2\bar a_3 - a_2^3 - \bar a_0\bar a_3^2 - \bar a_1^2\bar a_4 % [sic] \textquotedbl{}a_2^3\textquotedbl{} prints without overline, unlike the surrounding barred terms
\\
&= \{p_0\bar a_0+p_1\bar a_1+p_2\bar a_2+p_3\bar a_3+p_4\bar a_4\}^3 = p_a^3.
\end{aligned}
\]
Die entsprechende Invariante $J$ von $f+\lambda\varphi$, wobei $\varphi=\alpha_x^4$, wird
\[
J = j + \lambda\,\delta j + \lambda^2\,\delta^2 j + \lambda^3\,\delta^3 j
\]
und die Coefficienten $\delta j$, $\delta^2 j$, $\delta^3 j$ sind bezw. gleich den Polaren:
\[
3(p_a^3)_\alpha,\quad 3(p_a^3)_{\alpha^2},\quad (p_a^3)_{\alpha^3}.
\]
Bezeichnen wir der Kürze halber obige Determinante durch
ihr Diagonalglied $(\bar a_0\,\bar a_2\,\bar a_4)$, so werden entsprechend diese drei Po\-
laren durch
%CONT>>


% ===== pdf idx 78, printed 64.0 =====
% !PRINTED 64
\weberpage{64}
%<<CONT
\[
(p_a^3)_\alpha = \frac{1}{3}\sum\frac{\partial j}{\partial\bar a_k}\bar\alpha_k = \frac{1}{3}\left\{(\bar\alpha_0\bar a_2\bar a_4)+(\bar a_0\bar\alpha_2\bar a_4)+(\bar a_0\bar a_2\bar\alpha_4)\right\}
\]
\[
(p_a^3)_{\alpha^2} = \frac{1}{3\cdot 2}\sum\frac{\partial j}{\partial\bar a_k\partial\bar a_\lambda}\bar\alpha_k\bar\alpha_\lambda = \frac{2}{6}\left\{(\bar\alpha_0\bar\alpha_2\bar a_4)+(\bar\alpha_0\bar a_2\bar\alpha_4)+(\bar a_0\bar\alpha_2\bar\alpha_4)\right\}
\]
% ?? uncertain: exact placement of overlines (barred vs. unbarred a/alpha) inside the two bracketed triple-product sums above; reproduced by the pattern implied by the printed exponent structure, not fully legible at this resolution
\[
(p_a^3)_{\alpha^3} = \frac{1}{3\cdot 2\cdot 1}\sum\frac{\partial j}{\partial\bar a_k\partial\bar a_\lambda\partial\bar a_\nu}\cdot\bar\alpha_k\bar\alpha_\lambda\bar\alpha_\nu = (\bar\alpha_0\bar\alpha_2\bar\alpha_4) = ((\varphi,\,\varphi)^2,\,\varphi)^4
\]
dargestellt sein. Demnach wird
\[
\begin{aligned}
J = ((f,\,f)^2,\,f)^4 &+ \lambda\left\{(\bar a_0\bar\alpha_2\bar a_4)+(\bar a_0\bar a_2\bar\alpha_4)+(\bar\alpha_0\bar a_2\bar a_4)\right\}\\
&+ \lambda^2\left\{(\bar\alpha_0\bar\alpha_2\bar a_4)+(\bar\alpha_0\bar a_2\bar\alpha_4)+(\bar a_0\bar\alpha_2\bar\alpha_4)\right\} + \lambda^3((\varphi,\,\varphi)^2,\,\varphi)^4.
\end{aligned}
\]
Hierin sind auch die beiden mittleren Coefficienten Invarianten, deren
symbolische Darstellung wir sofort kennen lernen werden.

60. \textit{Darstellung des Aronhold'schen Processes $\delta i$ in seiner Wirkung
auf ein symbolisches Product.} Schon das letzte Beispiel führte uns zur
Frage: Wie äussert sich der Aronhold'sche Process an einem sym\-
bolischen Product? Die Antwort hierauf will ich zunächst an dem
gleiche Beispiel erläutern. Die gegebene Invariante
\[
j = (ab)^2(bc)^2(ac)^2 = 6\left\{\bar a_0\bar a_2\bar a_4+2\bar a_1\bar a_2\bar a_3-\bar a_2^3-\bar a_0\bar a_3^2-\bar a_1^2\bar a_4\right\} \tag{1}
\]
ist vom dritten Grade in den Coefficienten. Denken wir uns dement\-
sprechend die biquadratische Form $f=r_x^4$ einen Moment auf drei Arten
unsymbolisch dargestellt:
\[
\left.
\begin{aligned}
f &= \bar a_0 x_1^4+4\bar a_1 x_1^3x_2+\cdots \bar a_4 x_2^4\\
f &= \bar b_0 x_1^4+4\bar b_1 x_1^3x_2+\cdots \bar b_4 x_2^4\\
f &= \bar c_0 x_1^4+4c_1 x_1^3x_2+\cdots \bar c_4 x_2^4
\end{aligned}
\right\}\ r_i = \bar a_i = \bar b_i = \bar c_i,
\]
so können wir aus dem symbolischen Product $j=(ab)^2(bc)^2(ac)^2$ noch
zu einer andern unsymbolischen Darstellung $\bar j$ von $j$ gelangen, als der
in (1) gegebenen, wenn wir die Symbolreihe
\[
\begin{aligned}
a_1^4,\ a_1^3a_2,\ a_1^2a_2^2\ldots &\ \text{durch}\ \bar a_0\bar a_1\bar a_2\ldots\\
\text{und ebenso:}\quad b_1^4,\ b_1^3b_2,\ b_1^2b_2^2\ldots &\ \text{durch}\ \bar b_0\bar b_1\bar b_2\ldots\\
c_1^4,\ c_1^3c_2,\ c_1^2c_2^2\ldots &\ \text{durch}\ \bar c_0\bar c_1\bar c_2\ldots
\end{aligned}
\]
ersetzen. Der so erhaltene Ausdruck $\bar j$ ist aber nun linear in den
Grössen $\bar a_i$, $\bar b_i$, $\bar c_i$ und geht für $\bar a_i=\bar b_i=\bar c_i=\bar r_i$ in den unter (1) ge\-
gebenen über. Der partielle Differentialquotient von $\bar j$ nach $\bar r_i$ ist in
Folge dessen durch
\[
\frac{\partial\bar j}{\partial\bar r_i} = \frac{\partial\bar j}{\partial\bar a_i} + \frac{\partial\bar j}{\partial\bar b_i} + \frac{\partial\bar j}{\partial\bar c_i} \tag{2}
\]
%CONT>>


% ===== pdf idx 79, printed 65.0 =====
% !PRINTED 65
\weberpage{65}
dargestellt, und der Aronhold'sche Process erhält den analytischen Ausdruck
\weq{3}{\begin{aligned}
\partial \bar j &= \sum_0^4 \frac{\partial \bar j}{\partial \bar r_i}\,\bar\alpha_i\\
&= \sum_0^4 \frac{\partial \bar j}{\partial \bar a_i}\,\bar\alpha_i + \sum \frac{\partial \bar j}{\partial \bar b_i}\,\bar\alpha_i + \sum \frac{\partial \bar j}{\partial \bar c_i}\,\bar\alpha_i.
\end{aligned}}
Durch die erste der drei Operationen rechts wird aber in den in $\bar\alpha_i$ linearen Ausdruck $\bar j$, an Stelle des Coefficienten $\bar a_i$ der Coefficient $\bar\alpha_i$ oder also in $j$ an Stelle des Symboles $a_1^{4-i}a_2^{\,i}$ das Symbol $\alpha_1^{4-i}\alpha_2^{\,i}$ eingeführt. Wenn man daher zum symbolischen Product zurückkehrt, so erhält man an Stelle von $(ab)^2(bc)^2(ac)^2$ einfach $(\alpha b)^2(bc)^2(\alpha c)^2$.
Ebenso liefert die zweite und dritte Operation rechts in (3)
\[
(a\alpha)^2(\alpha c)^2(ac)^2, \quad (ab)^2(b\alpha)^2(a\alpha)^2.
\]
Weil aber alle drei Resultate identisch gleich sind, so hat man:
\[
\delta j = 3(ab)^2(b\alpha)^2(a\alpha)^2 = \big((ab)^2a_x^2b_x^2,\ \alpha_x^4\big)^4 = 3\big((f,f)^2,\varphi\big)^4.
\]
61. In derselben Weise findet man:
\[
\begin{aligned}
\delta^2 j &= 3(a\alpha)^2(\alpha\beta)^2(a\beta)^2 = \big((\varphi,\varphi)^2,f\big)^4\\
\delta^3 j &= (\alpha\beta)^2(\alpha\gamma)^2(\beta\gamma)^2 = \big((\varphi,\varphi)^2,\varphi\big)^4.
\end{aligned}
\]
Denn in Gleichung (3) sind die drei Glieder rechts linear in den Grössen $\bar a_i, \bar b_i, \bar c_i$; resp. $\bar\alpha_i, \bar a_i, \bar c_i$; $\bar\alpha_i, \bar a_i, \bar b_i$. Setzen wir nun der Einfachheit halber
\[
\begin{aligned}
\sum \frac{\partial \bar j}{\partial \bar a_i}\,\bar\alpha_i &= k_a\\
\sum \frac{\partial \bar j}{\partial \bar b_i}\,\bar\alpha_i &= k_b\\
\sum \frac{\partial \bar j}{\partial \bar c_i}\,\bar\alpha_i &= k_c
\end{aligned}
\]
und unterwerfen die Grössen $k_a$, $k_b$, $k_c$ wiederum dem Aronhold'schen Process, so erhalten wir entsprechend der Gleichung (3)
\[
\begin{aligned}
\delta k_a &= \sum_0^4 \frac{\partial k_a}{\partial \bar b_\varkappa}\,\bar\beta_\varkappa + \sum_0^4 \frac{\partial k_a}{\partial \bar c_\varkappa}\,\bar\beta_\varkappa\\
\delta k_b &= \sum_0^4 \frac{\partial k_b}{\partial \bar a_\varkappa}\,\bar\beta_\varkappa + \sum \frac{\partial k_b}{\partial \bar c_\varkappa}\,\bar\beta_\varkappa\\
\delta k_c &= \sum_0^4 \frac{\partial k_c}{\partial \bar a_\varkappa}\,\bar\beta_\varkappa + \sum \frac{\partial k_c}{\partial \bar b_\varkappa}\,\bar\beta_\varkappa.
\end{aligned}
\]


% ===== pdf idx 80, printed 66.0 =====
% !PRINTED 66
\weberpage{66}
Ersetzen wir in diesen drei Gleichungen $k_a$ durch $\displaystyle\sum \frac{\partial \bar j}{\partial \bar a_i}\,\bar\alpha_i$ etc. und addiren, so kommt:
\[
\sum_{a,\ b,\ c} \delta k_r = 2\left\{ \sum \frac{\partial \bar j}{\partial \bar a_i\,\partial \bar b_\varkappa}\,\bar a_i\bar\beta_\varkappa + \sum \frac{\partial \bar j}{\partial \bar a_i\,\partial \bar c_\varkappa}\,\bar a_i\bar\beta_\varkappa + \sum \frac{\partial \bar j}{\partial \bar b_i\,\partial \bar c_\varkappa}\,\bar a_i\bar\beta_\varkappa \right\}.
\]
Die drei Summen rechts sind linear in $\bar\alpha_i\bar\beta_\varkappa\bar c_\lambda$, resp. $\bar\alpha_i\bar\beta_\varkappa\bar b_\lambda$, $\bar\alpha_i\bar\beta_\varkappa\bar a_\lambda$ und führen demnach wie vorhin eindeutig zurück auf die symbolischen Producte $(\alpha\beta)^2(\alpha c)^2(\beta c)^2$, resp. $(\alpha\beta)^2(\alpha b)^2(\beta b)^2$, $(\alpha\beta)^2(\alpha a)^2(\beta b)^2$, welche alle drei die nämliche Bedeutung haben. Demnach ist
\[
\sum_{a,\ b,\ c} \partial k_r = 6(a\alpha)^2(a\beta)^2(\alpha\beta)^2.
\]
Es ist aber
\[
\sum_{a,\ b,\ c} \partial k_r = 2\sum \frac{\partial j}{\partial \bar r_i\,\partial \bar r_k}\,\bar\alpha_i\bar\alpha_k
\]
und also gemäss der Definition des Aronhold'schen Processes
\[
\delta^2 j = \frac{1}{2}\sum \partial k_r = 3(a\alpha)^2(a\beta)^2(\alpha\beta)^2 = \big((\alpha\beta)^2\alpha_x^2\beta_x^2,\ a_x^4\big)^4.
\]
Ebenso findet man:
\[
\delta^3 j = \frac{1}{3\cdot 2}\sum \frac{\partial \bar j}{\partial \bar a_i\,\partial \bar b_k\,\partial \bar c_\lambda}\,\bar\alpha_i\bar\beta_k\bar\gamma_\lambda = (\alpha\beta)^2(\alpha\gamma)^2(\beta\gamma)^2 = \big((\alpha\beta)^2\alpha_x^2\beta_x^2,\ \gamma_x^4\big).
\]
Die im letzten Beispiele zu berechnende Invariante hat also die Form
\[
J = \big((f,f)^2,f\big)^4 + 3\lambda\big((f,f)^2,\varphi\big)^4 + 3\lambda^2\big((\varphi,\varphi)^2,f\big)^4 + \big((\varphi,\varphi)^2,\varphi\big)^4.
\]
Genau in derselben Weise, wie wir hier an dem einzelnen Beispiele die Wirkung des Aronhold'schen Processes auf ein symbolisches Product studirt haben, so können wir im allgemeinen Falle vorgehen. Wir finden alsdann:
\[
\delta^\varrho i = \frac{1}{\varrho!}\sum \frac{\partial i}{\partial \bar a_1\,\partial \bar b_2\,\partial \bar c_3\cdots \partial \bar m_\varrho}\cdot \bar\alpha_1\bar\beta_2\bar\gamma_3\cdots \bar\mu_\varrho,
\]
wobei sich die Summe über alle Variationen von $n$ Elementen zur $\varrho^{\text{ten}}$ Klasse ohne Wiederholung erstreckt. In Worten ausgedrückt lautet diese Formel:

\glqq Ein symbolisches Product wird $\varrho$ mal deltairt, wenn man je $\varrho$ Symbole $a$, $b$, $c\dots$ durch $\varrho$ Symbole $\alpha$, $\beta$, $\gamma\dots$ so oftmal ersetzt, als sich $n$ Elemente zur $\varrho^{\text{ten}}$ Klasse combiniren lassen und alsdann die Summe aller so erhaltenen Producte bildet.\grqq{}

62. \textit{Zweiter Fall: Die beiden Formen $f$ und $\varphi$ sind nicht von einander unabhängig.} Wir haben bisher angenommen, dass $f$ und $\varphi$ in keinerlei Weise durch ihre Coefficienten mit einander in Beziehung
%CONT>>


% ===== pdf idx 81, printed 67.0 =====
% !PRINTED 67
\weberpage{67}
%<<CONT
stehen. Sobald nun zwischen den beiden Formen $f$ und $\varphi$ ein Abhängigkeitsverhältniss etwa der Art existirt, dass die Coefficienten von $\varphi$ Functionen der Coefficienten von $f$ sind, so ist der Aronhold'sche Process kein Iterationsprocess mehr, d.\,h. $\delta^2 i$ wird nicht in derselben Weise aus $\delta i$ erhalten, wie $\delta i$ aus $i$. Denn wenden wir in diesem Falle auf $\delta i$ den Aronhold'schen Process an, so erhalten wir:
\[
\begin{aligned}
\delta\{\delta i\} &= \delta\left\{ \frac{\partial i}{\partial \bar a_0}\bar\alpha_0 + \frac{\partial i}{\partial \bar a_1}\bar\alpha_1 + \frac{\partial i}{\partial \bar a_2}\bar\alpha_2 + \cdots + \frac{\partial i}{\partial \bar a_n}\bar\alpha_n \right\}\\
&= \frac{\partial^2 i}{\partial \bar a_0^{\,2}}\bar\alpha_0^{\,2} + \frac{\partial^2 i}{\partial \bar a_1^{\,2}}\bar\alpha_1^{\,2} + \cdots + 2\frac{\partial^2 i}{\partial \bar a_0\,\partial \bar a_1}\bar\alpha_0\bar\alpha_1 + 2\frac{\partial^2 i}{\partial \bar a_0\,\partial \bar a_2}\bar\alpha_1\bar\alpha_2 + \cdots % [sic] index pattern as printed
\\
&\quad \cdots + \frac{\partial i}{\partial \bar a_0} \left\{ \frac{\partial \bar\alpha_0}{\partial \bar a_0}\bar\alpha_0 + \frac{\partial \bar\alpha_0}{\partial \bar a_1}\bar\alpha_1 + \cdots + \frac{\partial \bar\alpha_0}{\partial \bar a_n}\bar\alpha_n \right\} + \cdots\\
&\quad \cdots + \frac{\partial i}{\partial \bar a_n} \left\{ \frac{\partial \bar\alpha_n}{\partial \bar a_0}\bar\alpha_0 + \frac{\partial \bar\alpha_n}{\partial \bar a_1}\bar\alpha_1 + \cdots + \frac{\partial \bar\alpha_n}{\partial \bar a_n}\bar\alpha_n \right\},
\end{aligned}
\]
oder
\[
\delta\{\delta i\} = \sum_{\substack{k=0\\ \sigma=0}}^{\substack{k=n\\ \sigma=n}} \frac{\partial^2 i}{\partial \bar a_k\,\partial \bar a_\sigma}\cdot \bar\alpha_k\bar\alpha_\sigma + \sum_0^n \frac{\partial i}{\partial \bar a_k}\sum_0^n \frac{\partial \bar\alpha_i}{\partial \bar a_k}\bar\alpha_k.
\]
Nun ist aber nach Gleichung (II) Nr. 57 das erste Glied rechts identisch mit $2!\,\delta^2 i$, und anderntheils ist nach der Definition des Aronhold'schen Processes:
\[
\sum_0^n \frac{\partial \bar\alpha_i}{\partial \bar a_k}\bar\alpha_k = \delta\alpha_i.
\]
Demnach erhält man:
\[
2!\,\delta^2 i = \delta(\delta i) - \sum_0^n \frac{\partial i}{\partial \bar a}\,\delta\bar\alpha_k.
\]
In derselben Weise findet man weiter:
\[
\begin{aligned}
3!\,\delta^3 i &= \delta(2!\,\delta^2 i) - 2\sum_0^n \frac{\partial i}{\partial \bar a_k\,\partial \bar a_i}\bar\alpha_k\,\delta\alpha_i\\
4!\,\delta^4 i &= \delta(3!\,\delta^3 i) - 3\sum_0^u \frac{\partial i}{\partial \bar a_k\,\partial \bar a_i\,\partial \bar a_\varrho}\bar\alpha_k\bar\alpha_i\,\delta\bar\alpha_\varrho\\
&\ \ \cdots\cdots\cdots\cdots\cdots\cdots\cdots\cdots\cdots\cdots\cdots\cdots\cdots\cdots\\
k!\,\delta^k i &= \delta\big((k-1)!\,\delta^{k-1}i\big) - (k-1)\sum_0^n \frac{\partial i}{\partial \bar a_{i_1}\,\partial \bar a_{i_2}\cdots \partial \bar a_{i_{k-1}}}\bar\alpha_{i_1}\bar\alpha_{i_2}\cdots \delta\bar\alpha_{i_{k-1}}.
\end{aligned}
\]
Dies ist die Recursionsformel für die Coefficienten $\delta^k i$ in der Entwicklung
\[
J = i + \lambda\delta i + \lambda^2\delta^2 i + \cdots \lambda^\nu \delta^\nu i,
\]
%CONT>>


% ===== pdf idx 82, printed 68.0 =====
% !PRINTED 68
\weberpage{68}
%<<CONT
sobald die beiden Formen $f$ und $\varphi$ zu einander in covariantem Verhältnisse stehen. Ich werde im nächsten Paragraphen Nr. 72 noch in ausführlicher Weise auf diese Recursionsformel zurückkommen. Hier will ich nur noch ein Beispiel zur Erläuterung anfügen, um auch für diesen Fall die Wirkung des Aronhold'schen Processes an einem symbolischen Product zu zeigen.

63. \textit{Beispiel.} Wir hatten in Nr. 37 jene drei Covarianten einer cubischen Form $f=a_x^3$ kennen gelernt, durch welche sich in Verbindung mit $f$ alle übrigen rational und ganz darstellen lassen. Sie waren (vergl. auch Nr. 9 und 39):
\[
\begin{aligned}
f &= a_x^3\\
\Delta &= (f,f)^2 = (ab)^2a_xb_x\\
Q &= \big(f,(f,f)^2\big) = (c\Delta)c_x^2\Delta_x = (ab)^2(cb)c_x^2a_x\\
R &= (\Delta,\Delta)^2 = (ab)^2(cd)^2(ac)(bd).
\end{aligned}
\]
Die Covariante $Q$ ist wie $f$ vom dritten Grade in den Variabeln, und man kann sich demnach wieder die Aufgabe stellen, irgend eine dieser drei Covarianten für die Form
\[
F = f + \lambda Q
\]
zu berechnen, etwa gerade die Form $Q_{f+\lambda Q}$ selbst. Sie wird dann dargestellt sein durch
\weq{I}{Q_{f+\lambda Q} = Q + \lambda\delta Q + \lambda^2\delta^2 Q + \lambda^3\delta^3 Q.}
Die Coefficienten $\delta Q$, $\delta^2 Q$, $\delta^3 Q$ lassen sich durch die eben aufgestellten Recursionsformeln berechnen. Zu dem Zwecke ist es nöthig, zunächst die Werthe von
\[
\delta f,\ \delta\Delta,\ \delta Q,\ \delta R
\]
zu ermitteln. Man erhält nach Nr. 60 und 61
\weq{1}{\delta f = \sum \frac{\partial f}{\partial \bar a_i}\bar Q_i = Q}
\weq{2}{\delta\Delta = \delta\{(ab)^2a_xb_x\} = 2(aQ)^2a_xQ_x = (f,Q)^2 = \Psi}
\weq{3}{\begin{aligned}
\delta Q &= \delta\{(c\Delta)c_x^2\Delta_x\} = (c\Psi)c_x^2\Psi_x + (Q\Delta)Q_x^2\Delta_x = \big(f,(f,Q)^2\big) + (Q,\Delta)\\
&= \big(f,(f,Q)^2\big) + (f,\Delta),\Delta) % [sic] parenthesis as printed
\end{aligned}}
oder weil die Functionaldeterminante
\[
\big((f,\Delta),\Delta\big) = -\frac{1}{2}(\Delta,\Delta)^2\cdot f + \frac{2}{3}(f,\Delta)^2\cdot\Delta,
\]
[vergl. Nr. 52 (II)], so wird endlich wegen $(\Delta,\Delta)^2=R$,
\[
\delta Q = \big(f,(f,Q)^2\big) - \frac{1}{2}R\cdot f + \frac{2}{3}\cdot(f,\Delta)^2\cdot\Delta.
\]
Ferner ist:
\weq{4}{\delta R = \delta(\Delta,\Delta)^2 = 2(\Delta,\Psi)^2 = 2\big(\Delta,(f,Q)^2\big).}


% ===== pdf idx 83, printed 69.0 =====
% !PRINTED 69
\weberpage{69}
Nun ist aber, wie wir in der Theorie der cubischen Formen sehen werden [vergl. Nr. 145 (I) und Nr. 146 (II)], oder wie man sich auch unschwer direct unsymbolisch berechnet:

\weq{5}{
(f,\, Q)^2 = 0 \quad \text{und} \quad (f,\, \Delta)^2 = 0.
}
Daher liefert der Aronhold'sche Process:

\weq{II}{
\left.\begin{aligned}
\delta f &= Q\\
\delta \Delta &= 0\\
\delta Q &= -\frac{1}{2}\,R\cdot f\\
\delta R &= 0
\end{aligned}\right\} .
}
64. Die Coefficienten $\delta^2 Q$ und $\delta^3 Q$ in (I) erhält man nun durch die folgenden Betrachtungen. Es ist nach Nr. 62:

\[
2!\,\delta^2 Q = \delta(\delta Q) - \sum \frac{\partial Q}{\partial \bar a_k}\,\delta \bar Q_k .
\]
Da aber $\delta Q = -\dfrac{1}{2}\,R\,.\,f$, und demnach $\delta \bar Q_k = -\dfrac{1}{2}\,R \cdot \bar a_k$ sein muss, so geht die letzte Gleichung über in:

\[
2!\,\delta^2 Q = -\frac{1}{2}\,R \cdot \delta f - \frac{1}{2}\,f \cdot \delta R + \frac{R}{2} \sum \frac{\partial Q}{\partial \bar a_k}\, \bar a_k .
\]
Nach den Gleichungen (II) ist $\delta f = Q$, $\delta R = 0$, und nach dem Euler'schen Satze ist, weil $Q$ vom dritten Grade in den Coefficienten $\bar a_k$:

\[
\sum \frac{\partial Q}{\partial \bar a_k}\, \bar a_k = 3\,Q .
\]
Daher wird:

\weq{III}{
\delta^2 Q = -\frac{1}{4}\,R \cdot Q + \frac{3}{4}\,R \cdot Q = \frac{1}{2}\,R \cdot Q .
}
Tragen wir endlich diese Werthe von $2!\,\delta^2 Q$ und $\delta \bar Q_k$ in die Gleichung

\[
3!\,\delta^3 Q = \delta(2!\,\delta^2 Q) - 2 \sum \frac{\partial Q}{\partial \bar a_k \partial \bar a_i}\, \bar Q_k\, \delta \bar Q_k
\]
ein, so kommt:

\[
\begin{aligned}
3!\,\delta^3 Q &= \delta(R \cdot Q) + R \sum \frac{\partial Q}{\partial \bar a_k \partial \bar a_i}\, \bar Q_k \cdot \bar a_i\\[4pt]
&= -\frac{R^2}{2}\cdot f + R \sum \frac{\partial\!\left(\dfrac{\partial Q}{\partial \bar a_k}\right)}{\partial \bar a_i}\, \bar a_i \cdot \bar Q_k .
\end{aligned}
\]
Weil aber $\dfrac{\partial Q}{\partial \bar a_k}$ vom zweiten Grade in den Coefficienten und daher nach dem Euler'schen Satze:
%CONT>>


% ===== pdf idx 84, printed 70.0 =====
% !PRINTED 70
\weberpage{70}
%<<CONT
\[
\sum \frac{\partial\!\left(\dfrac{\partial Q}{\partial \bar a_k}\right)}{\partial \bar a_i}\, a_i \cdot \bar Q_k = 2 \sum \frac{\partial Q}{\partial \bar a_k}\, \bar Q_k = 2\,\delta Q = -R f,
\]
so wird

\[
\delta^3 Q = -\frac{1}{4}\,R^2 f .
\]
Die cubische Covariante der Form $f + \lambda Q$ hat demnach die Form:

\weq{IV}{
Q_{f+\lambda Q} = Q - \frac{R}{2}\,f \lambda + \frac{R}{2}\,Q \lambda^2 - \frac{R^2}{4}\,f \lambda^3 ,
}
oder

\weq{V}{
Q_{f+\lambda Q} = \left(1 + \frac{R}{2}\cdot \lambda^2\right)\left(Q - \frac{R}{2}\,f \cdot \lambda\right) .
}

\begin{center}
\S\ 6. \textbf{Die Combinanten.}
\end{center}

65. \textit{Definition der Combinante.} Der zuletzt betrachtete Aronhold'sche Process ist auch noch dadurch von Bedeutung geworden, dass durch ihn eine Reihe simultaner Covarianten $i$ eines Systems von Formen $f, \varphi, \psi \ldots$ gleichen Grades in $x$ gekennzeichnet werden kann, welche bei vielen Untersuchungen, insbesondere der Geometrie, eine hervorragende Rolle spielen. Es sind das die unter dem Namen \gesperrt{Combinanten} bekannten Formen, welche wir folgender Weise definiren können:

\glqq{}Sind $f, \varphi, \psi \ldots$ irgend welche Formen $n^{\text{ten}}$ Grades, so ist eine simultane Covariante $i$ derselben Combinante, sobald sie bei Anwendung des Aronhold'schen Processes identisch verschwindet.\grqq{}

Die Combinanten genügen also, wenn wir der Einfachheit halber nur ein System von zwei Formen $f = a_x^n$ und $\varphi = \alpha_x^n$ voraussetzen, der Differentialgleichung

\weq{I}{
\delta i = \frac{\partial i}{\partial \bar a_0}\, \bar \alpha_0 + \frac{\partial i}{\partial \bar a_1}\, \bar \alpha_1 + \cdots + \frac{\partial i}{\partial \bar a_n}\, \bar \alpha_n = 0 .
}
66. \textit{Eigenschaft der Combinante, wenn das gegebene System $f, \varphi, \psi \ldots$ aus unabhängigen Formen besteht.} Sind nun die gegebenen Formen $f$ und $\varphi$ von einander völlig unabhängig, so ergiebt sich aus dieser Definition sofort eine wichtige Eigenschaft für die Combinanten. Sie kann in dem Satze ausgesprochen werden:

\glqq{}Die Combinante $i$ von $f$ und $\varphi$ ändert sich nicht, wenn man $f$ durch $f + \lambda \varphi$ ersetzt, d.\,h., wenn $J$ die entsprechende Form von $f + \lambda \varphi$ und $f$ ist, so besteht die Gleichung:

\[
i = J .
\]
\grqq{}


% ===== pdf idx 85, printed 71.0 =====
% !PRINTED 71
\weberpage{71}
Denn in diesem Falle ist $J$ eine Function der Grössen $\bar a, \bar \alpha, \lambda$ und zwar in der Verbindung $\bar a_k + \lambda \bar \alpha_k$. Wendet man also auf $J$ den Aronhold'schen Process an, indem man $J$ nach $\bar a_k + \lambda \bar \alpha_k$ differentiirt und die Incremente durch $\bar \alpha_k$ ersetzt, so muss nach Definition $J$ der Differentialgleichung genügen:

\[
\frac{\partial J}{\partial(\bar a_0 + \lambda \bar \alpha_0)}\, \bar \alpha_0 + \frac{\partial J}{\partial(\bar a_1 + \lambda \bar \alpha_1)}\, \bar \alpha_1 + \cdots = \delta J = 0 .
\]
Die linke Seite dieser Gleichung ist aber nichts anderes als der Differentialquotient von $J$ nach $\lambda$; demnach ist $\dfrac{\partial J}{\partial \lambda} = 0$, d.\,h. $J$ von $\lambda$ unabhängig. Entwickelt man daher wie im vorigen Paragraphen $J$ nach Potenzen von $\lambda$

\weq{II}{
J = i + \lambda\,\delta i + \lambda^2\,\delta^2 i + \cdots + \lambda^\nu\,\delta^\nu i ,
}
so müssen die Coefficienten der Potenzen von $\lambda$ einzeln verschwinden, was man auch bei der Unabhängigkeit von $f$ und $\varphi$ direct aus $\delta i = 0$ hätte folgern können. Gleichung (II) reducirt sich also auf $J = i$, wie behauptet war.

67. \textit{Beispiele.} Solche Combinanten sind alle \gesperrt{ungeraden Ueberschiebungen} zweier Formen

\[
f = a_x^n, \quad \varphi = \alpha_x^n ,
\]
also die Covarianten

\[
\pi = (a\alpha)\, a_x^{n-1}\, \alpha_x^{n-1}, \quad (a\alpha)^3\, a_x^{n-3}\, \alpha_x^{n-3}, \quad \cdots \quad (a\alpha)^{2\nu+1}\, a_x^{n-(2\nu+1)}\, \alpha_x^{n-(2\nu+1)} .
\]
Denn durch Process $\delta \pi$ erhält man:

\[
\begin{aligned}
\delta \pi &= \delta\!\left[(a\alpha)^{2\nu+1}\, a_x^{n-(2\nu+1)}\, \alpha_x^{n-(2\nu+1)}\right]\\
&= (\beta\alpha)^{2\nu+1}\, \beta_x^{n-(2\nu+1)}\, \alpha_x^{n-(2\nu+1)} .
\end{aligned}
\]
Dieses symbolische Product verschwindet aber identisch, da es durch Vertauschung der gleichwerthigen Symbole $\alpha$ und $\beta$ nur sein Zeichen ändert. Aus denselben Gründen folgt, dass die geraden Ueberschiebungen $(f, \varphi)^{2\mu}$ keine Combinanten sein können.

Die Gesammtheit aller Combinanten zweier Formen $f = a_x^n$ und $\varphi = \alpha_x^n$ ergiebt sich wie Gordan im Bd. V der Math. Annalen gezeigt hat, aus der einen Form

\[
P = f(x)\,\varphi(y) - \varphi(x)\,f(y) = \begin{vmatrix} a_x^n, & \alpha_x^n \\ a_y^n, & \alpha_y^n \end{vmatrix} .
\]
Er nannte sie \glqq{}Fundamentalcombinante\grqq{} von $f$ und $\varphi$. Dass sie in der That eine Combinante ist, erkennen wir unmittelbar, wenn wir $P$ in der Gestalt schreiben:
%CONT>>


% ===== pdf idx 86, printed 72.0 =====
% !PRINTED 72
\weberpage{72}
%<<CONT
\[
P = \begin{vmatrix} \bar a_0 & \bar a_1 & \bar a_2 & \cdots & \bar a_n \\ \bar \alpha_0 & \bar \alpha_1 & \bar \alpha_2 & \cdots & \bar \alpha_n \end{vmatrix} \cdot \begin{vmatrix} x_1^n, & \binom{n}{1} x_1^{n-1} x_2, & \cdots & x_2^n \\ y_1^n, & \binom{n}{1} y_1^{n-1} y_2, & \cdots & y_2^n \end{vmatrix}
\]
\[
= \sum \binom{n}{1}\binom{n}{k} (\bar a_i\, \bar \alpha_k) \left\{ x_1^{n-k} x_2^i\, y_1^{n-i}\, y_2^k - x_1^{n-k} x_2^k\, y_1^{n-i}\, y_2^i \right\},
\]
eine Darstellung, wie wir sie in Bd. I Nr. 151 gegeben haben. Denn der Aronhold'sche Process verwandelt die Determinanten $(\bar a_i\, \bar \alpha_k)$ in die verschwindenden Determinanten $(\bar \alpha_i\, \bar \alpha_k)$. Wir werden im nächsten Paragraphen sehen, wie sich diese Form $P$ nach Potenzen von $(xy)$ entwickeln lässt, deren Coefficienten gerade die oben erwähnten ungeraden Ueberschiebungen $(f, \varphi)^{2\nu+1}$ sind.

In gleicher Weise hat Gordan a. a. O. gezeigt, dass die Fundamentalcombinante von $p$ Formen $f_1, f_2, \ldots f_p$ dargestellt ist durch:

\[
P = \begin{vmatrix}
f_1(x), & f_2(x), & f_3(x) & \cdots & f_p(x) \\
f_1(y), & f_2(y), & f_3(y) & \cdots & f_p(y) \\
f_1(z), & f_2(z), & & \cdots & \\
\vdots & \vdots & & \ddots & \\
f_1(u), & & \cdots & & f_p(u)
\end{vmatrix} .
\]
Sind z.\,B. $f = a_x^2$, $\varphi = b_x^2$, $\psi = c_x^2$ drei quadratische Formen, so ist ihre Fundamentalcombinante

\[
P = \begin{vmatrix} a_x^2, & b_x^2, & c_x^2 \\ a_y^2, & b_y^2, & c_y^2 \\ a_z^2, & b_z^2, & c_z^2 \end{vmatrix}
= \begin{vmatrix} a_1^2, & a_1 a_2, & a_2^2 \\ b_1^2, & b_1 b_2, & b_2^2 \\ c_1^2, & c_1 c_2, & c_2^2 \end{vmatrix} \cdot
\begin{vmatrix} x_1^2, & 2 x_1 x_2, & x_2^2 \\ y_1^2, & 2 y_1 y_2, & y_2^2 \\ z_1^2, & 2 z_1 z_2, & z_2^2 \end{vmatrix}
\]
\[
= (ab)(ac)(bc) \cdot (xy)(xy)(yz) . % [sic] -- printed thus; expected (xy)(xz)(yz) given the later list of identical covariants (xy),(xz),(yz)
\]
Abgesehen von den identischen Covarianten $(xy)$, $(xz)$, $(yz)$ besitzen demnach drei quadratische Formen nur \gesperrt{eine} Combinante, die schon früher berechnete Form (vergl. Nr. 7)

\[
R = (ab)(ac)(bc) .
\]
So ist a. a. O. auch gezeigt, dass für drei Formen dritten Grades $f = a_x^3$, $\varphi = b_x^3$, $\psi = c_x^3$ die Fundamentalcombinante $P$ sich auf die eine Form

\[
Q = (ab)(ac)(bc)\, a_x\, b_x\, c_x
\]
reducirt. Alle Covarianten von $Q$ sind dann selbstverständlich Combinanten von $f$, $\varphi$ und $\psi$.

68. \textit{Ueber eine formale Eigenschaft aller Combinanten.} Aus der Bedingung $\delta i = 0$, welcher eine simultane Covariante genügen muss, damit sie Combinante ist, kann man schliessen, dass diese Combinanten
%CONT>>


% ===== pdf idx 87, printed 73.0 =====
% !PRINTED 73
\weberpage{73}
%<<CONT
auch unsymbolisch vor den übrigen Covarianten sich auszeichnen. Denn dieser Bedingung genügen beispielsweise sicher alle jene Covarianten zweier Formen $f=a_x^n$, $\varphi=\alpha_x^n$, welche die Coefficienten $\overline a_i$, $\overline\alpha_i$ derselben nur in den zweigliedrigen Determinanten $(\overline a_i\ \overline\alpha_k)$ enthalten; dieselben sind demnach stets Combinanten. Also ist auch die Resultante $R_{f,\varphi}$ dieser beiden Formen Combinante, da sie sich in der Bézout'schen Determinantenform gerade als Function dieser Grössen $(\overline a_i\ \overline\alpha_k)$ darstellt. In der That haben wir auch im ersten Bande für sie die aus $\delta i=0$ resultirende Eigenschaft $J=i$ bewiesen, indem wir dort (vergl. § 11 Nr. 146) zeigten, dass $R_{f,\varphi}=R_{f+\lambda\varphi,\varphi}$ ist.

In derselben Weise schliessen wir, dass die simultanen Covarianten von $p$ Formen $f_1, f_2, \cdots f_p$, welche die Coefficienten derselben nur in den $p$-gliedrigen Determinanten $(\overline a_{i_1}^{(1)}\overline a_{i_2}^{(2)}\overline a_{i_3}^{(3)}\cdots\overline a_{i_p}^{(n)})$ enthalten, % [sic] exponent printed (n), not (p) as in the quoted Satz below
Combinanten sein müssen. Denn der Aronhold'sche Process bewirkt, dass jede dieser Determinanten zwei gleiche Zeilen oder Colonnen erhält; in Folge dessen verschwinden sie sämmtlich und mit ihnen die betreffende In- oder Covariante.

Verbinden wir nun aber mit diesen Schlüssen den Gordan'schen Satz, dass aus der Fundamentalcombinante
\[
P = (a_{1x}^n,\ a_{2y}^n,\ a_{3z}^n,\ \dots)
\]
überhaupt alle Combinanten von $p$ Formen $f_i$ abgeleitet werden können, einen Satz, in Bezug auf dessen Beweis ich auf die mehrmals schon citirte Abhandlung in Bd. V der Math. Ann. verweisen muss, so ergiebt sich hieraus:

\glqq Jede Combinante von $p$ binären Formen $f_i$ gleichen Grades in $x$ ist eine solche Function der wirklichen Coefficienten von $f_i$, welche dieselben nur in den Determinantenverbindungen $(\overline a_{i_1}^{(1)}\overline a_{i_2}^{(2)}\overline a_{i_3}^{(3)}\cdots\overline a_{i_p}^{(p)})$ enthält.\grqq{}

69. \textit{Covarianten mit Combinanteneigenschaft binärer von einander abhängiger Formen.} Sind nun die gegebenen Formen nicht, wie wir bisher annahmen, von einander unabhängig, sondern beispielsweise zu einander covariant, so sind zwar immer noch jene Co- und Invarianten $i$, welche der Gleichung $\delta i=0$ genügen, Combinanten der gegebenen Formen; aber dieselben haben alsdann nicht mehr die Eigenschaft $J=i$. Wir können alsdann zur Zeit ihre Eigenschaften überhaupt nicht allgemein discutiren.

Nur ein specieller Fall ist bisher im binären wie im ternären Gebiete untersucht. Das sind die Combinanten einer Form $f$ und ihrer Covariante gleichen Grades $\varphi$, für den Fall, dass, wenn $M$ eine
%CONT>>


% ===== pdf idx 88, printed 74.0 =====
% !PRINTED 74
\weberpage{74}
%<<CONT
Invariante von $f$, zwischen den beiden Formen die Wechselbeziehung stattfindet:
\weq{I}{\begin{aligned} \delta f &= \varphi \\ \delta \varphi &= M\cdot f \end{aligned}}
Ich habe bereits bei Gelegenheit der Darstellung des Aronhold'schen Processes ein diesbezügliches Beispiel gegeben. Die Form $f=a_x^3$ hat eine Covariante gleichen Grades $Q=(ab)^2(cb)a_x c_x^2$ und wir sahen damals, dass
\[
\begin{aligned}
\delta f &= Q\\
\delta Q &= -\frac{R}{2}\cdot f,
\end{aligned}
\]
wobei $R$ die Invariante $(\varDelta,\varDelta)^2$ der Form $f$ ist. Ebenso werden wir später bei den Formen vierten Grades in der Form $f=a_x^4$ und ihrer Hesse'schen Covariante $H=(ab)^2a_x^2b_x^2$ zwei covariante Formen erhalten, welche in derselben Wechselbeziehung stehen, nämlich
\[
\begin{aligned}
\delta f &= H\\
\delta H &= \frac{i}{3}\cdot f,
\end{aligned}
\]
wobei $i$ die Invariante $(ab)^4$ von $f$ bedeutet.

Besitzt nun in einem solchen Falle $f$ eine Invariante oder Covariante $i$, so beschaffen, dass $\delta i=0$, so haben wir, wenn auch nicht direct die gleichen, so doch analoge Bedingungen vor uns, wie wir sie ursprünglich zur Definition der Combinante voraussetzten. Wir stellen uns demnach auch die analoge Frage: In welcher Beziehung steht die entsprechende Invariante $J$ von $f+\lambda\varphi$ zur ursprünglichen $i$ von $f$?

In dem ersten der eben angeführten Beispiele haben wir bereits die Hesse'sche Form $\varDelta$ von $f=a_x^3$ als eine solche Form erkannt, für welche $\delta\varDelta=0$ (vergl. Nr. 63, II). Es ist nun aber nicht mehr wie bei den eigentlichen Combinanten $\varDelta_{f+\lambda\varphi}=\varDelta$; vielmehr erhalten wir:
\[
\varDelta_{f+\lambda\varphi} = \varDelta + \lambda\,\delta\varDelta + \lambda^2\delta^2\varDelta,
\]
oder weil $\delta\varDelta=0$, und demnach (vergl. Nr. 62)
\[
\begin{aligned}
2\,\delta^2\varDelta &= \delta(\delta\varDelta) - \sum \frac{\partial\varDelta}{\partial\overline a_k}\,\delta\overline\alpha_k\\
&= \frac{R}{2}\sum \frac{\partial\varDelta}{\partial\overline a_k}\,\overline a_k = R\cdot\varDelta,
\end{aligned}
\]
so kommt:
\weq{II}{\varDelta_{f+\lambda\varphi} = \left(1 + \frac{R}{2}\,\lambda^2\right)\cdot\varDelta.}
Was sich aber hier in dem speciellen Beispiel ergeben hat, lässt sich auch allgemein beweisen, sobald die oben gegebenen Bedingungen
%CONT>>


% ===== pdf idx 89, printed 75.0 =====
% !PRINTED 75
\weberpage{75}
%<<CONT
(I) erfüllt sind, nämlich die Gleichung:
\[
J_{f+\lambda\varphi} = u\cdot i,
\]
wobei $u$ eine Function des Parameters $\lambda$ und jener Invarianten ist, die sich aus $M$ durch fortgesetztes Deltairen ergeben.

70. \textit{Beweis der Relation $J_{f+\lambda\varphi}=u\cdot i$.} Wir können diese aus den Bedingungen (I) in Nr. 69 resultirende Eigenschaft auf zwei Arten beweisen. Ich will sie zunächst direct durch Rechnung ableiten. Es ist, wie wir schon in Nr. 56 (I) erwähnt haben:
\weq{I}{J_{f+\lambda\varphi} = i + \lambda\delta i + \lambda^2\delta^2 i + \cdots + \lambda^\nu \delta^\nu i,}
wobei für die Coefficienten nach Nr. 62 die Recursionsformel gilt:
\[
\delta^k i = \frac{(k-1)!}{k!}\,\delta(\delta^{k-1}i) - \frac{k-1}{k!}\sum \frac{\partial i}{\partial\overline a_{i_1}\partial\overline a_{i_2}\cdots\partial\overline a_{i_{k-1}}}\cdot \overline\alpha_{i_1}\overline\alpha_{i_2}\cdots \delta\overline\alpha_{i_{k-1}}.
\]

Da aber nach Voraussetzung $\delta\varphi = M\cdot f$ und demnach
\[
\delta\overline\alpha_{i_k} = M.\overline a_{i_k},
\]
so wird:
\[
\begin{aligned}
\delta^k i &= \frac{1}{k}\,\delta(\delta^{k-1}i) - \frac{(k-1)M}{k!}\sum \overline\alpha_{i_1}\overline\alpha_{i_2}\cdots\overline\alpha_{i_{k-2}}\cdot \frac{\partial\left(\dfrac{\partial i}{\partial\overline a_{i_1}\partial\overline a_{i_2}\cdots\partial\overline a_{i_{k-2}}}\right)}{\partial\overline a_{i_{k-1}}}\cdot a_{i_{k-1}}\\
&= \frac{1}{k}\,\delta(\delta^{k-1}i) - \frac{(\nu-k+2)\cdot(k-1)M}{k!}\sum \frac{\partial i}{\partial\overline a_{i_1}\partial\overline a_{i_2}\cdots\partial\overline a_{i_{k-2}}}\cdot \overline\alpha_{i_1}\overline\alpha_{i_2}\cdots\overline\alpha_{i_{k-2}}\\
&= \frac{1}{k}\,\delta(\delta^{k-1}i) - \frac{(\nu-k+2)(k-1)M\cdot(k-2)!}{k!}\,\delta^{k-2}i,
\end{aligned}
\]
oder endlich:
\weq{II}{\delta i^k = \frac{1}{k}\,\delta(\delta^{k-1}i) - \frac{\nu-k+2}{k}\,M\cdot\delta^{k-2}i.}
Legen wir hier der Grösse $k$ der Reihe nach alle Werthe $1, 2, 3\dots$ etc. bei, so erhalten wir:
\[
\begin{aligned}
\delta i &= 0\\
\delta^2 i &= \frac{1}{2}\,\delta(\delta i) - \frac{\nu}{2}\,M\delta^0 i = -\frac{\nu}{2}\,M\cdot i\\
\delta^3 i &= \frac{1}{3}\,\delta(\delta^2 i) - \frac{\nu-1}{3}\,M\delta i = -\frac{\nu}{6}\,i\cdot \delta M\\
\delta^4 i &= \frac{1}{4}\,\delta(\delta^3 i) - \frac{\nu-2}{4}\,M\delta^2 i = -\frac{\nu}{24}\,i\cdot \delta^2 M + \frac{\nu(\nu-2)}{8}\,M^2\cdot i, \quad \text{etc.}
\end{aligned}
\]

Man erkennt, dass rechts stets wieder der Factor $i$ auftritt, gerade weil $\delta i=0$. Der Coefficient jeder Potenz von $\lambda$ enthält sonach diesen Factor.


% ===== pdf idx 90, printed 76.0 =====
% !PRINTED 76
\weberpage{76}
Die Gleichung (I) kann also in der That auf die Form gebracht werden:
\[
J_{f+\lambda\varphi} = u\cdot i,
\]
wo $u$ eine Function der Grössen $\lambda$, $M$ und jener Invarianten ist, in welche $H$ durch den Aronhold'schen Process übergeht.

71. \textit{Zweiter Beweis für diese Relation.} Die Eigenschaft $J_{f+\lambda\varphi}=u\cdot i$ geht auch direct aus der Differentialgleichung $\partial i=0$ durch Integration hervor, wenn wir dieselbe nur vorher geeignet umformen. Zu dem Zwecke müssen wir uns aber zunächst die Frage vorlegen: Was wird aus der Covariante $\varphi$ von $f$, wenn $f$ in $f+\lambda\varphi$ übergeht, oder mit andern Worten: Was ist $\varphi_{f+\lambda\varphi}$?

Die Antwort auf die Frage giebt wiederum die Identität
\weq{I}{\varphi_{f+\lambda\varphi} = \varphi + \lambda\delta\varphi + \lambda^2\delta^2\varphi + \cdots + \lambda^\varrho\delta^\varrho\varphi.}
Die Coefficienten $\delta^k\varphi$ gehen aus der Recursionsformel
\[
\delta^k\varphi = \frac{1}{k}\,\delta(\delta^{k-1}\varphi) - \frac{\varrho-k+2}{k}\,M\delta^{k-2}\varphi
\]
hervor, die ich in Nr. 70 entwickelt habe. Ersetzen wir hierin $k$ durch $2, 3, \dots \varrho$, so erhalten wir, da
\[
\delta\varphi = M.f,
\]
der Reihe nach:
\[
\begin{aligned}
\delta^2\varphi &= \frac{1}{2}\,f\cdot\delta M + \frac{1-\varrho}{2}\,M\cdot\varphi,\\
\delta^3\varphi &= \frac{1}{6}\,f\cdot\{\delta^2 M + 3(1-\varrho)M^2\} + \frac{2-\varrho}{6}\,\delta M\cdot\varphi., \quad \text{etc.}
\end{aligned}
\]

Es ergiebt sich, dass jeder Coefficient $\delta^k\varphi$ eine \gesperrt{lineare Function} der beiden Formen $f$ und $\varphi$ ist, was a priori einzusehen war, da durch den Aronhold'schen Process $f$ in $\varphi$ und $\varphi$ in $f$ verwandelt wird. Die obige Identität (I) lässt sich demnach in der Gestalt [vergl. Nr. 64 (IV)]
\[
\varphi_{f+\lambda\varphi} = A\cdot f + B\cdot\varphi
\]
schreiben, worin $A$ und $B$ Functionen der Grössen $\lambda$, der Invariante $M$ und der daraus durch Deltairen entstehenden Formen sind. Wenn also $f$ durch $f+\lambda\varphi$ ersetzt wird, so geht unter den Voraussetzungen $\delta f=\varphi$, $\delta\varphi=M\cdot f$, die Covariante $\varphi$ von $f$ in eine Form desselben Büschels $f+\lambda\varphi$ über, und die nämliche Stellung, welche eine Covariante oder Invariante $i$ von $f$ zu den Formen $f$ und $\varphi$ einnimmt, behauptet auch $J$ in Bezug auf $f+\lambda\varphi$ und $\varphi_{f+\lambda\varphi}=Af+B\varphi$.

Schreiben wir der Homogenität halber statt $f+\lambda\varphi$ in der Folge $\lambda_1 f+\lambda_2\varphi$, so besitzt diese Form die Coefficienten
\[
\overline a_k\lambda_1 + \overline\alpha_k\lambda_2;
\]
entsprechend besitzt $\varphi_{\lambda_1 f+\lambda_2\varphi}$ die Coefficienten:
\[
\overline a_k A + \overline\alpha_k B.
\]


% ===== pdf idx 91, printed 77.0 =====
% !PRINTED 77
\weberpage{77}
72. \textit{Umformung der Differentialgleichung $\delta i=0$.} Ist nun $i$ entweder direct eine rationale Combinante von $f$ und $\varphi$, oder auch nur schlechthin eine Covariante, rational in den Coefficienten von $f$ von der Eigenschaft, dass $\delta i=0$ (wie in dem Beispiel Nr.\,69 die Form $\varDelta$), so genügt auch die entsprechende Form $J$ von $\lambda_1 f+\lambda_2\varphi$ der Differentialgleichung
\[
\overline{\partial} J=0,
\]
wenn nur der Process $\overline{\partial} J$ auch in der entsprechenden Form:
\[
\overline{\partial} J=\frac{\partial J}{\partial(\lambda_1\overline{a}_0+\lambda_2\overline{\alpha}_0)}\cdot(A\overline{a}_0+B\overline{\alpha}_0)+\frac{\partial J}{\partial(\lambda_1\overline{a}_1+\lambda_2\overline{\alpha}_1)}\cdot(A\overline{a}_1+B\overline{\alpha}_1)+\cdots
\]
ausgeführt wird. Schreibt man die rechte Seite dieser Gleichung in der Gestalt:
\[
\begin{aligned}
&A\left\{\frac{\partial J}{\partial(\lambda_1\overline{a}_0+\lambda_2\overline{\alpha}_0)}\cdot\overline{a}_0+\frac{\partial J}{\partial(\lambda_1\overline{a}_1+\lambda_2\overline{\alpha}_2)}\cdot\overline{a}_1+\cdots\right\}+\\
&B\left\{\frac{\partial J}{\partial(\lambda_1\overline{a}_0+\lambda_2\overline{\alpha}_0)}\cdot\overline{\alpha}_0+\frac{\partial J}{\partial(\lambda_1\overline{a}_1+\lambda_2\overline{\alpha}_1)}\cdot\overline{\alpha}_1+\cdots\right\},
\end{aligned}
\]
so erkennt man die Factoren von $A$ und $B$ als die partiellen Differentialquotienten $\dfrac{\partial J}{\partial\lambda_1}$, $\dfrac{\partial J}{\partial\lambda_2}$. Die Differentialgleichung $\overline{\partial} J=0$ lässt sich also auch durch
\weq{II}{A\frac{\partial J}{\partial\lambda_1}+B\frac{\partial J}{\partial\lambda_2}=0}
darstellen. Diese Differentialgleichung hatte bereits Aronhold als jene Gleichung erkannt, welcher die Combinanten zweier Formen $f$ und $\varphi$ genügen, wenn $\varphi$ Covariante von $f$. Das war der Grund, weshalb Gordan den Process des Deltairens als \gesperrt{Aronhold'schen Process} bezeichnet. Aus ihr können wir nun direct die Combinanteneigenschaft
\[
J=u\cdot i
\]
herleiten.

73. Denn ist $g$ irgend eine particuläre Lösung dieser Gleichung
\[
A\frac{\partial\psi}{\partial\lambda_1}+B\frac{\partial\psi}{\partial\lambda_2}=0,
\]
so muss gemäss der Theorie der Differentialgleichungen die allgemeine Lösung $J$ eine ganze Function
\[
J=F(g)
\]
dieser speciellen Lösung $g$ sein.

Um nun eine solche Lösung $g$ zu finden, gehen wir von irgend einer Combinante von $f$ und $\varphi$ aus, etwa von der Functionaldeterminante $\vartheta=(f,\varphi)$. Es ist dann
\[
\begin{aligned}
\vartheta_{\lambda_1 f+\lambda_2\varphi}&=(\lambda_1 f+\lambda_2\varphi,\ Af+B\varphi)\\
&=(f,\varphi)\cdot\{\lambda_1 B-\lambda_2 A\}.
\end{aligned}
\]


% ===== pdf idx 92, printed 78.0 =====
% !PRINTED 78
\weberpage{78}
Die Form $\vartheta_{\lambda_1 f+\lambda_2\varphi}$ ist homogen in $\lambda_1$ und $\lambda_2$. Demnach auch die Form $\lambda_1 B-\lambda_2 A$, während $(f,\varphi)$ von diesen Parametern völlig unabhängig ist; da $\vartheta_{\lambda_1 f+\lambda_2\varphi}$ als Combinante der obigen Differentialgleichung genügt, so genügt ihr daher auch die specielle Form
\[
g=\lambda_1 B-\lambda_2 A.
\]
Dem Euler'schen Satz gemäss ist hierin:
\[
\begin{aligned}
A&=-g_2=-\frac{1}{\varrho}\,\frac{\partial g}{\partial\lambda_2},\\
B&=\ \ g_1=\ \ \ \frac{1}{\varrho}\,\frac{\partial g}{\partial\lambda_1}.
\end{aligned}
\]
Nun ist $J$ stets homogen in $\lambda_1$ und $\lambda_2$; es muss daher auch in der Gleichung
\[
J=F(g)
\]
die rechte Seite homogen sein, d.\,h. $F$ muss sich auf eine Potenz der particulären Lösung $g$ reduciren. Es ist also
\[
J=g^\mu\cdot\mathrm{Const.}
\]
Um die Constante zu bestimmen, setzen wir $\lambda_1=1$, $\lambda_2=0$; dann geht $J$ in $i$ über, und weil $\vartheta_{\lambda_1 f+\lambda_2\varphi}$ in $\vartheta$ sich verwandelt, auch $g$ in 1. Es ist demnach
\[
i=1^\mu\cdot\mathrm{Const.}
\]
und folglich:
\[
J=g^\mu\cdot i,
\]
was zu zeigen war. [Vergl. das Beispiel in Nr.\,69 (II).]

\begin{center}
\S.\,7.\\
\gesperrt{Binäre Formen mit zwei und mehr Reihen cogredienter Variabler. (Process der Reihenentwicklung.)}
\end{center}

74. \textit{Symbolische Darstellung dieser Formen.} Wir haben uns bisher im Allgemeinen beschränkt auf die Untersuchung von Formen mit einer einzigen Reihe Variabler $x_1$ und $x_2$. Schon die einfachen Polaren derselben führten uns indess auf Formen mit zwei Reihen Variabler $x_1x_2$, $y_1y_2$, während durch die gemischten Polaren sogar Formen mit mehr als zwei Reihen cogredienter Variabler in den Kreis unserer Betrachtungen hereingezogen wurden. Bei der Wichtigkeit, welche diese Formen für die ganze Invariantentheorie haben, wird es nunmehr nothwendig, dieselben einer eingehenden Betrachtung zu unterstellen und insbesondere die Haupteigenschaft derselben zu constatiren, dass alle Co- und Invarianten von Formen mit mehreren Reihen cogredienter Variabler sich durch die Co- und Invarianten von Formen mit nur einer Reihe Veränderlicher darstellen lassen. Dies gilt auch noch, wenn die Veränderlichen $x$ und $y$ nicht direct cogredient sind, sondern beliebig zusammenhängen; es existiren eben dann lineare
%CONT>>


% ===== pdf idx 93, printed 79.0 =====
% !PRINTED 79
\weberpage{79}
%<<CONT
Functionen $z$ der $y$, welche cogredient zu den Variabeln $x$ sind. Hier beschränken wir uns jedoch auf cogrediente Variable $x$ und $y$.

Es sei die Form $F=F(x;y)$ eine rationale ganze und homogene Function der beiden cogredienten Variabeln $x$ und $y$ und zwar sei dieselbe in $x$ homogen von der Dimension $m$, in $y$ homogen von der Dimension $n$. Jedes Glied dieser Function hat, von einem Coefficienten abgesehen, die Form
\weq{1}{x_1^{m-k}x_2^k\cdot y_1^{n-\lambda}y_2^\lambda.}
Die Zahl der Glieder ist $(n+1)(m+1)$, da jedes der $(n+1)$ Glieder einer allgemeinen Function $m^{\text{ten}}$ Grades in $x$ mit immer andern und andern Functionen $n^{\text{ten}}$ Grades in $y$ multiplicirt sein kann.

Um eine solche Form symbolisch darstellen zu können, braucht man auch zwei Reihen von Coefficienten, die eine um ein Glied hinsichtlich der Factoren $x$, die andere um es hinsichtlich der Factoren $y$ zu charakterisiren. Wir werden demnach als symbolischen Coefficienten des Gliedes (1) den Ausdruck wählen
\weq{2}{r_1^{m-k}r_2^k\cdot s_1^{n-\lambda}s_2^\lambda.}
Die Symbole $r$ und $s$ gehören unzertrennlich zusammen und der ganze Coefficient hat nur dann unsymbolische Bedeutung, wenn die Summe der Exponenten von $r_1$ und $r_2$ gleich $m$, die von $s_1$ und $s_2$ gleich $n$ ist. Die Potenz eines solchen Coefficienten werden wir durch ein Product von Ausdrücken (2) charakterisiren, in welchem die einzelnen zusammengehörigen Factoren durch obere Indices von den übrigen unterschieden sind. So ist die dritte Potenz des Coefficienten (2) dargestellt durch:
\[
r_1^{(1)m-k}r_2^{(1)k}s_1^{(1)n-\lambda}s_2^{(1)\lambda}\cdot r_1^{(2)m-k}r_2^{(2)k}s_1^{(2)n-\lambda}s_2^{(2)\lambda}\cdot r_1^{(3)m-k}r_2^{(3)k}s_1^{(2)n-\lambda}s_2^{(3)\lambda}.
% [sic] third factor exponent printed s_1^{(2)}, breaking the (1)(2)(3) pattern
\]
Multipliciren wir die Ausdrücke (1) und (2) mit einander, versehen sie sodann mit dem Product der Binomialcoefficienten
\weq{3}{\binom{m}{k}\binom{n}{\lambda}}
und summiren alle Glieder, die aus diesem einen hervorgehen, indem man den Grössen $k$ und $\lambda$ alle zulässigen Werthe beilegt, so entsteht die symbolische Darstellung
\[
F=(r_1x_1+r_2x_2)^m(s_1y_1+s_2y_2)^n=r_x^m s_y^n
\]
einer Form $F$ mit zwei Reihen cogredienter Variabler. In ihr kommt also erst der Verbindung von $r$ und $s$ eine reale Bedeutung zu, obwohl auch der Fall nicht ausgeschlossen zu werden braucht, dass $F$ in die beiden wirklichen Factoren $r_x^m$ und $s_y^n$ zerfällt.


% ===== pdf idx 94, printed 80.0 =====
% !PRINTED 80
\weberpage{80}
75. \textit{Die Elementarcovarianten von $F$.} Unsere Aufgabe ist nun, die Covarianten einer solchen Function mit zwei Reihen Variabler zu untersuchen. Unter denselben gehören jedenfalls wieder die symbolischen Producte, welche sich zusammensetzen:

$\alpha$) aus den Factoren erster Art:
\[
r_x,\ r_x^{(1)},\ \dots,\ s_x,\ s_x^{(1)},\ \dots;\quad r_y,\ r_y^{(1)},\ \dots,\ s_y,\ s_y^{(1)},
\]
$\beta$) aus den Klammerfactoren:
\[
(rs),\ (r^{(1)}r),\ (rs^{(1)}),\ (ss^{(1)}),\ \dots
\]
Soll das symbolische Product unsymbolische Bedeutung haben, dann müssen die Symbole $r^{(i)}$ in $m$ und die Symbole $s^{(i)}$ in $n$ Factoren vorhanden sein.

Die einfachsten unter diesen Covarianten sind durch jene symbolischen Producte dargestellt, die nur ein einziges Symbolenpaar enthalten, d.\,h. also deren Coefficienten lineare Functionen der Coefficienten der Stammform sind. Dieser Fall trat bei den Formen mit nur einer Reihe von Variabeln nicht auf, da ja die einfachsten Covarianten dort die Ueberschiebungen $(ab)^k a_x^{m-k}b_x^{n-k}$ waren, die bereits vom Grade 2 in den Coefficienten sind.

Eine solche Covariante von $F$ ist dargestellt durch das Product:
\weq{4}{(rs)^\lambda r_x^g r_y^h s_x^k s_y^l,}
wobei
\[
\begin{aligned}
\lambda+g+h&=m\\
\lambda+k+l&=n
\end{aligned}
\]
ist. Unter diesen Covarianten sind wieder die einfachsten jene, in welchen $k=l=0$ ist, also jene, die nur eine einzige Variabelnreihe enthalten. Ihr Symbol ist:
\[
E=(rs)^\lambda r_x^{m-\lambda}s_x^{n-\lambda}=(\overline{rs})^\lambda.
\]
Wir nennen sie in der Folge: \gesperrt{Elementarcovarianten der Form $F$}. Jede Form $F$ besitzt $n+1$ solcher Elementarcovarianten, wenn $n$ der kleinere der beiden Exponenten $m$ und $n$ ist, nämlich die $n$ Formen:
\[
\begin{aligned}
(\overline{rs})^0&=r_x^m s_x^n\\
(\overline{rs})^1&=(rs)\,r_x^{m-1}s_x^{n-1}\\
(\overline{rs})^2&=(rs)^2\,r_x^{m-2}s_x^{n-2}\\
&\ \cdots\cdots\cdots\cdots\\
(\overline{rs})^n&=(rs)^n\,r_x^{m-n}.
\end{aligned}
\]
Sie gehen durch Faltung aus den beiden Factoren $r_x^m$ und $s_x^n$ hervor.


% ===== pdf idx 95, printed 81.0 =====
% !PRINTED 81
\weberpage{81}

76. \textit{Zusammenhang der ursprünglichen Form mit ihren Elementarcovarianten.} Diese letzterwähnten Formen tragen ihren Namen deshalb, weil ihnen die wichtigste Rolle aller Covarianten von $F$ zugetheilt ist. Es gilt nämlich der Satz: \glqq \gesperrt{Die Form} $F$ \gesperrt{ist eine simultane Covariante ihrer Elementarcovarianten.}\grqq

Damit ist aber die Gesammtheit aller invarianten Formen von $F$ auf die simultanen Formen dieser Elementarcovarianten zurückgeführt, und da dieselben nur eine Reihe von Variabeln enthalten, so ist die Theorie einer Form mit mehr Reihen cogredienter Variabler übereinstimmend mit der von simultanen Formen einer einzigen Reihe Variabler.

Wir beweisen den Satz in einer allgemeineren Fassung, indem wir zeigen, dass jedes Product
\[ P = (rs)^\lambda r_x^{m-\lambda} s_y^{n-\lambda}, \quad \lambda = 0, 1, 2, \cdots n, \]
das aus Nr. 75 (4) für $r=k=0$ entsteht, eine lineare Function der mit Potenzen von $(xy)$ multiplicirten Polaren der Elementarcovarianten $E$ ist. Die Form $F = r_x^m s_y^n$ ist dann das erste dieser Producte $P$ für $\lambda = 0$. Der Grad einer jeden solchen mit $(xy)^\nu$ multiplicirten Polare der Elementarcovariante $E$
\[ E_{y^\mu}(xy)^\nu = (\overline{rs})_{y^\mu}^\lambda (xy)^\nu \]
muss in $x$ und $y$ mit dem von $F$ übereinstimmen, wenn $F$ sich in der angegebenen Weise darstellen lässt, d. h. es muss sein:
\[ m+n-2\lambda-\mu+\nu = m \]
und ferner
\[ \mu+\nu = n. \]

Ist alsdann die Form $P$ eine lineare Function dieser Formen $(xy)^\nu E_{y^\mu}$, so ist sie damit auch simultane Covariante derselben.

77. \textit{Beweis, dass $F$ eine lineare Function der Polaren $E_{y^\mu}$ ist.} Wir werden den Beweis des im vorhergehenden Absatze aufgestellten fundamentalen Satzes zunächst ohne rechnerische Operationen führen, indem wir in den einzelnen Schlüssen genau wie in der Polarentheorie vorgehen.

Wir können die Form
\[ P = (rs)^\lambda r_x^{m-\lambda} s_y^{n-\lambda} \]
betrachten als ein Glied der $(n-\lambda)^{\text{ten}}$ Polare von der Elementarcovariante
\[ E = (\overline{rs})^\lambda = (rs)^\lambda r_x^{m-\lambda} s_x^{n-\lambda}. \]

Nach den Sätzen der Polarentheorie lässt sich $P$ alsdann durch
%CONT>>


% ===== pdf idx 96, printed 82.0 =====
% !PRINTED 82
\weberpage{82}

%<<CONT
die Polare der Elementarcovariante $E$ plus einer Summe von Gliedern darstellen von der Form
\[ G = (xy)(rs)^{\lambda+1} r_x^g r_y^h s_x^k s_y^l. \]
Diese Glieder enthalten ausser dem Factor $(xy)$ auch noch einen Klammerfactor $(rs)$ mehr als die Covariante $E = (\overline{rs})^\lambda$. Sie können demnach als Glieder einer Polare von
\[ E_1 = (rs)^{\lambda+1} r_x^{m-\lambda-1} s_x^{n-\lambda-1} \]
aufgefasst werden, und zwar einer Polare niedrigeren Grades als die eben betrachtete. Alle diese Glieder lassen sich wieder durch die Polare der Elementarcovariante $E_1$ darstellen plus einer Summe von abermals einfacheren Producten, die wieder einen Factor $(xy)$ und einen Factor $(rs)$ mehr besitzen. Dieses Verfahren kann fortgesetzt werden, bis man auf die nullte Polare der $(n+1)^{\text{ten}}$ Elementarcovariante gelangt, womit das gewünschte Ziel erreicht ist. Die Form $P$ ist alsdann durch eine lineare Function der Polaren ihrer Elementarcovarianten dargestellt, die nach steigenden Potenzen von $(xy)$ fortschreitet.

Manche Untersuchungen lassen es wünschenswerth erscheinen, in dieser Reihe wiederum eine der Polaren $E_{y^\mu}$ durch eines ihrer Glieder und niedrigere Polaren zu ersetzen. Die äussere Form der Reihe bleibt hierbei unverändert, da sich die niedrigeren Polaren mit den bereits vorhandenen zusammenziehen lassen, so dass sich nur die numerischen Coefficienten der Reihe in anderer Weise darstellen.

Diese Entwicklung war unabhängig von dem Werthe $\lambda$; sie gilt also auch für $\lambda = 0$, d. h. für die Form
\[ F = r_x^m s_y^n. \]
Es bleibt nur noch die Aufgabe zu erledigen, für diese einfachste Form mit zwei Reihen Veränderlicher die Coefficienten der einzelnen Glieder in independenter Weise aufzustellen.

78. \textit{Die Reihenentwicklung von $F$ ist eindeutig.} Ehe wir jedoch zur Berechnung der Coefficienten schreiten, wollen wir noch zeigen, dass sich $F$ nur auf eine einzige Art in eine solche Reihe entwickeln lässt. Gäbe es nämlich zwei Entwicklungen für $F$, die nicht übereinstimmen:
\[ F = \sum C'_\mu \cdot E_{y^\mu} \cdot (xy)^\nu \]
und
\[ F = \sum C''_\mu \cdot E_{y^\mu} \cdot (xy)^\nu, \]
dann würde die Differenz beider eine Entwicklung für den Werth Null geben:
\[ 0 = \sum (C'_\mu - C''_\mu)\, E_{y^\mu} . (xy)^\nu . \]
% ?? uncertain: an isolated mark appears to the right of this display in the scan (possibly a print blemish or stray punctuation); not transcribed as text


% ===== pdf idx 97, printed 83.0 =====
% !PRINTED 83
\weberpage{83}

Nehmen wir nun an $C'_{\mu_i} - C''_{\mu_i}$ sei die erste nicht verschwindende Differenz in dieser Entwicklung, so können wir, nachdem wir dieselbe mit der niedrigsten Potenz von $(xy)$ dividirt haben, in ihr $x=y$ setzen. Dann verschwinden eben wegen des Factors $(xy)$ alle Glieder bis auf das erste. Man hätte also
\[ 0 = C'_{\mu_i} - C''_{\mu_i}, \ \text{oder} \ C'_{\mu_i} = C''_{\mu_i}. \]

Also können die beiden Entwicklungen nicht von einander verschieden sein.

79. \textit{Wirkung des $\Omega$-Processes auf ein Product $\varphi \cdot \psi$ zweier Formen.} Um nun die Coefficienten $C_\mu$ der Entwicklung zu bestimmen, müssen wir vorher noch eine Hilfsformel aufstellen, die uns lehrt, wie sich das Product $\varphi \cdot \psi$ zweier Formen mit zwei Reihen cogredienter Variabler unter dem Einfluss des Omegaprocesses verändert.

Es ist:
\[ \begin{aligned} \dfrac{\partial(\varphi\cdot\psi)}{\partial x_1} &= \varphi\,\dfrac{\partial\psi}{\partial x_1} + \psi\,\dfrac{\partial\varphi}{\partial x_1} \\ \dfrac{\partial(\varphi\cdot\psi)}{\partial x_2} &= \varphi\,\dfrac{\partial\psi}{\partial x_2} + \psi\,\dfrac{\partial\varphi}{\partial x_2}, \end{aligned} \]
und demnach:
\weq{1}{\dfrac{\partial^2(\varphi\cdot\psi)}{\partial x_1\partial y_2} = \varphi\,\dfrac{\partial^2\psi}{\partial x_1\partial y_2} + \dfrac{\partial\varphi}{\partial y_2}\dfrac{\partial\psi}{\partial x_1} + \dfrac{\partial\psi}{\partial y_2}\dfrac{\partial\varphi}{\partial x_1} + \psi\,\dfrac{\partial^2\varphi}{\partial x_1\partial y_2}}
\weq{2}{\dfrac{\partial^2(\varphi\cdot\psi)}{\partial x_2\partial y_1} = \varphi\,\dfrac{\partial^2\psi}{\partial x_2\partial y_1} + \dfrac{\partial\varphi}{\partial y_2}\dfrac{\partial\psi}{\partial x_1} + \dfrac{\partial\psi}{\partial y_1}\dfrac{\partial\varphi}{\partial x_2} + \psi\,\dfrac{\partial^2\varphi}{\partial x_2\partial y_1}.}
Setzt man nun der Kürze halber
\[ \dfrac{\partial\varphi}{\partial y_2}\dfrac{\partial\psi}{\partial x_1} - \dfrac{\partial\varphi}{\partial y_1}\dfrac{\partial\psi}{\partial x_2} = (\varphi\overset{\smile}{,}\psi), \quad \dfrac{\partial\psi}{\partial y_2}\dfrac{\partial\varphi}{\partial x_1} - \dfrac{\partial\psi}{\partial y_1}\dfrac{\partial\varphi}{\partial x_2} = (\psi\overset{\smile}{,}\varphi), \]
so erhält man, wenn $\varrho$ und $\mu$ die Grade von $x$ in $\varphi$ resp. $\psi$ bedeuten, $\sigma$ und $\nu$ die von $y$ in $\varphi$ resp. $\psi$, aus (1) und (2) durch Subtraction:
\weq{I}{(\varrho+\mu)(\sigma+\nu)\Omega(\varphi\cdot\psi) = \mu\cdot\nu\cdot\varphi\,\Omega(\psi) + (\varphi\overset{\smile}{,}\psi) + \varrho\cdot\sigma\cdot\psi\,\Omega(\varphi) + (\psi\overset{\smile}{,}\varphi).}
Ist speciell $\psi = (xy)^\lambda$, so ist
\[ \lambda^2\Omega(xy)^\lambda = 2\lambda(xy)^{\lambda-1} + \lambda(\lambda-1)(xy)^{\lambda-1} = \lambda(\lambda+1)(xy)^{\lambda-1}, \]
ferner wird
\[ (\varphi\overset{\smile}{,}(xy)^\lambda) = \dfrac{\partial\varphi}{\partial y_2}\lambda(xy)^{\lambda-1}y_2 + \dfrac{\partial\varphi}{\partial y_1}\lambda(xy)^{\lambda-1}y_1 \]
\[ ((xy)^\lambda\overset{\smile}{,}\varphi) = \lambda(xy)^{\lambda-1}x_1\dfrac{\partial\varphi}{\partial x_1} + \lambda(xy)^{\lambda-1}x_2\dfrac{\partial\varphi}{\partial x_2}, \]
folglich nach dem Euler'schen Satze:
\[ (\varphi\overset{\smile}{,}\psi) + (\psi\overset{\smile}{,}\varphi) = \lambda(\varrho+\sigma)(xy)^{\lambda-1}\varphi. \]
Daher erhalten wir aus Gleichung (I) für $\psi = (xy)^\lambda$:
\weq{II}{(\varrho+\lambda)(\sigma+\lambda)\Omega(\varphi\cdot(xy)^\lambda) = \varrho\cdot\sigma\cdot(xy)^\lambda\Omega(\varphi) + \lambda(\varrho+\sigma+\lambda+1)(xy)^{\lambda-1}\cdot\varphi.}


% ===== pdf idx 98, printed 84.0 =====
% !PRINTED 84
\weberpage{84}

80. \textit{Berechnung der Coefficienten $C_\mu$.} Die zuletzt erhaltene Formel benutzen wir nun, um die Coefficienten der Reihe
\weq{III}{(rs)^\lambda r_x^{m-\lambda} s_y^{n-\lambda} = \sum C_{\lambda,k} (\overline{rs})_{y^{n-k}}^k (xy)^{k-\lambda}, \quad k \geqq \lambda}
% ?? uncertain: printed relation symbol before "lambda" read as double-bar geqq; single-bar geq is also plausible at this resolution
zu bestimmen. Dass die einzelnen Glieder rechts die hier gegebene Form besitzen müssen, haben wir bereits in Nr. 76 und 77 entwickelt. Ihr Grad in $y$ ist
\[ n-k+k-\lambda = n-\lambda, \]
ihr Grad in $x$ ist
\[ m+n-2k-(n-k)+k-\lambda = m-\lambda, \]
wie es die linke Seite der Relation (III) auch erheischt.

Zunächst erhält man für $\lambda=k$ den ersten Coefficienten $C_{k,k}$, indem man $x=y$ setzt. Dabei geht die linke Seite in die Elementarcovariante
\weq{1}{E = (\overline{rs})^\lambda = (rs)^\lambda r_x^{m-\lambda} s_x^{n-\lambda}}
über, während die rechte Seite sich auf ihr erstes Glied
\weq{2}{C_{k,k}\left[(\overline{rs})_{y^{n-\lambda}}^\lambda\right]_{y=x} = C_{k,k}(\overline{rs})^\lambda}
reducirt. Durch Comparation von (1) mit (2) erhält man
\[ C_{k,k} = 1. \]

Um nun die weiteren Coefficienten zu finden, wenden wir auf beiden Seiten der Relation (III) den $\Omega$-Process an, was immer erlaubt ist, da dieselbe eine Identität darstellt. Die Rechnung erleichtert uns Formel (II), Nr. 79. In ihr tritt nur an Stelle von $\varphi$ die Polare $(\overline{rs})_{y^{n-k}}^k$, an Stelle von $\lambda$ der Exponent $k-\lambda$; ebenso haben wir ihren Grad $\varrho$ in $x$ durch $m-\lambda$, ihren Grad $\sigma$ in $y$ durch $n-\lambda$ zu ersetzen.

Der $\Omega$-Process auf eine Polare angewendet liefert aber null (vgl. Nr. 19). Es bleiben daher rechts nur die vom zweiten Terme der Formel (II) herrührenden Glieder unter dem Summenzeichen übrig, nämlich die Glieder:
\[ (k-\lambda)(m+n-k-\lambda+1)(xy)^{k-\lambda-1}\cdot(\overline{rs})_{y^{n-k}}^k. \]
Berücksichtigen wir nun, dass nach Nr. 20 (I)
\[ (m-\lambda)(n-\lambda)\Omega\!\left((rs)^\lambda r_x^{m-\lambda} s_y^{n-\lambda}\right) = (m-\lambda)(n-\lambda)(rs)^{\lambda+1} r_x^{m-\lambda-1} s_y^{n-\lambda-1}, \]
so liefert der $\Omega$-Process, angewendet auf die Reihe (III)
\weq{IV}{\begin{aligned}(m-\lambda)(n-\lambda)(rs)^{\lambda+1} r_x^{m-\lambda-1} s_y^{n-\lambda-1} \\ = \sum C_{\lambda,k}(k-\lambda)(m+n-k-\lambda+1)(xy)^{k-\lambda-1}(\overline{rs})_{y^{n-k}}^k.\end{aligned}}


% ===== pdf idx 99, printed 85.0 =====
% !PRINTED 85
\weberpage{85}
Ersetzen wir nun aber anderntheils in der Relation (III)
\[
\lambda \text{ durch } \lambda + 1,
\]
so erhalten wir für die linke Seite der Relation (IV) eine zweite Entwicklung
\weq{V}{(rs)^{\lambda+1} r_x^{m-\lambda-1} s_y^{n-\lambda-1} = \sum_{\lambda+1,\,k} (\overline{rs})^k_{y^{n-k}} (xy)^{k-\lambda-1}.}
Multipliciren wir diese Entwicklung mit $(m-\lambda)(n-\lambda)$ und compariren wir (IV) und (V), so erhalten wir die recurrente Coefficienten-Relation
\weq{VI}{(m-\lambda)(n-\lambda)\,C_{\lambda+1,\,k}=C_{\lambda,\,k}(k-\lambda)(m+n-k-\lambda+1),}
wobei
\[
C_{k,\,k}=1.
\]

81. Damit ist das Mittel gegeben, die Coefficienten einzeln zu berechnen. Setzen wir $k-\lambda=\varrho$, und somit $\lambda=k-\varrho$, so liefern die Substitutionen
\[
\varrho=1,\ \varrho=2,\ \dots\ \varrho=\varrho
\]
in die Formel
\[
(m-k+\varrho)(n-k+\varrho)\,C_{k+1-\varrho,\,k}=C_{k-\varrho,\,k}(m+n-2k+\varrho+1)\,\varrho
\]
der Reihe nach folgende Gleichungen:
\begin{align*}
(m-k+1)(n-k+1)\cdot 1 &= C_{k-1,\,k}(m+n-2k+2)\\
(m-k+2)(n-k+2)\,C_{k-1,\,k} &= C_{k-2,\,k}(m+n-2k+3)\cdot 2\\
(m-k+3)(n-k+3)\,C_{k-2,\,k} &= C_{\lambda-3,\,k}(m+n-2k+4)\cdot 3 % [sic] printed C_{\lambda-3,k}, apparent misprint for C_{k-3,k}
\\
\cdots\cdots\cdots\cdots &\cdots\cdots\cdots\cdots\\
\cdots\cdots\cdots\cdots &\cdots\cdots\cdots\cdots\\
(m-k+\varrho)(n-k+\varrho)\,C_{k-\varrho+1,\,k} &= C_{k-\varrho,\,k}(m+n-2k+\varrho+1)\,\varrho.
\end{align*}
Multiplicirt man die untereinander stehenden Gleichungen, so erhält man:
\begin{align*}
&(m-k+1)(m-k+2)\cdots(m-k+\varrho)(n-k+1)(n-k+2)\cdots(n-k+\varrho)\\
&= C_{k-\varrho,\,k}\cdot \varrho!\,(m+n-2k+2)(m+n-2k+3)\cdots(m+n-2k+\varrho+1),
\end{align*}
oder:
\weq{1}{\frac{(m-k+\varrho)!}{(m-k)!}\cdot\frac{(n-k+\varrho)!}{(n-k)!}=C_{k-\varrho,\,k}\cdot \varrho!\,\frac{(m+n-2k+\varrho+1)!}{(m+n-2k+1)!}.}
Nun ist aber:
\[
\frac{(m-k+\varrho)!}{(m-k)!}=\binom{m-k+\varrho}{\varrho}\cdot \varrho\,!
\]
und analoge Beziehungen gelten für die beiden andern Quotienten. Führt man also in (1) diese Werthe für die betreffenden Quotienten ein, so erhält man nach Division mit $(\varrho!)^2$:
\weq{VII}{\binom{m-k+\varrho}{\varrho}\binom{n-k+\varrho}{\varrho}=C_{k-\varrho,\,k}\binom{m+n-2k+\varrho+1}{\varrho}.}
Wir haben die Formel (VII) nicht in dieser Allgemeinheit nöthig, sondern nur für die Reihenentwicklung von
\[
F=r_x^m s_y^n,
\]
%CONT>>


% ===== pdf idx 100, printed 86.0 =====
% !PRINTED 86
\weberpage{86}
%<<CONT
also für den Fall
\[
\lambda=0,\ \text{d.\,h.}\ k=\varrho.
\]
Man erhält alsdann:
\[
\binom{m}{k}\binom{n}{k}=C_{0,\,k}\binom{m+n-k+1}{k},
\]
und demnach:
\[
F=r_x^m s_y^n=\sum C_{0,\,k}(\overline{rs})^k_{y^{n-k}}(xy)^k,
\]
oder
\weq{VIII}{r_y^m s_y^n=\sum \frac{\binom{m}{k}\binom{n}{k}}{\binom{m+n-k+1}{k}}\,(\overline{rs})^k_{y^{n-k}}(xy)^k.} % [sic] printed r_y^m s_y^n; preceding line has r_x^m s_y^n
Diese Entwicklung wurde von \gesperrt{Clebsch} und \gesperrt{Gordan} gleichzeitig entdeckt, und wir haben hier dieselbe auf dem von \gesperrt{Clebsch} eingeschlagenen Wege gegeben. Durch sie ist also der Beweis geliefert:

\begin{quote}
\glqq Jede Form mit zwei Reihen cogredienter Variabeln ist eine simultane Covariante ihrer Elementarcovarianten, also von Formen mit nur einer Variabelnreihe.\grqq{}
\end{quote}

82. \textit{Beispiel für die Clebsch-Gordan'sche Reihenentwicklung.} Wir wollen die Form
\[
F=r_x^2 s_y
\]
in eine Reihe entwickeln, und durch einfache Rechnung die Identität der Entwicklung mit der ursprünglichen Form nachweisen. Es ist:
\[
F=r_x^2 s_y=\sum \frac{\binom{2}{k}\binom{1}{k}}{\binom{4-k}{k}}\,(\overline{rs})^k_{y^{n-k}}(xy)^k,
\]
oder
\[
r_x^2 s_y=(\overline{rs})^0_y+\frac{2}{3}(\overline{rs})^1(xy).
\]
Also:
\[
F=\frac{1}{3}\left\{r_x^2 s_y+2r_x s_x r_y\right\}+\frac{2}{3}(rs)r_x(xy)
\]
\[
=\frac{1}{3}F+\frac{2}{3}r_x\left\{s_x r_y+(rs)(xy)\right\}.
\]
Nun ist aber:
\[
(rs)(xy)=r_x s_y - s_x r_y.
\]
Demnach wird:
\[
F=\frac{1}{3}F+\frac{2}{3}r_x\left\{s_x r_y+r_x s_y-s_x r_y\right\}
\]
\[
=\frac{1}{3}F+\frac{2}{3}r_x^2 s_y=F.
\]


% ===== pdf idx 101, printed 87.0 =====
% !PRINTED 87
\weberpage{87}
Weitere Beispiele. Es ist für die in $y$ linearen Formen $F$
\[
r_x s_y=(\overline{rs})^0_y+\frac{1}{2}(\overline{rs})^1(xy)
\]
\[
r_x^2 s_y=(\overline{rs})^0_y+\frac{2}{3}(\overline{rs})^1(xy)
\]
\[
r_x^m s_y=(\overline{rs})^0_y+\frac{m}{m+1}(\overline{rs})^1(xy).
\]
Ebenso erhalten wir für die in $y$ quadratischen Formen $F$
\[
r_x^2 s_y^2=(\overline{rs})^0_{y^2}+(\overline{rs})^1_y(xy)+\frac{1}{3}(\overline{rs})^2(xy)^2
\]
\[
r_x^4 s_y^2=(\overline{rs})^0_{y^2}+\frac{4}{3}(\overline{rs})^1_y(xy)+\frac{3}{5}(\overline{rs})^2(xy)^2
\]
\[
r_x^m s_y^2=(\overline{rs})^0_{y^2}+\frac{2m}{m+2}(\overline{rs})^1_y(xy)+\frac{m-1}{m+1}(\overline{rs})^2(xy)^2.
\]
Entwickeln wir endlich noch eine in $y$ biquadratische Form in ihre Reihe, etwa $r_x^4 s_y^4$, so erhalten wir:
\[
r_x^4 s_y^4=(\overline{rs})^0_{y^4}+2(\overline{rs})^1_{y^3}(xy)+\frac{12}{7}(\overline{rs})^2_{y^2}(xy)^2+\frac{4}{5}(\overline{rs})^3_y(xy)^3+\frac{1}{5}(\overline{rs})^4(xy)^4.
\]

83. \textit{Verallgemeinerung der Reihenentwicklung.} Wir können die eben gewonnene Reihenentwicklung noch auf etwas complicirtere Formen ausdehnen. Die zunächst liegenden Formen $F$ sind durch das symbolische Product dargestellt:
\[
r_x^\lambda s_y^m r_y^\varrho=F'.
\]
Um die Reihen für diese Form zu erhalten, brauchen wir nur die Formel (VIII) $\varrho$mal nach $x$ zu differentiiren und die Incremente durch $y$ zu ersetzen; dann erhalten wir:
\[
m(m-1)(m-2)\dots(m-\varrho+1)\,r_x^{m-\varrho}s_y^n r_y^\varrho
=\sum \frac{\binom{m}{k}\binom{n}{k}}{\binom{m+n-k+1}{k}}(m-k)(m-k+1)\dots(m-k-\varrho+1)(xy)^k(\overline{rs})^k_{y^{n-k+\varrho}}.
\]
Hierbei muss $\varrho$ kleiner als $(m-k)$ sein. Nach Division der Gleichung mit dem Coefficienten der linken Seite erhält man
\weq{1}{r_x^{m-\varrho} s_y^n r_y^\varrho=\sum \frac{\binom{m}{k}\binom{n}{k}}{\binom{m+n-k+1}{k}}\cdot\frac{(m-k)\dots(m-k-\varrho+1)}{m(m-1)\dots(m-\varrho+1)}\cdot(xy)^k(\overline{rs})^k_{y^{n-k+\varrho}}.}
Nun ist aber
\[
\binom{m}{k}\cdot\frac{(m-k)\dots(m-k-\varrho+1)}{m(m-1)\dots(m-\varrho+1)}=\frac{m(m-1)\dots(m-\varrho+1)(m-\varrho)\dots(m-\varrho-k+1)}{m(m-1)\dots(m-\varrho+1)\,.\,1\,.\,2\,.\,3\,.\,.\,k}=\binom{m-\varrho}{k}.
\]


% ===== pdf idx 102, printed 88.0 =====
% !PRINTED 88
\weberpage{88}
Setzen wir also $m-\varrho=\lambda$, so geht (1) über in:
\weq{IX}{r_x^\lambda s_y^n r_y^\varrho=\sum \frac{\binom{\lambda}{k}\binom{n}{k}}{\binom{\lambda+\varrho+n-k+1}{k}}\cdot(\overline{rs})^k_{y^{n-k+\varrho}}(xy)^k.}

84. \textit{$F$ ist ein Product von vier und mehr Factoren.} Während also für die Reihenentwicklung einer binären Form $F$ mit zwei Reihen Variabeln die Coefficienten noch allgemein angegeben werden können, wenn $F$ durch drei Factoren repräsentirt wird, ist uns das nicht mehr möglich, wenn die Zahl der Factoren die Anzahl drei überschreitet. Die figurirten Zahlen reichen alsdann nicht mehr aus, das Coefficientengesetz allgemein festzustellen. Die Reihe muss dann für jeden einzelnen Fall berechnet werden. Für den Fall von vier Factoren kann man sich etwa in folgender Weise behelfen. Es sei
\[
F=r_x^\lambda r_y^\varrho s_y^\sigma s_x^\mu
\]
und $\mu$ der kleinste der Exponenten. Erheben wir die Identität
\[
r_y s_x=\{r_x s_y-(rs)(xy)\}
\]
auf die $\mu^{\text{te}}$ Potenz, entwickeln rechts nach dem binomischen Lehrsatz und multipliciren alsdann mit $r_y^{\varrho-\mu}r_x^\lambda s_y^\sigma$, so erhalten wir eine Reihe, in welcher nur mehr Factoren wie $r_x^k s_y^\lambda$ neben den Potenzen von $(xy)^k$ und den Polaren der Elementarcovarianten auftreten. Diese beiden Factoren der einzelnen Glieder entwickelt man nun nach Reihe (VIII) und erhält so die allgemeine Darstellung.

Eine andere Methode ist die folgende. Man setzt
\[
r_x^\lambda s_x^\mu\cdot r_y^\varrho s_y^\sigma=R_x^{\lambda+\mu}\cdot S_y^{\varrho+\sigma}.
\]
Nun entwickelt man $R_x^{\lambda+\mu}\cdot S_y^{\varrho+\sigma}$ nach Reihe (VIII), ersetzt in ihr hierauf $R$ und $S$ in geeigneter Weise durch $r_x^\lambda s_x^\mu$ und $r_y^\varrho s_y^\sigma$ und wendet sodann auf die erhaltenen neuen Glieder wiederum die Entwicklung (VIII) an.

85. \textit{Symmetrische und alternirende Formen.} Diese Reihenentwicklungen werden durchgehends einfacher, sobald $F$ eine symmetrische oder eine alternirende Form ist, was natürlich von vornherein voraussetzt, dass $m=n$ sei. Im ersten Falle ändert sich nämlich $F=r_x^m s_y^m$ nicht, wenn $x$ mit $y$, oder was das nämliche ist, wenn $r$ mit $s$ vertauscht wird. Durch Vertauschung von $r$ mit $s$ verändern aber alle ungeraden Elementarcovarianten, also die Formen:
\[
E=(\overline{rs})^{2\lambda+1},
\]
ihr Zeichen. Sie müssen demnach, da die Function $F$ selbst unverändert
%CONT>>


% ===== pdf idx 103, printed 89.0 =====
% !PRINTED 89
\weberpage{89}
bleibt, identisch null sein. Die Reihe reducirt sich sonach auf die H\"alfte der Glieder.

Im zweiten Falle verschwinden aus dem gleichen Grunde die geraden Elementarcovarianten, also die Formen
\[
E = (rs)^{2\lambda},
\]
die ja bei Vertauschung von $r$ mit $s$ sich nicht \"andern, w\"ahrend $F$ als alternirende Function ihr Zeichen \"andert. Zu diesen Formen geh\"ort unter andern die in Nr. 67 aufgestellte Fundamentalcombinante
\[
f(x)\varphi(y) - \varphi(x)f(y) = P.
\]

86. \textit{Umkehrung der Reihe (VIII).} Ersetzt man in der Reihe (VIII) die Gr\"ossen $m$ und $n$ bezw. durch $m-\nu$, $n-\nu$ und multiplicirt auf beiden Seiten mit $(rs)^\nu \cdot (xy)^\nu$, so kommt:
\[
(rs)^\nu r_x^{m-\nu} s_y^{n-\nu} (xy)^\nu = \sum \frac{\binom{m-\nu}{k}\binom{n-\nu}{k}}{\binom{m+n-2\nu-k+1}{k}} (\overline{rs})_{y^{n-(\nu+k)}}^{\nu+k} (xy)^{k+\nu}. % ?? uncertain: exact stacked placement of the exponent n-(v+k) on the y-subscript of the barred (rs) term
\]

Legt man hier dem $\nu$ alle m\"oglichen Werthe bei von $\nu = 0$ bis $\nu = n$, so erh\"alt man $(n+1)$ Gleichungen, in denen rechts nur die $(n+1)$ Polaren der $(n+1)$ Elementarcovarianten $(rs)^k$ in linearer Verbindung auftreten. Dabei beginnt mit wachsendem $\nu$ die Entwicklung rechts mit einer immer h\"oheren Elementarcovariante, so dass also, wenn man diese $(n+1)$ Gleichungen nach der Determinantenmethode aufl\"ost, die Nennerdeterminante den Werth Eins besitzt, da ihr Diagonalglied nur Elemente vom Werthe 1 und links derselben nur Elemente vom Werthe 0 enth\"alt. Bei Aufl\"osung der Gleichungen nach diesen Polaren tritt also kein Nenner auf, und es wird sich somit der Werth einer Polare in der Form darstellen:
\weq{1}{(\overline{rs})^k_{y^{n-k}} (xy)^k = \sum C_{k,\,\nu} (rs)^\nu r_x^{m-\nu} s_y^{n-\nu} (xy)^\nu \cdot}

Die Glieder unter dem Summenzeichen rechts gehen, vom Factor $C_{k,\,\nu} (xy)^\nu$ abgesehen, aus $F = r_x^m s_y^n$ durch den $\Omega$-Process hervor (vgl. Nr. 20).

Auch diese Reihenentwicklung ist eindeutig, wie sich durch die gleichen Schl\"usse ergiebt, womit wir die Eindeutigkeit der Reihe (VIII) nachgewiesen haben. Wir d\"urfen also behufs Co\"efficientenbestimmung wiederum den $\Omega$-Process anwenden. Dabei beschr\"anken wir uns indess auf den Fall $k=0$; der andere ergiebt sich aus diesem, wenn man nur f\"ur $m$ und $n$ entsprechend andere Werthe einf\"uhrt.

87. \textit{Berechnung der Co\"efficienten.} F\"ur $k=0$, $\nu = k$ wird die Relation (1)
\weq{2}{(rs)^0_{y^n} = (r_x^m s_x^n)_{y^n} = \sum C_k (rs)^k r_x^{m-k} s_y^{n-k} (xy)^k \cdot}


% ===== pdf idx 104, printed 90.0 =====
% !PRINTED 90
\weberpage{90}
Da links eine Polare steht, so liefert der $\Omega$-Process auf dieser Seite Null. Auf der rechten Seite erhalten wir nach Formel (II) Nr. 79 zwei Summen, von denen wir die eine nach links transponiren. Es entsteht dann, da $\varrho = (m-k)$ und $\sigma = (n-k)$, die Beziehung:
\[
\left.\begin{aligned}
-\sum_{0}^{n} C_k\,(m-k)(n-k)(rs)^{k+1} r_x^{m-k-1} s_y^{n-k-1}(xy)^k \\
= \sum_{0}^{n} C_k\,k\,(m+n-k+1)(rs)^k r_x^{m-k} s_y^{n-k}(xy)^{k-1}
\end{aligned}\right\} \cdot (\mathrm{X})
\]
Da das erste Glied der Summe rechts verschwindet, so k\"onnen wir auf dieser Seite $k$ durch $k+1$ ersetzen und erhalten:
\[
-\sum C_k\,(m-k)(n-k)(rs)^{k+1} r_x^{m-k-1} s_y^{n-k-1}(xy)^k = \sum C_{k+1}\,(k+1)(m+n-k)(rs)^{k+1} r_x^{m-k-1} s_y^{n-k-1}(xy)^k \cdot
\]
Hieraus ergiebt sich durch Comparation gleich hoher Potenzen
\weq{3}{-C_k\,(m-k)(n-k) = C_{k+1}\,(k+1)(m+n-k).}
Nun folgt aber aus Reihe (2) f\"ur $k=0$, $C_0 = 1$. Demnach erhalten wir aus (2), indem wir der Gr\"osse $k$ alle Werthe von 0 bis $k$ beilegen:
\[
\begin{aligned}
m \cdot n &= C_1\,(m+n) \\
- C_1\,(m-1)(n-1) &= C_2\,(m+n-1)\cdot 2 \\
- C_2\,(m-2)(n-2) &= C_3\,(m+n-2)\cdot 3 \\
&\cdots\cdots\cdots\cdots\cdots\cdots\cdots\cdots \\
&\cdots\cdots\cdots\cdots\cdots\cdots\cdots\cdots \\
- C_{k-1}\,(m-k+1)(n-k+1) &= C_k\,(m+n-k+1)\,k.
\end{aligned}
\]
Durch Multiplication dieser Gleichungen mit einander erh\"alt man:
\[
(-1)^k \cdot \binom{m}{k} k! \binom{n}{k} k! = C_k \cdot k!\,k! \binom{m+n}{k}
\]
oder
\[
C_k = (-1)^k \frac{\binom{m}{k}\binom{n}{k}}{\binom{m+n}{k}}.
\]
Tr\"agt man diesen Werth von $C_k$ in die Relation (2) ein, so erh\"alt man:
\weq{XI}{(r_x^m s_x^n)_{y^n} = \sum (-1)^k \frac{\binom{m}{k}\binom{n}{k}}{\binom{m+n}{k}} (rs)^k r_x^{m-k} s_y^{n-k}(xy)^k \cdot}

\S.\,8. \textbf{Anwendungen des Processes der Reihenentwicklung.}

88. \textit{Zweck der Reihenentwicklungen.} Die Reihenentwicklungen dienen in erster Linie dazu, complicirte symbolische Producte in einfache \"uberzuf\"uhren, oder unbekannte Formen durch bekannte darzu-
%CONT>>


% ===== pdf idx 105, printed 91.0 =====
% !PRINTED 91
\weberpage{91}
%<<CONT
stellen. Sie werden so das wichtigste Hilfsmittel f\"ur die symbolische Rechnung, insbesondere bei Aufstellung der von einander linear unabh\"angigen Co- und Invarianten einer Form $n^{\mathrm{ten}}$ Grades. Wesentliche Vortheile bieten die hier entwickelten Formeln auch, wenn es sich darum handelt, lineare Relationen zwischen Ueberschiebungen herzustellen. In seinem Programme \glqq Ueber das Formensystem bin\"arer Formen\grqq{} (Teubner, 1875) hat Gordan aus der Hauptformel (VIII) noch mehrere andere durch Differentiation und einfache Substitutionen abgeleitet, welche ganz allgemein solche Relationen zwischen Ueberschiebungen mit drei Reihen von Symbolen bieten.

Im Folgenden wollen wir nun in einer Reihe von Anwendungen zeigen, inwiefern die Reihenentwicklungen zu einem so wirksamen Hilfsmittel der symbolischen Rechnung werden; zahlreiche specielle Beispiele werden wir im zweiten und dritten Teile dieses Bandes kennen lernen.

89. \textit{Die Ueberschiebnngen der Form $f = a_x^4$ \"uber ihre Hesse'sche Covariante $\Delta = (ab)^2 a_x^2 b_x^2$.} % [sic] \textquotedbl{}Ueberschiebnngen\textquotedbl{} printed thus, misprint for \textquotedbl{}Ueberschiebungen\textquotedbl{}
Als erste Anwendung will ich hier die Berechnung der vier Ueberschiebungen $(f,\Delta)^k$, $k=4,3,2,1$, so weit dies vortheilhaft, mit Reihenentwicklung geben.

Am einfachsten berechnet sich $(f,\Delta)^4$; man bedarf hierbei keiner Reihenentwicklung; vielmehr ergiebt sich unmittelbar aus der vierten Polare
\[
\Delta_{y^4} = (ab)^2 a_y^2 b_y^2,
\]
indem wir nach der in Nr. 35 gegebenen Theorie $y_1$ durch $c_2$, $y_2$ durch $-c_1$ ersetzen, der Werth:
\weq{1}{(\Delta, f)^4 = (ab)^2(ac)^2(bc)^2 = j.}
Zur Berechnung von $(f,\Delta)^3$ entwickeln wir ein Glied der dritten Polare von $\Delta$ nach der Reihe (IX) Nr. 83. Wir erhalten, da $k$ nur die Werthe 0, 1 annehmen kann:
\[
\begin{aligned}
(ab)^2 b_x a_y^2 b_y &= (ab)^2 \sum \frac{\binom{1}{k}\binom{2}{k}}{\binom{1+1+2-k+1}{k}} \cdot (\overline{ab})^{k}_{y^{2-k+1}} (xy)^k \\
&= \left[(ab)^2 a_x^2 b_x^2\right]_{y^3} + \frac{2}{4}\left[(ab)^3 a_x b_x\right](xy). % ?? uncertain: exact stacked exponent arrangement on the barred (ab) term in the line above
\end{aligned}
\]
Weil aber $(ab)^3 a_x b_x$ identisch null, und $(ab)^2 a_x^2 b_x^2 = \Delta$, so hat man
\weq{I}{(ab)^2 b_x a_y^2 b_y = \Delta_{y^3} = \Delta_x \Delta_y^3.}
Ersetzen wir wiederum $y_1$ durch $c_2$, $y_2$ durch $-c_1$ und multipliciren mit $c_x$, so kommt:
\weq{2}{(ab)^2 (ac)^2 (bc)\, b_x c_x = (\Delta, f)^3 = 0,}
%CONT>>


% ===== pdf idx 106, printed 92.0 =====
% !PRINTED 92
\weberpage{92}
%<<CONT
weil die linke Seite durch Vertauschung von $b$ mit $c$ nur ihr Zeichen \"andert. Vertauscht man in (I) $y$ mit $x$, so ergiebt sich aus dieser dritten Polare von $\Delta$ die erste, und damit durch die gleiche Substitution statt der dritten Ueberschiebung $(\Delta, f)^3$ die erste:
\weq{3}{(ab)^2 (bc)\, a_x^2 b_x c_x^3 = (\Delta, f) = t.}
Sie verschwindet nicht, und liefert somit eine Covariante sechsten Grades von $f = a_x^4$.

Es er\"ubrigt nur noch die zweite Ueberschiebung von $f$ \"uber $\Delta$ zu berechnen. Wir k\"onnten zu dem Zwecke wiederum ein Glied der zweiten Polare von $\Delta$ in eine Reihe entwickeln. Vortheilhafter ist es aber, diesen Process f\"ur ein Glied der dritten Polare von $(f,f)$ vorzunehmen, n\"amlich f\"ur das Glied $(ab)a_x^3 b_y^3$. Da diese Form $(f,f)$ sowohl, als auch $(f,f)^3$ identisch verschwinden, so ist in der durch Formel (VIII) dictirten Entwicklung
\[
(ab)\,a_x^3 b_y^3 = (ab) \sum \frac{\binom{3}{k}\binom{3}{k}}{\binom{7-k}{k}} (\overline{ab})_{y^{3-k}} (xy)^k
\]
die erste und dritte Elementarcovariante null, und sie reducirt sich demnach auf die f\"ur $k=1$ und $k=3$ sich ergebenden Glieder:
\[
\begin{aligned}
(ab)\,a_x^3 b_y^3 &= \frac{3}{2}\left\{(ab)^2 a_x^2 b_x^2\right\}_{y^2} (xy) + \frac{1}{4}(ab)^4 (xy)^3 \\
&= \frac{3}{2}\Delta_{y^2}\cdot (xy) + \frac{1}{4}(ab)^4 (xy)^3.
\end{aligned}
\]
Ersetzen wir hierin $y_1$ durch $-c_2$, $y_2$ durch $+c_1$ und multipliciren mit $c_x$, so kommt:
\[
-(ab)(bc)^3 a_x^3 c_x = \frac{3}{2}(\Delta, f)^2 + \frac{i}{4} \cdot f,
\]
wenn wir die Invariante $(ab)^4$ von $f$ wie schon fr\"uher mit $i$ bezeichnen. Vertauscht man links $b$ mit $c$, nimmt die halbe Summe der so entstehenden symbolischen Producte und ber\"ucksichtigt, dass
\[
(ab)c_x - (ac)b_x = -(bc)a_x,
\]
so kommt endlich
\[
\frac{i}{2} f = \frac{3}{2}(\Delta, f)^2 + \frac{i}{4} f,
\]
oder:
\weq{4}{(\Delta, f)^2 = \frac{i}{6} \cdot f.}
Wir erhalten somit das Resultat:
\[
(\Delta, f) = t, \quad (\Delta, f)^2 = \frac{i}{6} \cdot f, \quad (\Delta, f)^3 = 0, \quad (\Delta, f)^4 = j.
\]


% ===== pdf idx 107, printed 93.0 =====
% !PRINTED 93
\weberpage{93}
Es möge hier bemerkt sein, dass die fünf Formen
\[
f,\ t,\ \Delta,\ i,\ j
\]
überhaupt die einzigen selbstständigen Covarianten der Form $f$ sind (wie wir später noch beweisen werden), durch welche sich alle andern rational und ganz darstellen lassen.

90. \textit{Die aszygetischen Covarianten dritten Grades in den Coefficienten von $f=a_x^5$.} Wir stellen uns weiter die Aufgabe, alle aszygetischen, d. i. linear von einander unabhängigen Covarianten dritten Grades in den Coefficienten einer Form fünfter Ordnung zu berechnen. Sie müssen jedenfalls durch symbolische Producte von der Form:
\weq{I}{P = (ab)^\mu (bc)^\varkappa (ca)^\lambda a_x^{\varkappa_1} b_x^{\lambda_1} c_x^{\mu_1}}
repräsentirt werden können, wobei die Exponenten den Gleichungen
\[
5 = \lambda + \mu + \varkappa_1 = \mu + \varkappa + \lambda_1 = \varkappa + \lambda + \mu_1
\]
genügen müssen.

Jedes solche Product $P$ lässt sich aber durch Ueberschiebungen von Formen zweiten Grades in den Coefficienten mit solchen ersten Grades darstellen, wie aus unsern allgemeinen Betrachtungen Nr. 45 und Nr. 46 erhellt. Die einzige Form ersten Grades ist aber die Originalform $f=a_x^5$ selbst, und die Formen zweiten Grades sind die aus $(f, f)^\varrho$ für $\varrho = 2$ und 4 sich ergebenden nicht verschwindenden Covarianten:
\[
\varphi = (ab)^2 a_x^3 b_x^3,\ \ i = (ab)^4 a_x b_x.
\]
Jede Form $P$ muss sich demnach, von $P = f^3 = a_x^5 b_x^5 c_x^5$ abgesehen, durch eine Summe von der Gestalt
\weq{II}{P = \sum m_i \cdot (f,\ \varphi)^\varrho + \sum n_i (f,\ i)^\sigma,}
wo $m_i$, $n_i$ noch zu bestimmende Constante sind, darstellen lassen, und zu solchen Summen gelangen wir, wenn wir nach bereits früher gegebenen Methoden das Symbol $c$ in $P$ zunächst durch $y$ ersetzen, und nun die so entstehende Form $P_y$ in eine Reihe entwickeln. Die Coefficienten der Potenzen von $(xy)$ sind dann Polaren von $\varphi$ oder $i$. Da nun in (I) kein Exponent die Zahl 4 aus bekannten Gründen überschreiten kann, so giebt es eine endliche Zahl von linearen Relationen (II). Aus ihnen eliminiren wir so viele Formen dritten Grades $P$, $(f,\ \varphi)^\varrho$, $(f,\ i)^\sigma$, als nur möglich, und gelangen so zu den linear unabhängigen Formen dritten Grades in den Coefficienten.

91. In unserem Falle können wir indess die aszygetischen Formen dritten Grades a priori erkennen, so dass die Reihenentwicklungen, welche in allgemeineren Fällen zur Bildung der Relationen (II) nöthig
%CONT>>


% ===== pdf idx 108, printed 94.0 =====
% !PRINTED 94
\weberpage{94}
%<<CONT
werden, sich hier auf ein Minimum beschränken. Denn da die sechs Formen:
\[
f^3,\ f\cdot\varphi,\ (f,\ \varphi),\ f\cdot i,\ (f,\ i),\ (f,\ i)^2
\]
von immer andern und andern Grade 15, 11, 9, 7, 5, 3 in den Variabeln sind, so können zwischen ihnen keine \gesperrt{linearen Relationen} bestehen. Unsere Aufgabe besteht also nur darin, zu zeigen, dass in der That keine weiteren Formen von der erwähnten Eigenschaft existiren.

Zum Zwecke dieses Beweises schreiten wir von den einfachsten Producten $P$ allmälig zu der Berechnung complicirterer vor, indem wir den Exponenten $\mu$, $\varkappa$, $\lambda$ der Klammerfactoren von $P$ immer grössere und grössere Zahlenwerthe beilegen. Dabei können wir schon von vorn herein eine Grenze angeben, über welche hinaus diese Grössen $\mu$, $\varkappa$, $\lambda$ nicht wachsen können. Nehmen wir nämlich an --- was bei der Symmetrie in den Symbolen $a$, $b$, $c$ stets erlaubt ist ---
\[
P = (ab)^\mu (bc)^\varkappa (ca)^\lambda a_x^{\varkappa_1} b_x^{\lambda_1} c_x^{\mu_1}
\]
sei so geordnet, dass
\[
\mu > \varkappa \geqq \lambda,
\]
dann muss, wenn
\[
\mu + \varkappa + \lambda > 6
\]
ist, $\mu$ mindestens den Werth 3 besitzen. Jedes symbolische Product $P$ mit dem Factor $(ab)^3$ kann aber (vgl. Nr. 11) stets auf eines mit dem Factor $(ab)^4$ reducirt werden. Nach unsern allgemeinen Ueberschiebungssätzen lässt sich aber alsdann eine solche Form $P$ dritten Grades linear durch Ueberschiebungen darstellen mit $i = (ab)^4 a_x b_x$, und einer Form ersten Grades, also der Form $f$, d. h. $P$ lässt sich linear durch $i\cdot f$, $(i,\ f)$, $(i,\ f)^2$ ausdrücken, und ist somit keine aszygetische Form zu diesen dreien.

Die Summe $\mu + \varkappa + \lambda$ der drei Exponenten kann also im günstigsten Falle den Werth 6 erreichen, und dabei kann $\mu$ oder $\varkappa$ oder $\lambda$ den Werth 4 nicht überschreiten.

92. Wir wollen im Folgenden die nach diesen Beschränkungen übrig bleibenden Möglichkeiten untersuchen.

1) Es sei: $\mu + \varkappa + \lambda = 0$, d. h. $\mu = \varkappa = \lambda = 0$, dann ist:
\[
P_0 = a_x^5 b_x^5 c_x^5 = f^3.
\]

2) Ist $\mu + \varkappa + \lambda = 1$, so lässt sich ein hiezu gehöriges symbolisches Product $P_1$ immer reduciren (vgl. Nr. 11) auf ein Product $P_2$, für welches

3)
\[
\mu + \varkappa + \lambda = 2,
\]
%CONT>>


% ===== pdf idx 109, printed 95.0 =====
% !PRINTED 95
\weberpage{95}
%<<CONT
und zwar $\mu = 2$, $\varkappa = \lambda = 0$ ist. In diesem Falle wird
\[
P_2 = (ab)^2 a_x^3 b_x^3 c_x^5 = \varphi \cdot f.
\]

4) Wenn $\mu + \varkappa + \lambda = 3$, so kann entweder
\[
\mu = k = \lambda = 1 % [sic] printed \textquotedbl{}k\textquotedbl{} for \varkappa
\]
oder
\[
\mu = 2,\ \varkappa = 1,\ \lambda = 0
\]
oder
\[
\mu = 3,\ \varkappa = \lambda = 0.
\]
Die nach dem letzten der drei Fälle zu bildende Form ist identisch null. Die Form des ersten Falles dagegen lässt sich auf die des zweiten Falles reduciren, für welchen
\[
P_3 = (ab)^2 (bc) a_x^3 b_x^2 c_x^4.
\]
Ersetzt man hierin nach der allgemein gegebenen Methode $c$ durch $y$, und entwickelt nun $P_{3y}$ in eine Reihe, so kommt:
\[
(ab)^2 a_x^3 b_x^2 b_y = [(ab)^2 a_x^3 b_x^3]_y + \text{Const.} \times [(ab)^3 a_x^2 b_x^2] (xy) \cdot
\]
Die zweite Elementarcovariante rechts ist identisch Null, die erste ist $\varphi$; also wird, wenn man wieder $y$ durch $c$ ersetzt:
\[
P_3 = (\varphi,\ f).
\]

5) Ist ferner: $\mu + \varkappa + \lambda = 4$, so ist entweder
\[
\mu > 2,\quad \text{oder}\quad \mu = 2,\ \varkappa = \lambda = 1
\]
\[
\mu = 2,\ \varkappa = 2.
\]
Im ersten Falle ist entweder $P$ direct gleich $f \cdot i$, wenn $\mu = 4$, $k = \lambda = 0$ % [sic] printed \textquotedbl{}k\textquotedbl{} for \varkappa
oder, wenn $\mu = 3$, $k = 1$, $\lambda = 0$: % [sic] printed \textquotedbl{}k\textquotedbl{} for \varkappa
\[
P_4 = (ab)^3 (ac) a_x b_x^2 c_x^4.
\]
Vertauscht man $a$ mit $b$ und nimmt die halbe Summe, so kommt:
\[
P_4 = \tfrac{1}{2} (ab)^3 a_x b_x c_x^4 \{(ac) b_x - c_x (ab)\}
\]
\[
= \tfrac{1}{2} f \cdot i.
\]
Im zweiten Falle ist:
\[
P_5 = (ab)^2 (ac) (bc) a_x^2 b_x^2 c_x
\]
\[
= \tfrac{1}{3} a_x^2 b_x^2 c_x^2 (ab)(ac)(bc) \{c_x (ab) + a_x (bc) + b_x (ca)\} = 0.
\]
Im dritten Falle ist
\[
P_6 = (ab)^2 (ac)^2 a_x b_x^3 c_x^3.
\]
Um diese Form durch Ueberschiebungen auszudrücken, entwickeln wir $(ac) a_y^3 a_x c_x^4$ in eine Reihe und erhalten:
\[
(ac) a_y^3 a_x c_x^4 = [(ac) a_x^4 c_x^4]_{y^3} + c_1 [(ac)^2 a_x^3 c_x^3]_{y^2} (yx) + c_2 [(ac)^3 a_x^2 c_x^2]_y (yx)^2 + c_3 [(ac)^4 a_x c_x] (yx)^3.
\]


% ===== pdf idx 110, printed 96.0 =====
% !PRINTED 96
\weberpage{96}
Hierin verschwindet der erste und dritte Term rechts identisch. Ersetzen wir also in dieser Reihe $y$ durch $b$, so kommt
\[
(ab)^3 (ac) a_x b_x^2 c_x^4 = c_1 (\varphi,\ f)^2 + c_3 \cdot i \cdot f.
\]
Die linke Seite ist $P_4 = \tfrac{1}{2} f \cdot i$; der erste Term rechts ist aber bis auf einen Zahlenfactor $P_6 = (ab)^2 (ac)^2 a_x^3 b_x^3 c_x^3$. Denn die Polare
\[
[(ac)^2 a_x^3 c_x^3]_{y^2} = \varphi_{y^2} = \frac{3}{\binom{6}{2}} \left\{ 2 (ac)^2 c_x^3 a_x a_y^2 + 3 (ac)^2 a_x^2 c_x^2 a_y c_y \right\}
\]
geht durch die Substitution $y_1 = b_2$, $y_2 = -b_1$ in die Ueberschiebung über:
\[
(\varphi,\ f)^2 = \frac{1}{5} \left\{ 2 (ac)^2 c_x^3 a_x (ab)^2 b_x^3 + 3 (ac)^2 a_x^2 c_x^2 (ab)(cb) b_x^3 \right\}
\]
\[
= \frac{1}{5} \{ 2 P_6 + 3 P_5 \},\ \text{oder weil } P_5 = 0,
\]
\[
= \frac{2}{5} P_6.
\]
Demnach ist:
\[
P_6 = \varrho \cdot (f,\ \varphi)^2 = k \cdot i \cdot f,
\]
wo $\varrho$ und $k$ numerische Coefficienten. Zwischen den beiden Ueberschiebungen $(f,\ \varphi)^2$ und $i \cdot f$ besteht also eine lineare Relation.

6) Nehmen wir ferner an:
\[
\mu + \varkappa + \lambda = 5,
\]
so ist entweder:
\[
\mu = 5,\ \varkappa = 0,\ \lambda = 0,\ \text{also}\ P_7 = (ab)^5 c_x^5 = 0
\]
oder
\[
\mu = 4,\ \varkappa = 1,\ \lambda = 0,\ \text{also}\ P_8 = (ab)^4 (ac) b_x c_x^4 = (i,\ f)
\]
oder
\[
\mu = 3,\ \varkappa = 1,\ \lambda = 1,\ \text{also}\ P_9 = (ab)^3 (ac) (bc) a_x b_x c_x^3 = 0,
\]
da $P_9$ durch Vertauschung von $a$ mit $b$ sein Zeichen ändert,

oder
\[
\mu = 3,\ \varkappa = 2,\ \lambda = 0,\ \text{also}\ P_{10} = (ab)^3 (ac)^2 b_x^2 c_x^3.
\]
Dieses Product geht, wenn man $b$ mit $c$ vertauscht und die halbe Summe nimmt, in das Product
\[
P_{11} = (ab)^2 (ac)^2 (bc) b_x^2 c_x^2 a_x
\]
über, welches dem Werthsystem $\mu = 2$, $\varkappa = 2$, $\lambda = 1$ entspricht. Die Form $P_{11}$ verschwindet aber identisch, da sie nur ihr Zeichen ändert, bei Vertauschung von $b$ mit $c$.

7) Ist endlich:
\[
\mu + \varkappa + \lambda = 6,
\]
so sind die einzig möglichen Fälle:
\[
\mu = 4,\ \varkappa = 1,\ \lambda = 1,\ \text{also}\ P_{12} = (ab)^4 (ac) (bc) c_x^3 = (i,\ f)^2
\]
\[
\mu = 3,\ \varkappa = 2,\ \lambda = 1,\ \text{also}\ P_{13} = (ab)^3 (ac)^2 (bc) b_x c_x^2
\]
\[
\mu = 2,\ \varkappa = 2,\ \lambda = 2,\ \text{also}\ P_{14} = (ab)^2 (ac)^2 (bc)^2 a_x b_x c_x.
\]


% ===== pdf idx 111, printed 97.0 =====
% !PRINTED 97
\weberpage{97}
Die beiden Formen $P_{13}$ und $P_{14}$ führen auf dieselbe Ueberschiebung $(i,f)^2$ wie $P_{12}$. Denn vertauscht man in $P_{13}$ $a$ mit $b$ und nimmt die halbe Summe, so kommt:

\[
P_{13}=\frac{1}{2}(ab)^3(ac)(bc)c_x^2\{(ac)b_x-(bc)a_x\}=\frac{1}{2}(ab)^4(ac)(bc)c_x^3=\frac{1}{2}(i,f).
\]

Vertauscht man aber in $P_{13}$ $b$ mit $c$, so kommt durch Addition beider Producte:

\[
P_{13}=\frac{1}{2}(ab)^2(ac)^2(bc)b_xc_x\{(ab)c_x-(ac)b_x\}=\frac{1}{2}(ab)^2(ac)^2(bc)^2a_xb_xc_x,
\]
oder
\[
P_{14}=2P_{13}=(i,f)=(ab)^2(ac)^2(bc)^2a_xb_xc_x.
\]

Wir sind am Schlusse angelangt. Die sechs Formen $f^3$, $i\cdot f$, $(i,f)$, $(i,f)^2$, $\varphi\cdot f$, $(\varphi,f)$ sind in der That die einzigen aszygetischen Covarianten von $f$, die in den Coefficienten vom dritten Grade sind.

93. \textit{Reciprocitätsgesetz von Hermite.} Es ist nun sehr bemerkenswerth, das jedem aszygetischen System von Covarianten $P_i$ einer Form $f$ stets ein gleichfalls aszygetisches System von Covarianten $Q_i$ einer andern gewissen Form $\varphi$ entspricht, und zwar ist das Entsprechen nach Hermite durch folgendes Gesetz definirt:

\glqq Besitzt die Form $n^{\mathrm{ten}}$ Grades $f=a_x^n$ eine Covariante $P_i$ vom Grade $m$ in den Coefficienten, so besitzt die Form $m^{\mathrm{ten}}$ Grades $\varphi=A_x^m$ eine entsprechende Covariante $Q_i$ vom Grade $n$ in den Coefficienten und beide Covarianten sind von gleichem Grade $\lambda$ in $x$.\grqq{}

So entsprechen den vorhin ermittelten sechs aszygetischen Formen $P_i$ dritten Grades in den Coefficienten

\[
f^3,\ f\cdot\varphi,\ (f,\varphi),\ f\cdot i.\ (f,i),\ (f,i)^2
\]

der Form $f=a_x^5$ ebenfalls sechs aszygetische Formen $Q_i$ fünften Grades
%CONT>>


% ===== pdf idx 112, printed 98.0 =====
% !PRINTED 98
\weberpage{98}
%<<CONT
in den Coefficienten:

\[
\varphi^5,\ \varphi^3\cdot\Delta,\ \varphi^2\cdot Q,\ \varphi\cdot\Delta^2,\ Q\cdot\Delta,\ \varphi\cdot R
\]

der Form $\varphi=A_x^3$, wobei, wie wir früher gesehen haben (vgl. Nr. 37)

\[
\Delta=(AB)^2A_xB_x,\quad Q=(A\Delta)A_x^2\Delta_x,\quad R=(\Delta,\Delta)^2.
\]

Der Beweis dieses Reciprocitätsgesetzes kann auf folgende Weise geführt werden.

Es sei $P$ eine Covariante vom Grade $m$ in den Coefficienten von $f=a_x^n=b_x^n=$ etc. und vom Grade $\lambda$ in den Variabeln $x$. Es sei ferner $\varphi=A_x^m$ eine Form $m^{\mathrm{ten}}$ Grades in $x$, welche die $m$ Wurzeln $\alpha^{(k)}=\dfrac{\alpha_1^{(k)}}{\alpha_2^{(k)}}$ besitzen möge, und demgemäss sich auch in der Gestalt schreiben lässt (von einem constanten Factor abgesehen):

\[
\varphi=(\alpha_2^{(1)}x_1-\alpha_1^{(1)}x_2)(\alpha_2^{(2)}x_1-\alpha_1^{(2)}x_2)\cdots(\alpha_2^{(m)}x_1-\alpha_1^{(m)}x_2),
\]
oder
\[
\varphi=\alpha_{1x}\cdot\alpha_{2x}\cdot\alpha_{3x}\cdots\alpha_{mx}=A_x^m.
\]

Ersetzen wir alsdann in dem symbolischen Producte $P$ das Symbol $a$ durch $\alpha_1$, $b$ durch $\alpha_2$, \dots, $m$ durch $\alpha_m$, so geht dadurch $P$ zunächst in eine simultane Covariante $\overline{Q}$ der $m$ linearen Factoren $\alpha_{kx}$ über. Nun wird $P$ nicht geändert, wenn wir in ihm irgend welche zwei gleichberechtigte Symbole $a$ und $b$ vertauschen. Permutiren wir entsprechend in $\overline{Q}$ alle Wurzeln $\overline{\alpha}_k$, so gehört immer noch jedes $\overline{Q}$ zur nämlichen Form $P$. Aber die Summe aller dieser $\mu!$ Formen $\overline{Q}$ ist nun eine symmetrische Function der Wurzeln $\overline{\alpha}_k$ und lässt sich demnach rational in den Coefficienten von $\varphi$ darstellen. Da anderntheils jedes Glied $\overline{Q}$ gemäss seiner Ableitung Covarianteneigenschaft besass, so muss sie auch die Summe besitzen, d. h.

\[
Q=\sum{}^{\prime}\overline{Q}
\]

ist rationale Covariante von $\varphi=A_x^m$, und zwar vom Grade $\lambda$ in $x$, weil die Zahl der Factoren erster Art bei diesen Operationen unberührt blieb und vom Grade $n$ in den Coefficienten von $\varphi$, weil jede Wurzel $\overline{\alpha}_k$ in jedem $\overline{Q}$ gleich oft und zwar $n$ mal auftritt. (Vergl. Salmon, Lineare Transformationen, Art. 57 und 58, über Grad und Gewicht symmetr. Functionen der Wurzeln.) Es entspricht also in der That jeder Form $P$ von $f$ eine Form $Q$ von $\varphi$, und insbesondere einem aszygetischen System $P_i$ wieder ein aszygetisches System $Q_i$. So entspricht z. B. der Invariante

\[
j=(ab)^2(ac)^2(bc)^2
\]

dritten Grades in den Coefficienten der Form $f=a_x^4$ die Invariante
\[
R=(\alpha_1\alpha_2)^2(\alpha_1\alpha_3)^2(\alpha_2\alpha_3)^2,
\]
welche vierten Grades in den Coefficienten der Form
\[
\varphi=\alpha_{1x}\alpha_{2x}\alpha_{3x}=A_x^3
\]
ist, deren drei Wurzeln die Werthe $\alpha^{(1)}=\dfrac{\alpha_1^{(1)}}{\alpha_2^{(1)}}$, $\alpha^{(2)}=\dfrac{\alpha_1^{(2)}}{\alpha_2^{(2)}}$, $\alpha^{(3)}=\dfrac{\alpha_1^{(3)}}{\alpha_2^{(3)}}$ % ?? uncertain: the relation sign before the third fraction is faint/broken in the scan and reads more like a short dash than a full equals sign; transcribed as = by inference from the parallel construction
besitzen.

94. \textit{Darstellung einer in den Wurzeln von $\varphi=A_x^m$ ausgedrückten Covariante durch ein symbolisches Product; Beispiel.} Wir haben schon in \S.\,1 Nr. 8 darauf hingewiesen, dass jede symmetrische simultane Covariante der $m$ linearen Factoren $\alpha_{ix}$ von $\varphi$ auch eine \gesperrt{rationale} Covariante der Form $\varphi$ ist, und haben den einfachen Beweis dieser Thatsache eben gegeben. Die Reihenentwicklungen für Formen mit
%CONT>>


% ===== pdf idx 113, printed 99.0 =====
% !PRINTED 99
\weberpage{99}
%<<CONT
mehreren Reihen cogredienter Variablen setzen uns in den Stand, eine solche in den Wurzeln dargestellte Covariante direct in ein symbolisches Product umzuformen, was man allerdings auch gemäss der Theorie der symmetrischen Functionen auf einem Umwege über die unsymbolische Darstellung der Covariante erreichen könnte. Ehe wir die Methode allgemein entwickeln, will ich sie an einem Beispiele vorbereiten, und benutze hierzu die bereits vorhin angeführte simultane symmetrische Invariante

\[
R=(\alpha_1\alpha_2)^2(\alpha_1\alpha_3)^2(\alpha_2\alpha_3)^2
\]

der drei linearen Factoren $\alpha_{kx}$ von

\[
\varphi=\alpha_{1x}\alpha_{2x}\alpha_{3x}=A_x^3=B_x^3=\text{etc.}
\]

Wir wissen: $\varphi$ besitzt nur eine Invariante vierten Grades, nämlich die Discriminante (vergl. Nr. 37 und 39)

\[
R_1=(\Delta,\Delta)^2=(AB)^2(DC)^2(BC)(AD),
\]

und wir wollen nun zeigen, wie sich successive $R$ in $R_1$ überführen lässt. Zu dem Zwecke setzen wir zunächst:

\weq{1}{\alpha_{2x}\alpha_{3x}=p_x^2=p}

und erhalten so eine quadratische Hilfsform, deren Invariante durch

\weq{2}{(pp_1)^2=(\alpha_{2x}\alpha_{3x},\ \alpha_{2x}\alpha_{3x})^2=-\frac{1}{2}(\alpha_2\alpha_3)^2}

dargestellt ist. Unter Benutzung dieser Werthe (1) und (2) wird:

\weq{3}{\varphi=p_x^2\alpha_{1x}=A_x^3,}

\weq{4}{R=-2(pp_1)^2(\alpha_1\alpha_2)^2(\alpha_1\alpha_3)^2.}

Ueberschieben wir $p$ auch zweimal über $\alpha_{1x}$, so kommt:

\weq{5}{(p\alpha_1)^2=(\alpha_1\alpha_2)(\alpha_1\alpha_3)}

und demnach lässt sich Gleichung (4) auch schreiben:

\weq{6}{R=-2(pp_1)^2(p_2\alpha_1)^2(p_3\alpha_1)^2.}

In diesem Ausdrucke können wir nun successive die Symbole $A$, $B$, etc. der Form $\varphi$ einführen. Denn überschieben wir Gleichung (3) zweimal über $\alpha_{1x}$, so kommt:

\weq{7}{(p_2\alpha_1)^2\alpha_{1x}=3(A\alpha_1)^2A_x,}

und wenn wir hierin $x$ durch $p_3$ ersetzen:

\weq{8}{(p_2\alpha_1)^2(\alpha_1p_3)=3(A\alpha_1)^2(Ap_3);}

also:

\weq{9}{R=6(pp_1)^2(p_3\alpha_1)(A\alpha_1)^2(Ap_3).}

Damit ist eine erste Reduction von $R$ erzielt. Wir wollen nun weiter das Symbol $p_3$ eliminiren, und bilden dementsprechend die erste
%CONT>>


% ===== pdf idx 114, printed 100.0 =====
% !PRINTED 100
\weberpage{100}
%<<CONT
Ueberschiebung von $p_x^2\alpha_{1x}$ über $\alpha_{1x}$:

\weq{10}{2p_x(p\alpha_1)\alpha_{1x}=3A_x^2(A\alpha_1)=3B_x^2(B\alpha_1).}

Ersetzen wir hierin $x$ durch $A$, so kommt:

\[
2(Ap)(p\alpha_1)(A\alpha_1)=3(AB)^2(B\alpha_1),
\]

und demnach:

\[
R=9(pp_1)^2(A\alpha_1)(AB)^2(B\alpha_1),
\]

oder, weil $\Delta_x^2=(AB)^2A_xB_x$, und demnach $(\Delta\alpha_1)^2=(AB)^2(A\alpha_1)(B\alpha_1)$, so wird:

\[
R=9(pp_1)^2(\Delta\alpha_1)^2.
\]

Da aber nach dem Identitätssatze:

\[
\{(pp_1)(\Delta\alpha_1)\}^2=\{(p\Delta)(p_1\alpha_1)-(p_1\Delta)(p\alpha_1)\}^2
\]
\[
=2(p\Delta)^2(p_1\alpha_1)^2-2(p\Delta)(p_1\Delta)(p\alpha_1)(p_1\alpha_1),
\]

so wird endlich:

\weq{11}{R=18(p\Delta)^2(p_1\alpha_1)^2-18(p\Delta)(p_1\Delta)(p\alpha_1)(p_1\alpha_1).}

Für die beiden Terme rechts können wir nun aber zwei Relationen aufstellen, mit Hilfe deren wir das gewünschte Ziel erreichen. Der erste Term geht nämlich aus $p_y^2\alpha_{1x}$ hervor, wenn wir darin $y$ durch $\Delta$, $x$ durch $p_1$ ersetzen und mit $(\alpha_1p_1)$ multipliciren. Es ist aber

\[
p_y^2\alpha_{1x}=\{p_x^2\alpha_{1x}\}_{y2}+\frac{2}{3}\{(p\alpha_1)p_x\}_y(yx)
\]

nach Formel (VIII) Nr. 81, wobei nur $x$ mit $y$ vertauscht ist; und da nach Voraussetzung $p_x^2\alpha_x=B_x^3$, so wird, wenn wir die angegebene Substitution ausführen:

\[
(p\Delta)^2(p_1\alpha_1)^2=(B\Delta)^2(Bp_1)(\alpha_1p_1)+\frac{2}{3}(p\alpha_1)(p_1\alpha_1)(p\Delta)(p_1\Delta).
\]

Hierin ist der erste Term rechts identisch Null, da nach der Theorie der cubischen Formen jedes symbolische Product mit dem Klammerfactor $(B\Delta)^2$ verschwindet. [Vergl. Nr. 63, (5) u. Nr. 145, (I).] Demnach wird:

\weq{12}{(p\Delta)^2(p_1\alpha_1)^2=\frac{2}{3}(p\alpha_1)(p_1\alpha_1)(p\Delta)(p_1\Delta)}

und dies ist die erste der gesuchten Relationen. Um noch eine zweite zu erhalten, berücksichtigen wir, dass der zweite Term in Gleichung (11) aus $p_y(p\alpha_1)\alpha_{1x}$ hervorgeht, wenn wir darin $y$ durch $\Delta$, $x$ durch $p_1$ ersetzen und mit $(p_1\Delta)$ multipliciren. Es ist aber

\[
(p\alpha_1)p_y\alpha_{1x}=\{(p\alpha_1)p_x\alpha_{1x}\}_y+\frac{1}{2}\{(p\alpha_1)^2(yx)\}
\]

nach Formel (VIII) Nr. 81, und weil nach Gleichung (10) die Form $(p\alpha_1)p_x\alpha_{1x}$ gleich $\frac{3}{2}(B\alpha_1)B_x^2$ ist, so wird nach Ausführung der Substitution
%CONT>>


% ===== pdf idx 115, printed 101.0 =====
% !PRINTED 101
\weberpage{101}
%<<CONT
\[
-(p_1\Delta)(p\alpha_1)(p\Delta)(p_1\alpha_1) = \frac{3}{2}(B\alpha_1)(B\Delta)(Bp_1)(p_1\Delta) - \frac{1}{2}(p\alpha_1)^2(p_1\Delta)^2,
\]
oder, wenn wir für den zweiten Term rechts seinen Werth aus (12)
substituiren:
\weq{13}{(p\alpha_1)(p_1\alpha_1)(p\Delta)(p_1\Delta) = -\frac{9}{4}(B\alpha_1)(B\Delta)(Bp_1)(p_1\Delta).}
Die rechte Seite dieser Gleichung erhält man aus der Reihenentwicklung
\[
\alpha_{1x}p_{1x}p_{1y} = \{\alpha_{1x}p_{1x}^2\}_y + \frac{1}{2}\{(\alpha_1p_1)p_{1x}\}(xy),
\]
wenn man darin $x$ durch $B$, $y$ durch $\Delta$ ersetzt und mit $(B\Delta)$ multiplicirt. Berücksichtigt man noch, dass $p_{1x}^2\alpha_{1x} = C_x^3 = \varphi$, so ergiebt sich:
\[
(B\alpha_1)(Bp_1)(p_1\Delta)(B\Delta) = (BC)^2(B\Delta)(C\Delta) - \frac{1}{2}(B\Delta)^2(\alpha_1p_1)(p_1B).
\]
Hierbei verschwindet wieder der zweite Term rechts wegen $(B\Delta)^2$; substituirt man den sich so für die linke Seite ergebenden Werth $(BC)^2(B\Delta)(C\Delta)$, der nichts anderes ist als $(\Delta,\Delta)^2$, in (13), so kommt als zweite Relation:
\weq{14}{(p\alpha_1)(p_1\alpha_1)(p\Delta)(p_1\Delta) = -\frac{9}{4}(\Delta,\Delta)^2,}
also wegen (12):
\[
(p\Delta)^2(p_1\alpha_1)^2 = -\frac{3}{2}(\Delta,\Delta)^2,
\]
und demnach wegen (11):
\weq{15}{R = +\frac{27}{2}(\Delta,\Delta)^2 = (\alpha_1\alpha_2)^2(\alpha_1\alpha_3)^2(\alpha_2\alpha_3)^2,}
womit unsere Aufgabe erledigt ist.

95. \textit{Allgemeiner Fall.} Schon dieses einfache Beispiel zeigt, dass die Aufgabe, eine in symmetrischen Functionen der Wurzeldifferenzen gegebene Covariante direct in ein symbolisches Product umzuformen, complicirter Natur ist. Noch ungleich schwieriger gestalten sich die Verhältnisse in allgemeineren Fällen. Indess wollen wir doch --- theoretisch wenigstens --- diese Frage erledigen, die auch anderweitig von Interesse ist. Es sei:
\[
Q = \sum{}' \overline{Q}
\]
die gegebene Covariante der Form $(n+1)^{\text{ten}}$ Grades in $x$
\[
\varphi = \alpha_{1x}\cdot\alpha_{2x}\cdot\alpha_{3x}\cdots\alpha_{n+1,x} = A_x^{n+1} = B_x^{n+1} = \text{etc.}
\]
Dabei ist jedes Glied $\overline{Q}$ der Summe rechts ein Product von der Gestalt
\[
\overline{Q} = (\alpha_i\alpha_k)^{\lambda_1}(\alpha_i\alpha_\mu)^{\lambda_2}\cdots\alpha_{ix}^{\lambda_1}\alpha_{kx}^{\lambda_2}\alpha_{\mu x}^{\lambda_3}\cdots,
\]
wo jede Wurzel $\alpha_i$ in demselben Grad auftritt. Irgend eines der Glieder
%CONT>>


% ===== pdf idx 116, printed 102.0 =====
% !PRINTED 102
\weberpage{102}
%<<CONT
$\overline{Q}$ wollen wir als Repräsentant herausgreifen, an dem die Umformung gezeigt werden soll. Die Aufgabe ist alsdann, in $\overline{Q}$ successive durch Einführung von immer mehr und mehr Symbolen $A, B, C\ldots$ die Grössen $\alpha_i$ zu verdrängen. Hiezu benutzen wir wie im Beispiel die Hilfsform $n^{\text{ten}}$ Grades
\[
p = \alpha_{2x}\alpha_{3x}\alpha_{4x}\cdots\alpha_{n+1,x} = a_x^n = b_x^n = \text{etc.},
\]
welche sich von $\varphi$ nur um den Factor $\alpha_{1x}$ unterscheidet. Denken wir uns nun einen Augenblick in $Q = \sum{}' \overline{Q}$ die Wurzel $\alpha_1$ durch $y$ ersetzt, so ist die entstehende Form jedenfalls Covariante von $p$, und wir können annehmen, dass wir die Aufgabe, die wir für Formen $(n+1)^{\text{ten}}$ Grades lösen wollen, für Formen $n^{\text{ten}}$ Grades bereits erledigt haben, dass es uns also möglich ist, alle in symmetrischen Functionen der Wurzeln gegebenen Covarianten der Form niedrigeren Grades $p = p_x^n$ durch die Symbole $a, b, c\cdots$ von $p$ darzustellen. Ersetzen wir alsdann wiederum $y$ durch $\alpha_1$, so ist nunmehr $Q = \sum{}' \overline{Q}$ ausgedrückt in $\alpha_1$ und Symbolen $a, b, c$, d.\,h. $Q$ ist dargestellt als simultane Covariante der Formen $\alpha_{1x}$ und $p_x^n$. Die eben gemachte Annahme ist erlaubt, denn für Formen ersten Grades können wir sie jederzeit sofort realisiren.

96. Es handelt sich nun darum, vermöge der Identität
\weq{1}{A_x^{n+1} = a_x^n \cdot \alpha_{1x}}
die Symbole $a, b, c$ durch Symbole $A$ zu ersetzen und damit die Grösse $\alpha_1$ von selbst successive zu eliminiren. Zu dem Zwecke stellen wir zunächst nach dem in Nr.\,46 gegebenen Satze den Repräsentanten $\overline{Q}$ durch eine Summe von Ueberschiebungen dar
\weq{2}{\overline{Q} = \sum{}' (L,\,M)^\nu,}
wobei die Formen $L$ nur $\alpha_1$ und das Symbol $a$ enthalten, also die Gestalt
\weq{3}{L = (a\alpha_1)^k a_x^{n-k} \alpha_{1x}^\varrho}
besitzen, während die Formen $M$ irgend welche Covarianten von $p$ mit den Symbolen $b, c, d,\ldots$ bedeuten. Besitzt nun $L$ wirklich Factoren erster Art $\alpha_{1x}$, ist also $\varrho > 0$, so ergiebt sich durch $k$ malige successive Ueberschiebung von (1) über $\alpha_{1x}$:
\weq{4}{(n-k+1)(a\alpha_1)^k a_x^{n-k} \alpha_{1x} = (n+1)(A\alpha_1)^k A_x^{n-k+1},}
oder, wenn wir mit $\alpha_{1x}^{\varrho-1}$ multipliciren und mit (3) compariren:
\[
L = \frac{n+1}{n-k+1}(A\alpha_1)^k A_x^{n-k+1}\alpha_{1x}^{\varrho-1}.
\]
%CONT>>


% ===== pdf idx 117, printed 103.0 =====
% !PRINTED 103
\weberpage{103}
Dies ist eine erste Reductionsformel für alle jene Glieder der Gleichung (2), welche durch Ueberschiebung von Formen $L$ entstehen, die Factoren erster Art $\alpha_{1x}$ besitzen. Durch sie wird der Exponent von $\alpha_{1x}^\varrho$ um 1 vermindert und gleichzeitig das Symbol $A$ an Stelle von $a$ eingeführt. Besitzt dagegen eine solche Form $L$ keine Factoren erster Art $\alpha_{1x}$, in welchem Falle wir sie einstweilen \gesperrt{normirte Formen} $L$ nennen wollen, so kann für sie die Formel (4) nicht Reductionsformel sein; wir werden diese normirten Formen $L$ später reduciren. Alle Glieder $G_i$ der Summe (2) dagegen, welche durch Formel (4) reducirbar waren, können wir nun wieder nach dem schon einmal citirten Satze in Nr.\,46 darstellen in der Gestalt:
\[
G_i = \sum{}' (\overline{Q}_1,\, A_x^{n+1}),
\]
wobei nun mit $\overline{Q}_1$ wieder simultane Covarianten vom Typus $\overline{Q}$ bezeichnet sind, deren $L$ jedoch einen Factor $\alpha_{1x}$ weniger enthalten und auch niedrigeren Grades in den Coefficienten von $p$ sind. Diese Formen $\overline{Q}_1$ reduciren wir alsdann in der gleichen Weise wie die $\overline{Q}$ und indem wir diesen Process immer von Neuem wiederholen, werden dabei in die bezeichneten Glieder immer und mehr Symbole $A, B, C,\ldots$ eingeführt, während gleichzeitig ihr Grad sich successive reducirt, bis schliesslich überhaupt keine Factoren erster Art $\alpha_{1x}$ in den sich immer neuerdings ergebenden Formen $L$ auftreten.

97. Die weitere Reduction fällt alsdann zusammen mit der Reduction jener Glieder in (2), die aus normirten Formen $L$ durch Ueberschiebung gebildet werden können. Für diese ergiebt sich die Reduction auf folgende Weise. Wir stellen zunächst den Repräsentanten $\overline{Q}$, wenn in ihm bereits Symbole $A, B, C\ldots$ auftreten sollten, als Summe von Ueberschiebungen dar, gebildet von Formen $\overline{Q}_1$, welche nur $\alpha_1$ und Symbole $a, b, c\ldots$ enthalten, über Covarianten von $\varphi = A_x^{n+1}$. In den Formen $\overline{Q}_1$ trennen wir alsdann wieder $\alpha_1$ und die Symbole $a, b, c, \ldots$ in zwei Gruppen: eine erste, welche $\alpha_1$ und die Symbole $a$ und $b$ umfasst, eine zweite für die noch übrigen Symbole $c, d, e,\ldots$. Dann lässt sich $\overline{Q}_1$ in der Gestalt darstellen:
\[
\overline{Q}_1 = \sum{}' (N,\, R)^\nu,
\]
wo nun die Formen $N$ vom Typus
\[
N = (ab)^\mu (a\alpha_1)^k a_x^{n-k-\mu} b^{n-\mu}, \quad k > 1
\]
sind, während die Formen $R$ irgend welche Covarianten von $p$ bedeuten. Wie wir nun oben eine Reductionsformel für die Formen $L$ gegeben haben, so entwickeln wir jetzt eine solche für $N$. Die Form $N$
%CONT>>


% ===== pdf idx 118, printed 104.0 =====
% !PRINTED 104
\weberpage{104}
%<<CONT
geht nämlich bis auf den Factor $(-1)^{\mu+1}$ aus $b_y\alpha_{1y}b_x^{n-\mu}$ hervor, wenn wir darin $y$ durch $a$ ersetzen und mit $(a\alpha_1)^{k-1}a_x^{n-k-\mu}$ multipliciren. Nun ist aber nach (VIII) Nr.\,81
\[
b_y^\mu \alpha_{1y} b_x^{n-\mu} = \{b_x^n \alpha_{1x}\}_{y^{\mu+1}} + \frac{n-\mu}{n+1}(b\alpha_1)(xy)b_x^{n-\mu-1}b_y^\mu.
\]
Ersetzen wir darin die erste Elementarcovariante durch ihr Symbol $A_x^{n+1}$ nach (1) Nr.\,96 und führen die Substitution aus, so kommt:
\[
\begin{aligned}
(-1)^{\mu+1} N &= (Aa)^{\mu+1}(a\alpha_1)^{k-1} A_x^{n-\mu} a_x^{n-k-\mu} \\
&\quad - \frac{n-\mu}{n+1}(b\alpha_1)(a\alpha_1)^{k-1}(ba)^\mu b_x^{n-\mu-1} a_x^{n-k-\mu+1}.
\end{aligned}
\]
Das erste Glied rechts besitzt einen Factor $(a\alpha_1)$ weniger als $N$ und statt des Symboles $b$ das Symbol $A$, ist also reducirt gegenüber der Form $N$. Nicht aber das zweite Glied rechts, das eher complicirter genannt werden muss. Indess geht dieses zweite Glied, das wir, vom Factor $-\frac{n-\mu}{n+1}$ abgesehen, mit $N_1$ bezeichnen wollen, aus $(a\alpha)^{k-1} a_y^\mu a_x^{n-k-\mu+1}\alpha_y$ hervor, wenn wir darin $y$ durch $b$ ersetzen und mit $b_x^{n-\mu-1}$ multipliciren. Es ist aber nach (VIII) Nr.\,81
\[
(a\alpha_1)^{k-1} a_y^\mu \alpha_{1y} a_x^{n-k-\mu+1} = \{(a\alpha_1)^{k-1} a_x^{n-k+1} \alpha_{1x}\}_{y^{\mu+1}} + \frac{n-k-\mu+1}{n-k+3}(a\alpha_1)^k a_y^\mu a_x^{n-k-\mu}(xy).
\]
Ersetzen wir darin die erste Elementarcovariante vermöge der Gleichung (4) durch $\frac{n+1}{n-k}(A\alpha_1)^{k-1}A_x^{n-k}$ und führen alsdann die Substitution aus, so geht die linke Seite über in $N_1$ und das zweite Glied rechts in $N$. Schaffen wir dieses Glied auf die linke Seite, so kommt:
\weq{1}{(-1)^{\mu+1} N_1 - \frac{n-k-\mu+1}{n-k+3} N = \frac{n+1}{n-k}(A\alpha_1)^{k-1}(Ab)^{\mu+1} b_x^{n-\mu-1} A_x^{n-k-\mu}.}
Hiezu hatten wir:
\weq{2}{(-1)^{\mu+1} N - \frac{n-\mu}{n+1} N_1 = (Ab)^{\mu+1}(b\alpha_1)^{k-1} A_x^{n-\mu} b_x^{n-k-\mu},}
und diese beiden Gleichungen zusammen liefern in der That eine Reduction für die Formen $N$ und $N_1$. Durch fortgesetzte Wiederholung dieses Verfahrens wird hierdurch successive nicht nur die Anzahl der Klammerfactoren, welche $\alpha_1$ enthalten, vermindert, sondern auch die Anzahl der Symbole $a, b, c,\ldots$.

Diese beiden Reductionsformeln (1) und (2) führen so endlich neben Gliedern, die bereits völlig in den Symbolen $A, B, C,\ldots$ dargestellt sind, auf Glieder, die ausser den Symbolen $A, B, C,\ldots$
%CONT>>


% ===== pdf idx 119, printed 105.0 =====
% !PRINTED 105
\weberpage{105}
%<<CONT
nur mehr ein Symbol $a$ und einmal die Wurzel $\alpha_1$ enthalten, die
also entstanden gedacht werden können durch Ueberschiebung von
Formen
\weq{3}{L = (a\alpha_1)\,a_x^{n-1}}
mit symbolischen Producten in den Symbolen $A, B, C, \dots$ Auf sie
lassen sich Gleichungen (1) und (2) nicht mehr anwenden, da dieselben
die gleichzeitige Anwesenheit von mindestens zwei Symbolen
$a$ und $b$ voraussetzen. Allein diese Glieder verschwinden (vergl.
auch das Beispiel) identisch. Um dies theoretisch einzusehen, denke man
sich nur auch die übrigen Glieder $\bar Q$ in derselben Weise wie den
Repräsentanten $\bar Q$ behandelt. Dann scheidet sich schliesslich die
Covariante
\[
Q = \sum \bar Q
\]
in zwei Theile, deren erster nur mehr symbolische Producte in $A, B, C$
enthält und also für sich symmetrisch in den Wurzeln $\alpha_i$, und deren
zweiter Theil alle Ueberschiebungen mit den Formen $L$ vom Typus
(3) enthält. Da aber $Q$ selbst symmetrisch ist, so muss dieser Theil
ebenfalls für sich symmetrisch in den Wurzeln $\alpha_i$ sein, und kann also
seinen Werth nicht ändern bei Vertauschungen zweier Wurzeln. Ersetzen
wir nun in diesem Theile
\[
\sum (L, M)^\nu = \sum \big((a\alpha_1)\,a_x^{n-1}, M\big)^\nu
\]
$a_x^n$ durch seinen Wert $\alpha_{2x}\,\alpha_{3x}\,\alpha_{4x}\cdots\alpha_{n+1\,x}$, so kommt:
\[
\sum (L, M)^\nu = \sum \Big( \sum (\alpha_2\alpha_1)\,\alpha_{3x}\,\alpha_{4x}\cdots, M \Big)^\nu.
\]
Die rechte Seite ändert aber ihren Werth, wenn man $\alpha_2$ mit $\alpha_1$
vertauscht, d.\,h. sie kann nur eine verschwindende symmetrische
Function sein. Wir werden im dritten Bande eine Erweiterung der hier
gegebenen Entwicklungen auf ternäre Formen kennen lernen.

98. \textit{Ueber Covarianten dritten Grades in den Coefficienten einer
Form $f = a_x^n$.} Schon in dem ersten Paragraphen dieser Vorlesungen
haben wir gezeigt, dass ein symbolisches Product, dessen zumeist
auftretender Klammerfactor $(ab)^{2\mu-1}$ ist, sich stets durch eine Summe
anderer Producte darstellen lässt, die den Factor $(ab)^{2\mu}$ besitzen. Es
ist nun aber sehr bemerkenswerth, dass es auch eine Reihe symbolischer
Producte mit dem Factor $(ab)^{2\mu}$ giebt, welche in eine Summe
von andern Producten entwickelt werden können, die durchwegs
mindestens $(ab)^{2\mu+2}$ als Factor enthalten. Zu diesen gehören unter
andern alle geraden Ueberschiebungen:
\[
U = \big((f,f)^{2\mu}, f\big)^{2\mu+\varrho} \doteq \big((ab)^{2\mu}\,a_x^{n-2\mu}\,b_x^{n-2\mu},\, c_x^n\big)^{2\mu+\varrho},
\]
%CONT>>


% ===== pdf idx 120, printed 106.0 =====
% !PRINTED 106
\weberpage{106}
%<<CONT
wenn $\varrho \geqq 0$, $\mu \leqq \dfrac{n}{4}$ ist. Nur in einem besondern Falle, wenn
nämlich $n = 4\mu$ und $\varrho = 2\mu$ ist, wird diese Entwicklung unmöglich. Die
Ueberschiebung $U$ ist dann eine Invariante $j = (ab)^{2\mu}(ac)^{2\mu}(bc)^{2\mu}$.
Indem wir nun die Möglichkeit der erwähnten Entwicklungen von $U$
im allgemeinen Falle darthun, ergiebt sich alsdann der etwas allgemeinere
Satz:

\glqq Symbolische Producte mit dem Factor $(ab)^{2\mu}(ac)^{2\mu}$ lassen
sich stets — von etwaigen Gliedern mit dem Factor
\[
j = (ab)^{2\mu}(ac)^{2\mu}(bc)^{2\mu}
\]
abgesehen — in eine Summe anderer Producte entwickeln,
welche mindestens den Factor $(ab)^{2\mu+2}$ besitzen.\grqq{}

Der Beweis lässt sich folgendermassen durch Reihenentwicklung
führen. Das erste Glied der Ueberschiebung $U$ ist
\[
G_0 = (ab)^{2\mu}(ac)^{2\mu+\varrho}\,a_x^{n-4\mu-\varrho}\,b_x^{n-2\mu}\,c_x^{n-2\mu-\varrho}.
\]
Es entsteht aus dem Polarengliede: $(ab)^{2\mu}\,a_y^{2\mu+\varrho}\,a_x^{n-4\mu-\varrho}\,b_x^{n-2\mu}$
durch die Substitution $y_1 = c_2$, $y_2 = -c_1$; entwickeln wir aber
dasselbe in Reihe (VIII), so kommt:
\[
\begin{aligned}
\bar G_0 &= (ab)^{2\mu}\,a_y^{2\mu+\varrho}\,a_x^{n-4\mu-\varrho}\,b_x^{n-2\mu} \\
&= (f,f)^{2\mu}_{y^{2\mu+\varrho}} + \text{Const.}\times (f,f)^{2\mu+1}_{y^{2\mu+\varrho-1}}(yx) + \cdots,
\end{aligned}
\]
oder, weil rechts alle ungeraden Elementarcovarianten verschwinden:
\[
\bar G_0 = (f,f)^{2\mu}_{y^{2\mu+\varrho}} + \text{Const.}\times (f,f)^{2\mu+2}_{y^{2\mu+\varrho-2}}(yx)^2 + \cdots
\]
Jedes Glied rechts, vom ersten abgesehen, enthält den Factor
$(ab)^{2\mu+2}$, der auch erhalten bleibt, wenn wir nun $y_1$ durch $+c_2$ und
$y_2$ durch $-c_1$ ersetzen und mit $c_x^{n-2\mu-\varrho}$ multipliciren. Mit Rücksicht
auf den in Nr.\,43 erläuterten Modulbegriff können wir daher nach
ausgeführter Substitution diese letzte Gleichung schreiben:
\[
(ab)^{2\mu}(ac)^{2\mu+\varrho}\,a_x^{n-4\mu-\varrho}\,b_x^{n-2\mu}\,c_x^{n-2\mu-\varrho}
\equiv \text{Const.}\times \big((f,f)^{2\mu}, f\big)^{2\mu+\varrho} \bmod. (ab)^{2\mu+2}.
\]
Ist nun $\varrho > 0$, also mindestens gleich 1, so enthält nach dem in
Nr.\,11 gegebenen Satze auch das Glied links den Factor $(ab)^{2\mu+2}$, und
wir haben demnach als erstes Resultat:
\weq{I}{\big((f,f)^{2\mu}, f\big)^{2\mu+\varrho} \equiv 0 \bmod. (ab)^{2\mu+2}, \quad \mu < \dfrac{n}{4}, \quad \varrho > 0.}
Diese Gleichung gilt aber auch noch, wenn $\varrho = 0$, so lange nur
$n > 4\mu$, etwa $n = 4\mu + \lambda$. Zum Beweise gehen wir alsdann vom
symbolischen Producte $(ab)^{2\mu-1}\,a_y^{n-2\mu+1}\,b_x^{n-2\mu+1} = (ab)^{2\mu-1}\,a_y^{2\mu+\lambda+1}\,b_x^{2\mu+\lambda+1}$
%CONT>>


% ===== pdf idx 121, printed 107.0 =====
% !PRINTED 107
\weberpage{107}
%<<CONT
aus, und entwickeln dasselbe in Reihe (VIII). Dabei verschwinden
rechts wiederum alle ungeraden Ueberschiebungen, während die Elementarcovarianten
$(f,f)^{2\mu+2k}$, $k > 0$, von 0 verschieden sind.
Ersetzen wir alsdann wiederum $y_1$ durch $c_2$, $y_2$ durch $-c_1$ und multipliciren
mit $c_x^{n-2\mu-1}$, so kommt:
\[
(ab)^{2\mu-1}(ac)^{2\mu+\lambda+1}\,b_x^{2\mu+\lambda+1}\,c_x^{n-2\mu-1}
\doteq \text{Const.}\times \big((f,f)^{2\mu}, f\big)^{2\mu} \bmod. (ab)^{2\mu+2},
\]
oder, weil das Glied links wegen $(ac)^{2\mu+1}$ selbst $\equiv \bmod. (ab)^{2\mu+2}$, so ist:
\weq{II}{\big((f,f)^{2\mu}, f\big)^{2\mu} \equiv 0 \bmod. (ab)^{2\mu+2}, \quad \mu < \dfrac{n}{4}.}
Diese Gleichung gilt auch, wenn $n = 4\mu$, also $\lambda = 0$ ist. Es
fragt sich jetzt nur noch, besteht auch Gleichung (I) für $n = 4\mu$.
In diesem Falle tritt alsdann kein Factor $(ac)^{2\mu+\varrho}$ neben $(ab)^{2\mu}$ auf,
und wir können, um dies zu entscheiden, nicht mehr von dem dort
benutzten Gliede $\overline{\overline{G}}_0$ ausgehen. Wir benutzen vielmehr nun das symbolische
Product:
\[
P = (ab)^{2\mu-1}(ac)^{2\mu+1}(bc)^\varrho\,b_x^{2\mu-1-\varrho}\,c_x^{2\mu-1-\varrho},
\]
welches aus dem Polarengliede $\overline{\overline{G}}_0 = (ab)^{2\mu-1}\,a_y^{2\mu+1}\,b_y^\varrho\,b_x^{2\mu+1-\varrho}$ für $y_1 = c_2$,
$y_2 = -c_1$ hervorgeht. Die Reihenentwicklung liefert:
\[
\begin{aligned}
\overline{\overline{G}}_0 &= (f,f)^{2\mu-1}_{y^{2\mu+\varrho+1}} + \text{Const.}\,(f,f)^{2\mu}_{y^{2\mu+\varrho}}(yx) \\
&\qquad + \text{Const.}\times (f,f)^{2\mu+1}_{y^{2\mu+\varrho-1}}(yx)^2 + \cdots
\end{aligned}
\]
Hieraus entsteht nach ausgeführter Substitution:
\[
(ab)^{2\mu-1}(ac)^{2\mu+1}(bc)^\varrho\,b_x^{2\mu-1-\varrho}\,c_x^{2\mu-1-\varrho}
\equiv \text{C.}\,\big((f,f)^{2\mu}, f\big)^{2\mu+\varrho} \bmod. (ab)^{2\mu+2}.
\]
Diese Gleichung besteht, so lange $\varrho \leqq 2\mu - 1$ und demnach hat
man unter dieser Bedingung immer:
\weq{III}{\big((f,f)^{2\mu}, f\big)^{2\mu+\varrho} \equiv 0 \bmod. (ab)^{2\mu+2}, \quad n = 4\mu.}
Ist aber $\varrho = 2\mu$, so existirt kein symbolisches Product $P$, in
diesem Falle ist aber $\big((f,f)^{2\mu}, f\big)^{4\mu}$ die Invariante $j$, die sich somit
in kein anderes symbolisches Product umformen lässt. Nehmen wir
nun diese Invariante aber mit unter die Moduln auf, so können wir
die drei Gleichungen (I), (II), (III) in eine zusammenfassen:
\weq{IV}{U = \big((f,f)^{2\mu}, f\big)^{2\mu+\varrho} \equiv 0 \bmod. (ab)^{2\mu+2}, \; j,}
wobei nun $n \geqq 4\mu$, und $\varrho \leqq 2\mu$ sein kann.

Nun besitzt jede Ueberschiebung $U$ ein Glied
\[
G = (ab)^{2\mu}(ac)^{2\mu}\cdot Q,
\]
wo $Q$ die übrigen symbolischen Factoren enthält. Anderntheils ist
nach Nr.\,44 jedes Glied von $U$ durch $U$ selbst und Ueberschiebungen
%CONT>>


% ===== pdf idx 122, printed 108.0 =====
% !PRINTED 108
\weberpage{108}
%<<CONT
von Formen darstellbar, die aus $(f,f)^{2\mu} = (ab)^{2\mu}\,a_x^{n-2\mu}\,b_x^{n-2\mu}$, oder
$f = c_x^n$ durch Faltung entstehen. Die Faltung von $c_x^n$ liefert aber 0,
die von $(f,f)^{2\mu}$ einen weitern Klammerfactor $(ab)$. Demnach ist:
\[
G - \dot U \equiv 0 \bmod. (ab)^{2\mu+2}, % ?? uncertain: dot accent over U as printed, cause unclear
\]
oder, weil $U$ selbst $\equiv 0 \bmod. (ab)^{2\mu+2}$, $j$,
\[
G \equiv 0 \bmod. (ab)^{2\mu+2}, \; j.
\]
Alle symbolischen Producte aber mit dem Factor $(ab)^{2\mu}(ac)^{2\mu}$ kann
man sich aus $G$ durch Faltung mit anderen Formen entstanden denken.
Daher können wir schliesslich behaupten:
\[
(ab)^{2\mu}(ac)^{2\mu}\cdot R \equiv 0 \bmod. (ab)^{2\mu+2}, \; j,
\]
wo nun $R$ ein Product von irgend welchen symbolischen Factoren ist.
Der Eingangs aufgestellte Satz ist sonach bewiesen. Symbolische
Producte $G$ dieser Art liefern insbesondere die Ueberschiebungen
$\big((f,f)^{2\mu}, (f,f)^{2\mu}\big)^{2\mu}$; und demnach ist:
\[
U_1 = \big((f,f)^{2\mu}, (f,f)^{2\mu}\big)^{2\mu} \equiv 0 \bmod. (ab)^{2\mu+2}, \; j, % ?? uncertain: relation symbol faint in scan, read as \equiv
\]
da wie vorhin so auch hier $G - U_1 \equiv 0 \bmod. (ab)^{2\mu+2}$ ist. % ?? uncertain: relation symbol faint in scan, read as \equiv

99. \textit{Ueber die Cayley'sche respective Bézout'sche Methode der
Discriminantenbildung.} Die Discriminante einer Form $n^{\text{ten}}$ Grades ist
bekanntlich (vergl. auch Bd. I Nr.\,173) bis auf einen Zahlenfactor gleich
der Resultante der beiden ersten Differentialquotienten der Form.
Anstatt aber die Resultante aus diesen beiden Formen $(n-1)^{\text{ten}}$ Grades
direct zu bilden, wie dies Sylvester gethan hat, reducirten Bézout und
Cayley die zwei Gleichungen $(n-1)^{\text{ten}}$ Grades auf $(n-1)$ Gleichungen
$(n-2)^{\text{ten}}$ Grades und erniedrigten so wesentlich den Grad der Eliminationsdeterminante.
(Vergl. Bd. I Nr.\,150--152.) Wir wollen nun auch auf
symbolischem Wege diese Reduction vornehmen, um später bei der
Discriminantenbildung von Formen vierten, fünften und sechsten
Grades davon Gebrauch zu machen.

Angenommen die gegebene Form $f = a_x^n$ habe die Doppelwurzel
$\alpha$, dann verschwindet sowohl $f$, als auch $f_1$ und $f_2$ für $x = \alpha$, d.\,h.
es ist:
\weq{1}{(a\alpha)^n = 0}
\weq{2}{(a\alpha)^{n-1}\,a_x = 0.}
Aus der zweiten dieser Gleichungen folgt noch die dritte:
\weq{3}{(a\alpha)^{n-2}\,a_x^2 = \lambda \alpha^2.}
Denn bildet man die Functionaldeterminante dieser quadratischen
Form $(a\alpha)^{n-1}\,a_x^2 = c_x^2$ mit der linearen Form $\alpha_x$, so kommt:
\[
(c_x^2, \alpha_x) = \big((a\alpha)^{n-2}\,a_x^2, \alpha_x\big) = (a\alpha)^{n-1}\,\alpha_x = 0,
\]
%CONT>>


% ===== pdf idx 123, printed 109.0 =====
% !PRINTED 109
\weberpage{109}
%<<CONT
d. h. diese Functionaldeterminante verschwindet identisch. Aus dieser
Identität $(c\alpha)c_x = 0$ ergiebt sich unmittelbar durch Co\"efficientenvergleichung
:
\[
\bar{c}_0\alpha_2 - \bar{c}_1\alpha_1 = 0 \quad \text{und} \quad \bar{c}_1\alpha_2 - \bar{c}_2\alpha_1 = 0,
\]
oder:
\[
\frac{c_1}{c_0} = \frac{\alpha_2}{\alpha_1}, \quad \frac{c_2}{c_0} = \frac{\alpha_2^2}{\alpha_1^2},
\]
und damit die in (3) ausgesprochene Thatsache, dass die quadratische
Form $c_x^2$ proportional ist dem Quadrat der linearen Form $\alpha_x$.

Irgend zwei der drei Gleichungen (1), (2), (3) liefern also die
Bedingungen, dass $\alpha$ eine Doppelwurzel von $f$ ist; insbesondere gehen
aus der Identität (2) zwei Gleichungen $(n-1)^{\text{ten}}$ Grades in $\alpha$ hervor,
die wir nunmehr nach dem Vorgange von \gesperrt{Cayley} und \gesperrt{B\'ezout} auf
$(n-1)$ Gleichungen $(n-2)^{\text{ten}}$ Grades reduciren wollen. Zu dem Zwecke
ersetzen wir in (2) $x$ durch $b$, $\alpha$ durch $y$, multipliciren alsdann mit
$b_x^{n-1}$ und entwickeln die so entstehende Form in die Reihe (VIII).
Da in ihr alle ungeraden Ueberschiebungen verschwinden, so kommt:
\[
(ab)\,a_y^{n-1}b_x^{n-1} = \frac{n-1}{2}(f,f)^2_{y^{n-2}}(yx) + \frac{\left(\dfrac{n-1}{3}\right)^2}{\dbinom{2n-4}{3}}(f,f)^4_{y^{n-4}}(yx)^3 + \cdots
\]

Ersetzen wir hierin wiederum $y$ durch $\alpha$, so verschwindet wegen
(2) die linke Seite dieser Gleichung. Wir dividiren alsdann dieselbe
durch $\alpha_x$ und ersetzen die erste der rechts auftretenden Ueberschiebungen
durch ihr erstes Glied, indem wir uns die Ueberschiebung nach diesem
Glied und niedrigeren Ueberschiebungen entwickelt denken. Die letzteren
vereinigen wir mit den nämlichen ohnehin bereits auf der rechten
Seite vorhandenen entsprechenden Ueberschiebungen, während wir das
erste Glied von $((f,f)^2, \alpha_x^{n-2})^{n-2}$ auf die linke Seite schaffen; so erhalten
 wir:
\[
(ab)^2(a\alpha)^{n-2}b_x^{n-2} = c_0\left((f,f)^4,\alpha_x^{n-4}\right)^{n-4}\alpha_x^2 + c_1\left((f,f)^\sigma,\alpha_x^{n-6}\right)^{n-6}\cdot\alpha_x^4 + \cdots,
\]
wo nun $c_0, c_1, \ldots$ von den obigen Constanten verschieden sind. Diese
letzte Identität repräsentirt aber $(n-1)$ Gleichungen $(n-2)^{\text{ten}}$ Grades
in $\alpha$, und damit ist die gewünschte Reduction bewerkstelligt. Die
Determinante dieser $(n-1)$ homogenen Gleichungen für die $(n-1)$
Unbekannten $\alpha^{n-2}, \alpha^{n-3}, \ldots \alpha^2, \alpha, 1$ verschwindet identisch, und
da sie vom Grade $2(n-1)$ in den Co\"efficienten ist, so kann sie von
der Discriminante nur um einen numerischen Factor verschieden sein.


% ===== pdf idx 124, printed 110.0 =====
% !PRINTED 110
\weberpage{110}
\S.\,9. \textbf{Beweis, dass jede In- und Covariante durch ein symbolisches
Product darstellbar ist.}

100. \textit{Formulirung der Aufgabe für eine binäre Form.} Dass jedes
symbolische Product, sobald es nur die Symbole in geeigneter Zahl
enthält, eine simultane Invariante respective Covariante mehrerer gegebener
 Formen darstellt, haben wir bereits in \S.\,1 dieses Bandes erkannt
. Dies genügte für die bisherigen Untersuchungen. Die nun
folgenden Untersuchungen erheischen aber die Existenz des umgekehrten
 Satzes: Jede Invariante und Covariante einer Form oder eines
simultanen Systemes von Formen ist durch ein symbolisches Product
darstellbar; diesen Satz wollen wir nunmehr auf zwei verschiedenen
Wegen beweisen und nicht nur für binäre Formen, sondern auch für
Formen mit beliebig vielen Veränderlichen. Dabei können wir uns
im Beweise auf Invarianten einer einzigen Form beschränken. Denn
es wird sich einestheils zeigen, dass die Voraussetzung einer simultanen
Invariante mehrerer Formen den Beweis nur quantitativ und nicht
qualitativ ändern würde, und was anderntheils den entsprechenden Beweis
 für Covarianten betrifft, so geht derselbe wegen der Cogredienz
der Symbolreihen $a_2, -a_1; b_2, -b_1;$ etc., mit der Variabelnreihe $x_1,
x_2$ (vergl. Nr. 4) unmittelbar aus dem für Invarianten hervor. Ist
nämlich $\Phi$ eine simultane Covariante $m^{\text{ten}}$ Grades in $x$ der Formen
$f, \varphi, \psi, \chi \ldots$, und $g = g_1x_1 + g_2x_2$ irgend eine lineare Form, so ist
die Ueberschiebung $(\Phi, g^m)^m$ eine Invariante. Lässt sich nun diese Invariante
 durch ein symbolisches Product darstellen, so hat auch die
Covariante, die ja aus ihr durch die Substitution $x_1 = -g_2, x_2 = g_1$
wiederum hervorgeht, das analoge symbolische Product zum Repräsentanten
.

Es sei nun $i$ eine Invariante $\mu^{\text{ten}}$ Grades in den Co\"efficienten der
Form
\[
f = \bar{a}_0x_1^n + \binom{n}{1}\bar{a}_1x_1^{n-1}x_2 + \binom{n}{2}\bar{a}_2x_1^{n-2}x_2^2 + \cdots + \bar{a}_nx_2^n
\]
\[
= a_x^n = b_x^n = c_x^n = \ldots \text{etc.}
\]
Durch die Transformation
\[
x_1 = \xi_1y_1 + \eta_1y_2, \quad x_2 = \xi_2y_1 + \eta_2y_2
\]
vom Modul
\[
\Delta = (\xi\eta)
\]
geht diese Form $f$ über in die Form
\[
F = \bar{A}_0y_1^n + \binom{n}{1}\bar{A}_1y_1^{n-1}y_2 + \binom{n}{2}\bar{A}_2y_1^{n-2}y_2^2 + \cdots + \bar{A}_ny_2^n
% ?? uncertain: exponent on y_2 in the middle term of F, read by analogy with f above as y_2^2
\]
\[
= (a_\xi y_1 + a_\eta y_2)^n = (b_\xi y_1 + b_\eta y_2)^n = \ldots \text{etc.},
\]
%CONT>>


% ===== pdf idx 125, printed 111.0 =====
% !PRINTED 111
\weberpage{111}
%<<CONT
und wir wissen (vergl. Nr. 3), dass zwischen den transformirten wirklichen
 und symbolischen Co\"efficienten die analoge Beziehung:
\[
\bar{A}_k = a_\xi^{n-k}a_\eta^k
\]
besteht, wie zwischen den Co\"efficienten der ursprünglichen Form:
\[
\bar{a}_k = a_1^{n-k}a_2^k.
\]
Besitzt nun die Invariante der ursprünglichen Form $f$ den Ausdruck:
\[
i = \chi(\bar{a}),
\]
so muss gemäss der Definition einer Invariante der aus den Co\"efficienten
 $\bar{A}$ von $F$ gebildete Ausdruck
\[
J = \chi(\bar{A})
\]
bis auf eine Potenz des Transformationsmoduls dem ursprünglichen
gleich sein, d. h. die Bedingung, dass $i$ Invariante, stellt sich dar
durch die Gleichung
\weq{I}{J = i\cdot\Delta^\varrho = i\cdot(\xi\eta)^\varrho.}
Wir stellen nun die Behauptung auf: \glqq Die so definirte Function
\[
i = \chi(a)
\]
ist durch ein symbolisches Product oder eine Summe solcher Producte
darstellbar.\grqq{}

101. \textit{Umformung der Definitionsgleichung (I) einer Invariante.} Zunächst
 folgt aus der Identität (I), dass $i$ in den Co\"efficienten $\bar{a}$ homogen
und isobar vom Gewichte $\varrho$ sein muss. (Vergl. Nr. 8.) Führt man
daher in diese ganze homogene Function $i$ der Co\"efficienten $\bar{a}$ die
symbolischen Co\"efficienten
\[
\bar{a}_k = a_1^{n-k}a_2^k = b_1^{n-k}b_2^k = c_1^{n-k}c_2^k \ldots
\]
ein, so setzt sich $i$ aus einer Summe von Gliedern $G_k$ zusammen,
deren jedes die Form besitzt:
\[
G_k = \lambda_k \cdot a_1^{n-k}a_2^k \cdot b_1^{n-k}b_2^k \cdots r_1^{n-\pi}r_2^\pi \cdot t_1^{n-\pi}t_2^\pi \cdots s_1^{n-\nu}s_2^\nu \cdots \text{etc.}
\]
Hierbei ist $\lambda_k$ ein bestimmter Zahlenco\"efficient. Ersetzen wir hierin
nun $a_1, b_1, r_1, t_1 \ldots$ etc. durch $a_\xi, b_\xi, r_\xi t_\xi \ldots$ etc. und ebenso $a_2,
b_2, r_2, t_2 \ldots$ etc. durch $a_\eta, b_\eta, r_\eta, t_\eta \ldots$ etc., so erhalten wir das entsprechende
 Glied
\weq{1}{\bar{G}_k = \lambda_k a_\xi^{n-k}a_\eta^k \cdot b_\xi^{n-k}b_\eta^k \cdots r_\xi^{n-\pi}r_\eta^\pi \cdot t_\xi^{n-\pi}t_\eta^\pi \cdots s_\xi^{n-\nu}s_\eta \cdots \text{etc.}}
% ?? uncertain: last factor of (1), s_\xi^{n-\nu}s_\eta -- exponent on the final s-term is unclear in the scan, read by analogy with the r- and t-terms
der Invariante $J$. Es enthält $\varrho$ Factoren, die linear in $\xi$ und ebenso
viele, die linear in $\eta$ sind; daher können wir dasselbe kürzer bezeichnen
 mit
\[
\bar{G}_k = m^\varrho_{k,\xi}\,n^\varrho_{k,\eta}.
\]


% ===== pdf idx 126, printed 112.0 =====
% !PRINTED 112
\weberpage{112}
Die Invariante $J$ selbst ist dann dargestellt durch
\weq{2}{J = \sum_1^s \lambda_k m^\varrho_{k,\xi}\cdot n^\varrho_{k,\eta}.}
Die Summe rechts ist eine homogene Function mit zwei Reihen
variabler Parameter $\xi$ und $\eta$, und wir können sie stets symbolisch
durch
\weq{2a}{p^\varrho_\xi q^\varrho_\eta = (\mu_1m_\xi^{(1)} + \mu_2m_\xi^{(2)} + \cdots + \mu_sm_\xi^{(s)})^\varrho(\nu_1n_\eta^{(1)} + \nu_2n_\eta^{(2)} \cdots \nu_sn_\eta^{(s)})^\varrho}
% [sic] second bracket lacks a "+" before the last term, unlike the first bracket
repräsentiren, wenn wir nur die rechts auftretenden numerischen Grössen
$\mu_i$ und $\nu_i$ in Verbindung mit den sich ergebenden Polynomialco\"efficienten
 so bestimmen, dass ihre jeweiligen Producte mit den in (2) vorhandenen
 numerischen Grössen $\lambda_k$ übereinstimmen.

Führt man diesen Ausdruck $p^\varrho q^\varrho$ für $J$ in die Identität (I) Nr. 100
ein, so erhält man:
\weq{II}{p^\varrho_\xi q^\varrho_\eta = i\cdot(\xi\eta)^\varrho.}
An der so umgeformten Identität können wir nun in zweierlei
Weise den gewünschten Beweis erbringen. Einmal indem wir auf die
linke Seite (II) die Reihenentwicklung anwenden, das andere Mal
durch Anwendung des $\Omega$-Processes. Der erste Beweis lässt sich nicht
weiter verallgemeinern, ist aber insofern lehrreicher, als er zeigt, wie
sich direct durch die Reihenentwicklung die symbolische Darstellung
einer Invariante ergiebt. Der zweite Beweis ist dagegen der Erweiterung
 auf ein Gebiet mit beliebig vielen Variabeln fähig.

102. \textit{Erster Beweis.} Wir hatten in Nr. 81 für eine Form $F$ mit
zwei Reihen cogredienter Variabeln die Reihenentwicklung
\[
F = r_x^m s_y^n = \sum C_{0k}(\overline{rs})^k_{y^{n-k}}(xy)^k
\]
erhalten. Wenden wir dieselbe auf die linke Seite der Identität (II)
an, so kommt:
\weq{III}{p^\varrho_\xi q^\varrho_\eta = \sum_{k=0}^{k=\varrho} \frac{\binom{\varrho}{k}\binom{\varrho}{k}}{\binom{2\varrho-k+1}{k}}(\overline{pq})^k_{\eta^{\varrho-k}}(\xi\eta)^k.}
Diese Reihe ist eindeutig; wenn wir daher den Werth von $p^\varrho_\xi q^\varrho_\eta$
aus (III) in (II) eintragen, so müssen beide Seiten Glied um Glied
übereinstimmen. Da aber rechts alle $\varrho-1$ ersten Glieder verschwinden
, so folgt aus
\[
\sum \frac{\binom{\varrho}{k}\binom{\varrho}{k}}{\binom{2\varrho-k+1}{k}}\cdot(\overline{pq})^k_{\eta^{\varrho-k}}(\xi\eta)^k = i\cdot(\xi\eta)^\varrho
\]
direct:
%CONT>>


% ===== pdf idx 127, printed 113.0 =====
% !PRINTED 113
\weberpage{113}
\[
\frac{\binom{\varrho}{\varrho}\binom{\varrho}{\varrho}}{\binom{\varrho+1}{\varrho}}\,\overline{(p\,q)}^\varrho\,(\xi\,\eta)^\varrho = i\,(\xi\,\eta)^\varrho,
\]
oder
\[
\frac{1}{\varrho+1}(pq)^\varrho = i.
\]

Die Invariante $i$ stellt sich also zunächst als das symbolische Product $(pq)^\varrho$ dar. Nun ist aber gemäss der Gleichung (2a) Nr. 101:
\[
\begin{aligned}
p_1 &= \mu_1 m_1^{(1)} + \mu_2 m_1^{(2)} + \cdots + \mu_s m_1^{(s)}\\
p_2 &= \mu_1 m_2^{(1)} + \mu_2 m_2^{(2)} + \cdots + \mu_s m_2^{(s)}\\
q_1 &= \nu_1 n_1^{(1)} + \nu_2 n_1^{(2)} + \cdots + \nu_s n_1^{(s)}\\
q_2 &= \nu_1 n_2^{(1)} + \nu_2 n_2^{(2)} + \cdots + \nu_s n_2^{(s)}.
\end{aligned}
\]
Ersetzt man daher in $(pq)$ die Grössen $p_1\,p_2\,q_1\,q_2$ durch diese ihre Werthe in $m^{(k)}$ und $n^{(k)}$, so geht $(pq)$ über in $\sum \mu_k \nu_\lambda (mn)$ und demnach wird:
\[
(pq)^\varrho = \left\{\sum \mu_k \nu_\lambda (mn)\right\}^\varrho = \sum c_i \cdot (m_{k_1} n_{\lambda_1})^{\varrho_1} (m_{k_2} n_{\lambda_2})^{\varrho_2} \cdots -, % ?? uncertain: a small dot is printed above the second summation sign here; not reproducible as a diacritic on \sum, noted but not altered
\]
wobei $\varrho_1 + \varrho_2 + \varrho_3 + \cdots = \varrho$.

Jedes Product in der Summe rechts repräsentirt eine $\varrho^{\text{te}}$ Ueberschiebung je eines Productes von $\varrho$ linearen Factoren $a_\xi, b_\xi \ldots$, mit einem Product von $\varrho$ linearen Factoren $a_\eta, b_\eta\ldots$ Eine solche Ueberschiebung ist aber nach Nr. 38 (I) durch Aggregate von symbolischen Producten darstellbar, und damit auch die Form $(pq)^\varrho = (\varrho + 1) \cdot i$.

\textit{Anmerkung.} Wir haben diesen Beweis geführt für den Fall, dass $i$ eine Invariante von $f$ allein ist. Man erkennt aber, dass die Annahme, $i$ sei simultane Invariante von $f$, $\varphi$, $\psi$ etc., die ganze Beweisführung keineswegs ändern kann, da dieselbe lediglich eine Vermehrung in der Anzahl der Symbole bewirkt, von der der gegebene Beweis aber unabhängig ist. Daher ist auch jede simultane Invariante durch ein symbolisches Product darstellbar.

103. \textit{Zweiter Beweis.} Die Form $p_\xi^\varrho q_\eta^\varrho$ ist ebenso wie die Determinante $(\xi\eta)$ eine Form mit zwei Reihen binärer variabler Parameter. Wenden wir nun auf die Identität (II) Nr. 101 den Process
\[
\Omega(f) = \frac{\partial^2 f}{\partial \xi_1 \partial \eta_2} - \frac{\partial^2 f}{\partial \eta_1 \partial \xi_2}
\]
an, so erhalten wir gemäss der in Nr. 79 entwickelten Gleichung (II), insofern wir die rechte Seite differentiiren:
\[
\varrho^2 i \Omega(\xi\eta)^\varrho = i \varrho (\varrho + 1) (\xi\eta)^{\varrho - 1},
\]
%CONT>>


% ===== pdf idx 128, printed 114.0 =====
% !PRINTED 114
\weberpage{114}
%<<CONT
und insofern wir die linke Seite dem gleichen Processe unterwerfen:
\[
\begin{aligned}
\varrho^2 \cdot \Omega(p_\xi^\varrho q_\eta^\varrho) &= \varrho^2 p_1 q_2 \cdot p_\xi^{\varrho-1} q_\eta^{\varrho-1} - \varrho^2 p_2 q_1 \cdot p_\xi^{\varrho-1} q_\eta^{\varrho-1}\\
&= \varrho^2 (pq) p_\xi^{\varrho-1} q_\eta^{\varrho-1}.
\end{aligned}
\]
Die Identität (II) geht also über in:
\[
(\varrho + 1) \varrho \cdot (\xi\eta)^{\varrho-1} i = \varrho^2 (pq) p_\xi^{\varrho-1} q_\eta^{\varrho-1}.
\]
Unterwirft man diese neue Identität wiederholt dem gleichen Processe, so kommt:
\[
(\varrho + 1) \varrho \cdot \varrho (\varrho - 1) \cdot (\xi\eta)^{\varrho-2} \cdot i = \varrho^2 \cdot (\varrho - 1)^2 (pq)^2 p_\xi^{\varrho-2} q_\eta^{\varrho-2}
\]
\[
(\varrho + 1) \varrho (\varrho - 1) \cdot \varrho (\varrho - 1) (\varrho - 2) \cdot (\xi\eta)^{\varrho-3} \cdot i = \varrho^2 (\varrho - 1)^2 (\varrho - 2)^2 (pq)^3 p_\xi^{\varrho-3} q_\eta^{\varrho-3}
\]
\[
\cdot\ \cdot\ \cdot\ \cdot\ \cdot\ \cdot\ \cdot\ \cdot\ \cdot\ \cdot\ \cdot\ \cdot\ \cdot\ \cdot\ \cdot\ \cdot
\]
und schliesslich nach $\varrho$ maliger Ausführung des $\Omega$-Processes
\[
(\varrho + 1)!\, \varrho!\, i = (\varrho!)^2 \cdot (pq)^2,
\]
oder
\[
i = \frac{1}{\varrho + 1} (p\varrho)^\varrho, % [sic] printed as (p\varrho)^\varrho, not (pq)^\varrho as the context requires -- reproduced as printed
\]
ein Resultat, das mit dem im vorausgehenden Beweise gewonnenen übereinstimmt.

104. \textit{Verallgemeinerung des zweiten Beweises; vorbereitende Untersuchungen.} Wie schon erwähnt, lässt sich der zweite Beweis auf ein Gebiet von beliebig vielen Variabeln ausdehnen. Es wird genügen, denselben auf quinäres Gebiet zu erweitern, da Gedankengang und Rechnung, die sich in engster Weise an den ebengeführten Beweis für binäre Formen anschliessen, ganz ebenso im allgemeinen Falle sich gestalten. Sei also
\[
f = a_x^n = (a_1 x_1 + a_2 x_2 + a_3 x_3 + a_4 x_4 + a_5 x_5)^n = b_x^n = = \text{etc.} % [sic] double equals sign printed before \textquotedbl{}etc.\textquotedbl{}
\]
der symbolische Ausdruck einer ganzen homogenen Form mit fünf Variabeln, und $i = \chi(\bar a)$ sei irgend eine Invariante derselben, dargestellt in den unsymbolischen Coefficienten $\bar a$ von $f$.

Wenden wir auf $f$ die Substitutionen an
\weq{I}{
\left.\begin{aligned}
x_1 &= \xi_1 y_1 + \eta_1 y_2 + \zeta_1 y_3 + \lambda_1 y_4 + \tau_1 y_5\\
x_2 &= \xi_2 y_1 + \eta_2 y_2 + \zeta_2 y_3 + \lambda_2 y_4 + \tau_2 y_5\\
&\quad\cdot\ \cdot\ \cdot\ \cdot\ \cdot\ \cdot\ \cdot\ \cdot\ \cdot\ \cdot\ \cdot\ \cdot\ \cdot\ \cdot\\
x_5 &= \xi_5 y_1 + \eta_5 y_2 + \zeta_5 y_3 + \lambda_5 y_4 + \tau_5 y_5
\end{aligned}\right\},
}
so geht die ursprüngliche Function $f$ über in die Form
\[
\begin{aligned}
F &= (y_1 a_\xi + y_2 a_\eta + y_3 a_\zeta + y_4 a_\lambda + y_5 a_\tau) = A_y^n\\
&= B_y^n = C_y^n = = \text{etc.} % [sic] double equals sign printed before \textquotedbl{}etc.\textquotedbl{}
\end{aligned}
\]
Waren demnach in $f$ die nicht symbolischen Coefficienten $\bar a$ dargestellt durch
%CONT>>


% ===== pdf idx 129, printed 115.0 =====
% !PRINTED 115
\weberpage{115}
%<<CONT
\[
a_{k\varrho\mu\nu\sigma} = a_1^{n-k} a_2^\varrho a_3^\mu a_4^\nu a_5^\sigma, \quad k = \varrho + \mu + \nu + \sigma,
\]
so ist in der transformirten Form $F$ der entsprechende Coefficient
\[
A_{k\varrho\mu\nu\sigma} = a_\xi^{n-k} a_\eta^\varrho a_\zeta^\mu a_\lambda^\nu a_\tau^\sigma.
\]
Da ferner der Transformationsmodul durch die fünfreihige Determinante
\[
(\xi\eta\zeta\lambda\tau) =
\begin{vmatrix}
\xi_1 & \xi_2 & \xi_3 & \xi_4 & \xi_5\\
\eta_1 & \eta_2 & \eta_3 & \cdot & \cdot\\
\zeta_1 & \cdot & \cdot & \cdot & \cdot\\
\cdot & \cdot & \cdot & \cdot & \cdot\\
\tau_1 & \tau_2 & \tau_3 & \cdot & \tau_5
\end{vmatrix}
= \Delta % ?? uncertain: several matrix entries are printed only as a row of dots indicating omitted terms; the dot placement is reproduced as best legible, not a literal transcription of each dot's column
\]
dargestellt ist, so wird im quinären Gebiete der analytische Ausdruck für die Definition einer Invariante:
\weq{I}{J = i (\xi\eta\zeta\lambda\tau)^\varrho.}
Dabei ist also $i$ genau dieselbe homogene Function in den Symbolreihen $a_1 a_2 \ldots a_5$; $b_1 b_2 \ldots b_5$; $\ldots$ etc., welche $J$ in den Symbolreihen $a_\xi a_\eta a_\zeta a_\lambda a_\tau$; $b_\xi b_\eta b_\zeta b_\lambda b_\tau$; $\ldots$ etc. ist. Jedes einzelne Glied $\overline{G}_k$ von $J$ ist homogen von der $\varrho^{\text{ten}}$ Dimension jeder der fünf Variabeln $\xi\eta\zeta\lambda\tau$, und wir können es daher wie früher symbolisch darstellen durch
\[
\overline{G}_k = \lambda_k\, r_\xi^\varrho\, s_\eta^\varrho\, t_\zeta^\varrho\, u_\lambda^\varrho\, v_\tau^\varrho,
\]
wo $\lambda_k$ ein Zahlencoefficient.

Daher lässt sich auch die Summe aller dieser Glieder, d. h. die Invariante zunächst durch den symbolischen Ausdruck
\[
J = p_\xi^\varrho\, q_\eta^\varrho\, m_\zeta^\varrho\, n_\lambda^\varrho\, k_\tau^\varrho
\]
wiedergeben. Die Identität (I) kann also auch in der Form geschrieben werden:
\weq{II}{p_\xi^\varrho\, q_\eta^\varrho\, m_\zeta^\varrho\, n_\lambda^\varrho\, k_\tau^\varrho = i (\xi\eta\zeta\lambda\tau)^\varrho.}
105. \textit{Der $\Omega$-Process für eine Form mit beliebig vielen Variabeln.} Auf die zuletzt erhaltene Identität wenden wir nunmehr den $\Omega$-Process an. Derselbe lässt sich, von einem Zahlenfactor abgesehen, für fünf Reihen quinärer Variabler symbolisch durch die Determinante
\[
\Omega(f) =
\begin{vmatrix}
\dfrac{\partial f}{\partial \xi_1} & \dfrac{\partial f}{\partial \xi_2} & \dfrac{\partial f}{\partial \xi_3} & \dfrac{\partial f}{\partial \xi_4} & \dfrac{\partial f}{\partial \xi_5}\\[6pt]
\dfrac{\partial f}{\partial \eta_1} & \dfrac{\partial f}{\partial \eta_2} & \cdot & \cdot & \cdot\\[6pt]
\cdot & \cdot & \cdot & \cdot & \cdot\\[6pt]
\cdot & \cdot & \cdot & \cdot & \cdot\\[6pt]
\dfrac{\partial f}{\partial \tau_1} & \dfrac{\partial f}{\partial \tau_2} & \dfrac{\partial f}{\partial \tau_3} & \cdot & \dfrac{\partial f}{\partial \tau_5}
\end{vmatrix}
\]
%CONT>>


% ===== pdf idx 130, printed 116.0 =====
% !PRINTED 116
\weberpage{116}
%<<CONT
darstellen, analog wie wir ihn für Formen mit zwei Reihen binärer Variabler durch
\[
\Omega(f) =
\begin{vmatrix}
\dfrac{\partial f}{\partial \xi_1} & \dfrac{\partial f}{\partial \xi_2}\\[6pt]
\dfrac{\partial f}{\partial \eta_1} & \dfrac{\partial f}{\partial \eta_2}
\end{vmatrix}
\]
charakterisiren konnten. Dabei ist wie hier ein Determinantenglied nicht als Product von fünf ersten Differentialquotienten aufzufassen, sondern als ein einziger fünfter partieller Differentialquotient. Es fragt sich nun, in welcher Weise verändert sich die Identität (II), wenn wir die beiden Seiten derselben dem so definirten Processe $\Omega$ unterwerfen. Wir erledigen die Frage schrittweise. Zu dem Zwecke unterwerfen wir die Identität (II) zunächst den Operationen $\dfrac{\partial^2}{\partial \xi_1 \partial \eta_2}$ und $\dfrac{\partial^2}{\partial \xi_2 \partial \eta_1}$, indem wir uns alle übrigen Variabeln constant denken.

Aus der linken Seite erhalten wir durch die erste Operation:
\weq{1}{\frac{\partial^2}{\partial \xi_1 \partial \eta_2} (J_{\xi\eta}) = \varrho^2 \cdot p_\xi^{\varrho-1} q_\eta^{\varrho-1}\, m_\zeta^\varrho n_\lambda^\varrho k_\tau^\varrho \cdot p_1 q_2;}
die rechte Seite aber geht, abgesehen vom Factor $i$, über in
\[
\frac{\partial^2}{\partial \xi_1 \partial \eta_2} (\xi\eta\zeta\lambda\tau)^\varrho = \frac{\partial\left[\varrho \cdot \Delta^{\varrho-1} \cdot \dfrac{\partial \Delta}{\partial \xi_1}\right]}{\partial \eta_2},
\]
also
\weq{2}{\frac{\partial^2 (\Delta^\varrho)}{\partial \xi_1 \partial \eta_2} = \varrho \cdot (\varrho - 1) \cdot \Delta^{\varrho-2} \cdot \frac{\partial \Delta}{\partial \xi_1} \cdot \frac{\partial \Delta}{\partial \eta_2} + \varphi \cdot \Delta^{\varrho-1} \cdot \frac{\partial^2 \Delta}{\partial \xi_1 \partial \eta_2}. % [sic] coefficient of the second term reads \varphi in print, inconsistent with the parallel \varrho in eq. (3) below; reproduced as printed
}
Ganz ebenso liefert die Differentiation
\weq{3}{\frac{\partial^2 (\Delta^\varrho)}{\partial \xi_2 \partial \eta_1} = \varrho (\varrho - 1) \Delta^{\varrho-2} \cdot \frac{\partial \Delta}{\partial \xi_2} \cdot \frac{\partial \Delta}{\partial \eta_1} + \varrho \cdot \Delta^{\varrho-1} \cdot \frac{\partial^2 \Delta}{\partial \xi_2 \partial \eta_1},}
während die gleiche Operation auf der linken Seite ausgeführt, das Resultat:
\weq{4}{\frac{\partial^2 (J_{\xi\eta})}{\partial \xi_2 \partial \eta_1} = \varrho^2 \cdot p_\xi^{\varrho-1} \cdot q_\eta^{\varrho-1}\, m_\zeta^\varrho n_\lambda^\varrho k_\tau^\varrho \cdot p_2 q_1}
ergiebt. Subtrahirt man nun die vierte von der ersten Gleichung, so erhält man:
\weq{5}{\Omega(J_{\xi\eta}) = \varrho^2 p_\xi^{\varrho-1} q_\eta^{\varrho-1} m_\zeta^\varrho n_\lambda^\varrho k_\tau^\varrho (p_1 q_2 - q_1 p_2).}
Subtrahirt man ebenso die dritte Gleichung von der zweiten, so kommt
\weq{6}{\Omega(\Delta^\varrho) = \varrho (\varrho - 1) \Delta^{\varrho-2} \cdot \left\{\frac{\partial \Delta}{\partial \xi_1} \cdot \frac{\partial \Delta}{\partial \eta_2} - \frac{\partial \Delta}{\partial \xi_2} \cdot \frac{\partial \Delta}{\partial \eta_1}\right\} + \varrho \cdot \Delta^{\varrho-1} \left\{\frac{\partial^2 \Delta}{\partial \xi_1 \partial \eta_2} - \frac{\partial^2 \Delta}{\partial \xi_2 \partial \eta_1}\right\}.}


% ===== pdf idx 131, printed 117.0 =====
% !PRINTED 117
\weberpage{117}
Nun sind aber $\dfrac{\partial \Delta}{\partial \xi_i}$, $\dfrac{\partial \Delta}{\partial \eta_i}$, $\dfrac{\partial^2 \Delta}{\partial \xi_i\, \partial \eta_k}$ nichts anderes als Minoren erster, bezw. zweiter Ordnung von $\Delta$, welche zum Elemente $\xi_i$, $\eta_i$, bezw. zum Producte $\xi_i \eta_k$ geh\"oren, und man weiss aus der Lehre von den Determinanten einmal, dass
\[
\frac{\partial^2 \Delta}{\partial \xi_2\, \partial \eta_1} = -\, \frac{\partial^2 \Delta}{\partial \xi_1\, \partial \eta_2}
\]
% ?? uncertain: a faint mark beneath/after $\eta_2$ in this display on the scan; could not resolve it to a real diacritic, transcribed as plain $\eta_2$
und ebenso, dass $\dfrac{\partial \Delta}{\partial \xi_1} \cdot \dfrac{\partial \Delta}{\partial \eta_2} - \dfrac{\partial \Delta}{\partial \xi_2} \cdot \dfrac{\partial \Delta}{\partial \eta_1}$ nichts anderes ist als der Minor $n - 2 = 3^{\text{ter}}$ Ordnung der zu $\Delta$ adjungirten Determinante, dass also der Werth dieses Ausdruckes dargestellt ist durch (vgl. Nr. 88, Bd. I)
\[
\Delta \cdot \frac{\partial^2 \Delta}{\partial \xi_1\, \partial \eta_2}\,.
\]
Tr\"agt man also die so erhaltenen Werthe in (6), so reducirt sich der Ausdruck auf:
\[
\Omega(\Delta^{\varrho}) = \varrho(\varrho - 1)\cdot \Delta^{\varrho-1}\cdot \frac{\partial^2 \Delta}{\partial \xi_1\, \partial \eta_2} + 2\varrho \cdot \Delta^{\varrho-1}\cdot \frac{\partial^2 \Delta}{\partial \xi_1\, \partial \eta_2}\,,
\]
oder:
\weq{7}{\Omega(\Delta^{\varrho}) = \varrho(\varrho + 1)\cdot \Delta^{\varrho-1}\cdot \frac{\partial^2 \Delta}{\partial \xi_1\, \partial \eta_2}\,.}
Durch Comparation der mit $i$ multiplicirten Gleichung (7) mit der Gleichung (5) erh\"alt man sonach
\weq{I}{\varrho^2\, p_{\xi}^{\varrho-1} \cdot q_{\eta}^{\varrho-1} m_{\zeta}^{\varrho} n_{\lambda}^{\varrho} k_{\tau}^{\varrho} \cdot (pq) = i \cdot \varrho\, (\varrho + 1)\, \Delta^{\varrho-1}\cdot \frac{\partial^2 \Delta}{\partial \xi_1\, \partial \eta_2}\,.}

106. \textit{Wir untersuchen nun weiter die Wirkung des $\Omega$-Processes, wenn auch noch die Variable $\xi$ mit hereingezogen wird, so dass nur mehr $\eta$ und $\tau$ als constant betrachtet werden.} Zu dem Zwecke differentiiren wir die Gleichungen (I) von Nr. 105 partiell nach $\xi_3$ und erhalten auf der linken Seite:
\weq{8}{\varrho^3 \cdot p_{\xi}^{\varrho-1}\, q_{\eta}^{\varrho-1}\, m_{\zeta}^{\varrho-1}\, n_{\lambda}^{\varrho}\, k_{\tau}^{\varrho}\, (p_1 q_2 - q_1 p_2) \cdot m_3\,,}
auf der rechten Seite dagegen:
\weq{9}{i \cdot \varrho \cdot (\varrho + 1)(\varrho - 1) \cdot \Delta^{\varrho-2} \cdot \frac{\partial^2 \Delta}{\partial \xi_1\, \partial \eta_2} \cdot \frac{\partial \Delta}{\partial \zeta_3} + i \cdot \varrho \cdot (\varrho + 1) \cdot \Delta^{\varrho-1} \cdot \frac{\partial^3 \Delta}{\partial \xi_1\, \partial \eta_2\, \partial \zeta_3}\,.}
Hieraus entsteht durch cyklische Vertauschung der Zahlen 1, 2, 3 zun\"achst aus dem Producte (8)
\weq{10}{\varrho^3 \cdot p_{\xi}^{\varrho-1}\, q_{\eta}^{\varrho-1}\, m_{\zeta}^{\varrho-1}\, n_{\lambda}^{\varrho}\, k_{\tau}^{\varrho}\, (p_2 q_3 - q_2 p_3)\, m_1\,,}
und
\weq{11}{\varrho^3 \cdot p_{\xi}^{\varrho-1}\, q_{\eta}^{\varrho-1}\, m_{\zeta}^{\varrho-1}\, n_{\lambda}^{\varrho}\, k_{\tau}^{\varrho}\, (p_3 q_1 - q_3 p_1)\, m_2\,.}
Im Ausdrucke (9) bleibt bei dieser Vertauschung der zweite Term seinem Werthe nach unver\"andert, da $\dfrac{\partial^3 \Delta}{\partial \xi_1\, \partial \eta_2\, \partial \zeta_3}$, $\dfrac{\partial^3 \Delta}{\partial \xi_2\, \partial \eta_3\, \partial \zeta_1}$, $\dfrac{\partial^3 \Delta}{\partial \xi_3\, \partial \eta_1\, \partial \zeta_2}$ die n\"amlichen Unterdeterminanten inclusive Vorzeichen repr\"asentiren.
% ?? uncertain: a small speck is printed just before the third partial-derivative term above on the scan; treated as a printing blemish and not transcribed as a symbol


% ===== pdf idx 132, printed 118.0 =====
% !PRINTED 118
\weberpage{118}
In dem ersten Term ersetze ich vor der cyklischen Vertauschung das Product
\[
\Delta \cdot \frac{\partial^2 \Delta}{\partial \xi_1\, \partial \eta_2} \cdot \frac{\partial \Delta}{\partial \zeta_3}
\]
gem\"ass den Gesetzen \"uber die Minoren adjungirter Determinanten (vgl. Bd. I, § 7 Nr. 87 und 88) durch:
\weq{12}{\left(\frac{\partial \Delta}{\partial \xi_1} \cdot \frac{\partial \Delta}{\partial \eta_2} - \frac{\partial \Delta}{\partial \xi_2} \cdot \frac{\partial \Delta}{\partial \eta_1}\right) \cdot \frac{\partial \Delta}{\partial \zeta_3}\,,}
woraus durch cyklische Vertauschung die Ausdr\"ucke entstehen:
\weq{13}{\left(\frac{\partial \Delta}{\partial \xi_2} \cdot \frac{\partial \Delta}{\partial \eta_3} - \frac{\partial \Delta}{\partial \xi_3} \cdot \frac{\partial \Delta}{\partial \eta_2}\right) \cdot \frac{\partial \Delta}{\partial \zeta_1}}
\weq{14}{\left(\frac{\partial \Delta}{\partial \xi_3} \cdot \frac{\partial \Delta}{\partial \eta_1} - \frac{\partial \Delta}{\partial \xi_1} \cdot \frac{\partial \Delta}{\partial \eta_3}\right) \cdot \frac{\partial \Delta}{\partial \zeta_2}\,.}
Addirt man nun die durch cyklische Vertauschung f\"ur die linke Seite erhaltenen Resultate (8), (10), (11), so erh\"alt man:
\[
\Omega(J_{\xi\eta\zeta}) = \varrho^3 \cdot p_{\xi}^{\varrho-1}\, q_{\eta}^{\varrho-1}\, m_{\zeta}^{\varrho-1}\, n_{\lambda}^{\varrho}\, k_{\tau}^{\varrho} \cdot
\begin{vmatrix}
p_1 & p_2 & p_3\\
q_1 & q_2 & q_3\\
m_1 & m_2 & m_3
\end{vmatrix}\,.
\]
Auf der rechten Seite findet sich nach vollzogener Addition einmal das dreifache des zweiten Terms von (9), sodann aber als Factor von $i \cdot \varrho \cdot (\varrho+1)(\varrho-1)\, \Delta^{\varrho-3}$ die Determinante
\[
\begin{vmatrix}
\dfrac{\partial \Delta}{\partial \xi_1} & \dfrac{\partial \Delta}{\partial \xi_2} & \dfrac{\partial \Delta}{\partial \xi_3}\\[6pt]
\dfrac{\partial \Delta}{\partial \eta_1} & \dfrac{\partial \Delta}{\partial \eta_2} & \dfrac{\partial \Delta}{\partial \eta_3}\\[6pt]
\dfrac{\partial \Delta}{\partial \zeta_1} & \dfrac{\partial \Delta}{\partial \zeta_2} & \dfrac{\partial \Delta}{\partial \zeta_3}
\end{vmatrix}\,,
\]
welche als Minor zweiter Ordnung der zu $\Delta$ adjungirten Determinante den Werth besitzt
\[
\Delta^2 \cdot \frac{\partial^3 \Delta}{\partial \xi_1\, \partial \eta_2\, \partial \zeta_3}\,.
\]
Substituirt man also diesen Werth und setzt die beiden Seiten einander gleich, so erh\"alt man endlich:
\[
\begin{aligned}
\Omega(J_{\xi\eta\zeta}) &= \varrho^3 \cdot p_{\xi}^{\varrho-1}\, q_{\eta}^{\varrho-1}\, m_{\zeta}^{\varrho-1}\, n_{\lambda}^{\varrho}\, k_{\tau}^{\varrho} \cdot (pqm)\\
&= i \cdot \varrho \cdot (\varrho + 1)(\varrho - 1)\, \Delta^{\varrho-1} \cdot \frac{\partial^3 \Delta}{\partial \xi_1\, \partial \eta_2\, \partial \zeta_3}\\
&\quad + i \cdot \varrho \cdot (\varrho + 1) \cdot 3\, \Delta^{\varrho-1} \cdot \frac{\partial^3 \Delta}{\partial \xi_1\, \partial \eta_2\, \partial \zeta_3}
\end{aligned}
\]
oder:
% ?? uncertain: a stray printed dot appears before "oder" and again at the start of the next formula line on the scan; not transcribed as a symbol
\weq{II}{\varrho^3 \cdot p_{\xi}^{\varrho-1}\, q_{\eta}^{\varrho-1}\, m_{\zeta}^{\varrho-1}\, n_{\lambda}^{\varrho}\, k_{\tau}^{\varrho}\, (pqm) = i \cdot \varrho \cdot (\varrho + 1)(\varrho + 2) \cdot \Delta^{\varrho-1} \cdot \frac{\partial^3 \Delta}{\partial \xi_1\, \partial \eta_2\, \partial \zeta_3}}


% ===== pdf idx 133, printed 119.0 =====
% !PRINTED 119
\weberpage{119}
Aus dem Resultate (I) Nr. 105 und dem soeben erhaltenen (II) geht nun das Gesetz der Wirkung des $\Omega$-Processes auf eine Form mit beliebig vielen Variabeln hinreichend klar hervor. Man erkennt, dass dasselbe, ausgedehnt auf alle f\"unf Variablen, durch den Ausdruck gegeben ist:
\weq{III}{\varrho^5 \cdot p_{\xi}^{\varrho-1}\, q_{\eta}^{\varrho-1}\, m_{\zeta}^{\varrho-1}\, n_{\lambda}^{\varrho-1}\, k_{\tau}^{\varrho-1}\, (pqmnk) = i \cdot \varrho\, (\varrho + 1)(\varrho + 2)(\varrho + 3)(\varrho + 4)\, (\xi\eta\zeta\lambda\tau)^{\varrho-1}\,.}
Unterwirft man nun (III) abermals dem $\Omega$-Process, so kommt:
\[
\varrho^5 \cdot (\varrho - 1)^5 \cdot p_{\xi}^{\varrho-2}\, q_{\eta}^{\varrho-2}\, m_{\zeta}^{\varrho-2}\, n_{\lambda}^{\varrho-2}\, k_{\tau}^{\varrho-2}\, (pqmnk)^2 = i \cdot \varrho \cdot (\varrho - 1) \cdot (\varrho + 1) \cdot \varrho \cdot (\varrho + 2)(\varrho + 1) \dots (\xi\eta\zeta\lambda\tau)^{\varrho-2}\,.
\]
Wiederholt man daher den Process $\varrho$ mal, so ergiebt sich endlich die Identit\"at:
\[
(\varrho\,!)^5 \cdot (pqmnk)^{\varrho} = \varrho\, ! \, \frac{(\varrho+1)!}{1!}\, \frac{(\varrho+2)!}{2!}\, \frac{(\varrho+3)!}{3!}\, \frac{(\varrho+4)!}{4!}\, i
\]
oder ganz allgemein f\"ur $m$ beliebige Variable
\weq{IV}{(\varrho\,!)^m\, (pqmnk\dots)^{\varrho} = \varrho\, ! \, \frac{(\varrho+1)!}{1!}\, \frac{(\varrho+2)!}{2!} \dots \frac{(\varrho+m-1)!}{(m-1)!}\, i\,.}
Dadurch ist der Beweis erbracht, dass die Invariante $i$ einer Function von beliebig vielen Variabeln sich symbolisch durch ein Product von Klammerfactoren oder eine Summe solcher Producte darstellen l\"asst. Der Beweis gilt auch noch, wenn $i$ simultane Invariante mehrerer Formen mit beliebig vielen Variabeln ist.

\begin{center}
\textbf{§\,10. Partielle Differentialgleichungen f\"ur Co- und Invarianten.}
\end{center}

107. \textit{Aufstellung der Differentialgleichungen, denen eine Invariante $i$ einer einzigen Form $f = a_x^n$ gen\"ugt.} Wir hatten im § 1 als Definitionsgleichung einer Invariante die Gleichung aufgestellt:
\weq{I}{J = (\xi\eta)^{\varrho} \cdot i\,.}
Hiebei war $\varrho$ das Gewicht der Form $i$, und $J$ genau dieselbe ganze und homogene Function in den Coefficienten
\[
\overline{A}_k = a_{\xi}^{n-k} a_{\eta}^k
\]
der linear transformirten Form $F = A_y^n$, welche $i$ in den Coefficienten
\[
\overline{a}_k = a_1^{n-k} a_2^k
\]
der urspr\"unglichen Form $f = a_x^n$ war. Diese Definitionsgleichung (I)
%CONT>>


% ===== pdf idx 134, printed 120.0 =====
% !PRINTED 120
\weberpage{120}
%<<CONT
muss demnach die Quelle f\"ur die zu ermittelnden Differentialgleichungen bilden. Ihre rechte Seite giebt auch zugleich den Weg an, den wir hiebei einzuschlagen haben. Denn w\"ahrend $i$ von den vier willk\"urlichen Transformationsparametern $\xi_1\, \xi_2\, \eta_1\, \eta_2$ v\"ollig unabh\"angig ist, ist der andere Factor $\Delta^{\varrho} = (\xi\eta)^{\varrho}$ eine vollst\"andig bekannte Function derselben, f\"ur welche wir sehr leicht Differentialgleichungen aufstellen k\"onnen. Man hat n\"amlich, wie man unmittelbar erkennt:
\[
\left.
\begin{aligned}
\frac{\partial \Delta}{\partial \xi_1}\, \xi_1 + \frac{\partial \Delta}{\partial \xi_2}\, \xi_2 &= \Delta\\
\frac{\partial \Delta}{\partial \xi_1}\, \eta_1 + \frac{\partial \Delta}{\partial \xi_2}\, \eta_2 &= 0\\
\frac{\partial \Delta}{\partial \eta_1}\, \xi_1 + \frac{\partial \Delta}{\partial \eta_2}\, \xi_2 &= 0\\
\frac{\partial \Delta}{\partial \eta_1}\, \eta_1 + \frac{\partial \Delta}{\partial \eta_2}\, \eta_2 &= \Delta
\end{aligned}
\right\}\quad \text{(A)}
\]
Und diese vier Gleichungen zeigen uns nun, wie erw\"ahnt, die Richtung an, nach welcher wir behufs Aufstellung unserer Differentialgleichung zu gehen haben. Denn indem wir nun die Definitionsgleichung (I) den durch diese vier Gleichungen definirten Processen unterwerfen, erhalten wir unter Benutzung der unter (A) gewonnenen Resultate:
\[
\left.
\begin{aligned}
\frac{\partial J}{\partial \xi_1}\, \xi_1 + \frac{\partial J}{\partial \xi_2}\, \xi_2 &= \varrho \cdot J\\
\frac{\partial J}{\partial \xi_1}\, \eta_1 + \frac{\partial J}{\partial \xi_2}\, \eta_2 &= 0\\
\frac{\partial J}{\partial \eta_1}\, \xi_1 + \frac{\partial J}{\partial \eta_2}\, \xi_2 &= 0\\
\frac{\partial J}{\partial \eta_1}\, \eta_1 + \frac{\partial J}{\partial \eta_2}\, \eta_2 &= \varrho \cdot J
\end{aligned}
\right\}\quad \text{(I)}
\]
Nun enth\"alt aber $J$ die Gr\"ossen $\xi$ und $\eta$ nur in den festen Verbindungen:
\[
\left.
\begin{aligned}
\overline{A}_0 &= a_{\xi}^n\\
\overline{A}_1 &= a_{\xi}^{n-1} a_{\eta}\\
\overline{A}_2 &= a_{\xi}^{n-2} a_{\eta}^2\\
&\cdot\ \cdot\ \cdot\ \cdot\ \cdot\\
\overline{A}_k &= a_{\xi}^{n-k} a_{\eta}^k\\
&\cdot\ \cdot\ \cdot\ \cdot\ \cdot
\end{aligned}
\right\}\quad \text{(B)}
\]
und demnach sind in dem Gleichungssystem (I) die Differentiationen zun\"achst nach den Gr\"ossen $A_k$ auszuf\"uhren, so dass dieses System auch in der Form geschrieben werden kann:
%CONT>>


% ===== pdf idx 135, printed 121.0 =====
% !PRINTED 121
\weberpage{121}
\weq{II}{
\left.\begin{aligned}
\sum_{k=0}^{k=n} \frac{\partial J}{\partial \overline{A}_k}\left(\frac{\partial \overline{A}_k}{\partial \xi_1}\,\xi_1 + \frac{\partial \overline{A}_k}{\partial \xi_2}\,\xi_2\right) &= \varrho J\\
\sum \frac{\partial J}{\partial \overline{A}_k}\left(\frac{\partial \overline{A}_k}{\partial \xi_1}\,\eta_1 + \frac{\partial \overline{A}_k}{\partial \xi_2}\,\eta_2\right) &= 0\\
\sum \frac{\partial J}{\partial \overline{A}_k}\left(\frac{\partial \overline{A}_k}{\partial \eta_1}\,\xi_1 + \frac{\partial \overline{A}_k}{\partial \eta_2}\,\xi_2\right) &= 0\\
\sum \frac{\partial J}{\partial \overline{A}_k}\left(\frac{\partial \overline{A}_k}{\partial \eta_1}\,\eta_1 + \frac{\partial \overline{A}_k}{\partial \eta_2}\,\eta_2\right) &= \varrho J
\end{aligned}\right\}.
}

Da jedoch, wie in Folge der Beziehungen (B), theils direct der
Euler'sche Satz, theils auch ganz einfache Rechnung ergiebt:
\[
\begin{aligned}
\frac{\partial \overline{A}_k}{\partial \xi_1}\,\xi_1 + \frac{\partial \overline{A}_k}{\partial \xi_2}\,\xi_2 &= (n-k)\,\overline{A}_k\\
\frac{\partial \overline{A}_k}{\partial \xi_1}\,\eta_1 + \frac{\partial \overline{A}_k}{\partial \xi_2}\,\eta_2 &= (n-k)\,\overline{A}_{k+1}\\
\frac{\partial \overline{A}_k}{\partial \eta_1}\,\xi_1 + \frac{\partial \overline{A}_k}{\partial \eta_2}\,\xi_2 &= k\,\overline{A}_{k-1}\\
\frac{\partial \overline{A}_k}{\partial \eta_1}\,\eta_1 + \frac{\partial \overline{A}_k}{\partial \eta_2}\,\eta_2 &= k\,\overline{A}_k,
\end{aligned}
\]
so nehmen die Gleichungen (II) unmittelbar die Form an:
\weq{II$_a$}{
\left.\begin{aligned}
\sum (n-k)\cdot \overline{A}_k\ \ \frac{\partial J}{\partial \overline{A}_k} &= \varrho J\\
\sum (n-k)\cdot \overline{A}_{k+1}\,\frac{\partial J}{\partial \overline{A}_k} &= 0\\
\sum k\cdot \overline{A}_{k-1}\,\frac{\partial J}{\partial \overline{A}_k} &= 0\\
\sum k\cdot \overline{A}_k\ \ \frac{\partial J}{\partial \overline{A}_k} &= \varrho J
\end{aligned}\right\}.
}
Sie gelten für jeden Werth von $\xi$ und $\eta$; also auch für $\xi_1 = \eta_2 = 1$, $\xi_2 = \eta_1 = 0$.
Für diese Werthe geht aber $\overline{A}_k$ in $\overline{a}_k$ und demnach $J$ in $i$ über, so dass
endlich die gesuchten Differentialgleichungen in der Form sich ergeben:
\weq{III}{
\left.\begin{aligned}
n\,\frac{\partial i}{\partial \overline{a}_0}\,\overline{a}_0 + (n-1)\,\frac{\partial i}{\partial \overline{a}_1}\,\overline{a}_1 + (n-2)\,\frac{\partial i}{\partial \overline{a}_2}\,\overline{a}_2 + \cdots &= \varrho i\\
n\,\frac{\partial i}{\partial \overline{a}_0}\,\overline{a}_1 + (n-1)\,\frac{\partial i}{\partial \overline{a}_1}\,\overline{a}_2 + (n-2)\,\frac{\partial i}{\partial \overline{a}_2}\,\overline{a}_3 + \cdots &= 0\\
1\,\frac{\partial i}{\partial \overline{a}_1}\,\overline{a}_0 + \phantom{(n{-}1)}2\,\frac{\partial i}{\partial \overline{a}_2}\,\overline{a}_1 + \phantom{(n{-}2)}3\,\frac{\partial i}{\partial \overline{a}_3}\,\overline{a}_2 + \cdots &= 0\\
1\,\frac{\partial i}{\partial \overline{a}_1}\,\overline{a}_1 + \phantom{(n{-}1)}2\,\frac{\partial i}{\partial \overline{a}_2}\,\overline{a}_2 + \phantom{(n{-}2)}3\,\frac{\partial i}{\partial \overline{a}_3}\,\overline{a}_3 + \cdots &= \varrho i
\end{aligned}\right\}.
}


% ===== pdf idx 136, printed 122.0 =====
% !PRINTED 122
\weberpage{122}
Dieses System von partiellen Differentialgleichungen, denen jede Invariante
 $i$ einer Form $f = a_x^n$ genügt, gab bereits Cayley im 47.\ Bd.
des Crelle'schen Journals.

\textit{Anmerkung.} Addirt man in dem Systeme (III) die erste und
vierte Gleichung, so erhält man:
\[
\frac{2\varrho i}{n} = \overline{a}_0\,\frac{\partial i}{\partial \overline{a}_0} + \overline{a}_1\,\frac{\partial i}{\partial \overline{a}_1} + \overline{a}_2\,\frac{\partial i}{\partial \overline{a}_2} + \cdots \overline{a}_n\,\frac{\partial i}{\partial \overline{a}_n}.
\]
Anderntheils ist nach dem Euler'schen Satze, wenn $\mu$ der Grad der
Invariante $i$ in den Coefficienten von $f$ ist,
\[
\mu i = \overline{a}_0\,\frac{\partial i}{\partial \overline{a}_0} + \overline{a}_1\,\frac{\partial i}{\partial \overline{a}_1} + \overline{a}_2\,\frac{\partial i}{\partial \overline{a}_2} + \cdots \overline{a}_n\,\frac{\partial i}{\partial \overline{a}_n}.
\]
Hieraus folgt für das Gewicht $\varrho$ der Invariante $i$ die Beziehung:
\[
\varrho = \frac{\mu \cdot n}{2}.
\]
Da $\varrho$ stets eine ganze Zahl ist, so lehrt diese Gleichung überdies,
dass Formen $f$ ungeraden Grades in $x$ nur Invarianten $i$ geraden Grades
in den Coefficienten von $f$ besitzen können.

108. \gesperrt{Das System der Differentialgleichungen ist ein vollständiges.} Bezeichnen
 wir diese vier Gleichungen (III), indem wir noch in der Gleichung
(1) und (4) die Grösse $\varrho i$ auf die rechte Seite schaffen, mit $A_1 = 0$,
$A_2 = 0$, $A_3 = 0$, $A_4 = 0$, so können wir, wie Clebsch zuerst bemerkt
 hat, zeigen, dass dieses System von Differentialgleichungen ein
vollständiges ist, d.\,h. dass die Anwendung der Differentialgleichungen
auf einander in der Weise, wie sie z.\,B. durch
\[
A_1(A_2) - A_2(A_1) = 0
\]
angedeutet sein mag, zu keinen neuen Differentialgleichungen führt,
sondern nur auf lineare Combinationen der bereits vorhandenen. Es
wird genügen, dies an einem Beispiele zu zeigen, indem wir $i$ etwa
als eine Potenz der Discriminante $\overline{a}_0 \overline{a}_2 - \overline{a}_1{}^2$ einer quadratischen Form
annehmen. Wir haben dann:
\[
\begin{aligned}
A_1 &= 2\,\frac{\partial i}{\partial \overline{a}_0}\,\overline{a}_0 + \frac{\partial i}{\partial \overline{a}_1}\,\overline{a}_1 - 2i = 0\\
A_2 &= 2\,\frac{\partial i}{\partial \overline{a}_0}\,\overline{a}_1 + \frac{\partial i}{\partial \overline{a}_1}\,\overline{a}_2 \phantom{{}-2i{}} = 0.
\end{aligned}
\]
Folglich wird
%CONT>>


% ===== pdf idx 137, printed 123.0 =====
% !PRINTED 123
\weberpage{123}
%<<CONT
\[
\begin{aligned}
0 = {}& A_1(A_2) - A_2(A_1)\\
= {}& 2\overline{a}_0\,\frac{\partial\left[2\dfrac{\partial i}{\partial \overline{a}_0}\,\overline{a}_1 + \dfrac{\partial i}{\partial \overline{a}_1}\,\overline{a}_2\right]}{\partial \overline{a}_0} + \overline{a}_1\,\frac{\partial\left[2\dfrac{\partial i}{\partial \overline{a}_0}\,\overline{a}_1 + \dfrac{\partial i}{\partial \overline{a}_1}\,\overline{a}_2\right]}{\partial \overline{a}_1} - 2i\left[2\,\frac{\partial i}{\partial \overline{a}_0}\,\overline{a}_1 + \frac{\partial i}{\partial \overline{a}_1}\,\overline{a}_2\right]\\
& - 2\overline{a}_1\,\frac{\partial\left[2\dfrac{\partial i}{\partial \overline{a}_0}\,\overline{a}_0 + \dfrac{\partial i}{\partial \overline{a}_1}\,\overline{a}_1 - 2i\right]}{\partial \overline{a}_0} - \overline{a}_2\cdot\frac{\partial\left[2\dfrac{\partial i}{\partial \overline{a}_0}\,\overline{a}_0 + \dfrac{\partial i}{\partial \overline{a}_1}\,\overline{a}_1 - 2i\right]}{\partial \overline{a}_1}\\
= {}& 4\overline{a}_0\overline{a}_1\,\frac{\partial^2 i}{\partial \overline{a}_0{}^2} + 2\overline{a}_0\overline{a}_2\,\frac{\partial^2 i}{\partial \overline{a}_0 \partial \overline{a}_1} + 2\overline{a}_1{}^2\,\frac{\partial^2 i}{\partial \overline{a}_0 \partial \overline{a}_1} + \overline{a}_1\overline{a}_2\,\frac{\partial^2 i}{\partial \overline{a}_1{}^2}\\
& - 4\overline{a}_1 i\,\frac{\partial i}{\partial \overline{a}_0} - 2\overline{a}_2 i\,\frac{\partial i}{\partial \overline{a}_1}\\
& - 4\overline{a}_0\overline{a}_1\,\frac{\partial^2 i}{\partial \overline{a}_0{}^2} - 2\overline{a}_0\overline{a}_2\,\frac{\partial^2 i}{\partial \overline{a}_0 \partial \overline{a}_1} - 2\overline{a}_1{}^2\,\frac{\partial^2 i}{\partial \overline{a}_0 \partial \overline{a}_1} - \overline{a}_1\overline{a}_2\,\frac{\partial^2 i}{\partial \overline{a}_1{}^2}\\
& + 4\overline{a}_1\,\frac{\partial i}{\partial \overline{a}_0} + 2\overline{a}_2\,\frac{\partial i}{\partial \overline{a}_1}\\
= {}& -2i\left\{2\,\frac{\partial i}{\partial \overline{a}_0}\,\overline{a}_1 + \frac{\partial i}{\partial \overline{a}_1}\,\overline{a}_2\right\} + 2\left\{\frac{\partial i}{\partial \overline{a}_0}\,\overline{a}_1 + \frac{\partial i}{\partial \overline{a}_1}\,\overline{a}_2\right\}.
\end{aligned}
\]
Man erkennt, dass die ganze rechte Seite sich reducirt auf
\[
0 = -2i\,A_2 + 2A_2 = 2A_2\cdot(1-i),
\]
d.\,h. die Differentialgleichung $A_1(A_2) - A_2(A_1) = 0$ ist dieselbe, wie
$A_2 = 0$.

\indent Den allgemeinen Beweis hat Clebsch in seinen \glqq Binären Formen\grqq{}
Seite 310 gegeben.

\textit{Anmerkung.} Zu den vier Differentialgleichungen treten nur dann
neue hinzu, wenn es sich darum handelt, specielle Invarianten zu
charakterisiren. So hat Brioschi noch specielle Differentialgleichungen
aufgestellt, um Discriminanten und Resultanten zu charakterisiren.
(Vgl. auch die Untersuchungen Nöther's in Faà di Bruno, übersetzt
von Walter, Seite 275.)

109. \gesperrt{Differentialgleichungen für simultane Invarianten und Covarianten
.} Es hat nun keine Schwierigkeit, ein System von Differentialgleichungen
 für simultane Invarianten $i$ direct anzuschreiben. Ist
nämlich $i$ simultane Invariante der $s$ Formen
\[
f_1,\ f_2,\ f_3,\ \cdots f_s
\]
von den Graden $n_1$, $n_2$, $n_3$, $\ldots n_s$ in $x$, welche bezw. die Coefficienten
%CONT>>


% ===== pdf idx 138, printed 124.0 =====
% !PRINTED 124
\weberpage{124}
%<<CONT
\[
a_k^{(1)},\quad a_k^{(2)},\quad \overline{a}_k^{(3)},\quad \cdots a_k^{(s)}
\]
besitzen, so geht aus der ganzen Ableitung dieser Differentialgleichungen
 direct hervor, dass sie für diesen Fall die Gestalt annehmen
 müssen:
\weq{IV}{
\left.\begin{aligned}
\sum_{\lambda=1}^{\lambda=s}\left(n\cdot\frac{\partial i}{\partial \overline{a}_0{}^{(\lambda)}}\,\overline{a}_0^{(\lambda)} + (n-1)\,\frac{\partial i}{\partial \overline{a}_1{}^{(\lambda)}}\,\overline{a}_1^{(\lambda)} + (n-2)\,\frac{\partial i}{\partial \overline{a}_2{}^{(\lambda)}}\,\overline{a}_2^{(\lambda)}\ldots\right) &= \varrho\cdot i\\
\sum_{\lambda=1}^{\lambda=s}\left(n\cdot\frac{\partial i}{\partial \overline{a}_0{}^{(\lambda)}}\,\overline{a}_1^{(\lambda)} + (n-1)\,\frac{\partial i}{\partial \overline{a}_1{}^{(\lambda)}}\,\overline{a}_2^{(\lambda)} + \cdots\cdots\cdots\cdots\right) &= 0\\
\sum_{\lambda=1}^{\lambda=s}\left(1\cdot\frac{\partial i}{\partial \overline{a}_1{}^{(\lambda)}}\,a_0^{(\lambda)} + \phantom{(n-1)}2\cdot\frac{\partial i}{\partial \overline{a}_2{}^{(\lambda)}}\,a_1^{(\lambda)} + \cdots\cdots\cdots\right) &= 0\\
\sum_{\lambda=1}^{\lambda=s}\left(1\cdot\frac{\partial i}{\partial \overline{a}_1{}^{(\lambda)}}\,\overline{a}_1^{(\lambda)} + \phantom{(n-1)}2\cdot\frac{\partial i}{\partial \overline{a}_2{}^{(\lambda)}}\,\overline{a}_2^{(\lambda)} + \cdots\cdots\cdots\cdots\right) &= \varrho\cdot i
\end{aligned}\right\}.
}
% [sic] third row of (IV) prints a_0^{(\lambda)}, a_1^{(\lambda)} without the overline bar
% that the other rows carry on the corresponding letters; reproduced as printed.
Da man nun an Stelle jeder Covariante\footnotemark\ eine Invariante einführen
kann, bei deren Bildung das simultane System um ebensoviel lineare
Formen vermehrt ist, als Reihen von Veränderlichen in der Covariante
existiren, so genügt jede Covariante ebenfalls einem simultanen
Systeme (IV) von Differentialgleichungen. Will man in ihnen die
Glieder absondern, welche die Differentialquotienten nach den Veränderlichen
 enthalten, so denke man sich einen Augenblick die
Coefficienten $l^{(i)}$ der linearen Formen an Stelle der Variabeln eingeführt
. Dann treten links den vier Gleichungen entsprechend die
Summen auf:
\[
\sum l_1\,\frac{\partial i}{\partial l_1},\quad \sum l_2\,\frac{\partial i}{\partial l_1},\quad \sum l_1\,\frac{\partial i}{\partial l_2},\quad \sum l_2\,\frac{\partial i}{\partial l_2}
\]
oder, wenn wir nun wieder $l_1$ durch $x_2$, und $l_2$ durch $-x_1$ ersetzen:
\[
\sum x_2\,\frac{\partial i}{\partial x_2},\quad -\sum x_1\,\frac{\partial i}{\partial x_2},\quad -\sum x_2\,\frac{\partial i}{\partial x_1},\quad \sum x_1\,\frac{\partial i}{\partial x_1},
\]
oder endlich, wenn $i$ vom Grade $\sum \nu_i$ in den verschiedenen Variabelnreihen
 ist, gemäss dem Euler'schen Satze:
\[
\begin{aligned}
i\sum \nu_i - \sum x_1\,\frac{\partial i}{\partial x_1}, &\quad -\sum x_1\,\frac{\partial i}{\partial x_2},\\
-\sum x_2\,\frac{\partial i}{\partial x_1}, &\quad i\sum \nu_i - \sum x_2\,\frac{\partial i}{\partial x_2}.
\end{aligned}
\]
\footnotetext{Vgl. auch \gesperrt{Clebsch}, \glqq Binäre Formen\grqq, S. 315 und diese Vorles. Nr. 100.}


% ===== pdf idx 139, printed 125.0 =====
% !PRINTED 125
\weberpage{125}
Das System (IV) der Differentialgleichungen nimmt also, wenn man dort noch $\varrho+\Sigma\nu$ durch $\varrho$ ersetzt, für Covarianten auch die Formen an:

\[
\left.
\begin{aligned}
&\sum_{1}^{s}\sum_{0}^{n_s-1}(n_s-k)\frac{\partial i}{\partial \bar a_k^{(\lambda)}}\bar a_k^{(\lambda)} - \sum x_1\frac{\partial i}{\partial x_1} = i\cdot\varrho\\
&\sum_{1}^{s}\sum_{0}^{n_s-1}(n_s-k)\frac{\partial i}{\partial \bar a_k^{(\lambda)}}\bar a_{k+1}^{(\lambda)} - \sum x_1\frac{\partial i}{\partial x_2} = 0\\
&\sum_{1}^{s}\sum_{1}^{n_s} k\cdot\frac{\partial i}{\partial \bar a_k^{(\lambda)}}\bar a_{k-1}^{(\lambda)} - \sum x_2\frac{\partial i}{\partial x_1} = 0\\
&\sum_{1}^{s}\sum_{1}^{n_s} k\cdot\frac{\partial i}{\partial \bar a_k^{(\lambda)}}\bar a_k^{(\lambda)} - \sum x_2\frac{\partial i}{\partial x_2} = i\cdot\varrho
\end{aligned}
\right\}\ \text{(V)}
\]

Ist nun $i$ von den Graden $\mu_1,\mu_2\ldots\mu_s$ in den Co\"efficienten der Formen $f_1,f_2\cdots f_s$, so erhält man durch Addition der ersten und letzten Gleichung nach Anwendung des Euler'schen Satzes:

\[
\frac{1}{2}\left\{\sum_1^s \mu_i\nu_i - \sum \nu_i\right\} = \varrho.
\]

Ist insbesondere $s=1$, also $i$ Covariante einer einzigen Form $f=a_x^n$ mit einer Reihe von Variabeln, dann besteht wiederum zwischen Gewicht $\varrho$, dem Grade $\nu$ in $x$, dem Grade $\mu$ in den Co\"efficienten die Beziehung

\[
\varrho = \frac{1}{2}(\mu n - \nu).
\]

Daraus erkennt man, dass Formen ungeraden Grades keine Covarianten besitzen, die gleichzeitig ungeraden Grades in den Co\"efficienten und geraden Grades in den Variabeln sind, oder umgekehrt.

110. \gesperrt{Die Differentialgleichungen (IV) resp. (V) sind nicht nur nothwendig, sondern auch hinreichend zur Definition einer Invariante, resp. Covariante.} Wir haben bisher gezeigt, dass jede ganze homogene Function $i$, welche der Gleichung

\[
J = i\cdot(\xi\eta)^\varrho \tag{I}
\]

genügt, stets ein System von Differentialgleichungen befriedigt. Es lässt sich nun aber auch umgekehrt der Beweis liefern, dass jede Lösung $i$ der Differentialgleichungen (IV) eine Invariante ist, d. h. dass für sie die Relation (I) besteht, in welcher $J$ aus $i$ durch die Vertauschung von $a_1$ mit $a_\xi$, $a_2$ mit $a_\eta$ etc. hervorgegangen ist.


% ===== pdf idx 140, printed 126.0 =====
% !PRINTED 126
\weberpage{126}
Um diesen Beweis zu erbringen, nehmen wir an: $J$ sei eine ganze homogene Function der Co\"efficienten $\bar A_k = a_\xi^{n-k}a_\eta^k$ etc., welche dem Systeme (IV) genügt, oder also, was dasselbe ist, dem Systeme (I). In gleicher Weise sei $K$ eine beliebige zweite Lösung dieser Gleichungen (I). Dann bestehen zunächst die beiden Systeme von Gleichungen:

\[
\begin{array}{r@{\ }l@{\qquad}r@{\ }l}
& \dfrac{\partial J}{\partial \xi_1}\xi_1+\dfrac{\partial J}{\partial \xi_2}\xi_2=\varrho J
& & \dfrac{\partial K}{\partial \xi_1}\xi_1+\dfrac{\partial K}{\partial \xi_2}\xi_2=\varrho K\\[6pt]
(1)\quad & \dfrac{\partial J}{\partial \xi_1}\eta_1+\dfrac{\partial J}{\partial \xi_2}\eta_2=0
& (2)\quad & \dfrac{\partial K}{\partial \xi_1}\eta_1+\dfrac{\partial K}{\partial \xi_2}\eta_2=0\\[4pt]
& \cdots\cdots
& & \cdots\cdots\\[4pt]
& \dfrac{\partial J}{\partial \eta_1}\eta_1+\dfrac{\partial J}{\partial \eta_1}\eta_2=\varrho J
& & \dfrac{\partial K}{\partial \eta_1}\eta_1+\dfrac{\partial K}{\partial \eta_2}\eta_2=\varrho K.
\end{array}
\]
% [sic] printed second term of the last J-row as $\partial J/\partial\eta_1$ (not $\partial J/\partial\eta_2$), unlike the parallel K-row

Multipliciren wir das erste mit $K$, das zweite mit $J$ und subtrahiren je zwei in gleicher Zeile befindliche Gleichungen von einander, so bekommen wir, wenn wir die erhaltenen Differenzen noch mit $K^2$ dividiren:

\[
\frac{K\dfrac{\partial J}{\partial \xi_1}-J\dfrac{\partial K}{\partial \xi_1}}{K^2}\,\xi_1+\frac{K\dfrac{\partial J}{\partial \xi_2}-J\dfrac{\partial K}{\partial \xi_2}}{K^2}\,\xi_2=0
\]
\[
\frac{K\dfrac{\partial J}{\partial \xi_1}-J\dfrac{\partial K}{\partial \xi_1}}{K^2}\,\eta_1+\frac{K\dfrac{\partial J}{\partial \xi_2}-J\dfrac{\partial K}{\partial \xi_2}}{K^2}\,\eta_2=0\text{, etc.}
\]

Setzen wir hierin:

\[
\frac{J}{K} = L\ldots \tag{c}
\]

so kommt:

\[
\frac{\partial L}{\partial \xi_1}\xi_1+\frac{\partial L}{\partial \xi_2}\xi_2=0,\quad \frac{\partial L}{\partial \xi_1}\eta_1+\frac{\partial L}{\partial \xi_2}\eta_2=0 \tag{1}
\]

\[
\frac{\partial L}{\partial \eta_1}\xi_1+\frac{\partial L}{\partial \eta_2}\xi_2=0,\quad \frac{\partial L}{\partial \eta_1}\eta_1+\frac{\partial L}{\partial \eta_2}\eta_2=0. \tag{2}
\]

Jedes der beiden Gleichungssysteme (1) und (2) ist homogen in den Differentialquotienten; da die Determinante $(\xi\eta)$ nicht verschwindet, so haben sie nur die identischen Lösungen (vgl. Bd. I, Nr. 96 und 97)

\[
\frac{\partial L}{\partial \xi_1}=0,\quad \frac{\partial L}{\partial \xi_2}=0,\quad \frac{\partial L}{\partial \eta_1}=0,\quad \frac{\partial L}{\partial \eta_2}=0,
\]

d. h. die Grösse $L$ ist von $\xi$ und $\eta$ unabhängig. Nun ist aber $K$ eine beliebige specielle Lösung der Gleichungen (I), also etwa gleich $\Delta^\varrho$, von welcher wir wissen, dass sie eine solche ist. Benutzen wir dieselbe, so ist wegen (c)

\[
J = \Delta^\varrho \cdot L.
\]


% ===== pdf idx 141, printed 127.0 =====
% !PRINTED 127
\weberpage{127}
Setzen wir hier, da $\xi_1\eta_1\xi_2\eta_2$, beliebig sind, $\xi_1=1$, $\eta_2=1$, $\xi_2=0$, $\eta_2=0$, so geht $a_\xi$ über in $a_1$, $a_\eta$ in $a_2$, $\Delta$ in $1$, und $J$ in $i$, so dass wir erhalten
% [sic] the second condition is printed $\eta_2=0$, duplicating the earlier $\eta_2=1$

\[
L = i,
\]

d. h. jede Lösung der Differentialgleichung (IV) genügt der Relation

\[
J = \Delta^\varrho \cdot i,
\]

so dass dieselben in der That nur Invarianten definiren.

111. \gesperrt{Begriff der Evectante.} Wir können die Differentialgleichungen für Invarianten auch direct aufstellen, ohne dabei von der Transformation Gebrauch zu machen. Zu dem Zwecke führen wir aber zunächst einen neuen Begriff ein, nämlich den der Evectante.

Wir haben schon mehrfach von dem Satze Gebrauch gemacht, dass man in einem symbolischen Producte jederzeit $a_1$ durch $x_2$ und $a_2$ durch $-x_1$ ersetzen darf, wodurch die Invarianteneigenschaft desselben nicht aufgehoben wird. Nun ist nach dem Euler'schen Satze

\[
\mu i = \frac{\partial i}{\partial \bar a_0}\bar a_0+\frac{\partial i}{\partial a_1}\bar a_1+\frac{\partial i}{\partial \bar a_2}\bar a_2+\cdots \frac{\partial i}{\partial \bar a_n}\bar a_n.
\]
% [sic] second term printed $\partial i/\partial a_1$ without the bar, unlike the other terms

Ersetzt man jeden zweiten Factor $\bar a_k$ der Summe rechts durch sein Symbol $a_1^n$, $a_1^{n-1}a_2$, $a_1^{n-2}a_2^2$, und in diesen Symbolen wieder $a_1$ durch $x_2$, $a_2$ durch $-x_1$, so erhält man als Resultat dieser Substitutionen die Covariante

\[
p_x^n = x_2^n\frac{\partial i}{\partial \bar a_0} - x_2^{n-1}x_1\frac{\partial i}{\partial \bar a_1} + x_2^{n-2}x_1^2\frac{\partial i}{\partial \bar a_2} - \cdots
\]

oder

\[
p_x^n = \sum(-1)^k\cdot x_2^{n-k}\cdot x_1^k\cdot\frac{\partial i}{\partial \bar a_k}.
\]

Man nennt diese Covariante Evectante von $i$, und zwar die erste Evectante, insofern $i$ nur einmal nach den Co\"efficienten $\bar a_k$ differentiirt wurde.

Die $r^{\text{te}}$ Evectante entsteht alsdann durch $r$-malige Wiederholung des angegebenen Processes.

112. \gesperrt{Beispiel:} Die Discriminante der quadratischen Form $a_x^2$ ist

\[
(ab)^2 = 2(\bar a_0\bar a_2-\bar a_1^2)=i;
\]

also ist nach dem Euler'schen Satze:

\[
2i=\left[\frac{\partial i}{\partial \bar a_0}\bar a_0+\frac{\partial i}{\partial \bar a_1}\bar a_1+\frac{\partial i}{\partial \bar a_2}\bar a_2\right]=\left[\frac{\partial i}{\partial \bar a_0}a_1^2+\frac{\partial i}{\partial \bar a_1}a_1a_2+\frac{\partial i}{\partial \bar a_2}a_2^2\right]
\]

und somit, wenn man rechts $a_1$ durch $x_2$, $a_2$ durch $-x_1$ ersetzt:

\[
p_x^2 = 2(\bar a_2x_2^2+2\bar a_1x_1x_2+\bar a_0x_1^2)=2a_x^2=2b_x^2,
\]

d. h. die erste Evectante der Discriminante ist die Form selbst.


% ===== pdf idx 142, printed 128.0 =====
% !PRINTED 128
\weberpage{128}
2) Die Discriminante der cubischen Form $a_x^3$ ist

\[
i=(ab)^2(cd)^2(ad)(cb)=[\bar a_0^2\bar a_3^2+4\bar a_0\bar a_2^3-6\bar a_0\bar a_1\bar a_2\bar a_3+4\bar a_1^3\bar a_3-3\bar a_1^2a_2^2]2;
\]
% [sic] final term printed $a_2^2$ without the bar, unlike the other terms of the bracket

folglich ist

\[
4i=\frac{\partial i}{\partial \bar a_0}a_1^3+\frac{\partial i}{\partial \bar a_1}a_1^2a_2+\frac{\partial i}{\partial \bar a_2}a_1a_2^2+\frac{\partial i}{\partial \bar a_3}a_2^3;
\]

demnach

\[
\begin{aligned}
p_x^3 = {} & [(2\bar a_0a_3^2+4a_2^3-6a_1a_2a_3)x_2^3+(6a_0\bar a_2\bar a_3-12\bar a_1^2\bar a_3+6\bar a_1\bar a_2^2)x_2^2x_1\\
& {}+(12\bar a_0\bar a_2^2-6\bar a_0a_1a_3-6a_1^2\bar a_2)x_2x_1^2+(2\bar a_0^2\bar a_3-6\bar a_0a_1a_2+4a_1^3)x_1^3]2\\
= {} & 4(ab)^2(cb)c_x^2a_x.
\end{aligned}
\]

Diese Evectante ist aber nichts anderes als die Covariante dritten Grades von $f=a_x^3$, welche auch als Functionaldeterminante von $f$ mit der Hesse'schen Form von $f$ sich ergiebt. (Vgl. Nr. 37.)

113. \gesperrt{Symbolische Darstellung des Evectantenprocesses.} Betrachten wir die beiden gegebenen Beispiele näher, so sieht man: durch den Evectantenprocess wurde aus $(ab)^2$ die ursprüngliche Form $f=2a_x^2=2b_x^2$. Dieselbe entsteht auch direct aus $(ab)^2$, indem man einmal die Symbolreihe $a_2,a_1$ durch $x_1,-x_2$, und dann aber auch die Symbolreihe $b_2,b_1$ durch $x_1,-x_2$ ersetzt. Die erste Operation giebt ein Glied $g_1$, die zweite ein Glied $g_2$, und es ist

\[
f=g_1+g_2=b_x^2+a_x^2=2a_x^2=2b_x^2.
\]

Ebenso erhält man im zweiten Falle durch successive Vertauschung der vier Symbole $a,b,c,d$ mit den Variabeln $x_1$ und $-x_2$ vier Glieder $g_1g_2g_3g_4$, und zwar, da alle vier Symbole vollständig symmetrisch auftreten, so ist:

\[
g_1+g_2+g_3+g_4=4g_1=4g_2=4g_3=4g_4,
\]

also ist

\[
p_x^3=4(ab)^2(cb)c_x^2a_x.
\]

Wir schliessen: \glqq{}Die erste Evectante entsteht aus der Invariante $\mu^{\text{ten}}$ Grades, indem wir aus ihr $\mu$ Glieder $g_1g_2\ldots g_\mu$ dadurch bilden, dass wir immer je eine Symbolreihe durch die Variabelnreihe der $x$ ersetzen, und dann die Summe aller Glieder nehmen.\grqq{} (Vgl. auch die Wirkung des Aronhold'schen Processes an einem symbolischen Producte, Nr. 60.)

114. \gesperrt{Umformung der Differentialgleichungen für Invarianten.} Wir können nun die Differentialgleichungen (III) Nr. 107 für Invarianten auch in kürzerer Form schreiben. Vergleicht man in der Identität

\[
p_x^n=\sum(-1)^kx_1^kx_2^{n-k}\frac{\partial i}{\partial \bar a_k}
\]
%CONT>>


% ===== pdf idx 143, printed 129.0 =====
% !PRINTED 129
\weberpage{129}
%<<CONT
die Coefficienten gleich hoher Potenzen von $x$, so erhält man:
\[
\binom{n}{k} p_1^k p_2^{n-k} = \frac{\partial i}{\partial \bar a_k} (-1)^k.
\]
Tragen wir diese Werthe der Differentialquotienten in die Gleichungen (III) Nr.~107 ein, so kommt:
\[
\begin{aligned}
\varrho i &= \sum (n-k) \binom{n}{k} (-1)^k a_1^{n-k} a_2^k p_1^k p_2^{n-k}\\
&= \sum n \binom{n-1}{k} (-1)^k a_1^{n-k} a_2^k p_1^k p_2^{n-k}, \quad \text{etc.}
\end{aligned}
\]
oder:
\weq{VI}{
\left.
\begin{aligned}
\varrho i &= -\, n\, a_1 p_2\, (pa)^{n-1}\\
0 &= \phantom{-\,} n\, a_1 p_1\, (pa)^{n-1}\\
0 &= \phantom{-\,} n\, a_2 p_2\, (pa)^{n-1}\\
\varrho i &= +\, n\, a_2 p_1\, (pa)^{n-1}
\end{aligned}
\right\}
}

Dieses System (VI) lässt sich noch kürzer in die beiden folgenden Gleichungen zusammenfassen:
\weq{VII}{
\left.
\begin{aligned}
\frac{2\varrho i}{n} = \mu i &= (pa)^n\\
0 &= (pa)^{n-1} p_x a_x
\end{aligned}
\right\}
}

Die erste derselben ergiebt sich direct durch Addition der ersten und vierten Gleichung des Systemes (VI). Sie lehrt:
\glqq Die $n^{\text{te}}$ Ueberschiebung der ursprünglichen Form $f = a_x^n$ mit der ersten Evectante einer ihrer Invarianten $i$ liefert wiederum die Invariante $i$.\grqq{}

Die zweite der Gleichungen (VII) erhalten wir, wenn wir die vier Gleichungen (VI) der Reihe nach mit $-x_1x_2$, $x_1^2$, $x_2^2$, $x_1x_2$ multipliciren und addiren. Sie lehrt:
\glqq Die $(n-1)^{\text{te}}$ Ueberschiebung der ursprünglichen Form $f=a_x^n$ mit der ersten Evectante einer ihrer Invarianten ist identisch null.\grqq{}

115. \textit{Evectanten von Combinanten.} Ist $i$ eine Combinante der beiden Formen $f=a_x^n$, $\varphi=\alpha_x^n$, genügt also $i$ der Gleichung
\weq{1}{
\delta i = \sum \frac{\partial i}{\partial \bar a_k} \bar\alpha_k = \sum \frac{\partial i}{\partial \bar\alpha_k} \bar a_k = 0,
}
so liefert die $n^{\text{te}}$ Ueberschiebung der ersten Evectante von $i$ über eine der Grundformen nicht mehr $i$, wie im vorigen einfachen Falle. Vielmehr hat man nun:
\weq{2}{
(p\alpha)^n = (pa)^n = 0;
}
denn es ist:
\[
p_x^n = \sum (-1)^k \frac{\partial i}{\partial \bar a_k} x_1^k x_2^{n-k},
\]
also
\[
\frac{\partial i}{\partial \bar a_k} = (-1)^k \binom{n}{k} \bar p_{n-k}
\]


% ===== pdf idx 144, printed 130.0 =====
% !PRINTED 130
\weberpage{130}
Substituirt man den Werth dieser partiellen Differentialquotienten in (1), so erhält man:
\[
\sum_{k=n}^{k=0} (-1)^k \binom{n}{k} \bar p_{n-k} \bar\alpha_k = (\alpha p)^n = 0.
\]
Eine solche Combinante ist insbesondere die Resultante $R_{f,\varphi}$ der beiden Formen $f$ und $\varphi$. Für sie können wir die Richtigkeit der Gleichung (2) auch noch in anderer Weise zeigen, wodurch zugleich auf die Combinanteneigenschaft von $R_{f,\varphi}$ ein neues Licht geworfen wird. Man weiss nämlich (vergl. auch Bd.~I Nr.~164), dass die partiellen Differentialquotienten $\dfrac{\partial R_{f,\varphi}}{\partial \bar a_i}$ den Potenzen der gemeinschaftlichen Wurzel $\xi = -\dfrac{\xi_2}{\xi_1}$ von $f$ und $\varphi$ proportional sind. Ersetzt man also in der Evectante
\[
p_x^n = \sum (-1)^k \frac{\partial R_{f,\varphi}}{\partial \bar a_k} x_1^k x_2^{n-k}
\]
dementsprechend die Grössen $\dfrac{\partial R_{f,\varphi}}{\partial \bar a_k}$ durch $c\cdot(-1)^k\cdot\xi^{n-k}$, so kommt:
\[
p_x^n = \bar c \cdot (\xi_1 x_1 + \xi_2 x_2)^n = \bar c \cdot \xi_x^n.
\]
Folglich ist:
\[
(\varphi, p)^n = \bar c \cdot (\varphi, \xi^n)^n = \bar c \cdot \varphi(\xi).
\]
Da aber $\varphi(x)$ für die gemeinschaftliche Wurzel $x=\xi$ verschwindet, so ist $\varphi(\xi)=0$ und damit
\[
(\varphi, p)^n = 0.
\]
Ist nun aber $\chi = r_x^n$ neben $f=a_x^n$ und $\varphi=\alpha_x^n$ eine dritte Form, die gleichfalls für $x=\xi$ verschwindet, so ist auch
\[
(\chi, p)^n = \bar c \cdot (\chi, \xi^n)^n = 0.
\]
Es ist aber:
\[
(\chi, p)^n = \sum (-1)^k \binom{n}{k} \bar p_{n-k} \bar r_k = \sum \frac{\partial R_{f,\varphi}}{\partial \bar a_k} \cdot \bar r_k = \bar\delta R_{f,\varphi}.
\]
Demnach ist auch:
\[
\bar\delta R_{f,\varphi} = 0,
\]
wobei nun die Incremente durch die Coefficienten einer völlig fremden Form $\chi$ ersetzt werden, für welche nur die Bedingung besteht, dass sie mit $f$ und $\varphi$ eine gemeinschaftliche Wurzel $\xi$ besitzt. Wir werden später von diesem Satze Gebrauch machen.

116. \textit{Zweite Methode zur Aufstellung der Differentialgleichungen für Invarianten.} Das in Nr.~114 erhaltene System (VII) von Gleichungen ist nur eine andere Form des Systemes der Differentialgleichungen
%CONT>>


% ===== pdf idx 145, printed 131.0 =====
% !PRINTED 131
\weberpage{131}
%<<CONT
(III) Nr.~107 für Invarianten. Indem wir dasselbe direct aufstellen, erhalten wir einen zweiten Beweis, dass alle Invarianten gewissen Differentialgleichungen genügen müssen.

Die Aufgabe ist, die Richtigkeit der beiden Gleichungen nachzuweisen
\weq{A}{
\left.
\begin{aligned}
(p,f)^n &= \mu \cdot i\\
(p,f)^{n-1} &= 0
\end{aligned}
\right\}
}
Nun ist die erste Evectante $p$ von $i$ eine Summe symbolischer Producte $g$, welche entstehen, wenn jede der $\mu$ Symbolreihen $a$, resp. $b$, $c \ldots$ etc. durch $x$ ersetzt wird. Es ist also:
\weq{2}{
(p,f)^n = \sum_1^\mu (g,f)^n.
}
Anderntheils entsteht aber die $n^{\text{te}}$ Ueberschiebung zweier Formen, indem man umgekehrt die Variabeln $x_1 x_2$ durch Symbolreihen $a_2, -a_1$, resp. $b_2, -b_1$ etc. ersetzt. Beide Operationen hinter einander ausgeführt, der Evectantenprocess und der Ueberschiebungsprocess, ändern also das ursprüngliche symbolische Product nicht; jedes Glied der Summe (2) geht durch die Ueberschiebung wieder in $i$ über, so dass man in der That erhält:
\[
(p,f)^n = \mu \cdot i.
\]
Was nun die zweite Gleichung betrifft, so gestaltet sich ihr Beweis folgendermassen. Es ist wiederum zunächst
\[
(p,f)^{n-1} = \sum_1^\mu (g,f)^{n-1}.
\]
Jede der $\mu$ Ueberschiebungen rechts besitzt $n$ Terme. Denn ist z.\,B. $g_1$ jenes Glied der Evectante, das entstanden ist, indem man für die Symbolreihe $a$ die Variabeln $x$ einführte, so entsteht nach den Gesetzen der Ueberschiebung $(g_1,f)^{n-1}$ dadurch, dass man in immer andern und andern $(n-1)$ Factoren erster Art von $g_1$ umgekehrt $x_1$ und $x_2$ durch $a_2, -a_1$ ersetzt, was bei $n$ Elementen nur auf $n$ Arten möglich ist. Die Summe rechts enthält also $n \cdot \mu$ Terme, wobei $n \cdot \mu$ jedenfalls eine gerade Zahl ist, wie aus Nr.~107, Anmerkung hervorgeht. Jeder der $n \cdot \mu$ Terme enthält zwei und nur zwei Factoren erster Art, da $f$ sowohl als jedes Glied $g_i$ vom Grade $n$ in $x$ ist. Man kann einen solchen Term auch erhalten, indem man z.\,B. in der Invariante $i$ in einem Klammerfactor $(ba)$ die Symbolreihe $a_1, a_2$ durch $-x_2, +x_1$ ersetzt und das so entstehende symbolische Product mit $a_x$ multiplicirt, so dass also $(ba)$ durch $b_x a_x$ ersetzt wurde. Es tritt
%CONT>>


% ===== pdf idx 146, printed 132.0 =====
% !PRINTED 132
\weberpage{132}
%<<CONT
aber alsdann bei dieser Bildung der $n \cdot \mu$ Terme, die ja völlig symmetrisch erfolgt, jedenfalls auch noch ein zweiter Term neben dem eben erwähnten auf, der entsteht, wenn in demselben Klammerfactor $(ba)$ von $i$ die Symbolreihe $b_1, b_2$ durch $-x_2, x_1$ ersetzt und das so entstehende symbolische Product mit $b_x$ multiplicirt wird. Da dieser Term aber bis auf das Vorzeichen mit dem vorhin erwähnten übereinstimmt, so sieht man, dass je zwei der $n \cdot \mu$ Terme sich gegenseitig aufheben, d.\,h. dass
\[
(p,f)^{n-1} = \sum_1^\mu (g,f)^{n-1} = 0,
\]
was wir zu beweisen hatten.

Wir haben somit durch reine Invariantenprocesse die Existenz der Gleichungen (A) nachgewiesen, welche das System der Differentialgleichungen (III) nach jeder Richtung ersetzen, und da jener Satz in Nr.~108, dass dieses System von Differentialgleichungen ein vollständiges ist, zu seinem Beweise nur Processe der Invariantentheorie benutzt, so dürfte es jedenfalls möglich sein, auch ihn durch ein rein algebraisches Aequivalent zu ersetzen.

117. \textit{Umformung einer Relation $F(A_i)=0$ in}
\[
\sum \{(ab)c_x + (bc)a_x + (ca)b_x\} \varphi_\lambda = 0.
\]
Zum Schlusse noch eine Anwendung des Evectantenprocesses. Ich habe bereits in \S.\,1 an einem Beispiele gezeigt, wie sich eine in Coefficienten und Variabeln homogene Relation $F(A_i)=0$ von symbolischen Producten $A_i$ in die Gestalt bringen lässt
\[
\sum \{(ab)c_x + (bc)a_x + (ca)b_x\} = 0.
\]
Die Möglichkeit dieser Umformung wollen wir nun für eine beliebige Relation $F(A_i)=0$ nachweisen. Wir betrachten $F(A_i)$ als Covariante $\mu^{\text{ten}}$ Grades in den Coefficienten einer Form $f=a_x^n=b_x^n=$ etc.; da sie identisch verschwindet, so verschwinden auch die partiellen Differentialquotienten in der nach dem Euler'schen Satze gegebenen Darstellung:
\[
\mu F(A_i) = \sum \frac{\partial (F(A_i))}{\partial \bar a_k} \cdot \bar a_k,
\]
und demnach auch die erste Evectante dieser Covariante
\[
P_y^n = \sum (-1)^k \frac{\partial (F(A_i))}{\partial \bar a_k} \cdot y_1^k y_2^{n-k}.
\]
Nun erhält man diese Evectante auch, wenn man in $F(A_i)$ der Reihe nach das Symbol $a$, $b$, $c$ etc. durch $y$ ersetzt; es ist also:
\[
0 = P_y^n = F(A_i)_{a=y} + F(A_i)_{b=y} + \cdots + F(A_i)_{m=y} = \sum_{m=a,b,\ldots} F(A_i)_{m=y}.
\]


% ===== pdf idx 147, printed 133.0 =====
% !PRINTED 133
\weberpage{133}
Hiebei ist es gestattet, nachträglich rechts im zweiten, dritten, etc.,
letzten Posten, $a$ mit $b$, resp. $a$ mit $c$, etc. zu vertauschen, so dass
in keinem Gliede mehr das Symbol $a$ auftritt. Wir setzen nun voraus,
dass sich einfachere Relationen $F(A_i)$ — also insbesondere solche,
welche mindestens \gesperrt{ein} Symbol $a$, $b$, etc. weniger besitzen, oder
solche, welche durch Spaltung in zwei Factoren, deren einer mindestens
\gesperrt{ein} Symbolfactor $r_x$ oder $(pq)$ ist, sich auf einfachere Relationen
reduciren, — stets in die Form $\sum \{(ab)\,c_x + (bc)\,a_x + (ca)\,b_x\}\,\varphi_\lambda = 0$
bringen lassen. Der Beweis gestaltet sich dann für complicirtere Fälle
folgendermassen.

$P_y^n$ wird sich im allgemeinen in zwei Theile trennen lassen, einen
Theil $\mathbf{M} = (xy) \sum Q_i^{(1)}$, der also $(xy)$ zum Factor hat, und einen
Theil

\weq{1}{\mathbf{N} = \sum c_i\,G_i,}

der diesen Factor nicht besitzt. Sollten zufälliger Weise in $P_y^n = 0$
alle Glieder diesen Factor enthalten, so dividire man mit der höchsten
Potenz von $(xy)$ diese Identität.

Setzen wir alsdann in $P_y^n = 0$ einen Moment $y = x$, so verschwindet
wegen des Factors $(xx)$ der erste Theil, und demnach auch der zweite,
d.\,h. wir haben

\weq{2}{\mathbf{N} = \sum c_i\,H_i = 0,}

wenn wir die Glieder $G_i$ für $x = y$ mit $H_i$ bezeichnen. Da aber die
Glieder $H_i$ um das Symbol $a$ weniger Symbole besitzen als die ursprünglichen
Glieder in $F(A_i)$, so kann man nach Voraussetzung schreiben:

\weq{3}{\mathbf{N} = \sum c_i\,H_i = \sum \{(db)\,c_x + (bc)\,d_x + (cd)\,b_x\}\,\varphi_\lambda = 0.}

Bilden wir nun wieder die $n^{\text{te}}$ Polare von $\sum c_i\,H_i$, wobei die
Polare von $H_i$ mit $\overline{G}_i$ bezeichnet sei, dann ist, weil die Glieder $G_i$
in (1) Glieder dieser Polaren sein müssen [vergl. Nr. 25, (I)]:

\[
G_i = \overline{G}_i + (xy) \sum \varrho_i\,Q_i^{(2)},
\]

also

\[
\sum c_i\,G_i = \sum c_i\,\overline{G}_i + (xy) \sum \overline{\varrho}_i\,Q_i^{(3)},
\]

oder, weil $\sum c_i\,G_i = \mathbf{N}$, und $\sum c_i\,\overline{G}_i = \big[\sum c_i\,H_i\big]_{y^n}$,
mit Benutzung von Gleichung (1) und (3)

\[
\mathbf{N} = \Big[\sum \{(db)\,c_x + (bc)\,d_x + (cd)\,b_x\}\,\varphi_\lambda\Big]_{y^n} + (xy) \sum \overline{\varrho}_i\,Q_i^{(3)}.
\]


% ===== pdf idx 148, printed 134.0 =====
% !PRINTED 134
\weberpage{134}
Addiren wir hiezu:

\[
\mathbf{M} = (xy) \sum Q_i^{(1)},
\]

so erhalten wir:

\[
P_y^n = \Big[\sum \{(db)\,c_x + (bc)\,d_x + (cd)\,b_x\}\,\varphi_\lambda\Big]_{y^\mu} + (xy) \sum \lambda_i\,Q_i^{(4)}.
\]

Ersetzen wir endlich hierin $y$ wieder durch $a$, so geht $P_y^\mu$ in
$\sum F(A_i)$ über, wobei die späteren Posten $F(A_i)$ aus dem ersten durch
Vertauschung von $a$ mit $b$, resp. $c$, $d$, \dots\ etc. entstehen. Da sonach
alle diese Posten einander gleich sind und nur $\mu$ Symbole $a$, $b$, $c$ \dots\
$m$ vorhanden sind, so wird:

\[
\mu\,F(A_i) = \sum F(A_i) = \sum \{(db)\,c_x + (bc)\,d_x + (cd)\,b_x\}\,\varphi_\lambda + a_x \sum \overline{\lambda}_i\,Q_i^{(4)}.
\]

Weil nun $F(A_i)$ sowohl, als auch $\sum \{(db)\,c_x + (bc)\,d_x + (cd)\,b_x\}\,\varphi_\lambda$
identisch verschwinden, so verschwindet auch $\sum \overline{\lambda}_i\,Q_i^{(4)}$,
da aber wegen des abgesonderten Factors $a_x$ diese Relation
$\sum \overline{\lambda}_i\,Q_i^{(4)} = 0$ einen symbolischen Factor weniger besitzt
als die Relation $F(A_i) = 0$, so lässt sich nach Voraussetzung auch
$\sum \overline{\lambda}_i\,Q_i^{(4)}$ in die Form

\[
\sum \{(ab)\,c_x + (bc)\,a_x + (ca)\,b_x\}\,\overline{\varphi}_\lambda
\]

bringen und demnach hat man in der That:

\[
F(A_i) = \sum \{(ab)\,c_x + (bc)\,a_x + (ca)\,b_x\}\,\varphi_\lambda = 0.
\]

% [end-of-section ornament: short horizontal rule below the last display,
%  marking the close of "Erster Theil"; remainder of the page is blank]


% ===== pdf idx 149, printed 134.001 =====
% !PRINTED none
% Unnumbered part-opening leaf: large blank area at the head of the page
% (a short horizontal rule and two small printer's dots), followed by the
% part title, the part subtitle, the section heading, and the opening
% subsection text.

\begin{center}
Zweiter Theil.

\textbf{Die Formen zweiten, dritten und vierten Grades.}
\end{center}

\begin{center}
\textbf{§ 11. Die Form zweiten Grades: System einer und zweier}

\textbf{simultaner Formen.}
\end{center}

118. \textit{Ueberblick über die folgenden Untersuchungen.} Wir haben mit
den letzten Paragraphen die allgemeinen Untersuchungen über In- und
Covarianten binärer Formen (so weit sie sich nicht auf Systembildung
und associirte Formen beziehen) abgeschlossen und wenden uns nunmehr
den einzelnen Formen zu, indem wir mit der quadratischen Form
beginnen.

Die ersten Fragen, die wir uns hier stellen, sind die Fragen nach
dem vollen System einer quadratischen Form, oder zweier simultaner
Formen. Unter vollem System verstehen wir hier wie später die Gesammtheit
derjenigen Co- und Invarianten, durch welche sich alle übrigen
rational und ganz darstellen lassen. Dieses volle System aufzufinden,
wird immer eine unserer Hauptaufgaben sein. Eine allgemeine Methode,
das \gesperrt{kleinste volle System} aufzufinden, existirt nicht. Es muss immer
erst durch successive Begrenzung der Anzahl gezeigt werden, dass
wirklich keine der in das System aufgenommenen Formen überflüssig
ist, d.\,h. durch die übrigen sich rational und ganz darstellen lässt.
Diese Reduction der Formenanzahl stützt sich insbesondere auf einen
Begriff: den des \gesperrt{Reducenten}, dessen Definition im Laufe dieses
Paragraphen festgestellt wird.

Bei den quadratischen Formen lässt sich das volle System \textit{a priori}
angeben. Nachdem wir dies für zwei simultane Formen erledigt und
die Relationen zwischen diesen Formen aufgestellt haben, werden wir
in einigen Anwendungen davon Gebrauch machen. Hieran schliessen
wir die Betrachtungen über das volle System dreier quadratischer
Formen, und den mannigfachen Relationen die zwischen den Co- und
Invarianten desselben bestehen. Das System lässt sich dann leicht
allgemein für eine beliebige Anzahl quadratischer Formen erweitern.
Nach diesen Untersuchungen über allgemeine Formen zweiten Grades
%CONT>>


% ===== pdf idx 150, printed 136.0 =====
% !PRINTED 136
\weberpage{136}
%<<CONT
wenden wir uns speciellen zu, insbesondere solchen, die bei Auflösung
der Gleichungen vierten und fünften Grades eine wesentliche Rolle
spielen.

119. \textit{System einer quadratischen Form.} Es sei die quadratische
Form

\[
f = \overline{a}_0\,x_1^{\,2} + 2\,\overline{a}_1\,x_1 x_2 + \overline{a}_2\,x_2^{\,2}
\]

symbolisch dargestellt durch

\[
f = a_x^2 = b_x^2 = c_x^2 = \text{etc.}
\]

Die erste Ueberschiebung dieser Form über sich selbst ist

\[
(f,\,f) = (ab)\,a_x\,b_x\,;
\]

diese verschwindet aber identisch, da sie durch Vertauschung der völlig
gleichberechtigten Symbole $a$ und $b$ nur das Vorzeichen ändert, wie
wir bereits auch früher bemerkt haben. Die zweite Ueberschiebung
liefert die bekannte Discriminante

\[
D = (f,\,f)^2 = (ab)^2.
\]

Ihre unsymbolische Gestalt ist (vergl. auch Nr. 2)

\[
D = 2(\overline{a}_0\,\overline{a}_2 - \overline{a}_1^{\,2}).
\]

Die Ueberschiebung von $f$ über sich selbst führt somit nur auf eine
einzige Form, und da dieselbe Invariante ist, so können wir durch
diesen Process zu keinen weiteren Formen gelangen. Die Frage ist,
hat überhaupt die Form $f$ noch andere Co- und Invarianten und wenn
ja, lassen sich dieselben vielleicht als ganze rationale Functionen von
$f$ und $D$ darstellen.

Wir können zeigen, dass in der That beide Fragen bejaht werden
müssen, und bezeichnen deshalb $f$ und $D$ zusammen als volles System
der quadratischen Form.

Das allgemeinste symbolische Product nämlich, das eine Covariante
von $f$ darstellen kann, ist, weil höhere als erste Potenzen von $a_x$ und
$(ab)$ immer durch $f$ oder $D$ ersetzt werden können, dargestellt durch

\[
P = (ab)(ac) \dots (ik)\,a_x b_x c_x \dots i_x k_x.
\]

Greifen wir irgend einen Klammerfactor $(ab)$ aus diesem Producte
heraus, so müssen, soll $P$ wirkliche und nicht blos symbolische Bedeutung
haben, noch ausser ihm entweder die Factoren $a_x b_x$ oder die Factoren
$(ac)(bc)$ oder $(ad)(bc)$ \dots\ auftreten. Um diese Fälle in einen
zusammenzufassen, sagen wir

\[
Q = (ab)\,a_\xi\,b_\eta
\]

muss ein Theil des Productes $P$ sein, wenn $(ab)$ einer seiner Klammerfactoren
ist. Aber dieser Theil $Q$ lässt sich auf die bekannten Formen
%CONT>>


% ===== pdf idx 151, printed 137.0 =====
% !PRINTED 137
\weberpage{137}
%<<CONT
$f$ oder $D$ zurückführen. Vertauscht man nämlich $a$ und $b$, so kommt
\[
Q = -(ab)\,b_\xi a_\eta\,,
\]
also
\[
2Q = (ab)(a_\xi b_\eta - b_\xi a_\eta) = (ab)^2(\xi\eta)\,,
\]
oder
\[
Q = \frac{1}{2}\,D\cdot(\xi\eta).
\]

Das Product $P$ hat also die Discriminante $D$ als Factor, und die Ursache liegt in dem Auftreten des Klammerfactors $(ab)$. Durch Absonderung dieses Factors $D$ wird $P$ auf ein einfacheres Product $P_1$ zurückgeführt; es ist, wie wir uns in Zukunft ausdrücken werden, reducibel auf $D$ und $P_1$. In $P_1$ treten aber immer nur Klammerfactoren vom Typus $(ab)$ auf, an deren Stelle wieder die Discriminante $D$ eingeführt werden kann. Was übrig bleibt ist eine Constante mal einer Potenz der Form $f$. $P$ ist also reducibel auf $D^u\cdot f^v$; d. h. irgend ein symbolisches Product $P$ ist rational durch $D$ und $f$ darstellbar.

120. \textit{Auflösung der quadratischen Gleichung.} Eine Gleichung ist als gelöst zu betrachten, wenn sie in ihre Linearfactoren zerfällt. Diese Zerfällung kann man für
\[
f(x) = a_x^2 = b_x^2 = \ldots
\]
folgendermassen bewerkstelligen.

Quadrirt man die Identität
\[
(ab)(xy) = (a_x b_y - b_x a_y)\,,
\]
so kommt:
\[
D\cdot(xy)^2 = a_x^2 b_y^2 + b_x^2 a_y^2 - 2a_x a_y b_x b_y.
\]
Es ist aber $a_x a_y = b_x b_y = f_y$ und demnach wird:
\[
D\cdot(xy)^2 = 2f(x)\cdot f(y) - 2f_y^2\,,
\]
oder
\[
f(x)\cdot f(y) = f_y^2 + \frac{D}{2}\,(xy)^2.
\]
Hieraus folgt:
\[
f(x) = \frac{1}{f(y)}\left\{f_y + (xy)\sqrt{-\frac{D}{2}}\right\}\cdot\left\{f_y - (xy)\sqrt{-\frac{D}{2}}\right\}.
\]

Die rechte Seite dieser Gleichung besteht aus zwei in $x$ linearen Factoren, wobei $y$ nur als willkürlicher Parameter fungirt, so dass die Lösung als eine ganz allgemeine erscheint. Wollen wir die specielle für die gewöhnliche Gleichung zweiten Grades, so brauchen wir nur
\[
y_1 = 1\,, \quad y_2 = 0
\]
zu setzen und erhalten alsdann die bekannte Zerfällung:
\weq{I}{f(x) = \frac{1}{a_0}\left\{a_0 x_1 + a_1 x_2 - x_2\sqrt{-\frac{D}{2}}\right\}\left\{a_0 x_1 + a_1 x_2 + x_2\sqrt{-\frac{D}{2}}\right\}.}


% ===== pdf idx 152, printed 138.0 =====
% !PRINTED 138
\weberpage{138}
121. \textit{Darstellung der Discriminante in den Co\"efficienten der Linearfactoren von $f$.} Bezeichnen wir in (I) die beiden Klammern rechts mit $p_x$ und $q_x$, setzen wir also, indem wir $\dfrac{1}{a_0}$ in einen dieser Factoren hereinziehen,
\[
f = p\cdot q\,,
\]
dann ist
\[
D = -\frac{1}{2}\,(pq)^2.
\]

Denn $D$ ist als zweite Ueberschiebung von $f$ über sich selbst dargestellt durch (vergl. Nr. 38, I)
\[
\begin{aligned}
D &= (p_x q_x,\, p_x' q_x')^2 \\
&= \frac{1}{2}\,\{(pp')(qq') + (pq')(qp')\}.
\end{aligned}
\]

Da aber $p$ und $q$ lineare Formen sind, so sind ihre symbolischen Co\"efficienten gleich ihren wirklichen, und demnach ist:
\[
p = p'\,, \quad q = q'\,,
\]
also
\[
(pp') = 0\,, \quad (qq') = 0 \quad \text{und} \quad (pq') = -(qp') = +(pq).
\]
Daher wird:
\[
D = \frac{1}{2}\,(pq)(qp) = -\frac{1}{2}\,(pq)^2.
\]

Das gleiche Resultat erhält man auch durch folgende Ueberlegung. Die Form $D=(ab)^2$ geht aus $f=a_x^2$ hervor, wenn man $x$ durch $b$, d. h. $x_1$ durch $b_2$ und $x_2$ durch $-b_1$ ersetzt. Wenn also
\[
f = p_x q_x = a_x^2\,,
\]
dann ist:
\[
D = (pb)(qb) = (ab)^2.
\]

Das symbolische Product $(bp)(bq)$ geht aber aus der ersten Polare
\[
2f_y = 2b_x b_y = p_x q_y + q_x p_y
\]
hervor, wenn man $x$ durch $p$, $y$ durch $q$ ersetzt; dies liefert die Beziehung:
\[
2(bp)(bq) = (qp)(pq) = -(pq)^2\,,
\]
und somit:
\[
D = -\frac{1}{2}\,(pq)^2.
\]

122. \textit{Simultanes System zweier quadratischer Formen.} Das volle simultane System zweier quadratischer Formen
\[
f = a_x^2\,, \quad \varphi = \alpha_x^2
\]
lässt sich nach den vorausgegangenen Entwicklungen leicht a priori angeben. Es besteht einmal aus denjenigen In- und Covarianten, welche die einzelnen Formen selbst besitzen, also aus den vier Formen
%CONT>>


% ===== pdf idx 153, printed 139.0 =====
% !PRINTED 139
\weberpage{139}
%<<CONT
\[
\begin{aligned}
f &= a_x^2\,, & D_1 &= A_{ff} = (ab)^2\,, \\
\varphi &= \alpha_x^2\,, & D_2 &= A_{\varphi\varphi} = (\alpha\beta)^2\,,
\end{aligned}
\]
ferner aus jenen, welche durch Ueberschiebung von $f$ über $\varphi$ entstehen, also aus den zwei weitern Formen:
\[
\begin{aligned}
(f,\,\varphi) &= \vartheta_x^2 = \vartheta \\
(f,\,\varphi)^2 &= A_{f\varphi}.
\end{aligned}
\]

Da hiebei eine neue Covariante $\vartheta$ auftritt, so würden sich hieran jene Formen anschliessen, welche durch Ueberschiebung von $\vartheta$ über sich selbst und über $f$ und $\varphi$ sich ergeben. Indess entstehen aus diesen Ueberschiebungen keine neuen Formen; vielmehr lässt sich zeigen:

\glqq Das volle simultane System der beiden quadratischen Formen $f$ und $\varphi$ besteht aus den sechs Formen
\[
f,\ \varphi,\ A_{ff},\ A_{\varphi\varphi},\ \vartheta,\ A_{f\varphi};
\]
alle übrigen In- und Covarianten sind ganze und rationale Functionen derselben.\grqq{}

123. \textit{Gedankengang im Beweise dieses Satzes.} Das allgemeinste symbolische Product, das eine neue simultane Covariante von $f$ und $\varphi$ darstellen kann, wird die Symbole
\[
a,\ b,\ c,\ \ldots\,;\quad \alpha,\ \beta,\ \gamma,\ \ldots\,;\quad \vartheta,\ \vartheta_1,\ \vartheta_2,\ \ldots
\]
enthalten, da es aus den Formen $f$, $\varphi$, $\vartheta$ durch Ueberschiebung entstanden sein muss. Doch ist von vornherein klar, dass wir nur solche Producte zu untersuchen brauchen, in denen die Factoren erster Art, wie $a_x$, $\alpha_x$, $\vartheta_x$ nur in ersten Potenzen auftreten. Denn im andern Falle liessen sich ja sofort die bekannten Covarianten $f$, $\varphi$, $\vartheta$ des Systemes als Factoren absondern.

Was die Klammerfactoren betrifft, so können sie, da die Grundformen $f$, $\varphi$, $\vartheta$ selbst nur vom zweiten Grade sind, nur in ersten oder zweiten Potenzen auftreten. Die überhaupt denkbaren Klammerfactoren sind aber folgende sechs:
\weq{A}{(ab),\ (\alpha\beta),\ (a\alpha)\,;\quad (\vartheta\vartheta_1),\ (a\vartheta),\ (\alpha\vartheta).}
Wir zeigen nun zuerst:

\glqq Ein symbolisches Product, das einen solchen Factor quadratisch enthält, ist reducibel, d. h. es enthält Factoren, die sich aus den ursprünglichen Invarianten $A_{ff}$, $A_{\varphi\varphi}$, $A_{f\varphi}$ zusammensetzen.\grqq{}

In zweiter Linie beweisen wir:

\glqq Ein symbolisches Product, in welchem ein solcher Factor linear auftritt, ist reducibel, d. h. es enthält Factoren, welche rationale ganze Functionen der ursprünglichen sechs Co- und Invarianten sind.\grqq{}


% ===== pdf idx 154, printed 140.0 =====
% !PRINTED 140
\weberpage{140}
124. \textit{Begriff der Reducibilität und des Reducenten.} Ich will hiebei zunächst die Begriffe der Reducibilität und des Reducenten, wie wir sie für jetzt, wie für später n\"othig haben, fest umgrenzen. Ein symbolisches Product ist in drei Fällen als reducibel anzusehen: Erstens, wenn es identisch verschwindet; zweitens, wenn es bereits bekannte Formen, die in dem System schon aufgenommen sind, insbesondere Invarianten zu Factoren hat, oder gar durch solche Formen vollständig dargestellt werden kann; drittens, wenn es sich in andere symbolische Producte überführen lässt, die bereits als reducibel erkannt wurden, was sich oft dadurch schon ausweist, dass man es auf ein Product mit weniger bereits bekannten Symbolen zu reduciren vermag. (Vergl. auch Nr. 144.)

Von den in einem solchen reduciblen Producte auftretenden Klammerfactoren sind alsdann insbesondere jene von Bedeutung, kraft deren die Reducibilität eintritt. Sie erhalten den Namen \glqq Reducent\grqq{}, sobald jedes symbolische Product, welches sie besitzt, reducibel ist. Hiezu ist nur n\"othig, dass das ursprüngliche Product $P$, welches aus ihm und seinen Symbolen allein construirt ist, und alle aus $P$ durch Faltung entstehenden Producte $P_i$ reducible Formen sind. Denn ein beliebiges Product $\Pi$, das den betreffenden Reducenten besitzt, kann stets durch Ueberschiebungen dargestellt werden, welche diese reduciblen Formen $P$ und $P_i$ mit geeigneten andern Formen bilden. (Vergl. Nr. 46.) Und diese Ueberschiebungen sind reducibel, wenn $P$ und $P_i$ es waren. (Vergleiche auch Nr. 144.)

125. \textit{Beweis, dass alle Producte, welche einen der sechs Factoren (A) quadratisch enthalten, reducibel sind.} Da die Grössen $(ab)^2$, $(\alpha\beta)^2$, $(a\alpha)^2$ nicht anderes als die Invarianten $A_{ff}$, $A_{\varphi\varphi}$, $A_{f\varphi}$ des Systemes sind, so können wir von vornherein von Producten mit diesen Klammerfactoren absehen.

Die Grösse $(\vartheta\vartheta_1)^2$ ist die Discriminante $A_{\vartheta\vartheta}$ der Form $\vartheta$, wir können zeigen, dass sie durch die drei Invarianten des Systemes ausdrückbar ist.

Zu dem Zwecke wollen wir diese Discriminante noch durch andere symbolische Producte darstellen, aus deren Vergleichung alsdann die Reducibilität ersichtlich ist.

Zunächst ersetzen wir in $(\vartheta\vartheta_1)^2$ das Symbol $\vartheta_1$ durch die ursprünglichen Symbole $a$, $\alpha$, gemäss der Definition
\[
\vartheta_{1x}^2 = (a\alpha)\,a_x \alpha_x\,,
\]
und erhalten so:
\weq{1}{(\vartheta\vartheta_1)^2 = (a\alpha)(\vartheta\alpha)(\vartheta a).}


% ===== pdf idx 155, printed 141.0 =====
% !PRINTED 141
\weberpage{141}

Das symbolische Product rechts erhalten wir aber auch, wenn wir
ein Glied der ersten Polare von $\vartheta_x^2 = (b\beta)\,b_x\beta_x$ in die Reihe (VIII)
\S.\,7 entwickeln. Es ist nämlich:

\weq{2}{(b\beta)\,b_x\beta_y = \{(b\beta)\,b_x\beta_x\}_y + \frac{1}{2}\,(b\beta)^2(xy)\,,}

oder, indem wir nun $x$ durch $a$, $y$ durch $\alpha$ ersetzen und die Symbole
$\vartheta_x^2$, $A_{f\varphi}$ für $(b\beta)\,b_x\beta_x$, resp. $(b\beta)^2$ einführen und mit $(a\alpha)$
multipliciren:

\weq{3}{(b\beta)(ba)(\beta\alpha)(a\alpha) = (\vartheta a)(\vartheta\alpha)(a\alpha) + \frac{1}{2}\,A_{f\varphi}\cdot(a\alpha)^2\,.}

Hierin tritt bereits das symbolische Product der Gleichung (1)
auf, und da $(a\alpha)^2 = A_{f\varphi}$, so handelt es sich nur mehr darum, auch
die linke Seite dieser Gleichung (3) auf bekannte Formen zurückzuführen. Es ist aber, wenn wir darin $a$ mit $b$ vertauschen, und die
halbe Summe der so entstehenden Producte nehmen:

\[
(b\beta)(ba)(\beta\alpha)(a\alpha) = \frac{1}{2}\,(ab)(\alpha\beta)\{(a\alpha)(b\beta) - (b\alpha)(a\beta)\}
\]
oder

\weq{4}{= \frac{1}{2}\,(ab)^2(\alpha\beta)^2 = \frac{1}{2}\,A_{ff}\cdot A_{\varphi\varphi}\,,}

wie der Identitätssatz unmittelbar liefert. Demnach erhalten wir aus
Gleichung (3) und (4)

\weq{5}{\frac{1}{2}\,\{A_{ff}\cdot A_{\varphi\varphi} - A_{f\varphi}^2\} = (a\alpha)(\vartheta\alpha)(\vartheta a)\,,}

oder durch Vergleichung mit (1)

\weq{6}{(\vartheta\vartheta_1)^2 = \frac{1}{2}\,\{A_{ff}A_{\varphi\varphi} - A_{f\varphi}^2\}\,,}

d. h.: \glqq Jedes symbolische Product mit dem Factor $(\vartheta\vartheta_1)^2$ ist
reducibel.\grqq{}
Dies gilt aber auch für symbolische Producte mit den Factoren
$(a\vartheta)^2$ und $(\alpha\vartheta)^2$. Denn es ist:

\weq{7}{
\left.
\begin{aligned}
0 &= A_{f\vartheta} = (a\vartheta)^2 = (b\alpha)(ab)(a\alpha) \\
0 &= A_{\varphi\vartheta} = (\alpha\vartheta)^2 = (a\beta)(\alpha a)(\alpha\beta)
\end{aligned}
\right\}\,,
}

wie man unmittelbar erkennt, da die symbolischen Producte rechts
durch Vertauschung gleichberechtigter Symbole nur ihr Vorzeichen
ändern. Wir haben somit den Satz bewiesen:

\begin{quote}
\glqq Alle symbolischen Producte, welche einen der sechs möglichen
Klammerfactoren \gesperrt{quadratisch} enthalten, sind reducibel.\grqq{}
\end{quote}

Es erübrigt also nur noch, dieselbe Eigenschaft für symbolische
Producte nachzuweisen, welche diese Factoren linear besitzen.

126. \textit{Beweis, dass alle Producte, welche einen der sechs Factoren (A)
in Nr. 123 linear enthalten, reducibel sind.}


% ===== pdf idx 156, printed 142.0 =====
% !PRINTED 142
\weberpage{142}

Die Factoren $(ab)$ und $(\alpha\beta)$ haben wir bereits in Nr. 119 als Reducenten
erkannt; genau dieselbe Eigenschaft besitzt $(\vartheta\vartheta_1)$ aus gleichen Gründen.
Es bleiben also nur noch Producte mit den linearen Klammerfactoren
$(a\alpha)$, $(a\vartheta)$, $(\alpha\vartheta)$ zu untersuchen. Ein Product, das den Factor $(a\vartheta)$
oder $(\alpha\vartheta)$ besitzt, muss auch noch die Factoren $a_\xi\vartheta_\eta$ resp.
$\alpha_\xi\vartheta_\eta$ enthalten, wobei die Grössen $\xi$, $\eta$ für irgend welche andere
Symbole oder Variable stehen mögen. Wir brauchen also nur zu zeigen, dass

\[
(a\vartheta)\,a_\xi\,\vartheta_\eta \quad \text{und} \quad (\alpha\vartheta)\,\alpha_\xi\,\vartheta_\eta
\]

reducibel sind. Dies erkennen wir aber auf folgende Weise. Da $(a\vartheta)$
und $(\alpha\vartheta)$ die Klammerfactoren der Functionaldeterminanten $(f,\ \vartheta)$
und $(\varphi,\ \vartheta)$ sind, so sind sie sicher Reducenten, wenn diese Formen
und die aus ihnen durch Faltung entstehenden reducibel sind. Es
sind aber $(f,\ \vartheta)$ und $(\varphi,\ \vartheta)$ Functionaldeterminanten von Functionaldeterminanten, und als solche nach den früher gegebenen allgemeinen
Theorien stets durch einfachere Formen darstellbar. In der That
liefert Nr. 52, (I) und (II) direct:

\[
(f,\ \vartheta) = \frac{1}{2}\,\{A_{f\varphi}\cdot f - A_{ff}\cdot\varphi\}
\]

\[
(\vartheta,\ \varphi) = \frac{1}{2}\,\{A_{f\varphi}\cdot\varphi - A_{\varphi\varphi}\cdot f\}\,.
\]

Die aus $(f,\ \vartheta)$ und $(\vartheta,\ \varphi)$ durch Faltung entstehenden Formen
$A_{f\vartheta}$ und $A_{\varphi\vartheta}$ sind aber gleichfalls reducibel, da sie verschwinden;
demnach sind $(a\vartheta)$ und $(\alpha\vartheta)$ in der That Reducenten.

Man kann die Reducibilität der Formen $(a\vartheta)\,a_\xi\vartheta_\eta$ und $(\alpha\vartheta)\,\alpha_\xi\vartheta_\eta$
auch in anderer Weise einfach zeigen. Aus den Formeln Nr. 51, (II)
erhält man nämlich direct:

\[
(a\vartheta)\,\vartheta_\eta\,a_\xi = (f,\ \vartheta)_\eta = \frac{1}{2}\,\{A_{f\varphi}\cdot f_\xi\cdot f_\eta - A_{ff}\ \cdot\varphi_\xi\cdot\varphi_\eta\}
\]

\[
-(\alpha\vartheta)\,\vartheta_\eta\,\alpha_\xi = (\vartheta,\ \varphi)_\eta = \frac{1}{2}\,\{A_{f\varphi}\cdot\varphi_\xi\cdot\varphi_\eta - A_{\varphi\varphi}\cdot f_\xi\cdot f_\eta\}\,.
\]

Die symbolischen Producte rechts enthalten neben den bekannten
drei Invarianten des Systemes nur mehr ein Symbol $a$, respective
$\alpha$, also \gesperrt{weniger} Symbole wie die Ausdrücke links, wodurch deren
Reducibilität nachgewiesen ist.

Es erübrigt noch, die Reducibilität von Producten mit dem Klammerfactor $(a\alpha)$, oder, was dasselbe ist, mit dem Theilproduct $(a\alpha)\,a_\xi\alpha_\eta$
nachzuweisen. Sie ergiebt sich direct aus Nr. 125, (2), wonach:

\[
(a\alpha)\,a_\xi\,\alpha_\eta = \vartheta_\xi\,\vartheta_\eta + \frac{1}{2}\,A_{f\varphi}\,(\xi\eta)\,.
\]


% ===== pdf idx 157, printed 143.0 =====
% !PRINTED 143
\weberpage{143}

Die rechte Seite enthält nur reducible Glieder, das zweite wegen
des Factors $A_{f\varphi}$, das erste, da wir für jeden Klammerfactor mit dem
Symbol $\vartheta$ bereits die Eigenschaft nachgewiesen haben, ein Product,
dem er angehört, reducibel zu machen.

Wir sind also damit zu dem Resultate gelangt, dass jedes symbolische Product, welches irgend einen der sechs möglichen Klammerfactoren linear enthält, auf die sechs bekannten Formen reducirt
werden kann.

\begin{quote}
\glqq Demnach sind alle sechs Klammerfactoren

\[
(ab),\ (\alpha\beta),\ (a\alpha);\quad (a\vartheta),\ (\alpha\vartheta),\ (\vartheta\vartheta)
\]

Reducenten, und es lassen sich alle Co- und Invarianten der
Formen $f$ und $\varphi$ ganz und rational durch die Formen

\[
f,\ \varphi,\ \vartheta,\ A_{ff},\ A_{\varphi\varphi},\ A_{f\varphi}
\]

darstellen.\grqq{}
\end{quote}

127. \textit{Zusammenfassung der gewonnenen Relationen.} Ich will nur
noch tabellarisch auf die wichtigeren Relationen, die bei dieser Untersuchung sich ergeben haben, aufmerksam machen, indem ich sie im
Folgenden zusammenstelle. Ist $f = a_x^2$, $\varphi = \alpha_x^2$, $\vartheta = (a\alpha)\,a_x\,\alpha_x$, so
hat man:

\[
\begin{array}{rll}
(1) & 2A_{\vartheta\vartheta} = A_{ff}A_{\varphi\varphi} - A_{f\varphi}^2 & \text{(Nr. 125)} \\[4pt]
(2) & A_{f\vartheta} = 0,\quad A_{\varphi\vartheta} = 0 & \text{(Nr. 125)} \\[4pt]
(3) & 2(\vartheta,\ f) = A_{ff}\cdot\varphi - A_{f\varphi}\cdot f & \text{(Nr. 126)} \\[4pt]
    & 2(\varphi,\ \vartheta) = A_{\varphi\varphi}\cdot f - A_{f\varphi}\cdot\varphi. &
\end{array}
\]

Hiezu tritt:

\[
(4)\quad -2\vartheta^2 = A_{ff}\varphi^2 - 2A_{f\varphi}f\varphi + A_{\varphi\varphi}f^2\,,
\]

eine Relation (4), welche zwar hier nicht neuerdings abgeleitet wurde,
die sich aber unmittelbar aus der allgemeineren Nr. 48, (I) ergiebt.

128. \textit{Erste Anwendung.} Ist $\vartheta$ identisch Null, so sind $f$ und
$\varphi$ einander proportional. Wir zeigen zuerst, dass $\vartheta = 0$ die hinreichende Bedingung ist, damit die Gleichung besteht

\[
f = \varrho \cdot \varphi\,,
\]

wo $\varrho$ ein Proportionalitätsfactor. Es ist nämlich

\[
2\vartheta_x\vartheta_y = (a\alpha)\{a_x\alpha_y + \alpha_x a_y\}\,,
\]

und weil nach dem Productsatze:

\[
(xy) = \frac{1}{(a\alpha)}\{a_x\alpha_y - \alpha_x a_y\}\,,
\]

so folgt durch Multiplication beider Gleichungen

\[
2(xy)\vartheta_x\vartheta_y = a_x^2\alpha_x^2 - a_y^2\alpha_x^2\,.
\]


% ===== pdf idx 158, printed 144.0 =====
% !PRINTED 144
\weberpage{144}

Wenn nun $\vartheta = 0$, dann ist auch $\vartheta_y = 0$, also wird in diesem Falle:

\[
a_x^2\,\alpha_y^2 = a_y^2\,\alpha_x^2\,,
\]

oder

\[
a_x^2 = \frac{a_y^2}{\alpha_y^2}\cdot\alpha_x^2 = \varrho\cdot\alpha_x^2\,.
\]

Die Bedingung $\vartheta = 0$ ist aber nicht nur hinreichend, sondern
auch nothwendig. Denn wenn

\[
\alpha_x^2 = \frac{1}{\varrho}\,a_x^2 = \frac{1}{\varrho}\,b_x^2\,,
\]

dann ist

\[
\vartheta = (a_x^2,\ \alpha_x^2) = \frac{1}{\varrho}\,(a_x^2,\ b_x^2) = \frac{1}{\varrho}\,(ab)\,a_x\,b_x = 0\,.
\]

129. \textit{Zweite Anwendung.} Ist $A_{\vartheta\vartheta}$ identisch null, so besitzen $f$ und
$\varphi$ einen gemeinsamen linearen Factor. Diese Bedingung ist nothwendig; denn wenn $f$ und $\varphi$ einen gemeinsamen linearen Factor besitzen, also etwa

\[
f = p_x\,q_x
\]

\[
\varphi = p_x\,r_x
\]

ist, dann wird $\vartheta$ ein reines Quadrat. Denn man hat nach Nr. 38, (I)

\[
\vartheta = (p_x\,q_x,\ p_x\,r_x) = \frac{1}{4}\,\{(pp)\,q_x\,r_x + (pr)\,q_x\,p_x + (qp)\,p_x\,r_x + (qr)\,p_x^2\}\,.
\]

Da aber $(pp) = 0$ und $(pr)\,q_x + (qp)\,r_x = (qr)\,p_x$, so wird

\[
\vartheta = \frac{1}{4}\,\{(qr)\,p_x^2 + (qr)\,p_x^2\}\,,
\]

oder

\[
\vartheta = \frac{1}{2}\,(qr)\cdot p_x^2\,.
\]

Die Functionaldeterminante $\vartheta$ ist also das Quadrat des gemeinsamen
linearen Factors $p_x$ und demnach $A_{\vartheta\vartheta}$ als Discriminante identisch null.

Die Bedingung $A_{\vartheta\vartheta} = 0$ ist aber auch hinreichend. Denn wenn

\[
A_{\vartheta\vartheta} = 0\,,
\]

dann ist

\[
\vartheta = m_x^2\,,
\]

wo $m_x$ eine lineare Form ist. Die zweite Ueberschiebung von $f = a_x^2$
über $\vartheta = m_x^2$ ist sodann

\[
A_{f\vartheta} = (am)^2 = a_m^2\,.
\]

Wir wissen aber, dass $A_{f\vartheta}$ wie $A_{\varphi\vartheta}$ identisch verschwinden; es
ist also sowohl

\[
a_m^2 = 0\,,
\]

als auch

\[
\alpha_m^2 = 0\,,
\]
%CONT>>


% ===== pdf idx 159, printed 145.0 =====
% !PRINTED 145
\weberpage{145}
%<<CONT
d. h. die Formen $a_x^2$ und $\alpha_x^2$ werden identisch null für $x=m$, oder was das
Nämliche ist, sie besitzen einen gemeinschaftlichen linearen Factor $m_x$.

Da also $A_{\vartheta\vartheta}=0$ die nothwendige und hinreichende Bedingung ist,
dass $f$ und $\varphi$ einen gemeinsamen Factor besitzen, so muss $A_{\vartheta\vartheta}$
der Resultante $R_{f,\varphi}$ beider Formen proportional sein. Die
Proportionalitätsconstante kann nur ein Zahlenfactor sein, da sowohl $A_{\vartheta\vartheta}$
als auch $R_{f,\varphi}$ vom vierten Grade in den Coefficienten beider Formen ist. Man
findet durch ein beliebiges Beispiel leicht:

\[
R_{f,\varphi} = -2A_{\vartheta\vartheta} = A_{f\varphi}^2 - A_{ff}A_{\varphi\varphi}.
\]

\textit{Anmerkung.} Die Identität:

\[
2\vartheta = (qr)\cdot p_x^2
\]
lehrt zugleich, wie man den gemeinschaftlichen Factor ermitteln kann. Denn vergleicht
man in der ersten Polare

\[
2\vartheta_x\vartheta_y = (qr)p_xp_y
\]
die Coefficienten von $y_1$ und $y_2$, so erhält man:

\[
\begin{gathered}
2\vartheta_x\vartheta_1 = (qr)p_xp_1\\
2\vartheta_x\vartheta_2 = (qr)p_xp_2,
\end{gathered}
\]
und somit durch Division:

\weq{1}{\vartheta_x\vartheta_1 : \vartheta_x\vartheta_2 = p_1:p_2}

für jeden Werth von $x$. Die beiden ersten partiellen Differentialquotienten von
$\vartheta$ sind sonach den Coefficienten des gemeinschaftlichen linearen Factors von
$f$ und $\varphi$ proportional; dies gilt insbesondere beispielsweise für $x_1=1,
x_2=0$, für welchen Fall Gleichung (1) in

\weq{2}{\bar\vartheta_0 : \bar\vartheta_1 = p_1 : p_2}

übergeht. Es ist aber nach Nr. 47, (II)

\[
\vartheta = \begin{vmatrix} \bar a_0 & \bar a_1 & \bar a_2 \\ \bar\alpha_0 & \alpha_1 & \alpha_2 \\ x_1^2 & -x_1x_2 & x_2^2 \end{vmatrix} = \bar\vartheta_0 x_1^2 + 2\bar\vartheta_1 x_1x_2 + \bar\vartheta_2 x_2^2,
\]
wobei

\[
\bar\vartheta_0 = (\bar a_1\bar\alpha_2), \quad \bar\vartheta_1 = \frac{1}{2}(\bar a_0\bar\alpha_2), \quad \bar\vartheta_2=(\bar a_0\bar\alpha_1).
\]

Durch Substitution dieser Werthe in (2) erhalten wir:

\weq{3}{\frac{p_1}{p_2} = \frac{2(\bar a_1\bar\alpha_2)}{(\bar a_0\bar\alpha_2)}}

130. \textit{Dritte Anwendung. Kanonische Form für zwei simultane
quadratische Formen.} Wir stellen uns die Aufgabe, zwei solche lineare
%CONT>>


% ===== pdf idx 160, printed 146.0 =====
% !PRINTED 146
\weberpage{146}
%<<CONT
Formen $r_x$ und $s_x$ zu ermitteln, dass durch dieselbe Transformation $f$ und
$\varphi$ gleichzeitig die Gestalt erhalten:

\weq{I}{\left.\begin{aligned} f &= \lambda_1 r_x^2 + \lambda_2 s_x^2\\ \varphi &= \mu_1 r_x^2 + \mu_2 s_x^2\end{aligned}\right\}}

Die erste Untersuchung zur Lösung solcher Aufgaben ist immer: Wie gestaltet sich
das simultane Formensystem, wenn wir $f$ und $\varphi$ bereits in dieser kanonischen
Form zu Grunde legen? Es ergiebt sich sofort, dass $\vartheta$ das Product der beiden
linearen Formen $r_x$ und $s_x$ ist, bis auf einen constanten Factor. Denn man erhält:

\[
\begin{gathered}
\vartheta = (f,\varphi) = (\lambda_1 r_x^2 + \lambda_2 s_x^3,\ \mu_1 r_x^2 + \mu_2 s_x^2) % [sic]
\\
= \lambda_1\mu_1(r_x^2,r_x^2) + \lambda_1\mu_2(r_x^2,s_x^2) + \lambda_2\mu_1(s_x^2,r_x^2) + \lambda_2\mu_2(s_x^2,s_x^2)
\\
= \lambda_1\mu_2(rs)r_xs_x - \lambda_2\mu_1(rs)r_xs_x
\\
= (\lambda\mu)(rs)r_xs_x.
\end{gathered}
\]

Wir lernen daraus: Die Factoren $r_x$ und $s_x$, durch welche sich $f$ und $\varphi$
gleichzeitig in die kanonische Form (I) bringen lassen, sind den linearen Factoren
von $\vartheta$ proportional. Damit also überhaupt diese Darstellung möglich ist,
muss $\vartheta$ wirklich zwei verschiedene lineare Factoren besitzen, d. h. $f$ und
$\varphi$ dürfen keinen Factor gemeinsam haben.

Die Einführung dieser beiden linearen Factoren von $\vartheta$ ist aber auch
hinreichend, um $f$ und $\varphi$ gleichzeitig als Summe zweier Quadrate darzustellen.
Denn sei

\[
\vartheta = r_x \cdot s_x,
\]
dann ist nach Nr. 126 (2):

\weq{1}{\left.\begin{aligned} A_{f\vartheta} &= (as)(ar) = 0\\ A_{\varphi\vartheta} &= (\alpha s)(\alpha r) = 0\end{aligned}\right\}}

Nun hat man gemäss dem Identitätssatze:

\[
\begin{gathered}
a_x(rs) = r_x(as) - s_x(ar)\\
\alpha_x(rs) = r_x(\alpha s) - s_x(\alpha r).
\end{gathered}
\]

Also, indem man beide Identitäten quadrirt, und berücksichtigt, dass wegen der
Relationen (1) die mittleren Glieder rechts verschwinden:

\weq{II}{\left.\begin{aligned} f\cdot(rs)^2 &= r_x^2(as)^2 + s_x^2(ar)^2\\ \varphi\cdot(rs)^2 &= r_x^2(\alpha s)^2 + s_x^2(\alpha r)^2\end{aligned}\right\}}

wodurch in der That $f$ und $\varphi$ als Summen zweier Quadrate dargestellt sind.

Um die Transformation von $f$ und $\varphi$ wirklich auszuführen, wird man sonach
entweder die Form $\vartheta$ in ihre Linearfactoren spalten oder
%CONT>>


% ===== pdf idx 161, printed 147.0 =====
% !PRINTED 147
\weberpage{147}
%<<CONT
einfacher in folgender Weise verfahren. Man multiplicirt die Gleichungen (II) mit
1 resp. $\lambda$ und addirt dieselben; dann erhält man:

\[
f + \lambda\varphi = r_x^2\left\{\frac{(as)^2+\lambda(\alpha s)^2}{(rs)^2}\right\} + s_x^2\left\{\frac{(ar)^2+\lambda(\alpha r)^2}{(rs)^2}\right\}.
\]

Nun lässt sich $\lambda$ so bestimmen, dass der Coefficient von $r_x^2$ oder von
$s_x^2$ verschwindet; in diesem Falle geht dann $f+\lambda\varphi$ gerade in das
Quadrat der linearen Form $s_x$ resp. $r_x$ (bis auf einen constanten Factor) über.
Damit aber $f+\lambda\varphi$ ein reines Quadrat sei, muss die Discriminante:

\[
(f+\lambda\varphi,\ f+\lambda\varphi)^2 = (f,f)^2 + 2\lambda(f,\varphi)^2 + \lambda^2(\varphi,\varphi)^2
\]
identisch verschwinden. Dies liefert eine quadratische Gleichung für $\lambda$,
nämlich:

\[
A_{ff} + 2\lambda A_{f\varphi} + \lambda^2 A_{\varphi\varphi} = 0.
\]

Sind dann $\lambda_1$ und $\lambda_2$ die Wurzeln dieser Gleichung, so ist:

\[
\begin{gathered}
f + \lambda_1\varphi = m_x^2\\
f + \lambda_2\varphi = n_x^2,
\end{gathered}
\]
und somit:

\[
\begin{gathered}
f = \frac{\lambda_2}{\lambda_2-\lambda_1}m_x^2 - \frac{\lambda_1}{\lambda_2-\lambda_1}n_x^2\\
\varphi = \frac{1}{\lambda_1-\lambda_2}m_x^2 - \frac{1}{\lambda_1-\lambda_2}n_x^2.
\end{gathered}
\]

\begin{center}
\textbf{\S\ 12. Die Form zweiten Grades: System dreier und mehr simultaner Formen.}
\end{center}

131. \textit{Das simultane System dreier quadratischer Formen.} Das volle System
dreier quadratischer Formen $f,\varphi,\psi$ besteht aus folgenden Bildungen.

1) Die Originalformen selbst:

\[
\begin{gathered}
f = a_x^2 = b_x^2 = \dots\\
\varphi = \alpha_x^2 = \beta_x^2 = \dots\\
\psi = r_x^2 = s_x^2 = \dots;
\end{gathered}
\]

2) ihre sechs zweiten Ueberschiebungen:

\[
\begin{gathered}
(f,f)^2 = A_{ff}, \quad (\varphi,\varphi)^2 = A_{\varphi\varphi}\\
(f,\varphi)^2 = A_{f\varphi}, \quad (\varphi,\psi)^2 = A_{\varphi\psi}\\
(f,\psi)^2 = A_{f\psi}, \quad (\psi,\psi)^2 = A_{\psi\psi};
\end{gathered}
\]

3) ihre drei Functionaldeterminanten (schiefe Covarianten):

\[
\begin{gathered}
(\varphi,\psi) = \vartheta_1 = \vartheta_{1x}^2\\
(\psi,f) = \vartheta_2 = \vartheta_{2x}^2\\
(f,\varphi) = \vartheta_3 = \vartheta_{3x}^2;
\end{gathered}
\]


% ===== pdf idx 162, printed 148.0 =====
% !PRINTED 148
\weberpage{148}
4) ihre Combinante (die einzige schiefe Invariante):

\[
(a\alpha)(\alpha r)(ar) = \begin{vmatrix} \bar a_0 & \bar a_1 & \bar a_2 \\ \bar\alpha_0 & \bar\alpha_1 & \bar\alpha_2 \\ \bar r_0 & \bar r_1 & \bar r_2 \end{vmatrix} = R_{123},
\]
welch letztere wir bereits in § 1 kennen gelernt haben. Ausser diesen dreizehn
Formen existirt keine, die sich nicht als ganze rationale Function derselben
darstellen liesse.

132. \textit{Alle Producte mit quadratischen Klammerfactoren sind auf diese
dreizehn Formen reducibel.} Irgend ein symbolisches Product nämlich, das eine neue
Covariante der drei Formen repräsentiren soll, kann nur folgende Klammerfactoren

\weq{A}{\begin{gathered}
1)\quad (ab)\quad (\alpha\beta)\quad (rs)\quad (a\alpha)\quad (ar)\quad (\alpha r)\\
(a\vartheta_3)\quad (a\vartheta_2)\quad (\alpha\vartheta_1)\quad (\alpha\vartheta_3)\quad (r\vartheta_2)\quad (r\vartheta_1)\\
(\vartheta_1\vartheta_1)\quad (\vartheta_2\vartheta_2)\quad (\vartheta_3\vartheta_3);\\
2)\quad (a\vartheta_1)\quad (\alpha\vartheta_2)\quad (r\vartheta_3)\quad (\vartheta_1\vartheta_2)\quad (\vartheta_1\vartheta_3)\quad (\vartheta_2\vartheta_3)
\end{gathered}}

besitzen und zwar im höchsten Falle in der zweiten Potenz. Was die fünfzehn
Factoren der ersten Gruppe betrifft, so sahen wir bereits bei Aufstellung des
simultanen Systemes zweier quadratischer Formen, dass jedes Product, in welchem
sie quadratisch oder linear auftreten, auf die Formen des Systemes reducibel ist.
Unsere Aufgabe ist also nur noch nachzuweisen, dass auch die Klammerfactoren der
zweiten Gruppe Reducenten sind. Wir betrachten wiederum zuerst Producte, deren
Klammerfactoren den Exponenten 2 haben. Die Factoren

\[
(a\vartheta_1)^2,\ (\alpha\vartheta_2)^2,\ (r\vartheta_3)^2,\ (\vartheta_1\vartheta_2)^2,\ (\vartheta_1\vartheta_3)^2,\ (\vartheta_2\vartheta_3)^2
\]
gehören zunächst folgenden Ueberschiebungen an:

\weq{1}{(f,(\varphi,\psi))^2,\qquad (\varphi,(\psi,f))^2,\qquad (\psi,(f,\varphi))^2}

\weq{2}{((\varphi,\psi),(\psi,f))^2,\quad ((\varphi,\psi),(f,\varphi))^2,\quad ((\psi,f),(f,\varphi))^2.}

Es ist aber:

\weq{I}{\begin{aligned} (f,(\varphi,\psi))^2 &= (a_x^2,(\alpha r)\alpha_xr_x)^2 = (a\alpha)(\alpha r)(ar) = R_{123}\\ &= (\varphi,(\psi,f))^2 = (\psi,(f,\varphi))^2,\end{aligned}}

und folglich führen die Ueberschiebungen der Gruppe (1) auf keine neue Form, d. h.
Producte mit den Factoren $(a\vartheta_1)^2$, $(\alpha\vartheta_2)^2$,
$(r\vartheta_3)^2$ sind reducibel.

Was die Ueberschiebungen der Gruppe (2) betrifft, so können wir deren
Reducibilität auf folgende Art beweisen. Ich will zuerst den Weg der directen
Rechnung einschlagen. Es ist, wenn $f=a_x^2$, $\varphi=\alpha_x^2$, $\psi=r_x^2$,
$\chi=\mu_x^2$ irgend vier quadratische Formen darstellen:
%CONT>>


% ===== pdf idx 163, printed 149.0 =====
% !PRINTED 149
\weberpage{149}

\[
\begin{aligned}
U &= ((f,\varphi),(\psi,\chi))^2 = (a\alpha)(r\mu)(a_x\alpha_x,\ r_x\mu_x)^2\\
&= \frac{1}{2}(a\alpha)(r\mu)\{(ar)(\alpha\mu)+(a\mu)(\alpha r)\}\\
&= \frac{1}{2}(G_1+G_2).
\end{aligned}
\]
Die Differenz der beiden Glieder rechts ist:
\[
G_1-G_2=(a\alpha)(r\mu)\{(ar)(\alpha\mu)-(a\mu)(\alpha r)\},
\]
oder weil:
\[
(ar)(\alpha\mu)-(a\mu)(\alpha r)=(a\alpha)(r\mu)
\]
(Identit\"atssatz), so erh\"alt man:
\[
G_1-G_2=(a\alpha)^2(r\mu)^2=A_{f\varphi}A_{\psi\chi}.
\]
Addiren wir die beiden Gleichungen:
\[
\begin{aligned}
U &= \frac{1}{2}(G_1+G_2)\\
\frac{1}{2}A_{f\varphi}A_{\psi\chi} &= \frac{1}{2}(G_1-G_2),
\end{aligned}
\]
so erh\"alt man:
\[
U+\frac{1}{2}A_{f\varphi}A_{\psi\chi}=(a\alpha)(r\mu)(ar)(\alpha\mu),
\]
oder, wenn wir auf das symbolische Product rechts den Productsatz anwenden:
\[
U+\frac{1}{2}A_{f\varphi}A_{\psi\chi}=\frac{1}{2}\{A_{f\varphi}A_{\psi\chi}+A_{f\psi}A_{\varphi\chi}-A_{f\chi}A_{\varphi\psi}\},
\]
also
\weq{II}{U=((f,\varphi),(\psi,\chi))^2=\frac{1}{2}\{A_{f\psi}A_{\varphi\chi}-A_{f\chi}A_{\varphi\psi}\}.}
F\"ur $\chi=f$ erh\"alt man hieraus
\[
((f,\varphi),(\psi,f))^2=\frac{1}{2}(A_{f\psi}A_{f\varphi}-A_{ff}A_{\varphi\psi})=A_{\vartheta_3\vartheta_2}.
\]
Man kann das gleiche Resultat rascher durch den \gesperrt{Aronhold}'schen Process erlangen. F\"uhrt man n\"amlich in die Identit\"at
\[
\begin{aligned}
2(\vartheta_3,\vartheta_3)^2 &= A_{ff}A_{\varphi\varphi}-A_{f\varphi}^2\\
&= (ab)^2(\alpha\beta)^2-(a\alpha)^2(b\beta)^2
\end{aligned}
\]
durch einmaliges Deltairen statt des Symboles $\alpha$ das Symbol $r$ ein, so kommt, weil
\[
\begin{aligned}
\delta(\vartheta_3) &= \delta(f,\varphi)=(f,\psi)=-\vartheta_2\\
-4(\vartheta_3,\vartheta_2)^2 &= 2(ab)^2(\alpha r)^2-2(a\alpha)^2(br)^2,
\end{aligned}
\]
oder
\weq{III}{A_{\vartheta_3\vartheta_2}=-\frac{1}{2}\{A_{ff}A_{\varphi\psi}-A_{f\varphi}A_{f\psi}\}=(\vartheta_3\vartheta_2)^2.}
Analoge Relationen erh\"alt man f\"ur $A_{\vartheta_3\vartheta_1}$ und $A_{\vartheta_1\vartheta_2}$. Folglich sind Producte, welche die Factoren $(\vartheta_1\vartheta_2)^2$, $(\vartheta_1\vartheta_3)^2$, $(\vartheta_2\vartheta_3)^2$ besitzen, auf Formen des Systems reducibel.


% ===== pdf idx 164, printed 150.0 =====
% !PRINTED 150
\weberpage{150}

133. \gesperrt{Producte, welche die Klammerfactoren (A) linear enthalten, sind reducibel.} Es bleiben also nur mehr Producte zu untersuchen, in denen kein Klammerfactor in h\"oherer als in der ersten Potenz auftritt. Auch hier brauchen wir die Reducibilit\"at nur mehr f\"ur Producte nachzuweisen, in welchen die letzten sechs Klammerfactoren
\[
(a\vartheta_1)\quad(\alpha\vartheta_2)\quad(r\vartheta_3)\quad(\vartheta_1\vartheta_2)\quad(\vartheta_1\vartheta_3)\quad(\vartheta_2\vartheta_3)
\]
linear auftreten. Der Beweis gestaltet sich genau wie bei zwei simultanen quadratischen Formen. Wir haben nur zu zeigen, dass die Producte
\[
(a\vartheta_1)a_x\vartheta_{1y},\ \ldots,\ (\vartheta_1\vartheta_2)\vartheta_{1x}\vartheta_{2y},\ \ldots
\]
reducibel sind. Die hiezu n\"othigen Formeln gehen direct aus der Relation (2) Nr. 125
\[
(f,\varphi)_y=(a\alpha)a_x\alpha_y-\frac{1}{2}A_{f\varphi}(xy)
\]
hervor. Ersetzt man hierin n\"amlich die quadratische Form $\varphi$ durch die quadratische Form $\vartheta_1$, so kommt:
\weq{I}{(f,\vartheta_1)_y=(a\vartheta_1)a_x\vartheta_{1y}-\frac{1}{2}A_{f\vartheta_1}(xy);}
und substituirt man hierin noch $\vartheta_2$ f\"ur $f$, so kommt auch:
\weq{II}{(\vartheta_2,\vartheta_1)_y=(\vartheta_2\vartheta_1)\vartheta_{2x}\vartheta_{1y}-\frac{1}{2}A_{\vartheta_2\vartheta_1}(xy).}
Beide Relationen lehren, dass die Producte
\[
(a\vartheta_1)a_x\vartheta_{1y}\quad\text{und}\quad(\vartheta_2\vartheta_1)\vartheta_{2x}\vartheta_{1y}
\]
reducibel sind; denn die Formen $A_{f\vartheta_1}$, $A_{\vartheta_2\vartheta_1}$ sind reducibel gem\"ass den Entwicklungen in Nr. 132, und die Formen $(f,\vartheta_1)$ und $(\vartheta_2,\vartheta_1)$ sind als Functionaldeterminanten von Functionaldeterminanten reducibel, und zwar hat man, wie aus der allgemeinen Formel (III) Nr. 53 direct hervorgeht:
\weq{III}{(f,\vartheta_1)=(f,(\varphi,\psi))=-\frac{1}{2}\{A_{f\varphi}\psi-A_{f\psi}\varphi\},}
und ebenso:
\[
(\vartheta_2,\vartheta_1)=((\psi,f),\vartheta_1)=-\frac{1}{2}\{A_{f\vartheta_1}\psi-A_{\psi\vartheta_1}f\},
\]
oder weil:
\[
A_{\psi\vartheta_1}=(r\vartheta_1)^2=(\psi,(\varphi,\psi))^2=0
\]
und
\[
A_{f\vartheta_1}=(a\vartheta_1)^2=R_{123}:
\]
so wird:
\weq{IV}{(\vartheta_2,\vartheta_1)=-\frac{1}{2}R_{123}\psi.}
Die Relationen (I) und (II) gehen also unter Benutzung dieser Werthe (III) und (IV) \"uber in:
%CONT>>


% ===== pdf idx 165, printed 151.0 =====
% !PRINTED 151
\weberpage{151}
%<<CONT
\weq{V}{
\begin{aligned}
(a\vartheta_1)a_x\vartheta_{1y} &= \frac{1}{2}\{A_{f\psi}\varphi_y\varphi_x-A_{f\varphi}\psi_y\psi_x+R_{123}(xy)\}\\
(\vartheta_1\vartheta_2)\vartheta_{2x}\vartheta_{1y} &= \frac{1}{2}\{R_{123}\psi_y\psi_x+\frac{1}{2}(xy)(A_{\psi\psi}A_{f\varphi}-A_{\psi\varphi}A_{\psi f})\}
\end{aligned}
}
Somit sind alle Producte, welche irgend einen der vorhandenen Klammerfactoren \gesperrt{linear} enthalten, reducibel. Durch Faltung derselben entstehen Producte mit quadratischen Klammerfactoren, die gleichfalls auf die Formen des Systemes zur\"uckgef\"uhrt werden k\"onnen. Demnach sind alle Klammerfactoren Reducenten und das simultane System dreier quadratischer Formen besteht aus keinen weiteren als den Eingangs aufgez\"ahlten dreizehn Formen.

134. \gesperrt{Relationen zwischen den Formen des Systemes.} Ich will zun\"achst der Uebersichtlichkeit halber wiederum die s\"ammtlichen, bisher gewonnenen Relationen tabellarisch ordnen, und durch einige neue erg\"anzen, die sich direct aus den allgemeinen Formeln \"uber Functionaldeterminanten ergeben.

1) Relationen zwischen den zweiten Ueberschiebungen.
\[
\alpha)\quad
\left.\begin{aligned}
A_{f\vartheta_2} &= A_{f\vartheta_3} = 0\\
A_{\varphi\vartheta_1} &= A_{\varphi\vartheta_3} = 0\\
A_{\psi\vartheta_1} &= A_{\psi\vartheta_2} = 0 % ?? uncertain: source shows a small mark on this \vartheta_2 subscript (dust speck or prime-like artifact); reproduced per the pattern of the other two rows
\end{aligned}\right\}\ \text{vergl. Nr. 127 (2),}
\]
\[
\beta)\quad
\left.\begin{aligned}
2A_{\vartheta_1\vartheta_1} &= -R_{\varphi\psi} = A_{\varphi\varphi}A_{\psi\psi}-A_{\varphi\psi}^2\\
2A_{\vartheta_2\vartheta_2} &= -R_{f\psi} = A_{ff}A_{\psi\psi}-A_{f\psi}^2\\
2A_{\vartheta_3\vartheta_3} &= -R_{f\varphi} = A_{ff}A_{\varphi\varphi}-A_{f\varphi}^2
\end{aligned}\right\}\ \text{vergl. Nr. 127 (1),}
\]
\[
\gamma)\quad
\left.\begin{aligned}
2A_{\vartheta_1\vartheta_2} &= -A_{\psi\psi}A_{f\varphi}+A_{\psi f}A_{\psi\varphi}\\
2A_{\vartheta_1\vartheta_3} &= -A_{\varphi\varphi}A_{f\psi}+A_{\varphi f}A_{\varphi\psi}
\end{aligned}\right\}\ \text{vergl. Nr. 132 (III).}
\]
\[
2A_{\vartheta_2\vartheta_3} = -A_{ff}A_{\varphi\psi}+A_{f\varphi}A_{f\psi}
\]
2) Relationen zwischen der Combinante und den zweiten Ueberschiebungen:
\[
\alpha)\quad R_{123}=A_{f\vartheta_1}=A_{\varphi\vartheta_2}=A_{\psi\vartheta_3},\ \text{vergl. Nr. 132 (I),}
\]
\[
\beta)\quad 2R_{123}^2=
\begin{vmatrix}
A_{ff} & A_{f\varphi} & A_{f\psi}\\
A_{\varphi f} & A_{\varphi\varphi} & A_{\varphi\psi}\\
A_{\psi f} & A_{\psi\varphi} & A_{\psi\psi}
\end{vmatrix}.
\]
Diese Relation ist eine unmittelbare Folge der allgemeinen Relation (II) Nr. 49. Man erkennt auch, dass die rechten Seiten der Relationen $\beta)$ und $\gamma)$ in (1) nichts anderes als die Minoren dieser Determinante sind. Aus dieser Relation kann man noch zahlreiche andere ableiten, wie sie die Lehrs\"atze \"uber adjungirte Determinanten
%CONT>>


% ===== pdf idx 166, printed 152.0 =====
% !PRINTED 152
\weberpage{152}
%<<CONT
(vergl. Bd. I Nr. 87, 88) liefern; ich will hierauf indess nicht weiter eingehen.

3) Relationen zwischen den Covarianten des Systems.

$\alpha)$ Die Relationen f\"ur die Functionaldeterminanten $(f,\vartheta_i)$, $(\varphi,\vartheta_i)$, $(\psi,\vartheta_i)$, $(\vartheta_i,\vartheta_k)$. Sie ergeben sich sofort aus der allgemeinen Formel Nr. 53 (III). Aus ihr erh\"alt man beispielsweise:
\[
\left.\begin{aligned}
2(f,\vartheta_1) &= -A_{f\varphi}\psi+A_{f\psi}\varphi\\
2(f,\vartheta_2) &= -A_{f\psi}f+A_{ff}\psi\\
2(f,\vartheta_3) &= -A_{ff}\varphi+A_{f\varphi}f
\end{aligned}\right\}\ \text{vergl. auch Nr. 127 (3)}
\]
und ebenso
\[
\left.\begin{aligned}
2(\vartheta_2,\vartheta_1) &= -R_{123}\psi\\
2(\vartheta_1,\vartheta_3) &= -R_{123}\varphi\\
2(\vartheta_3,\vartheta_2) &= -R_{123}f
\end{aligned}\right\}\ \text{vergl. Nr. 133 (IV).}
\]
$\beta)$ Darstellung der Functionaldeterminanten $\vartheta_i$ durch gerade Formen:
\[
\left.\begin{aligned}
-2\vartheta_1^2 &= A_{\varphi\varphi}\psi^2-2A_{\varphi\psi}\varphi\psi+A_{\psi\psi}\varphi^2\\
-2\vartheta_2^2 &= A_{ff}\psi^2-2A_{f\psi}f\psi+A_{\psi\psi}f^2\\
-2\vartheta_3^2 &= A_{ff}\varphi^2-2A_{f\varphi}f\varphi+A_{\varphi\varphi}f^2
\end{aligned}\right\}\ \text{Nr. 48 (I).}
\]
$\gamma)$ Allgemeine Relation f\"ur drei quadratische Formen:
\[
\begin{vmatrix}
A_{ff} & A_{f\varphi} & A_{f\psi} & f\\
A_{f\varphi} & A_{\varphi\varphi} & A_{\varphi\psi} & \varphi\\
A_{f\psi} & A_{\varphi\psi} & A_{\psi\psi} & \psi\\
f & \varphi & \psi & 0
\end{vmatrix}=0
\quad\text{und}\quad
\begin{vmatrix}
A_{\vartheta_1\vartheta_1} & A_{\vartheta_2\vartheta_1} & A_{\vartheta_3\vartheta_1} & \vartheta_1\\
A_{\vartheta_1\vartheta_2} & A_{\vartheta_2\vartheta_2} & A_{\vartheta_3\vartheta_2} & \vartheta_2\\
A_{\vartheta_1\vartheta_3} & A_{\vartheta_2\vartheta_3} & A_{\vartheta_3\vartheta_3} & \vartheta_3\\
\vartheta_1 & \vartheta_2 & \vartheta_3 & 0
\end{vmatrix}=0.
\]
Die beiden letzten Relationen gehen wieder unmittelbar aus der allgemeinen Formel (III) Nr. 50 hervor.

135. \gesperrt{Weitere Relationen zwischen den Formen des Systemes.} Bei der eingreifenden Bedeutung der Theorie quadratischer Formen f\"ur andere Untersuchungen will ich diese Tabelle noch durch einige weitere Relationen erg\"anzen, die wir im Folgenden ableiten wollen. Aus den Identit\"aten
\[
\begin{aligned}
\bar a_0 x_1^2 + 2\bar a_1 x_1 x_2 + a_2 x_2^2 - f &= 0 % [sic] third coefficient printed without a bar, unlike the barred triples in the following three rows
\\
\bar\alpha_0 x_1^2 + 2\bar\alpha_1 x_1 x_2 + \bar\alpha_2 x_2^2 - \varphi &= 0\\
\bar r_0 x_1^2 + 2\bar r_1 x_1 x_2 + \bar r_2 x_2^2 - \psi &= 0\\
y_2^2 x_1^2 - 2y_1y_2x_1x_2 + y_1^2x_2^2 - (xy)^2 &= 0
\end{aligned}
\]
eliminiren wir die vier Gr\"ossen $x_1^2, 2x_1x_2, x_2^2, -1$ und erhalten die nothwendig verschwindende Determinante:
%CONT>>


% ===== pdf idx 167, printed 153.0 =====
% !PRINTED 153
\weberpage{153}

\[
\begin{vmatrix}
\overline{a_0} & \overline{a_1} & \overline{a_2} & f \\
\overline{\alpha_0} & \overline{\alpha_1} & \overline{\alpha_2} & \varphi \\
\overline{r_0} & \overline{r_1} & \overline{r_2} & \psi \\
y_2{}^2 & -y_1y_2 & y_1{}^2 & (xy)^2
\end{vmatrix} = 0.
\]

Ordnen wir sie nach der letzten Colonne, so erhalten wir

\[
R\,(xy)^2 = (\overline{\alpha_0}\,\overline{r_1}\,y_1{}^2)\cdot f - (\overline{a_0}\,\overline{r_1}\,y_1{}^2)\cdot\varphi + (\overline{a_0}\,\overline{\alpha_1}\,y_1{}^2)\cdot\psi\,.
\]

Nun ist aber:

\[
(\overline{\alpha_0}\,\overline{r_1}\,y_1{}^2) =
\begin{vmatrix}
\overline{\alpha_0} & \overline{\alpha_1} & \overline{\alpha_2} \\
\overline{r_0} & \overline{r_1} & \overline{r_2} \\
y_1{}^2 & -y_1y_2 & y_1{}^2
\end{vmatrix}
=
\begin{vmatrix}
\alpha_1{}^2 & \alpha_1\alpha_2 & \alpha_2{}^2 \\
r_1{}^2 & r_1r_2 & r_2{}^2 \\
y_2{}^2 & -y_1y_2 & y_1{}^2
\end{vmatrix}
\]
\[
= (\alpha_1 r_2 - \alpha_2 r_1)(\alpha_1 y_1 + \alpha_2 y_2)(r_1 y_1 + r_2 y_2)
\]
\[
= (\alpha r)\,\alpha_y r_y = \vartheta_1(y)\,.
\]

Demnach besteht folgende Relation:

\weq{I}{R\,(xy)^2 = f(x)\,\vartheta_1(y) + \varphi(x)\,\vartheta_2(y) + \psi(x)\,\vartheta_3(y)}

oder auch, wenn wir $x$ mit $y$ vertauschen,

\weq{II}{R\,(xy)^2 = f(y)\,\vartheta_1(x) + \varphi(y)\,\vartheta_2(x) + \psi(y)\,\vartheta_3(x)\,.}

Beide bilden wieder eine Quelle neuer Relationen. F\"ur $y_i = x_i$ ergiebt sich aus beiden Gleichungen (I) und (II) die Relation:

\weq{III}{f\vartheta_1 + \varphi\vartheta_2 + \psi\vartheta_3 = 0\,.}

Ersetzen wir dagegen in (II) $y_1$ durch $-a_2 (= -b_2)$, $y_2$ durch $a_1 (= -b_1)$ etc., so kommt:

\[
\left.
\begin{aligned}
R f &= A_{ff}\vartheta_1 + A_{f\varphi}\vartheta_2 + A_{f\psi}\vartheta_3 \\
R \varphi &= A_{f\varphi}\vartheta_1 + A_{\varphi\varphi}\vartheta_2 + A_{\varphi\psi}\vartheta_3 \\
R \psi &= A_{f\psi}\vartheta_1 + A_{\varphi\psi}\vartheta_2 + A_{\psi\psi}\vartheta_3
\end{aligned}
\right\}\,. \tag{IV$_\text{a}$}
\]

Ebenso findet man aus der Relation (I), indem man darin die Variable $y$ durch die Coefficienten der Form $\vartheta_i$ ersetzt:

\[
\left.
\begin{aligned}
R \vartheta_1 &= A_{\vartheta_1\vartheta_1} f + A_{\vartheta_1\vartheta_2}\varphi + A_{\vartheta_1\vartheta_3}\psi \\
R \vartheta_2 &= A_{\vartheta_2\vartheta_1} f + A_{\vartheta_2\vartheta_2}\varphi + A_{\vartheta_2\vartheta_3}\psi \\
R \vartheta_3 &= A_{\vartheta_3\vartheta_1} f + A_{\vartheta_3\vartheta_2}\varphi + A_{\vartheta_3\vartheta_3}\psi
\end{aligned}
\right\}\,. \tag{IV$_\text{b}$}
\]

In derselben Weise, wie wir seinerzeit (Nr. 48) die Relation f\"ur $2\vartheta^2(x)$
%CONT>>


% ===== pdf idx 168, printed 154.0 =====
% !PRINTED 154
\weberpage{154}

%<<CONT
ermittelt haben, k\"onnen wir nun auch Relationen f\"ur $2\vartheta^2_{1x}\vartheta^2_{1y}$ etc. berechnen. Es ist n\"amlich:

\[
\vartheta_1(x) = (\varphi,\ \psi) =
\begin{vmatrix}
\overline{\alpha_0} & \overline{\alpha_1} & \overline{\alpha_2} \\
\overline{r_0} & \overline{r_1} & \overline{r_2} \\
x_2{}^2 & -x_1x_2 & x_2{}^2
\end{vmatrix}
\]
\[
2\vartheta_1(y) =
\begin{vmatrix}
\alpha_2 & -2\alpha_1 & \overline{\alpha_0} \\
\overline{r_2} & -2\overline{r_1} & \overline{r_0} \\
y_1{}^2 & 2y_1y_2 & y_2{}^2
\end{vmatrix}\,.
\]

Durch Multiplication beider Determinanten erhalten wir:

\[
2\vartheta_1(x)\vartheta_1(y) =
\begin{vmatrix}
A_{\varphi\varphi} & A_{\varphi\psi} & \varphi(x) \\
A_{\psi\varphi} & A_{\psi\psi} & \psi(x) \\
\varphi(y) & \psi(y) & (xy)^2
\end{vmatrix}\,. \tag{V}
\]

Analoge Relationen hat man f\"ur $2\vartheta_2(x)\vartheta_2(y)$, $2\vartheta_3(x)\vartheta_3(y)$. R\"andern wir ferner die Determinante

\[
R = (\overline{a_0}\,\overline{\alpha_1}\,\overline{r_2})
\]

einmal vertical mit Elementen $0$ und horizontal mit $x_1{}^2$, $-x_1x_2$, $x_2{}^2$, ein zweites Mal, nachdem wir in derselben Determinante die Colonnen gest\"urzt und die zweite Colonne mit $-2$ multiplicirt haben, vertical mit $0$, horizontal mit $y_1{}^2$, $2y_1y_2$, $y_2{}^2$ und bilden das Product der so erhaltenen zwei viergliedrigen Determinanten, so kommt als Resultat:

\[
\begin{vmatrix}
A_{ff} & A_{f\varphi} & A_{f\psi} & f(y) \\
A_{\varphi f} & A_{\varphi\varphi} & A_{\varphi\psi} & \varphi(y) \\
A_{\psi f} & A_{\psi\varphi} & A_{\psi\psi} & \psi(y) \\
f(x) & \varphi(x) & \psi(x) & (xy)^2
\end{vmatrix} = 0\,. \tag{VI}
\]

R\"andern wir endlich die Determinante $R$ durch zwei Saumreihen, n\"amlich einmal vertical durch die Doppelreihe $0\ 0$, und horizontal durch die Zeilen

\[
\begin{aligned}
x_2{}^2 & -x_1x_2 & x_1{}^2 & 0 & 0 \\
y_2{}^2 & -y_1y_2 & y_1{}^2 & 0 & 0\,,
\end{aligned}
\]

ein zweites Mal, nachdem wir die Colonnen gest\"urzt und die zweite Colonne mit $-2$ multiplicirt haben, vertical wieder durch die Doppelreihe $0\ 0$, aber horizontal durch die Zeilen

\[
\begin{aligned}
x_1{}^2 & 2x_1x_2 & x_2{}^2 & 0 & 0 \\
y_1{}^2 & 2y_1y_2 & y_2{}^2 & 0 & 0\,,
\end{aligned}
\]
%CONT>>


% ===== pdf idx 169, printed 155.0 =====
% !PRINTED 155
\weberpage{155}

%<<CONT
so entsteht durch Multiplication der beiden f\"unfgliedrigen Determinanten die letzte Relation:

\[
\begin{vmatrix}
A_{ff} & A_{f\varphi} & A_{f\psi} & f(x) & f(y) \\
A_{\varphi f} & A_{\varphi\varphi} & A_{\varphi\psi} & \varphi(x) & \varphi(y) \\
A_{\psi f} & A_{\psi\varphi} & A_{\psi\psi} & \psi(x) & \psi(y) \\
f(x) & \varphi(x) & \psi(x) & 0 & (xy)^2 \\
f(y) & \varphi(y) & \psi(y) & (xy)^2 & 0
\end{vmatrix} = 0\,. \tag{VII}
\]

Damit m\"ogen die Relationen, welche zwischen den Formen des Systemes stattfinden, ihren Abschluss finden. Wir werden in mannigfachen Untersuchungen von ihnen Gebrauch machen.

136. \textit{Das simultane System von $n$ quadratischen Formen.} Aus dem System der drei Formen $f$, $\varphi$ und $\psi$ l\"asst sich nun leicht das System beliebig vieler quadratischer Formen erschliessen. Bezeichnen wir die Formen mit $f_i$, so besteht dasselbe

1) aus den $n$ Formen $f_1, f_2, f_3 \ldots f_n$ selbst,

2) aus den $\dfrac{n(n-1)}{2}$ Functionaldeterminanten $\vartheta_{ik} = (f_i, f_k)$,

3) aus den $\dfrac{n(n+1)}{2}$ Invarianten $A_{f_i f_k} = (f_i, f_k)^2$,\,*)

4) aus den $\dfrac{n\cdot(n-1)(n-2)}{1\cdot 2\cdot 3}$ Combinanten $R_{ikm} = (f_i, (f_k, f_m))^2 = (f_k, (f_m, f_i))^2 = (f_m, (f_i, f_k))^2$.\,*)

Man k\"onnte noch die Frage aufwerfen, ob nicht die Ueberschiebungen der Functionaldeterminanten zu neuen Formen f\"uhren, n\"amlich zweier solcher Functionaldeterminanten, welche zusammen vier verschiedene Grundformen enthalten, wie

\[
((f_i, f_k), (f_\varrho, f_\sigma))^\lambda\,, \qquad \lambda = 1,\ 2\,. \tag{I}
\]

Doch ist auch deren Reducibilit\"at bereits nachgewiesen. Die erste Ueberschiebung ist n\"amlich reducibel, als Functionaldeterminante von Functionaldeterminanten, die zweite nach Nr. 132 (II). Das System besteht also in der That nur aus den oben angef\"uhrten Formen. Zwischen diesen bestehen nat\"urlich wieder eine grosse Zahl von Relationen, welche den im Vorhergehenden aufgestellten entsprechen,

\noindent\rule{4cm}{0.4pt}

*) Legt man, wie Herr Study es gethan hat, bei der symbolischen Darstellung einer Form $f$ nicht wie wir lineare Factoren zu Grunde, sondern quadratische, so lassen sich alle Invarianten von $f = a_x^{2n}$ darstellen durch simultane Invarianten dieser beiden Formengruppen vom Typus $(ab)^2$ und $(ab)(ac)(bc)$.
%CONT>>


% ===== pdf idx 170, printed 156.0 =====
% !PRINTED 156
\weberpage{156}

%<<CONT
von denen ich hier nur diejenigen anf\"uhren will, die sich auf die unter (I) angegebenen Ueberschiebungen beziehen. Man hat:

\[
((f_i, f_k), (f_\varrho, f_\sigma))^0 = \vartheta_{ik}\cdot\vartheta_{\varrho\sigma}
= -\frac{1}{2}\left\{ A_{i\varrho} f_k f_\sigma - A_{i\sigma} f_k f_\varrho + A_{k\sigma} f_i f_\varrho - A_{k\varrho} f_i f_\sigma \right\}\,, \qquad \text{Nr. 48 (II)}
\]

\[
((f_i, f_k), (f_\varrho, f_\sigma))^1 = (\vartheta_{ik},\ \vartheta_{\varrho\sigma})
= \frac{1}{2}\left\{ R_{i\varrho\sigma} f_k - R_{k\varrho\sigma} f_i \right\} \qquad \text{Nr. 53 (III)}
\]

\[
((f_i, f_k), (f_\varrho, f_\sigma))^2 = (\vartheta_{ik},\ \vartheta_{\varrho\sigma})^2
= \frac{1}{2}\left\{ A_{i\varrho} A_{k\sigma} - A_{i\sigma} A_{k\varrho} \right\}\,. \qquad \text{Nr. 132 (II)}
\]

\[
(A_{\varrho\varrho},\ A_{kk},\ A_{\sigma\sigma},\ A_{ii}) = 0\,. \qquad \text{Nr. 50.}
\]

\begin{center}
\S.\,13. \textbf{Specielle quadratische Formen.}
\end{center}

137. \textit{Conjugirte quadratische Formen.} Bei der Cayley'schen Methode der Aufl\"osung einer Gleichung vierten Grades mit Hilfe der Invariantentheorie (vgl. Nr. 175) zeigen sich drei quadratische Formen $\varphi$, $\psi$, $\chi$ (die quadratischen Factoren der Covariante $t$), deren zweite Ueberschiebungen $A_{\varphi\psi}$, $A_{\varphi\chi}$, $A_{\psi\chi}$ identisch verschwinden. Wir wollen allgemein drei solche quadratische Formen $f_0$, $f_1$, $f_2$ conjugirt nennen, sobald ihre simultanen Invarianten $A_{f_0f_1}$, $A_{f_0f_2}$, $A_{f_1f_2}$ identisch null sind. Das volle System dreier conjugirter quadratischer Formen enth\"alt nicht mehr 13, sondern nur 7 selbst\"andige Formen. Die Bedingung $A_{f_if_k} = 0$ bewirkt n\"amlich, dass die Functionaldeterminanten
\[
\vartheta_0 = (f_1, f_2)\,,\qquad \vartheta_1 = (f_0, f_2)\,,\qquad \vartheta_2 = (f_0, f_1)
\]
den Originalformen $f_0$, resp. $f_1$, $f_2$ proportional werden. Denn nach Formel Nr. 134, 3 $(\alpha)$ ist:

\[
2(f_0,\ \vartheta_0) = -A_{f_0f_1}\cdot f_2 + A_{f_0f_2}\cdot f_1\,.
\]

Da aber die beiden Invarianten rechts nach Voraussetzung verschwinden, so ist:

\[
(f_0,\ \vartheta_0) = 0
\]

und ebenso $(f_1,\ \vartheta_1)$, $(f_2,\ \vartheta_2)$. Das Verschwinden dieser Functionaldeterminanten lehrt aber nach Nr. 128, dass die beiden Formen, aus denen sie gebildet wurden, einander proportional sind. Man hat demnach

\[
\begin{aligned}
\vartheta_0 &= c_0 \cdot f_0 \\
\vartheta_1 &= c_1 \cdot f_1 \\
\vartheta_2 &= c_2 \cdot f_2\,.
\end{aligned}
\]

138. \textit{Kanonische Form dreier conjugirter Formen.} Man kann sich fragen: Welchen algebraischen Ausdruck besitzen drei conjugirte
%CONT>>


% ===== pdf idx 171, printed 157.0 =====
% !PRINTED 157
\weberpage{157}
%<<CONT
Formen, wenn man eine derselben in ihrer einfachsten kanonischen
Form $f_0 = 2x_1x_2$ zu Grunde legt? Zur Beantwortung dieser Frage
nehmen wir an $f_1 = \alpha_x^2$ sei eine zu $f_0 = a_x^2$ conjugirte Form, dann
muss, weil nach Voraussetzung
\[
A_{f_0f_1} = \bar a_0\bar a_2 - 2\bar a_1\bar a_1 + \bar a_0\bar a_2 = 0, \quad \bar a_0 = \bar a_2 = 0, \quad \bar a_1 = 1,
\]
auch $\alpha_1 = 0$ sein, d. h. die betreffende Form $f_1$ muss dem Büschel
angehören:
\[
F = \bar\alpha_0 x_1^2 + \bar\alpha_2 x_2^2.
\]
Sind nun $f_1$ und $f_2$ zwei Formen des Büschels $F$, die nicht nur zu $f_0$,
sondern auch unter einander conjugirt sind, so müssen ausserdem noch
deren Coefficienten der Bedingung genügen:
\[
A_{f_1f_2} = \alpha_0\bar\alpha_2' + \bar\alpha_0'\bar\alpha_2 = 0.
\]
Hieraus ergiebt sich, dass die drei conjugirten Formen in diesem Falle
die Gestalt annehmen:
\[
\begin{aligned}
f_0 &= 2x_1x_2 \\
f_1 &= \bar\alpha_0 x_1^2 + \bar\alpha_2 x_2^2 \\
f_2 &= (\bar\alpha_0 x_1^2 - \bar\alpha_2 x_2^2)\dfrac{\bar\alpha_0'}{\bar\alpha_0}.
\end{aligned}
\]
\textit{Anmerkung.} Drei specielle Formen dieser Art erhalten wir, wenn
wir hierin $\bar\alpha_2 = 1$ setzen und überdies $\bar\alpha_0$ und $\bar\alpha_0'$ so bestimmen, dass
die Discriminanten $A_{f_1f_1}$, $A_{f_2f_2}$ denselben Werth $-2$ annehmen, den
auch $A_{f_0f_0}$ besitzt. Unter diesen Voraussetzungen wird
\[
\bar\alpha_0 = -1, \quad \bar\alpha_0' = i
\]
und somit werden diese drei speciellen Formen:
\[
\begin{aligned}
f_0 &= \phantom{-}2x_1x_2 \\
f_1 &= -(x_1^2 - x_2^2) \\
f_2 &= +i(x_1^2 + x_2^2).
\end{aligned}
\]
Für diese drei Formen berechnen sich die Proportionalitätsconstanten
in den am Schlusse von Nr. 137 erwähnten Gleichungen:
\[
\vartheta_0 = c_0f_0, \quad \vartheta_1 = c_1f_1, \quad \vartheta_2 = c_2f_2
\]
sehr einfach. Denn gerade wegen dieser Beziehungen reducirt sich
die allgemeine Relation Nr. 134 (2, $\alpha$)
\[
R_{012} = A_{f_0\vartheta_0} = A_{f_1\vartheta_1} = A_{f_2\vartheta_2}
\]
auf
\[
R_{012} = c_0A_{f_0f_0} = c_1A_{f_1f_1} = c_2A_{f_2f_2},
\]
%CONT>>


% ===== pdf idx 172, printed 158.0 =====
% !PRINTED 158
\weberpage{158}
%<<CONT
oder, indem wir die Zahlencoefficienten einführen, auf
\[
\left|\begin{array}{ccc} 0 & 1 & 0 \\ -1 & 0 & 1 \\ -i & 0 & i \end{array}\right| = -2c_0 = -2c_1 = -2c_3. % [sic]
\]
Daraus erhalten wir:
\[
c_0 = c_1 = c_2 = \sqrt{-1} = i.
\]
Das Product dieser drei quadratischen Formen ist eine Form sechsten
Grades
\[
t = \mathrm{Const.}\ x_1x_2(x_1^4 - x_2^4),
\]
welche Klein (vgl. Vorles. über das Ikosaeder § 10) Octaeder nennt.
Man kann nämlich die Gauss'sche Zahlenebene so auf eine im Nullpunkt
berührende Kugel projiciren, dass die sechs Wurzeln der Gleichung
$t = 0$ gerade die Eckpunkte eines in die Kugel einbeschriebenen
Octaeders werden. Wir werden auf das Octaeder noch später zu
sprechen kommen (vgl. § 19).

139. \textit{Die quadratischen Formen des Würfels.} Wir hatten uns in
Nr. 138 die Aufgabe vorgelegt, quadratische Formen aufzustellen,
deren zweite Ueberschiebungen $A_{f_if_k}$ $(i \gtrless k)$ verschwinden. Die Frage
kann man verallgemeinern: Lässt sich eine Gruppe von quadratischen
Formen, die keinen gemeinschaftlichen Factor besitzen, so bestimmen,
dass ihre sämmtlichen zweiten Ueberschiebungen $A_{f_if_k}$ $(i \gtrless k)$ einen
festen Werth $\varrho$ besitzen?

Zur Beantwortung wählen wir wiederum als eine der quadratischen
 Formen die Form
\[
f_0 = 2x_1x_2.
\]
Wir wissen, dass ihre Discriminante den Werth $-2$ besitzt, und wollen
annehmen, dass auch die übrigen Formen $f_i$ $(i = 1, 2 \dots n)$ eine
Normalform besitzen mögen, so dass ihre Discriminante
\[
A_{f_if_i} = -2
\]
sei. Dann lautet die Frage: Welches sind unter diesen Voraussetzungen
 die Formen $f_i$, deren zweite Ueberschiebungen $A_{f_if_k}$ $(i \gtrless k)$
sowohl unter einander als mit $f_0$ den festen Werth $\varrho$ besitzen? Es
mag von vornherein bemerkt sein, dass dieser Werth $\varrho$ von $\pm 2$ verschieden
 sein muss. Denn da die gesuchten Formen keinen gemeinschaftlichen
 Factor besitzen sollen, so müssen die Resultanten
\[
R_{f_if_k} \gtrless 0
\]
sein. Nun ist aber nach Nr. 134 (1, $\beta$)
%CONT>>


% ===== pdf idx 173, printed 159.0 =====
% !PRINTED 159
\weberpage{159}
%<<CONT
\[
R_{f_if_k} = \left|\begin{array}{cc} A_{f_if_k} & A_{f_if_i} \\ A_{f_kf_k} & A_{f_if_k} \end{array}\right| = \left|\begin{array}{cc} A_{f_if_k} & -2 \\ -2 & A_{f_if_k} \end{array}\right| = A_{f_if_k}^2 - 4 \gtrless 0.
\]
d. h. $A_{f_if_k}$ muss von $\pm 2$ verschieden sein.

140. Sei nun $f_1 = x^2 + 2bx + c$, wobei wir noch der Einfachheit
 halber den Coefficienten von $x^2$ gleich 1 annahmen, eine erste
Form von den verlangten Eigenschaften, und
\[
f_2 = \alpha x^2 + 2\beta x + \gamma
\]
eine zweite, dann müssen nach unseren Voraussetzungen die Beziehungen
 bestehen:
\weq{I}{
\left.\begin{aligned}
A_{f_1f_1} &= 2(c - b^2) = -2 \\
A_{f_0f_1} &= -2b = \varrho \\
A_{f_2f_2} &= 2(\alpha\gamma - \beta^2) = -2 \\
A_{f_0f_2} &= -2\beta = \varrho \\
A_{f_1f_2} &= \gamma + c\alpha - 2b\beta = \varrho
\end{aligned}\right\} .
}
Dies sind gerade fünf Gleichungen mit fünf Unbekannten. Aus den
ersten beiden bestimmen sich die Coefficienten von $f_1$ linear; aus den
drei letzten dagegen erhält man eine quadratische Gleichung für $\alpha$
oder $\gamma$, und erkennt somit, dass es ausser $f_0 f_1 f_2$ noch eine weitere
Form $f_3$ giebt, so dass die fünf Invarianten
\[
A_{f_0f_1}, \quad A_{f_0f_2}, \quad A_{f_0f_3}, \quad A_{f_1f_2}, \quad A_{f_1f_3}
\]
den Werth $\varrho$ besitzen. Das gestellte Problem verlangt aber, dass
auch $A_{f_2f_3} = \varrho$ sei. Da jedoch die Coefficienten von $f_2$ und $f_3$ durch
die obigen fünf Gleichungen schon völlig bestimmt sind, so kann
diese Forderung nur mehr eine Bedingung für $\varrho$ sein. Nehmen wir
also an, $\varrho$ sei eine positive Grösse und fragen uns, welchen numerischen
 Werth muss sie besitzen, damit allgemein
\[
A_{f_if_k} = \varrho
\]
ist. Die Beantwortung der Frage liefert uns die Relation Nr. 136 (4),
welche zwischen den zweiten Ueberschiebungen von vier Formen bestehen
 muss, und die wir allgemein in Nr. 50 abgeleitet haben. Dieselbe
 war:
\[
\left|\begin{array}{cccc}
(f,f)^2, & (f,\varphi)^2, & (f,\psi)^2, & (f,\chi)^2 \\
(\varphi,f)^2, & (\varphi,\varphi)^2, & (\varphi,\psi)^2, & (\varphi,\chi)^2 \\
(\psi,f)^2, & (\psi,\varphi)^2, & (\psi,\psi)^2, & (\psi,\chi)^2 \\
(\chi,f)^2, & (\chi,\varphi)^2, & (\chi,\psi)^2, & (\chi,\chi)^2
\end{array}\right| = 0.
\]


% ===== pdf idx 174, printed 160.0 =====
% !PRINTED 160
\weberpage{160}
Setzen wir in ihr die Werthe $+\varrho$ und $-2$ der zweiten Ueberschiebungen
 ein, so kommt:
\[
\left|\begin{array}{cccc}
-2, & \varrho, & \varrho, & \varrho \\
\varrho, & -2, & \varrho, & \varrho \\
\varrho, & \varrho, & -2, & \varrho \\
\varrho, & \varrho, & \varrho, & -2
\end{array}\right| = 0,
\]
oder, wenn man $-\dfrac{2}{\varrho} = \sigma$ setzt und geeignet reducirt:
\[
(1 - \sigma)^3 \left|\begin{array}{cccc}
\sigma & 1 & 1 & 1 \\
1 & -1 & 0 & 0 \\
1 & 0 & -1 & -1 \\
1 & 0 & 0 & -1
\end{array}\right| = 0,
\]
oder endlich: $(1 - \sigma)^3(\sigma + 3) = 0.$

141. Der cubische Factor $(1 - \sigma)^3 = 0$ liefert $\varrho = -2$ als Lösung;
diese haben wir gerade ausgeschlossen. Es bleibt als einzige Lösung
\[
\sigma = -3
\]
also
\[
\varrho = +\frac{2}{3}.
\]
Unter Benutzung dieses Werthes gehen die fünf Coefficientenbedingungen
(I) über in:
\[
\begin{aligned}
A_{f_1f_1} &= \phantom{-}c - b^2 = -1 \\
A_{f_0f_1} &= -b = \phantom{-}\dfrac{1}{3} \\
A_{f_2f_2} &= \phantom{-}\alpha\gamma - \beta^2 = -1 \\
A_{f_0f_2} &= -\beta = \phantom{-}\dfrac{1}{3} \\
A_{f_1f_2} &= \gamma + c\alpha - 2b\beta = \phantom{-}\dfrac{2}{3}.
\end{aligned}
\]
Daraus folgt zunächst $b = -\dfrac{1}{3}$, $c = -\dfrac{8}{9}$, $\beta = -\dfrac{1}{3}$. Demnach
erhält man für $\alpha$ und $\gamma$:
\weq{II}{
\left.\begin{aligned}
\gamma - \dfrac{8}{9}\alpha &= +\dfrac{8}{9} \\
\alpha\gamma &= -\dfrac{8}{9}
\end{aligned}\right\} .
}
Denken wir uns in $f_2$ die Constante $\gamma$ in der Form $-\dfrac{8}{9}\gamma'$, dann
gehen die Gleichungen (II) in die symmetrisch gebauten:
\[
\gamma' + \alpha = -1
\]
\[
\alpha\gamma' = \phantom{-}1
\]


% ===== pdf idx 175, printed 161.0 =====
% !PRINTED 161
\weberpage{161}
%<<CONT
über, woraus man erkennt, dass $\alpha$ und $\gamma'$ die beiden dritten Wurzeln $\varepsilon$ und $\varepsilon^2$ der Einheit sind. Die vier gesuchten Formen $f$ sind demnach:
\[
\begin{aligned}
f_0 &= 2\,x_1x_2\\
f_1 &= x_1^2 - \frac{2}{3}\,x_1x_2 - \frac{8}{9}\,x_2^2\\
f_2 &= \varepsilon\,x_1^2 - \frac{2}{3}\,x_1x_2 - \frac{8}{9}\,\varepsilon^2 x_2^2\\
f_3 &= \varepsilon^2\,x_1^2 - \frac{2}{3}\,x_1x_2 - \frac{8}{9}\,\varepsilon\,x_2^2\,.
\end{aligned}
\]
Damit ist die Aufgabe gelöst, quadratische Formen $f$ so zu bestimmen, dass ihre zweiten Ueberschiebungen $A_{ik}$ durchwegs den Werth $+\varrho$ besitzen.

\textit{Anmerkung.} Das Product der vier quadratischen Formen liefert eine Form achten Grades:
\[
W = \text{Const.}\ xy\,(x^6 - 8x^3y^3 - 8y^6)\,,
\]
wenn wir der Einfachheit halber $\frac{2}{3}x_2=y$, $x_1=x$ setzen. Klein nennt diese Form (vgl. Vorles. über das Ikosaeder \S.\,5 und 10) Würfel, da man wiederum die Gauss'sche Zahlenebene so auf eine im Nullpunkt derselben berührende Kugel projiciren kann, dass die acht Wurzeln von $W=0$ auf der Kugel acht Punkte bestimmen, die gerade die acht Ecken eines der Kugel einbeschriebenen Würfels bilden. Ueber diese Form $W$ vergleiche auch \S.\,19.

142. \textit{Relationen zwischen den vier Formen $f_i$.} Man erkennt direct, dass die Summe aller vier Formen identisch verschwindet; demnach hat man als erste Relation:
\weq{1}{f_0+f_1+f_2+f_3=0\,.}
Aus ihr folgt, dass die Summe irgend zweier derselben bis auf das Vorzeichen gleich der Summe der beiden andern ist. Bezeichnen wir eine solche Summe mit $\varphi_{ik}$, so hat man:
\weq{2}{\begin{aligned}
\varphi_{01} &= -\varphi_{23}\\
\varphi_{02} &= -\varphi_{13}\\
\varphi_{03} &= -\varphi_{12}\,.
\end{aligned}}
Da also die sechs Summen zu je zweien der vier Formen $f_i$ nur drei verschiedene absolute Werthe annehmen, so müssen sie sich als Wurzeln einer Gleichung sechsten Grades darstellen lassen, die sich durch die einfache Substitution
\[
\varphi_{ik}^2 = z
\]
auf eine' cubische $\Phi=0$ reduciren lässt. % [sic] stray mark printed after \textquotedbl{}eine\textquotedbl{}


% ===== pdf idx 176, printed 162.0 =====
% !PRINTED 162
\weberpage{162}
Die drei Formen $\varphi_{01}$, $\varphi_{02}$, $\varphi_{03}$ sind drei conjugirte quadratische Formen; denn man hat:
\weq{3}{\begin{aligned}
(\varphi_{01},\varphi_{02})^2 &= 0\\
(\varphi_{02},\varphi_{03})^2 &= 0\\
(\varphi_{03},\varphi_{01})^2 &= 0\,.
\end{aligned}}
Dies ergiebt sich sofort durch Rechnung. So ist beispielsweise:
\[
\begin{aligned}
(\varphi_{01},\varphi_{02})^2 &= (f_0+f_1,\,f_0+f_2)^2 = A_{00}+A_{01}+A_{02}+A_{12}\\
&= -2 + \frac{2}{3} + \frac{2}{3} + \frac{2}{3} = 0\,.
\end{aligned}
\]
Es ist ferner bemerkenswerth, dass auch eine lineare Relation zwischen den Quadraten der vier Formen $f_i$ existirt. Man verificirt nämlich sehr leicht durch Rechnung:
\weq{4}{f_0^2+f_1^2+f_2^2+f_3^2=0\,.}

\textit{Anmerkung.} 1. Wegen der Relationen (1) und (4) reducirt sich die Gleichung:
\[
(\xi-f_0)(\xi-f_1)(\xi-f_2)(\xi-f_3)=0
\]
auf die drei Glieder:
\[
\xi^4+A\xi+B=0\,.
\]
Durch die Substitution
\[
\xi=\eta\sqrt[3]{A}
\]
geht sie über in die Gleichung
\weq{I}{\eta^4+\eta+c=0\,,}
eine Gleichung vierten Grades, die nur einen einzigen Parameter $C$ besitzt. Sie hat eine cubische Resolvente, die oben erwähnte Gleichung $\Phi=0$, die natürlich ebenfalls nur diesen einen Parameter $C$ enthalten kann. Man könnte diesen Umstand benutzen — wenn man nicht aus praktischen Gründen davon absehen müsste — die allgemeine Gleichung vierten Grades aufzulösen. Zu dem Zwecke bringt man sie zunächst — was immer möglich ist — auf die Form (I), vergleicht die Parameter und löst nun die Gleichung (I) mit Hilfe der cubischen Resolvente, aus deren Wurzeln $\varphi_{ik}$ sich leicht die Wurzeln $f_i$ der vorgelegten Gleichung bestimmen lassen. (Vgl. auch \S.\,19.)

\textit{Anmerkung.} 2. Herr Brill hat im 20. Bande der Math. Annalen (Binäre Formen und Gleichung sechsten Grades) gleichfalls gezeigt: Man kann immer vier quadratische Formen so finden, dass zwischen ihnen sowohl als zwischen ihren Quadraten je eine lineare Relation besteht, also dass man, wenn $f_1$ $f_2$ $f_3$ $f_4$ die vier quadratischen Formen, die Beziehungen hat:
\weq{1}{f_1+f_2+f_3+f_4=0}
\weq{2}{\alpha_1f_1^2+\alpha_2f_2^2+\alpha_3f_3^2+\alpha_4f_4^2=0\,.}


% ===== pdf idx 177, printed 163.0 =====
% !PRINTED 163
\weberpage{163}
Er zeigt ferner, dass, wenn man irgend drei der vier Formen durch ihre drei Functionaldeterminanten $\vartheta_{ik}$ ersetzt, zwischen dieser und der vierten Form ebenfalls analoge Relationen bestehen. Da dies auf vier verschiedene Arten möglich ist, so haben folgende fünf Formenquadrupel
\[
\begin{array}{cccc}
f_1 & f_2 & f_3 & f_4\\
f_1 & \vartheta_{23} & \vartheta_{34} & \vartheta_{42}\\
\vartheta_{13} & f_2 & \vartheta_{34} & \vartheta_{41}\\
\vartheta_{12} & \vartheta_{24} & f_3 & \vartheta_{41}\\
\vartheta_{12} & \vartheta_{23} & \vartheta_{31} & f_4
\end{array}
\]
die nämliche Eigenschaft, gleichzeitig durch Relationen von der Form (1), (2) aneinander gebunden zu sein. Ist also ein Quadrupel gefunden, so erhält man aus ihm stets vier andere. Hiebei kann man jedes der fünf Quadrupel in irgend einer Permutation zu Grunde legen, um auf dieselben vier andern zu gelangen. Man hat also $24\cdot5=120$ Operationen zur Verfügung, um von einem Quadrupel auf ein anderes der fünf zu gelangen, und insofern bilden die fünf Quadrupel eine Gruppe im Sinne der Gleichungstheorie.

142. % [sic] section number \textquotedbl{}142.\textquotedbl{} repeats the one used earlier in this fragment for \textquotedbl{}Relationen zwischen den vier Formen f_i\textquotedbl{} — reproduced as printed
\textit{Die quadratischen Formen des Ikosaeders.} Indem wir der Invariante $A_{f_2f_3}$ die Bedingung auferlegten, dass auch sie in Uebereinstimmung mit den fünf übrigen den constanten Werth $+\varrho$ besitzen möge, gelangten wir zu vier quadratischen Formen, die durch diese Bedingung völlig bestimmt waren. Es stand uns damals aber frei, für $A_{f_2f_3}$ den Werth $\pm\varrho$ zu wählen, und wir wollen nun untersuchen, zu welchem Resultate die Wahl $A_{f_2f_3}=-\varrho$ führt. In diesem Falle wird die Bedingung für den numerischen Werth von $\varrho$ dargestellt durch die Determinante
\[
\begin{vmatrix}
-2 & \varrho & \varrho & \varrho\\
\varrho & -2 & \varrho & \varrho\\
\varrho & \varrho & -2 & -\varrho\\
\varrho & \varrho & -\varrho & -2
\end{vmatrix} = 0\,.
\]
Wir ersetzen in ihr wiederum $-\dfrac{2}{\varrho}$ durch $\sigma$ und erhalten alsdann
\[
\begin{aligned}
\begin{vmatrix}
\sigma & 1 & 1 & 1\\
1 & \sigma & 1 & 1\\
1 & 1 & \sigma & -1\\
1 & 1 & -1 & \sigma
\end{vmatrix}
&=
\begin{vmatrix}
\sigma-1, & 1-\sigma, & 0, & 0\\
1, & \sigma, & 1, & 1\\
0, & 0, & (\sigma+1), & -(\sigma+1)\\
1, & 1, & -1, & \sigma
\end{vmatrix}\\[4pt]
&= (\sigma^2-1)
\begin{vmatrix}
1 & -1 & 0 & 0\\
1 & \sigma & 1 & 1\\
0 & 0 & 1 & -1\\
1 & 1 & -1 & \sigma
\end{vmatrix} = 0\,.
\end{aligned}
\]


% ===== pdf idx 178, printed 164.0 =====
% !PRINTED 164
\weberpage{164}
Der Factor $\sigma^2-1=0$ führt wiederum auf den Werth $\varrho=\pm 2$, den wir aus gleichen Gründen wie früher verwerfen müssen. Der andere Factor reducirt sich auf
\[
(\sigma+1)(\sigma-1)-4=0
\]
oder
\[
\sigma=\pm\sqrt5\,.
\]
Die Wahl des Vorzeichens steht uns frei. Wir nehmen zunächst
\[
\sigma=+\sqrt5
\]
und erhalten alsdann
\[
\varrho=\frac{-2}{\sqrt5}\,.
\]
Die Coefficientenrelationen (I) in Nr. 140 gehen dann, wenn wir diesmal den noch willkürlichen Coefficienten von $x^2$ in $f_1$ gleich $\dfrac{2}{\sqrt5}$ wählen, anstatt ihn wie früher gleich 1 zu setzen, über in:
\weq{III}{\left.\begin{aligned}
\frac{2}{\sqrt5}\,c - b^2 &= -1\\
b &= \frac{1}{\sqrt5}\\
\alpha\gamma-\beta^2 &= -1\\
\beta &= \frac{1}{\sqrt5}\\
\frac{2}{\sqrt5}\,\gamma+c\alpha-2b\beta &= \frac{-2}{\sqrt5}
\end{aligned}\right\}}
Hieraus folgt zunächst:
\[
b=\frac{1}{\sqrt5}\,,\qquad \beta=\frac{1}{\sqrt5}\,,\qquad c=\frac{-2}{\sqrt5}\,.
\]
Die übrigen Gleichungen gehen durch Substitution dieser Werthe über in:
\[
\frac{2}{\sqrt5}\,\gamma-\frac{2}{\sqrt5}\,\alpha-\frac{2}{5}=\frac{-2}{\sqrt5}
\]
\[
\alpha\gamma=\frac{-4}{5}\,.
\]
Wenn wir hier wiederum, um die Gleichungen symmetrischer zu gestalten, in $f_2$ den Coefficienten $\alpha$ durch $\dfrac{2}{\sqrt5}\alpha'$ und den Coefficienten $\gamma$ durch $\dfrac{-2}{\sqrt5}\gamma'$ ersetzen, so erhalten wir:
\weq{IV}{\left.\begin{aligned}
\alpha'+\gamma' &= \frac{-1+\sqrt5}{2}\\
\alpha'\gamma' &= 1
\end{aligned}\right\}}


% ===== pdf idx 179, printed 165.0 =====
% !PRINTED 165
\weberpage{165}
Nun sind bekanntlich die fünften Wurzeln der Einheit
\[
\begin{aligned}
\varepsilon &= \dfrac{-1+\sqrt5}{4} + i\,\dfrac{\sqrt{10+2\sqrt5}}{4}, &
\varepsilon^3 &= \dfrac{-1-\sqrt5}{4} - i\,\dfrac{\sqrt{10-2\sqrt5}}{4} \\
\varepsilon^2 &= \dfrac{-1-\sqrt5}{4} + i\,\dfrac{\sqrt{10-2\sqrt5}}{4}, &
\varepsilon^4 &= \dfrac{-1+\sqrt5}{4} - i\,\dfrac{\sqrt{10+2\sqrt5}}{4} \\
\varepsilon^5 &= 1. &&
\end{aligned}
\]
Demnach kann man die Gleichungen (IV) auch in der Form schreiben
\weq{V}{\begin{aligned} \alpha' + \gamma' &= \varepsilon + \varepsilon^4 \\ \alpha'\,\gamma' &= \varepsilon\cdot\varepsilon^4 \end{aligned}}
oder, wenn man auch in den Gleichungen (IV) der $\sqrt5$ das negative Zeichen giebt, was nach unseren Entwicklungen erlaubt ist,
\weq{VI}{\begin{aligned} \alpha' + \gamma' &= \varepsilon^2 + \varepsilon^3 \\ \alpha'\,\gamma' &= \varepsilon^2\cdot\varepsilon^3. \end{aligned}}
Die beiden Gleichungen (V) und (VI) aber lehren:

\glqq Die Coefficienten $\alpha'$ und $\gamma'$ der quadratischen Form $f_2$ sind je ein Paar $\varepsilon^\nu$ und $\varepsilon^\mu$ der fünften Einheitswurzeln, wenn man dieselben so auswählt, dass die Exponentensumme $\nu+\mu=5$ ist.\grqq{}

Es existiren demnach vier Formen $f_2$, welche die verlangten Eigenschaften besitzen. Wir erhalten also sechs quadratische Formen von der Beschaffenheit, dass ihre fünfzehn zweiten Ueberschiebungen $A_{f_if_k}$ $(i\gtrless k)$ bis auf das Vorzeichen übereinstimmen. Sie sind, wenn wir die Formen $f_1,\dots f_5$ noch mit $\dfrac{\sqrt5}{2}$ multipliciren, dargestellt durch:
\[
\begin{aligned}
f_0 &= 2x_1x_2 \\
\frac{\sqrt5}{2}f_1 &= x_1^2 + x_1x_2 - x_2^2 \\
\frac{\sqrt5}{2}f_2 &= \varepsilon\,x_1^2 + x_1x_2 - \varepsilon^4 x_2^2 \\
\frac{\sqrt5}{2}f_3 &= \varepsilon^2 x_1^2 + x_1x_2 - \varepsilon^3 x_2^2 \\
\frac{\sqrt5}{2}f_4 &= \varepsilon^3 x_1^2 + x_1x_2 - \varepsilon^2 x_2^2 \\
\frac{\sqrt5}{2}f_5 &= \varepsilon^4 x_1^2 + x_1x_2 - \varepsilon\, x_2^2.
\end{aligned}
\]
\textit{Anmerkung.} Das Product dieser sechs quadratischen Formen bezeichnete Klein mit dem Namen \glqq Ikosaeder\grqq{} (vergl. a. a. O.), aus den nämlichen Gründen, die ihn zu den Bezeichnungen \glqq Oktaeder\grqq{} und \glqq Würfel\grqq{} veranlassten. Die Gleichung des Ikosaeders ist demnach
%CONT>>


% ===== pdf idx 180, printed 166.0 =====
% !PRINTED 166
\weberpage{166}
%<<CONT
\[
x_1x_2\big(x^5 - (\varepsilon+\varepsilon^4)^5 x_2^5\big)\big(x^5 - (\varepsilon^2+\varepsilon^3)^5 x_2^5\big) = 0\,, % [sic] printed \textquotedbl{}x^5\textquotedbl{} without subscript in the first two factors
\]
oder:
\[
x_1x_2\big(x_1^{10} + 11\,x_1x_2 - x_2^{10}\big) = 0.
\]
Ich möchte noch darauf aufmerksam machen, dass auch zwischen den sechs quadratischen Formen specielle Relationen bestehen müssen, und zwar vier an der Zahl, da sie nur von den zwei Parametern $x_1$ und $x_2$ abhängen. In der That findet man leicht durch directe Rechnung, wenn man die erste Form $f_0$ noch mit $\dfrac{\sqrt5}{2}$ multiplicirt und nun $\dfrac{\sqrt5}{2}f_i$ mit $\psi_i$ bezeichnet, die vier Identitäten:
\weq{1}{\sum_{i=0}^{i=5}\psi_i^2 = 0\,,}
\weq{2}{\sum_{i=0}^{i=5}\psi_i^4 = 0\,,}
\weq{3}{\sum_{i=0}^{i=5}\psi_i^6 = 0\,,}
\weq{4}{\sum_{1}^{5}\psi_i^3 = \dfrac{(1+\sqrt5)^3}{(1-\sqrt5)^2\sqrt5}\Big(\sum_1^5 \psi_i\Big)^3.}
Aus den ersten drei Relationen geht sofort hervor, dass die Gleichung
\[
\prod_{i=0}^{i=5}(z-\psi_i^2) = 0
\]
die Form haben wird
\[
z^6 + az^2 + bz + c = 0.
\]
Durch die Transformation $z = \sqrt[4]{a}\cdot y$ geht sie über in die Form
\weq{5}{y^6 + y^2 + b'y + c' = 0.}
Aber auch diese beiden Constanten $b'$ und $c'$ sind noch von einander abhängig wegen der noch unbenutzten Relation (4). Sie sind beide Functionen ein und desselben Parameters $\lambda$, so dass also (5) eine Gleichung sechsten Grades mit einem \gesperrt{einzigen} Parameter darstellt. Wir werden später auf solche Gleichungen noch zu sprechen kommen. Nimmt man statt der sechs quadratischen Formen ihre ersten Polaren $f_{i_y}$, so bestehen immer noch vier analoge Identitäten, wie die oben angeführten. Aber die Gleichung sechsten Grades, welche die Quadrate dieser Formen zu Wurzeln hat, besitzt nunmehr zwei Parameter. Es ist dies dieselbe Gleichung sechsten Grades, welche Jacobi zwischen Thetareihen im 2. Bd. des Crelle'schen Journals erhalten hat. Setzt man nämlich:
%CONT>>


% ===== pdf idx 181, printed 167.0 =====
% !PRINTED 167
\weberpage{167}
%<<CONT
\[
\Theta(q) = 1 + 2q + 2q^4 + 2q^9 + 2q^{16} + 2q^{25} + \cdots
\]
und
\[
\varepsilon = e^{\frac{2i\pi}{5}}\,,
\]
so wird
\[
\Theta(\varepsilon^\varrho\cdot q) = 1 + 2\varepsilon^\varrho q + 2\varepsilon^{4\varrho}q^4 + 2\varepsilon^{4\varrho}q^9 + 2\varepsilon^\varrho q^{16} + 2q^{25} + \cdots\,,
\]
oder
\[
\Theta(\varepsilon^\varrho q) = \varepsilon^\varrho A_0 + A_1 + \varepsilon^{4\varrho}A_2\,,
\]
wobei
\[
\begin{aligned}
A_0 &= 2q + 2q^{16} + 2q^{36} + 2q^{81} + \cdots \\
A_1 &= 1 + 2q^{25} + 2q^{100} + 2q^{225} + \cdots \\
A_2 &= 2q^4 + 2q^9 + 2q^{49} + 2q^{64} + \cdots\,.
\end{aligned}
\]
Setzen wir aber $A_0 = x_1y_1$, $A_1 = (x_1y_2 + x_2y_1)$, $A_2 = -x_2y_2$, so haben wir:
\[
\begin{aligned}
\Theta(q) &= f_{1y} \\
\Theta(\varepsilon q) &= f_{2y} \\
\Theta(\varepsilon^2 q) &= f_{3y} \\
\Theta(\varepsilon^3 q) &= f_{4y} \\
\Theta(\varepsilon^4 q) &= f_{5y}\,,
\end{aligned}
\]
und überdies:
\[
\Theta(q^{25}) = f_{0y}\,.
\]
Es lassen sich also in der That diese sechs Thetareihen mit den Polaren der sechs quadratischen Formen $f_0, f_1, \dots f_5$ identificiren.

\begin{center}
\S.\,14.\\[2pt]
\textbf{Die Form dritten Grades.}
\end{center}

\textbf{143.} \textit{Einleitende Bemerkungen.} Wir haben bereits in den Beispielen Nr. 37 und Nr. 39, welche den Ueberschiebungsprocess erläutern sollten, eine Reihe von Co- und Invarianten der Form dritten Grades kennen gelernt. Bezeichnet man die cubische Form mit
\[
\begin{aligned}
f &= \bar a_0 x_1^3 + 3\bar a_1 x_1x_2 + 3\bar a_2 x_1x_2^2 + \bar a_3 x_2^3 \\
&= a_x^3 = b_x^3 = c_x^5 = \text{etc.}\,,
\end{aligned}
\]
% [sic] printed exponent "c_x^5"; the other symbolic powers in the line read a_x^3 = b_x^3
so haben die a. a. O. gebildeten Covarianten die symbolischen Producte
\[
\begin{aligned}
(f,f)^2 &= (ab)^2\,a_xb_x = \Delta \\
(\Delta,\Delta)^2 &= (ab)^2(cd)^2(ad)(bc) = A_{\Delta\Delta} = R \\
(f,\Delta) &= (ab)^2(ca)\,b_x c_x^2 = Q
\end{aligned}
\]
zum Repräsentanten. Schon damals bemerkten wir, dass mit diesen vier Formen
\[
f,\ \Delta,\ R,\ Q
\]
das volle System von $f$ abgeschlossen ist, d. h. dass alle übrigen Co- und Invarianten von $f$ sich rational und ganz durch dieselben dar-
%CONT>>


% ===== pdf idx 182, printed 168.0 =====
% !PRINTED 168
\weberpage{168}
%<<CONT
stellen lassen, und es ist unsere erste Aufgabe, diese Behauptung nachzuweisen. Da nun alle weiteren Co- und Invarianten nur durch Ueberschiebung der drei Covarianten $f$, $\Delta$, $Q$ entstehen können und somit durch symbolische Producte dargestellt sein müssen, die irgend welche der sechs Klammerfactoren
\[
(ab),\ (a\Delta),\ (\Delta\Delta_1),\ (aQ),\ (\Delta Q),\ (QQ_1)
\]
zu Factoren haben, so besteht unsere Aufgabe darin, den Nachweis zu liefern, dass jeder dieser sechs Klammerfactoren \gesperrt{Reducent} ist (vergl. Nr. 124), d. h. dass jedes symbolische Product, welches solche Klammerfactoren besitzt, entweder identisch verschwindet oder auf Producte mit Factoren $f$, $\Delta$, $R$, $Q$ reducirt werden kann. Indem ich nun diesen Beweis antrete, will ich nur den Beweis für die Reducenteneigenschaft des Klammerfactors $(ab)$ nochmals in voller Ausführlichkeit entwickeln. Da der Gedankengang bei den folgenden Beweisen derselbe bleibt, so wird alsdann eine grössere Kürze gestattet sein.

\textbf{144.} \textit{Der Factor $(ab)$ ist Reducent.} Jedes symbolische Product $P$, das den Factor $(ab)$ besitzt, muss in irgend welchen Verbindungen auch noch zwei weitere Symbole $a$ und zwei Symbole $b$ enthalten, wenn es eine wirkliche Covariante von $f$ darstellen soll. Wir drücken dies am allgemeinsten aus, wenn wir sagen:
\[
\pi = (ab)\,a_\xi a_x b_\eta b_y
\]
muss ein Factor des zu untersuchenden Productes $P$ sein, wobei $\xi$, $x$, $\eta$, $y$ durch irgend welche Symbole ersetzt sein können. Dieses Theilproduct $\pi$ kann als eine gemischte Polare der Form
\[
\Pi = (ab)\,a_x^2 b_y^2
\]
angesehen werden und ist mit dieser gleichzeitig reducibel. Die Form $\Pi$ ist aber ein Glied der zweiten Polare von
\[
(f,f) = (ab)\,a_x^2 b_x^2
\]
und lässt sich somit darstellen durch diese Polare $(f,f)_{y^2}$ selbst plus Gliedern, welche den Factor $(ab)^2(xy)$ besitzen (vergl. Nr. 25). $\Pi$ ist also reducibel, sobald $(f,f)_{y^2}$ und \gesperrt{auch} die Glieder mit dem Factor $(ab)^2(xy)$, oder was dasselbe ist, sobald $(f,f)$ und $(ab)^2a_xb_y$ reducibel sind. Der Term $(ab)^2a_xb_y$ ist jedoch ein Glied der ersten Polare von
\[
(f,f)^2 = (ab)^2a_xb_x
\]
und lässt sich somit ersetzen durch diese plus Gliedern mit dem Factor $(ab)^3(xy)$. Die Reducibilität von $\Pi$ ist also zurückgeführt auf die drei Bedingungen, dass $(f,f)$, $(f,f)^2$, $(f,f)^3$ reducibel sind, d. h.
%CONT>>


% ===== pdf idx 183, printed 169.0 =====
% !PRINTED 169
\weberpage{169}
%<<CONT
auf bekannte Formen führen. Es ist aber
\[
(f,\,f) = 0, \quad (f,\,f)^2 = \Delta, \quad (f,\,f)^3 = 0,
\]
und demnach ist $\Delta$ reducibel, und $(ab)$ Reducent.
% ?? uncertain: printed glyph before "reducibel" is small and could read as \pi; read here as \Delta by mathematical context with (f,f)^2=\Delta established immediately above

\textit{Anmerkung.} Ich möchte noch bemerken, dass sich diese Darstellung der Reducibilität mit der in Nr. 124 gegebenen vollständig deckt. Wir belegten damals Klammerfactoren mit dem Namen Reducenten, sobald das Product, welches sie allein besass, sammt allen daraus durch Faltung entstehenden reducibel war. In der That sind wir hier zum gleichen Resultate gelangt, denn $(f,\,f)^2$ und $(f,\,f)^{3}$ entstehen aus $(f,\,f)$ durch Faltung.

145. \textit{Die Factoren $(a\Delta)$ und $(\Delta\Delta_1)$ sind Reducenten.} Die erste Ueberschiebung von $f$ über $\Delta$ lieferte die Covariante $Q$; die zweite Ueberschiebung aber verschwindet identisch (vergl. auch Nr. 62, Beispiele). Denn ersetzt man in
\[
\Delta_x^2 = (ab)^2 a_x b_x
\]
$x$ durch $c$, so kommt:
\[
(c\Delta)^2 = (ab)^2 (ac)(bc),
\]
und somit:
\[
(f,\,\Delta)^2 = (c\Delta)^2 c_x = (ab)^2 (ac)(bc) c_x.
\]
Vertauscht man rechts einmal $c$ mit $b$, dann auch $c$ mit $a$ und addirt die drei so erhaltenen Ausdrücke für $(f,\,\Delta)^2$, so ergiebt sich:
\[
(f,\,\Delta)^2 = \frac{1}{3}(ab)(ac)(bc)\{(ab)c_x - (ac)b_x + (bc)a_x\},
\]
oder:
\weq{I}{(f,\,\Delta)^2 = 0,}
da der in der Klammer befindliche Ausdruck identisch verschwindet. Weitere Ueberschiebungen von $f$ über $\Delta$ existiren nicht; $(f,\,\Delta) = Q$, $(f,\,\Delta)^2 = 0$, folglich ist $(a\Delta)$ Reducent.

Ebenso ist $(\Delta\Delta_1)$ Reducent, weil $(\Delta,\,\Delta) = 0$, $(\Delta,\,\Delta)^2 = A_{\Delta\Delta}$ ist.

146. \textit{Der Factor $(aQ)$ ist Reducent.} Nach den Entwicklungen in Nr. 144 haben wir nur den Werth der drei Ueberschiebungen:
\[
(f,\,Q), \quad (f,\,Q)^2 \quad \text{und} \quad (f,\,Q)^3
\]
zu untersuchen. Wir schlagen hiebei ganz den regulären Gang des Ueberschiebungsprocesses ein, indem wir zunächst die erste Polare von $Q$ bilden. Es wird:
\weq{1}{3\,Q_x^2 Q_y = (c\Delta)\{c_x^2 \Delta_y + 2\Delta_x c_x c_y\} = G_1 + 2G_2.}
Wir können den Ausdruck rechts vereinfachen. Die Differenz aus der Polare und einem ihrer Glieder, nämlich
\[
3\,Q_x^2 Q_y - 3G_2 = (c\Delta)\{c_x^2 \Delta_y + 2\Delta_x c_x c_y\} - 3(c\Delta)\Delta_x c_x c_y,
\]
%CONT>>


% ===== pdf idx 184, printed 170.0 =====
% !PRINTED 170
\weberpage{170}
%<<CONT
hat den Factor $(c\Delta)^2(xy)$. Ein Product mit dem Factor $(c\Delta)^2$ ist aber identisch Null, da $(f,\,\Delta)^2 = 0$, und ein solches Product aus dieser Form durch Ueberschiebung entstanden gedacht werden kann. Folglich sind die beiden Glieder in (1) einander gleich, d. h. es ist:
\[
(c\Delta)c_x^2 \Delta_y = (c\Delta)\Delta_x c_y c_x
\]
und demnach wird die Polare:
\weq{2}{Q_x^2 Q_y = (c\Delta)c_x^2 \Delta_y = (c\Delta)c_x c_y \Delta_x.}
Hieraus folgt für $y_1 = -b_2$, $y_2 = b_1$:
\[
(bQ)Q_x^2 b_x^2 = (c\Delta)(bc)c_x \Delta_x b_x^2,
\]
oder, indem man $b$ mit $c$ vertauscht und beide Ausdrücke addirt:
\[
(f,\,Q) = \frac{1}{2}(bc)c_x b_x \Delta_x \big((c\Delta)b_x - (b\Delta)c_x\big)
= -\frac{1}{2}(bc)^2 b_x c_x \Delta_x^2,
\]
oder
\weq{I}{(f,\,Q) = -\frac{1}{2}\Delta^2.}
Bildet man ferner die erste Polare von $\Delta_x$, so kommt:
\weq{3}{\Delta_y \Delta_x = (ab)^2 a_x b_y.}
Wir ersetzen hierin $x$ durch $\Delta'$ und erhalten nach Multiplication mit $\Delta'_y$:
\[
(\Delta\Delta')\Delta_y \Delta'_y = (ab)^2(a\Delta')b_y \Delta'_y = (ab)^2(a\Delta)b_y \Delta_y.
\]
Substituiren wir dagegen in (2) für $x$ das Symbol $b$, so kommt:
\[
(Qb)^2 Q_y b_y = (cb)^2(c\Delta)b_y \Delta_y.
\]
Durch Comparation dieser beiden Gleichungen ergiebt sich:
\[
(Qb)^2 Q_y b_y = (\Delta\Delta')\Delta_y \Delta'_y.
\]
Rechts steht die erste Ueberschiebung von $\Delta_y^2$ über sich selbst, sie verschwindet identisch; also hat man:
\weq{II}{(f,\,Q)^2 = 0.}
Setzt man endlich in $Q = (c\Delta)c_x^2 \Delta_x$ das Symbol $b$ an die Stelle von $x$, so erhält man:
\[
(Qb)^3 = (c\Delta)(cb)^2(\Delta b).
\]
Dasselbe ergiebt sich, wenn man in $\Delta = (cb)^2 c_x b_x$ an Stelle von $x$ das Symbol $\Delta$ substituirt, nämlich:
\[
(\Delta,\,\Delta)^2 = (c\Delta)(cb)^2(b\Delta).
\]
Durch Comparation ergiebt sich also:
\weq{III}{(f,\,Q)^3 = -(Qb)^3 = (\Delta,\,\Delta)^2 = A_{\Delta\Delta}.}
Die drei Ueberschiebungen von $f$ über $Q$ sind somit auf bekannte Formen reducibel und demnach ist $(aQ)$ Reducent.


% ===== pdf idx 185, printed 171.0 =====
% !PRINTED 171
\weberpage{171}
147. \textit{Die Factoren $(\Delta Q)$ und $(QQ')$ sind Reducenten.} Der Nachweis, dass auch diese letzten Factoren Reducenten sind, lässt sich nun auf Grund der eben gewonnenen Beziehungen einfach führen. Da nach Nr. 146 (2):
\weq{1}{Q_x^2 Q_y = (c\Delta)c_x^2 \Delta_y,}
so hat man hieraus für $y_1 = -\Delta_2$, $y_2 = +\Delta_1$
\[
(\Delta Q)Q_x^2 \Delta_x = (c\Delta)(\Delta'\Delta)c_x^2 \Delta'_x
= \frac{1}{2}(\Delta'\Delta)c_x^2 \{(c\Delta)\Delta'_x - (c\Delta')\Delta_x\}
= +\frac{1}{2}(\Delta'\Delta)^2 \cdot c_x^3,
\]
oder
\weq{IV}{(\Delta,\,Q) = \frac{1}{2}\cdot A_{\Delta\Delta}\cdot f.}
Setzt man in derselben Relation (1) statt $x$ das Symbol $\Delta (= \Delta')$, so kommt:
\[
(Q\Delta)^2 Q_y = (c\Delta)(c\Delta')^2 \Delta_y = 0,
\]
da rechts der Reducent $(c\Delta)^2$ auftritt (vergl. Nr. 145 und 146). Demnach ist:
\weq{V}{(\Delta,\,Q)^2 = 0}
und folglich ist $(\Delta Q)$ Reducent.

Substituirt man endlich $Q'$ in diese Relation (1) für die Variable $x$, und multiplicirt mit $Q'_y (= Q_y)$, so erhält man:
\[
(QQ')^2 Q_y Q'_y = (c\Delta)(cQ)^2 \Delta_y Q_y.
\]
Der symbolische Ausdruck rechts besitzt aber den Werth $\frac{1}{2}A_{\Delta\Delta}\cdot\Delta$. Denn die erste Polare der verschwindenden Ueberschiebung $(f,\,Q)^2 = (aQ)^2 a_x Q_x$ ist:
\[
(f,\,Q)_y^2 = (aQ)^2\{a_x Q_y + Q_x a_y\} = 0;
\]
anderntheils ist:
\[
(aQ)^2\{a_x Q_y - Q_x a_y\} = (aQ)^3(xy).
\]
Folglich erhält man durch Addition:
\[
(aQ)^2 a_x Q_y = \frac{1}{2}(aQ)^3(xy),
\]
und wenn man hierin $x$ durch $\Delta$ ersetzt, und $(aQ)^3$ durch seinen Werth $A_{\Delta\Delta}$, so kommt endlich:
\[
(a\Delta)(aQ)^2 \Delta_y Q_y = \frac{1}{2}A_{\Delta\Delta}\cdot \Delta_y^2;
\]
also ist in der That:
\weq{VI}{(Q,\,Q)^2 = \frac{1}{2}A_{\Delta\Delta}\cdot \Delta.}


% ===== pdf idx 186, printed 172.0 =====
% !PRINTED 172
\weberpage{172}
Weil aber auch aus bekannten Gründen $(Q,\,Q) = 0$ und $(Q,\,Q)^3 = 0$, so ist $(Q,\,Q')$ Reducent.

Demnach sind alle überhaupt möglichen Klammerfactoren Reducenten, und es besteht somit der Satz:

\begin{quote}
\glqq Das volle System der Form dritten Grades $f$ besteht aus den vier Formen $f$, $\Delta$, $A_{\Delta\Delta}$, $Q$.\grqq{}
\end{quote}

148. \textit{Tabelle der Ueberschiebungswerthe dieser vier Formen.} Ich will die in diesem Beweise gewonnenen Resultate noch einmal tabellarisch ordnen. Ist
\[
f = a_x^3,
\]
so hat man als symbolische Darstellungen der Grundformen des Systemes:
\[
(f,\,f)^2 = \Delta \phantom{{}={}} = (ab)^2 a_x b_x
\]
\[
(\Delta,\,\Delta)^2 = A_{\Delta\Delta} = R \phantom{{}={}} = (ab)^2(cd)^2(ad)(bc)
\]
\[
(f,\,\Delta) = Q \phantom{{}={}} = (c\Delta)c_x^2 \Delta_x = (ab)^2(cb)c_x^2 a_x.
\]
Als Ueberschiebungswerthe erhielten wir:
\weq{1}{(f,\,f) = 0, \quad (f,\,f)^2 = \Delta, \quad (f,\,f)^3 = 0}
\weq{2}{(\Delta,\,\Delta) = 0, \quad (\Delta,\,\Delta)^2 = A_{\Delta\Delta} = R}
\weq{3}{(Q,\,Q) = 0, \quad (Q,\,Q)^2 = \frac{1}{2}A_{\Delta\Delta}\cdot \Delta, \quad (Q,\,Q)^3 = 0}
\weq{4}{(f,\,\Delta) = Q, \quad (f,\,\Delta)^2 = 0}
\weq{5}{(\Delta,\,Q) = \frac{1}{2}A_{\Delta\Delta}\cdot f, \quad (\Delta,\,Q)^2 = 0}
\weq{6}{(f,\,Q) = -\frac{1}{2}\Delta^2, \quad (f,\,Q)^2 = 0, \quad (f,\,Q)^3 = A_{\Delta\Delta}.}
Unter den Covarianten des Systemes befindet sich eine schiefe Form, nämlich die Covariante $Q$. Ihr Quadrat muss sich durch die übrigen geraden Formen darstellen lassen. Um dies zu bewerkstelligen, berücksichtigen wir, dass sie Functionaldeterminante von $f$ und $\Delta$ ist. Indem wir die allgemeine Formel für Functionaldeterminanten benutzen (Nr. 48, I), gelangen wir zu der bereits von Cayley entdeckten Relation zwischen den vier Grundformen des Systemes:
\[
2Q^2 =
\begin{vmatrix}
(f,\,f)^2, & (f,\,\Delta)^2, & f \\
(f,\,\Delta)^2, & (\Delta,\,\Delta)^2, & \Delta \\
f, & \Delta, & 0
\end{vmatrix}
=
\begin{vmatrix}
\Delta, & 0, & f \\
0, & R, & \Delta \\
f, & \Delta, & 0
\end{vmatrix}
\]
oder:
\weq{I}{2Q^2 + \Delta^3 + Rf^2 = 0.}


% ===== pdf idx 187, printed 173.0 =====
% !PRINTED 173
\weberpage{173}
Wir werden bei Auflösung der Gleichung dritten Grades von dieser Relation Gebrauch machen.

149. \textit{Das Formensystem von $Q + \lambda f$.} Wir sahen, dass die Form $f$ eine Covariante gleichen Grades in den Variabeln besitzt, nämlich die Form $Q$ und können uns demnach die Frage vorlegen: welchen Werth haben die Grundformen des Büschels
\[
Q + \lambda f,
\]
also die Formen
\[
\Delta_{Q+\lambda f}, \quad R_{Q+\lambda f}, \quad Q_{Q+\lambda f}.
\]
Wir können zu diesem Zwecke den Weg directer Ueberschiebung oder auch den \gesperrt{Aronhold'schen} Process benutzen. Wir wollen hier den ersten einschlagen, und erhalten in einfacher Weise:
\weq{1}{\Delta_{Q+\lambda f} = (Q+\lambda f,\, Q+\lambda f)^2 = (Q,\,Q)^2 + 2\lambda(Q,\,f)^2 + \lambda^2(f,\,f)^2 = \frac{1}{2}R\cdot\Delta + \lambda^2\Delta = \left(\lambda^2 + \frac{R}{2}\right)\Delta}
\weq{2}{R_{Q+\lambda f} = \left(\left(\lambda^2+\frac{R}{2}\right)\Delta,\, \left(\lambda^2+\frac{R}{2}\right)\Delta\right)^2 = \left(\lambda^2+\frac{R}{2}\right)^2(\Delta,\,\Delta)^2 = \left(\lambda^2+\frac{R}{2}\right)^2 R}
\weq{3}{Q_{Q+\lambda f} = \left(Q+\lambda f,\, \left(\lambda^2+\frac{R}{2}\right)\Delta\right) = \left(\lambda^2+\frac{R}{2}\right)\left\{(Q,\,\Delta) + \lambda(f,\,\Delta)\right\} = \left(\lambda^2+\frac{R}{2}\right)\left\{-\frac{R}{2}f + \lambda Q\right\}.}

150. \textit{Beispiele zur Uebung.} Wir hatten zur Berechnung der Tabelle in Nr. 148 lediglich den Ueberschiebungsprocess benutzt. Der Umstand, dass $f$ eine Covariante $Q$ besitzt, die gleichfalls dritten Grades in den Variabeln ist, legt den Gedanken nahe, mit Hilfe des verwandten \gesperrt{Aronhold'schen} Processes solche Ueberschiebungen in vielleicht einfacherer Weise zu ermitteln. In der That können wir, indem wir eine beliebige Ueberschiebung von bekanntem Werthe deltairen, auch Werthe von andern Ueberschiebungen berechnen. Ich erinnere zunächst daran, dass wir in Nr. 63, unter der Voraussetzung
\weq{I}{\delta f = Q,}
folgende Beziehungen erhielten:
\weq{II}{\delta Q = -\frac{1}{2}Rf, \quad \delta\Delta = 0, \quad \delta R = 0.}
Deltairen wir nun die beiden Relationen
\[
\begin{aligned}
(f,\,\Delta)^2 &= 0\\
(f,\,Q)^2 &= 0,
\end{aligned}
\]
so erhalten wir:
\weq{1}{\delta(f,\,\Delta)^2 = (\delta f,\,\Delta)^2 + (f,\,\delta\Delta)^2 = 0,}
%CONT>>


% ===== pdf idx 188, printed 174.0 =====
% !PRINTED 174
\weberpage{174}
%<<CONT
also wegen den Beziehungen (I) und (II)
\[
(Q,\,\Delta)^2 = 0
\]
\weq{2}{\delta(f,\,Q)^2 = (\delta f,\,Q)^2 + (f,\,\delta Q)^2 = 0,}
oder:
\[
(Q,\,Q)^2 + \left(f,\,-\frac{1}{2}R\cdot f\right)^2 = 0,
\]
oder:
\[
(Q,\,Q)^2 = \frac{R}{2}\cdot\Delta.
\]
Man erkennt, wie auf diesem Wege oft in sehr einfacher Weise Ueberschiebungen berechnet werden können. Wir werden bei Gelegenheit der Theorie biquadratischer Formen noch ausgiebiger von dieser Methode Gebrauch machen.

151. \textit{Erste Anwendung: Die Discriminante der Form $f$.} Bekanntlich ist die Discriminante einer Form darstellbar als Resultante der gleich null gesetzten ersten Differentialquotienten dieser Form. (Vergl. Bd. I, Nr. 173 und 178.) In unserm Falle hier sind die beiden ersten Differentialquotienten von $f = a_x^3 = b_x^3$ zwei quadratische Formen
\[
f_1 = a_1 a_x^2
\]
\[
f_2 = b_2 b_x^2,
\]
% ?? uncertain: exponent on $a_x$ in the printed formula for $f_1$ was hard to resolve against $f_2$'s $b_x^2$ — transcribed both as quadratic (degree 2) forms, consistent with "zwei quadratische Formen"
und die Resultante derselben ist nach Nr. 129 gleich der Discriminante ihrer Functionaldeterminante $\vartheta$, nämlich:
\weq{1}{2A_{\vartheta\vartheta} = -R_{f_1f_2}.}
Nun ist die Functionaldeterminante von $f_1$ und $f_2$ dargestellt durch
\[
\vartheta = (ab)\,a_x b_x a_1 b_2,
\]
oder, wenn man $a$ mit $b$ vertauscht:
\[
\vartheta = -(ab)\,a_x b_x b_2 a_1,
\]
folglich:
\[
2\vartheta = (ab)^2 a_x b_x = \Delta,
\]
d. h.: \glqq Die Discriminante der Form $f$ ist also bis auf einen Zahlenfactor gleich der Discriminante von $\Delta$.\grqq{}
Man hat:
\[
A_{\vartheta\vartheta} = \frac{1}{4}A_{\Delta\Delta},
\]
also wegen Relation (1)
\[
-\frac{1}{2}A_{\Delta\Delta} = R_{f_1f_2}.
\]
Verschwindet die Discriminante von $f$, dann hat also sowohl $f$ als $\Delta$ zwei gleiche Wurzeln. Wir werden sofort sehen, dass in diesem Falle $\Delta$ ein rationaler Factor von $f$ ist, und dass die Form $Q$ den
%CONT>>


% ===== pdf idx 189, printed 175.0 =====
% !PRINTED 175
\weberpage{175}
%<<CONT
Doppelfactor der beiden Formen $f$ und $\Delta$ dreifach enthält, also ein reiner Cubus ist.

152. \textit{Zweite Anwendung: Auflösung der cubischen Gleichung.} Die Methoden, eine Gleichung dritten Grades aufzulösen, werden naturgemäss zunächst darauf gerichtet sein, sie auf eine reine cubische Gleichung, deren Wurzeln uns ja bekannt sind, zurückzuführen. Hiezu bieten uns die vorausgegangenen Entwicklungen zwei Relationen, einmal die \gesperrt{Cayley'sche} Identität
\weq{1}{2Q^2 + \Delta^3 + Rf^2 = 0,}
dann aber auch die Relation
\weq{2}{\Delta_{Q+\lambda f} = \left(\frac{R}{2}+\lambda^2\right)\cdot\Delta.}
Wir gehen zunächst von der zweiten Gleichung aus. Sobald nämlich die \gesperrt{Hesse'sche} Form $\Delta_{Q+\lambda f}$ von $Q+\lambda f$ identisch verschwindet, ist nach Nr. 55 die Form $Q+\lambda f$ ein reiner Cubus. Dies tritt also immer ein, sobald
\weq{3}{\frac{R}{2}+\lambda^2 = 0,}
und damit haben wir eine quadratische Resolvente der cubischen Form. Setzen wir ihre beiden Wurzeln:
\[
+\sqrt{-\frac{R}{2}} = \lambda_1
\]
\[
-\sqrt{-\frac{R}{2}} = \lambda_2,
\]
so sind die beiden Formen:
\[
Q+\lambda_1 f, \quad Q+\lambda_2 f
\]
die dritten Potenzen zweier linearer Formen, die wir mit $\alpha_x$ resp. $\beta_x$ bezeichnen wollen. Für jene Werthe von $x$ aber, welche das Verschwinden von $f$ bewirken, besteht wegen
\weq{4}{\left.\begin{aligned} Q+\lambda_1 f &= \alpha_x^3\\ Q+\lambda_2 f &= \beta_x^3 \end{aligned}\right\}}
die Beziehung
\[
\alpha_x^3 = \beta_x^3,
\]
d. h.: \glqq Die drei Wurzeln von $f$ bestimmen sich beziehungsweise aus den drei Gleichungen
\[
\begin{aligned}
\alpha_x - \beta_x &= 0\\
\alpha_x - \varepsilon\beta_x &= 0\\
\alpha_x - \varepsilon^2\beta_x &= 0.
\end{aligned}
\]\grqq{}
Demnach lässt sich $f$ bis auf eine Constante als Product folgender drei linearer Factoren schreiben:
%CONT>>


% ===== pdf idx 190, printed 176.0 =====
% !PRINTED 176
\weberpage{176}
%<<CONT
\weq{5}{\begin{aligned}
f = c\cdot\Bigg\{&\sqrt[3]{Q+\sqrt{-\frac{R}{2}}\cdot f} - \sqrt[3]{Q-\sqrt{-\frac{R}{2}}\cdot f}\Bigg\}\times\\
\Bigg\{&\sqrt[3]{Q+\sqrt{-\frac{R}{2}}\cdot f} - \varepsilon\sqrt[3]{Q-\sqrt{-\frac{R}{2}}\cdot f}\Bigg\}\times\\
\Bigg\{&\sqrt[3]{Q+\sqrt{-\frac{R}{2}}\cdot f} - \varepsilon^2\sqrt[3]{Q-\sqrt{-\frac{R}{2}}\cdot f}\Bigg\}.
\end{aligned}}
Es ist hiebei nicht nöthig, direct die dritte Wurzel aus $Q+\sqrt{-\frac{R}{2}}\cdot f$ auszuziehen; vielmehr braucht man nur die erste Polare zu bilden
\[
\begin{aligned}
Q_x^2 Q_y + \lambda_1 f_x^2 f_y &= \alpha_x^2\alpha_y\\
Q_x^2 Q_y + \lambda_2 f_x^2 f_y &= \beta_x^2\beta_y
\end{aligned}
\]
und nun etwa $x_1 = 1,\ x_2 = 0$ zu setzen. Die sich so ergebenden Werthe $\alpha_y$ und $\beta_y$ sind bis auf Constante gleich den dritten Wurzeln aus dem vollständigen Cubus $Q+\lambda_i f$.

153. Wir können nun aber auch die \gesperrt{Cayley'sche} Relation zur Auflösung der Gleichung dritten Grades benutzen, welche uns zugleich auch Aufschluss über die Bedeutung der beiden linearen Formen $\alpha_x$ und $\beta_x$ giebt, auf die wir im Vorausgehenden geführt wurden. Wir schreiben zu dem Zwecke die Relation (I) Nr. 148 in der Form:
\[
\Delta^3 = -\{2Q^2 + Rf^2\},
\]
oder auch
\[
\Delta^3 = -2\left\{Q+f\sqrt{-\frac{R}{2}}\right\}\left\{Q-f\sqrt{-\frac{R}{2}}\right\}.
\]
Auf der linken Seite befindet sich nun ein reiner Cubus, und rechts das Product zweier Formen dritten Grades, welche, da im allgemeinen $f$ und $Q$ keinen gemeinschaftlichen Factor besitzen werden, ebenfalls keinen Theiler haben. Demnach muss jeder Factor rechts einzeln ein voller Cubus sein, etwa:
\[
\begin{aligned}
\frac{1}{2}\left\{Q+f\sqrt{-\frac{R}{2}}\right\} &= \alpha_x^3\\
\frac{1}{2}\left\{Q-f\sqrt{-\frac{R}{2}}\right\} &= \beta_x^3,
\end{aligned}
\]
wo nun $\alpha_x$ und $\beta_x$ bis auf den Factor $\frac{1}{2}$ mit den in Nr. 152 gleich bezeichneten linearen Factoren übereinstimmen. Wir haben demnach
\[
\Delta^3 = -8\alpha_x^3\beta_x^3,
\]
oder
\weq{1}{\Delta = -2\alpha_x\beta_x.}


% ===== pdf idx 191, printed 177.0 =====
% !PRINTED 177
\weberpage{177}

Hiebei ist rechts noch eine der dritten Einheitswurzeln als Factor zu denken.

Wenn wir nun die Form $f$ so transformiren, dass an Stelle der Variabeln $x_1$ und $x_2$ die Linearfactoren von $\Delta$ eintreten, so geht $f=0$ in eine reine cubische Gleichung über. Es ist nämlich nach Nr. 148:
\[
(f,\Delta)^2=a_x(a\alpha)(a\beta)=0,
\]
d. h.: \glqq Jedes symbolische Product mit dem Factor $(a\alpha)(a\beta)$ hat den Werth null.\grqq{} Erhebt man daher die Identität
\[
a_x(\alpha\beta)=\alpha_x(a\beta)-\beta_x(a\alpha)
\]
auf die dritte Potenz, so ergiebt sich:
\weq{2}{a_x^3(\alpha\beta)^3=\alpha_x^3(a\beta)^3-\beta_x^3(a\alpha)^3=c\cdot f.}

Die Form $f$ ist also unter Zugrundelegung der Linearfactoren von $\Delta$ in eine reine cubische Form transformirt, und man erkennt zugleich, dass es hiebei wegen der überall auftretenden dritten Potenzen gleichgiltig ist, welche dritte Einheitswurzel wir der rechten Seite der Gleichung (1) als Factor geben.

Zur Auflösung der cubischen Gleichung:
\[
f=ax_1^{\,3}+3bx_1x_2^{\,2}+3cx_1x_2^{\,2}+dx_2^{\,3}=0
\] % [sic] both middle terms print the exponent 2 on x_2 as shown
hat man demnach folgenden Weg einzuschlagen.

Man bilde zunächst die Hesse'sche Form von $f$
\[
\frac{1}{2}\Delta=(ac-b^2)x_1^{\,2}+(ad-bc)x_1x_2+(bd-c^2)x_2^{\,2}
\]
und suche die beiden Wurzeln derselben:
\[
\begin{matrix}\lambda_i\\ \mu_i\end{matrix}
=\frac{(bc-ad)\pm\sqrt{(ad-bc)^2-4(bd-c^2)(ac-b^2)}}{2(ac-b^2)}.
\]

Es ist alsdann:
\[
\frac{1}{2}\Delta=\frac{1}{\mu_1}(x_1\mu_1-x_2\lambda_1)(x_1\mu_2-x_2\lambda_2).
\]
Mit Hilfe der Transformation
\[
\begin{aligned}
y_1&=\mu_1x_1-\lambda_1x_2=\alpha_x\\
y_2&=\mu_2x_1-\lambda_2x_2=\beta_x
\end{aligned}
\]
geht $f$ in die Form über
\[
F=Ay_1^{\,3}+Dy_2^{\,3}=0,
\]
die nun durch Ausziehen der dritten Wurzel gelöst werden kann.

154. \textit{Schlussbemerkungen.} Sind die Coefficienten von $f$ reell und ist $R=A_{\Delta\Delta}$ negativ, so besitzt $f$ eine reelle und zwei imaginäre Wurzeln. Dies erkennt man direct aus der Darstellung (5) der Form $f$ in Nr. 152.


% ===== pdf idx 192, printed 178.0 =====
% !PRINTED 178
\weberpage{178}

Ist $R=0$, so besitzt $f$ zwei gleiche Wurzeln und ebenso $\Delta$. In diesem Falle ist natürlich die in Nr. 153 gegebene Transformation von $f$ unmöglich.$^{\ast}$ % ?? uncertain: raised mark after \textquotedbl{}unmöglich.\textquotedbl{} -- no footnote text visible in the provided bands for this page
Die Hesse'sche Form ist dann, wie schon früher erwähnt, Factor von $f$. Denn da in jedem Falle nach Nr. 152 (4)
\[
\begin{aligned}
Q+\sqrt{-\frac{R}{2}}\,f&=\alpha_x^3\\
Q-\sqrt{-\frac{R}{2}}\,f&=\beta_x^3,
\end{aligned}
\]
so ist demnach:
\[
\begin{aligned}
2f\cdot\sqrt{-\frac{R}{2}}&=\alpha_x^3-\beta_x^3=(\alpha_x-\beta_x)\cdot(\alpha_x-\varepsilon\beta_x)\cdot(\alpha_x-\varepsilon^2\beta_x)\\
2Q&=\alpha_x^3+\beta_x^3\\
\Delta&=\varrho\cdot\alpha_x\beta_x,
\end{aligned}
\]
wo $\varrho$ eine Constante.

Ist nun $R=0$, also $\alpha_x=\beta_x$, so wird
\[
\begin{aligned}
Q&=\alpha_x^3\\
\Delta&=\varrho\alpha_x^2.
\end{aligned}
\]
Es ist aber
\[
(f,\Delta)=Q,
\]
oder:
\[
\varrho(a\alpha)a_x^2\alpha_x=\alpha_x^3,
\]
also nach Division mit dem linearen Factor $\alpha_x$:
\[
\varrho(a\alpha)a_x^2=\alpha_x^2.
\]
Daraus schliesst man, dass $\alpha_x$ der Doppelfactor von $f$. [Vergl. Nr. 99, (3).] Ueberdies sieht man, dass, im Falle $R=0$, $Q$ die dritte Potenz des Doppelfactors von $f$ ist.

Ist endlich $\Delta$ selbst identisch null, so ist nach Nr. 55 die Form $f$ der Cubus einer linearen Form.

\begin{center}
\S\ 15. \textbf{Das Formensystem der Form vierten Grades.}
\end{center}

155. \textit{Die fundamentalen In- und Covarianten.} Wie bei den Formen dritten Grades, so liefern auch hier die Ueberschiebungen von $f$ über sich selbst und über ihre Hesse'sche Covariante das vollständige System der Form vierten Grades, d. h. die Gesammtheit jener Formen, durch welche sich alle übrigen Co- und Invarianten ganz und rational ausdrücken lassen. Wir haben diese Ueberschiebungen $(f,f)^\lambda$, $(f,\Delta)^\lambda$ bereits früher, insbesondere auch in § 8 unter den Anwendungen des Processes der Reihenentwicklung, kennen gelernt, und wollen nur noch einmal die dort gewonnenen Resultate zusammenfassen.


% ===== pdf idx 193, printed 179.0 =====
% !PRINTED 179
\weberpage{179}

Bezeichnen wir mit
\[
\begin{aligned}
f&=\bar a_0x_1^{\,4}+4\bar a_1x_1^{\,3}x_2+6\bar a_2x_1^{\,2}x_2^{\,2}+4\bar a_3x_1x_2^{\,3}+\bar a_4x_2^{\,4}\\
&=a_x^4=b_x^4=c_x^4=\dots\ \text{etc.}
\end{aligned}
\]
die gegebene Form vierten Grades, so sind die erwähnten Formen durch folgende symbolische Producte dargestellt:
\[
\begin{aligned}
(f,f)^2&=(ab)^2a_x^2b_x^2=\Delta\\
(f,f)^4&=(ab)^4\phantom{a_x^2b_x^2}=i\ {}^{\ast})\\
(f,\Delta)&=(a\Delta)a_x^3\Delta_x^3=(ab)^2(cb)c_x^3a_x^2b_x=t\\
(f,\Delta)^4&=(a\Delta)^4\phantom{a_x^2b_x^2}=(ab)^2(bc)^2(ca)^2=j\ {}^{\ast}).
\end{aligned}
\]
Diese fünf Formen:
\[
f,\ \Delta,\ t,\ i,\ j
\]
bilden, wie schon erwähnt, das vollständige System der Form $f$, und mit dem Beweise dieser Behauptung wollen wir uns nun im Folgenden beschäftigen.

156. \textit{Gedankengang des Beweises.} Irgend eine Co- oder Invariante von $f$ kann nur durch Ueberschiebung der Formen $f$, $\Delta$, $t$ über sich selbst oder über einander entstehen. Legen wir hiebei $t$ in der Form $(a\Delta)a_x^3\Delta_x^3$ zu Grunde, so enthält jedes symbolische Product, das bei diesem Ueberschiebungsprocess entsteht, nur Symbole $a$ und $\Delta$, und zwar ausser in Factoren erster Art $a_x$, $\Delta_x$ nur in den drei allein möglichen Klammerfactoren:
\[
(ab),\ (a\Delta),\ (\Delta\Delta_1).
\]
Gelingt es uns also zu zeigen, dass alle Formen, welche einen dieser drei Klammerfactoren enthalten, reducibel sind, vergl. Nr. 124, so ist damit auch bewiesen, dass ausser den erwähnten fünf Formen keine weitern existiren können, die sich nicht rational und ganz durch die gegebenen $f$, $\Delta$, $t$, $i$, $j$ ausdrücken lassen.

Wir werden demnach zuerst die zunächst liegenden einfachen Ueberschiebungen von $f$ über $\Delta$, und $\Delta$ über $\Delta$ berechnen, wobei

\par\noindent\rule{4cm}{0.4pt}\par
\noindent $^{\ast})$ Es ist:
\[
i=(ab)^4=a_1^4b_2^4-4a_1^3a_2\cdot b_2^3b_1+6a_1^2a_2^2b_1^2b_2^2-\cdots\ \text{etc.}=2(\bar a_0\bar a_4-4\bar a_1\bar a_3+3\bar a_2^2),
\]
\[
j=(ab)^2(bc)^2(ca)^2=\begin{vmatrix}a_1^2 & a_1a_2 & a_2^2\\ b_1^2 & b_1b_2 & b_2^2\\ c_1^2 & c_1c_2 & c_2^2\end{vmatrix}^2=6\begin{vmatrix}\bar a_0 & \bar a_1 & \bar a_2\\ \bar a_1 & \bar a_2 & \bar a_3\\ \bar a_2 & \bar a_3 & \bar a_4\end{vmatrix}
\]
\[
=6(\bar a_0\bar a_2\bar a_4-\bar a_0\bar a_3^2-\bar a_1^2\bar a_4-\bar a_2^3+2\bar a_1\bar a_2\bar a_3)\ \text{(vergl. Bd. I, Nr. 40).}
\]


% ===== pdf idx 194, printed 180.0 =====
% !PRINTED 180
\weberpage{180}
%<<CONT
sich als erstes Resultat ergiebt, dass sich die so entstehenden Formen in der That durch die fünf Grundformen rational und ganz ausdrücken lassen.

Sodann zeigen wir auf Grund dieser Thatsache, dass die Klammerfactoren
\[
(ab),\ (a\Delta)^2,\ (a\Delta)(b\Delta),\ (\Delta\Delta_1)
\]
Reducenten sind, und daraus ergiebt sich alsdann, dass jedes durch Ueberschiebung von $f$, $\Delta$ und $t$ entstehende Product reducibel ist.

Die obenerwähnten Klammerfactoren $(ab)$ und $(\Delta\Delta_1)$ sind unter diesen Reducenten, und was symbolische Producte mit dem Factor $(a\Delta)$ betrifft, so haben sie entweder diesen ganz allein und dann nothwendig auch noch die Factoren $a_x^3\cdot\Delta_x^3$, d. h. sie besitzen $t$ als Factor; oder sie haben als nächst einfachsten Fall die Factoren $(a\Delta)^2$, resp. $(a\Delta)(b\Delta)$, und von diesen wollen wir ja gerade beweisen, dass sie Reducenteneigenschaft besitzen.

157. \textit{Die Ueberschiebungen $(f,\Delta)^\lambda$, $(\Delta,\Delta)^\lambda$.} Den ersten Theil des Beweises haben wir bereits in früheren Untersuchungen zum grössten Theil erledigt, oder können ihn wenigstens auf Grund der dort gewonnenen Resultate rasch zum Ziele führen.

Die Ueberschiebungen $(f,\Delta)$, $(f,\Delta)^4$ haben wir in das System aufgenommen, und $(\Delta,\Delta)$, $(\Delta,\Delta)^3$ verschwinden identisch; die Formen $(f,\Delta)^2$, $(f,\Delta)^3$ wurden in § 8 berechnet als erste Anwendung des Processes der Reihenentwicklung, wobei wir die Werthe
\[
\begin{aligned}
(f,\Delta)^2&=\frac{i}{6}\cdot f\\
(f,\Delta)^3&=0
\end{aligned}
\]
erhielten. Es erübrigt also nur noch die Berechnung von
\[
(\Delta,\Delta)^2\ \text{und}\ (\Delta,\Delta)^4.
\]

158. \textit{Berechnung von $(\Delta,\Delta)^2$.} Die Ueberschiebung $(\Delta,\Delta)^2$ geht aus der zweiten Polare $\Delta_{y^2}$ hervor, wenn wir darin $y$ durch $\Delta_1$ ersetzen, und mit $\Delta_{1x}^2$ multipliciren. Diese zweite Polare ist aber das erste Glied in der Reihenentwicklung der Form
\[
(ab)a_x^3b_y^3=\frac{3}{2}\Delta_{y^2}(xy)+\frac{1}{4}i(xy)^3.
\]
Daher ist:
\weq{1}{(ab)(b\Delta)^3a_x^3\Delta_x=\frac{3}{2}(\Delta,\Delta)^2+\frac{1}{4}i\cdot\Delta.}

Hiebei entsteht das symbolische Product links, wenn wir in $(b\Delta)^3b_y\Delta_x$, $y$ durch $a$ ersetzen und mit $a_x^3$ multipliciren. Entwickeln
%CONT>>


% ===== pdf idx 195, printed 181.0 =====
% !PRINTED 181
\weberpage{181}

wir aber diese Form nach Reihe (VIII) \S.\,7, so kommt:

\weq{1}{(b\Delta)^3 b_y \Delta_x = [(b\Delta)^3 b_x \Delta_x]_y + \frac{1}{2}(b\Delta)^4(yx),}

oder weil:
\[
\begin{gathered}
(b\Delta)^3 b_x \Delta_x = (f,\Delta)^3 = 0 \\
(b\Delta)^4 = (f,\Delta)^4 = j,
\end{gathered}
\]
so ergiebt sich nach Ausführung der erwähnten Substitution

\weq{2}{(ab)(b\Delta)^3 \Delta_x a_x^3 = \frac{1}{2} j\cdot f,}

und demnach durch Comparation von (1) und (2)

\weq{3}{(\Delta,\Delta)^2 = \frac{1}{3}\cdot j\cdot f - \frac{i}{6}\Delta.}

159. \textit{Berechnung von $(\Delta,\Delta)^4$.} Der Werth dieser Invariante wird am einfachsten erhalten, wenn wir in der Ueberschiebung
\[
(f,\Delta)^2 = \frac{i}{6}\cdot a_x^4 = (a\Delta)^2 a_x^2 \Delta_x^2
\]
$x$ durch $b$, und gleichzeitig in der Identität
\[
\Delta_x^4 = (ab)^2 a_x^2 b_x^2
\]
$x$ durch $\Delta$ ersetzen. Wir erhalten alsdann:
\[
\begin{gathered}
\frac{i}{6}\cdot i = (a\Delta)^2(ab)^2(\Delta b)^2 \\
(\Delta,\Delta)^4 = (ab)^2(a\Delta)^2(b\Delta)^2,
\end{gathered}
\]
und demnach durch Comparation:

\weq{4}{(\Delta,\Delta)^4 = \frac{i^2}{6}.}

Hiemit ist der erste Theil unserer Aufgabe erledigt. Es zeigte sich, dass in der That die Ueberschiebungen $(f,\Delta)^\lambda$ und $(\Delta,\Delta)^\lambda$ keine neuen Formen erzeugen. Im zweiten Theile haben wir nunmehr nachzuweisen, dass
\[
(ab),\ (a\Delta)^2,\ (\Delta\Delta_1),\ (a\Delta)(b\Delta)
\]
Reducenten sind.

160. \textit{Die Reducenten.} Auf Grund der nunmehr berechneten Ueberschiebungswerthe ist nicht schwer nachzuweisen, dass zunächst $(ab)$, $(a\Delta)^2$, $(\Delta\Delta_1)$ Reducenteneigenschaft besitzen. Denn entwickelt man die entsprechenden einfachsten symbolischen Producte
\[
(ab)a_x^3 b_y^3,\ (a\Delta)^2 a_x^2 \Delta_y^2,\ (\Delta\Delta_1)\Delta_x^3\Delta_{1y}^3
\]
in die Reihe (VIII) \S.\,7, so findet man, dass die rechts auftretenden Elementarcovarianten sämmtlich Formen des Systems sind. Man erhält nämlich:
%CONT>>


% ===== pdf idx 196, printed 182.0 =====
% !PRINTED 182
\weberpage{182}

%<<CONT
\weq{1}{(ab)a_x^3 b_y^3 = \frac{3}{2}\Delta_{y^2}(xy) + \frac{1}{4} i(xy)^3}

\weq{2}{(a\Delta)^2 a_x^2 \Delta_y^2 = \frac{i}{6} f_{y^2} + \frac{1}{3} j\cdot(xy)^2}

\weq{3}{(\Delta\Delta_1)\Delta_x^3\Delta_{1y}^3 = \frac{1}{2}\left(j\cdot f_{y^2} - \frac{i}{2}\Delta_{y^2}\right)(xy) + \frac{1}{24} i^2 (xy)^3.}

Nun kann jedes symbolische Product, das $(ab)$, oder $(a\Delta)^2$, oder $(\Delta\Delta_1)$ zum Factor hat, durch Ueberschiebung mit den Formen (1), resp. (2), (3) oder deren Polaren gewonnen werden (vergl. \S.\,3) und die Gleichungen (1), (2), (3) lehren daher:

1) Alle Formen, welche $(ab)$ zum Factor haben, lassen sich reduciren auf Formen, welche an Stelle der beiden Symbole $a$ und $b$ nur mehr das eine $\Delta$ besitzen und auf Formen, die $i$ zum Factor haben.

2) Alle Formen, welche $(a\Delta)^2$, $(\Delta\Delta_1)$ besitzen, lassen sich auf Formen reduciren, welche $i$ oder $j$ zum Factor haben,

und daher sind die erwähnten Klammerfactoren in der That Reducenten. (Vergl. Nr. 124.)

161. Um nun die gleiche Eigenschaft auch für $(a\Delta)(b\Delta)$ nachzuweisen, denken wir uns wieder das einfachste symbolische Product, das $(a\Delta)(b\Delta)$ zu Klammerfactoren besitzt, nämlich
\[
P = (a\Delta)(b\Delta)a_x^3 b_y^3 \Delta_z^2
\]
in eine nach Potenzen von $(xy)$ fortschreitende Reihe entwickelt. Dann enthalten alle Elementarcovarianten, abgesehen von der ersten, den Factor $(ab)$ und sind demnach reducibel. Nun ist aber nach dem Productsatz diese erste Elementarcovariante:
\[
(a\Delta)(b\Delta)a_x^3 b_x^3 = \frac{1}{2} a_x^2 b_x^2 \{(b\Delta)^2 a_x^2 + (a\Delta)^2 b_x^2 - \Delta_x^2(ab)^2\},
\]
und somit, weil rechts nur Reducenten sich befinden, gleichfalls reducibel. Es ist also auch $P$ reducibel und $(a\Delta)(b\Delta)$ Reducent. Damit sind wir aber am Schlusse unseres Beweises angelangt. Jedes symbolische Product, das eine Co- oder Invariante der biquadratischen Form darstellen soll, kann zunächst reducirt werden auf eine Summe von Producten, die $i$ oder $j$ zum Factor haben, oder wenigstens auf Producte, mit weniger Symbolen $a$ oder $\Delta$. Durch Wiederholung der Reduction wird schliesslich jedes symbolische Product reducirt auf eine rationale ganze Function der Formen $f$, $\Delta$, $t$, $i$, $j$.

162. \textit{Die Fundamentalrelation zwischen den fünf Formen des Systemes.} Zwischen den fünf Formen des Systemes besteht eine Relation, die
%CONT>>


% ===== pdf idx 197, printed 183.0 =====
% !PRINTED 183
\weberpage{183}

%<<CONT
wir sofort nach allgemeinen Gesetzen berechnen können. Die Covariante $t$ sechsten Grades ist nämlich Functionaldeterminante von $f$ und $\Delta$, und demnach ist ihr Quadrat durch die geraden Formen des Systemes darstellbar. Man erhält nach Nr. 48, (I):
\[
2t^2 = \left|
\begin{array}{ccc}
\Delta, & \dfrac{i}{6} f, & f \\[4pt]
\dfrac{i}{6} f, & \dfrac{j}{6} f - \dfrac{i}{6}\Delta, & \Delta \\[4pt]
f, & \Delta, & 0
\end{array}
\right|
\]
oder:
\[
2t^2 = \frac{i}{2} f^2 \cdot \Delta - \Delta^3 - \frac{j}{3} f^3.
\]

163. \textit{Relationen zur Berechnung der Ueberschiebung von $t$ über sich selbst und über $f$ und $\Delta$.} Aus den letzten Untersuchungen geht hervor, dass die Ueberschiebungen von $t$ über sich selbst und über $f$ und $\Delta$ zu keinen neuen Covarianten Veranlassung geben können. Wenn wir im Folgenden sie dennoch berechnen, so geschieht dies einmal, weil sich uns hier ein lehrreiches Beispiel für die Behandlung derartiger Aufgaben darbietet, dann aber auch, weil gerade die Ueberschiebungen von $t$ über sich selbst interessante Aufschlüsse über den Charakter dieser Covariante gewähren. Am besten wählt man zur Bildung der durch diese Ueberschiebungen entstehenden Formen als Ausgangspunkt die Form:
\[
P = a_y^4 \Delta_x^4 - a_x^4 \Delta_y^4 = p_y^4 q_x^4,
\]
welche wir schon früher als Combinante charakterisirt haben. Entwickelt man dieselbe nach Potenzen von $(yx)$, so erhält man

\weq{1}{P = \sum \frac{\binom{4}{\varrho}\binom{4}{\varrho}}{\binom{8-\varrho+1}{\varrho}} \cdot \overline{(pq)}^{\,\varrho}_{y^{4-\varrho}} (yx)^\varrho.}

Von den Elementarcovarianten dieser Form $P$ verschwinden aber alle bis auf die zweite:
\[
(\overline{pq})^1 = [(a\Delta)a_x^3 \Delta_x^3 - (\Delta a)a_x^3\Delta_x^3] = 2(a\Delta)a_x^3\Delta_x^3 = 2t_x^6.
\]

Der Coefficient dieses Gliedes der Reihe (1) ist $\dfrac{\binom{4}{1}\binom{4}{1}}{\binom{8}{1}} = 2$; also erhalten wir:

\weq{I}{a_y^4\Delta_x^4 - a_x^4\Delta_y^4 = 4t_x^3 t_y^3(yx);}

daraus ergeben sich durch Polarisation nach $y$ die zwei weitern Formeln:

\weq{II}{a_y^3 a_x \Delta_x^4 - a_x^4 \Delta_y^3\Delta_x = 3t_x^4 t_y^2(yx)}

\weq{III}{a_y^2 a_x^2\Delta_x^4 - a_x^4\Delta_y^2\Delta_x^2 = 2t_x^5 t_y(yx),}
%CONT>>


% ===== pdf idx 198, printed 184.0 =====
% !PRINTED 184
\weberpage{184}

%<<CONT
wobei also:

\weq{IV}{(a\Delta)a_x^3\Delta_x^3 = t_x^6.}

Auf den rechten Seiten dieser vier Formeln befinden sich Polaren von $t$, von der nullten angefangen bis zur dritten. Je nachdem wir also für $y$ die Symbole $a$, $\Delta$, $t$ in dieselben eintragen, erhalten wir der Reihe nach die gewünschten Ueberschiebungen, ausgedrückt in Symbolen von $a$ und $\Delta$ allein.

164. \textit{Die ersten Ueberschiebungen.} Hiezu benutzen wir Formel III, wir setzen darin zunächst $y_1 = b_2$, $y_2 = -b_1$ und multipliciren mit $b_x^2$; dann erhält man:

\weq{1}{(ab)_x^2 a_x^2 b_x^2 \cdot \Delta - f\cdot(b\Delta)^2 b_x^2\Delta_x^2 = 2(t,f) = \Delta^2 - \frac{i}{6}\cdot f^2,}

sodann ersetzen wir $y_1$ durch $\Delta'$, $y_2$ durch $\Delta_2'$ und bekommen:

\weq{2}{(a\Delta')^2 a_x^2 \Delta_x'^2 \cdot \Delta - f\cdot(\Delta\Delta')^2\Delta_x^2\Delta_x'^2 = 2(t,\Delta) = \frac{1}{3} f(i\Delta - jf).}

Hiezu tritt als erste Ueberschiebung von $t$ über sich selbst:

\weq{3}{(t,t) = 0.}

165. \textit{Die zweiten Ueberschiebungen.} Wir ersetzen in Formel (II) zuerst $y$ durch $b$, dann $y$ durch $\Delta'$; beide Male erhalten wir als Resultat null; daher

\weq{4}{(t,f)^2 = 0}

\weq{5}{(t,\Delta)^2 = 0.}

Substituiren wir aber für $y$ das Symbol $t'$, so entsteht:
\[
(at)^3 a_x t_x^3 \cdot \Delta - f\cdot(\Delta t)^3 \Delta_x t_x^3 = 3(t,t')^2 t_x^4 t_x'^4,
\]
oder
\[
(f,t)^3 \cdot \Delta - (\Delta,t)^3 f = (t,t')^2.
\]
Führen wir für die sofort zu berechnenden Ueberschiebungen $(f,t)^3$ und $(\Delta,t)^3$ ihre aus (7) und (8) hervorgehenden Ausdrücke ein, so kommt:
\[
\frac{1}{4}\{jf - i\Delta\}\Delta - \frac{1}{4}\left\{\frac{i^2}{6} f - j\Delta\right\}f = 3(t,t')^2,
\]
oder

\weq{6}{-\frac{1}{12}\left\{i\Delta^2 - 2j\Delta f + \frac{i^2}{6} f^2\right\} = (t,t')^2.}

166. \textit{Die dritten Ueberschiebungen.} Sie gehen direct aus der Combinante $P$ selbst hervor [Gleichung (I)], indem man der Reihe nach $y$ durch $b$, $\Delta'$, $t'$ ersetzt. Man erhält:

\weq{7}{4(t,f)^3 = i\Delta - jf}

\weq{8}{4(t,\Delta)^3 = j\Delta - \frac{i^2}{6} f}

\weq{9}{4(t,t')^3 = 0.}


% ===== pdf idx 199, printed 185.0 =====
% !PRINTED 185
\weberpage{185}

167. \textit{Die vierten und fünften Ueberschiebungen.} Wir benutzen zu diesen Bildungen eine andere Ausgangsform.

Gemäss der Reihenentwicklung (IX), \S.\,7 ist:

\[
(a\Delta)\,\Delta_x^2\,a_y^3\,\Delta_y = \sum \frac{\binom{2}{k}\binom{3}{k}}{\binom{6-k+1}{k}}\,(f,\Delta)^k_{y^{3-k+1}}\,(xy)^k,
\]
oder
\[
(a\Delta)\,a_y^3\,\Delta_y\,\Delta_x^2 = t_{y^4} + \frac{i}{6}\cdot f_{y^3}(yx).
\]

Ersetzen wir hier $y$ durch $b$, so folgt:

\[
(a\Delta)(ab)^3(\Delta b)\Delta_x^2 = (t,f)^4 + \frac{i}{6}(f,f)^3.
\]

Nun ist aber $(f,f)^3=0$, und ebenso verschwindet die linke Seite, da sie durch Vertauschung von $a$ mit $b$ ihr Zeichen ändert, daher:

\weq{10}{(t,f)^4 = 0,}
und ebenso
\weq{11}{(t,\Delta)^4 = 0,}

wie aus (10) direct durch den Aronhold'schen Process (vergl. auch \S.\,16) hervorgeht. Beide Resultate sind auch a priori einleuchtend. Denn diese Ueberschiebungen würden quadratische Covarianten erzeugen, die aber unmöglich durch $f, \Delta, t$ sich rational ausdrücken lassen.

Die vierte Ueberschiebung von $t$ über sich selbst gewinnt man wiederum aus der Formel (II), Nr. 163, wenn man darin zuerst $x$ mit $y$ vertauscht, also aus:

\[
a_y^4\,\Delta_x^3\,\Delta_y - a_x^3\,a_y\,\Delta_y^4 = 3\,t_y^4\,t_x^2\,(yx).
\]

Man substituirt $t$ für $y$ und erhält:

\[
(at)^4(\Delta t)\,\Delta_x^3\,t_x - a_x^3(at)(\Delta t)^4 t_x = 3\,(t,t')^4\,t_x^2\,t_x'^{\,2}.
\]

Die beiden Glieder links verschwinden einzeln wegen der Symbolfactoren $(at)^4$, resp. $(\Delta t)^4$, vermöge welcher dieselben durch Ueberschiebung mit $(f,t)^4=0$, resp. $(\Delta t)^4=0$ entstanden gedacht werden können.

Somit habeu wir die wichtige Relation % [sic]

\weq{12}{(t,t)^4 = 0.}

Durch sie wird $t$ als eine ganz specielle Form ihres Grades charakterisirt. Wir werden in \S.\,19 solche Formen eingehender studiren. Hier möge nur so viel bemerkt sein, dass gerade wegen dieser Relation $(t,t)^4=0$
%CONT>>


% ===== pdf idx 200, printed 186.0 =====
% !PRINTED 186
\weberpage{186}

%<<CONT
nach den in \S.\,20 entwickelten Gesetzen das ganze Formensystem einer solchen Form $t$ sich auf die vier Formen reducirt

\[
t,\ (t,t)^2,\ ((t,t)^2,t),\ (t,t)^6.
\]

Insbesondere müssen daher alle Invarianten Potenzen von $(t,t)^6$ sein, welche, wie der nächste Absatz zeigt, den Werth $\frac{1}{4}\left(j^2-\frac{i^3}{6}\right)$ besitzt.

168. \textit{Die fünften und sechsten Ueberschiebungen.} Die fünften Ueberschiebungen sind durchwegs null, da ja $\Delta$ und $f$ nur vom vierten Grade sind und $(t,t)^5$ sein Zeichen ändert bei Vertauschung von $t$ mit $t'$.

Es bleibt sonach nur noch $(t,t)^6$ zu berechnen. Wir bilden von

\[
t_x^6 = (a\Delta)\,a_x^3\,\Delta_x^3
\]

die sechste Polare und ersetzen $y$ durch $t$:

\[
(t,t)^6 = (a\Delta)(at)^3(\Delta t)^3 = -(at)^3(a\Delta)(t\Delta)^3.
\]

Der Ausdruck rechts kann auch betrachtet werden als vierte Ueberschiebung von $\Delta$ über $(f,t)^3=(at)^3a_x t_x^3$; daher hat man:

\[
(t,t)^6 = -((f,t)^3,\Delta)^4 = -\frac{1}{4}\big((i\Delta-jf),\Delta\big)^4 = \frac{1}{4}\left(j^2-\frac{i^3}{6}\right).
\]

Wir werden alsbald sehen, dass die berechnete Grösse rechts nichts anderes ist als die Discriminante $R$ der Form $f=a_x^4$, abgesehen von einem Zahlenfactor.

\begin{center}
\S.\,16.

\gesperrt{Anwendung des Aronhold'schen Processes auf das Formensystem der Form vierten Grades.}
\end{center}

169. \textit{Wirkung des Processes auf die Grundformen des Systemes.} Wie die cubische Form eine Covariante $Q$ gleichen Grades besitzt, so haben wir auch im System der biquadratischen Form eine Covariante, die Hesse'sche Form $\Delta$, die gleichen Grades ist mit der Originalform $f$. Wir können daher in derselben Weise wie dort unter Einführung der Symbole $\Delta$, symbolische Producte, welche Covarianten von $f=a_x^4$ repräsentiren, dem Aronhold'schen Processe unterwerfen, und so wichtige Aufschlüsse über die Eigenschaften und den inneren Zusammenhang der Formen des Systems erhalten.

Ist $P$ irgend eine Covariante von $f=a_x^4$, so ist in unserm Falle der Aronhold'sche Process dargestellt durch:

\weq{I}{\delta P = \sum \frac{\partial P}{\partial \bar a_i}\,\bar\Delta_i,}
%CONT>>


% ===== pdf idx 201, printed 187.0 =====
% !PRINTED 187
\weberpage{187}

%<<CONT
wobei $\bar a_i$ und $\bar\Delta_i$ die unsymbolischen Co\"efficienten von $f$ resp. $\Delta$ bedeuten. Ersetzen wir nun in dieser Gleichung (I) die Covariante $P$ der Reihe nach durch $f, \Delta, i, t, j$, so erhalten wir:

\weq{1}{\delta f = \Delta;}

\weq{2}{\delta \Delta = \delta(f,f)^2 = 2(f,\delta f)^2 = 2(f,\Delta)^2 = \frac{i}{3}\cdot f.}

Wir haben also hier wiederum die wechselseitige Beziehung, die wir schon mehrmals bei analogen Untersuchungen betont haben. Ebenso wie $\delta f$ auf $\Delta$ führt, so gelangt man umgekehrt durch $\delta\Delta$ wieder auf $f$ zurück.

Es ist weiter:

\weq{3}{\delta i = \delta(f,f)^4 = 2(f,\delta f)^4 = 2(f,\Delta)^4 = 2j;}

\weq{4}{\delta t = \delta(f,\Delta) = (\delta f,\Delta)+(f,\delta\Delta) = (\Delta,\Delta)+i\cdot(f,f) = 0.}

Das letzte Resultat war a priori einzusehen, da $t$ als erste Ueberschiebung von $f$ und $\Delta$ eine Combinante dieser beiden Formen ist. (Vergl. \S.\,6.) Die allgemeinste Combinante des Systemes $f, \Delta$ ist die Form

\[
C = a_y^4\Delta_x^4 - a_x^4\Delta_y^4,
\]

wie wir bei Gelegenheit der Definition von Combinanten \S.\,6 erwähnt haben.

Da aber nach Gleichung Nr. 163 (I), $t$ die einzige Elementar-
%CONT>>


% ===== pdf idx 202, printed 188.0 =====
% !PRINTED 188
\weberpage{188}

%<<CONT
covariante von $C$ ist, so treten beide Combinanten $t$ wie $C$ gleichberechtigt im System auf. Endlich ist:

\weq{5}{\delta j = \delta(f,\Delta)^4 = (\delta f,\Delta)^4+(f,\delta\Delta)^4 = (\Delta,\Delta)^4+\frac{i}{3}(f,f)^4 = \frac{i^2}{6}+\frac{i^2}{3} = \frac{i^2}{2}.}

Damit ist an den Grundformen des Systemes die Wirkung des $\delta$-Processes studirt. Man kann die eben gewonnenen Resultate wiederum verwenden, aus einfacheren Ueberschiebungen schwierigere zu berechnen. So erhalten wir z.\,B. den Werth von $(\Delta,\Delta)^4$, wenn wir

\[
(f,\Delta)^2 = \frac{i}{6}f
\]

dem $\delta$-Process unterwerfen. Es ergiebt sich:

\[
(\Delta,\Delta)^2 + \frac{i}{3}(f,f)^2 = \frac{j}{3}f + \frac{i}{6}\Delta,
\]
also
\[
(\Delta,\Delta)^2 = \frac{1}{3}jf - \frac{1}{6}i\Delta,
\]

was mit dem früher gewonnenen Resultat, Nr. 158, (3), übereinstimmt.

170. \textit{Die Combinanten der Form $\varphi=kf+\lambda\Delta$.} Wir haben seinerzeit eine Form $C$ als Combinante zweier binärer Formen $\chi$ und $\psi$ definirt, sobald sie der Differentialgleichung

\[
\delta C = \sum \frac{\partial C}{\partial \bar\chi_i}\cdot \psi_i = 0
\]

genügt. Ist nun $P$ irgend eine Combinante von $f=a_x^4$ und $\Delta=\Delta_x^4$, also, wenn wir uns in $P$ die Co\"efficienten von $\Delta$ durch ihre Werthe in den Co\"efficienten von $f$ ersetzt denken, eine Combinante von $f$, so besitzt die Form

\[
\varphi = kf + \lambda\Delta
\]

eine entsprechende Combinante $P_\varphi$, welche man erhält, wenn man in $P$ die Co\"efficienten $\bar a_i$ von $f$ durch die Co\"efficienten $k\bar a_i+\lambda\bar\Delta_i$ von $\varphi$ ersetzt. Insbesondere entspricht der einzigen Fundamentalcombinante $t=(f,\Delta)$ eine Form $t_\varphi$ von $\varphi$, aus der alle Combinanten dieser Form $\varphi$ entspringen. Die Berechnung derselben ist also zurückgeführt auf die Berechnung von $t_\varphi$. Nach den allgemeinen Entwicklungen in \S.\,6 hat man aber:

\weq{II}{t_\varphi = g^\varrho\cdot t.}

Hiebei ist $\varrho$ eine noch zu bestimmende ganze Zahl; ferner ist

\[
g = kB - \lambda A,
\]

worin $A$ und $B$ die Co\"efficienten in der Gleichung

\[
\Delta_{kf+\lambda\Delta} = Af+B\Delta,
\]

sind. Es handelt sich also zunächst darum, $\Delta_{kf+\lambda\Delta}=\Delta_\varphi$ zu berechnen.

171. \textit{Berechnung von $\Delta_\varphi$.} Wir ermitteln diese Form direct durch Ueberschiebung. Man erhält:

\[
\begin{aligned}
\Delta_\varphi &= (kf+\lambda\Delta,\,kf+\lambda\Delta)^2\\
&= k^2\cdot(f,f)^2 + 2k\lambda\cdot(f,\Delta)^2 + \lambda^2\cdot(\Delta,\Delta)^2\\
&= k^2\cdot\Delta + \frac{i}{3}f\cdot k\lambda + \lambda^2\left(\frac{j}{3}f-\frac{i}{6}\Delta\right),
\end{aligned}
\]
also:
\[
\Delta_\varphi = \Delta\left\{k^2-\frac{i}{6}\lambda^2\right\} + f\left\{\frac{k\lambda i}{3}+\frac{j}{3}\lambda^2\right\}.
\]

Es ist sehr bemerkenswerth, dass die Coefficienten von $\Delta$ und $f$ auf der rechten Seite dieser Gleichung nichts anderes sind, als die
%CONT>>


% ===== pdf idx 203, printed 189.0 =====
% !PRINTED 189
\weberpage{189}
%<<CONT
ersten Differentialquotienten der Form
\weq{III}{g = k^3 - \frac{i}{2}\,k\lambda^2 - \frac{j}{3}\,\lambda^3}
nach $k$ beziehungsweise $\lambda$. Die Differentiation ergiebt nämlich unmittelbar
\[
\frac{1}{3}\,\frac{\partial g}{\partial k} = g_1 = k^2 - \frac{i}{6}\lambda^2
\]
\[
\frac{1}{3}\,\frac{\partial g}{\partial \lambda} = g_2 = -\frac{k\lambda i}{3} - \frac{j}{3}\lambda^2 .
\]
Die Form $g$ selbst aber ist eine cubische Resolvente der Gleichung
vierten Grades $f = a_x^4 = 0$, auf welche die Lagrange'sche Methode der
Auflösung dieser Gleichung direct führt*). Wir können also die Form
$\varDelta_\varphi$ einfacher darstellen durch
\[
\varDelta_\varphi = g_1\varDelta - g_2 f = B\varDelta + Af,
\]
und die Combinante $t_\varphi$ durch:
\[
t_\varphi = \left(k^3 - \frac{i}{2}\,k\lambda^2 - \frac{j}{3}\,\lambda^3\right)\cdot t .
\]
Denn die Constante $\varrho$ muss nothwendig den Werth $1$ besitzen,
da $t$ nur vom dritten Grade in den Coefficienten von $f$ und demnach
$t_\varphi$ auch nur vom dritten Grade in $k$ und $\lambda$ sein kann.

172. \textit{Berechnung der Formen $i_\varphi$ und $j_\varphi$.} Die letzten Betrachtungen
haben uns veranlasst, die beiden Formen $\varDelta$ und $t$ des Systemes von
$\varphi = kf + \lambda\varDelta$ zu berechnen. Wir wollen zur Vervollständigung auch
noch die beiden übrigen Formen $i$ und $j$ dieser Form $\varphi$ ermitteln,
und schlagen zu dem Zwecke wiederum den Weg der directen Ueberschiebung ein. Gemäss der Definition von $i$ als vierte Ueberschiebung
der Grundform über sich selbst wird:
\[
i_\varphi = (kf+\lambda\varDelta,\ kf+\lambda\varDelta)^4 = k^2(f,f)^4 + 2k\lambda(f,\varDelta)^4 + \lambda^2(\varDelta,\varDelta)^4
\]
\[
= k^2 i + 2k\lambda j + \lambda^2\,\frac{i^2}{6}.
\]
Genau dieselbe Form erhält man als quadratische (also Hesse'sche)
Covariante $\tau$ der cubischen Form $g$. Da nämlich
\[
g = k^3 - \frac{i}{2}\,k\lambda^2 - \frac{j}{3}\,\lambda^3,
\]

*) Sind $x_1\,x_2\,x_3\,x_4$ die Wurzeln von $f=0$, dann ist jene Gleichung, welche
die Grössen $y_1 = x_1x_2+x_3x_4$, $y_2 = x_1x_3+x_2x_4$, $y_3 = x_1x_4+x_2x_3$ zu Wurzeln
hat, cubische Resolvente von $f$. Aus ihr geht durch lineare Transformation die
Form $g$ hervor.
%CONT>>


% ===== pdf idx 204, printed 190.0 =====
% !PRINTED 190
\weberpage{190}
%<<CONT
so wird die Determinante der zweiten Differentialquotienten:
\[
\frac{1}{2}\,\tau_g = \frac{1}{36}
\begin{vmatrix}
\dfrac{\partial^2 g}{\partial k^2} & \dfrac{\partial^2 g}{\partial k\,\partial \lambda} \\
\dfrac{\partial^2 g}{\partial k\,\partial \lambda} & \dfrac{\partial^2 g}{\partial \lambda^2}
\end{vmatrix}
=
\begin{vmatrix}
k, & \dfrac{-i}{6}\lambda \\
\dfrac{-i}{6}\lambda, & -\left(\dfrac{ik}{6}+\dfrac{j}{3}\lambda\right)
\end{vmatrix}
\]
\[
\tau_g = -\frac{1}{3}\left\{ i k^2 + 2k\lambda j + \frac{i^2}{6}\lambda^2 \right\};
\]
daher
\[
i_\varphi = -3\tau_g .
\]
Wir erhalten endlich:
\[
j_\varphi = (kf+\lambda\varDelta,\ g_1\varDelta - g_2 f)^4 = k g_1 j - k g_2 i + \lambda g_1\,\frac{i^2}{6} - \lambda g_2 j
\]
\[
= g_1\left(kj + \frac{\lambda i^2}{6}\right) - g_2(ki+\lambda j).
\]
Auch dieser Ausdruck rechts lässt sich als cubische Covariante $q$ der
Form $g$ herstellen; denn es ist
\[
q_g = (g,\ \tau_g) = \frac{1}{6}
\begin{vmatrix}
3g_1, & 3g_2 \\
-\dfrac{2}{3}(ik+\lambda j), & -\dfrac{2}{3}\left(kj+\dfrac{i\lambda}{6}\right)
\end{vmatrix}
\]
\[
= -\frac{1}{3}\left\{ g_1\left(kj+\frac{i\lambda}{6}\right) - g_2(ik+\lambda j) \right\}.
\]
Daher:
\[
j_\varphi = -3(g,\ \tau_g) = -3 q_g .
\]
Setzt man in die Relation, welche zwischen den In- und Covarianten
 der cubischen Form existirt [vgl. Nr. 148 (I)]
\[
q_g^2 = -\frac{1}{2}(\tau_g^3 + R_g g^2)
\]
die Werthe $q_g$, $\tau_g$, $R_g$, $g$ ein, ausgedrückt in den Formen $j$ und $i$,
so erhält man nach leichter Umformung auch noch
\[
i_\varphi^3 - 6 j_\varphi^2 = g^2(i^3 - 6j^2)
\]
eine Relation, welche nach Nr. 170 (II) diese Verbindung als Combinante
 charakterisirt. Das ganze Formensystem der Form $\varphi$ ist also
dargestellt durch:
\[
\varphi = kf+\lambda\varDelta, \quad \varDelta_\varphi = g_1\varDelta - g_2 f, \quad t_\varphi = g\cdot t,
\]
\[
i_\varphi = -3\tau_g, \quad j_\varphi = -3q_g,
\]
wobei $g = k^3 - \dfrac{ik\lambda^2}{2} - \dfrac{j}{3}\lambda^3$, und $g_1 = \dfrac{1}{3}\dfrac{\partial g}{\partial k}$, $g_2 = \dfrac{1}{3}\dfrac{\partial g}{\partial \lambda}$ ist.


% ===== pdf idx 205, printed 191.0 =====
% !PRINTED 191
\weberpage{191}
\begin{center}
\S.\,17.\\[2pt]
\textbf{Die Gleichung vierten Grades.}
\end{center}

173. \textit{Die Discriminante von $f=a_x^4$.} Ehe wir uns mit der Auflösung
 der Gleichung vierten Grades und mit den Eigenschaften ihrer
Wurzeln eingehender beschäftigen, wollen wir zunächst jene Invariante
der Form vierten Grades berechnen, welche als Discriminante bezeichnet
wird. Sie muss sich nach dem Vorausgegangenen rational durch $i$
und $j$ darstellen lassen, da diese beiden Formen die einzigen selbstständigen
 Invarianten des Systemes sind.

Gesetzt nun, die Form $f=a_x^4$ besitze zwei gleiche Wurzeln und
sei dementsprechend dargestellt durch
\weq{1}{f = a_x^4 = \alpha_x^2\cdot p_x^2 ,}
wo $\alpha_x$ der lineare Doppelfactor, und $p_x^2$ irgend eine quadratische Form
ist, dann muss für diese Form die Discriminante identisch verschwinden
. Es muss also in diesem Falle zwischen den beiden Fundamentalinvarianten
 $i$ und $j$ eine Relation bestehen, die sich von selbst
einstellen muss, sobald wir $i$ und $j$ für $f=\alpha_x^2\cdot p_x^2$ berechnen. Beide
Formen werden in diesem Falle durch einen gewissen Parameter $\varrho$
sich darstellen lassen, durch dessen Elimination sich die gewünschte
Relation ergiebt.

Nun hat die zweite Ueberschiebung von $f=\alpha_x^2\cdot p_x^2$ über $\alpha_x^2$ den
Werth:
\weq{2}{(a\alpha)^2 a_x^2 = \frac{1}{6}(p\alpha)^2\cdot\alpha_x^2 = \varrho\cdot\alpha_x^2}
und daher erhalten wir, indem wir diese Gleichung einmal über
$f=b_x^4$ schieben, als erste Ueberschiebung von $f$ über $\alpha$:
\weq{3}{(ab)(a\alpha)^2 a_x b_x^3 = -\varrho\cdot\alpha\cdot(f,\ \alpha).}
Die linke Seite enthält den Reducenten $(ab)$ und muss daher auf eine
Ueberschiebung von $\varDelta$ über $\alpha$ führen. In der That erhält man durch
Anwendung des Identitätssatzes:
\[
(ab)(a\alpha)^2 a_x b_x^3 = (ab)(a\alpha)\,a_x b_x^2\{(b\alpha)a_x + (ab)\alpha_x\},
\]
oder, weil hier das erste Glied rechts durch Vertauschung von $a$ mit
$b$ nur sein Zeichen ändert und sonach verschwindet:
\weq{4}{(ab)(a\alpha)^2 a_x b_x^3 = (ab)^2(a\alpha)\,a_x b_x^2\alpha_x = \alpha\cdot(\varDelta,\ \alpha).}
Subtrahirt man nun beide Gleichungen (3) und (4), so erhält man
nach Absonderung des Factors $\alpha_x$:
\weq{5}{(\varDelta+\varrho\cdot f,\ \alpha) = 0.}
Da nun die Invarianten $i$ und $j$ durch viermalige Ueberschiebung von
$\varDelta$ und $f$ über sich selbst und über einander entstehen, so brauchen
%CONT>>


% ===== pdf idx 206, printed 192.0 =====
% !PRINTED 192
\weberpage{192}
%<<CONT
wir nur, um die gewünschten Relationen für $i$ und $j$ zu erhalten,
die Gleichung (5) dreimal über $\varDelta$ und über $f$, oder besser über die Combination
 $\varDelta+\mu f$ zu schieben, wo $\mu$ jeden beliebigen Werth hat. Wir
erhalten dann zunächst:
\[
U = ((\varDelta+\varrho f,\ \alpha),\ \varDelta+\mu f)^3 = ((P,\ \alpha),\ Q)^3 = 0.
\]
Ersetzen wir hier einen Moment $\alpha$ durch $y$ und entwickeln das so
sich ergebende Polarenglied: $(PQ)^3 P_y Q_x$ in eine Reihe, so kommt,
wenn wir darin wiederum für $y$ das Symbol $\alpha$ einführen:
\[
U = 0 = ((\varDelta+\varrho f,\ \varDelta+\mu f)^3,\ \alpha) + \frac{1}{2}(\varDelta+\varrho f,\ \varDelta+\mu f)^4\cdot\alpha,
\]
oder, weil alle dritten Ueberschiebungen $(\varDelta,\varDelta)^3$, $(\varDelta,f)^3$, $(f,f)^3$ verschwinden
:
\[
(\varDelta+\varrho f,\ \varDelta+\mu f)^4 = 0,
\]
d.\,h.
\weq{6}{\frac{i^2}{6} + (\varrho+\mu)j + \varrho\mu\cdot i = 0.}
Da diese Identität für jeden Werth $\mu$ gilt, so müssen die Beziehungen
bestehen
\[
\frac{i^2}{6} + \varrho\cdot j = 0
\]
\[
i\varrho + j = 0.
\]
Durch Elimination des Parameters $\varrho$ erhält man
\[
R_f = j^2 - \frac{i^3}{6} = 0
\]
und in dieser Relation muss die linke Seite als Discriminante der
Form $f$ angesehen werden, da diese Relation nur entstanden ist unter
der Bedingung des Vorhandenseins einer Doppelwurzel; sie ist in der
That auch vom Grade $2(n-1)=6$ in den Coefficienten von $f$.

\textit{Anmerkung.} Berechnet man die Discriminante $R_g$ der cubischen
Resolvente
\[
g = k^3 - \frac{i}{2}\,k\lambda^2 - \frac{j}{3}\,\lambda^3
\]
nach den Regeln der niederen Algebra (vgl. auch Bd. I Nr. 183), so
ergiebt sich
\[
R_g = \frac{1}{9}\left(j^2 - \frac{i^3}{6}\right).
\]
Ebenso haben wir gesehen, dass die Invariante
\[
(t,\ t)^6 = \frac{1}{4}\left(j^2 - \frac{i^3}{6}\right).
\]

174. \textit{Die Cayley'sche Auflösung der Gleichung vierten Grades.} Die
invariantentheoretische Auflösung der Gleichung
\[
f = \bar a_0 x_1^4 + 4\bar a_1 x_1^3 x_2 + 6\bar a_2 x_1^2 x_2^2 + 4\bar a_3 x_1 x_2^3 + \bar a_4 x_2^4 = 0
\]
%CONT>>


% ===== pdf idx 207, printed 193.0 =====
% !PRINTED 193
\weberpage{193}
%<<CONT
knüpft Cayley an die in Nr. 162 aufgestellte Relation
\[
2t^2 = \frac{i}{2}\, f^2\, \Delta - \Delta^3 - \frac{j}{3}\, f^3,
\]
die wir auch in der Form schreiben können
\[
-2t^2 = \Delta^3 - \frac{i}{2}\, f^2\, \Delta - \frac{j}{3}\, (-f)^3.
\]
Vergleicht man nun die rechte Seite dieser Relation mit der cubischen Resolvente
\[
g = k^3 - \frac{i}{2}\, \lambda^2 k - \frac{j}{3}\, \lambda^3
\]
der Gleichung $f = 0$, so erkennt man, dass dieselbe daraus hervorgeht, wenn man $k$ durch $\Delta$ und $\lambda$ durch $-f$ ersetzt. Man hat also:
\[
-2t^2 = g(\Delta,\, -f)
\]
oder, wenn $k_1$, $k_2$, $k_3$ die drei Wurzeln von $g=0$ sind,
\weq{1}{-2t^2 = (\Delta+k_1f)(\Delta+k_2f)(\Delta+k_3f).}
Ist nun $f$ eine allgemeine Form ihres Grades, sind also ihre vier Wurzeln von einander verschieden, so ist die Discriminante $R$ von $f$ nicht null; folglich sind in diesem Falle keine zwei der drei Wurzeln $k_1$, $k_2$, $k_3$ einander gleich (da $R$ auch Discriminante der Resolvente $g$ ist); anderntheils können $f$ und $\Delta$ im Allgemeinen keinen Factor gemeinsam haben, wenn $R \gtrless 0$*)\footnote{Vgl. auch Nr.~177, Anmerkung.}.

Die drei Factoren der Gleichung (I) sind also in Bezug auf ihre linearen Factoren von einander verschieden, und da die linke Seite ein völliges Quadrat ist, so muss jeder der drei Factoren rechts gleichfalls das Quadrat einer Form zweiten Grades sein.

Wir können daher setzen
\weq{2}{\left.\begin{aligned}
\Delta + k_1 f &= \varphi^2 \\
\Delta + k_2 f &= \psi^2 \\
\Delta + k_3 f &= \chi^2
\end{aligned}\right\}}
wo $\varphi$, $\psi$ und $\chi$ drei quadratische Formen sind.

Alsdann wird die Covariante
\[
t = \frac{1}{\sqrt{-2}} \cdot \varphi \cdot \psi \cdot \chi.
\]
Die drei quadratischen Formen $\varphi$, $\psi$, $\chi$ sind aber conjugirte Formen (vgl. Nr. 137); denn die Functionaldeterminanten je zweier derselben sind immer der dritten proportional. So ist z. B.
%CONT>>


% ===== pdf idx 208, printed 194.0 =====
% !PRINTED 194
\weberpage{194}
%<<CONT
\[
\begin{aligned}
(\Delta+k_2f,\ \Delta+k_3f) &= (\Delta,\Delta)+(k_2-k_3)(f,\Delta)+k_2k_3(f,f) \\
&= (k_2-k_3)\cdot t \\
&= \frac{k_2-k_3}{\sqrt{-2}} \cdot \varphi \cdot \psi \cdot \chi.
\end{aligned}
\]
Anderntheils ist
\[
(\psi^2,\ \chi^2) = (\psi,\chi)\cdot\psi\cdot\chi.
\]
Durch Comparation beider Gleichungen ergiebt sich
\weq{3}{\left.\begin{aligned}
(\psi,\chi) &= \frac{k_2-k_3}{\sqrt{-2}}\,\varphi \\
\text{und analog}\quad (\chi,\varphi) &= \frac{k_3-k_1}{\sqrt{-2}}\,\psi \\
(\varphi,\psi) &= \frac{k_1-k_2}{\sqrt{-2}}\,\chi
\end{aligned}\right\}.}
Die Covariante $t$ ist also eine ganz specielle Form sechsten Grades, welche unter Adjunction der drei Wurzeln einer cubischen Gleichung in drei conjugirte quadratische Factoren gespalten werden kann.

175. Gerade dieser Umstand gestattet aber die Wurzeln der Gleichung $f=0$ in einfacher Weise zu berechnen. Durch Subtraction je zweier der Gleichungen (2) erhalten wir nämlich
\weq{4}{\left.\begin{aligned}
(k_1-k_2)f &= \varphi^2-\psi^2 = (\varphi-\psi)(\varphi+\psi) \\
(k_2-k_3)f &= \psi^2-\chi^2 = (\psi-\chi)(\psi+\chi) \\
(k_3-k_1)f &= \chi^2-\varphi^2 = (\chi-\varphi)(\chi+\varphi)
\end{aligned}\right\}.}
Wir schliessen aus diesen Gleichungen (4), dass je eine Summe und Differenz je zweier solcher quadratischer Formen alle Wurzeln von $f$ enthalten, dass daher die Differenzen
\[
\begin{aligned}
&\varphi-\psi \quad \text{und} \quad \psi-\chi \\
\text{oder}\quad &\psi-\chi \quad \text{und} \quad \chi-\varphi \\
\text{oder}\quad &\chi-\varphi \quad \text{und} \quad \varphi-\psi
\end{aligned}
\]
oder die entsprechenden Summenpaare stets einen gemeinsamen linearen Factor von $f$ besitzen. Denn enthielten sie keinen solchen gemeinsamen Factor, so müsste sich einmal $f$ darstellen lassen z. B. durch
\[
(k_1-k_2)f = (\varphi-\psi)(\varphi+\psi),
\]
dann aber auch durch
\[
\varrho \cdot f = (\varphi-\psi)(\varphi-\chi).
\]
Also müsste bis auf eine Constante $\chi = -\psi$ sein, und analog $\varphi = -\chi$, d. h. alle drei quadratischen Formen müssten übereinstimmen und demnach $t$ Doppelwurzeln besitzen. Dann verschwindet aber die Discriminante von $t$, welche nur eine Potenz der Discriminante $R$ von $f$ sein kann,
%CONT>>


% ===== pdf idx 209, printed 195.0 =====
% !PRINTED 195
\weberpage{195}
%<<CONT
wie wir in Nr. 167 gesehen haben. Nach unserer Voraussetzung ist aber $R$ von null verschieden, und folglich haben in der That je zwei solche oben angeführte Summen- oder Differenzenpaare einen gemeinsamen Factor.

Nun hatten wir früher (Nr. 128) den Satz erhalten: Besitzen zwei quadratische Formen einen gemeinsamen linearen Factor, so enthält ihre Functionaldeterminante denselben quadratisch. Also ist in der Relation
\[
\sqrt{-2}\cdot(\varphi-\psi,\ \varphi-\chi) = (k_3-k_1)\psi+(k_1-k_2)\chi+(k_2-k_3)\varphi
\]
die rechte Seite das vollständige Quadrat des gesuchten linearen Factors von $f$. Wir können diese rechte Seite auch in die Determinantenform bringen
\[
-D = \begin{vmatrix} 1 & k_1 & \varphi \\ 1 & k_2 & \psi \\ 1 & k_3 & \chi \end{vmatrix}.
\]
Um den Factor selbst zu erhalten, brauchen wir nur die erste Polare zu bilden; dann stellt
\[
-D_y = \begin{vmatrix} 1 & k_1 & \varphi_y \\ 1 & k_2 & \psi_y \\ 1 & k_3 & \chi_y \end{vmatrix}
\]
für $y_2=1$ und $y=0$ einen linearen Factor von $f$ dar. Nun hat bereits Euler gezeigt, dass, wenn $x_1=\sqrt{t}+\sqrt{u}+\sqrt{v}$ eine Lösung der Gleichung vierten Grades ist, die drei andern Lösungen daraus dadurch erhalten werden, dass man je zwei der drei Grössen $\sqrt{t}$, $\sqrt{u}$, $\sqrt{v}$ mit negativem Vorzeichen nimmt. Die vier linearen Factoren der Gleichung vierten Grades sind also:
\[
\begin{vmatrix} 1, & k_1, & (+\sqrt{\Delta+k_1f})_y \\ 1, & k_2, & (+\sqrt{\Delta+k_2f})_y \\ 1, & k_3, & (+\sqrt{\Delta+k_3f})_y \end{vmatrix},
\quad
\begin{vmatrix} 1, & k_1, & (+\sqrt{\Delta+k_1f})_y \\ 1, & k_2, & (-\sqrt{\Delta+k_2f})_y \\ 1, & k_3, & (-\sqrt{\Delta+k_3f})_y \end{vmatrix},
\]
\[
\begin{vmatrix} 1, & k_1, & (-\sqrt{\Delta+k_1f})_y \\ 1, & k_2, & (+\sqrt{\Delta+k_2f})_y \\ 1, & k_3, & (-\sqrt{\Delta+k_3f})_y \end{vmatrix},
\quad
\begin{vmatrix} 1, & k_1, & (-\sqrt{\Delta+k_1f})_y \\ 1, & k_2, & (-\sqrt{\Delta+k_2f})_y \\ 1, & k_3, & (+\sqrt{\Delta+k_3f})_y \end{vmatrix}.
\]
\textit{Anmerkung.} Eine zweite und dritte Methode der Auflösung einer Gleichung vierten Grades wird in den nächsten Paragraphen gegeben werden.

176. \textit{Bedingung, dass $f=0$ drei gleiche Wurzeln enthält.} An die Auflösung der Gleichung vierten Grades knüpfen sich naturgemäss die
%CONT>>


% ===== pdf idx 210, printed 196.0 =====
% !PRINTED 196
\weberpage{196}
%<<CONT
Fragen: Unter welchen Bedingungen wird es eintreten, dass zwei, drei, vier der eben ermittelten Wurzelgrössen übereinstimmen? Was den ersten und letzten Fall betrifft, so sind die Bedingungen hiefür bereits ermittelt.

Die Gleichung $f=0$ hat zwei gleiche Wurzeln, sobald die Discriminante
\[
j^2 - \frac{i^3}{6} = 0.
\]
Die Gleichung $f=0$ hat vier gleiche Wurzeln, wenn die Hesse'sche Form identisch verschwindet, also wenn
\[
\Delta = (ab)^2 a_x^2 b_x^2 = 0 \quad \text{(vgl. Nr. 55)}.
\]
Es bleibt also noch zu untersuchen, unter welcher Bedingung $f=0$ drei gleiche Wurzeln besitzt. Zu dem Zwecke wählen wir $f=a_x^4$ in einer Normalform, in welcher von vorn herein drei Wurzeln als gleich angenommen sind, nämlich
\[
f = 4x_1^3 x_2.
\]
Die Coefficienten $a_0\, a_1\, a_2\, a_3\, a_4$ sind also der Reihe nach: $0, 1, 0, 0, 0$. Trägt man die Werthe dieser Coefficienten in die Invarianten $i$ und $j$ ein, so findet sich:
\[
i = 0,\quad j = 0.
\]
Wir fragen uns nun, ob die beiden Bedingungen $i=0$, $j=0$ auch hinreichend sind, damit $f=0$ drei gleiche Wurzeln besitze. Da in Folge des Verschwindens dieser Invarianten auch die Discriminante $R=0$ wird, so hat $f$ sicher zwei gleiche Wurzeln und kann also in diesem Falle in der Form dargestellt werden
\[
f = a_0x_1^4 + 4a_1x_1^3x_2 + 6a_2x_1^2x_2^2.
\]
Bildet man aber für diese Form die Invariante $i$, so reducirt sie sich auf $a_2^2=0$, d. h. auch $a_2$ muss nach Voraussetzung verschwinden; dann ist aber $f$ von der Form
\[
f = a_0x_1^4 + 4a_1x_1^3x_2,
\]
d. h. $f=0$ besitzt drei gleiche Wurzeln, und $j=0$, $i=0$ sind also die nothwendigen und hinreichenden Bedingungen dafür.

177. \textit{Bedingung, dass $f=a_x^4$ ein vollständiges Quadrat ist.} Es kann nun endlich noch der Fall eintreten, dass $f$ zwei Paare gleicher Wurzeln enthält, oder mit anderen Worten in das Quadrat einer Form zweiten Grades zerfällt. Um auch hiefür die Kriterien zu gewinnen, nehmen wir $f$ in einer entsprechenden Normalform
\[
f = cx_1^2 x_2^2
\]
%CONT>>


% ===== pdf idx 211, printed 197.0 =====
% !PRINTED 197
\weberpage{197}
und bilden die Hesse'sche Covariante derselben. Sie wird (vgl. Nr. 54)
\[
\Delta=\frac{1}{18}\begin{vmatrix}cx_2{}^2 & 2cx_1x_2\\ 2cx_1x_2 & cx_1{}^2\end{vmatrix}=-\frac{1}{6}c^2x_1{}^2x_2{}^2=\varrho\cdot f.
\]
In diesem Falle ist sonach $\Delta$ proportional der Form $f$. Man kann auch sagen: Wenn $f$ ein vollständiges Quadrat ist, dann bestehen zwischen den Coefficienten von $f$ und $\Delta$ die Relationen:
\[
\overline{a_0}=\varrho\,\overline{\Delta_0},\quad \overline{a_1}=\varrho\,\overline{\Delta_1},\quad \overline{a_2}=\varrho\,\overline{\Delta_2},\quad \overline{a_3}=\varrho\,\overline{\Delta_3},\quad \overline{a_4}=\varrho\,\overline{\Delta_4}.
\]
Wir fragen uns nun: Ist umgekehrt, wenn $f=\varrho\cdot\Delta$, diese Bedingung hinreichend, damit $f$ zwei Paare gleicher Wurzeln enthält? Um die Frage zu entscheiden, bilden wir
\[
(yx)^2\cdot\Delta_x^4=(yx)^2(ab)^2a_x^2b_x^2=(a_xb_y-b_xa_y)^2a_x^2b_x^2=2\begin{vmatrix}f, & f_y\\ f_y, & f_{y^2}\end{vmatrix}.
\]
Ist nun $x=\alpha$ irgend eine Wurzel von $f$, so verschwindet in der Gleichung:
\weq{I}{(yx)^2\Delta_x^4=2\begin{vmatrix}f, & f_y\\ f_y, & f_{y^2}\end{vmatrix},}
wenn wir dieselbe für $x$ eintragen, zunächst $f$ selbst, dann aber auch $\Delta$ wegen der Proportionalität mit $f$. Sie reducirt sich also auf:
\[
0=f_y^2\quad\text{oder}\quad f_y=0,
\]
d.\,h. auch die erste Polare von $f$ verschwindet, wenn $x=\alpha$, für jeden Werth von $y$. Dies ist aber nur möglich, wenn in derselben
\[
f_y=y_1\frac{\partial f}{\partial x_1}+y_2\frac{\partial f}{\partial x_2}
\]
die Differentialquotienten für $x=\alpha$ einzeln verschwinden. Dann aber ist diese Wurzel $x=\alpha$ eine Doppelwurzel. Wenn aber jede Wurzel $\alpha$ von $f$ Doppelwurzel ist, so ist $f$ ein vollständiges Quadrat.

\textit{Anmerkung.} Aus der eben benutzten Relation (I) geht direct hervor, dass $\Delta$ und $f$ keine Wurzeln gemein haben können, wenn $R\gtrless 0$. Denn ersetzen wir darin $x$ der Reihe nach durch die vier Wurzeln $\alpha_1\,\alpha_2\,\alpha_3\,'\alpha_4$ von $f$ und multipliciren die vier Identitäten, so kommt: % [sic] Komma vor $\alpha_4$ wie gedruckt
\[
\Delta(\alpha_1)\cdot\Delta(\alpha_2)\cdot\Delta(\alpha_3)\cdot\Delta(\alpha_4)=16\cdot\frac{f_y^2(\alpha_1)\cdot f_y^2(\alpha_2)\cdot f_y^2(\alpha_3)\cdot f_y^2(\alpha_4)}{(y\alpha_1)^2\cdot(y\alpha_2)^2\cdot(y\alpha_3)^2\cdot(y\alpha_4)^2}
\]
oder für $y_2=0$, $y_1=1$ erhalten wir (vgl. auch Bd. I Nr. 168, 173)
\[
(-1)^{4\cdot4}R_{\Delta,f}=16\cdot\varrho\left\{\left(\frac{\partial f}{\partial x}\right)_{x=\alpha_1}\cdot\left(\frac{\partial f}{\partial x}\right)_{x=\alpha_2}\cdot\left(\frac{\partial f}{\partial x}\right)_{x=\alpha_3}\cdot\left(\frac{\partial f}{\partial x}\right)_{x=\alpha_4}\right\}^2
\]
oder
\[
R_{\Delta,f}=16\cdot\varrho R_{f_1,f}^2=\mu\cdot R^2,\quad\text{d.\,h.}
\]
\glqq die Resultante von $\Delta$ und $f$ ist bis auf einen Zahlenfactor gleich dem Quadrat der Discriminante $R$ von $f$.\grqq{}


% ===== pdf idx 212, printed 198.0 =====
% !PRINTED 198
\weberpage{198}
\S.\,18. \textbf{Normalform von $f=a_x^4$; Transformation in die Normalform.}

178. \textit{Das Formensystem der Normalform.} Bei den Formen vierter Ordnung ist vorzugsweise ein Ausdruck derselben als canonische Form oder Normalform bezeichnet worden, nämlich jener, in welchem alle ungeraden Potenzen von $x$ fehlen, also
\[
f=x_1{}^4+6mx_1{}^2x_2{}^2+x_2{}^4.
\]
Wir studiren zunächst die algebraischen Ausdrücke des Systemes dieser Form, um alsdann umgekehrt aus dessen Eigenschaften die Mittel zu gewinnen, eine beliebige Form vierten Grades auf ihre Normalform zu transformiren.

Da die Coefficienten von $f$ der Reihe nach die Werthe haben
\[
1,\ 0,\ m,\ 0,\ 1,
\]
so wird die Hesse'sche Covariante derselben repräsentirt durch
\[
\Delta=2\begin{vmatrix}x_1{}^2+mx_2{}^2, & 2mx_1x_2\\ 2mx_1x_2 & mx_1{}^2+x_2{}^2\end{vmatrix}=2m(x_1{}^4+x_2{}^4)+2(1-3m^2)x_1{}^2x_2{}^2.
\]
Die rechte Seite dieser Gleichung enthält gleichfalls nur gerade Potenzen, also tritt auch $\Delta$ in der Normalform auf. Die Invarianten
\[
i=2(\overline{a_0}\,\overline{a_4}-4\overline{a_1}\,\overline{a_3}+3\overline{a_2}{}^2),\quad j=6\begin{vmatrix}\overline{a_0} & \overline{a_1} & \overline{a_2}\\ \overline{a_1} & \overline{a_2} & \overline{a_3}\\ \overline{a_2} & \overline{a_3} & \overline{a_4}\end{vmatrix}
\]
erhalten für die canonische Form von $f$ die Werthe:
\[
i=2(1+3m^2),\quad j=6m(1-m)^2.
\]
Endlich die Functionaldeterminante $t=(f,\ \Delta)$ ergiebt sich in der Form
\[
t=\frac{1}{16}\begin{vmatrix}\dfrac{\partial f}{\partial x_1}, & \dfrac{\partial f}{\partial x_2}\\ \dfrac{\partial\Delta}{\partial x_1}, & \dfrac{\partial\Delta}{\partial x_2}\end{vmatrix}=2\begin{vmatrix}x_1{}^3+3mx_1x_2{}^2, & x_2{}^3+3mx_1{}^2x_2\\ mx_1{}^3+\frac{1}{2}(1-3m^2)x_1x_2{}^2, & mx_2{}^3+\frac{1}{2}(1-3m^2)x_1{}^2x_2\end{vmatrix}
\]
\[
=(1-9m^2)x_1x_2(x_1{}^4-x_2{}^4)=(1-9m^2)\cdot(x_1{}^2-x_2{}^2)(x_1{}^2+x_2{}^2)x_1x_2
\]
\[
=(1-9m^2)\varphi\cdot\chi\cdot\psi.
\]
Die Covariante $t$ zerfällt also unter Zugrundelegung der obigen Normalform von $f$ nicht nur a priori in ihre drei conjugirten quadratischen Formen, sondern diese quadratischen Formen stellen sich hiebei selbst in der einfachsten Normalform ein; insbesondere reducirt sich $\varphi$ auf das Product der Linearfactoren $x_1\cdot x_2$. Dieser Umstand legt die Frage nahe, ob nicht umgekehrt durch die Transformation
%CONT>>


% ===== pdf idx 213, printed 199.0 =====
% !PRINTED 199
\weberpage{199}
%<<CONT
\[
\begin{aligned}
y_1&=r_x\\
y_2&=s_x,
\end{aligned}
\]
wobei $r_x$ und $s_x$ die Linearfactoren der Form $\varphi$ einer in allgemeiner Darstellung gegebenen Form $f$ sind, diese Form $f$ auf die oben erwähnte Normalform gebracht werden kann.

179. \textit{Transformation der Form $f$ in die Normalform.} In der That erreichen wir durch diese Transformation das gewünschte Ziel. Zunächst geht aus der Annahme $\varphi=r_x\cdot s_x$ auch die Gestalt der beiden andern conjugirten Formen hervor, nämlich (vgl. Nr. 137)
\weq{1}{\psi=r_x^2-s_x^2,\ \chi=r_x^2+s_x^2}
und wir wissen aus der Auflösung der Gleichung vierten Grades, dass, wenn wir setzen:
\weq{2}{\varphi^2=\Delta+k_1f,}
die Quadrate von $\psi$ und $\chi$ dargestellt sind durch
\weq{3}{\psi^2=\Delta+k_2f}
\weq{4}{\chi^2=\Delta+k_3f,}
wo $k_1$, $k_2$, $k_3$ die Wurzeln der cubischen Resolvente $g=0$ sind. Nun können wir mit Hilfe der Identität
\[
a_x(rs)=(as)r_x-(ar)s_x
\]
den Ausdruck $f\cdot(rs)^4=a_x^4(rs)^4=[a_x(rs)]^4$ direct nach dem binomischen Lehrsatze entwickeln und erhalten
\[
f\cdot(rs)^4=(as)^4r_x^4-4(as)^3(ar)r_x^3s_x+6(as)^2(ar)^2r_x^2s_x^2-4(as)(ar)^3r_xs_x^3
\]
\[
+(ar)^4s_x^4.
\]
Wir haben also nur zu zeigen, dass in Folge der Beziehungen (1) (2) (3) und (4) das zweite und vierte Glied dieses Ausdruckes rechts verschwindet; dann ist in der That $f$ auf die gewünschte Normalform gebracht.

Zu dem Zwecke untersuchen wir die Ueberschiebungen $(\varphi,\ r_x^2)^2$ und $(\psi,\ r_xs_x)^2$. Beide sind null; denn nach Nr. 38 erhalten wir:
\[
(\varphi,\ r_x^2)^2=(r_x\cdot s_x,\ r_x\cdot r_x)^2=(rr)(sr)=0
\]
und
\[
(\psi,\ \varphi)^2=(r_x^2-s_x^2,\ r_x\cdot s_x)^2=(rr)(rs)-(sr)(ss)=0.
\]
Daraus folgt aber unmittelbar, dass auch:
\weq{5}{(\varphi^2,\ r_x^3\cdot s_x)^4=((\varphi^2,\ r_x^2)^2,\ r_xs_x)^2=0}
und
\weq{6}{(\psi^2,\ r_x^3\cdot s_x)^4=((\psi^2,\ \varphi)^2,\ r_x^2)^2=0}
ist, oder, indem wir beide Relationen combiniren, dass auch
\weq{7}{(\varphi^2-\psi^2,\ r_x^3s_x)^4=0.}


% ===== pdf idx 214, printed 200.0 =====
% !PRINTED 200
\weberpage{200}
Nun ist aber wegen der Gleichungen (2) und (3)
\[
\varphi^2-\psi^2=(k_1-k_2)f;
\]
also ist auch gemäss der Relation (7):
\weq{8}{(k_1-k_2)(f,\ r^3s)^4=(k_1-k_2)(ar)^3(as)=0}
und ebenso folgt, wenn wir in (8) $r$ mit $s$ vertauschen,
\[
(f,\ s^3r)^4=(ar)(as)^3=0.
\]
Die Coefficienten des zweiten und vierten Gliedes in der binomischen Entwicklung von $f\cdot(rs)^4$ verschwinden also, und daher hat man in der That:
\weq{A}{f\cdot(rs)^4=(ar)^4s_x^4+6(ar)^2(as)^2r_x^2s_x^2+(as)^4r_x^4.}
\textit{Anmerkung über die Verwerthung der Normalform bei Auflösung der Gleichung vierten Grades.} Die eben zum Abschluss gebrachten Untersuchungen geben eine weitere Methode, die Gleichung vierten Grades aufzulösen. Man berechnet zunächst die cubische Resolvente derselben
\[
g=k^3-\frac{i}{2}k-\frac{j}{3}
\]
und bestimmt die drei Wurzeln $k_1\,k_2\,k_3$ dieser Gleichung. Irgend eine derselben liefert wegen der Relation
\[
\Delta+k_1f=\varphi^2
\]
die quadratische Form $\varphi=\sqrt{\Delta+k_1f}$; man setzt dieselbe gleich null und bestimmt ihre linearen Factoren $r_x$ und $s_x$, mittelst welcher $f$ in die Normalform (A) transformirt wird. Dieselbe enthält nur mehr gerade Potenzen und kann also durch Quadratwurzeln gelöst werden.

180. \textit{Berechnung des Parameters $m$ der Normalform aus der allgemeinen Form $f=a_x^4$.} Die Herstellung der Normalform von $f$ kann auch noch in anderer Weise erreicht werden, indem wir direct den Parameter $m$ berechnen. Man erhält denselben aus einer Gleichung dritten Grades, die man als eine zweite cubische Resolvente von $f=a_x^4$ auffassen kann. Sind nämlich $i'$ und $j'$ die aus $f=a_x^4$ direct gebildeten Invarianten der biquadratischen Form, so bestehen zwischen ihnen und den Invarianten
\weq{1}{i=2(1+3m^2),\quad j=6m(1-m^2)}
der transformirten Form
\[
f=x^4+6mx^2y^2+y^4
\]
die Relationen:
\weq{2}{\Delta^4i'=i\quad\text{und}\quad\Delta^6j'=j,}
wenn $\Delta$ der Modul der Transformation ist, durch welche die Form $f=a_x^4$ übergeht in $f=x^4+6mx^2y^2+y^4$.


% ===== pdf idx 215, printed 201.0 =====
% !PRINTED 201
\weberpage{201}
Bilden wir nun den Quotienten
\weq{3}{\frac{i^3}{j^2}=\frac{\Delta^{12}i'^3}{\Delta^{12}j'^2}=\frac{i'^3}{j'^2},}
der somit vom Modul unabh\"angig ist, so gewinnen wir damit unter Benutzung der Relationen (1) eine Gleichung f\"ur den Parameter $m$, n\"amlich:
\weq{4}{\frac{i'^3}{j'^2}=\frac{2}{9}\cdot\frac{(1+3m^2)^3}{m^2(1-m^2)^2}.}
Der Ausdruck links ist eine bekannte Constante, und $m$ l\"asst sich somit durch Aufl\"osung einer cubischen Gleichung ermitteln.

Man nennt den Ausdruck $\dfrac{i^3}{j^2}$, da er, wie aus (3) erhellt, bei beliebiger linearer Transformation auch nicht einmal um eine Potenz des Moduls sich \"andert, \gesperrt{absolute Invariante}.

181. \textit{Berechnung des Moduls der Transformation aus dem Parameter $m$ der Normalform.} Sobald der Parameter $m$ bestimmt ist, so k\"onnen wir den Modul $\Delta=(rs)$ der Transformation
\[
\left.\begin{aligned}
x_1 &= r_1 x + r_2 y\\
x_2 &= s_1 x + s_2 y
\end{aligned}\right\},
\]
durch welche $f=a_x^4$ \"ubergeht in
\[
f = x^4 + 6mx^2y^2 + y^4,
\]
berechnen. Denn es ist wegen (2) und (1)
\[
\Delta^2\cdot\frac{j'}{i'} = \frac{3m(1-m^2)}{1+3m^2}
\]
oder
\weq{5}{\Delta^2 = \frac{3m(1-m^2)}{1+3m^2}\cdot\frac{i'}{j'}.}
Setzt man nun in Gleichung (4)
\[
m^2 = \frac{3i'\varrho+2j'}{3i'\varrho-6j'},
\]
so geht sie \"uber in die cubische Resolvente $g$, n\"amlich:
\[
g(\varrho) = \varrho^3 - \frac{i'}{2}\varrho - \frac{j'}{3} = 0.
\]
Die Gleichung (5) aber geht durch dieselbe Substitution \"uber in:
\[
\Delta^2 = -\frac{\varrho}{2m}.
\]
Man erh\"alt also f\"ur das Quadrat des Moduls $\Delta$ drei Werthe. Dass sich nur das Quadrat desselben bestimmt, kann dadurch erkl\"art werden, dass die Normalform unge\"andert bleibt, wenn man $m$ durch $-m$ und gleichzeitig $y$ durch $y\sqrt{-1}$ ersetzt. (Vergl. auch \gesperrt{Clebsch}, bin\"are Formen, Seite 168.)


% ===== pdf idx 216, printed 202.0 =====
% !PRINTED 202
\weberpage{202}
182. \textit{Rationale Transformation der biquadratischen Form.} Sowohl die Transformation mit Hilfe der linearen Factoren $r_x$, $s_x$ von $\varphi_x^2$, als auch jene, die sich auf die directe Berechnung des Parameters $m$ st\"utzte, konnte nur auf irrationalem Wege, n\"amlich durch L\"osung einer cubischen Resolvente hergestellt werden. \gesperrt{Hermite} hat im Crelle'schen Journal Bd. 52 eine rationale Methode angegeben, die wir im Folgenden klarlegen wollen. Bezeichnen wir den Quotienten $\dfrac{\Delta}{f}$ mit $y$, so ist nach den Regeln der Differentialrechnung
\[
dy = (f\cdot d\Delta - \Delta\cdot df):f^2,
\]
oder
\[
f^2\,dy = \begin{vmatrix} f, & df\\ \Delta, & d\Delta \end{vmatrix} = 4\begin{vmatrix} x_1f_1+x_2f_2, & f_1dx_1+f_2dx_2\\ x_1\Delta_1+x_2\Delta_2, & \Delta_1dx_1+\Delta_2dx_2 \end{vmatrix},
\]
die Determinante rechts ist aber das Product
\[
4\begin{vmatrix} x_1, & x_2\\ dx_1, & dx_2 \end{vmatrix}\cdot\begin{vmatrix} f_1, & f_2\\ \Delta_1, & \Delta_2 \end{vmatrix},
\]
dessen zweiter Factor die Functionaldeterminante $(f,\Delta)=t$ ist. Wir haben also die Beziehung:
\[
\frac{1}{4}f^2\,dy = t\cdot(x_1dx_2-x_2dx_1).
\]
Quadriren wir auf beiden Seiten und multipliciren mit $(-2)$, so folgt wegen der Identit\"at
\[
-2t^2 = \Delta^3 - \frac{i}{2}\Delta f^2 + \frac{i}{3}f^3 % [sic] printed \textquotedbl{}i/3\textquotedbl{} here, \textquotedbl{}j/3\textquotedbl{} in the line below
\]
\[
-\frac{1}{8}f^4(dy)^2 = \left(\Delta^3-\frac{i\Delta}{2}f^2+\frac{j}{3}f^3\right)(x_1dx_2-x_2dx_1)^2.
\]
Dividiren wir mit $f^3$ die ganze Gleichung und ziehen die Wurzel, so erhalten wir wegen $\dfrac{\Delta}{f}=y$
\[
\frac{dy\sqrt{-\dfrac{1}{8}}}{\sqrt{y^3-\dfrac{i}{2}y+\dfrac{j}{3}}} = \frac{x_1dx_2-x_2dx_1}{\sqrt{\bar a_0 x_1^4+4\bar a_1 x_1^3x_2+\cdots \bar a_4 x_2^4}}.
\]
Setzt man hier noch $y=\dfrac{j}{i}z$, so kommt:
\[
\frac{1}{i}\sqrt{\frac{-j}{8}}\cdot\frac{dz}{\sqrt{\dfrac{1}{3}-\dfrac{z}{2}+\dfrac{j^2}{i^3}z^3}} = \frac{x_1dx_2-x_2dx_1}{\sqrt{\bar a_0 x_1^4+4\bar a_1 x_1^3x_2\cdots \bar a_4 x_2^4}}.
\]
Damit ist der Ausdruck unter dem Wurzelzeichen rechts auf rationalem Wege in einen Ausdruck dritten Grades transformirt, der nur mehr einen Parameter, n\"amlich die absolute Invariante enth\"alt.


% ===== pdf idx 217, printed 203.0 =====
% !PRINTED 203
\weberpage{203}
183. \textit{Das Doppelverh\"altniss der vier Wurzeln von $a_x^4=0$.} Wir haben im Vorhergehenden eine Invariante $\dfrac{i^3}{j^2}$ kennen gelernt, die sich bei linearer Transformation auch nicht einmal um einen Zahlenfactor \"andert. Dieselbe Eigenschaft besitzt bekanntlich auch das Doppelverh\"altniss von vier Punkten, die man sich durch $a_x^4=0$ repr\"asentirt denken kann. Um nun den Zusammenhang der absoluten Invariante mit dem Doppelverh\"altniss zu studiren, w\"ahlen wir die Form vierter Ordnung in der \gesperrt{Legendre'schen} Normalform:
\[
\begin{aligned}
f &= (1-x^2)(1-k^2x^2) = 1-(1+k^2)x^2+k^2x^4\\
  &= x_1^4-(1+k^2)x_1^2x_2^2+k^2x_2^4,
\end{aligned}
\]
deren Coefficienten der Reihe nach dargestellt sind durch
\[
1,\ 0,\ \frac{-(1+k^2)}{6},\ 0,\ k^2.
\]
Dann ist:
\[
i = \frac{1}{6}(12k^2+(1+k^2)^2), \quad j = \frac{1}{36}(k^2+1)((k^2+1)^2-36k^2).
\]
Demnach wird:
\[
\frac{i^3}{j^2} = 6\cdot\frac{(12k^2+(1+k^2)^2)^3}{((1+k^2)^2-36k^2)^2(k^2+1)^2}.
\]
Dividirt man Z\"ahler und Nenner mit $k^6$, so kommt:
\[
\frac{i^3}{j^2} = 6\cdot\frac{\left(12+\left(k+\dfrac{1}{k}\right)^2\right)^3}{\left(k+\dfrac{1}{k}\right)^2\left(\left(k+\dfrac{1}{k}\right)^2-36\right)^2}.
\]
Substituirt man hier $12\lambda$ f\"ur $\left(k+\dfrac{1}{k}\right)^2$, so erh\"alt man:
\weq{1}{\frac{i^3}{j^2} = 6\cdot\frac{(1+\lambda)^3}{\lambda(\lambda-3)^2}}
Nun sind andernteils die vier linearen Factoren der Form:
\[
(x_1^2-x_2^2)(x_1^2-k^2x_2^2)=0
\]
dargestellt durch:
\[
x_1-x_2=0,\quad x_1+x_2=0,\quad x_1-kx_2=0,\quad x_1+kx_2=0,
\]
also ihr Doppelverh\"altniss durch:
\[
\varrho = \frac{1+k}{1-k}:\frac{-1+k}{-1-k} = \frac{1+k}{1-k}\cdot\frac{1+k}{1-k} = \frac{(1+k)^2}{(1-k)^2} = \frac{\left(k+\dfrac{1}{k}\right)^2}{\left(k-\dfrac{1}{k}\right)^2}.
\]
Daher:
\[
\varrho = \frac{3\lambda}{3\lambda-1}, \quad \text{oder} \quad \lambda = \frac{\varrho}{3(\varrho-1)}.
\]
Das Doppelverh\"altniss steht also in einem directen Zusammenhange
%CONT>>


% ===== pdf idx 218, printed 204.0 =====
% !PRINTED 204
\weberpage{204}
%<<CONT
mit der absoluten Invariante, und zwar erh\"alt man durch Substitution von
\[
\lambda = \frac{\varrho}{3(\varrho-1)}
\]
in (1) die Beziehung:
\weq{2}{\frac{i^3}{j^2} = 6\cdot\frac{3\left(1+\dfrac{\varrho}{3(\varrho-1)}\right)^3(\varrho-1)}{\varrho\left(\dfrac{\varrho}{3(\varrho-1)}-3\right)} = 18\cdot\frac{(4\varrho-3)^3}{\varrho(9-8\varrho)^2}.}
Es ergeben sich sonach aus (2) drei Werthe f\"ur $\varrho$ resp. $\lambda$, und somit sechs Werthe f\"ur $k$ wegen der Beziehung $\left(k+\dfrac{1}{k}\right)^2=12\lambda$.

\begin{center}
\textbf{\S.\,19. Ueber die Formen, f\"ur welche $(f,f)^4=0$.}
\end{center}

184. \textit{Einleitende Bemerkungen.} Wir haben gesehen, dass in dem Formensystem einer Form vierten Grades eine Covariante $t$ auftritt, sechsten Grades in den Variabeln und dritten Grades in den Coefficienten, welche zwei bemerkenswerthe Eigenschaften besitzt: einmal, dass sie in drei conjugirte quadratische Factoren zerf\"allt, also durch Wurzelziehen aufl\"osbar ist, sodann, dass ihre vierte Ueberschiebung \"uber sich selbst identisch verschwindet. Es liegt die Frage nach dem Zusammenhange dieser beiden Eigenschaften nahe, und \gesperrt{Clebsch} hatte sich die Aufgabe gestellt, zu untersuchen, ob jede Form sechsten Grades, deren vierte Ueberschiebung \"uber sich selbst verschwindet, in drei conjugirte quadratische Factoren zerf\"allt. Es zeigte sich (vgl. \gesperrt{Clebsch}, \glqq Bin\"are Formen\grqq{}, \S.\,111), dass in der That die eine Eigenschaft die andere involvirt, und somit jede derartige Gleichung sechsten Grades aufl\"osbar ist.

Wir stellen uns hier nach dem Vorgange von \gesperrt{Brioschi} und \gesperrt{Wedekind} die allgemeine Frage: Welche Eigenschaften besitzt eine beliebige Form $f$ $n^{\text{ten}}$ Grades, f\"ur welche $(f,f)^4$ identisch verschwindet?

Hierbei muss nat\"urlich der Grad $n$ von $f>3$ sein. Die Untersuchung wird ergeben, dass der Bedingung $(f,f)^4=0$ ausser den Formen mit $(n-1)$fachem Factor nur ganz specielle Formen gen\"ugen, welche \gesperrt{Klein} mit den Namen Tetraeder, Octaeder und Ikosaeder belegt hat.

Von diesen Formen haben wir bereits in \S.\,12 gesehen, wie sie sich, sobald sie in einer gewissen Normalform vorliegen, in einfacher Weise in quadratische Factoren spalten lassen.

185. \textit{Die Coefficientenrelationen, die aus $(f,f)^4=0$ hervorgehen.} Die Identit\"at $(f,f)^4=0$ liefert eine Reihe von Relationen f\"ur die
%CONT>>


% ===== pdf idx 219, printed 205.0 =====
% !PRINTED 205
\weberpage{205}
%<<CONT
Coefficienten der Form $f$, die wir in erster Linie zu untersuchen
haben, um die Eigenschaften solcher Formen $f$ kennen zu lernen.
Diese Relationen sind nat\"urlich nicht alle von einander unabh\"angig.

Es sei:
\[
f = C_0 x_1^n + \binom{n}{1} C_1 x_1^{n-1}x_2 + \binom{n}{2} C_2 x_1^{n-2}x_2^2 + \cdots + C_n x_2^n = a_x^n
\]
\[
C_k = a_1^{n-k}a_2^k,
\]
dann ist:
\[
\begin{aligned}
(f,f)^4 &= (ab)^4 a_x^{n-4}b_x^{n-4}\\
&= (ab)^4 (a_1x_1+a_2x_2)^{n-4}(b_1x_1+b_2x_2)^{n-4}\\
&= (ab)^4 \sum_{\lambda,\mu}\binom{n-4}{\mu}\binom{n-4}{\lambda}(a_1x_1)^{n-\lambda-4}(a_2x_2)^\lambda(b_1x_1)^{n-\mu-4}(b_2x_2)^\mu\\
&= 0.
\end{aligned}
\]
Da diese Identit\"at f\"ur jeden Werth von $x$ gilt, so m\"ussen die
Coefficienten der einzelnen Potenzen von $x_2$ f\"ur sich verschwinden.
Greifen wir aus dieser Entwicklung den Coefficienten von $x_2^\varrho$
heraus, wobei also $\varrho = \lambda + \mu$, oder $\mu = \varrho - \lambda$.
Er ist dargestellt durch:
\[
\begin{aligned}
0 &= (ab)^4 \sum \binom{n-4}{\varrho-\lambda}\binom{n-4}{\lambda} a_1^{n-\lambda-4}a_2^\lambda b_1^{n-\varrho+\lambda-4}b_2^{\varrho-\lambda}\\
&= \sum 2\left(a_1^4b_2^4 - 4a_1^3a_2b_1b_2^3 + 3a_1^2a_2^2b_1^2b_2^2\right)\binom{n-4}{\varrho-\lambda}\binom{n-4}{\lambda} a_1^{n-\lambda-4}a_2^\lambda b_1^{n-\varrho+\lambda-4}b_2^{\varrho-\lambda},
\end{aligned}
\]
oder
\weq{1}{0 = \sum \binom{n-4}{\varrho-\lambda}\binom{n-4}{\lambda}\left\{C_\lambda C_{\varrho+4-\lambda} - 4C_{\lambda+1}C_{\varrho+3-\lambda} + 3C_{\lambda+2}C_{\varrho+2-\lambda}\right\},}
wobei die Summen \"uber alle Werthe $\lambda$ zu erstrecken sind, f\"ur
welche die figurirten Zahlen noch Bedeutung haben. Hieraus gehen alle
Relationen hervor, welchen die Coefficienten von $f = a_x^n$ gen\"ugen
m\"ussen, wenn wir der Reihe nach $\varrho = 0$ bis $\varrho = 2n-8$ setzen.
Ihre Anzahl ist also $2n-7$.

Insbesondere erkennt man, dass die Coefficienten von Formen $f$ mit
$n$-fachen, oder $(n-1)$-fachen Wurzeln immer diese $(2n-7)$ Relationen
befriedigen. Denn nehmen wir diese Formen in der Normalform, so kann
man sie darstellen durch
\weq{2}{\text{und}\quad\left.\begin{aligned} f &= x^n\\ f &= x^{n-1}y\end{aligned}\right\}}
Im ersten Falle ist $C_0 = 1$, $C_1 = C_2 \cdots = C_n = 0$;
im zweiten Falle ist $C_1 = 1$, $C_0 = C_2 \cdots = C_n = 0$,
und diese Werthe der $C_i$ befriedigen immer die Gleichungen (1). Die
Formen (2) sind also selbstverst\"andliche Formen, f\"ur welche
$(f,f)^4 = 0$.


% ===== pdf idx 220, printed 206.0 =====
% !PRINTED 206
\weberpage{206}
186. \textit{Formen $f$ mit zwei- und mehrfachen Wurzeln.} Nehmen wir
nun an, dass die Form $f$ eine $\nu$-fache Wurzel, mindestens aber
eine Doppelwurzel besitze. Man kann sie dann stets so transformiren,
dass die $\nu$ ersten Coefficienten verschwinden. Wir d\"urfen also in
diesem Falle $f$ in einer Normalform zu Grunde legen, in welcher
\weq{I}{C_0 = C_1 = C_2 = C_3 = \cdots = C_{\nu-1} = 0,\quad \nu > 1.}
Soll nun f\"ur diese Form $f$ auch die Relation
\[
(f,f)^4 = 0
\]
bestehen, dann k\"onnen wir zeigen, dass in diesem Falle im Allgemeinen
alle weitern Coefficienten, $C_\nu$, $C_{\nu+1}\cdots C_{n-2}$,
gleichfalls verschwinden.

Zum Beweise gen\"ugt es, wenn wir zeigen, dass unter den
Voraussetzungen:
\weq{A}{(f,f)^4 = 0,}
\weq{B}{C_0 = C_1 = C_2 = \cdots = C_{\nu-1} = 0,}
auch
\[
C_\nu = 0
\]
ist. Nehmen wir n\"amlich an $C_\nu \gtrless 0$, dann muss in jener
von $2n-7$ Relationen, in welcher verm\"oge der Bedingungen (B) die
Gr\"osse $C_\nu^2$ allein auftritt, ihr Zahlencoefficient identisch
null sein. Es l\"asst sich aber zeigen, dass dieser $\gtrless 0$ ist
und sonach die Annahme $C_\nu \gtrless 0$ falsch sein muss.

Ersetzen wir n\"amlich in der Relation (1) Nr. 185 $\varrho$ durch
$2\nu-4$, so geht sie \"uber in
\weq{II}{0 = \sum_{\lambda=0}^{\lambda=2\nu-4}\binom{n-4}{\lambda}\binom{n-4}{2\nu-4-\lambda}\left\{C_\lambda C_{2\nu-\lambda} - 4C_{\lambda+1}C_{2\nu-\lambda-1} + 3C_{\lambda+2}C_{2\nu-\lambda-2}\right\}.}
In jedem Gliede dieser Relation ist die Summe der Indices der
Coefficienten gleich $2\nu$. Wenn daher der erste Factor eines jeden
Gliedes einen Index besitzt, der gr\"osser ist als $\nu$, so enth\"alt
nothwendig der andere Factor $C$ einen Index $< \nu$; diese Glieder
verschwinden also wegen der Bedingungen (B) und es bleiben nur die
Glieder $C_\nu^2$ \"ubrig. Solche Glieder treten aber in einer
Relation nur drei auf, n\"amlich nur, wenn $\lambda = \nu$, dann wird
der erste Term der Klammer $C_\nu^2$, ferner wenn $\lambda = \nu-1$,
wobei der zweite, und endlich, wenn $\lambda = \nu-2$, wobei der
dritte Term in $C_\nu^2$ \"ubergeht.

Dabei ist f\"ur $\lambda = \nu$ der Zahlencoefficient von $C_\nu^2$
dargestellt durch $\binom{n-4}{\nu}\binom{n-4}{\nu-4}$, und in den
beiden andern F\"allen durch
\[
-4\binom{n-4}{\nu-1}\binom{n-4}{\nu-3},\quad\text{resp.}\quad 3\binom{n-4}{\nu-2}\binom{n-4}{\nu-2},
\]
%CONT>>


% ===== pdf idx 221, printed 207.0 =====
% !PRINTED 207
\weberpage{207}
%<<CONT
so dass die Coefficientenrelation (II) \"ubergeht in:
\[
C_\nu^2\left\{\binom{n-4}{\nu}\binom{n-4}{\nu-4} - 4\binom{n-4}{\nu-1}\binom{n-4}{\nu-3} + 3\binom{n-4}{\nu-2}\binom{n-4}{\nu-2}\right\} = 0.
\]
Da wir nun angenommen hatten $C_\nu \gtrless 0$, so m\"usste der
zweite Factor verschwinden. Das ist aber im Allgemeinen nicht der
Fall, vielmehr besitzt derselbe den Werth
\[
\frac{6(n-3)(n-2)\binom{n-4}{\nu-4}\binom{n-4}{\nu-2}}{\nu(\nu-1)(\nu-2)(\nu-3)},
\]
wie wir sofort zeigen k\"onnen. Es ist n\"amlich:
\[
\begin{aligned}
\binom{m}{\lambda}\binom{m}{\lambda+4} - \binom{m}{\lambda+1}\binom{m}{\lambda+3} &= \binom{m}{\lambda}\binom{m}{\lambda+3}\left\{\frac{m-\lambda-3}{\lambda+4} + 1 - \frac{m-\lambda}{\lambda+1} - 1\right\}\\
&= \binom{m}{\lambda}\binom{m}{\lambda+3}(m+1)\left\{\frac{1}{\lambda+4} - \frac{1}{\lambda+1}\right\}\\
&= \frac{-3(m+1)}{(\lambda+4)(\lambda+1)}\binom{m}{\lambda}\binom{m}{\lambda+3}.
\end{aligned}
\]
Ferner ist:
\[
\begin{aligned}
\binom{m}{\lambda+1}\binom{m}{\lambda+3} - \binom{m}{\lambda+2}^2 &= \binom{m}{\lambda}\binom{m}{\lambda+2}\left\{\frac{m-\lambda-2}{\lambda+3} + 1 - \frac{m-\lambda-1}{\lambda+2} - 1\right\}\\
&= \frac{-(m+1)}{(\lambda+3)(\lambda+2)}\binom{m}{\lambda+1}\binom{m}{\lambda+2}.
\end{aligned}
\]
Multipliciren wir diese Gleichung mit $-3$ und addiren sie zur
vorhergehenden, dann erhalten wir:
\[
\begin{aligned}
\binom{m}{\lambda}\binom{m}{\lambda+4} - 4\binom{m}{\lambda+1}\binom{m}{\lambda+3} + 3\binom{m}{\lambda+2}^2 &= 3(m+1)\left\{\frac{\binom{m}{\lambda+1}\binom{m}{\lambda+2}}{(\lambda+3)(\lambda+2)} - \frac{\binom{m}{\lambda}\binom{m}{\lambda+3}}{(\lambda+1)(\lambda+4)}\right\}\\
&= \frac{3(m+1)\binom{m}{\lambda}\binom{m}{\lambda+2}}{(\lambda+3)(\lambda+1)}\left\{\frac{m-\lambda}{\lambda+2} + 1 - \frac{m-\lambda-2}{\lambda+4} - 1\right\}\\
&= \frac{6(m+1)(m+2)\binom{m}{\lambda}\binom{m}{\lambda+2}}{(\lambda+1)(\lambda+2)(\lambda+3)(\lambda+4)}.
\end{aligned}
\]
Ersetzen wir hierin $m$ durch $n-4$ und $\lambda+4$ durch $\nu$, so
geht diese Gleichung \"uber in
\weq{III}{\binom{n-4}{\nu-4}\binom{n-4}{\nu} - 4\binom{n-4}{\nu-3}\binom{n-4}{\nu-1} + 3\binom{n-4}{\nu-2}^2 = \frac{6(n-3)(n-2)\binom{n-4}{\nu-4}\binom{n-4}{\nu-2}}{\nu(\nu-1)(\nu-2)(\nu-3)},}
was wir zu beweisen hatten. Die Gr\"osse rechts ist stets von null
verschieden, sobald $n > 3$, wie wir vorausgesetzt, und $\nu$ solche
Werthe
%CONT>>


% ===== pdf idx 222, printed 208.0 =====
% !PRINTED 208
\weberpage{208}
%<<CONT
hat, dass \"uberhaupt dieser Relation noch eine Bedeutung zukommt.
Dies trifft nicht mehr ein, wenn $\nu = 2, 3, n-2, n-3$; und diese
F\"alle sind noch speciell zu untersuchen.

Im Allgemeinen verschwindet der Zahlencoefficient von $C_\nu^2$
nicht, folglich muss $C_\nu$ selbst verschwinden.

187. \textit{Untersuchung dieser Grenzf\"alle.} Besitzt $\nu$ den
Werth 2, dann kann wegen der figurirten Zahlen in der Relation
\weq{I}{0 = \sum_{\lambda=0}^{\lambda=2\nu-4}\binom{n-4}{\lambda}\binom{n-4}{2\nu-4-\lambda}\left\{C_\lambda C_{2\nu-\lambda} - 4C_{\lambda+1}C_{2\nu-\lambda-1} + 3C_{\lambda+2}C_{2\nu-\lambda-2}\right\}}
die Gr\"osse $\lambda$ den Werth 0 nicht \"ubersteigen. Die Relation
geht also, weil nach Voraussetzung $C_0 = C_1 = C_{\nu-1} = 0$,
\"uber in
\[
0 = 3\binom{n-4}{0}\binom{n-4}{0}C_2^2 = 3C_2^2,
\]
folglich:
\[
C_2 = 0.
\]
Besitzt $\nu$ den Werth 3, so kann $\lambda$ die Werthe 0, 1, 2
annehmen. In der Relation verschwinden dann, weil nach Voraussetzung
$C_0 = C_1 = C_2 = C_{\nu-1} = 0$ alle Glieder bis auf die Glieder
$C_3^2$, deren eines f\"ur $\lambda = 1$ aus $3C_{\lambda+2}C_{2\nu-\lambda-2}$
und deren anderes aus $-4C_{\lambda+1}C_{2\nu-\lambda-1}$ f\"ur
$\lambda = 2$ hervorgeht.

Der Zahlencoefficient von $C_3^2$ wird also:
\[
3\binom{n-4}{1}\binom{n-4}{1} - 4\binom{n-4}{2}\binom{n-4}{0} = -3(n-4)(n-2).
\]
Dieser Zahlencoefficient kann in der That verschwinden, wenn $n = 4$,
oder $n = 2$ ist. Im ersten Falle ist aber der Coefficient $C_3$
bereits der vorletzte Coefficient $C_{n-1}$, von dem wir \"uberhaupt
nicht behauptet haben, dass er verschwinden muss, und im zweiten
Falle existirt kein Coefficient $C_3$ mehr.

Besitzt $\nu$ den Werth $n-3$, so geht f\"ur diesen Fall die Relation
(I) \"uber in
\[
0 = \sum \binom{n-4}{\lambda}\binom{n-4}{2n-10-\lambda}\left\{C_\lambda C_{2n-6-\lambda} - 4C_{\lambda+1}C_{2n-7-\lambda} + 3C_{\lambda+2}C_{2n-8-\lambda}\right\}.
\]
Hierin muss $n-4 \geqq \lambda \geqq n-6$ sein, gem\"ass der beiden
figurirten Zahlen $\binom{n-4}{\lambda}$ resp.
$\binom{n-4}{2n-10-\lambda}$. Die Gr\"osse $\lambda$ kann also nur
die Werthe $n-4$, $n-5$, $n-6$ annehmen, und da nach Voraussetzung
alle Coefficienten von $C_{\nu-1} = C_{n-4}$ incl. abw\"arts
verschwinden, so reducirt sich die Relation auf die beiden Glieder
\[
-4\binom{n-4}{n-4}\binom{n-4}{n-6}C_{n-8}^2 + 3\binom{n-4}{n-5}\binom{n-4}{n-5}C_{n-3}^2 = 0,
\]
oder, was dasselbe ist, auf:
\[
\left\{-4\binom{n-4}{2} + 3\binom{n-4}{1}^2\right\}C_{n-3}^2 = 0.
\]


% ===== pdf idx 223, printed 209.0 =====
% !PRINTED 209
\weberpage{209}
Der Coefficient von $C^2_{n-3}$ verschwindet aber nur für $n=4$ oder
$n=2$, wie wir soeben gesehen haben. Im ersten Falle ist aber
ohnehin auch $C_{n-3}=C_1$ nach Voraussetzung null, und im zweiten
Falle existirt $C_{n-3}$ überhaupt nicht.

Besitzt endlich $\nu$ den Werth $n-2$, so kann wegen der figurirten
Zahlen in der Relation:
\[
0=\sum\binom{n-4}{\lambda}\binom{n-4}{2n-8-\lambda}\left\{C_\lambda C_{2n-4-\lambda}-4C_{\lambda+1}C_{2n-5+\lambda}+3C_{\lambda+2}C_{2n-6-\lambda}\right\}
\]
die Grösse $\lambda$ von $n-4$ nicht verschieden sein; und da alle Coefficienten
von $C_{\nu-1}=C_{n-3}$ incl. abwärts verschwinden, so reducirt sich diese
Relation auf
\[
3\binom{n-4}{n-4}\binom{n-4}{n-4}C_{n-2}^2=3C_{n-2}^2=0,
\]
d. h.
\[
C_{n-2}=0.
\]
Es muss also auch dieser Coefficient nach unsern Voraussetzungen
verschwinden, während $C_{n-1}$ und $C_n$ von null verschieden sein
können. Ist dies aber der Fall, so reducirt sich $f$ auf
\[
f=x_2^{n-1}(C_{n-1}x_1+C_nx_2)=\xi^{n-1}\eta;
\]
$f$ ist also dann eine der selbstverständlichen Formen.

188.\ \textit{Formen $f$ ohne mehrfache Wurzeln.} Wir sahen soeben, dass
keine Form $f$ mit mehrfachen Wurzeln existirt, abgesehen von den
selbstverständlichen Formen, welche der Bedingung $(f,f)^4=0$ genügen
. Es erübrigt also noch, die Formen mit einfachen Wurzeln zu
untersuchen. Hiebei setzen wir $f$ in einer Normalform voraus, in
welcher
\weq{1}{C_0=0,\quad C_1=1,\quad C_2=0,}
was unter Adjunction einer Wurzel von $f=0$ immer zu erreichen
möglich ist, ohne der Allgemeinheit der Form Eintrag zu thun. Es
lässt sich alsdann zeigen:

\glqq Ist der Grad $n$ der Form $f$ von 4, 6 oder 12 verschieden,
so verschwindet, sobald $(f,f)^4=0$ ist, unter der Voraussetzung
 (1) nicht nur $C_2$, sondern auch $C_3, C_4\cdots C_n$, d. h. die
Form $f$ ist eine der selbstverständlichen Formen $f=x^{n-1}y$.\grqq{}

Zum Beweise zeigen wir wiederum, dass, wenn alle Coefficienten
$C_3, C_4, \ldots C_{\nu-1}, \nu>1$, verschwinden, auch der Coefficient $C_\nu$ verschwinden
 muss. Wir haben nur zu sorgen, dass in der allgemeinen
Coefficientenrelation die Summe der Indices jedes Gliedes gleich
$\nu+1$ ist, damit stets ein Coefficient $C_i, i<\nu$, mit einem Coefficienten
$C_k, k>\nu$, verbunden auftritt, abgesehen von den Gliedern $C_\nu\cdot C_1$.


% ===== pdf idx 224, printed 210.0 =====
% !PRINTED 210
\weberpage{210}
Während alsdann die ersten Glieder nach Voraussetzung verschwinden
müssen, werden zunächst die Grössen $C_\nu C_1 \gtrless 0$ sein, so dass, wenn
gleichwohl die Relation $(f,f)^4=0$ bestehen soll, der Zahlencoefficient
von $C_\nu C_1$ nicht von null verschieden sein kann. Damit nun aber die
Summe der Indices gleich $\nu+1$ sei, setzen wir in der allgemeinen
Relation $\varrho=\nu-3$; sie geht dann über in
\[
0=\sum\binom{n-4}{\nu-3-\lambda}\binom{n-4}{\lambda}\left\{C_\lambda C_{\nu+1-\lambda}-4C_{\lambda+1}C_{\nu-\lambda}+3C_{\lambda+2}C_{\nu-1-\lambda}\right\}
\]
und reducirt sich gemäss den vorausgegangenen Bemerkungen auf
\weq{2}{\binom{n-4}{\nu-4}\binom{n-4}{1}C_1C_\nu-4\binom{n-4}{\nu-3}\binom{n-4}{0}C_1C_\nu=0,}
welche zwei Glieder aus ihr für $\lambda=1$ und $\lambda=0$ entstehen.

Die Gleichung (2) geht aber wegen $C_1=1$ über in
\[
\binom{n-4}{\nu-4}\left\{n-4-4\cdot\frac{n-\nu}{\nu-3}\right\}C_\nu=0,
\]
oder
\[
\frac{n}{\nu-3}\cdot\binom{n-4}{\nu-4}\left\{\nu-7+\frac{12}{n}\right\}C_\nu=0.
\]
Verschwindet nun der Zahlencoefficient von $C_\nu$ nicht, so muss
$C_\nu$ verschwinden. Dieser Zahlencoefficient kann aber nur dann gleich
null werden, wenn $n$ ein Theiler von 12 ist, also wenn $n$ gleich 4,
6 oder 12 ist. Es können also ausser den selbstverständlichen Formen
mit mehrfachen Wurzeln in der That Formen vierten, sechsten und
zwölften Grades auftreten, welche der Bedingung $(f,f)^4=0$ genügen.
Wir wollen daher im Folgenden diese speciellen Fälle $n=4, 6, 12$
noch eingehender untersuchen.

189.\ \textit{Die Formen des Tetraeders und Octaeders.} Wir setzen hier
wie im Folgenden voraus, dass die Form vierten Grades $f$, welche
der Bedingung
\[
(f,f)^4=0
\]
genügt, wiederum in einer Normalform gegeben sei, in welcher
\[
C_0=0,\quad C_1=1,\quad C_2=0
\]
ist. Dann liefert die Bedingung $(f,f)^4=i=0$ unmittelbar
\[
C_0C_4-4C_1C_3+3C_2{}^2=0,
\]
oder
\[
-4C_1C_3=0,
\]
also
\[
C_3=0.
\]
Die betreffende Form $f$ hat daher die Gestalt:
\weq{1}{f=y(4x^3+C_4y^3).}


% ===== pdf idx 225, printed 211.0 =====
% !PRINTED 211
\weberpage{211}
Ich möchte hiebei erwähnen, dass man gewöhnlich die Tetraederform
 in folgender Normalform
\[
f=x^4+6mx^2y^2+y^4
\]
darstellt (vergl. Nr. 180), wobei $m$ wegen $i=(f,f)^4=0$ vermöge
Gleichung (4) Nr. 180 den Werth
\[
m=\frac{1}{3}\sqrt{-3}
\]
besitzt. Es ist dann
\weq{2}{f=x^4\pm 2\sqrt{-3}\cdot x^2y^2+y^4.}
Die Form (1) geht in die Form (2) über durch eine Transformation,
deren Modul $\varDelta$ der Gleichung
\[
\varDelta^2=-\frac{\varrho}{2m},
\]
genügt, wobei $\varrho=\varepsilon^k\sqrt[3]{-\frac{1}{3}j}=\varepsilon^k\sqrt[3]{-2C_4}$, und $\varepsilon^k$ eine dritte Wurzel
der Einheit ist. (Vergl. auch Nr. 181.)

\textit{Anmerkung.} Das Product der beiden Formen
\[
f=x^4\pm 2\sqrt{-3}\cdot x^2y^2+y^4
\]
liefert die Würfelform $x^8+14x^4y^4+y^8$. (Vergl. Nr. 141.)

Um die Coefficienten der Form sechsten Grades, welche der Bedingung
 $(f,f)^4=0$ genügt, zu berechnen, gehen wir wieder zurück auf
die allgemeine Coefficientenrelation, die sich für $n=6$ in der Form
\[
0=\sum\binom{2}{\varrho-\lambda}\binom{2}{\lambda}\left\{C_\lambda C_{\varrho+4-\lambda}-4C_{\lambda+1}C_{\varrho+3-\lambda}+3C_{\lambda+2}C_{\varrho+2-\lambda}\right\}
\]
darstellt. Wir erhalten aus ihr vier Relationen, für $\varrho=0,1,2,3$,
wobei $\lambda$ entsprechend die Werthe $0; 0,1; 0,1,2; 1,2$ annehmen
kann. Es ergiebt sich, wenn wiederum $C_0=0$, $C_1=1$, $C_2=0$ gesetzt
 wird:
\weq{1}{C_0C_4-4C_1C_3+3C_2{}^2=0,}
also $C_3=0;$
\weq{2}{2\{C_0C_5-4C_1C_4+3C_2C_3\}+2\{C_1C_4-4C_2C_3+3C_3C_2\}=0,}
also $(-2\cdot4+2)C_1C_4=0,$

d. h. $C_4=0;$
\weq{3}{\{C_0C_6-4C_1C_5+3C_2C_4\}+4\{C_1C_5-4C_2C_4+3C_3{}^2\}+\{C_2C_4-4C_3{}^2+3C_4C_2\}=0,}
oder $(-4+4)C_1C_5=0,$

d. h. $C_5$ braucht nicht zu verschwinden;
\weq{4}{2\{C_1C_6-4C_2C_5+3C_3C_4\}+2\{C_2C_5-4C_3C_4+3C_4C_3\}=0,}
d. h. $C_6=0.$


% ===== pdf idx 226, printed 212.0 =====
% !PRINTED 212
\weberpage{212}
Also hat das Octaeder die Form:
\[
xy(x^4+C_5y^4)=f,
\]
oder für $y_1=\sqrt[4]{C_5}\,y$
\[
f=\frac{1}{\varrho}xy_1(x^4+y_1{}^4),
\]
wo $\varrho$ eine Constante ist. (Vergl. auch Nr. 138.)

190.\ \textit{Die Ikosaederform.} In derselben Weise wie vorhin, berechnet
 man der Reihe nach die Coefficienten der Form $f=a_x^{12}$, welche
der Bedingung $(f,f)^4=0$ genügt. Die allgemeine Coefficientenrelation
 wird hiebei
\[
0=\sum\binom{8}{\lambda}\binom{8}{\varrho-\lambda}\left\{C_\lambda C_{\varrho+4-\lambda}-4C_{\lambda+1}C_{\varrho+3-\lambda}+3C_{\lambda+2}C_{\varrho+2-\lambda}\right\}
\]
und es genügt, wenn wir hieraus die zehn Relationen bilden, die sich
für $\varrho=0$ bis $\varrho=9$ ergeben. Unter der Voraussetzung $C_0=0$,
$C_1=1$, $C_2=0$ findet man alsdann:

Für $\varrho=0$ ist der Coefficient von $C_1C_3$ die Zahl $-4$, also $C_3=0$.

\begin{flushleft}
\text{\textquotedbl}\quad $\varrho=1$ \quad \text{\textquotedbl}\quad\text{\textquotedbl}\quad\quad\text{\textquotedbl}\quad\text{\textquotedbl}\quad $C_1C_4$ \quad \text{\textquotedbl}\quad\text{\textquotedbl}\quad $-24$, \quad \text{\textquotedbl}\quad $C_4=0$\\
\text{\textquotedbl}\quad $\varrho=2$ \quad \text{\textquotedbl}\quad\text{\textquotedbl}\quad\quad\text{\textquotedbl}\quad\text{\textquotedbl}\quad $C_1C_5$ \quad \text{\textquotedbl}\quad\text{\textquotedbl}\quad $-48$, \quad \text{\textquotedbl}\quad $C_5=0$\\
\text{\textquotedbl}\quad $\varrho=3$ \quad \text{\textquotedbl}\quad\text{\textquotedbl}\quad\quad\text{\textquotedbl}\quad\text{\textquotedbl}\quad $C_1C_6$ \quad \text{\textquotedbl}\quad\text{\textquotedbl}\quad $0$, \quad \text{\textquotedbl}\quad $C_6\begin{smallmatrix}\leqq\\[-2pt]\geqq\end{smallmatrix}0$ % ?? uncertain: exact form of stacked leqq/geqq symbol
\\
\text{\textquotedbl}\quad $\varrho=4$ \quad \text{\textquotedbl}\quad\text{\textquotedbl}\quad\quad\text{\textquotedbl}\quad\text{\textquotedbl}\quad $C_1C_7$ \quad \text{\textquotedbl}\quad\text{\textquotedbl}\quad $168$, \quad \text{\textquotedbl}\quad $C_7=0$\\
\text{\textquotedbl}\quad $\varrho=5$ \quad \text{\textquotedbl}\quad\text{\textquotedbl}\quad\quad\text{\textquotedbl}\quad\text{\textquotedbl}\quad $C_1C_8$ \quad \text{\textquotedbl}\quad\text{\textquotedbl}\quad $336$, \quad \text{\textquotedbl}\quad $C_8=0$\\
\text{\textquotedbl}\quad $\varrho=6$ \quad \text{\textquotedbl}\quad\text{\textquotedbl}\quad\quad\text{\textquotedbl}\quad\text{\textquotedbl}\quad $C_1C_9$ \quad \text{\textquotedbl}\quad\text{\textquotedbl}\quad $344$, \quad \text{\textquotedbl}\quad $C_9=0$\\
\text{\textquotedbl}\quad $\varrho=7$ \quad \text{\textquotedbl}\quad\text{\textquotedbl}\quad\quad\text{\textquotedbl}\quad\text{\textquotedbl}\quad $C_1C_{10}$ \quad \text{\textquotedbl}\quad\text{\textquotedbl}\quad $192$, \quad \text{\textquotedbl}\quad $C_{10}=0$.
\end{flushleft}

Für $\varrho=8$ ergiebt sich die Relation:
\[
\begin{aligned}
(8\cdot8-4)C_1C_{11}+\{28^2-4\cdot56^2+3\cdot70^2\}C_6{}^2&=0\\
&=60\{C_1C_{11}+49C_6{}^2\}=0,
\end{aligned}
\]
oder
\[
C_{11}+49C_6{}^2=0.
\]
Der Coefficient $C_{11}$ ist also linear von $C_6{}^2$ abhängig. Für $\varrho=9$
erhält man endlich:
\[
8C_1C_{12}=0,\quad\text{also}\quad C_{12}=0.
\]
Die Ikosaederform ist demnach dargestellt durch:
\[
f=12C_1x^{11}y+12\cdot11\cdot7C_6{}^2x^6y^6-49\cdot12C_6{}^2xy^{11},
\]
oder wegen $C_1=1$
\[
\frac{1}{12}f=xy(x^{10}+11\cdot7C_6x^5y^5-49C_6{}^2xy^{10}).
\]
Setzen wir hierin noch $y_1=\sqrt[5]{7C_5}\cdot y$, so erhalten wir die be-
%CONT>>


% ===== pdf idx 227, printed 213.0 =====
% !PRINTED 213
\weberpage{213}
%<<CONT
\[
\frac{\sqrt[5]{7C_0}}{12}\cdot f = xy_1(x^{10}+11\,x^5y_1^{5}-y_1^{10}).
\]

Vergleiche auch Nr. 142, Anmerkung.

191. \textit{Allgemeine Betrachtungen über die regulären Körper.} Wir haben somit gesehen, dass die Formen
\[
\begin{aligned}
f &= xy(x^{10}+11\,x^5y^5-y^{10})\\
f &= xy(x^4+y^4)\\
f &= x^4 \pm 2\,i\sqrt{3}\,x^2y^2+y^4
\end{aligned}
\]
und die daraus durch lineare Transformation hervorgehenden Formen die einzigen sind, für welche $(f,f)^4$ identisch verschwindet. Das Formensystem dieser Formen reducirt sich nach den in \S.\,20 zu entwickelnden Gesetzen direct auf die drei Covarianten $f$, $(f,f)^2=H$, $((f,f)^2,f)=T$ und eine Invariante $A$. Da nun $T$ eine schiefe Covariante, so muss, weil sich das Quadrat derselben stets durch gerade Formen ausdrücken lässt, und aus Gründen der Homogeneität in den Variabeln für das Ikosaeder eine lineare Relation zwischen $f^5$, $H^3$, $T^2$, für das Octaeder eine ebensolche zwischen $f^4$, $H^3$, $T^2$, und für das Tetraeder eine solche zwischen $f^3$, $H^3$, $T^2$ bestehen. In allen drei Gleichungen treten natürlich noch gewisse Potenzen der Invariante $A$ auf, damit auch die Homogeneität in den Coefficienten gewahrt bleibt\footnote{Zu den regulären Körpern gelangt man auch, wenn man diese Relationen als Definitionsgleichungen auffasst. Halphen hat gezeigt (siehe \glqq Mémoire sur la réduction des équations différentielles linéaires aux formes intégrables\grqq{}, Paris 1883): Die Formen, welche der Gleichung
\[
c_1f^{\lambda}+c_2H^{\nu}+c_3T^{\mu}=0
\]
genügen, sind ausser den allgemeinen cubischen und biquadratischen Formen noch die regulären Körper und die aus diesen Formen durch die Substitution $x_1=\varphi(y)$, $x_2=\psi(y)$ hervorgehenden Formen. Er nennt die cubischen, biquadratischen und regulären Formen zum Unterschied von den übrigen \glqq primitive Formen\grqq{}. Diese primitiven Formen, sowie die aus ihnen durch die erwähnte Transformation hervorgehenden genügen auch den Gleichungen:
\weq{I}{(f,T)^2=0,\quad (T,H)^2=0,\quad (H,f)^2=0}
\weq{II}{(f,T)=\varrho_i H^{\lambda},\quad (T,H)=\varrho_2\cdot f^{\mu},\quad (H,f)=\varrho_3\cdot T^{\nu}.} % [sic] \textquotedbl{}ρ_i\textquotedbl{} so gedruckt, vermutlich für ρ_1 (Inkonsistenz mit ρ_2, ρ_3)
}.

Denken wir uns nun jede der Relationen mit $H^3$ dividirt, so sind
\[
\frac{Cf^5}{H^3},\quad \frac{C_1f^4}{H^3},\quad \frac{C_2f^3}{H^3}
\]
Formen nullten Grades in den Variabeln, und wenn $C$, $C_1$, $C_2$ geeignete Potenzen der betreffenden Invarianten, auch nullten Grades in den Coefficienten. Man nennt diese Formen von der Dimension null in den Variabeln und Coefficienten die Parameter des Ikosaeders, resp.
%CONT>>


% ===== pdf idx 228, printed 214.0 =====
% !PRINTED 214
\weberpage{214}
%<<CONT
Octaeders oder Tetraeders und bezeichnet sie mit $\varrho$. Man könnte ebenso den andern Quotienten $\dfrac{T^2}{H^3}$ als Parameter einführen; denn vermöge der linearen Relation ist der eine bekannt, wenn der andere Quotient gegeben ist.

Der Parameter $\varrho$ wird jedesmal null, wenn $f$ verschwindet; denn $H$ und $f$ können im Allgemeinen keinen Factor gemeinschaftlich haben, da ja sowohl die Discriminante von $f$ als auch die Resultante von $H$ und $f$ eine Potenz der einzigen im System nicht verschwindenden Invariante $A$ sein müssen.

Wenn es also gelingt, diese Functionen
\[
\frac{Cf^5}{H^3}=\varrho,\quad \frac{C_1f^4}{H^3}=\varrho,\quad \frac{C_2f^3}{H^3}=\varrho
\]
umzukehren, also $x=\dfrac{x_1}{x_2}$ in Function von $\varrho$ darzustellen, etwa
\[
x=\frac{x_1}{x_2}=\vartheta(\varrho),
\]
so ist damit auch unmittelbar die Auflösung des Ikosaeders, Octaeders und Tetraeders gegeben. Die Functionen $\vartheta(\varrho)$ nennt Klein die Irrationalitäten des Ikosaeders, Octaeders und Tetraeders. Die Berechnung der Ikosaederirrationalität verlangt die Auflösung einer Gleichung sechzigsten Grades, und ihre Darstellung ist nur auf transcendentem Wege möglich. Dagegen können die beiden andern durch Wurzelziehen hergestellt werden; wir wollen uns daher mit der Berechnung derselben nun beschäftigen und zwar einmal für die Normalformen des Octaeders und Tetraeders, dann aber auch für die allgemeinen Formen sechsten beziehungsweise achten Grades.

192. \textit{Berechnung der Octaederirrationalität für die Normalform.} Die Covariante $H$ der Form
\[
f=xy(x^4+y^4)
\]
besitzt den Werth
\[
H=x^8+14\,x^4y^4+y^8. \quad \text{(Vgl. Nr. 189, Anm.)}
\]
Der Octaederparameter ist also, wenn wir die Invariante $C_1$ in die Grösse $\varrho$ hereinziehen, dargestellt durch
\[
\frac{x^4y^4(x^4+y^4)^4}{(x^8+14\,x^4y^4+y^8)^3}=\varrho.
\]
Die Grösse $\varrho=\dfrac{x}{y}$ genügt also einer Gleichung vierundzwanzigsten Grades. Um sie zu lösen, dividiren wir zunächst Nenner und Zähler der linken Seite mit $x^{12}y^{12}$ und erhalten, wenn wir im Zähler noch quadriren:
\[
\frac{\left(\dfrac{x^4}{y^4}+2+\dfrac{y^4}{x^4}\right)^2}{\left(\dfrac{x^4}{y^4}+14+\dfrac{y^4}{x^4}\right)^3}=\varrho,
\]
%CONT>>


% ===== pdf idx 229, printed 215.0 =====
% !PRINTED 215
\weberpage{215}
%<<CONT
oder, wenn wir nun
\weq{I}{\frac{x^4}{y^4}+\frac{y^4}{x^4}=\eta-14}
setzen:
\weq{II}{(\eta-12)^2=\varrho\cdot\eta^3.}
Hieraus lässt sich $\eta$ berechnen. Die Gleichung (I) liefert alsdann eine quadratische Gleichung für $\dfrac{x^4}{y^4}=\xi^4$, so dass sich endlich $\xi$ in der Form darstellt
\weq{III}{\xi=\varepsilon^k\sqrt[4]{\frac{\eta_i-14\pm\sqrt{\eta_i^{2}-28\eta_i+192}}{2}}=\vartheta(\varrho),}
wobei $\varepsilon^k$ eine der vier Wurzeln $\pm\sqrt{\pm i}$ und $\eta_i$ eine der drei Wurzeln von (II) ist.

193. \textit{Berechnung der Tetraederirrationalität für die Normalform.} Wir bezeichnen das Tetraeder
\[
f=x^4\pm 2\,i\sqrt{3}\,x^2y^2+y^4
\]
mit $u$, wenn das obere Vorzeichen, mit $v$, wenn das untere Vorzeichen gemeint ist. Die Hesse'sche Form $H$ von $u$ fällt dann mit $v$, und umgekehrt die Hesse'sche Form $H$ von $v$ mit $u$ bis auf eine Constante zusammen. Die Gleichung für die Tetraederirrationalität ist also dargestellt durch
\[
\frac{u^3}{v^3}=\varrho,
\]
oder
\[
\frac{u}{v}=\frac{(x^4+2\,i\sqrt{3}\,x^2y^2+y^4)}{(x^4-2\,i\sqrt{3}\,x^2y^2+y^4)}=\sqrt[3]{\varrho}.
\]

Hieraus folgt, wenn $\varepsilon$ die dritte Wurzel der Einheit bezeichnet, nach bekannten Proportionsgesetzen:
\weq{1}{\frac{\varepsilon u-\varepsilon^2v}{\varepsilon^2u-\varepsilon v}=\frac{\varepsilon\sqrt[3]{\varrho}-\varepsilon^2}{\varepsilon^2\sqrt[3]{\varrho}-\varepsilon}.}
Weil aber
\[
\varepsilon-\varepsilon^2=i\sqrt{3},
\]
so geht die linke Seite der Gleichung (1) über in
\[
-\left(\frac{x^2+y^2}{x^2-y^2}\right)^2.
\]
Demnach erhält man:
\[
\frac{x^2+y^2}{x^2-y^2}=\pm\sqrt{-\frac{\varepsilon\sqrt[3]{\varrho}-\varepsilon^2}{\varepsilon^2\sqrt[3]{\varrho}-\varepsilon}},
\]
oder nach bekannten Methoden
\[
\frac{x^2}{y^2}=\frac{1\pm\sqrt{-\dfrac{\varepsilon\sqrt[3]{\varrho}-\varepsilon^2}{\varepsilon^2\sqrt[3]{\varrho}-\varepsilon}}}{-1\pm\sqrt{-\dfrac{\varepsilon\sqrt[3]{\varrho}-\varepsilon^2}{\varepsilon^2\sqrt[3]{\varrho}-\varepsilon}}}.
\]


% ===== pdf idx 230, printed 216.0 =====
% !PRINTED 216
\weberpage{216}
Die Tetraederirrationalität ist also schliesslich dargestellt durch:
\[
\xi=\frac{x}{y}=\pm\sqrt[4]{\frac{\pm\sqrt{\varepsilon^2\sqrt[3]{\varrho}-\varepsilon}\pm\sqrt{-\varepsilon\sqrt[3]{\varrho}+\varepsilon^2}}{\pm\sqrt{-\varepsilon\sqrt[3]{\varrho}+\varepsilon^2}\mp\sqrt{\varepsilon^2\sqrt[3]{\varrho}-\varepsilon}}}=\vartheta(\varrho). % ?? uncertain: exact arrangement of the four ± / ∓ signs in this nested radical is hard to resolve at this resolution; sign pattern reproduced as best read
\]

194. \textit{Auflösung der allgemeinen Octaedergleichung.} Es sei nun eine Form
\[
f=\bar a_0x_1^{6}+6\,\bar a_1x_1^{5}x_2+\cdots \bar a_6x_2^{6}
\]
gegeben, welche die Eigenschaft besitzt, dass
\[
(f,f)^4=0
\]
ist. Wir bilden den zugehörigen Octaederparameter
\[
\frac{Cf^4}{H^3}=\varphi(x)=\varrho\cdot
\]
Die gebrochene Function $\varphi(x)$ ist vom nullten Grade in $x$, und wenn $C=(f,f)^6$, auch vom nullten Grade in den Coefficienten. Bei einer linearen Transformation dieses Parameters hebt sich, da im Nenner und Zähler gleichviel Klammerfactoren sich befinden, die Transformationsdeterminante weg.

Denkt man sich daher die Transformation von $f$ in die Normalform ausgeführt, wobei $f(x)$ in $f'(y)$, $H(x)$ in $H'(y)$, $\varphi(x)$ in $\varphi'(y)$ und $C$ in $C'$ übergeht, so besteht direct die Gleichung
\weq{1}{\varphi(x)=\varphi'(y)=\varrho.}
Die Grösse $y$ ist die Octaederirrationalität, $y=\vartheta(\varrho)$, oder auch
\weq{2}{y=\vartheta(\varphi(x)).}
Ersetzen wir nun die Variable $x$ durch einen festen numerischen Werth $\xi$, dem der Werth $\eta$ von $y$ entsprechen möge, so geht diese Gleichung über in:
\weq{3}{\eta=\vartheta(\varphi(\xi)).}
Ist insbesondere $\alpha$ eine Wurzel der Gleichung $\varphi(\xi)=\varrho=0$, so ist der dem $\alpha$ entsprechende Werth von $y$ dargestellt durch
\weq{4}{\beta=\vartheta(\varphi(\alpha)).}
Da nun $\varphi'(y)$ durch lineare Transformation aus $\varphi(x)$ hervorgegangen ist, so entspricht demnach dem Werthe
\[
\xi+\lambda\alpha \quad \text{der Werth} \quad \eta+\lambda\beta
\]
und die Gleichung (1) geht durch Substitution dieser entsprechenden Werthe über in
\weq{5}{\varphi(\xi+\lambda\alpha)=\varphi'(\eta+\lambda\beta)}
oder
\[
\frac{C\cdot f^4(\xi+\lambda\alpha)}{H^3(\xi+\lambda\alpha)}=\frac{C'\cdot f'^4(\eta+\lambda\beta)}{H'^3(\eta+\lambda\beta)}.
\]


% ===== pdf idx 231, printed 217.0 =====
% !PRINTED 217
\weberpage{217}

Es müssen demnach Zähler und Nenner beiderseits einander proportional sein, d.\,h. es müssen die Gleichungen bestehen:

\[
\begin{aligned}
C\cdot f^4(\xi+\lambda\alpha) &= m\cdot C'\cdot f'(\eta+\lambda\beta)\\
H^3(\xi+\lambda\alpha) &= m\cdot H'^3(\eta+\lambda\beta).
\end{aligned}
\]
Hieraus folgt durch Vergleichung der Co\"efficienten von $\lambda^1$

\[
\begin{aligned}
C\cdot f^3(\xi)\cdot f_\alpha(\xi) &= m\cdot C'\cdot f'^3(\eta)\cdot f'_\beta(\eta)\\
H^2(\xi)H_\alpha(\xi) &= m\cdot H'^2(\eta)\cdot H'_\beta(\eta).
\end{aligned}
\]
Dividirt man die beiden Gleichungen, so kommt:

\weq{6}{C\cdot f^3(\xi)\cdot f(\xi)_\alpha\cdot H'^2(\eta) H'_\beta(\eta) - C'\cdot f'^3(\eta)f'_\beta(\eta)\cdot H^2(\xi) H_\alpha(\xi) = 0.}

Hierin ist $\xi$ eine Constante, etwa $\xi_1=1$, $\xi_2=0$. Berechnet man zu diesem Werthe von $\xi$ die Grösse $\eta$ vermöge $\eta=\vartheta(\varphi(\xi))$, so liefert, weil auch $\beta=\vartheta(\varphi(\alpha))=\vartheta(0)$ eine völlig bekannte Grösse, die Relation (6) eine lineare Bestimmungsgleichung für die Wurzel $\alpha$ der Gleichung $\varphi(x)=0$, deren Wurzeln mit den Wurzeln von $f(x)=0$ übereinstimmen.

Es ist also auf diese Weise möglich, die Wurzeln der allgemeinen Octaedergleichung rational durch die Wurzeln der Normalform und bekannte Grössen darzustellen.

\textit{Anmerkung.} Nach Art der Gleichung (6) lassen sich auch die Wurzeln der allgemeinen Ikosaederform durch die Wurzeln der Normalform und bekannte Grössen ausdrücken.

195. \textit{Auflösung der Gleichung vierten Grades.} Ist die gegebene Gleichung vierten Grades $f=0$ von vornherein eine Tetraedergleichung, d.\,h. genügt $f$ der Bedingung $(f,f)^4=(ab)^4=i=0$, so kann man genau das nämliche Verfahren zur Auflösung derselben einschlagen, das wir vorhin bei der Octaedergleichung verfolgten.

Man kann indess auch die Auflösung einer beliebigen Gleichung vierten Grades

\[
f = a_x^4 = 0
\]
mit Hilfe der Tetraederirrationalität bewerkstelligen. Der Gedanke, der uns hierbei leitet, entspringt aus der Thatsache, dass man, von einer Originalform $f$ ausgehend, stets zwei Tetraederformen sich verschaffen kann, nämlich

\[
F = \Delta+\lambda_1 f,\quad F_1=\Delta+\lambda_2 f,
\]
in denen $\Delta$ das bekannte Symbol für die Hesse'sche Covariante und
%CONT>>


% ===== pdf idx 232, printed 218.0 =====
% !PRINTED 218
\weberpage{218}
%<<CONT
$\lambda_1$ und $\lambda_2$ die Wurzeln der Gleichung\footnote{Man bemerke, dass die Discriminante dieser Gleichung auch die Discriminante der biquadratischen Form $f=a_x^4$ selbst ist. (Vgl.\ Nr.\ 173.)}

\weq{1}{i_{\Delta+\lambda f} = \frac{i^2}{6} + 2\lambda j + i\lambda^2 = 0}

bedeuten.

Da nun $F$ und $F_1$ Tetraederformen sind, so müssen sie sich durch lineare Transformation in die Formen $u = y_1^4 + 2\sqrt{-3}\cdot y_1^2y_2^2+y_2^4$ und $v=y_1^4-2\sqrt{-3}y_1^2y_2^2+y_2^4$ bringen lassen. Alsdann ist:

\weq{2}{\Delta+\lambda_1 f = r\cdot u}

\weq{3}{\Delta+\lambda_2 f = s\cdot v,}

wobei $r$ und $s$ noch zu ermittelnde Constante sind. Zu ihrer Berechnung können wir unter anderm die Hesse'schen Formen von (2) und (3) benutzen. Wir erhalten:

\[
\begin{aligned}
\Delta_{\Delta+\lambda_1 f} &= r^2\cdot(u,\,u)^2\\
&= r^2\cdot\frac{48\sqrt{-3}}{72}\cdot v = r^2\cdot\frac{2\sqrt{-3}}{3}\cdot v.
\end{aligned}
\]
Es ist aber nach Nr.\ 171

\weq{4}{\Delta_{\Delta+\lambda_1 f} = \left(\lambda_1^2-\frac{i}{6}\right)\Delta+\left(\frac{j}{3}+\frac{i}{3}\lambda_1\right)\cdot f.}

Demnach hat man:

\[
\left(\lambda_1^2-\frac{i}{6}\right)\left\{\Delta+\frac{1}{3}\cdot\frac{j+\lambda_1 i}{\lambda_1^2-\frac{i}{6}}f\right\} = \frac{2r^2\sqrt{-3}}{3}\cdot v.
\]
Diese Gleichung ist eine Identität; es muss also insbesondere

\[
\Delta+\frac{1}{3}\cdot\frac{j+\lambda_1 i}{\lambda_1^2-\frac{i}{6}}f = \Delta+\lambda_2 f
\]
und

\weq{5}{\left(\lambda_1^2-\frac{i}{6}\right)s = \frac{2r^2\sqrt{-3}}{3}}

sein, vermöge der Gleichung (3). Ersetzt man in (5) die Grösse $\frac{i}{6}$ durch $\lambda_1\lambda_2$, was nach Gleichung (1) erlaubt ist, so kommt:

\[
\lambda_1(\lambda_1-\lambda_2)s = r^2\cdot\frac{2}{3}\sqrt{-3}.
\]
Entsprechend ergiebt sich bei Bildung der Hesse'schen Form von Gleichung (3)

\weq{6}{\lambda_2(\lambda_1-\lambda_2)r = s^2\cdot\frac{2}{3}\sqrt{-3}.}


% ===== pdf idx 233, printed 219.0 =====
% !PRINTED 219
\weberpage{219}

Aus diesen beiden letzten Gleichungen folgt durch Division:

\[
\frac{\lambda_1}{\lambda_2} = \frac{r^3}{s^3}
\]
und damit ist der Quotient $\frac{r}{s}$ bis auf eine dritte Wurzel der Einheit vollständig bestimmt.

Behufs Auflösung der allgemeinen Gleichung vierten Grades verfahren wir nun ganz analog wie bei Auflösung der allgemeinen Octaedergleichung.

Wir lösen zuerst die Tetraederform $\Delta+\lambda_1 f$, deren Hesse'sche Form durch $\lambda_1(\lambda_1-\lambda_2)\{\Delta+\lambda_2 f\}$ dargestellt ist. Es ist alsdann wie dort:

\weq{7}{C\cdot\left(\frac{\Delta+\lambda_1 f}{\Delta+\lambda_2 f}\right)^3 = \varphi(x) = \varrho,}

wobei nun $C = j_{\Delta+\lambda_1 f}:(\lambda_1-\lambda_2)\lambda_1$ ist. Wegen der Gleichungen (2) und (3) erhält man sonach auch:

\[
\frac{C\cdot r^3}{s^3}\cdot\frac{u^3(y)}{v^3(y)} = \varphi(y) = \varrho\,.
\]
Die Tetraederirrationalität ist

\[
y = \vartheta(\varrho) = \vartheta(\varphi(x))\,.
\]
Jedem Werth $\xi$ von $x$ entspricht ein Werth $\eta$ von $y$. Ist insbesondere $\alpha$ eine Wurzel der Gleichung $\varphi(x)=\varrho=0$, so entspricht ihr ein Werth

\[
\beta = \vartheta(\varphi(0))\,.
\]
Die lineare Gleichung zwischen $\alpha$ und $\beta$ erhält man wie früher, nur dass man hier einfacher die erste Polare von $\frac{u}{v}$ statt von $\frac{u^3}{v^3}$ nehmen kann. Sie geht wie früher wegen

\[
\left(\frac{\Delta+\lambda_1 f}{\Delta+\lambda_2 f}\right)^3 = \frac{r^3}{s^3}\cdot\frac{u^3}{v^3}
\]
aus der ersten Polare

\[
\left\{\frac{\Delta(\xi)+\lambda_1 f(\xi)}{\Delta(\xi)+\lambda_2(f\xi)}\right\}_\alpha = \varepsilon^\varkappa\sqrt[3]{\frac{\lambda_1}{\lambda_2}}\cdot\left\{\frac{u(\eta)}{v(\eta)}\right\}_\beta
\]
hervor, wobei $\varepsilon^\varkappa$ eine dritte Wurzel der Einheit ist. Die sich hieraus ergebende Wurzel $\alpha$ ist eine Wurzel der Gleichung $\Delta+\lambda_1 f=0$. Dagegen erhält man für $\varrho=C$, also $\beta=\vartheta(\varphi(C))$ die Wurzeln der Gleichung $f=0$; denn wenn $f$ verschwindet, geht (7) in die Relation

\[
C=\varrho
\]
über.


% ===== pdf idx 234, printed 220.0 =====
% !PRINTED 220
\weberpage{220}

Durch die eben entwickelte Methode für die Auflösung der Gleichung vierten Grades könnte man demnach, sobald man im Besitze von Tafeln für die Tetraederirrationalität

\[
y = \vartheta(\varrho)
\]
wäre, die Wurzeln jeder biquadratischen Gleichung durch zweimaliges Aufschlagen der Tafel und sonst rationales Rechnen ermitteln.

% [divider: short centred horizontal rule at foot of page, marking end of subsection; nothing else printed below it]


% ===== pdf idx 235, printed 220.001 =====
% !PRINTED none
% Part-title / section-opening page: large centred "Dritter Theil." heading, subtitle,
% and the "\S. 20" section heading, followed directly by running numbered body text
% (Nr. 196, Nr. 197). No running-head page number is visible in any of the three bands.

\begin{center}
\textbf{Dritter Theil.}\\[4pt]
\textbf{Systeme einzelner und simultaner Formen h\"oheren Grades.}
\end{center}

\begin{center}
\S.\,20.\\
\gesperrt{Ueberschiebung von Systemen.}
\end{center}

196. \textit{Begriff des Systemes.} Wenn wir eine Anzahl symbolischer Ausdr\"ucke in ihrer Gesammtheit bezeichnen wollen, so werden wir dasselbe kurz mit dem Namen \glqq System von Formen\grqq{} belegen. Der Begriff eines Formensystems ist damit in seinen weitesten Grenzen gegeben. Die Formen desselben k\"onnen einzelne symbolische Producte sein mit oder ohne Zahlenfactoren. Sie k\"onnen aber auch ganze Aggregate symbolischer Producte repr\"asentiren; nur wollen wir in diesem Falle voraussetzen, dass alle Glieder des Aggregates die gleiche Anzahl von Symbolen der etwa auftretenden Formen enthalten. Das ganze System kann alle Covarianten und Invarianten eines simultanen Systems bin\"arer Formen $f_1, f_2, f_3 \ldots f_n$ umfassen; es kann sich aber auch auf irgend eine durch gewisse Eigenschaften definirte Gruppe derselben beschr\"anken. Dabei ist nicht nothwendig, dass alle Co- und Invarianten eines solchen Systems von einander unabh\"angig sind; doch werden wir im allgemeinen wenigstens die Annahme machen, dass keine eine \gesperrt{ganze} Function von andern Formen des Systemes ist.

197. \textit{Ueberschiebung zweier Systeme.} Denken wir uns nun zwei solche Formensysteme mit je einer endlichen Anzahl von Formen gegeben, einmal das System
\weq{I}{\overset{a_1}{A_1(x)}, \quad \overset{a_2}{A_2(x)}, \quad \overset{a_3}{A_3(x)} \ldots \overset{a_r}{A_r(x)},}
sodann ein zweites System
\weq{II}{\overset{b_1}{B_1(x)}, \quad \overset{b_2}{B_2(x)}, \quad \overset{b_3}{B_3(x)} \ldots \overset{b_s}{B_s(x)},}
wobei mit $A_i(x)$, $B_i(x)$ irgend welche symbolische Producte oder ganze Aggregate von solchen bezeichnet sind, und die dar\"uber geschriebenen Zahlen $a_i$, $b_i$ ihren Grad in $x$ ausdr\"ucken.


% ===== pdf idx 236, printed 222.0 =====
% !PRINTED 222
\weberpage{222}
Wir stellen uns die Aufgabe, beide Systeme \"ubereinander zu schieben. Wie wir fr\"uher die Ueberschiebung zweier Formen dadurch bildeten, dass wir an dem Product derselben alle m\"oglichen Faltungen des betreffenden Ueberschiebungsgrades ausf\"uhrten, so werden wir auch jetzt eine Ueberschiebung der beiden Systeme erhalten, indem wir irgend ein Product aus den Formen des Systemes (I) mit einem Product von Formen des Systemes (II) falten.

Sei
\[
U = A_1{}^{\alpha_1} A_2{}^{\alpha_2} \ldots A_r{}^{\alpha_r}
\]
ein Product aus den Formen des ersten,
\[
V = B_1{}^{\beta_1} B_2{}^{\beta_2} \ldots B_s{}^{\beta_s}
\]
ein Product aus den Formen des zweiten Systemes, wobei die Exponenten $\alpha_i$, $\beta_i$ irgend welche ganze Zahlen, Null nicht ausgeschlossen, bedeuten.

Das Product $U$ ist in $x$ vom Grade
\[
a_1 \alpha_1 + a_2 \alpha_2 + \cdots a_r \alpha_r;
\]
das Product $V$ vom Grade
\[
b_1 \beta_1 + b_2 \beta_2 + \cdots b_s \beta_s.
\]
Um eine beliebige Ueberschiebung $(U, V)^{\nu}$ zu bilden, haben wir (vgl. Nr. 38) das Product $U \cdot V$ auf alle m\"oglichen Arten $\nu$ mal zu falten und die erhaltene Summe von symbolischen Producten durch die Anzahl ihrer Glieder zu dividiren. Jedes dieser Glieder besitzt einmal die \gesperrt{alten}, von vornherein dem ganzen Aggregat gemeinschaftlichen Klammerfactoren, die ja an der Faltung nicht betheiligt sind, dann aber auch \gesperrt{$\nu$ neue} Klammerfactoren, in deren jedem ein Symbol des Systemes der $A_i(x)$ mit einem Symbol des Systemes der $B_i(x)$ verbunden ist; ausserdem enthalten diese Glieder im Allgemeinen noch Factoren \gesperrt{erster Art}. Ist der Grad eines Gliedes der Ueberschiebung in den Factoren erster Art aus dem Systeme (I) noch $\varrho$, und ebenso in den Factoren erster Art aus dem Systeme (II) noch $\sigma$, dann m\"ussen f\"ur jedes Glied und damit auch f\"ur die ganze Ueberschiebung $(U, V)^{\nu}$ die Gleichungen bestehen:
\weq{1}{a_1 \alpha_1 + a_2 \alpha_2 + a_3 \alpha_3 \cdots + a_r \alpha_r = \varrho + \nu}
\weq{2}{b_1 \beta_1 + b_2 \beta_2 + b_3 \beta_3 \cdots + b_s \beta_s = \sigma + \nu.}
Ist umgekehrt ein Zahlensystem $\alpha_i$, $\beta_i$, $\nu$ gegeben, das diese Gleichungen befriedigt, so k\"onnen wir uns auch dazu eine Ueberschiebung der Producte $U$ und $V$ verschaffen. Denn durch die Exponenten $\alpha_i$ ist das Product $U$, durch die Exponenten $\beta_i$ das Product
%CONT>>


% ===== pdf idx 237, printed 223.0 =====
% !PRINTED 223
\weberpage{223}
%<<CONT
$V$ eindeutig bestimmt, und zwei solche Producte haben nur Eine Ueberschiebung vom Grade $\nu$.

198. \textit{Umgrenzung der zu betrachtenden Ueberschiebungen.} Die beiden Systeme k\"onnen nat\"urlich auf diese Weise zu unz\"ahligen Ueberschiebungen Veranlassung geben. Denn die beiden Gleichungen (1) und (2) sind lineare diophantische Gleichungen mit den Unbekannten $\alpha_i$, $\beta_i$, $\nu$, und zwar solche (vgl. Bd.\ I, Nr. 187), die immer positive ganzzahlige L\"osungen besitzen. Jeder Ueberschiebung entspricht aber eine solche L\"osung, und umgekehrt sahen wir soeben, geh\"ort jeder positiven ganzzahligen L\"osung eine und nur eine Ueberschiebung zu. Von dieser unendlich grossen Anzahl von Ueberschiebungen interessiren uns aber nun jene, durch welche wesentlich neue Formen entstehen. Es werden n\"amlich wie fr\"uher bei der Systembildung einfacher Formen die Ueberschiebungen in zwei Gruppen zu theilen sein: in eine erste Gruppe von Ueberschiebungen, bei denen die Faltung des Productes $U \cdot V$ kein zerfallendes Glied liefert, und in eine zweite, welche solche Glieder enthalten. Hierbei nennen wir ein Glied $G$ ein \gesperrt{zerfallendes}, wenn man dasselbe in zwei Theile $G_1$ und $G_2$ spalten kann, von denen der eine, also auch der andere, ein Glied einer fr\"uheren, niedrigeren Ueberschiebung ist, oder auch ein Product urspr\"unglicher Formen. Die Ueberschiebungen der zweiten Gruppe aber k\"onnen keine wesentlich \gesperrt{neuen} Formen liefern, da sich ja nach den allgemeinen S\"atzen in \S.\,3 jede solche Ueberschiebung durch dieses Glied und niedrigere Ueberschiebungen, also durch Bekanntes ausdr\"ucken l\"asst. (Siehe als Beisp. etwa Nr. 211.) Wir beschr\"anken uns also auf das Studium der ersten Ueberschiebungsgruppe, deren Formen wir in diesem Paragraphen mit $C_i$ bezeichnen werden. Von diesen Formen werden wir in der Folge vier Lehrs\"atze aufstellen, mit deren Hilfe die Endlichkeit des Formensystems einer bin\"aren Form sich einfach beweisen l\"asst. Hiebei bezeichnen wir als Form $C_i$ nur ein einziges \"ubrigens willk\"urliches Glied der betreffenden Ueberschiebung, die nach unsern allgemeinen Ueberschiebungss\"atzen durch ein solches Glied vollst\"andig definirt ist.

199. \textit{Erster Lehrsatz:} \gesperrt{\glqq Ist das System der Formen $A_i$ sowohl als auch das System der Formen $B_i$ ein endliches, so ist auch das durch Ueberschiebung entstandene System der Formen $C_i$ endlich.\grqq{}}

Der Beweis wird mit Hilfe des im ersten Bande dieser Vorlesungen \S.\,15 bewiesenen Satzes erbracht, wonach ein simultanes System linearer diophantischer Gleichungen nur eine endliche Zahl einfacher L\"osungen besitzt. Wir sahen n\"amlich soeben, dass jeder
%CONT>>


% ===== pdf idx 238, printed 224.0 =====
% !PRINTED 224
\weberpage{224}
%<<CONT
L\"osung der beiden diophantischen Gleichungen (1) und (2) Nr. 197 eine und nur eine Ueberschiebung entspricht. K\"onnen wir also zeigen, dass jeder zusammengesetzten L\"osung eine Ueberschiebung mit zerfallenden Gliedern zugeh\"ort und umgekehrt, dann k\"onnen die Ueberschiebungen mit nicht zerfallenden Gliedern, also unsere Formen $C_i$, nur den einfachen L\"osungen entsprechen. Diese sind aber nur in endlicher Zahl vorhanden, also auch die Formen $C_i$.

Angenommen aber, die Ueberschiebung $(U, V)^{\nu}$ enthalte ein Glied
\[
G = G_1 \cdot G_2,
\]
wobei $G_1$ und $G_2$ Glieder niedrigerer Ueberschiebungen vom Grade $\nu_1$ bezw. $\nu_2$ sind, so dass $\nu = \nu_1 + \nu_2$, dann l\"asst sich zu jedem der beiden Glieder je eine L\"osung $L_1 = (\alpha_i^{(1)}, \beta_i^{(1)}, \nu_1)$, beziehungsweise $L_2 = (\alpha_i^{(2)}, \beta_i^{(2)}, \nu_2)$ angeben, welche diese niedrigeren Ueberschiebungen v\"ollig bestimmen. Dabei ist alsdann, wenn $L = (\alpha_i, \beta_i, \nu)$ die L\"osung ist, durch welche die Ueberschiebung mit dem Gliede $G = G_1 \cdot G_2$ erhalten wurde:
\[
\begin{aligned}
\nu' &= \nu_1 + \nu_2 \\
\alpha_i &= \alpha_i^{(1)} + \alpha_i^{(2)} \\
\beta_i &= \beta_i^{(1)} + \beta_i^{(2)},
\end{aligned}
\]
d.\,h. die L\"osung $L$ l\"asst sich zusammensetzen aus den L\"osungen $L_1$ und $L_2$. Umgekehrt: Entsprechen den L\"osungen $L_1$ und $L_2$ bezw. die Ueberschiebungen $(U_1, V_1)^{\nu_1}$ und $(U_2, V_2)^{\nu_2}$, so entspricht der zusammengesetzten L\"osung $L = L_1 + L_2$ nothwendig die Ueberschiebung:
\[
(U_1 U_2, \ V_1 V_2)^{\nu_1+\nu_2}
\]
und diese enth\"alt zerfallende Glieder. Denn unter den Gliedern, die aus dem Producte
\[
U_1\, U_2\, V_1\, V_2
\]
hervorgehen, indem man dasselbe auf alle m\"oglichen Arten $(\nu_1+\nu_2)$-mal faltet, treten auch Producte von Gliedern auf, die entstehen, wenn man erst $U_1 V_1$ $\nu_1$-mal und dann $U_2 V_2$ $\nu_2$-mal faltet, die also Glieder der niedrigeren Ueberschiebungen
\[
(U_1, V_1)^{\nu_1} \text{ und } (U_2, V_2)^{\nu_2}
\]
darstellen. Es k\"onnen also in der That die Ueberschiebungen \gesperrt{ohne zerfallende Glieder nur den einfachen L\"osungen entsprechen}, und deren Anzahl ist eine endliche.

200. \textit{Definition eines vollst\"andigen Systemes.} Wir nennen nun ein System von Formen $A_i$ ein \gesperrt{vollst\"andiges}, sobald aus einer einzelnen Form oder einem Producte derselben durch Faltung keine wesentlich neue Form entstehen kann, sondern immer nur Formen, welche ganze rationale Functionen der $A_i$ sind.


% ===== pdf idx 239, printed 225.0 =====
% !PRINTED 225
\weberpage{225}
Solche vollständige Systeme sind die bereits früher aufgestellten
Systeme der Form zweiter, dritter oder vierter Ordnung. Jede
Ueberschiebung $(U,\ U')^\nu$ zweier Producte $U$ und $U'$ der Formen
$A_i$ ist also darstellbar in der Form
\[ (U,\ U')^\nu = \boldsymbol{F}(A_1,\ A_2,\ \ldots\ A_r) = \boldsymbol{F}(A_i), \]
wo $\boldsymbol{F}$ eine ganze rationale Function der Formen $A_i$ bedeutet. So ist
z. B. für die Form vierter Ordnung
\[ (f^{\alpha_1}\cdot \Delta^{\alpha_2}\cdot t^{\alpha_3}\cdot i^{\alpha_4}\cdot j^{\alpha_5},\ \ f^{\beta_1}\cdot \Delta^{\beta_2}\cdot t^{\beta_3}\cdot i^{\beta_4}\cdot j^{\beta_5})^\nu = \boldsymbol{F}(f,\ \Delta,\ t,\ i,\ j). \]

201. \textit{Zweiter Lehrsatz.} Sind die Systeme $(A)$ und $(B)$
vollständige Formensysteme mit den Formen
\weq{I}{ A_1,\ A_2,\ A_3,\ \ldots\ A_r }
\weq{II}{ B_1,\ B_2,\ B_3,\ \ldots\ B_s, }
gelten also für diese die Bedingungen
\weq{1}{ (U,\ U')^\nu = \boldsymbol{F}(A_i) }
\weq{2}{ (V,\ V')^\nu = \boldsymbol{\Phi}(B_i), }
dann führt wiederum die Ueberschiebung derselben auf ein drittes
System selbstständiger Formen
\weq{III}{ C_1,\ C_2,\ C_3,\ \ldots\ C_p, }
zu welchem selbstverständlich als nullte Ueberschiebungen auch die
Formen der beiden ersten Systeme gehören.

Von diesem System können wir behaupten:

\glqq Das System der Formen $C_i$ ist ein vollständiges, sobald die
ursprünglichen Systeme $(A)$ und $(B)$, aus denen dasselbe durch
Ueberschiebung entstanden ist, vollständig sind.\grqq{}

Unsere Aufgabe besteht darin, den Nachweis zu führen, dass auch
für das System der Formen $C_i$, wie für die beiden ursprünglichen
Systeme, die Gleichung besteht
\weq{3}{ (\Gamma,\ \Gamma')^\nu = \boldsymbol{\Psi}(C_i), }
wenn $\Gamma$ und $\Gamma'$ beliebige Producte der $C_i$ sind.

Nun entsteht jede dieser Ueberschiebungen (3) durch Faltung der
in $\Gamma$ enthaltenen Formen $C_i$ mit jenen, die sich in $\Gamma'$
befinden. Diese Formen $C_i$ sind aber aus Producten $U_i$ der Formen
$A_i$ durch Faltung mit Producten $V_i$ der Formen $B_i$ hervorgegangen.
Also kann man sich auch die Form $\boldsymbol{\Psi}$ direct durch Faltung eines
Productes $U\cdot V$ entstanden denken, wobei einmal Faltungen an den
Factoren $A_i$ von $U$, sodann Faltungen an den Factoren $B_i$ von $V$
und endlich Faltungen der Factoren $A_i$ mit Factoren $B_i$ auszuführen
sind. Bezeichnen wir das durch Faltungen der ersten Art aus $U$
entstehende Product mit $\overline U$, das durch Faltungen der zweiten
Art aus $V$ entstehende mit
%CONT>>


% ===== pdf idx 240, printed 226.0 =====
% !PRINTED 226
\weberpage{226}
%<<CONT
$\overline V$, dann entsteht $\boldsymbol{\Psi}$ aus dem Producte $\overline U\cdot\overline V$
durch Faltungen der dritten Art, und kann somit als Glied der
Ueberschiebung
\[ (\overline U,\ \overline V)^\varrho,\quad \varrho \leqq \nu, \]
angesehen werden. Daher lässt sich $\boldsymbol{\Psi}$ durch diese und niedrigere
Ueberschiebungen dieser Art nach \S.\,3 Nr. 44 darstellen, d. h. es ist
\[ (\Gamma,\ \Gamma')^\nu = \sum (\overline U,\ \overline V)^{\varrho_i}. \]
Da aber $\overline U$ durch Faltung von Formen $A_i$, $\overline V$ durch Faltung
von Formen $B_i$ entstanden ist, so kann nach Voraussetzung
\weq{4}{ \left.\begin{aligned} \overline U &= \sum (U'',\ U')^\nu = \boldsymbol{F}(A_i) \\ \overline V &= \sum (V'',\ V')^\mu = \boldsymbol{\Phi}(B_i) \end{aligned}\right\} }
gesetzt werden, und demnach ist
\[ (\Gamma,\ \Gamma')^\nu = \sum \big(\boldsymbol{F}(A_i),\ \boldsymbol{\Phi}(B_i)\big)^{\varrho_i}; \]
weil aber die Functionen $\boldsymbol{F}$ und $\boldsymbol{\Phi}$ rational und ganz in den
Formen $A_i$ bezw. $B_i$ sind, also
\[ \boldsymbol{F} = \sum A_1^{\alpha_1} A_2^{\alpha_2} \ldots A_r^{\alpha_r} = \sum U \]
\[ \boldsymbol{\Phi} = \sum B_1^{\beta_1} B_2^{\beta_2} \ldots B_s^{\beta_s} = \sum V, \]
und weil die Ueberschiebung von Summen gleich der Summe der
Ueberschiebungen ihrer Posten ist, so wird endlich
\weq{5}{ (\Gamma,\ \Gamma')^\nu = \sum (U,\ V)^{\varrho_i}. }
Die Ueberschiebungen $(U,\ V)^{\varrho_i}$ sind aber nun entweder
Ueberschiebungen der ersten Gruppe, d. h. (vergl. Nr. 198) sie enthalten
keine zerfallenden Glieder und wir können sie direct durch die Formen
$C_i$ ersetzen, oder aber sie enthalten zerfallende Glieder
$G_i = G_1\cdot G_2$; dann ersetzen wir sie zunächst nach Nr. 44 durch
$G_1\cdot G_2$ plus Glieder \gesperrt{niedrigerer Ueberschiebungen}
$(U,\ V)^\lambda$, wo $\lambda < \varrho_i$. Gleichzeitig führen wir alsdann
auch für die Factoren $G_1$ wie $G_2$ die ihnen einzeln entsprechenden
niedrigeren Ueberschiebungen nach demselben Satze Nr. 44 ein. Damit
ist eine erste Reduction der rechten Seite der Gleichung hergestellt.
Diese Reduction werden wir nöthigenfalls so lange fortsetzen, bis nur
mehr nullte Ueberschiebungen, d. i. Producte der $C_i$ auftreten. Die
rechte Seite der letzten Gleichung lässt sich also in der That auf
eine Function der Formen $C_i$ successive reduciren und die
Behauptung:
\[ (\Gamma,\ \Gamma')^\nu = \boldsymbol{\Psi}(C_i) \]
ist sonach erwiesen.


% ===== pdf idx 241, printed 227.0 =====
% !PRINTED 227
\weberpage{227}
202. \textit{Definition relativ vollständiger Systeme.} Es kann nun
aber auch der Fall eintreten, dass durch Faltung des Productes $U$
der Formen $A_i$, die einem bestimmten Systeme (I) angehören, Formen
$\overline U$ entstehen, welche der Gleichung
\weq{1}{ \overline U = \boldsymbol{F}(A_i) + \sum G\cdot W }
genügen. Dabei ist $\boldsymbol{F}(A_i)$ wieder eine ganze rationale Function
der $A_i$, $\sum G\cdot W$ aber ein Aggregat von symbolischen Producten,
deren jedes mindestens einen aus mehreren von vornherein bezeichneten
Klammerfactoren $G$ (oder auch ein Product solcher) zum Factor
besitzt. Diese Gruppe von Klammerfactoren $G$ nennen wir \glqq
Moduln\grqq{} (vergl. Nr. 43) des Formensystemes $(A)$, und sagen
alsdann:

\glqq Das Formensystem $(A)$ ist relativ vollständig in Bezug auf
die Moduln $G$.\grqq{}

Wir drücken dies durch die Gleichung (1) aus oder auch durch die
Congruenz:
\weq{2}{ \overline U \equiv \boldsymbol{F}(A_i)\ \mathrm{mod.}\ G. }
So bilden z. B. im Formensystem der Form vierter Ordnung $f=a_x^4$
die Formen
\[ f = a_x^4,\qquad t = (ab)^2(ca)a_xb_x^2c_x^3 \]
\[ \Delta = (ab)^2 a_x^2 b_x^2,\quad j = (ab)^2(bc)^2(ca)^2 \]
für sich ein relativ vollständiges System modulo $(ab)^4$, insofern
alle aus dem Product
\[ f^\lambda\cdot \Delta^\mu\cdot t^\nu\cdot j^\sigma = \Gamma \]
durch Faltung entstehenden Formen ganze Functionen dieser vier
Formen sind bis auf Glieder mit dem Factor $(ab)^4$. Wir drücken
dies aus durch
%CONT>>


% ===== pdf idx 242, printed 228.0 =====
% !PRINTED 228
\weberpage{228}
%<<CONT
\[ \overline\Gamma = (\Gamma,\ \Gamma')^\nu = \boldsymbol{F}(f,\ \Delta,\ t,\ j) + \sum (ab)^4 W, \]
oder
\[ \overline\Gamma = (\Gamma,\ \Gamma')^\nu \equiv \boldsymbol{F}(f,\ \Delta,\ t,\ j)\ \mathrm{mod.}\ (ab)^4.\footnote{Es dürfte hiebei die selbstverständliche Bemerkung erlaubt sein, dass, wenn $i=(ab)^4$ identisch null ist, dieses relativ vollständige System zu einem absolut vollständigen wird. (Vergl. auch Nr. 209.)} \]

203. \textit{Dritter Lehrsatz.} Sind nun zwei Systeme $(A)$ und $(B)$
gegeben, deren jedes relativ vollständig mod. $G$ ist, dann können
wir wiederum durch Ueberschiebung derselben ein drittes System $(C)$
erhalten. Von diesem System gilt der Satz:

\glqq Das System der Formen $C_i$ ist relativ vollständig in Bezug
auf die Moduln $G$, sobald es auch die beiden ursprünglichen
Systeme $(A)$ und $(B)$ sind, aus denen das System $(C)$ durch
Ueberschiebung entstanden ist.\grqq{}

Wir haben wiederum, wenn $\Gamma$ und $\Gamma'$ zwei beliebige
Producte der Formen $C_i$ sind, die Richtigkeit der Gleichung:
\[ (\Gamma,\ \Gamma')^\nu = \boldsymbol{\Psi}(C_i) + \sum W\cdot G, \]
nachzuweisen. Der Beweis gestaltet sich analog wie der des zweiten
Lehrsatzes. Zunächst schliessen wir ganz in derselben Weise wie
dort, dass sich die Form $(\Gamma,\ \Gamma')^\nu$ als Glied der
Ueberschiebung
\[ (\overline U,\ \overline V)^\varrho,\quad \varrho \leqq \nu, \]
darstellen lässt, wobei $\overline U$ und $\overline V$ Formen sind, die
aus $U$ resp. $V$ durch Faltung der Formen $A_i$ resp. $B_i$
entstehen. Es lässt sich demnach $(\Gamma,\ \Gamma')^\nu$ wiederum
durch eine Summe von Ueberschiebungen vom Typus $(\overline U,\ \overline V)^\varrho$
darstellen, so dass man hat
\[ (\Gamma,\ \Gamma')^\nu = \sum (\overline U,\ \overline V)^{\varrho_i}, \]
Nun ist aber nach Voraussetzung
\[ \overline U = \boldsymbol{F}(A_i) + \sum G\cdot W \]
\[ \overline V = \boldsymbol{\Phi}(B_i) + \sum G\cdot W', \]
also wird
\[ (\Gamma,\ \Gamma')^\nu = \sum \big(\boldsymbol{F}(A_i) + \sum G\cdot W,\ \ \boldsymbol{\Phi}(B_i)+\sum G\cdot W'\big)^{\varrho_i}, \]
oder nach Nr. 32, $(\beta)$
\[ = \sum \big(\boldsymbol{F}(A_i),\ \boldsymbol{\Phi}(B_i)\big)^{\varrho_i} \ + \sum \big(\boldsymbol{F}(A_i),\ \sum G\cdot W'\big)^{\varrho_i} \]
\[ {}+ \sum \big(\sum G\cdot W,\ \boldsymbol{\Phi}(B_i)\big)^{\varrho_i} + \sum \big(\sum G\cdot W,\ \sum G\cdot W'\big)^{\varrho_i}. \]

Der erste Term rechts reducirt sich in derselben Weise wie früher
auf $\sum (U,\ V)^{\varrho_i}$; die übrigen Terme auf Formen, welche die
Moduln $G$, die ja als Klammerfactoren nicht gefaltet werden können,
zu Factoren haben. Daher können wir die letzte Gleichung schreiben
\[ (\Gamma,\ \Gamma')^\nu = \sum (U,\ V)^{\varrho_i} + \sum G\cdot Q. \]

Die Ueberschiebungen $(U,\ V)^{\varrho_i}$ führen aber entweder wieder
auf Formen $C_i$ oder sind Ueberschiebungen mit zerfallenden
Gliedern $G_i$ und wir können alsdann genau wie in Nr. 201 die
allmähliche Reduction auf Formen $C_i$ vornehmen. Also ist in der
That
\[ (\Gamma,\ \Gamma')^\nu \equiv \boldsymbol{\Psi}(C_i)\ \mathrm{mod.}\ G. \]


% ===== pdf idx 243, printed 229.0 =====
% !PRINTED 229
\weberpage{229}
204. \textit{Vierter Lehrsatz.} Von den Moduln, in Bezug auf welche ein
erstes System (A) relativ vollständig ist, können einige als Klammerfactoren
von Formen $B_i$ eines zweiten Systemes (B) auftreten. Wir wollen diese
Moduln mit $H$, die übrigen des Systemes (A) mit $G$ bezeichnen. Es möge
nun das System (A) der Gleichung genügen:
\weq{1}{\bar U = F(A_i) + \sum H\cdot W + \sum G\cdot W'.}

Nehmen wir alsdann an, dass im System (B), das wir gleichfalls relativ
vollständig mod. $G$ voraussetzen, zu jedem der Moduln $H$ eine
entsprechende Form $B_i$ vorhanden sei, die neben Factoren erster Art nur
diesen Modul als Klammerfactor hat, dann genügt System (B) immer noch der
Gleichung
\weq{2}{\bar V = \Phi(B_i) + \sum G\cdot W'',}
wobei allerdings unter den Klammerfactoren von $W''$ auch Moduln $H$
vorhanden sein können, die wir jedoch hier neben den Moduln $G$ nicht
weiter hervorheben wollen.

Durch Ueberschiebung beider Systeme (A) und (B) entsteht ein drittes
System (C). Von ihm gilt der Satz:

\glqq Genügt das System (A) der Gleichung (1), das System (B) aber, das
je eine Covariante besitzen möge, die je einem Modul $H$ von (A)
entspricht, der Gleichung (2), so genügt auch das System (C) der
Gleichung
\[ (\Gamma,\ \Gamma')^\nu = F(C_i) + \sum G\cdot W''', \]
d. h. es ist relativ vollständig modulo $G$.\grqq{}

Ehe ich den Beweis antrete, will ich den Fall durch ein Beispiel
illustriren. Im System der Form $f=a_x^6$ bilden, wie wir demnächst sehen
werden, die drei Formen
\[ f = a_x^6,\quad H = (ab)^2 a_x^2 b_x^2,\quad t = (ab)^2(ca) a_x^3 b_x^4 c_x^5 \]
für sich ein in Bezug auf die beiden Moduln $(ab)^4$ und $(ab)^6$ relativ
vollständiges System (A), d. h. es ist:
\[ \bar U = F(f,\ H,\ t) + \sum (ab)^4\cdot W + \sum (ab)^6\cdot W'. \]

Stellen wir diesem System (A) das volle System (B) der biquadratischen
Covariante $k=(f,f)^4=(ab)^4a_x^2b_x^2$ von $f$ gegenüber
\weq{B}{\left.\begin{aligned}
B_1&=k, & B_2&=(k,\ k)^2, & B_3&=(k,\ k)^4\\
B_4&=(k,\ (k,\ k)^2), & B_5&=((k,\ k)^2,\ k)^2
\end{aligned}\right\},}
so ist dasselbe gerade ein System mit einer Form (nämlich der Grundform
$B_1=k$ selbst), die den Modul $(ab)^4$ besitzt, und das im übrigen,
%CONT>>


% ===== pdf idx 244, printed 230.0 =====
% !PRINTED 230
\weberpage{230}
%<<CONT
wie der nächste Paragraph schon lehren wird, relativ vollständig ist
modulo $(ab)^6$. Durch Ueberschiebung beider Systeme entsteht das System
(C), welches der Gleichung genügt:
\[ (\Gamma,\ \Gamma')^\nu = F(C_i) + \sum (ab)^6\cdot W'''. \]

Wir werden sehen, dass dieses System (C) zusammen mit der Invariante
$(ab)^6$ selbst, gerade das volle System von $f=a_x^6$ bildet.

Der Beweis unseres Satzes gestaltet sich nun wie die beiden
vorhergehenden. Es ist wiederum zunächst:
\[
\begin{aligned}
(\Gamma,\ \Gamma')^\nu &= \sum (\bar U,\ \bar V)^{\varrho_i}, \quad \varrho_i \leqq \nu\\
&= \sum \bigl(F(A_i) + \sum H\cdot W + \sum G\cdot W',\ \ \Phi(B_i) + \sum G\cdot W''\bigr)^{\varrho_i},
\end{aligned}
\]
oder:
\weq{3}{(\Gamma,\ \Gamma')^\nu = \sum \bigl(F(A_i),\ \Phi(B_i)\bigr)^{\varrho_i} + \sum \bigl(H\cdot W,\ \Phi(B_i)\bigr)^{\varrho_i} + \sum G\cdot Q.}

Hierin führt der erste Term rechts in bekannter Weise auf Formen $C_i$;
der dritte besitzt nur Formen, welche Moduln $G$ zu Factoren haben. Es
erübrigt also nur die Discussion des zweiten Termes. Nun kann $H\cdot W$
als Glied der Ueberschiebung $(B_i,\ U_1)^\lambda$ erhalten werden, eben
weil der Modul $H$ auch einer Form $B_i$ als Klammerfactor angehört
(vergl. Nr. 45), und lässt sich folglich durch solche Ueberschiebungen
darstellen. Man kann daher statt des zweiten Termes auch schreiben:
\[ \sum \sum \bigl((B_i,\ \bar U_1)^{\lambda_k},\ \Phi(B_i)\bigr)^{\varrho_i}. \]

Die Glieder dieser Ueberschiebungen entstehen durch Faltung aus Producten
\[ B_i \cdot \bar U_1 \cdot \Phi(B_i), \]
also aus Producten, die eine Form $B_i$ mehr als das ursprüngliche $\bar
V$ und demnach eine Form $A_i$ weniger als das ursprüngliche Product
$\bar U$ in der Gleichung $(\Gamma,\ \Gamma')^\nu = \sum (\bar U,\ \bar
V)^{\varrho_i}$ enthalten. Ersetzen wir nun das gegenüber $\bar U$
reducirte Product $\bar U_1$ durch seinen Werth $F_1(A_i) + \sum H\cdot
W_1 + \sum G\cdot W_1'$, wobei in jedem dieser drei Terme eine Form $A_i$
weniger betheiligt ist als in (1), so gelangen wir genau durch dieselben
Schlüsse wie bisher zu einer abermaligen Reduction des zweiten Termes in
(3), indem derselbe hiedurch auf Producte zurückgeführt wird, die aus
\[ B_k \cdot B_i \cdot \bar U_2 \cdot \Phi(B_i) \]
durch Faltung entstehen. Hierin besitzt das Product $\bar U_2$ abermals
einen Factor $A_i$ weniger als $\bar U$. Durch fortgesetzte Wiederholung
dieses Verfahrens wird schliesslich der Ausdruck
%CONT>>


% ===== pdf idx 245, printed 231.0 =====
% !PRINTED 231
\weberpage{231}
%<<CONT
\[ \sum \bigl(H\cdot W,\ \Phi(B_i)\bigr)^{\varrho_i} \]
in (3) auf eine Summe von symbolischen Producten reducirt, die nur mehr
durch Faltung von Producten der $B$, also durch Faltung einer Grösse
$\bar V$ entstehen. Dies liefert aber nach (2):
\[ \bar V = \Phi(B_i) + \sum G\cdot W'', \]
oder, weil die Formen $B_i$ zu den Formen $C_i$ gehören:
\[ \bar V = \Phi(C_i) + \sum G\cdot W''. \]

Somit ist auch dieser Term $\equiv \Phi(C_i)$ mod. $G$ und damit die
ganze rechte Seite der Gleichung (3), was zu beweisen war.

§ 21. \gesperrt{Beweis für die Endlichkeit des Formensystems einer
binären Form.}

205. \textit{Gedankengang des Beweises.} Der soeben bewiesene Satz
bildet nun die Grundlage des von Gordan zum ersten Male im 69. Bande des
Borchardt'schen Journals bewiesenen Fundamentalsatzes:

\glqq Jede binäre Form besitzt ein endliches System von Covarianten,
durch welche sich alle unendlich vielen übrigen rational und ganz
darstellen lassen.\grqq{}

Ehe wir hier zum Beweise schreiten, will ich versuchen, dessen
Gedankengang klar zu legen.

Denken wir uns sämmtliche Covarianten von $f=a_x^n$ gefunden und in den
Symbolen $a,\ b,\ c\ldots$ der Form $f$ dargestellt. Jedes dieser
symbolischen Producte können wir nach § 1, Nr. 11 so umformen, dass die
in ihm zumeist auftretende Potenz eines Klammerfactors, etwa des
Factors $(ab)$, eine gerade ist. Ist dann $g$ die grösste in
$\dfrac{n}{2}$ enthaltene ganze Zahl, so können wir sämmtliche
Covarianten von $f$ in folgende $g$ Klassen eintheilen, wobei die
späteren Klassen alle Covarianten der früheren umfassen. Die erste
Klasse enthält alle Covarianten mit dem zumeist auftretenden
Klammerfactor $(ab)^0$; wir bezeichnen sie als System $(A^{(0)})$ und
erkennen, dass es nur aus der Originalform $f=a_x^n$ allein bestehen
kann.

Die zweite Klasse enthält alle Covarianten, deren zumeist auftretender
Klammerfactor höchstens $(ab)^2$ ist; wir bezeichnen sie als System
$(A^{(1)})$.

Das System $(A^{(2)})$ wird sodann alle Covarianten umfassen, deren
zumeist auftretender Klammerfactor höchstens $(ab)^4$ ist etc., und
endlich das System $(A^{(g)})$ alle Covarianten von $f$ überhaupt, da ja
keine
%CONT>>


% ===== pdf idx 246, printed 232.0 =====
% !PRINTED 232
\weberpage{232}
%<<CONT
Covariante einen höheren Klammerfactor als $(ab)^{2g}$ besitzen kann,
insoferne eben $2g=n$, oder $n-1$ ist. \gesperrt{Jedes solche System
$(A^{(k)})$ ist relativ vollständig in Bezug auf den Modul $(ab)^{2k+2}$,
und alle höheren Potenzen von $(ab)$.} Denn durch Faltung eines
Productes von Formen $A_i^{(k)}$ können nur symbolische Producte
entstehen, deren zumeist auftretender Klammerfactor $(ab)^{2\varrho}$,
$\varrho \leqq k$, ist, und symbolische Producte mit dem Factor
$(ab)^{2\varrho}$, $\varrho > k$.

Da nun aber das System $(A^{(k)})$ sämmtliche Covarianten enthalten
soll, deren zumeist auftretender Klammerfactor höchstens $(ab)^{2k}$
ist, so kann die erste Sorte von Producten nur wieder auf Formen
$A_i^{(k)}$ führen, und es ist demnach in der That:
\[ \bar U^{(k)} \equiv F(A_i^{(k)}) \ \text{mod. } (ab)^{2k+2},\ (ab)^{2k+4}\ \text{etc.,} \]
und insbesondere
\[ \bar U^{(g)} \equiv F(A_i^{(g)}) \ \text{mod. } (ab)^{2g+1}, \]
wobei $(A^{(g)})$ das volle System von $f=a_x^n$.\footnote{*)
\textit{Beispiel.} Bei Formen vierter Ordnung besteht das System
$(A^{(0)})$ aus der Form
\[ f = a_x^4; \]
das System $(A^{(1)})$ aus den Formen:
\[ f,\ \Delta,\ t,\ j; \]
das System $(A^{(2)})$ aus den Formen:
\[ f,\ \Delta,\ t,\ j,\ i, \]
und es ist:
\[
\begin{aligned}
\bar U^{(0)} &\equiv F(f) + \sum (ab)^2\cdot W\\
\bar U^{(1)} &\equiv F(f,\ \Delta,\ t,\ j) + \sum (ab)^4\cdot W \qquad\text{(vergl. Nr. 202)}\\
\bar U^{(2)} &\equiv F(f,\ \Delta,\ t,\ j,\ i).
\end{aligned}
\]}

Hiebei ist das erste System $(A^{(0)})$ endlich; denn es besteht aus der
Form $f$ allein. Gelingt es uns nun durch Ueberschiebung des Systemes
$(A^{(0)})$ mit einem andern gleichfalls endlichen Systeme $(B^{(0)})$
das System $(A^{(1)})$ zu erhalten, und weiter durch Ueberschiebung des
Systemes $(A^{(1)})$, das nun nothwendig nach Nr. 199 gleichfalls
endlich ist, mit einem endlichen System $(B^{(1)})$ das System
$(A^{(2)})$ zu erzeugen etc. etc., so ist der Beweis der Endlichkeit
geliefert. Solche brauchbare Systeme $(B^{(k)})$ liefert aber sofort der
bereits in Nr. 43 geschilderte Process, und ich gehe nun dazu über,
diese Bildsysteme aufzustellen.

206. \textit{Aufstellung der Bildsysteme $(B^{(k)})$.} Es sei $A_i^{(k)}$
irgend eine Form des Systemes $(A^{(k)})$; sie sei durch eine gewisse
Anzahl Faltungen
%CONT>>


% ===== pdf idx 247, printed 233.0 =====
% !PRINTED 233
\weberpage{233}
%<<CONT
aus einer Potenz $f^\varrho = a_x^n \cdot b_x^n \cdot c_x^n \cdots r_x^n$ gewonnen worden. Nehmen
wir nun \gesperrt{genau dieselben} Faltungen an der $\varrho^{\mathrm{ten}}$ Potenz der Covariante
$\varphi = (ab)^{2k+2} a_x^{n-2k-2} b_x^{n-2k-2}$ vor, was (vergl. Nr. 43) auf mehrere
Weisen möglich ist, so entsteht wiederum eine Covariante von $f$, die
wir mit $B_i^{(k)}$ bezeichnen. Wir sagen: die Form $B_i^{(k)}$ ist nach dem
\gesperrt{Modell} $A_i^{(k)}$ aus der Covariante $\varphi = (f,f)^{2k+2}$ gebildet worden. Auf
dieselbe Weise können wir nach jeder Form des Systemes $(A^{(k)})$ ein
Bild herstellen, so lange der Grad der Covariante $\varphi$ grösser oder
gleich $n$ ist, und wir nennen alsdann die Gesammtheit dieser Formen
ein \gesperrt{System} $(B^{(k)})$. Nur wenn der Grad $2n - 4k - 4$ dieser Form
$\varphi$ kleiner wird als $n$, kann diese Operation unausführbar werden, da
alsdann $\varphi^\varrho$ möglicher Weise nicht mehr die nöthige Anzahl Factoren
erster Art besitzt, durch deren Faltung das Bild $B_i^{(k)}$ nach dem Muster
von $A_i^{(k)}$ entstehen soll. In diesem Falle aber nehmen wir alsdann
die vollen Systeme der Formen niedrigerer Ordnung, also der Formen
$(n-1)^{\mathrm{ten}}$, $(n-2)^{\mathrm{ten}}$ etc. Grades zum Muster für die Bildung des
Systemes $(B^{(k)})$. Da wir die Endlichkeit des Systemes einer Form
$n^{\mathrm{ter}}$ Ordnung beweisen wollen, so dürfen wir die vollen Systeme der
Formen niedrigerer Ordnung als endliche annehmen, wie wir dies ja
auch in der That bei Formen ersten bis vierten Grades gesehen haben.

Um zu wiederholen: Ist $k \leqq \frac{n}{4} - 1$, so stellen wir das System
$(B^{(k)})$ nach dem Modelle der vorhin Nr. 205 definirten Systeme $(A^{(k)})$
an der Form $\varphi = (f,f)^{2k+2}$ her. Alle spätern Systeme $(B^{(k)})$ haben
die vollen Systeme $(\overline{A}^{(k)})$ der Formen niedrigeren Grades zum Muster,
die wir als endliche Systeme voraussetzen. Die Systeme $(A^{(k)})$ von
$k = 0$ bis $k \leqq \frac{n}{4} - 1$ genügen, wie wir in Nr. 205 gesehen, der
Congruenz:
\weq{1}{\overline{U}^{(k)} \equiv F(A_i^{(k)}) \bmod (ab)^{2k+2},\ (ab)^{2k+4},\ \text{etc.};}
die späteren Systeme $(\overline{A}^{(k)})$ aber, als endliche Systeme, der Gleichung
\weq{2}{\overline{U}^{(k)} = F(\overline{A}_i^{(k)}).}
207. \textit{Eigenschaft der Bildsysteme $(B^{(k)})$.} Während so die Modellsysteme
, wenn sie nicht a priori endlich sind, der Congruenz (1) genügen
, befriedigen sämmtliche Bildsysteme die Congruenz:
\weq{3}{\overline{V}^{(k)} \equiv F(A_i^{(k)}) \bmod (ab)^{2k+4},\ (ab)^{2k+6} \ldots} % ?? uncertain: printed symbol at start of this equation is obscured by what looks like a handwritten correction mark over the scan; read as $\overline{V}^{(k)}$ by analogy with the parallel formula on p. 234
Dies erkennt man für die Systeme $(B^{(k)})$, $k > \frac{n}{4} - 1$, unmittelbar.
Denn ihre Modellsysteme genügen der Gleichung (2), d. h. zwischen
%CONT>>


% ===== pdf idx 248, printed 234.0 =====
% !PRINTED 234
\weberpage{234}
%<<CONT
den Formen des Systemes und einer Form $\overline{U}^{(k)}$, die aus ihnen durch
Faltung entsteht, besteht die Relation:
\[
\overline{U}^{(k)} - F(A_i^{(k)}) = 0.
\]
Jedes Bild derselben ausgeführt an $\varphi = (f,f)^{2k+2}$ genügt alsdann
nach Nr. 43, (3) der Relation:
\[
\overline{V}^{(k)} - F(B_i^{k}) \equiv 0 \bmod (ab)^{2k+4},
\]
d. h. es ist:
\[
\overline{V}^{(k)} \equiv F(B_i^{k}) \bmod (ab)^{2k+4}.
\]
Für die übrigen Systeme $(B^{(k)})$ von $k = 0$ bis $k \leqq \frac{n}{4} - 1$ ergiebt
 sich die Richtigkeit unserer Behauptung auf folgendem Wege.
Wir schreiben wieder die Congruenz (1), welcher ihre Modelle genügen,
in der Gleichungsform:
\weq{4}{\overline{U}^{(k)} - F(A_i^{(k)}) - \sum{}'(ab)^{2k+2} W - \sum{}'(ab)^{2k+4} W' - \cdots = 0,}
und erkennen in ihr eine Relation von symbolischen Producten, deren
Glieder sich weder gegenseitig aufheben, noch einzeln verschwinden.
Von einer solchen Relation haben wir in Nr. 43 gezeigt, dass ihr
Bild hergestellt an der Form $\varphi = (f,f)^{2k+2}$ sich darstellen lässt durch
symbolische Producte, die mindestens $(ab)^{2k+4}$ als Klammerfactor besitzen
.

Nun ist das Bild von $\overline{U}^{(k)}$ die Form $\overline{V}^{(k)}$, das Bild von $F(A_i^{(k)})$
die Function $F(B_i^{(k)})$; dabei seien die Formen $B_i^{(k)}$, die ja auf verschiedene
 Arten nach dem Modell $A_i^{(k)}$ hergestellt werden können, so
gewählt, dass die Bilder des dritten Termes $\sum{}'(ab)^{2k+2} W$, und auch
der späteren, die Form annehmen
\weq{5}{\sum{}'(ab)^{2k+2}(aa_1)^{2k+2}\overline{W},\ \text{etc.};}
dies ist immer möglich, weil ebenso wie der dritte Term in (4) durch
mindestens $(2k+2)$-malige Faltung aus $f^\varrho = a_x^n b_x^n c_x^n \cdots r_x^n$, auch das
Bild durch mindestens $(2k+2)$-malige Faltung aus
\[
\varphi^\varrho = (ab)^{2k+2} a_x^{n-2k-2} b_x^{n-2k-2} \cdot (a_1 b_1)^{2k+2} a_{1x}^{n-2k-2} b_{1x}^{n-2k-2}
\]
entstehen muss. Allein alle symbolischen Producte, wie sie in der
Summe (5) auftreten, lassen sich durch andere Producte darstellen,
die mindestens den Factor $(ab)^{2k+4}$ besitzen, wie wir in Nr. 98 ausführlich
 gezeigt haben. Es ist also für $k \leqq \frac{n}{4} - 1$ zunächst:
%CONT>>


% ===== pdf idx 249, printed 235.0 =====
% !PRINTED 235
\weberpage{235}
%<<CONT
\[
\begin{gathered}
\overline{V}^{(k)} - F(B_i^{(k)}) - \sum{}'(ab)^{2k+2}(aa_1)^{2k+2}\overline{W} - \sum{}'(ab)^{2k+4}(aa_1)^{2k+4}\overline{W}' \cdots \\
\equiv 0 \bmod (ab)^{2k+4},
\end{gathered}
\]
oder weil vom dritten Terme links incl. $ab$ alle Producte den nämlichen
 Modul $(ab)^{2k+4}$ als Factor besitzen:
\[
\overline{V}^{(k)} \equiv F(B_i^{(k)}) \bmod (ab)^{2k+4},\ (ab)^{2k+6}\ \text{etc.}
\]
208. \textit{Rückblick.} Wir haben so die Gesammtheit aller Covarianten
von $f$ in doppelter Weise eingetheilt: Einmal in Systeme $(A^{(k)})$ relativ
vollständig modulo $(ab)^{2k+2}$, $(ab)^{2k+4}$, etc. Dann in Systeme $(B^{(k)})$,
welche die gleiche Anzahl von Formen enthalten, wie das
entsprechende Modell $(A^{(k)})$, welche ferner relativ vollständig sind
modulo $(ab)^{2k+4}$, $(ab)^{2k+6} \ldots$ etc., und welche endlich jedesmal mindestens
eine Form enthalten, die den niedrigsten Modul $(ab)^{2k+2}$ des Modellsystemes
 $A^{(k)}$ zum Factor hat. Denn $(B^{(k)})$ hat immer die jeweilige
Form $\varphi = (f,f)^{2k+2}$ als Originalform. Je zwei Systeme $(A^{(k)})$ und
$(B^{(k)})$ sind also gerade zwei Systeme, wie sie der vierte Lehrsatz
Nr. 204 voraussetzt und durch Ueberschiebung derselben entsteht
sonach ein relativ vollständiges Covariantensystem von $f$, das als
niedrigsten Modul den Modul $(ab)^{2k+4}$ besitzt, d. h. das System
$(A^{(k+1)})$ der Form $f$. Und darin liegt der Hauptschritt des Beweises für
den an die Spitze dieses Paragraphen gestellten Lehrsatzes.

209. \textit{Beweis der Endlichkeit des Formensystemes.} Nachdem wir
nun gesehen haben, wie durch die Hilfssysteme $(B^{(k)})$ aus einem System
$(A^{(k)})$ das nächst höhere System $(A^{(k+1)})$ vermittelst Ueberschiebung
erzeugt wird, erübrigt uns nur mehr, dass wir uns von der Endlichkeit
 der einzelnen relativ vollständigen Systeme $(A^{(k)})$ überzeugen.
Daraus folgt dann von selbst die Endlichkeit des Systemes $(A^{(\varrho)})$,
d. i. des vollen Systemes von $f$. Nun besteht aber das System $(A^{(0)})$
aus der Form $f$ allein und demnach das Bildsystem $(B^{(0)})$ aus
$\varphi = (ab)^2 a_x^{n-2} b_x^{n-2}$ allein; durch ihre Ueberschiebungen entsteht das
System $(A^{(1)})$, das somit endlich ist, weil es die Systeme $(A^{(0)})$ und
$(B^{(0)})$ waren. Nach dem Muster $(A^{(1)})$ bilden wir $(B^{(1)})$; es enthält
dieselbe endliche Anzahl von Formen wie $(A^{(1)})$ und durch Ueberschiebung
 beider entsteht das endliche System $(A^{(2)})$ etc. etc. So gelangen
 wir schliesslich zum vollen System $(A^{(\varrho)})$, von dem wir nun
erkennen, dass es in der That endlich sein muss. In diesem Systeme
$(A^{(\varrho)})$ treten freilich im Allgemeinen noch viele überflüssige Formen
auf, d. h. solche, die sich rational und ganz durch die übrigen des
Systemes ausdrücken lassen. Dieselben zu erkennen bleibt immer eine
der Hauptaufgaben bei der wirklichen Systembildung. Sollte eine
%CONT>>


% ===== pdf idx 250, printed 236.0 =====
% !PRINTED 236
\weberpage{236}
%<<CONT
Form $f = a_x^n$ so beschaffen sein, dass $(f,f)^{2k+2+2\varrho}$, $\varrho = 1, 2, \ldots$
identisch verschwindet, dann bildet natürlich bereits das System
$(A^{(k+1)})$ das vollständige System von $f$. (Vergl. die Systeme der
regulären Körper \S\ 19).*)

\begin{center}
\textbf{\S\ 22. Das Formensystem der Form fünfter Ordnung.}
\end{center}

210. \textit{Aufstellung der Systeme $(A^{(k)})$ und $(B^{(k)})$.} Die eben entwickelte
 Theorie wollen wir nunmehr durch die Aufstellung des Formensystemes
 einer Form fünfter Ordnung erläutern. Das Formensystem
$A^{(0)}$ ist dargestellt durch die Form
\[
(A^{(0)}) \qquad f = a_x^5;
\]
das Formensystem $(B^{(0)})$ durch die Form
\[
(B^{(0)}) \qquad \varphi = (ab)^2 a_x^3 b_x^3 = (f,f)^2.
\]
Die Ueberschiebung beider Systeme liefert das System $(A^{(1)})$. Die
Formen dieses Systemes sind also enthalten in dem Ausdruck
\[
U = \left(f^\alpha,\ ((f,f)^2)^\beta\right)^\lambda.
\]
Es handelt sich nun zunächst darum, jene Ueberschiebungen auszuscheiden
, welche zerfallende oder reducible Glieder enthalten. Hieher
gehören in erster Linie alle jene, für welche $\lambda > 1$ ist, welche Werthe
auch $\alpha > 0$, $\beta > 0$ haben mögen. Denn diese Ueberschiebungen enthalten
 Glieder mit den Klammerfactoren $(ab)^2 (ac)^2$. Solche symbolische
Producte sind aber bereits (nach Nr. 98) $\equiv 0 \bmod (ab)^4$.

$\lambda$ kann also nur den Werth 1 besitzen; in diesem Falle darf aber
wiederum $\alpha$ oder $\beta$ nicht grösser als 1 sein, da ja alsdann schon $f$
und $\varphi$ allein den zur einmaligen Faltung nöthigen Factor erster
Art liefern, während die übrigen Factoren $f$ und $\varphi$ nur multiplicativ
zum Faltungsglied hinzutreten und so dieses Glied zu einem zerfallenden
 Gliede machen würden.

Die einzige neue Form, welche sonach durch die Ueberschiebung
der Systeme $(A^{(0)})$ und $(B^{(0)})$ entsteht, ist die Form
%CONT>>

\bigskip
\noindent\rule{4cm}{0.4pt}

\noindent *) Im 100. Bande des Crelle'schen Journals hat neuerdings auch Herr
\gesperrt{Mertens} einen Beweis für die Endlichkeit des Formensystems erbracht.
Indem er voraussetzt, dass das System $S$ der Form $p = a_x^n$ endlich
sei, bildet er zunächst das simultane System von $p$ und einer linearen
Form $\alpha_{1x}$. Mit Hilfe von diophantischen Gleichungen zeigt er, dass
dieses System eine endliche Anzahl von Formen $B$ enthält, welche
in den Coefficienten von $p$ und $\alpha_1$ vom gleichen Grade sind.
Nun identificirt er die Form $f = A_x^{n+1}$ mit $p \cdot \alpha_1 = p_x^n \cdot \alpha_{1x}$, indem
er $\alpha_1$ als den einen und $p_x^n$ als das Product $\alpha_{2x}\alpha_{3x}\cdots\alpha_{n+1,x}$ der
übrigen linearen Factoren von $f$ betrachtet.


% ===== pdf idx 251, printed 237.0 =====
% !PRINTED 237
\weberpage{237}
\[
t = (f,\, (f,f)^2) = (ab)^2(ca)\, a_x^2 b_x^3 c_x^4
\]
und das System $(A^{(1)})$ besteht somit aus den drei Formen
\[
f,\ (f,f)^2,\ (f,\,(f,f)^2).
\]
Dies galt schon bei den Formen dritten und viertes Grades und % [sic]
gilt auch noch bei h\"oheren Formen. F\"ur die regul\"aren K\"orper ist
das System $(A^{(1)})$, abgesehen von einer Invariante, das volle System.

Nach diesem Modelle w\"are nun an der Form
\[
i = (ab)^4 a_x b_x = (f,f)^4
\]
das System $(B^{(1)}) \equiv 0 \bmod. (ab)^5$ zu bilden. Da jedoch $i$ bereits eine
Form zweiter Ordnung ist, so muss man das System einer Form
$f = a_x^2$ zum Muster nehmen, das nur aus den Formen $a_x^2$ und $(ab)^2$
besteht. Demnach enth\"alt $(B^{(1)})$ die beiden Formen
\[
(B^{(1)}) \qquad i,\ (i,i)^2 = (ab)^4(cd)^4(ac)(bd) = A.
\]
Dieses System ist $\equiv 0 \bmod. (ab)^5$. Durch Ueberschiebung desselben
 mit $(A^{(1)})$ erh\"alt man somit das vollst\"andige System $(A^{(2)})$ der
Form $f = a_x^5$.

211. \textit{Ueberschiebung der Systeme $(A^{(1)})$ und $(B^{(1)})$.} Unsere n\"achste
Aufgabe besteht nun wiederum darin, aus den Ueberschiebungen
$(U,V)^\lambda$ jene auszuscheiden, welche zerfallende Glieder enthalten. Zu
dem Zwecke untersuchen wir zuerst die Ueberschiebungen einzelner
Formen aus beiden Systemen, sodann die Ueberschiebungen von Producten
 $U$ mit Producten $V$. Wir beginnen mit den Ueberschiebungen:
\[
(f^\mu,\, i^\varrho)^\lambda.
\]
Ist hier $\mu = 1$, so ergeben sich von $\lambda = 1$ bis $\lambda = 5$ Ueberschiebungen
 ohne zerfallende Glieder, n\"amlich die Formen
\[
(f,i),\ (f,i)^2,\ (f,i^2)^3,\ (f,i^2)^4,\ (f,i^3)^5,
\]
wovon man sich leicht \"uberzeugt, wenn man das Product
\[
f \cdot i^\varrho
\]
bezw. 1, 2, 3, 4, 5-mal auf alle m\"oglichen Arten faltet.

Bei sechsmaliger Faltung des Productes
\[
f^2 \cdot i^3 = c_x^5 \cdot d_x^5 \cdot (ab)^4 a_x b_x \cdot (a_1 b_1)^4 a_{1x} b_{1x} \cdot (a_2 b_2)^4 a_{2x} b_{2x}
\]
entstehen dagegen zerfallende Glieder; so z.\,B. wenn man die vier ersten
Factoren $c_x$ mit den vier Factoren $a_x b_x a_{1x} b_{1x}$, und zwei Factoren $d_x$
mit $a_{2x} b_{2x}$ faltet. Das so entstehende Glied der Ueberschiebung
$(f^2, i^3)^6$ ist nichts anderes als das Product zweier Glieder niedriger
Ueberschiebungen, n\"amlich der Ueberschiebungen
\[
(f, i^2)^4 \ \text{und}\ (f, i)^2.
\]


% ===== pdf idx 252, printed 238.0 =====
% !PRINTED 238
\weberpage{238}
Ganz ebenso zeigt man, dass die Ueberschiebungen
\[
(f^2, i^4)^7,\ (f^2, i^4)^8,\ (f^2, i^5)^9
\]
zerfallende Glieder enthalten. Dagegen liefert die Ueberschiebung
\[
(f^2, i^5)^{10},
\]
wie man leicht wieder durch den Faltungsprocess erkennt, keine zerfallenden
 Glieder; sie ist eine Invariante vom Grade 12 in den Coefficienten
 von $f$.

Wenn man nun weitere Potenzen von $i$ mit $f^3, f^4, \ldots$ etc. zur
Faltung bringt, so sieht man, dass immer Glieder hiebei gebildet
werden k\"onnen, welche diese Invariante zum Factor haben. Wir erhalten
 sonach von dieser Seite aus keine neuen Formen mehr.

\textit{Die Ueberschiebungen $(\varphi^\mu, i^\varrho)^\lambda$.} Nicht zerfallende Glieder enthalten
 die Ueberschiebungen:
\[
(\varphi, i),\ (\varphi, i)^2,\ (\varphi, i^2)^3,\ (\varphi, i^2)^4,\ (\varphi, i^3)^5,\ (\varphi, i^3)^6,
\]
wie man wieder unmittelbar aus den entsprechenden Faltungsprocessen
erkennt. Hiebei liefert die letzte Ueberschiebung eine Invariante vom
Grade 8 in den Coefficienten von $f$.

Bei allen weitern Ueberschiebungen von Potenzen der Form $\varphi$
mit Potenzen von $i$ kann man sich stets durch Faltung Glieder verschaffen
, welche diese Invariante zum Factor haben, so dass die Reihe
der Formen, die hier entstehen konnten, ebenfalls abgeschlossen ist.

\textit{Die Ueberschiebungen $(t^\mu, i^\varrho)^\lambda$.} Hier enthalten alle ungeraden
Ueberschiebungen von $t$ \"uber $i, i^2, i^3, i^4, i^5$, ausgenommen die neunte
Ueberschiebung, stets reducible Glieder.

Denn setzen wir einen Augenblick
\[
i = (ab)^4 a_x b_x = i_x^2 \qquad
t = (cd)^2 (ce)\, c_x^2 d_x^3 e_x^4,
\]
so tritt bei diesen ungeraden Ueberschiebungen stets ein Glied auf
mit den Faltungen
\[
(ci)(ce)\, i_x c_x.
\]
Dieser Theil des Gliedes kann aber nach dem Productsatz in die
Form gebracht werden:
\[
(ci)(ce)\, i_x c_x = \frac{1}{2} \left\{ (ci)^2 e_x + (ce)^2 i_x^2 - (ie)^2 c_x^2 \right\},
\]
und wenn man daher die rechte Seite f\"ur die linke in das betreffende
Glied substituirt, so spaltet es sich in drei zerfallende Glieder.


% ===== pdf idx 253, printed 239.0 =====
% !PRINTED 239
\weberpage{239}
Dagegen tritt bei den geraden Ueberschiebungen kein solches
Glied auf, weil hiebei jedesmal alle Factoren erster Art $i_x$, und
keines bei der neunten Ueberschiebung, weil hier alle Factoren erster
Art $t_x$ aufgebraucht werden.

Wir erhalten sonach aus den Ueberschiebungen mit $t_x^9$ die Formen
\[
(t, i)^2,\ (t, i^2)^4,\ (t, i^3)^6,\ (t, i^4)^8,\ (t, i^5)^9.
\]
Weitere Ueberschiebungen brauchen hier nicht mehr untersucht
zu werden, da das Quadrat der Form $t$, welche ja Functionaldeterminante
 ist, durch niedrigere Formen sich ausdr\"ucken l\"asst.

212. \textit{Die Ueberschiebungen von Producten von Formen.} Es er\"ubrigt
also nur noch die Ueberschiebungen der Form $i^\varrho$ \"uber das Product
der Formen
\[
f^\alpha \cdot \varphi^\beta \cdot t^\gamma
\]
zu studiren.

Hiebei kann aber, da $i$ quadratisch ist, also bei Faltung sich
seine Linearfactoren h\"ochstens auf zwei Formen vertheilen lassen,
stets einer der Exponenten gleich null genommen werden. Von den
sich so ergebenden drei verschiedenen Producten
\[
t^\gamma \cdot \varphi^\beta,\ \ f^\alpha \cdot \varphi^\beta,\ \ f^\alpha \cdot t^\gamma
\]
k\"onnen die beiden ersten unber\"ucksichtigt bleiben, da sie je eine Form
geraden und je eine ungeraden Grades enthalten und somit stets die
Linearfactoren von $i^\varrho$ bei Faltung zun\"achst auf alle Factoren erster
Art von $\varphi$ vertheilt werden k\"onnen, so dass eine Potenz der Invariante
$(\varphi, i^3)^6$ als Factor auftritt.

Sonach bleiben nur noch die Ueberschiebungen
\[
(f^\alpha \cdot t^\gamma,\, i^\varrho)^\lambda
\]
zu studiren. Dabei muss $2\varrho$ sowohl als $\lambda$ den Grad 9 von $t$ \"ubersteigen
, da ja sonst a priori zerfallende Glieder in der Ueberschiebung
auftreten, n\"amlich die bereits in das System aufgenommenen Ueberschiebungen
 von $t$ \"uber $i^\varrho$.

Setzen wir in erster Linie $\alpha = \gamma = 1$. Die Ueberschiebungen
\[
(f \cdot t, i^5)^{10},\ (f \cdot t, i^6)^{11},\ (f \cdot t, i^6)^{12},\ (f \cdot t, i^7)^{13}
\]
enthalten Glieder, welche in Producte von Gliedern zerfallen, die bereits
betrachteten Ueberschiebungen angeh\"oren. Um nur ein Beispiel anzuf
\"uhren: die Ueberschiebung $(f \cdot t, i^6)^{11}$ enth\"alt ein Glied, in welchem
die sechs Factoren erster Art von $i^3$ mit sechs Factoren von $t$ zur
Faltung kommen, w\"ahrend von den \"ubrigen sechs Factoren f\"unf mit
$f$ gefaltet werden. Dieses Glied ist also ein Product zweier Glieder,
%CONT>>


% ===== pdf idx 254, printed 240.0 =====
% !PRINTED 240
\weberpage{240}
%<<CONT
welche den Ueberschiebungen
\[
(t, i^3)^6,\ \text{bezw.}\ (f, i^2)^5
\]
angeh\"oren.

Dagegen liefert die Ueberschiebung
\[
(f \cdot t, i^7)^{14}
\]
eine neue Invariante, die wir mit $R$ bezeichnen wollen. Ueberschiebungen
 h\"oherer Potenzen von $f$ und $t$ brauchen nicht betrachtet
zu werden; denn jede solche Ueberschiebung enth\"alt immer Glieder,
die $R$ als Factor haben.

213. \textit{Ueberblick \"uber das vollst\"andige Formensystem von $f = a_x^5$.}
Wir haben somit das ganze Formensystem der Form f\"unfter Ordnung
aufgestellt. Es besteht aus den 23 Formen:
\[
\begin{aligned}
1) &\quad f,\ \varphi = (f,f)^2,\ t = (f,\,(f,f)^2),\ i = (f,f)^4,\ A = (i,i)^2 \\
2) &\quad (f,i)^1,\ (f,i)^2,\ (f,i^2)^3,\ (f,i^2)^4,\ (f,i^3)^5 \\
3) &\quad (\varphi,i)^1,\ (\varphi,i)^2,\ (\varphi,i^2)^3,\ (\varphi,i^2)^4,\ (\varphi,i^3)^5,\ (\varphi,i^3)^6 \\
4) &\quad (t,i)^2,\ (t,i^2)^4,\ (t,i^3)^6,\ (t,i^4)^8,\ (t,i^5)^9 \\
5) &\quad (f^2,i^5)^{10},\ (f \cdot t, i^7)^{14}.
\end{aligned}
\]
Es ist freilich noch nicht gezeigt, dass nicht einige dieser Formen
durch die \"ubrigen sich rational und ganz ausdr\"ucken lassen; doch
darf die Annahme gemacht werden, dass dieses System nicht mehr
weiter reducibel ist, da alle bisherigen diesbez\"uglichen Untersuchungen
keine gegentheiligen Resultate zu Tage gef\"ordert haben.

Unter diesen 23 Formen befinden sich drei fundamentale Invarianten,
n\"amlich:
\[
(i,i)^2,\ (\varphi, i^3)^6,\ (f^2, i^5)^{10},
\]
vom beziehungsweise vierten, achten und zw\"olften Grade in den Coefficienten
. Hiezu tritt noch die Invariante
\[
(f \cdot t, i^7)^{14}.
\]
Sie ist vom Grade 18 in den Coefficienten und hat die Eigenschaft
, bei der Transformation
\[
x_1 = y_1,\quad x_2 = -y_2,
\]
deren Modul $\Delta = -1$ ist, das Zeichen zu \"andern. Denn das symbolische
 Product
\[
\bigl(f \cdot (f,\,(f,f)^2),\, i^7\bigr)^{14}
\]
enth\"alt 90 Buchstaben, 70 herr\"uhrend von $i^7$, und viermal je f\"unf
herr\"uhrend von $f$; dieselben vertheilen sich sonach auf 45 Klammerfactoren
, und da bei dieser Transformation jeder solcher Klammerfactor
sein Zeichen \"andert, so \"andert dasselbe die ganze Form, d.\,h. sie ist
eine \glqq schiefe Invariante\grqq{} (vergl. Nr. 8).


% ===== pdf idx 255, printed 241.0 =====
% !PRINTED 241
\weberpage{241}

214. \textit{Uebernommene Formen.} Die Formen $(f, f)^2$, $(f, (f, f)^2)$,
$(f, f)^4$, welche in diesem Systeme auftreten, haben wir bereits auch
im System der Form vierten Grades vorgefunden, und es ist die Frage,
ob sich unter den Covarianten der Form fünften Grades nicht auch
eine Covariante

\weq{1}{j = (ab)^2 (ac)^2 (bc)^2 a_x b_x c_x}

befindet, welche der Invariante

\weq{2}{j = (ab)^2 (ac)^2 (bc)^2}

der Form vierten Grades entspricht. Erinnern wir uns, dass die
Invariante $j$ unsymbolisch die Form hatte:

\weq{3}{j = 6 \begin{vmatrix} \overline{a_0} & \overline{a_1} & \overline{a_2} \\ \overline{a_1} & \overline{a_2} & \overline{a_3} \\ \overline{a_2} & \overline{a_3} & \overline{a_4} \end{vmatrix},}

so können wir nunmehr für die Form fünften Grades die analoge
Determinante bilden:

\weq{4}{j = 6 \begin{vmatrix} \overline{a_0} & \overline{a_1} & \overline{a_2} & \overline{a_3} \\ \overline{a_1} & \overline{a_2} & \overline{a_3} & \overline{a_4} \\ \overline{a_2} & \overline{a_3} & \overline{a_4} & \overline{a_5} \\ -x_2^3 & x_1 x_2^2 & -x_1^2 x_2 & x_1^3 \end{vmatrix}.}

Setzt man hierin statt der wirklichen Coefficienten die symbolischen,
so erhält man in derselben Weise, wie wir früher die Determinante
(3) in (2) umgeformt hatten, für die Determinante (4) das symbolische
Product (1).

Diese Covariante $j$ befindet sich aber nicht unmittelbar im
System, und es ist die nächste Frage, welche Stellung sie zu ihm
einnimmt. Zu dem Zwecke werden wir sie so umformen, dass sie
den Factor $(ac)^4$ erhält, was nach Nr. 98 stets möglich ist. Es ist:

\[ a_x (bc) = b_x (ac) - c_x (ab), \]

also

\[ j = (ab)^2 (ac)^3 (bc) b_x^2 c_x - (ab)^3 (ac)^2 (bc) b_x c_x^2. \]

Vertauscht man im zweiten Terme $b$ mit $c$, so kommt:

\[ j = 2 (ab)^2 (ac)^3 (bc) b_x^2 c_x. \]

Vertauscht man hierin $a$ mit $c$ und nimmt die halbe Summe der so
entstehenden Ausdrücke, so erhalten wir endlich

\[ j = -(ac)^4 (ab) (cb) b_x^3. \]

Das symbolische Product rechts wird aber durch die Ueberschiebung
erhalten:

\[ ((ac)^4 a_x c_x, b_x^5)^2 = (i, f)^2, \]
%CONT>>


% ===== pdf idx 256, printed 242.0 =====
% !PRINTED 242
\weberpage{242}

%<<CONT
d. h. $j$ ist bis auf den Factor $-1$ nichts anderes, als die bereits im
Systeme aufgenommene Ueberschiebung $(f, i)^2$.

Wir erkennen also die bemerkenswerthe Thatsache, dass sämmtliche
Formen

\weq{1}{f, \ \Delta, \ t, \ i, \ j}

des vollständigen Systemes vierter Ordnung ihre entsprechenden
Formen im System von $f = a_x^5$ besitzen, wobei die gegenseitige Beziehung
 durch die je zwei solchen Formen gemeinschaftlichen charakteristischen
 Klammerfactoren hergestellt ist. Wir nennen diese Formen

\weq{2}{f, \ (f, f), \ (f, (f, f)^2), \ (f, f)^4, \ -(f, i)^2}

\gesperrt{übernommene Formen}. Sie gehen aus dem Systeme (1) einfach
dadurch hervor, dass man die Formen desselben jedem Symbole entsprechend
 mit je einem Factor erster Art multiplicirt. Diese merkwürdige
 Thatsache wiederholt sich indessen später nicht mehr; das
System der Form fünften Grades kann nicht mehr in das System der
Form sechsten Grades in dieser Weise übernommen werden.

215. \textit{Zweite Definition der Form $j$.} Es möge gleich an dieser
Stelle eine die Form $j$ geradezu definirende Eigenschaft derselben
erwähnt werden. Wenn man nämlich die Ueberschiebungen von $f$ mit
$j$ bildet, so zeigt sich, dass die dritte Ueberschiebung $(f, j)^3$ identisch
verschwindet. Denn man hat:

\[ (f, j)^3 = -(ab)^3 (bi)^2 a_x^2. \]

Setzt man hierin

\[ a_x (bi) = b_x (ai) - i_x (ab), \]

so kommt:

\[ -(ab)^3 (bi) (ai) a_x b_x + (ab)^4 (bi) a_x i_x = (f, j)^3. \]

Beide Terme links verschwinden; der erste, weil er bei Vertauschung
von $a$ mit $b$ nur sein Zeichen ändert; der zweite, weil er nichts anderes
 darstellt, als die erste Ueberschiebung der quadratischen Form
$i = (ab)^4 a_x b_x$ über sich selbst. Es ist also in der That:

\[ (f, j)^3 = 0. \]

Durch diese Eigenschaft ist aber $j$ vollständig charakterisirt, d. h.
ist $\psi$ eine beliebige cubische Form, so beschaffen, dass

\[ (f, \psi)^3 = 0, \]

dann ist stets:

\[ \psi = j, \]

bis auf einen constanten Factor.

Ist nämlich

\[ (f, \psi)^3 = a_x^2 (a\psi)^3 = a_x^2 (a_1 \psi_2 - \psi_1 a_2)^3 = 0 \]
%CONT>>


% ===== pdf idx 257, printed 243.0 =====
% !PRINTED 243
\weberpage{243}

%<<CONT
so müssen die Coefficienten von $x$ einzeln verschwinden, d. h.
die Gleichungen bestehen:

\[
\begin{aligned}
a_1^2 (a_1 \psi_2 - \psi_1 a_2)^3 &= \overline{a_0}\, \overline{\psi_3} - 3 \overline{a_1}\, \overline{\psi_2} + 3 \overline{a_2}\, \overline{\psi_1} - \overline{a_3}\, \overline{\psi_0} = 0 \\
a_1 a_2 (a_1 \psi_2 - \psi_1 a_2)^3 &= \overline{a_1}\, \overline{\psi_3} - 3 \overline{a_2}\, \overline{\psi_2} + 3 \overline{a_3}\, \overline{\psi_1} - \overline{a_4}\, \overline{\psi_0} = 0 \\
a_2^2 (a_1 \psi_2 - \psi_1 a_2)^3 &= \overline{a_2}\, \psi_3 - 3 \overline{a_3}\, \overline{\psi_2} + 3 \overline{a_4}\, \overline{\psi_1} - \overline{a_5}\, \overline{\psi_0} = 0. % [sic] third row prints $\psi_3$ without the bar carried by the other rows
\end{aligned}
\]

Bestimmt man hieraus die Verhältnisse der $\psi_i$ und trägt deren
Werthe in

\[ \psi = \overline{\psi_0} x_1^3 + 3 \overline{\psi_1} x_1^2 x_2 + 3 \overline{\psi_2} x_1 x_2^2 + \overline{\psi_3} x_2^3 \]

ein, so erhält man gerade die Covariante $j$.

\textit{Anmerkung.} Alle Formen, welche den Factor $(aj)^3$ besitzen, verschwinden
 identisch, denn sie können durch Ueberschiebung mit

\[ (f, j)^3 = (aj)^3 a_x^2 \]

entstanden gedacht werden. Der Klammerfactor $(aj)^3$ ist demnach
\gesperrt{Reducent}.

216. \textit{Ersetzung der Formen des Systemes durch eines ihrer Glieder.}
Während nun die übernommenen Formen bereits einen einfachen symbolischen
 Ausdruck besitzen, der für die weiteren Rechnungen tauglich
 ist, sind die meisten der übrigen Formen des Systemes durch
Aggregat symbolischer Producte dargestellt, das die symbolische Rechnung
 sehr erschweren würde. Es ist daher angezeigt, hier von dem
Satze Gebrauch zu machen, dass sich jede Ueberschiebung durch eines
ihrer Glieder und niedrigere Ueberschiebungen darstellen lässt. Indem
wir alsdann von den niedrigeren Ueberschiebungen, als bekannten
Bildungen, absehen, können wir direct das betreffende Glied, das ja
alle charakteristischen Merkmale der Ueberschiebung besitzt, welcher
es angehört, an Stelle der Ueberschiebung selbst in das Formensystem
einfügen. Insbesondere wollen wir nun im Folgenden die quadratischen
und linearen Covarianten, sowie die Invarianten durch eines ihrer
Glieder ersetzen, soweit die betreffenden Ueberschiebungen nicht ohnehin
 schon durch ein einfaches symbolisches Product dargestellt sind.

217. \textit{Die drei quadratischen Covarianten.} Als erste quadratische
Covariante erhielten wir:

\[ i = i_x^2 = (ab)^4 a_x b_x. \]

Zu einer zweiten giebt die Hesse'sche Form $\tau$ der cubischen Covariante
$j$ Veranlassung:

\[ \tau = \tau_x^2 = (j j_1)^2 j_x j_{1x}. \]

Diese Form $\tau$ lässt sich auch noch in anderer Weise darstellen. Es ist:

\weq{1}{(f, (j, i))^3 = (a_x^5, (ji) j_x^2 i_x)^3 = a_x^2 (aj)^2 (ai) (ji).}
%CONT>>


% ===== pdf idx 258, printed 244.0 =====
% !PRINTED 244
\weberpage{244}

%<<CONT
Setzt man nach dem Identitätssatze:

\weq{2}{a_x (ji) = -i_x (aj) + j_x (ai)}

und führt diesen Werth von $a_x (ji)$ in (1) ein, so kommt, da der
erste Term wegen des Reducenten $(aj)^3$ null ist:

\[ (f, (ji))^3 = a_x j_x (aj)^2 (ai)^2. \]

Der Ausdruck rechts kann als zweite Ueberschiebung von

\[ a_x^3 (ai)^2 \quad \text{über} \quad j_x^3 \]

angesehen werden, oder auch, weil $a_x^3 (ai)^2 = -j_x^3$, als zweite Ueberschiebung
 von $-j_x^3$ über $j_x^3$. Es ist also:

\[ (f, (j, i))^3 = -(j j_1)^2 j_x j_{1x} = -\tau. \]

Die dritte quadratische Covariante endlich, die wir mit in unsere
nächsten Betrachtungen hereinziehen, ist die Functionaldeterminante
von $i$ und $\tau$, nämlich

\[ \vartheta = (i, \tau) = (i \tau) i_x \tau_x. \]

Es fragt sich nun, in welchem Zusammenhange $\tau$ und $\vartheta$ mit den
beiden andern im Systeme auftretenden quadratischen Covarianten
$(\varphi, i^2)^4$ und $(\varphi, i^3)^5$ stehen.

Was die Covariante $\tau$ betrifft, so erkennt man sie leicht als ein
Glied der Ueberschiebung $(\varphi, i^2)^4$. Denn da $j = -(ai)^2 a_x^3$, so ist

\[ \tau = (j j')^2 j_x j'_x = ((ai)^2 a_x^3, (bi')^2 b_x^3)^2 = (ab)^2 (ai)^2 (bi')^2 a_x b_x. \]

Der Term rechts kann aber durch viermalige Faltung des Productes

\[ \varphi \cdot i^2 = (ab)^2 a_x^3 b_x^3 \cdot i_x^2 \cdot i_x'^2 \]

erhalten werden, wodurch die Behauptung erwiesen ist.

Ganz ebenso ist $\vartheta$ ein Glied der Covariante $(\varphi, i^3)^5$. Denn

\[ (i, \tau) = (i_x''^2, (ab)^2 (ai)^2 (bi')^2 a_x b_x) = -(ab)^2 (ai)^2 (bi')^2 (ai'') b_x i_x''. \]

Der Ausdruck rechts geht durch fünfmalige Faltung aus dem Product

\[ (ab)^2 a_x^3 b_x^3 \cdot i_x^2 \cdot i_x'^2 \cdot i_x''^2 \]

hervor, ist also in der That ein Glied von $(\varphi, i^3)^5$.

218. \textit{Die fundamentalen Invarianten.} Wir hatten im Systeme zunächst
 drei gerade fundamentale Invarianten erhalten; es waren dies
die Formen

\[ (i, i)^2, \quad (\varphi, i^3)^6, \quad (f^2, i^5)^{10}. \]

Die erste derselben behalten wir im Folgenden bei und bezeichnen
sie mit

\[ A_{ii} = A = (i, i)^2. \]

An Stelle der zweiten setzen wir die zweite Ueberschiebung von $i$
über $\tau$, d. i. die Invariante

\[ B = (i, \tau)^2. \]


% ===== pdf idx 259, printed 245.0 =====
% !PRINTED 245
\weberpage{245}
Sie ist gleichfalls vom achten Grade in den Co\"efficienten und man erkennt sie als ein Glied der Ueberschiebung $(\varphi,\ i^3)^6$. Denn es ist
\[
(i,\ \tau)^2 = (i_x^{\prime\prime 2},\ (ab)^2(ai)^2(bi')^2 a_x b_x)^2 = (ab)^2(ai)^2(bi')^2(ai'')(bi''),
\]
ein Term, der auch durch sechsmalige Faltung aus $(ab)^2 a_x^3 b_x^3 \cdot i_x^2 i_x^{\prime 2} i_x^{\prime\prime 2}$ erhalten wird. Die dritte Invariante des Systemes ersetzen wir durch die Discriminante von $j$, oder (was das n\"amliche ist) durch die Discriminante von $\tau$; wir erhalten
\[
(\tau,\ \tau)^2 = C.
\]
Ihren Zusammenhang mit $(f^2, i^5)^{10}$ werden wir sp\"ater kennen lernen. (Vgl. Nr. 226). Da $\tau$ vom sechsten Grade ist in den Co\"efficienten von $f$, so ist $C$ vom zw\"olften Grade wie $(f^2, i^5)^{10}$.

219. \textit{Die vier linearen Covarianten.} Wir hatten im System als einfachste lineare Covariante erhalten:
\weq{1}{(f,\ i^2)^4 = (a_x^5,\ i_x^2 \cdot i_x^{\prime 2})^4 = (ai)^2(ai')^2 a_x\,.}
Wir bezeichnen sie im Folgenden mit $\alpha$, und man erkennt aus dem Factor $(ai)^2$, dass $\alpha$ betrachtet werden kann als zweite Ueberschiebung von $j = -(ai)^2 a_x^3$ \"uber $i$, d.\,h. man hat:
\[
\alpha = -(ji)^2 j_x\,.
\]
Als zweite lineare Covariante trat im Systeme die Form $(f^3, i)^5$ auf. Wir ersetzen sie durch
\weq{2}{\beta = (i,\ \alpha)\,.}
Beide Formen stimmen \"uberein, denn man hat einestheils:
\[
(i,\ \alpha) = (i_x^2,\ (i'a)^2(i''a)^2 a_x) = (ia)(i'a)^2(i''a)^2 i_x
\]
und anderntheils
\[
-(f,\ i^3)^5 = (i^3,\ f)^5 = (i_x^2 i_x^{\prime 2} i_x^{\prime\prime 2},\ a_x^5)^5 = (ia)(i'a)^2(i''a)^2 i_x\,.
\]
Die dritte lineare Covariante $\gamma$ wird von der Ueberschiebung
\[
(t,\ i^4)^8 = ((ab)^2(ca)\,a_x^2 b_x^3 c_x^4,\ i_{1x}^2 i_{2x}^2 i_{3x}^2 i_{4x}^2)^8
\]
geliefert. Wir benutzen von ihr n\"amlich das Glied
\[
\gamma = (ab)^2(ac)(ai_1)^2(bi_2)^2 b_x (ci_3)^2(ci_4)^2\,.
\]
F\"uhren wir rechts das Symbol $\alpha_x = (ci_3)^2(ci_4)^2 c_x$ ein, so geht $\gamma$ \"uber in
\[
\gamma = (ab)^2(a\alpha)(ai_1)^2(bi_2)^2 b_x\,.
\]
Das symbolische Product rechts ist aber nichts anderes als die erste Ueberschiebung von
\[
\tau = (jj')^2 j_x j'_x = (ab)^2(ai)^2(bi')^2 a_x b_x
\]
\"uber $\alpha_x$;
%CONT>>


% ===== pdf idx 260, printed 246.0 =====
% !PRINTED 246
\weberpage{246}
%<<CONT
daher k\"onnen wir die lineare Covariante $\gamma$ darstellen durch
\weq{3}{\gamma = (\tau,\ \alpha)\,.}
Was endlich die vierte lineare Covariante $(t, i^5)^9$ betrifft, welche im Systeme auftritt, so ersetzen wir sie zun\"achst durch ihr Glied:
\[
(ab)^2(ac)(ia)^2(i_1b)^2(i_2c)^2(i_3c)^2(i_4b)\,i_{4x} = G\,.
\]
Verm\"oge der Identit\"at: $(ac)i_{4x} = (i_4c)a_x - (i_4a)c_x$ geht dasselbe \"uber in:
\[
\begin{aligned}
G &= (ab)^2(ia)^2(i_1b)^2(i_2c)^2(i_3c)^2(i_4c)(i_4b)\,a_x \\
  &\quad - (ab)^2(ia)^2(i_1b)^2(i_4a)(i_4b)\cdot(i_2c)^2(i_3c)^2\,c_x = G_1 - G_2\,.
\end{aligned}
\]
Der Term $G_2$ besitzt den Factor $\alpha_x = (i_2c)^2(i_3c)^2 c_x$, w\"ahrend sein anderer Factor nichts anderes ist als die Invariante
\[
B = (\tau,\ i)^2 = ((ab)^2(ai)^2(bi_1)^2 a_x b_x,\ i_{4x}^2)^2\,.
\]
Wir k\"onnen daher von diesem Terme $G_2 = B \cdot \alpha$, als einem Producte bekannter Formen, absehen, und nunmehr $G_1$ als die einzuf\"uhrende lineare Invariante $\delta$ definiren:
\[
\delta = (ab)^2(ia)^2(i_1b)^2(i_2c)^2(i_3c)^2(i_4c)(i_4b)\,a_x\,.
\]
Ber\"ucksichtigen wir, dass $(i_2c)^2(i_3c^2)(i_4c)i_{4x} = \beta_x$, so k\"onnen wir diese Covariante $\delta$ (vom Vorzeichen abgesehen) auch darstellen durch
\[
\begin{aligned}
\delta &= (ab)^2(ia^2)(i_1b)^2(b\beta)\,a_x \\
       &= ((ab)^2(ia)^2(i_1b)^2 b_x a_x,\ \beta_x) = ((jj_1)^2 j_x j_{1x},\ \beta_x),
\end{aligned}
\]
oder endlich durch
\weq{4}{\delta = (\tau,\ \beta)\,.{}^{*)}}

220. \textit{Die schiefe Invariante.} Wir greifen aus der Ueberschiebung $(f\cdot t,\ i^7)^{14}$, welche wir als schiefe Invariante erkannt haben, jenes Glied heraus, welches durch viermalige Faltung von $a_x^5$ mit $i_1^2$, also durch Bildung von $\alpha$, sodann durch weitere achtmalige Faltung von $b_x^2 c_x^3 d_x^4$ mit $i_2^4$, also durch Bildung von $\gamma$, und endlich durch zweimalige Faltung von $\alpha \cdot \gamma$ mit $i_3$ entsteht. Dieses Glied
\[
R = (i_x^2,\ \alpha_x \gamma_x)^2 = (i_x^2,\ (\tau\alpha)\tau_x \cdot \alpha_x)^2
\]
oder
\weq{5}{R = (i\tau)(\tau\alpha)(i\alpha)}
f\"uhren wir als schiefe Invariante an Stelle von $(f\cdot t,\ i^7)^{14}$ in das System ein.

Wir k\"onnen das symbolische Product (5) noch in verschiedener Weise durch Ueberschiebung erzeugen; so ist
\weq{6}{(i\tau)(\tau\alpha)(i\alpha) = -((\tau\alpha)\tau_x,\ (i\alpha)i_x) = -(\gamma,\ \beta) = (\beta\gamma) = -(\tau\alpha)(\tau\beta)}

${}^{*)}$ \gesperrt{Clebsch} f\"uhrt in seinen \glqq Bin\"aren Formen\grqq{} als vierte lineare Covariante die Form $(\vartheta,\ \alpha)$ ein.
%CONT>>


% ===== pdf idx 261, printed 247.0 =====
% !PRINTED 247
\weberpage{247}
%<<CONT
oder
\weq{7}{(i\tau)(\tau\alpha)(i\alpha) = ((i\alpha)(i\tau)\tau_x,\ \alpha_x) = -(\delta,\ \alpha) = (\alpha\delta)}
oder
\weq{8}{(i\tau)(\tau\alpha)(i\alpha) = ((i\tau)i_x\tau_x,\ \alpha_x^2)^2 = (\vartheta,\ \alpha^2)^2\,.}

221. \textit{Tabelle der wichtigsten Formen.} Ehe wir nun zu Anwendungen \"ubergehen, wollen wir noch einmal jene Formen zusammenstellen, die wir in den folgenden Untersuchungen vorzugsweise benutzen werden. Es sind dies

1) die \"ubernommenen Formen
\[
\begin{aligned}
f &= a_x^5 \\
\varphi &= (ab)^2 a_x^3 b_x^3 \qquad t = (ab)^2(ca)\,a_x^2 b_x^3 c_x^4 \\
i &= (ab)^4 a_x b_x \qquad\ \ -j = (ai)^2 a_x^3,
\end{aligned}
\]

2) die Covarianten:
\[
\begin{aligned}
\tau &= (jj_1)^2 j_x j_{1x} = -(f,\ (j,\ i))^3 \\
\vartheta &= (i\tau)i_x\tau_x \\
\alpha &= (ai)^2(ai_1)^2 a_x = -(ji)^2 j_x \\
\beta &= (i\alpha)i_x \\
\gamma &= (\tau\alpha)\tau_x \\
\delta &= (\tau\beta)\tau_x\,,
\end{aligned}
\]

3) die Invarianten:
\[
A = (i,\ i)^2,\quad B = (i,\ \tau)^2,\quad C = (\tau,\ \tau)^2,
\]
\[
R = (\beta\gamma) = (\alpha\delta) = (\vartheta\alpha)^2 = (i\tau)(\tau\alpha)(i\alpha) = -(\tau\alpha)(\tau\beta)\,.
\]
Hiezu treten noch

4) drei Hilfsinvarianten
\[
\begin{aligned}
M &= (i\alpha)^2 = (\beta\alpha) \\
N &= (\tau\alpha)^2 = (\gamma\alpha) \\
P &= (\tau\beta)^2 = (\delta\beta),
\end{aligned}
\]
welche bezw. vom 12., 16. und 20. Grade in den Co\"efficienten sind, w\"ahrend $A$, $B$, $C$ bezw. die Grade 4, 8, 12 besitzen. Wir werden die Darstellung der Formen $M$, $N$, $P$ durch die drei Fundamentalinvarianten in den folgenden Paragraphen kennen lernen.

\begin{center}
\textbf{\S.\,23.}\\
\textbf{Relationen zwischen den In- und Covarianten der Form f\"unften Grades.}
\end{center}

222. \textit{Die Quadrate der linearen Covarianten $\beta$, $\gamma$, $\delta$.} Die im vorigen Paragraphen eingef\"uhrten linearen Covarianten $\beta$, $\gamma$, $\delta$ haben wir in Functionaldeterminantenform dargestellt; ihre Quadrate m\"ussen
%CONT>>


% ===== pdf idx 262, printed 248.0 =====
% !PRINTED 248
\weberpage{248}
%<<CONT
sich also durch einfachere Formen ausdr\"ucken lassen. In der That, erhebt man zun\"achst die Identit\"at
\[
(ii_1)\alpha_x = (\alpha i_1)i_x - (\alpha i)i_{1x}
\]
in das Quadrat, so erh\"alt man:
\[
(ii_1)^2 \cdot \alpha^2 = (\alpha i_1)^2 \cdot i + (\alpha i)^2 \cdot i_1 - 2(\alpha i_1)(\alpha i)i_{1x}i_x
\]
oder:
\weq{1}{A \cdot \alpha^2 = 2M \cdot i - 2\beta^2\,.}
Ebenso findet man unter Benutzung der Identit\"aten
\[
\begin{aligned}
(\tau\tau_1)\alpha &= (\alpha\tau_1)\tau_x - (\alpha\tau)\tau_{1x} \\
(\tau\tau_1)\beta &= (\beta\tau_1)\tau_x - (\beta\tau)\tau_{1x}
\end{aligned}
\]
die Relationen:
\weq{2}{C \cdot \alpha^2 = 2N \cdot \tau - 2\gamma^2}
\weq{3}{C \cdot \beta^2 = 2P \cdot \tau - 2\delta^2\,.}

223. \textit{Die simultanen Invarianten der vier linearen Covarianten.} Die vier linearen Formen $\alpha$, $\beta$, $\gamma$, $\delta$ liefern ferner sechs simultane Invarianten:
\[
(\beta\alpha),\quad (\gamma\alpha),\quad (\delta\beta);\quad (\delta\alpha),\quad (\delta\gamma),\quad (\gamma\beta)\,.
\]
Die drei ersten derselben haben wir bereits als Hilfsinvarianten $M$, $N$, $P$ eingef\"uhrt; die Formen $(\delta\alpha)$ und $(\gamma\beta)$ stimmen mit $-R$ \"uberein; es bleibt also nur noch die Form $(\delta\gamma)$ zu berechnen. Nun ist:
\weq{1}{(\delta,\ \gamma) = (\tau\beta)(\tau\gamma) = ((\tau\gamma)\tau_x,\ \beta) = ((\tau_1\alpha)(\tau\tau_1)\tau_x,\ \beta)\,.}
Gem\"ass dem Identit\"atssatze ist aber
\[
(\tau_1\alpha)\tau_x - (\tau\alpha)\tau_{1x} = -(\tau\tau_1)\alpha_x\,.
\]
Vertauscht man also in (1) $\tau$ mit $\tau_1$ und nimmt die halbe Summe, so kommt:
\weq{2}{(\delta\gamma) = \left(-\tfrac{1}{2}(\tau\tau_1)^2 \alpha_x,\ \beta_x\right) = \tfrac{1}{2} M \cdot C\,.}
Damit ist auch die sechste dieser simultanen Invarianten auf bekannte Formen reducirt. Es tritt an uns nunmehr die Aufgabe heran, die hiebei benutzten Hilfsinvarianten $M$, $N$, $P$ als ganze Functionen der drei Fundamentalinvarianten darzustellen. Was die letzte dieser drei Formen betrifft, so k\"onnen wir sie unmittelbar auf die beiden andern reduciren. Denn \"uberschiebt man die Relation (1) Nr.\,222 zweimal \"uber $\tau$, so findet man unmittelbar:
\weq{3}{P = \tfrac{1}{2}\left\{2BM - AN\right\}\,.}
Demnach bleibt uns nur mehr die Aufgabe, $M$ und $N$ als Functionen von $A$, $B$ und $C$ darzustellen.


% ===== pdf idx 263, printed 249.0 =====
% !PRINTED 249
\weberpage{249}

224. \textit{Berechnung der Form N.} Das Princip, das uns bei Aufstellung von Relationen leitet, lässt sich dahin formuliren: Ein und dieselbe Form in zweifacher Weise durch Ueberschiebung darzustellen. Nun war $N$ definirt durch

\weq{1}{N = (\tau\alpha)^2 = (jj_1)^2(j_1\alpha)(j\alpha).}

Es handelt sich also nur darum, das symbolische Product $(jj_1)^2(j_1\alpha)(j\alpha)$ durch Ueberschiebung anderer Formen zu gewinnen. Der symbolische Factor $(j\alpha)$ veranlasst uns, die erste Ueberschiebung von $j$ über $\alpha = -(ji)^2j_x$ zu berechnen.

Man erhält:
\[
\begin{aligned}
(j,\alpha) &= -j_x^2(jj_1)(j_1i)^2\\
&= -j_x(jj_1)(j_1i)\{j_{1x}(ji) - i_x(jj_1)\}\\
&= -j_xj_{1x}(jj_1)(j_1i)(ji) + j_xi_x(jj_1)^2(j_1i).
\end{aligned}
\]
Der erste Term rechts ist null; denn er ändert durch Vertauschung von $j$ mit $j_1$ nur sein Zeichen. Der zweite Term ist aber nichts anderes, als $-\vartheta$; denn man hat:
\[
-\vartheta = (\tau,\,i) = \big((jj_1)^2j_xj_{1x},\,i_x^2\big) = j_xi_x(jj_1)^2(j_1i).
\]
Daher hat man zunächst die Relation

\weq{2}{(j,\alpha) = -\vartheta.}

Ueberschiebt man aber diese Gleichung zweimal über sich selbst, so kommt:

\weq{3}{(j\alpha)(j_1\alpha)(jj_1)^2 = (\vartheta,\,\vartheta)^2.}

Es ist aber nach der Theorie der quadratischen Formen

\weq{4}{2(\vartheta,\,\vartheta)^2 = A_{ii}A_{\tau\tau} - A_{i\tau}^2 = AC - B^2.}

und demnach durch Comparation von (1), (3) und (4)
\[
N = \frac{1}{2}(AC-B^2).
\]

225. \textit{Hilfssatz zur Berechnung von M.} Es erübrigt noch, die Form $M$ zu berechnen; zu dem Zwecke beweisen wir zuerst die Hilfsrelation
\[
(f,\,\tau)^2 = -\left\{\frac{2}{3}Aj + i\alpha\right\},
\]
indem wir von einem der symbolischen Producte rechts, etwa von $Aj$, ausgehen und dasselbe geeignet umformen. Es ist
\[
-\frac{2}{3}Aj = -\frac{2}{3}(ab)^4(ai)(bi)\cdot j_x^3.
\]
Gemäss dem Identitätssatze ist
\[
(ab)^3j_x^3 = (aj)^3b^3 - 3(aj)^2(bj)b_x^2a_x + 3(aj)(bj)^2b_xa_x^2 - (bj)^3a_x^3,
\]
%CONT>>


% ===== pdf idx 264, printed 250.0 =====
% !PRINTED 250
\weberpage{250}

%<<CONT
also, da der erste und letzte Term rechts wegen des Reducenten $(aj)^3$ verschwinden:
\[
\begin{aligned}
-\frac{2}{3}Aj &= \left\{2(ab)(aj)^2(bj)b_x^2a_x - 2(ab)(aj)(bj)^2b_xa_x^2\right\}(ai)(bi)\\
&= 2(ab)(aj)(bj)(ai)(bi)b_xa_x\{(aj)b_x - (bj)a_x\}\\
&= 2(ab)^2(aj)(bj)j_x\cdot(ai)(bi)b_xa_x.
\end{aligned}
\]
Unter Benutzung des Productsatzes
\[
(ai)(bi)a_xb_x = \frac{1}{2}\left\{(ai)^2b_x^2 + (bi)^2a_x^2 - (ab)^2i_x^2\right\}
\]
geht das letzte symbolische Product über in
\[
\begin{aligned}
-\frac{2}{3}Aj &= (ab)^2(aj)(bj)(ai)^2b_x^2j_x\\
&\quad + (ab)^2(aj)(bj)(bi)^2a_x^2j_x - (ab)^4(aj)(bj)j_x\cdot i.
\end{aligned}
\]
Die beiden ersten Terme rechts repräsentiren die nämliche Form; der zweite Term zerfällt in $(i,\,j)^2\cdot i = -\alpha\cdot i$. Demnach wird
\[
-\frac{2}{3}Aj = 2(ai)^2(ab)^2(aj)(bj)b_x^2j_x + i\cdot\alpha.
\]
Der symbolische Factor $(ai)^2$ des ersten Termes rechts veranlasst uns, da ja $(ai)^2a_x^3 = -j_x^3$ das Symbol $j_1' = -(ai)^2a_1$, $j_2' = -(ai)^2a_2$ einzuführen. Dadurch wird
\[
\begin{aligned}
-\frac{2}{3}Aj &= -2(j'b)^2(j'j)(bj)b_x^2j_x + i\alpha\\
&= -\{(j'b)j_x - (jb)j_x'\}(j'b)(j'j)(bj)b_x^2 + i\alpha\\
&= -(j'j)^2(j'b)(bj)b_x^3 + i\alpha\\
&= \big(f,\,(jj')^2j_xj_x'\big)^2 + i\alpha,
\end{aligned}
\]
oder, weil ja $\tau = (jj')^2j_xj_x'$:

\weq{1}{(f,\,\tau)^2 = -\left\{\frac{2}{3}Aj + i\alpha\right\}.}

226. \textit{Berechnung der Form M.} Schiebt man nun diese Relation
\[
-3(a\tau)^2a_x^3 = 2Aj + 3i\alpha
\]
dreimal über die Form
\[
(j,\,i) = (ji)j_x^2i_x,
\]
so entsteht links:
\[
-3(a\tau)^2(ji)(aj)^2(ai) = +3(a\tau)(aj)^2(ai)\{(ai)(\tau j) + (ja)(\tau i)\},
\]
d.\,i., weil wegen $(aj)^3a_x^2 = 0$ der zweite Term rechts verschwindet:
\[
\begin{aligned}
&= 3(a\tau)(\tau j)(aj)^2(ai)^2 = -3(ai)^2(aj)^2(a\tau)(j\tau)\\
&= -3\big((ai)^2(aj)^2a_xj_x,\,\tau_x^2\big)^2.
\end{aligned}
\]


% ===== pdf idx 265, printed 251.0 =====
% !PRINTED 251
\weberpage{251}

Nun ist aber $-(ai)^2a_x^3 = j$; daher lässt sich diese Ueberschiebung auch schreiben

\weq{1}{= +3\big((jj_1)^2j_xj_{1x},\,\tau\big)^2 = 3(\tau,\,\tau)^2 = 3C.}

Auf der rechten Seite erhalten wir durch die dreimalige Ueberschiebung: 1) das Glied

\weq{2}{(2Aj,\,(j,\,i))^3 = 2A\cdot(j_1i)(jj_1)^2(ji) = 2A\cdot(\tau,\,i)^2 = 2A\cdot B,}

2) das Glied
\[
(3i_1\cdot\alpha,\,(ji)j_x^2i_x)^3 = (ji)(i_1j)^2(\alpha i) + 2(ji)(\alpha j)(i_1j)(i_1i).
\]
Der zweite Term rechts ist null; denn er ändert durch Vertauschung von $i$ mit $i_1$ nur sein Zeichen. Der erste Term kann aber dargestellt werden durch

\weq{3}{\begin{aligned}
(ji_1)^2(ji)(\alpha i) &= \big((ji_1)^2j_x,\,(\alpha i)i_x\big)\\
&= (-\alpha_x,\,(\alpha i)i_x)\\
&= -(\alpha i)^2 = -M.
\end{aligned}}

Die gesuchte Relation wird demnach
\[
3C = 2AB - M.
\]
Fassen wir die bisher gewonnenen Resultate zusammen, so erhalten wir:

\weq{4}{(\tau\alpha)^2 = N = \frac{1}{2}(AC-B^2)}

\weq{5}{(i\alpha)^2 = M = 2AB - 3C}

\weq{6}{(\tau\beta)^2 = P = \frac{1}{4}(9AB^2 - A^2C - 12BC).}

Die Gleichnng (5) lehrt auch, dass wir berechtigt waren, $(\tau,\,\tau)^2 = C$ an Stelle von $(f^2,\,i^5)^{10}$ einzuführen. % [sic] \textquotedbl{}Gleichnng\textquotedbl{} printed for \textquotedbl{}Gleichung\textquotedbl{}
Denn $(i\alpha)^2$ ist das Glied
\[
G = \big((ai_1)^2(ai_2)^2a_x\cdot(bi_3)^2(bi_4)^2b_x,\,i_{5x}^2\big)^2
\]
dieser Ueberschiebung $(f^2,\,i^5)^{10}$.

227. \textit{Die Ueberschiebungen der quadratischen Formen $i,\,\tau,\,\vartheta$.} Eine weitere Reihe von Relationen, deren wir später bedürfen, liefern die gegenseitigen Ueberschiebungen der quadratischen Formen des Systemes, sowie die Ueberschiebungen dieser Formen über die vier linearen Formen. Sie sind leicht zu ermitteln, da die simultanen Systeme dieser Formen niedrigern Grades nach den früheren Entwicklungen bekannt sind. Es ist überhaupt, wie wir schon bisher sehen konnten, nützlich, das Studium der Systeme höherer Formen dadurch zu vereinfachen, dass man insbesondere die Ueberschiebungen der linearen und quadratischen Formen des Systemes in den Vordergrund stellt.

Die drei quadratischen Formen des Systemes waren $i,\,\tau,\,\vartheta$.


% ===== pdf idx 266, printed 252.0 =====
% !PRINTED 252
\weberpage{252}

Da $\vartheta$ Functionaldeterminante der beiden ersten Formen ist, so haben wir nach der Theorie der quadratischen Formen unmittelbar folgende Relationen (vgl. Nr. 127):

\weq{1}{(i,\,\vartheta) = \frac{1}{2}\left\{\cdot B\cdot i - A\cdot\tau\right\}} % [sic] stray dot before B as printed

\weq{2}{(\tau,\,\vartheta) = \frac{1}{2}\left\{C\cdot i - B\cdot\tau\right\}}

\weq{2a}{\left.\begin{aligned}
(i,\,\vartheta)^2 &= A_{i\vartheta} = 0\\
(\tau,\,\vartheta)^2 &= A_{\tau\vartheta} = 0
\end{aligned}\right\}}

\weq{3}{-2\vartheta^2 = A\tau^2 - 2Bi\tau + Ci^2.}

Aus der letzten dieser Relationen geht ohne weiteres die Beziehung hervor, welche zwischen dem Quadrat der schiefen Invariante $R$ und den übrigen fundamentalen Invarianten besteht. Denn es war
\[
R = (\vartheta,\,\alpha^2)^2.
\]
Ersetzt man also in dieser Relation $x$ durch $\alpha$, so kommt:

\weq{4}{-2R^2 = AN^2 - 2BMN + CM^2.}

228. \textit{Die Ueberschiebungen von $\vartheta$ über $\alpha,\,\beta,\,\gamma,\,\delta$.} Von den Ueberschiebungen der quadratischen Formen $i,\,\tau,\,\vartheta$ will ich, soweit wir sie nicht ohnehin schon berechnet haben, noch jene Formen ermitteln, welche sich bei einmaliger Ueberschiebung der Form $\vartheta$ über $\alpha,\,\beta,\,\gamma,\,\delta$ ergeben. Alle weiteren Ueberschiebungen mögen der selbstständigen Uebung zugewiesen sein.

Um nun aber $(\vartheta,\,\alpha),\,(\vartheta,\,\beta),\,(\vartheta,\,\gamma),\,(\vartheta,\,\delta)$ zu berechnen, benutzen wir einfach die schon bei Gelegenheit der quadratischen Formen gegebene Reihenentwicklung [vgl. Nr. 125 (2)]:
\[
\vartheta_x\vartheta_y = (i\tau)i_x\tau_y - \frac{B}{2}(xy) = (i\tau)i_y\tau_x - \frac{B}{2}(yx).
\]
Wir ersetzen in diesen Gleichungen $x$ resp. $y$ der Reihe nach durch $\alpha,\,\beta,\,\gamma,\,\delta$ und erhalten nach einfacher Umformung:

\weq{1}{(\vartheta,\,\alpha) = (i\tau)(i\alpha)\tau_x - \frac{B}{2}\alpha = -(\tau\beta)\tau_x - \frac{B}{2}\alpha = -\delta - \frac{B}{2}\alpha}

\weq{2}{\begin{aligned}
(\vartheta,\,\beta) &= (i\tau)(i\beta)\tau_x - \frac{B}{2}\beta = (i\tau)(ii_1)(i_1\alpha)\tau_x - \frac{B}{2}\beta\\
&= \frac{A}{2}(\tau\alpha)\tau_x - \frac{B}{2}\beta = +\frac{A}{2}\gamma - \frac{B}{2}\beta
\end{aligned}}

\weq{3}{(\vartheta,\,\gamma) = (i\tau)(\tau\gamma)i_x + \frac{B}{2}\gamma = (i\tau)(\tau\tau_1)(\tau_1\alpha)i_x + \frac{B}{2}\gamma = -\frac{C}{2}\beta + \frac{B}{2}\gamma}

\weq{4}{(\vartheta,\,\delta) = (i\tau)(\tau\delta)i_x + \frac{B}{2}\delta = (i\tau)(\tau\tau_1)(\tau_1\beta)i_x + \frac{B}{2}\delta = +\frac{AC}{4}\alpha + \frac{B}{2}\delta.}


% ===== pdf idx 267, printed 253.0 =====
% !PRINTED 253
\weberpage{253}

In ähnlicher Weise erhält man unter Benutzung dieser vier Relationen
die Ueberschiebungen von $\vartheta$ über Producte der linearen Covarianten.

Damit will ich diese Untersuchungen vorläufig abschliessen. Wir
werden im folgenden Paragraphen noch eine Reihe anderer Ueberschiebungen
berechnen, so weit wir derselben zur Lösung der dort
gestellten Aufgabe bedürfen. Insbesondere werden wir an geeigneter
Stelle die Ueberschiebungen von $f$ über die linearen Covarianten $\alpha$, $\beta$,
$\gamma$, $\delta$ durch Formen des Systemes und durch Ueberschiebungen derselben
darstellen.

\begin{center}
\S.\,24.

\textbf{Typische Darstellung der Form fünften Grades.}
\end{center}

229. \textit{Begriff der typischen Darstellung.} Unter typischer Darstellung
einer Form $f = a_x^n$ versteht man im Allgemeinen eine derartige
Darstellung derselben, dass die Coefficienten Invarianten und
die beiden Variabeln irgend zwei Covarianten des Systemes sind. Der
neue Ausdruck für $f$ ist alsdann der unveränderliche Typus für alle
Formen, die aus $f$ durch lineare Transformation hervorgehen; und umgekehrt:
Haben zwei Formen $f$ die nämliche typische Form, so können
sie stets durch lineare Transformation in einander übergeführt werden.

Für Formen ungeraden Grades kann die typische Darstellung unter
anderm auf folgendem Wege erreicht werden. Dieselben besitzen im
allgemeinen Falle und wenn $n > 3$ stets mindestens zwei lineare
Covarianten $\alpha_x$ und $\beta_x$ mit nicht verschwindender Determinante $(\alpha\beta)$,
wie \gesperrt{Clebsch} (vgl. \glqq Binäre Formen\grqq{} \S.\,90, Seite 357) gezeigt hat, und
diese werden als neue Variable für die typische Darstellung benutzt.
Erhebt man nämlich die Identität:
\[
a_x(\alpha\beta) = \alpha_x(\alpha\beta) - \beta_x(a\alpha)
\]
auf die $n^{\mathrm{te}}$ Potenz, so kommt:
\[
a_x^n \cdot (\alpha\beta)^n = \alpha_x^n(\alpha\beta)^n - \binom{n}{1}\alpha_x^{n-1}\beta_x(\alpha\beta)^{n-1}(a\alpha) + \cdots \pm \beta_x^n(a\alpha)^n
\]
oder:
\weq{I}{(\alpha\beta)^n \cdot f = (f,\ \beta^n)^n \cdot \alpha_x^n - \binom{n}{1}(f,\ \beta^{n-1}\alpha)^n \cdot \alpha_x^{n-1}\beta_x
+ \binom{n}{2}(f,\ \beta^{n-2}\alpha^2)^n \cdot \alpha_x^{n-2}\beta_x^2 - \cdots}

Hierin sind aber die Variabeln Covarianten, die Coefficienten Invarianten,
und man hat nur noch die Aufgabe zu lösen, diese Invarianten:
\weq{II}{(f,\ \beta^n)^n,\quad (f,\ \beta^{n-1}\alpha)^n,\quad (f,\ \beta^{n-2}\alpha^2)^n \cdots (f,\ \alpha^n)^n,\quad (\alpha\beta)^n}
durch die fundamentalen Invarianten des Systemes zu ersetzen\ldots{} Man
%CONT>>


% ===== pdf idx 268, printed 254.0 =====
% !PRINTED 254
\weberpage{254}

%<<CONT
hätte auch, wie dies \gesperrt{Hermite} gethan hat, statt der linearen Formen
$\alpha$ und $\beta$ die linearen Factoren einer quadratischen Covariante des
Systemes benutzen können. Dadurch ist man freilich genöthigt, Irrationalitäten
einzuführen; aber die typischen Coefficienten sind alsdann
vom niedrigeren Grade in den Coefficienten der ursprünglichen Form.

Um nun die Form $f = a_x^5$ typisch darzustellen, benutzen wir zunächst
die beiden linearen Covarianten
\[
\alpha_x = (f,\ i^2)^4, \quad \beta_x = -(f,\ i^3)^5.
\]
Wir erhalten alsdann, weil $(\alpha\beta) = -M = 3C - 2AB$:
\[
(-M)^5 \cdot f = \alpha_x^5 \cdot (f,\ \beta^5)^5 - 5\alpha_x^4\beta_x \cdot (f,\ \beta^4\alpha)^5
+ 10\alpha_x^3\beta_x^2(f,\ \beta^3\alpha^2)^5 - \cdots - \beta_x^5(f,\ \alpha^5)^5
\]
und unsere Aufgabe besteht nun darin, die sechs Invarianten
\[
(f,\ \beta^5)^5,\quad (f,\ \beta^4\alpha)^5,\quad (f,\ \beta^3\alpha^2)^5,\quad (f,\ \beta^2\alpha^3)^5,\quad (f,\ \beta\alpha^4)^5,\quad (f,\ \alpha^5)^5
\]
durch die Invarianten $A$, $B$, $C$, $R$ des Systemes darzustellen.

230. \textit{Berechnung der Ueberschiebungen $(f, \alpha^{\varrho})^{\varrho}$.} Zu dem Zwecke
haben wir zunächst die fünf Ueberschiebungen von $f$ über $\alpha$ durch
andere Ueberschiebungen darzustellen, eine Aufgabe, die wir hier,
wenn auch nicht in voller Ausführlichkeit, erledigen wollen. Man
findet zunächst:
\weq{I}{(f,\ \alpha) = 2\left((f,\ j)^2,\ i\right)}
\weq{II}{(f,\ \alpha^2)^2 = \frac{3}{2}\,\alpha\tau - Ai\alpha + \lambda j,}
wobei $\lambda$ den Werth $B - \frac{2}{3}A^2$ besitzt. Um zur ersten Gleichung zu
gelangen, hat man in $(a\alpha)a_x^4$ die Form $\alpha_x$ durch ihren Werth
$-(ji)^2 j_x$ zu ersetzen und das so entstehende symbolische Product durch
den Identitätssatz umzuformen. Indem man dann anderntheils die Polare
%CONT>>


% ===== pdf idx 269, printed 255.0 =====
% !PRINTED 255
\weberpage{255}

%<<CONT
$(f,\ j)_y^2$ berechnet und darin $y$ durch $i$ ersetzt, kommt man durch abermalige
Anwendung derselben Identität direct zu Gleichung (I).

Complicirter gestaltet sich die Herstellung der zweiten Relation.
Man benutzt die eben gewonnene Relation (I) und überschiebt sie
nochmals über die Form $\alpha$, indem man der Einfachheit halber für
$(f,\ j)^2$ zunächst die Bezeichnung $p_x^4$ einführt. Dies liefert
\[
4(f,\ \alpha^2)^2 = 6p_x^2 i_x (p\,i)(p\alpha) i_x + 2p_x^3(p\,i)(i\alpha),
\]
oder, indem man im zweiten Terme rechts $p_x(i\alpha)$ durch $i_x(p\alpha) - \alpha_x(p\,i)$
ersetzt:
\weq{1}{4(f,\ \alpha^2)^2 = 8(p\,i)(p\alpha)p_x^2 i_x - 2(p\,i)^2 \alpha_x p_x^2.}

Man hat nun zu berücksichtigen, dass die beiden Glieder der schon
mehrmals benutzten Polare $(f,\ j)_y^2$ einander gleich sind, da ihre Differenz
wegen des sich einstellenden Factors $(aj)^3$ verschwindet, dass also
\weq{2}{p_x^3 p_y = (aj)^2 a_x^3 j_y}
\weq{3}{p_x^3 p_y = (aj) a_x^2 j_x a_y.}
Mit Hilfe dieser beiden Gleichungen lässt sich nun leicht in (1) das
Symbol $p$ wiederum eliminiren, da auch die beiden Glieder von $p_x^2 p_y^2$
in (3) identisch gleich sind, und man findet so:
\weq{4}{4(f,\ \alpha^2)^2 = 8(aj)^2(ai)(j\alpha)a_x^2 i_x - 2(aj)^2(ai)^2 a_x j_x \alpha_x \cdot}

Hierin ist der zweite Term incl. Vorzeichen nichts anderes als $+2\tau \cdot \alpha$.
Der erste Term aber geht durch die Substitution $i_x(aj) = j_x(ai) - a_x(j\,i)$
über in
\weq{5}{8(aj)^2(ai)(j\alpha)a_x^2 i_x = 8(aj)(ai)^2(j\alpha)j_x a_x^2 - 8(aj)(ai)(j\alpha)(j\,i)a_x^3,}
wobei [vgl. auch Nr. 221, und Nr. 224 (2)]
\[
(aj)(ai)^2(j\alpha)j_x a_x^2 = \left(-(ai)^2 a_x^3,\ -(j\alpha)j_x^2\right) = (j,\ \vartheta) = \frac{\tau \cdot \alpha}{2}
\]
\[
-(aj)(ai)(j\alpha)(j\,i)a_x^3 = \left(a_x^5,\ -(j\alpha)(j\,i)i_x j_x\right)^2 = (f,\ (\vartheta,\ i))^2.
\]
Demnach wird
\[
4(f,\ \alpha^2)^2 = 6\alpha\tau + 8(f,\ (\vartheta,\ i))^2.
\]
Führen wir hierin nun den in Nr. 227 (1) berechneten Werth von
$(\vartheta,\ i)$ ein und eliminiren die sich hiebei einstellende Form $(f,\ \tau)^2$ mittelst
Gleichung (1) Nr. 225, so ergiebt sich in der That die oben gegebene
Relation (II). Aus den beiden Relationen (I) und (II) gehen
nun die übrigen Ueberschiebungen direct hervor. Man überschiebt (II)
abermals über $\alpha$ und findet:
\weq{III}{(f,\ \alpha^3)^3 = \alpha \cdot \gamma - \frac{2}{3}A \cdot \alpha \cdot \beta - \lambda \cdot \vartheta}
%CONT>>


% ===== pdf idx 270, printed 256.0 =====
% !PRINTED 256
\weberpage{256}

%<<CONT
und hieraus durch weitere Ueberschiebungen
\weq{IV}{(f,\ \alpha^4)^4 = \mu \cdot \alpha + \lambda \cdot \delta}
\weq{V}{(f,\ \alpha^5)^5 = -\lambda \cdot R,}
wobei $\mu$ den Werth $\frac{1}{2}N - \frac{1}{3}AM + \lambda\frac{B}{2}$ besitzt.

231. \textit{Berechnung der typischen Coefficienten.} Die soeben berechneten
Ueberschiebungen gestatten uns nun in Verbindung mit den in
Nr. 222 gegebenen Relationen für die Quadrate der Functionaldeterminanten
$\beta^2$, $\gamma^2$, $\delta^2$ in einfacher Weise nicht nur die typischen Coefficienten,
sondern überhaupt alle Ueberschiebungen von $f$ über ein Product
$\alpha^{\varkappa}\beta^{\lambda}\gamma^{\mu}\delta^{\nu}$, $\varkappa + \lambda + \mu + \nu = 5$, durch die Invarianten des
Systemes darzustellen. Insbesondere ist:
\[
(f,\ \alpha^5)^5 = -\lambda \cdot R
\]
\[
(f,\ \alpha^4\beta)^5 = \left((f,\ \alpha^4)^4,\ \beta\right) = -\mu \cdot M + \lambda P
\]
\[
(f,\ \alpha^3\beta^2)^5 = \left((f,\ \alpha^3)^3,\ \beta^2\right)^2 = (f,\ \alpha^3 \cdot i)^5 M - (f,\ \alpha^5)^5 \frac{A}{2}
\]
\[
{} = +R\left(M + \frac{\lambda}{2}A\right), \quad \text{[weil: } (f,i)^2 = -j,\ -(j,\alpha)^3 = (\vartheta,\alpha^2)^2 = R\text{]}
\]
\[
(f,\ \alpha^2\beta^3)^5 = \left((f,\ \alpha^2)^2,\ \beta\left(Mi - \frac{A\alpha^2}{2}\right)\right)^3
\]
\[
{} = (f,\ \alpha^2\beta i)^5 M - (f,\ \alpha^4\beta)^5 \frac{A}{2}, \quad \text{[weil: } (-j,\alpha^2\beta)^3 = (\vartheta,\alpha\beta)^2\text{]}
\]
\[
{} = M\left(\frac{A}{2}N - \frac{B}{2}M\right) - \frac{A}{2}(\lambda P - \mu M), \quad \text{[vgl. Nr. 228 (1)]}
\]
\[
(f,\ \alpha \cdot \beta^4)^5 = -AMR - \frac{A^2}{4} \cdot \lambda \cdot R
\]
\[
(f,\ \beta^5)^5 = (f,\ \beta i^2)^5 M^2 - (f,\ \beta\alpha^2 i)^5 A \cdot M + (f,\ \beta\alpha^4)^5 \frac{A^2}{4}, \quad [-(j,i)^2 = \alpha]
\]
\[
{} = -M^3 + AM\left(\frac{A}{2}N - \frac{B}{2}M\right) + \frac{A^2}{4}(\lambda P - \mu M).
\]

232. \textit{Typische Darstellung der Covariante $j$.} Ehe wir die typischen
Darstellungen der Form $f$ weiter verfolgen, wollen wir auch für die
Covariante $j$ eine solche geben, indem wir hiezu die linearen Covarianten
$\alpha$ und $\delta$ benutzen. Nach dem binomischen Lehrsatze erhält
man:
\[
(\alpha\delta)^3 j_x^3 = (j\delta)^3 \alpha_x^3 - 3(j\delta)^2(j\alpha)\alpha_x^2\delta_x + 3(j\delta)(j\alpha)^2\alpha_x\delta_x - (j\alpha)^3\delta_x^3.
\]
Hierin ist:

1) $(j\alpha)^3 = ((j,\alpha),\alpha^2)^2 = (-\vartheta,\alpha^2)^2 = -R$ (vgl. Nr. 224 (2), 220 (8)),

2) $(j,\alpha^2\delta)^3 = ((j,\alpha),\alpha\delta)^2 = (-\vartheta,\alpha\delta)^2 = (-(\vartheta,\alpha),\delta)$
\[
{} = \left(\delta + \frac{B}{2}\alpha,\ \delta\right) = +\frac{B}{2} \cdot R \quad \text{(vgl. Nr. 228 (1)),}
\]

3) $(j,\alpha\delta^2)^3 = (-\vartheta,\delta^2)^2 = \left(-\vartheta,\ P\tau - C \cdot \frac{\beta^2}{2}\right)^2 \quad \text{(vgl. Nr. 222 (3)),}$
\[
{} = -P \cdot (\vartheta,\tau)^2 + \frac{C}{2}(\vartheta,\beta^2)^2,
\]
oder, weil nach Nr. 227 (2a) die Form $(\vartheta,\tau)^2 = 0$, so kommt:
\[
(j,\ \alpha\delta^2)^3 = \frac{C}{2}\left((\vartheta,\beta),\beta\right) = \frac{C}{2}\left(\frac{A}{2}\gamma - \frac{B}{2}\beta,\ \beta\right) = -\frac{CA}{4}R,
\]

4) $(j,\delta^3)^3 = \left(j,\ \delta\left(P\tau - \frac{C}{2}\beta^2\right)\right)^3 = P \cdot (j,\delta\tau)^3 - \frac{C}{2}(j,\delta\beta^2)^3,$

oder, weil $(j,\tau)^2$ gemäss der Theorie cubischer Formen verschwindet:
\[
(j,\ \delta^3)^3 = -\frac{C}{2}\left(j,\ \delta\left(Mi - \frac{A\alpha^2}{2}\right)\right)^3 = +\frac{CR}{2}\left(M + \frac{AB}{4}\right).
\]


% ===== pdf idx 271, printed 257.0 =====
% !PRINTED 257
\weberpage{257}
Die typische Darstellung von $j$ wird also, da $(\alpha\delta) = +\mathbf{R}$, nach
Division mit $\mathbf{R}$

\weq{I}{
\mathbf{R}^2\cdot j = \delta^3 + \frac{3\mathbf{B}}{2}\cdot\delta^2\alpha + \frac34\,\mathbf{A}\mathbf{C}\cdot\delta\alpha^2 + \frac{\mathbf{C}}{2}\left(\mathbf{M}+\frac{\mathbf{A}\mathbf{B}}{4}\right)\cdot\alpha^3.
}

Nehmen wir an, die cubische Form $j$ besitze, gleich null gesetzt, drei
verschiedene Wurzeln $m_1\,m_2\,m_3$, was ja im Allgemeinen auch der
Fall sein wird, da sonst die Invariante $\mathbf{C} = (\tau,\tau)^2$ des Systems
verschwinden würde, so können wir schreiben:

\weq{II}{
\mathbf{R}^2\cdot j = (\delta - m_1\alpha)(\delta - m_2\alpha)(\delta - m_3\alpha).
}

233. \textit{Darstellung der Form $f$ durch drei fünfte Potenzen.} Eine
der wichtigsten typischen Darstellungen der Form $f$ ist die typische
Darstellung durch fünfte Potenzen. Der Umstand, dass

\weq{1}{
(f,\,j)^3 = 0,
}

zeigt, dass es immer möglich ist, $f$ in die Gestalt

\[
f = k_1\,p_x^5 + k_2\,q_x^5 + k_3\,r_x^5
\]

zu bringen, wo $p_x$, $q_x$, $r_x$ die drei linearen als von einander
verschieden angenommenen Factoren der cubischen Form $j$ sind. Denn
nach den Ueberschiebungsgesetzen (vergl. Nr. 38) ist stets:

\[
(k_1p_x^5 + k_2q_x^5 + k_3r_x^5,\ p_xq_xr_x)^3 = 0,
\]

und diese Gleichung ist ja gerade die Definitionsgleichung der Co\-
variante $j = p_xq_xr_x$ (vgl. Nr. 215).

Legen wir $j$ in der typischen Darstellung (II) zu Grunde, dann ist

\[
p = \delta - m_1\alpha
\]
\[
q = \delta - m_2\alpha
\]
\[
r = \delta - m_3\alpha,
\]

und demnach die typische Form von $f$

\weq{2}{
f = k_1(\delta - m_1\alpha)^5 + k_2(\delta - m_2\alpha)^5 + k_3(\delta - m_3\alpha)^5.
}

Die Berechnung der Constanten $k_1\,k_2\,k_3$ kann in folgender Weise
vorgenommen werden.

Wir überschieben Gleichung (2) fünfmal über $\alpha^5$, resp. $\alpha^4\delta$,
$\alpha^3\delta^2$ und erhalten:

\weq{3}{
\left.
\begin{aligned}
(f,\ \alpha^5)^5 &= -\ (k_1 + k_2 + k_3)\ \mathbf{R}^5\\
(f,\ \alpha^4\delta)^5 &= -\ 5(k_1m_1 + k_2m_2 + k_3m_3)\ \mathbf{R}^5\\
(f,\ \alpha^3\delta^2)^5 &= -\ 10(k_1m_1^2 + k_2m_2^2 + k_3m_3^3)\ \mathbf{R}^5
\end{aligned}
\right\}
}
% [sic] letzter Exponent m_3^3 so gedruckt (statt m_3^2, vgl. die beiden vorigen Glieder)

Die Werthe der linken Seite dieser drei Gleichungen ergeben sich
unmittelbar aus den Gleichungen (III), (IV), (V) Nr. 230, und da die
Determinante dieses Systemes (3) von linearen Gleichungen nicht
%CONT>>


% ===== pdf idx 272, printed 258.0 =====
% !PRINTED 258
\weberpage{258}
%<<CONT
verschwindet (wegen $(\tau_1\tau)^2 \gtrless 0$), so lassen sich aus ihm die drei
Constanten $k_1$, $k_2$, $k_3$ eindeutig berechnen.

234. \textit{Typische Darstellung von $f$, wenn $\mathbf{C} = 0$.} Sind nun aber die
drei Wurzeln der Covariante $j = 0$ nicht von einander verschieden,
sind vielmehr zwei derselben einander gleich, dann ist die Darstellung
durch fünfte Potenzen allein nicht mehr möglich. Immer noch aber
liefert die typische Darstellung von $f$ unter Zugrundelegung der Linear\-
factoren von $j$ ein sehr einfaches Resultat, wovon man sich sofort
überzeugen kann. In diesem Falle ist nämlich:

\weq{1}{
j = p_x^2\,q_x
}

und demnach

\weq{2}{
(f,\,j)^3 = (ap)^2(aq)\,a_x^2 = 0.
}

Erhebt man also die Identität

\[
a_x(pq) = p_x(aq) - q_x(ap)
\]

auf die fünfte Potenz, so verschwinden in der binomischen Entwicklung
rechts alle Glieder, in denen die Klammerfactoren $(ap)^2(aq)$ auftreten.
Demnach reducirt sich dieselbe auf

\weq{3}{
(pq)^5\cdot f = (aq)^5p_x^5 - 5(ap)(aq)^4q_xp_x^4 - (ap)^5q_x^5,
}

oder also auf

\weq{4}{
(pq)^5\cdot f = c_1p_x^5 + 5c_2p_x^4q_x + c_3q_x^5,
}

wo $c_1\,c_2\,c_3$ gewisse Constante sind, die wir nun berechnen wollen.

235. \textit{Berechnung der Coefficienten der Gleichung (4).} Zu dem
Zwecke gehen wir wiederum von der typischen Darstellung der Co\-
variante $j$ aus. [Vergl. Nr. 232, (I).] Da aber $j$ nach Voraussetzung
einen linearen Doppelfactor besitzt, und demnach ihre Discriminante
$\mathbf{C} = (\tau,\tau)^2 = 0$ ist, so reducirt sich dieselbe auf

\[
\mathbf{R}^2\cdot j = \delta^2\left(\delta + \frac32\,\mathbf{B}\alpha\right).
\]

Setzen wir also

\weq{5}{
\varrho = \delta + \frac32\,\mathbf{B}\alpha,
}

so hat man:

\weq{5a}{
\mathbf{R}^2\cdot j = \delta^2\cdot\varrho,
}

und

\weq{5b}{
(\delta\varrho) = \frac32\,\mathbf{B}(\delta\alpha) = -\,\frac32\,\mathbf{B}\cdot\mathbf{R}.
}

Die Gleichung (4) geht also über in

\[
(\delta\varrho)^5\cdot f = (a\varrho)^5\delta_x^5 - 5(a\delta)(a\varrho)^4\delta_x^4\varrho_x - (a\delta)^5\varrho_x^5,
\]

oder

\weq{6}{
(\delta\varrho)^5f = c_1\delta^5 + 5c_2\delta^4\varrho + c_3\varrho^5.
}


% ===== pdf idx 273, printed 259.0 =====
% !PRINTED 259
\weberpage{259}
Die Berechnung der Coefficienten $c_i$ bewerkstelligt man nun genau
in derselben Weise wie im allgemeineren Falle Nr. 233. Wir bilden
die fünften Ueberschiebungen dieser Gleichung (6) über $\alpha^5$, $\alpha^4\delta$, $\alpha^3\delta^2$,
und erhalten dadurch drei Gleichungen zur Bestimmung der Grössen $c_i$.

Es ergiebt sich nämlich dadurch, wenn wir einen Augenblick den
Factor $(\delta\varrho)^5$ mit $m\cdot\mathbf{R}^5$ bezeichnen [siehe Gleichung (5b)]:

\weq{6a}{
\left.
\begin{aligned}
\mathbf{R}^5\cdot m\cdot(a\alpha)^5 &= c_1(\delta\alpha)^5 + 5c_2(\delta\alpha)^4(\varrho\alpha) + c_3(\varrho\alpha)^5\\
\mathbf{R}^5\cdot m\cdot(a\alpha)^4(a\delta) &= 5c_2(\delta\alpha)^4(\varrho\delta) + c_3(\varrho\alpha)^4(\varrho\delta)\\
\mathbf{R}^5\cdot m\cdot(a\alpha)^3(a\delta)^2 &= c_3(\varrho\alpha)^3(\varrho\delta)^2
\end{aligned}
\right\}
}

Nun ist gemäss den Relationen (IV) und (V) in Nr. 230

\weq{7}{
(a\alpha)^5 = -\,\lambda\mathbf{R},
}

\[
(a\alpha)^4(a\delta) = \ \mu\mathbf{R}.
\]

Um aber $(f,\,\alpha^3\delta^2)^5$ zu berechnen, benutzen wir am einfachsten
die Relation $\delta^2 = \mathbf{P}\tau - \dfrac{\mathbf{C}}{2}\beta^2$ (vergl. Nr. 222), oder weil $\mathbf{C} = 0$, die
Relation:

\[
\delta^2 = \mathbf{P}\tau.
\]

Dann wird:

\[
\begin{aligned}
(a\alpha)^3(a\delta)^2 &= (f,\,\alpha^3\mathbf{P}\tau)^5 = \mathbf{P}((f,\tau)^2,\,\alpha^3)^3,\\
&= \mathbf{P}\left(\left(-i\alpha - 2\,\frac{\mathbf{A}}{3}\,j\right),\ \alpha^3\right)^3 \qquad \text{[vergl. Nr. 225, (1)]}\\
&= \frac23\,\mathbf{P}\mathbf{A}(-j,\,\alpha^3)^3 = \frac23\,\mathbf{P}\mathbf{A}(-(j,\alpha),\,\alpha^2)^2 \qquad \text{[Nr. 224, (2)]}\\
&= \frac23\,\mathbf{A}\mathbf{P}(\vartheta,\,\alpha)^2 = \frac23\,\mathbf{A}\mathbf{P}\mathbf{R},
\end{aligned}
\]

oder, weil für $\mathbf{C} = 0$ die Invariante $\mathbf{P} = \dfrac94\,\mathbf{A}\mathbf{B}^2$ wird [vgl. Nr. 226, (6)]:

\weq{9}{
(a\alpha)^3(a\delta)^2 = \frac32\,\mathbf{A}^2\mathbf{B}^2\mathbf{R}.
}

Anderntheils ist:

\[
(\delta\alpha) = -\,\mathbf{R}, \quad (\varrho\alpha) = -\,\mathbf{R}, \quad (\varrho\delta) = +\,\frac32\,\mathbf{B}\cdot\mathbf{R}.
\]

Demnach ergeben sich durch Eintragen dieser Werthe aus den
Gleichungen (6a) die folgenden Gleichungen:

\weq{10}{
\left.
\begin{aligned}
m\cdot\lambda\cdot\mathbf{R} &= c_1 + 5c_2 + c_3\\
m\cdot\mu\cdot\mathbf{R} &= {}+\,\frac{15}{2}\,\mathbf{B}\cdot c_2 + \frac32\,\mathbf{B}\cdot c_3\\
-\,\frac32\,m\cdot\mathbf{A}^2\cdot\mathbf{B}^2\cdot\mathbf{R} &= {}+\,\frac94\,\mathbf{B}^2\cdot c_3
\end{aligned}
\right\}
}

Hieraus findet sich zunächst

\[
c_3 = -\,\frac23\,\mathbf{A}^2\cdot m\cdot\mathbf{R}.
\]


% ===== pdf idx 274, printed 260.0 =====
% !PRINTED 260
\weberpage{260}
Indem man alsdann $\mu$ und $\lambda$ durch ihre Werthe (Nr. 230) in $\mathbf{A}$
und $\mathbf{B}$ ersetzt, findet man ferner

\[
c_2 = c_1 = \frac{\mathbf{B}}{6}\cdot m\cdot\mathbf{R}.
\]

Die typische Darstellung der Form $f$ wird demnach für $\mathbf{C} = 0$

\weq{11}{
6\mathbf{R}^4\cdot f = \mathbf{B}\delta^5 + 5\mathbf{B}\delta^4\varrho - 4\mathbf{A}^2\varrho^5.
}

Setzen wir $f = 0$, so haben wir in dieser Darstellung die
\glqq Bring'sche\grqq{} Normalform der Gleichung fünften Grades. Durch die
Substitutionen

\[
\frac{5\mathbf{B}}{-4\mathbf{A}^2} = \mathbf{Z}, \quad \frac{\varrho}{\delta} = t\sqrt[4]{\mathbf{Z}}
\]

geht sie in die von Hermite benutzte Gestalt über

\[
t^5 - t - \frac{1}{5\sqrt[4]{\mathbf{Z}}} = 0.
\]

Sobald also $\mathbf{C} = 0$, lässt sich immer eine lineare Transformation
angeben, wodurch die allgemeine Gleichung fünften Grades in diese
Normalform mit dem einzigen Parameter $X = \dfrac{1}{5\sqrt[4]{\mathbf{Z}}}$ übergeht. Da\-
gegen kann die allgemeine Gleichung fünften Grades nur durch eine
Transformation vierten Grades in diese Normalform gebracht werden.

236. \textit{Bedeutung der Relation $\mathbf{C} = 0$.} Es ist interessant, die Be\-
deutung der Relation $\mathbf{C} = 0$ für das Formensystem noch weiter zu
untersuchen. Man findet alsdann insbesondere zweierlei: 1) Die Co\-
varianten $j$ und $i$ besitzen, sobald $\mathbf{C} = 0$, einen gemeinsamen linearen
Factor, oder mit andern Worten: die Invariante $\mathbf{C}$ ist bis auf einen
Zahlenfactor gleich der Resultante von $j$ und $i$; 2) dieser gemeinsame
Factor von $i$ und $j$ ist nichts anderes als die oben eingeführte lineare
Covariante $\varrho$.

Indem wir die diesbezüglichen Untersuchungen hier anfügen, ver\-
folgen wir einen doppelten Zweck: einmal an einem Beispiel zu zeigen,
wie durch symbolische Rechnung die Resultante zweier Covarianten
abgeleitet wird, dann aber auch, um durch die sich ergebenden Re\-
sultate ein neues Licht auf das Wesen der Bring'schen Normalform
zu werfen.

Wir beginnen damit, die Resultante von $j$ und $i$ zu berechnen.

Angenommen, die beiden linearen Factoren der quadratischen
Form $i$ seien

\[
i = r_x\,s_x;
\]

sie besitze also die Wurzeln $\xi_1 = \dfrac{-r_2}{r_1} = r$, $\xi_2 = \dfrac{-s_2}{s_1} = s$; dann ist,
%CONT>>


% ===== pdf idx 275, printed 261.0 =====
% !PRINTED 261
\weberpage{261}
%<<CONT
von einem Zahlenfactor abgesehen, die Resultante von $j(x)$ und $i(x)$
bekanntlich dargestellt durch (vergl. auch Bd.\ I Nr.\ 168)
\[
\begin{aligned}
R_{j,\,i} &= j(r)\cdot j(s) = (jr)^3\cdot(j's)^3 \\
&= (jr)^2(j's)^2\{(jj')(rs)+(js)(j'r)\} \\
&= (jr)^2(j's)^2(jj')(rs)+(jr)^2(j'r)(j's)^2(js) \\
&= (jr)(j's)(jj')(rs)\{(jj')(rs)+(js)(j'r)\}+(jr)^2(j'r)(j's)^2(js) \\
&= (jj')^2(jr)(j's)(rs)^2+(jr)(js)(j'r)(j's)(jj')(rs) \\
&\qquad {}+(jr)^2(js)(j's)^2(j'r).
\end{aligned}
\]

Der erste Term rechts zerf\"allt in die Factoren
\[
(rs)^2 = -\,2(i,\,i)^2 = -\,2A
\]
und
\[
(jj')^2(jr)(j's) = (\tau,\,rs)^2 = (\tau,\,i)^2 = B.
\]
Der zweite Term verschwindet identisch, da er durch Vertauschung
von $j$ mit $j'$ nur sein Zeichen \"andert; der dritte Term endlich zer-
f\"allt wiederum in die beiden Factoren
\[
\begin{aligned}
(jr)^2(js) &= (j,\,rs\cdot r)^3 = (j,\,i\cdot r)^3 = ((j,\,i),\,r) \\
&= -(\alpha,\,r) = -(\alpha r),
\end{aligned}
\]
und
\[
(j's)^2(jr) = -(\alpha s).
\]
Durch Substitution dieser Werthe der einzelnen Glieder erhalten wir
\[
\begin{aligned}
R_{j,\,i} &= -\,2AB+(\alpha r)(\alpha s) \\
&= -\,2AB+(\alpha^2,\,i)^2 \\
&= -\,2AB+M,
\end{aligned}
\]
oder endlich, weil $M=2AB-3C$:
\weq{1}{R_{j,\,i} = -\,3C.}

Die Resultante von $j$ und $i$ stimmt also in der That mit der
Invariante $C$ bis auf einen numerischen Factor überein.

237. \textit{Berechnung des gemeinsamen Factors.} Nehmen wir nun an,
die Invariante $C$ verschwinde, und der gemeinsame lineare Factor von
$j$ und $i$ sei $r_x$, dann k\"onnen wir uns die Aufgabe stellen, diesen gemeinsamen
 Factor, der sich bekanntlich auf rationalem Wege berechnen
 l\"asst, durch die linearen Covarianten des Systemes auszudr\"ucken.

Um diese Aufgabe zu l\"osen, gehen wir von der Form aus
\[
F = (jr)^3 s_x^3 - (js)^3 r_x^3,
\]
welche sich f\"ur den Fall, dass $r$ oder $s$ der gemeinschaftliche Factor
ist, wegen $(jr)^3=0$, resp.\ $(js)^3=0$ auf den Cubus dieses gemeinschaftliches
 Factors reducirt. Nun ist aber
%CONT>>


% ===== pdf idx 276, printed 262.0 =====
% !PRINTED 262
\weberpage{262}
%<<CONT
\[
\begin{aligned}
F &= \{(jr)s_x-(js)r_x\}\{(jr)^2s_x^2+(jr)(js)r_xs_x+(js)^2r_x^2\} \\
&= (rs)j_x\{[(jr)s_x-(js)r_x]^2+3(jr)(js)r_xs_x\} \\
&= (rs)j_x\{j_x^2(rs)^2+3(jr)(js)r_xs_x\} \\
&= \{j\cdot(rs)^2+3(jr)(js)j_xr_xs_x\}(rs).
\end{aligned}
\]
Der Factor $(rs)$ ist eine von null verschiedene Constante; der
erste Term in der Klammer hat den Werth $=-\,2A\cdot j$, der zweite
Term den Werth $3(j,\,i)^2\cdot i=-\,3\alpha\cdot i$. Also wird
\[
F = \text{Const.}\{-\,2Aj-3\alpha i\},
\]
oder, wenn wir den Factor 3 herausziehen und mit der Constanten
vereinigen:
\[
F = \text{Const.}\left\{-\frac{2}{3}\,Aj-\alpha i\right\},
\]
d.\,i. wegen Nr.\ 225, (1)
\[
= \text{Const.}\,(f,\,\tau)^2.
\]

Wir haben also den Satz:

\glqq Besitzen $j$ und $i$ einen gemeinsamen Factor, ist also $C=0$,
so ist die zweite Ueberschiebung von $f$ und $\tau$ proportional
dem Cubus dieses gemeinsamen Factors.\grqq{}

Schiebt man diesen Cubus $(f,\,\tau)^2$ zweimal \"uber irgend eine beliebige
 Form zweiten Grades, so erh\"alt man den gemeinschaftlichen
Factor linear. Wir w\"ahlen hiezu die Form $\tau$ und erhalten
\[
((f,\,\tau)^2,\,\tau)^2 = \text{Const.}\left(-\frac{2}{3}\,Aj-\alpha i,\,\tau\right)^2;
\]
nach der Theorie der cubischen Formen aber ist $(j,\,\tau)^2=0$; demnach
wird, unter Vernachl\"assigung der Gr\"osse: $-\dfrac{3}{2}\cdot\text{Const.}$,
\[
\begin{aligned}
((f,\,\tau)^2,\,\tau)^2 &= +3(\alpha i,\,\tau)^2 = ((i\tau)^2\alpha+2(i\tau)(\alpha\tau)i_x) \\
&= B\cdot\alpha+2(i\tau)\{(i\tau)\alpha_x-(i\alpha)\tau_x\} \\
&= B\cdot\alpha+2(i\tau)^2\alpha-2(i\tau)(i\alpha)\tau_x \\
&= 3B\cdot\alpha+2(\tau,\,(i\alpha)i_x) \\
&= 3B\cdot\alpha+2\delta = +3B\cdot\alpha+2\delta,
\end{aligned}
\]
also wegen Nr.\ 235, (5)
\[
((f,\,\tau)^2,\,\tau)^2 = \frac{1}{2}\,\varrho\,.
\]

Es ergiebt sich daher:

\glqq Ist $C=0$, so ist die lineare Covariante $\varrho$ der gemeinschaftliche
 Factor von $j$ und $i$, und die Gleichung f\"unften Grades
geht durch die Substitution $x_1=\varrho$, $x_2=\delta$ in die Bring'sche
Normalform \"uber.\grqq{}


% ===== pdf idx 277, printed 263.0 =====
% !PRINTED 263
\weberpage{263}
238. \textit{Typische Darstellung der Form f\"unften Grades f\"ur $B=0$.}
Unter den verschiedenen Normalformen der Gleichung f\"unften Grades
beansprucht neben der Bring'schen Normalform auch noch eine
von Brioschi herr\"uhrende ein hervorragendes Interesse. Auf diese
Brioschi'sche Normalform wird man direct gef\"uhrt, sobald man die
typische Darstellung einer solchen Form $f$ unternimmt, deren Invariante
$B$ identisch verschwindet. Man hat nur statt der fr\"uher benutzten
linearen Covarianten $\alpha$ und $\beta$, resp.\ $\delta$ und $\varrho$ in diesem Falle die
Formen $\alpha$ und $\gamma$ f\"ur die betreffende typische Darstellung zu Grunde
zu legen. Wir wollen im Folgenden diese Normalform von $f=a_x^5$
f\"ur den Fall $B=0$ entwickeln, m\"ussen aber zu dem Zwecke zun\"achst
die Ueberschiebung $(f,\,\gamma)$ durch andere Formen des Systemes darstellen
, um alsdann in einfacher Weise die typischen Coefficienten
\[
(f,\,\alpha^5)^5,\quad (f,\,\alpha^4\gamma)^5,\quad (f,\,\alpha^3\gamma^2)^5,\quad (f,\,\alpha^2\gamma^3)^5,\quad (f,\,\alpha\gamma^4)^5,\quad (f,\,\gamma^5)^5
\]
berechnen zu k\"onnen.

239. \textit{Hilfssatz f\"ur die typische Darstellung von $f$ f\"ur $B=0$.} Wir
behaupten: Die lineare Ueberschiebung $(f,\,\gamma)$ l\"asst sich in der Form
darstellen
\weq{1}{(f,\,\gamma) = \frac{3}{2}\,\tau^2 + B\cdot(f,\,j)^2.}

Man erkennt schon aus dieser Relation deren Tragweite. Denn
f\"ur $B=0$ reducirt sich die rechte Seite auf das Quadrat der Form
$\tau$, und alle typischen Coefficienten --- von dem bereits bekannten
$(f,\,\alpha^5)^5$ abgesehen --- ergeben sich dann durch Ueberschiebung dieser
Form $\tau$ \"uber Potenzen von linearen Formen*).

Um diese Relation (1) zu beweisen, schreiben wir sie zun\"achst
in der Form
\[
(f,\,\gamma) - \frac{\tau^2}{2} = \tau^2 + B\cdot(f,\,j)^2
\]
und formen nun jede Seite dieser Gleichung durch symbolische Rechnung
geeignet um.

Wir erhalten dann f\"ur die linke Seite, da nach der Theorie der
cubischen Formen $(j,\,(j,\,\tau)) = -\dfrac{\tau^2}{2} = (j'\tau)(j\tau)j_x^2j_x'^2$:
\[
(f,\,\gamma) - \frac{\tau^2}{2} = (a\gamma)a_x^4 + (j'\tau)(j\tau)j_x^2j_x'^2
\]

\rule{4cm}{0.4pt}

\noindent *) Man kann $(f,\,\gamma)$ als Polare $f_y$ f\"ur $y=\gamma$ betrachten. Im Falle $B=0$
besitzt also die betreffende Form $f$ eine erste Polare, die ein volles Quadrat ist.
Umgekehrt: Weiss man von einer Form f\"unften Grades, dass sie eine erste Polare
besitzt, die ein volles Quadrat ist, so ist f\"ur diese Form $B=0$, und diese Polare
ist mit $(f,\,\gamma)$ identisch.


% ===== pdf idx 278, printed 264.0 =====
% !PRINTED 264
\weberpage{264}
%<<CONT
und, wenn wir im ersten Gliede $\gamma$ durch $(\tau\alpha)\tau_x$ oder wegen $\alpha_x=-(ji)^2j_x$
durch $-(ji)^2(\tau j)\tau_x$ ersetzen und im zweiten Gliede an Stelle von $j_x'^3$
das gleichwerthige Symbol $-(ai)^2a_x^3$ einf\"uhren:
\[
(f,\,\gamma) - \frac{\tau^2}{2} = -\,a_x^4(a\tau)(\tau j)(ji)^2 - a_x^2(ai)^2(a\tau)(j\tau)j_x^2,
\]
oder
\weq{2}{(f,\,\gamma) - \frac{\tau^2}{2} = a_x^2(a\tau)(j\tau)\left|\begin{array}{cc} a_x^2, & (ai)^2 \\ j_x^2, & (ji)^2 \end{array}\right|.}
F\"ur die rechte Seite dagegen erhalten wir successive
\[
\begin{aligned}
B\cdot(f,\,j)^2+\tau^2 &= (i\tau)^2\cdot(aj)^2a_x^3j_x+\tau_x^2\cdot(jj')^2j_xj_x' \\
&= a_x^3j_x\{(a\tau)(ij)+(j\tau)(ai)\}^2-\tau_x^2\cdot(ai)^2(aj)^2j_xa_x.
\end{aligned}
\]
Ersetzen wir im zweiten Term rechts $(ai)j_x$ durch $(aj)i_x+(ji)a_x$,
dann reducirt sich derselbe wegen $(aj)^3a_x^2=0$ auf $\tau_x^2\cdot(ai)(aj)^3(ji)a_x^2$.
Wenn wir also darin $(aj)\tau_x$ durch $(\tau j)a_x-(\tau a)j_x$ ersetzen, so geht
die letzte Gleichung \"uber in:
\[
B\cdot(f,\,j)^2+\tau^2 = a_x^3j_x\cdot\{(a\tau)(ij)+(j\tau)(ai)\}^2 - (ai)(ji)a_x^2\{(\tau j)a_x-(\tau a)j_x\}^2.
\]
Nun ist $(j,\,\tau)^2=0$, also verschwindet im ersten Binom rechts
das dritte, im zweiten Binom das erste Glied, w\"ahrend sich die beiden
mittleren Glieder gegenseitig aufheben. Es bleibt demnach:
\[
B\cdot(f,\,j)^2+\tau^2 = a_x^3j_x\cdot(a\tau)^2(ij)^2 - (ai)(ji)a_x^2\cdot(\tau a)^2j_x^2,
\]
oder
\weq{3}{B\cdot(f,\,j)^2+\tau^2 = a_x^2(a\tau)^2\left|\begin{array}{cc} a_xj_x, & (ai)(ji) \\ j_x^2, & (ji)^2 \end{array}\right|.}
Subtrahiren wir nun Gleichung (2) von Gleichung (3), so erhalten
 wir:
\[
B\cdot(f,\,j)^2 + \frac{3}{2}\tau^2 - (f,\,\gamma) = a_x^2(a\tau)^2\left|\begin{array}{cc} a_xj_x, & (ai)(ji) \\ j_x^2, & (ji)^2 \end{array}\right|
\]
\[
{}- a_x^2(a\tau)(j\tau)\left|\begin{array}{cc} a_x^2, & (ai)^2 \\ j_x^2, & (ji)^2 \end{array}\right|
\]
\[
= a_x^2\cdot\left|\begin{array}{ccc} (a\tau)^2, & a_x^2, & (ai)^2 \\ (a\tau)(j\tau), & a_xj_x, & (ai)(ji) \\ 0, & j_x^2, & (ji)^2 \end{array}\right|,
\]
oder wegen $0=(j,\,\tau)^2=(j\tau)^2j_x$
\weq{4}{B\cdot(f,\,j)^2 + \frac{3}{2}\tau^2 - (f,\,\gamma) = a_x^2\left|\begin{array}{ccc} (a\tau)^2, & a_x^2, & (ai)^2 \\ (a\tau)(j\tau), & a_xj_x, & (ai)(ji) \\ (j\tau)^2, & j_x^2, & (ji)^2 \end{array}\right|.}


% ===== pdf idx 279, printed 265.0 =====
% !PRINTED 265
\weberpage{265}
Die Determinante rechts verschwindet aber identisch; denn sie
ist durch Multiplication der beiden Determinanten

\[
\begin{vmatrix} \tau_1^{2}, & -\tau_1\tau_2, & \tau_2^{2} \\ x_1^{2}, & x_1x_2, & x_2^{2} \\ i_1^{2}, & -i_1i_2, & i_2^{2} \end{vmatrix}
\cdot
\begin{vmatrix} a_1^{2}, & 2a_1a_2, & a_2^{2} \\ a_1j_1, & a_1j_2+j_1a_2, & a_2j_2 \\ j_1^{2}, & 2j_1j_2, & j_2^{2} \end{vmatrix}
\]

entstanden. Die erste dieser Determinanten besitzt den Werth (vergl.
Bd. I Nr. 40, Anmerkung)

\[ \tau_x\cdot(\tau i)\cdot i_x, \]

die zweite besitzt den Werth:

\[ (aj)^3. \]

Die rechte Seite der Gleichung (4) geht also über in

\[ (aj)^3a_x^2\cdot(\tau i)\tau_x i_x, \]

und diese Grösse ist null, da der eine Factor $(aj)^3a_x^2=(f,j)^3$ identisch
verschwindet. Also hat man:

\[ B\cdot(f,j)^2+\frac{3}{2}\tau^2=(f,\gamma) \]

wie oben behauptet war.

240. \textit{Berechnung der typischen Coefficienten, wenn $B=0$.} Die
typische Darstellung der Form $f=a_x^5$ für $B=0$ ergiebt sich nun,
indem wir die Identität:

\[ (\alpha\gamma)a_x=(\alpha\gamma)\alpha_x-(a\alpha)\gamma_x \]

auf die fünfte Potenz erheben, und die typischen Coefficienten alsdann
in Function der Fundamentalinvarianten darstellen. Zu letzterem
Zwecke berücksichtigen wir, dass wegen $B=0$ folgende Beziehungen
bestehen:

\[
\left.\begin{aligned}
(f,\gamma) &= \frac{3}{2}\tau^2, \quad (\tau,\alpha)^2 = -(\alpha\gamma) = N = \frac{AC}{2} \\
(i,\alpha)^2 &= M = -3C
\end{aligned}\right\}
\begin{array}{l} \text{vgl.} \quad \text{Nr. 221} \quad (1) \\ \phantom{\text{vgl.}} \quad \text{Nr. 226} \quad (2) \end{array}
\]

\weq{3}{(\tau,\gamma)=(\tau',(\tau,\alpha))=-\frac{C}{2}\cdot\alpha \quad \text{(Identit\"atssatz)}}

\weq{4}{(\tau,\gamma^2)^2=((\tau,\gamma),\gamma)=-\frac{C}{2}(\alpha\gamma)=\frac{C}{2}\cdot N=\frac{AC^2}{4}\quad[\text{nach (3)}]}

\weq{5}{(\tau\gamma)(\tau\alpha)=-\frac{C}{2}(\alpha_1\alpha)=0\quad[\text{nach (3)}]}

\weq{6}{\lambda=-\frac{2}{3}A^2. \quad (\text{nach Nr. 230})}

Dann erhalten wir der Reihe nach folgende Gleichungen:

\[ (a\alpha)^5=\lambda\cdot R=-\frac{2}{3}A^2\cdot R \quad \text{(vergl. Nr. 230)} \]

\[ (a\alpha)^4(a\gamma) = ((f,\gamma),\alpha^4)^4=\frac{3}{2}(\tau^2,\alpha^4)^4=\frac{3}{2}(\tau\alpha)^2\cdot(\tau\alpha)^2=\frac{3A^2\cdot C^2}{8} \]


% ===== pdf idx 280, printed 266.0 =====
% !PRINTED 266
\weberpage{266}
\[ (a\alpha)^3(a\gamma)^2=((f,\gamma),\alpha^3\gamma)^4=\frac{3}{2}(\tau^2,\alpha^3\gamma)=\frac{3}{2}(\tau\alpha)^2\cdot(\tau\gamma)(\tau\alpha)=0 \]

\begin{align*}
(a\alpha)^2(a\gamma)^3 &= ((f,\gamma),\alpha^2\gamma^2)^4=\frac{3}{2}(\tau^2,\alpha^2\gamma^2)^4= \begin{pmatrix}\text{nach Ueberschiebungs-}\\ \text{gesetzen}\end{pmatrix}\\
&=\frac{3}{2}\cdot\frac{1}{6}\left\{(\tau\alpha)^2\cdot(\tau'\gamma)^2+4(\tau\gamma)(\tau\alpha)\cdot(\tau'\gamma)(\tau'\alpha)+(\tau'\alpha)^2(\tau'\gamma)^2\right\}\\
&=\frac{1}{4}\left\{[(\tau\alpha)(\tau'\gamma)-(\tau\gamma)(\tau'\alpha)]^2+6(\tau\gamma)(\tau\alpha)\cdot(\tau'\gamma)(\tau'\alpha)\right\}\\
&=\frac{1}{4}(\tau\tau')^2\cdot(\alpha\gamma)^2=\frac{1}{4}C\cdot N^2=\frac{A^2C^3}{16}
\end{align*}

\[ (a\alpha)(a\gamma)^4=((f,\gamma),\alpha\gamma^3)^4=\frac{3}{2}(\tau^2,\alpha\gamma^3)^4=\frac{3}{2}(\tau\gamma)^2\cdot(\tau\gamma)(\tau\alpha)=0 \]

\[ (a\gamma)^5=((f,\gamma),\gamma^4)^4=\frac{3}{2}(\tau^2,\gamma^4)=\frac{3}{2}(\tau\gamma)^2\cdot(\tau\gamma)^2=\frac{3A^2C^4}{32}. \]

Tragen wir die soeben berechneten Werthe der typischen Coefficienten
 in die Gleichung

\[ (\alpha\gamma)^5\cdot a_x^5=\{\alpha_x(\alpha\gamma)-\gamma_x(a\alpha)\}^5 \]

ein, so erhalten wir:

\[ -\frac{A^5\cdot C^5}{32}\cdot f = \frac{3}{32}A^2C^4\alpha^5+\frac{10}{16}A^2C^3\alpha^3\gamma^2+\frac{15}{8}A^2C^2\alpha\gamma^4-\frac{2}{3}A^2R\gamma^5. \]

Dividiren wir diese Gleichung mit $A^2C^2$ und setzen

\[ \frac{\alpha}{\gamma}\sqrt{-\frac{3}{2}C}=\xi, \]

so kommt endlich nach leichter Reduction:

\[ -\frac{1}{4}\cdot\frac{A^3C^4}{\gamma^5}\cdot\sqrt{-\frac{3C}{2}}\cdot f=\xi^5-10\xi^3+45\xi-\frac{16\sqrt{-\frac{3}{2}}}{\sqrt{C^3}}\cdot R. \]

Die Form fünften Grades erscheint hiemit abermals in einer
Normalform, die nur einen Parameter besitzt, und diese Normalform
wird die Brioschi'sche genannt. Eine ganz elementare Methode,
die allgemeine Gleichung fünften Grades durch eine Tschirnhaustransformation
 in die Brioschi'sche Form zu bringen, hat Gordan
neuerdings in den Math. Ann. Bd. XXVIII publicirt.

\begin{center}
\textbf{\S\ 25. Die Discriminante und die Auflösung der Gleichung}\\
\textbf{fünften Grades.}
\end{center}

241. \textit{Erste Darstellung der Discriminante durch die Fundamentalinvarianten
.} Aus der in Nr. 235, (11) gegebenen Darstellung der Form
$f$ lässt sich unmittelbar der Werth der Discriminante $\Delta$ von $f$ erschliessen
. Die Discriminante muss nämlich als Resultante von $\dfrac{\partial f}{\partial x_1}$
und $\dfrac{\partial f}{\partial x_2}$ vom Grade 8 in den Coefficienten sein; deshalb kann in der
%CONT>>


% ===== pdf idx 281, printed 267.0 =====
% !PRINTED 267
\weberpage{267}
%<<CONT
ganzen und rationalen Function von Invarianten, durch welche sie
sich ausdrücken lassen muss, die Invariante $C$ zwölften Grades in
den Coefficienten nicht auftreten. $\Delta$ ist sonach unabhängig von $C$
und ändert sich daher nicht, wenn $C=0$ gesetzt wird. Die Discriminante
der Form

\[ 6R^4f = B\delta^5+5B\delta^4\varrho-4A^2\varrho^5=0 \]

ist also auch die Discriminante der allgemeinen Form. Es ist aber

\[ \frac{\partial f}{\partial\delta} = 5\{B\delta^4+4B\delta^3\varrho\}=0 \]

\[ \frac{\partial f}{\partial\varrho} = 5\{B\delta^4-4A^2\varrho^4\}=0. \]

Die erste der beiden Gleichungen liefert

\[ \frac{\delta}{\varrho} = -4. \]

Substituirt man diesen Werth in die zweite, so ergiebt sich als
Resultat der Elimination von $\dfrac{\delta}{\varrho}$:

\[ \Delta = 64B-A^2. \]

242. \textit{Zweite Darstellung der Discriminante.} Wir wollen im Folgenden
 noch eine zweite Methode zur Berechnung der Discriminante geben,
die an die allgemeine Form $f=a_x^5$ anknüpft und auf denselben Grundsätzen
 beruht, wie früher die Berechnung der Discriminante einer
Form vierten Grades.

Wie damals nehmen wir wieder an, die Form $f$ besitze einen
linearen Doppelfactor $p_x$, sei also dargestellt durch

\[ f = p^2\cdot l = a_x^5, \]

wo $l$ irgend eine cubische Form bedeute. Indem wir uns alsdann die
Frage stellen, auf welche Invariantenrelation $F(A,B)=0$ führt diese
Annahme, werden wir naturgemäss auf die eine geführt, deren linke
Seite wir als Discriminante bezeichnen müssen.

Die Beantwortung der Frage lässt sich mannigfach bewerkstelligen.
Es handelt sich in der Hauptsache darum, die Invarianten $A$ und $B$
der allgemeinen Form mit dem simultanen System unserer Form $f$
und $p$ in Beziehung zu setzen. Hiezu liefern uns die in Nr. 99 aus
der Bedingung $f=p^2\cdot l$ abgeleiteten Gleichungen die geeigneten
Mittel. Wir erhielten damals:

\weq{1}{(ap)^5=0, \quad (ap)^4a_x=0}

\weq{2}{(ap)^3a_x^2=\lambda\cdot p^2}

\weq{3}{(ab)^2(ap)^3b_x^3=\frac{1}{2}(i,p)\cdot p^2,}


% ===== pdf idx 282, printed 268.0 =====
% !PRINTED 268
\weberpage{268}
wobei Gleichung (3) aus $(ap)^4a_x=0$ durch die a. a. O. angegebene
Reihenentwicklung entsteht. Man könnte nun direct aus diesen
Gleichungen und der Voraussetzung $a_x^5=p^2l$ die Formen $A=(i,i)^2$,
$B=(\tau,i)^2$ zu berechnen suchen; dabei würde man finden, dass beide
Formen $A$ und $B$ sich als Functionen von $(ip)^2$ und $(\tau p)^2$ darstellen
lassen. Indem wir daher hier von vornherein nach der Berechnung
dieser beiden Grössen streben, gelangen wir auf kürzerem Wege
zum Ziele.

Um den Werth von $(ip)^2$ zu ermitteln, überschieben wir (2) zweimal
über $f$ und erhalten:

\[ (ab)^2(ap)^3b_x^3 = \lambda\cdot(f,p^2)^2, \]

oder durch Comparation mit (3)

\weq{4}{(i,p)\cdot p^2=2\lambda\cdot(f,p^2)^2.}

Indem wir aber diese Gleichung abermals mit $p$ überschieben,
kommt:

\[ (i,p)^2\cdot p^2 = 6\lambda\cdot(f,p^3)^3, \]

oder endlich wegen (2)

\weq{5}{6\lambda^2=(ip)^2.}

Zur Berechnung von $(\tau p)^2$ stellen wir zunächst eine Entwicklung für
das Product $i\cdot p^4$ her, indem wir die Identität $(ab)p_x=(ap)b_x-(bp)a_x$
auf die vierte Potenz erheben, und nachträglich mit $a_xb_x$ multipliciren.
Dann kommt nach einfacher Reduction, weil $(ap)^4a_x=0$:

\[ i\cdot p^4 = -8(ap)^3a_x^2\cdot(bp)b_x^4+6(ap)^2a_x^3\cdot(bp)^2b_x^3, \]

oder:

\[ i\cdot p^4 = -8\lambda p^2\cdot(f,p)+6(f,p^2)^2\cdot(f,p^2)^2 \]

und demnach wegen Gleichung (4)

\weq{6}{i\cdot p^4 = -8\lambda p^2\cdot(f,p)+\frac{6}{4\lambda^2}\cdot p^4\cdot(ip)(i_1p_1)i_xi_{1x}.}

Das symbolische Product $2(ip)(i_1p_1)i_xi_{1x}$ lässt sich mit Hilfe des
Productsatzes (vergl. Nr. 11) leicht in die Form $12i\lambda^2-Ap^2$ bringen,
so dass alsdann die Gleichung (6) die Form annimmt:

\weq{7}{8ip^2 = 8\lambda\cdot(f,p)+\frac{3}{4\lambda^2}Ap^4,}

wie man sieht, eine Gleichung für die einzige uns noch unbekannte
Ueberschiebung der Form $f$ über $p$. Aus ihr lässt sich zunächst
$(j,p)=(-(f,i)^2,p)$ berechnen, womit der erste Schritt zur Darstellung
 von $(\tau p)^2$ gemacht ist, da ja $\tau=(j,j)^2$. Wir überschieben
nämlich (7) zweimal über $i$ und erhalten:

\[ \frac{8}{6}\{Ap^2+4(ii_1)(pi_1)i_xp_x+i\cdot(pi_1)^2\} = -8\lambda\cdot(j,p)+\frac{3}{4\lambda^2}A\cdot p^2\cdot(pi)^2, \]
%CONT>>


% ===== pdf idx 283, printed 269.0 =====
% !PRINTED 269
\weberpage{269}
%<<CONT
oder, weil nach dem Identit\"atssatze das mittlere Glied links mit $2\,A\,p^2$ \"ubereinstimmt und $(i\,p)^2 = 6\,\lambda^2$, so kommt:
\weq{8}{8\,\lambda\cdot(j,\,p) = \frac{1}{2}\,A\,p^2 - 8\,\lambda^2\cdot i.}

Ueberschiebt man aber diese Gleichung zweimal \"uber sich selbst, so erhalten wir nach Division mit $16\,\lambda^2$:
\weq{9}{4\,(\tau p)^2 = \lambda^2\cdot A.}

Ueberschiebt man dagegen Gleichung (8) zweimal \"uber $\tau$, so kommt, weil nach der Theorie der cubischen Formen $(j,\,\tau)^2$ verschwindet:
\[
0 = \frac{1}{2}\,A\cdot(\tau p)^2 - 8\,\lambda^2\cdot B,
\]
oder wegen (9)
\weq{10}{0 = A^2 - 64\,B.}

Wenn also $f$ einen Doppelfactor besitzt, so besteht zwischen $A$ und $B$ die Relation (10), d.\,h. $A^2-64B$ ist die Discriminante von $f$.

243. \textit{Aufl\"osung der Gleichung f\"unften Grades.} Die Aufl\"osung der Gleichung f\"unften Grades ist bekanntlich durch algebraische Methoden allein nicht mehr m\"oglich. Die Grenze, bis zu welcher die Algebra vorgehen kann, haben wir bereits kennen gelernt: sie ist durch die Transformation der allgemeinen Gleichung in eine Form mit einem einzigen Parameter $Z$ gekennzeichnet. Dieselbe kann unter anderm stets nach Klein (vergl. \glqq Vorlesungen \"uber das Ikosaeder\grqq{} \S.\,10, Seite 252) in rationaler Weise auf die in Nr.\,192 erw\"ahnte Ikosaedergleichung $\dfrac{C_2\cdot f^3}{H^3} = \varrho$ zur\"uckgef\"uhrt werden, deren Wurzeln $\xi = \dfrac{x_1}{x_2} = \vartheta(\varrho)$ sich mit Hilfe elliptischer Modulfunctionen darstellen lassen. (Vergl. Klein: \glqq Vorlesungen \"uber das Ikosaeder\grqq{} \S.\,7, Seite 132.) Wenn wir uns also hier mit Aufl\"osung von Gleichungen f\"unften Grades besch\"aftigen, so m\"ussen wir uns auf specielle F\"alle beschr\"anken, bei welchen verm\"oge gewisser Invariantenrelationen die Aufl\"osung sich ohne transcendente Methoden bewerkstelligen l\"asst. Aus der grossen Mannigfaltigkeit der hier gegebenen M\"oglichkeiten besprechen wir hier nur einige wichtige F\"alle. Hieher geh\"oren in erster Linie die Formen $f$, f\"ur welche $R=0$, $M\gtrless 0$, $N\gtrless 0$; denn w\"ahrend f\"ur $B=0$ resp. $C=0$ nur Reductionen der Form $f$ auf gewisse, besonders interessante Normalformen eintreten und die Bedingung $A=0$ \"uberhaupt keine bemerkenswerthe Typik zur Folge hat, zerf\"allt f\"ur $R=0$ die Form $f$ rational in einen linearen und einen biquadratischen Factor. Hieher geh\"oren ferner die F\"alle:
%CONT>>


% ===== pdf idx 284, printed 270.0 =====
% !PRINTED 270
\weberpage{270}
%<<CONT
\[
\text{II)}\quad \alpha = 0, \text{ und eine der drei Formen } A,\ B,\ C\gtrless 0,
\]
\[
\text{III)}\quad \alpha = 0,\ B=0,\ C=0,\ A_1=0,\ j\gtrless 0,
\]
\[
\text{IV)}\quad \alpha = 0,\ A=B=C=0,\ j=0.
\]

244. \textit{Erster Fall.} Wir betrachten zun\"achst den Fall, in welchem die schiefe Invariante verschwindet, also den Fall
\[
R=0,\ M\gtrless 0,\ N\gtrless 0.
\]

Da alle schiefen Invarianten des Systemes $R$ zum Factor haben m\"ussen, so verschwinden dieselben durchwegs.

Solche schiefe Invarianten sind z.\,B.
\[
(f,\ \alpha)^5,\ \ (f,\ \alpha^3\beta^2)^5,\ \ (f,\ \alpha\beta^4)^5,
\]
ebenso
\[
(f,\ \gamma)^5,\ \ (f,\ \gamma^3\alpha^2)^5,\ \ (f,\ \gamma\alpha^4)^5.
\]
Wenn wir also entweder
\[
a_x(\beta\alpha) = \beta_x(a\alpha) - \alpha_x(a\beta),
\]
oder
\[
a_x(\gamma\alpha) = \gamma_x(a\alpha) - \alpha_x(a\gamma)
\]
auf die f\"unfte Potenz erheben, so verschwinden stets drei Glieder in der binomischen Entwicklung rechts, so dass wir erhalten:
\[
f\cdot(\beta\alpha)^5 = c_1\alpha^5 + c_2\alpha^3\beta^2 + c_3\alpha\beta^4,
\]
oder
\[
f\cdot(\gamma\alpha)^5 = \bar c_1\alpha^5 + \bar c_2\alpha^3\gamma^2 + \bar c_3\alpha\gamma^4,
\]
wo $c_1\,c_2\,c_3$, $\bar c_1\,\bar c_2\,\bar c_3$ wie in Nr.\,235 zu berechnen sind. $(\beta\alpha)$ ist dabei gleich $M$, $(\gamma\alpha) = N$; da diese Invarianten nicht verschwinden, so sind die rechten Seiten in der That Entwicklungen von $F$. Setzen wir sie gleich null, so erhalten wir in beiden F\"allen von vornherein eine Wurzel von $f$, n\"amlich $\alpha_x=0$. Der andere Factor ist dann eine biquadratische Gleichung, die sich leicht auf eine quadratische reduciren l\"asst.

245. \textit{Zweiter Fall.} Wir nehmen an, dass die beiden Invarianten $M$ und $N$ verschwinden, woraus von selbst hervorgeht, dass auch $R$ null wird [vergl. Nr.\,227, (4)], also dass
\[
M=0,\ N=0,\ R=0.
\]
Die Invariante $C$ sei von null verschieden.

Wenn aber $M=N=0$, dann ist, wie wir nun beweisen wollen, auch $\alpha=0$, und folglich existirt in diesem Falle \"uberhaupt keine lineare Covariante.

Nehmen wir zun\"achst an, $\alpha$ w\"are von null verschieden. Es ist nun:
\weq{11}{i_x^2\tau_y^2 - \tau_x^2 i_y^2 = (i_x\tau_y - \tau_x i_y)(i_x\tau_y + \tau_x i_y) = 2\,\vartheta_x\vartheta_y(xy).}


% ===== pdf idx 285, printed 271.0 =====
% !PRINTED 271
\weberpage{271}
Ersetzen wir $y$ durch $\alpha$, so kommt:
\[
i_x^2(\tau\alpha)^2 - \tau_x^2(i\alpha)^2 = 2(\vartheta\alpha)\vartheta_x\alpha_x.
\]
Weil aber
\[
(\tau\alpha)^2 = N = 0,
\]
\[
(i\alpha)^2 = M = 0,
\]
so folgt:
\[
(\vartheta,\ \alpha) = 0,
\]
d.\,h. $\vartheta$ hat den Factor $\alpha$. Da nun aber
\[
j_x^2(j\alpha) = -\vartheta, \quad \text{[vergl. Nr.\,224, (2)]}
\]
so ist auch
\[
(j\alpha)^2 j_x = ((j\alpha)j_x^2,\ \alpha) = -(\vartheta,\ \alpha) = 0.
\]
Daraus folgt, dass $j$ den linearen Factor $\alpha$ mindestens quadratisch besitzt, dass also, wenn $\varrho$ irgend ein linearer Factor
\[
j = \alpha^2\cdot\varrho,
\]
und folglich (vergl. Nr.\,154)
\[
\tau = \alpha^2\cdot c,
\]
wo $c$ eine Constante ist.

Erinnern wir uns nun, dass [vergl. Nr.\,230, (II) und 225, (1)]
\weq{12}{(f,\ \alpha^2)^2 = A\cdot(f,\ \tau)^2 + Bj + \frac{3}{2}\,\alpha\tau,}
folglich
\[
(f,\ \alpha^2)^2 = Ac\cdot(f,\ \alpha^2)^2 + B\alpha^2\varrho + \frac{3}{2}\,c\alpha^3;
\]
und bedenken wir ferner, dass
\[
0 = (f,\ j)^3 = (f,\ \varrho\alpha^2)^3 = ((f,\ \alpha^2),\ \varrho),
\]
so folgt, wenn wir (10) dreimal \"uber $\varrho$ schieben:
\[
((f,\ \alpha^2),\ \varrho^3)^3 = Ac\,((f,\ \alpha^2)^2,\ \varrho^3)^3 + B\,(\alpha^2\varrho,\ \varrho^3)^3 + \frac{3}{2}\,c\,(\alpha\varrho)^3.
\]
Nun ist aber $(\alpha^2\varrho,\ \varrho^3)^3 = 0$, und wegen
\[
((f,\ \alpha^2),\ \varrho) = 0
\]
sowohl das Glied links als auch das erste Glied rechts gleich null; es bleibt also noch
\[
(\alpha\varrho)^3 = 0,
\]
d.\,h.
\[
\varrho = c_1\alpha,
\]
folglich
\[
j = c_1\alpha^3.
\]
Da aber auch $i$ den Factor $\alpha$ besitzen muss wegen
\[
M = (i\alpha)^2 = 0,
\]
so ist auch
\[
(j,\ i)^2 = -\alpha = 0,
\]
was der Voraussetzung widerspricht.

Die Bedingung $M=0$, $N=0$ f\"uhrt also stets $\alpha=0$ mit sich, wie umgekehrt aus $\alpha=0$ auch sofort das Verschwinden von $M$ und $N$ sich ergiebt.


% ===== pdf idx 286, printed 272.0 =====
% !PRINTED 272
\weberpage{272}
246. Wir nehmen also an, die lineare Covariante $\alpha$ verschwinde, doch von den drei fundamentalen Invarianten $A$, $B$, $C$ sei wenigstens eine von null verschieden.

Wegen $\alpha=0$ ist zun\"achst
\[
\vartheta = -(j,\ \alpha) = 0,
\]
und daher [vergl. Nr.\,227, (1)]
\[
2(\vartheta,\ i) = A\tau - Bi = 0.
\]
Es ist sonach $\tau$ proportional der Covariante $i$. Und da nun
\[
(f,\ \alpha) = 2((f,\ j)^2,\ i) = 0, \quad \text{[vergl. Nr.\,230, (I)]}
\]
so muss, so lange $A=(i,\ i)^2\gtrless 0$, $(f,\ j)^2$ eine Potenz von $i$ sein, d.\,h. es muss die Relation bestehen:
\[
(aj)^2 a_x^3 j_x = c\cdot i^2,
\]
oder weil
\[
A\tau = Bi,
\]
auch
\[
(aj)^2 a_x^3 j_x = c_1\tau^2,
\]
wo $c_1$ eine Constante.

Schiebt man diese Relation einmal \"uber $j$, so erh\"alt man, da die beiden Glieder der ersten Polare von $(aj)^2 a_x^3 j_x$ einander gleich sein m\"ussen --- ihre Differenz hat den Factor $(aj)^3$ ---
\weq{13}{(aj)^2(jj_1)a_x^3 j_{1x}^2 = c_1(\tau^2,\ j) = c_1\tau\cdot(\tau,\ j).}
Vertauscht man links $j$ mit $j_1$ und nimmt die halbe Summe, so kommt:
\[
\frac{1}{2}\,a_x^3(jj_1)\{(aj)^2 j_{1x}^2 - (aj_1)^2 j_x^2\}
\]
\[
= \frac{1}{2}\,a_x^3(jj_1)\{(aj)j_{1x} - (aj_1)j_x\}\{(aj)j_{1x} + (aj_1)j_x\}
\]
\[
= -\frac{1}{2}\,a_x^4(jj_1)^2(aj)j_{1x} + \frac{1}{2}\,a_x^4(jj_1)^2(aj_1)j_x.
\]
Beide Glieder sind aber ein und dasselbe symbolische Product, welches nichts anderes darstellt, als die Functionaldeterminante von $f$ und $\tau=(jj_1)^2 j_x j_{1x}$. Wir haben also:
\[
(jj_1)^2(aj)a_x^4 j_{1x} = (f,\ \tau),
\]
und folglich l\"asst sich Gleichung (13) in der Form schreiben:
\[
(f,\ \tau) + c_1\tau(j,\ \tau) = 0,
\]
oder
\[
(f + c_1\tau j,\ \tau) = 0.
\]
Es m\"usste also, da $\tau = \dfrac{B}{A}\cdot i$ von null verschieden ist % ?? uncertain: possible diacritic/bar over A in this fraction, resolution insufficient to confirm
und keine zwei gleichen Factoren besitzt, $f+c_1\tau j$ eine Potenz von $\tau$ sein. Dies ist aber unm\"oglich, da $f$ und $j$ ungeraden Grades. Daher bleibt nur
\[
f + c_1\tau j = 0,
\]
%CONT>>


% ===== pdf idx 287, printed 273.0 =====
% !PRINTED 273
\weberpage{273}
d. h., wenn die lineare Covariante $\alpha$ verschwindet, aber nicht alle
fundamentalen Invarianten, so zerfällt die Form fünften Grades in ihre
Covariante $j$ und deren Hesse'sche $\tau$:
\[
f = c_2 \cdot j \cdot \tau .
\]

247. \textit{Dritter Fall.} Wir nehmen nun an, dass neben $\alpha = 0$ auch
die Invarianten $B$ und $C$ verschwinden mögen, während die Invariante $A$
noch willkürlich sei.

Aus der Relation [vgl. Nr. 230 (II)]
\[
(f,\alpha^2)^2 = \frac{3}{2}\,\alpha\tau - Ai\alpha + \left(B - \frac{2}{3}A^2\right) j
\]
folgt alsdann für $\alpha = 0$ die Beziehung:
\[
-\frac{2}{3}A^2 j = 0.
\]
Daraus ergiebt sich, dass entweder $A$ oder $j$ verschwinden muss. Nehmen wir an
\weq{1}{j \gtrless 0,\quad A = 0.}
Dann ist einmal $(f,\tau)^2 = 0$, weil ja nach Nr. 225:
\weq{1}{(f,\tau)^2 = -i\alpha - \frac{2}{3}Aj,}
und ebenso [vgl. Nr. 245 (11)]:
\weq{2}{i_x^2\tau_y^2 - i_y^2\tau_x^2 = 0,}
da $\vartheta = -(j,\alpha) = 0$. Ersetzen wir in (2) $y$ durch $a$, so kommt:
\[
i\cdot(f,\tau)^2 - \tau\cdot(f,i)^2 = 0,
\]
also wegen $(f,\tau)^2 = 0$
\[
\tau\cdot(f,i)^2 = 0 = -j\cdot\tau .
\]
Da $j$ von null verschieden sein soll, so folgt daraus
\[
\tau = 0,
\]
d. h. aber, die Covariante $j$ ist ein reiner Cubus:
\[
j = \varrho_x^3 .
\]
Und weil nun in jedem Falle:
\[
(f,j)^3 = 0,
\]
so ist nunmehr
\[
(a\varrho)^3 a_x^2 = 0,
\]
und daraus erkennt man, dass $f$ den Factor $\varrho$ dreimal enthält. Wenn
also: $\alpha = 0$, $B = 0$, $C = 0$, $A = 0$ und $j \gtrless 0$, dann ist
\[
f = \varrho_x^3 \cdot l,
\]
wo $l$ quadratische Form ist.

248. Setzen wir nun aber voraus, dass
\weq{2}{A \gtrless 0,\quad j = 0;}
%CONT>>


% ===== pdf idx 288, printed 274.0 =====
% !PRINTED 274
\weberpage{274}
%<<CONT
dann besitzt zunächst $i$ wegen $A = A_{ii} \gtrless 0$ zwei verschiedene
lineare Factoren
\[
i = \varrho \cdot \sigma .
\]
Aber aus:
\[
-j = (f,i)^2 = (a\varrho)(a\sigma)a_x^3 = 0
\]
folgt, dass jedes symbolische Product mit dem Reducenten
\[
(a\varrho)(a\sigma)
\]
verschwindet. Wenn man daher die Identität
\[
a_x(\varrho\sigma) = \varrho_x(a\sigma) - \sigma_x(a\varrho)
\]
auf die fünfte Potenz erhebt, so reducirt sich die binomische Entwicklung
rechts auf ihr erstes und letztes Glied, und wir erhalten
\[
f\cdot(\varrho\sigma)^5 = c_1\varrho^5 + c_2\sigma^5,
\]
d. h. $f$ ist in diesem Falle eine binomische Gleichung.

249. \textit{Vierter Fall.} Nehmen wir endlich an, dass neben $\alpha = 0$,
$B = 0$, $C = 0$ auch
\weq{3}{j = 0 \text{ und } A = 0}
sei, dann ist zunächst $i$ ein volles Quadrat, also
\[
i = \varrho_x^2,
\]
und weil
\[
-j = (f,i)^2 = (a\varrho)^2 a_x^3 = 0
\]
so muss $f$ den linearen Factor $\varrho$ viermal enthalten. Daraus geht
aber umgekehrt hervor, dass $i$ selbst verschwinden muss. Denn es ist
\weq{1}{\varrho_x^4(ab)^4 a_x b_x = \{b_x(a\varrho) - a_x(b\varrho)\}^4 a_x b_x .}
Der Ausdruck rechts verschwindet aber wegen der Reducenten
\[
(a\varrho)^2,\quad (b\varrho)^2,
\]
und folglich auch
%CONT>>


% ===== pdf idx 289, printed 275.0 =====
% !PRINTED 275
\weberpage{275}
%<<CONT
\[
\varrho_x^4(ab)^4 a_x b_x = i^3 = 0.
\]
Wenn also $\alpha = 0$, $B = C = A = 0$, und $j = 0$, dann verschwindet auch $i$.

Fragen wir uns nun aber umgekehrt, welche Eigenschaft besitzt $f$, wenn $i$
verschwindet, so geht aus dieser einen Bedingung zunächst das Verschwinden
aller Invarianten hervor, so dass $f$ jedenfalls mindestens einen
quadratischen Factor besitzt, den wir mit $\varrho$ bezeichnen wollen.
Bilden wir nun:
\weq{2}{\varrho_x^4(ab)^4(a\varrho)(b\varrho) = \{b_x(a\varrho) - a_x(b\varrho)\}^4(a\varrho)(b\varrho),}
so ist der Ausdruck links gleich dem Product $\varrho_x^4\cdot(i,\varrho^2)^2 = 0$.
Die rechte Seite reducirt sich aber wegen $(a\varrho)^4 a_x = (b\varrho)^4 b_x = 0$ auf
\weq{3}{(a\varrho)^3 a_x^2 \cdot (b\varrho)^3 b_x^2 = \{(a\varrho)^3 a_x^2\}^2 = 0,}
d. h. $f$ enthält $\varrho$ cubisch. Aus diesem Grunde erhalten wir aus der
Identität
\[
0 = \varrho_x^4\cdot(ab)^4(b\varrho)a_x = \{b_x(a\varrho) - a_x(b\varrho)\}^4(b\varrho)a_x
\]
die Beziehung:
\weq{4}{(a\varrho)^2 a_x^3 = 0,}
welche lehrt, dass $f$ den Factor $\varrho$ viermal besitzt. Verschwindet
also die Form $(f,f)^4$ identisch, so reducirt sich $f$ auf
$\varrho_x^4\cdot\sigma_x$, wo $\varrho$ und $\sigma$ lineare Formen sind,
ein Resultat, das wir schon in der allgemeinen Untersuchung § 19 erhalten
haben.

§ 26. \textbf{Das Formensystem der Form sechsten Grades $f = a_x^6$.}

250. \textit{Aufstellung der Systeme $(A^{(i)})$ und $(B^{(i)})$.} Wie bei
der Form fünften Grades, so ist auch hier das System $(A^{(0)})$ durch
$f = a_x^6$, das System $(B^{(0)})$ durch $(f,f)^2 = (ab)^2 a_x^4 b_x^4$
dargestellt, und man zeigt wie dort, dass das aus beiden durch
Ueberschiebung entstehende System $(A^{(1)})$ nur die drei Formen enthält:
\[
f,\quad (f,f)^2,\quad ((f,f)^2, f).
\]
Dieses System $(A^{(1)})$ ist relativ vollständig mod. $(ab)^4$. Wir
überschieben dasselbe mit dem Bildsystem resp. wirklichen System
$(B^{(1)})$ der Form:
\[
k = (f,f)^4 = (ab)^4 a_x^2 b_x^2 .
\]
Da diese Form $k$ vierten Grades ist, so ist ihr System bereits bekannt.
Es besteht aus den fünf Formen:
\[
k,\quad (k,k)^2,\quad ((k,k)^2, k),\quad i_k = B,\quad j_k = C,
\]
und ist relativ vollständig mod. $(ab)^6$. Die beiden Formen $i_k$ und
$j_k$ sind Invarianten und als solche an der Ueberschiebung nicht
betheiligt. Das vollständige Formensystem $(A^{(2)})$ von $f$ besteht
demnach aus den durch Ueberschiebung der Covarianten
\[
f,\quad (f,f)^2,\quad ((f,f)^2, f)
\]
über die Covarianten
\[
k,\quad (k,k)^2,\quad ((k,k)^2, k)
\]
sich ergebenden Formen, zu welchen ausserdem die drei Invarianten
\[
A = (f,f)^6,\quad B = (k,k)^4,\quad C = ((k,k)^2, k)^4
\]
hinzutreten.

251. \textit{Methode der Ueberschiebung beider Systeme.} Es ist nun
vortheilhaft, diesen Ueberschiebungsprocess nicht direct auf dem
allgemeinen Wege vorzunehmen, wie ihn die aus je drei Formen bestehenden
Systeme dictiren würden. Vielmehr bietet die folgende Methode wesentliche
Vereinfachungen. Denken wir uns das volle System von $f$ gebildet, so
können wir die symbolischen Producte, durch welche die Formen des
Systemes dargestellt sind, in zwei Klassen theilen.


% ===== pdf idx 290, printed 276.0 =====
% !PRINTED 276
\weberpage{276}
Bei den Faltungen, welche die Ueberschiebungen des Systemes $(A^{(1)})$
über das System $(B^{(1)})$ vorschreiben, entstehen nämlich nur
Klammerfactoren vom Typus $(ak)$. Alle sich so ergebenden Formen des
Systemes werden demnach mindestens die erste und höchstens die vierte
Potenz dieses Klammerfactors $(ak) = (bk_1) = (ck_2)$ etc. enthalten.
Fassen wir nun alle Formen zusammen, welche die höchste Potenz $(ak)^4$
enthalten, so kann man sich dieselben entstanden denken durch
Ueberschiebung mit der Form
\[
l = (ak)^4 a_x^2,
\]
während sämmtliche Formen, welche die drei ersten Potenzen dieses
Factors enthalten, für sich ein modulo $(ak)^4$ relativ vollständiges
System $S$ bilden. Dieses System $S$ suchen wir zuerst zu ermitteln;
indem wir dasselbe alsdann mit der Form $l$ überschieben, erhalten wir
das volle System $(A^{(2)})$ der Form $f$.

Auf diese Weise ist die Arbeit getheilt in zwei übersichtlichere und
einfachere Aufgaben, einmal weil das mod. $(ak)^4$ relativ vollständige
System viel kleiner ist als das gesammte System, dann aber auch, weil
die Ueberschiebungen über eine quadratische Form $l$ leicht discutirbar
sind.

252. \textit{Aufstellung der Reducenten; erster Reducent $(ab)^3$.} Ehe
wir jedoch die erste Aufgabe in Angriff nehmen, das System $S$ zu
ermitteln, wollen wir versuchen, uns über die Reducenten, sei es nun des
Systemes $S$, oder sei es des vollen Systemes $(A^{(2)})$ zu orientiren.
Dabei erinnere ich, dass ein Klammerfactorenproduct als Reducent
bezeichnet wird, sobald das aus ihm allein construirte Stammproduct,
sowie alle daraus durch Faltung entstehenden reducibel sind. (Vergl.
Nr. 124 und 144.)

So ist $(ab)^3$ Reducent. Denn das Stammproduct $(ab)^3 a_x^3 b_x^3$ ist
reducibel, weil es verschwindet, und die aus ihm durch Faltung
hervorgehenden von null verschiedenen Formen
\[
k = (ab)^4 a_x^2 b_x^2,\quad A = (ab)^6
\]
sind bekannte in das System $(A^{(2)})$ bereits aufgenommene Formen.

253. \textit{Zweiter Reducent $(ak)^3$.} Aus diesem Reducenten $(ab)^3$
lässt sich sofort ein zweiter gewinnen. Denn ersetzen wir in dem
symbolischen Product:
\[
(ab)^3 a_y^3 b_y b_x^2 = k_y^3 k_x (yx)
\]
$y$ durch $c$, so erhalten wir:
\[
(ab)^3(ac)^3(bc) b_x^2 c_x^2 = (kc)^3 c_x^3 k_x^3 = (ka)^3 a_x^3 k_x .
\]
%CONT>>


% ===== pdf idx 291, printed 277.0 =====
% !PRINTED 277
\weberpage{277}

Da die linke Seite verschwindet, wie man durch Vertauschung von $b$ mit $c$
erkennt, so ist auch das Product:
\weq{1}{(ak)^3 a_x^3 k_x = 0 = (f,\,k)^3,}
d.\,h. reducibel; aus ihm entsteht durch Faltung die zum System $(A^{(2)})$
gehörige Form
\[
l = (ak)^4 a_x^2 ;
\]
und somit ist $(ak)^3$ Reducent.

Besitzt also ein symbolisches Product den Factor $(ab)^3$ oder $(ak)^3$,
so lässt es sich im ersten Falle, von Producten mit dem Invariantenfactor
$A$ ganz abgesehen, auf Ueberschiebungen mit der Form $k$, im zweiten
Falle auf Ueberschiebungen mit der Form $l$ reduciren. Die irreduciblen
Formen dieser beiden Ueberschiebungsgruppen werden wir alsbald kennen
lernen.

254. \textit{Dritter Reducent $(k\,k')$.} Durch Combination der beiden
Reducenten $(ab)^3$ und $(ak)^3$ ergiebt sich ein dritter. Denn die
Polarenglieder $(ab)^3 a_y^3 b_x^3$ und $(ak)^3 a_y^3 k_x$ ihrer beiden
Stammproducte lassen sich in die Reihen entwickeln
\weq{1}{(ab)^3 a_y^3 b_x^3 = \frac{3}{2}\,k_{y^2}(yx) + \frac{1}{4}\,A(yx)^3}
\weq{2}{(ak)^3 a_y^3 k_x = \frac{3}{4}\,l_{y^2}(yx).}
Ersetzt man in (1) $\dot{y}$ durch $k$, in (2) aber durch $b$ und
vergleicht die beiden sich so ergebenden Ueberschiebungsrelationen, so
kommt direct
\weq{I}{(f,\,l)^2 = 2\varDelta + \frac{A}{3}\,k,}
wobei $\varDelta$ die Hesse'sche Form von $k$ ist. Wir werden später die
Form $(f,\,l)^2$, wie überhaupt alle Ueberschiebungen von $f$ mit
$l^\varrho$, $[\varrho = 1,\,2,\,3]$, in das System $(A^{(2)})$ aufnehmen.
Die Covariante $\varDelta$ ist dann eine reducible Form vermöge Gleichung
(I). Daraus folgt, dass drittens $(k\,k')$ Reducent ist. Denn
\[
(k\,k')\,k_x^3\,k_x'^{\,3} = 0
\]
und die aus diesem Stammproduct durch Faltung entstehenden sind
$\varDelta$ und die Invariante $B$. Jedes symbolische Product mit dem
Factor $(k\,k')$ lässt sich also reduciren auf solche, die entweder die
Invarianten $B$ bezw. $A$ zu Factoren haben, oder auf Ueberschiebungen
mit dem Factor $(al)^2$. Die letzteren werden wir aber, soweit wir nicht
noch deren abermalige Reduction zeigen, im System vorfinden und somit
ist $(k\,k')$ Reducent auf diese Formen.


% ===== pdf idx 292, printed 278.0 =====
% !PRINTED 278
\weberpage{278}

255. \textit{Vierter Reducent $(ab)^2(ak)^2$.} Einen vierten Reducenten
gewinnen wir durch folgende Betrachtung. Es ist nach Nr. 81 (VIII):
\weq{1}{(ak)^3 a_y a_x^2 k_x = \frac{1}{4}\,l\cdot(yx)\ \text{(wegen }(f,\,k)^3=0\text{)}}
und nach Nr. 83 (IX):
\weq{2}{(ab)\,a_y^3 a_x^2 b_x^5 = \frac{3}{2}\,(f,f)^2_{y^2}(yx) + \frac{5}{28}\,k\cdot(yx)^3.}
Ersetzen wir in (1) $y$ durch $b$, in (2) aber durch $k$, so kommt durch
Comparation der gleichen rechten Seiten:
\weq{II}{\frac{3}{2}\,\big((f,f)^2,\,k\big)^2 = \frac{1}{4}\,l\cdot f - \frac{5}{28}\,k^2.}
Es ist also $\big((f,f)^2,\,k\big)^2$ eine reducible Form, und daraus
ergiebt sich $(ab)^2(ak)^2$ als Reducent. Denn das Stammproduct
\[
P = (ab)^2(ak)^2\,a_x^2\,b_x^4\,k_x^2
\]
ist ein Glied der Ueberschiebung $\big((f,f)^2,\,k\big)^2$ und als
solches durch diese plus einer Summe niedrigerer Ueberschiebungen
darstellbar, die mindestens den Reducenten $(ab)^3$ zum Factor haben,
da sie ja aus $(f,f)^2$, resp. $k$ durch Faltung entstehen müssen (vergl.
Nr. 42). Faltet man in dem Product $P$ den Factor $a_x$ mit $b_x$, so
entsteht $(ab)^3$, $a_x$ mit $k_x$, so kommt $(ak)^3$, und endlich $b_x$
mit $k_x$, so ergeben sich mittelst des Identitätssatzes beide
Reducenten. Es ist also $P$ und jedes aus ihm durch Faltung entstehende
Product reducibel, und folglich $(ab)^2(ak)^2$ Reducent.

256. \textit{Fünfter Reducent $(ak)^2(al)^2$.} Da $(f,\,l)^2 = 2\varDelta
+ \frac{A}{3}\,k$, so erhalten wir einen weitern Reducenten durch
Ueberschiebung dieser Form mit $k$. Es ist nämlich:
\[
\big((f,l)^2,\,k\big)^2 = (ak)^2(al)^2\,a_x^2\,k_x^2 =
\Big(2\varDelta + \frac{A}{3}\,k,\ k\Big)^2.
\]
Weil aber nach der Theorie der biquadratischen Formen
\[
(\varDelta,\,k)^2 = \frac{B}{6}\,k,
\]
so erhält man
\[
(ak)^2(al)^2\,a_x^2\,k_x^2 = \frac{1}{3}\,B\cdot k + \frac{1}{3}\,A\cdot\varDelta.
\]
Das symbolische Product links ist also reducibel. Durch einmalige
Faltung desselben entsteht der Reducent $(ak)^3$ und demnach ist
$(ak)^2(al)^2$ Reducent.

257. \textit{Sechster Reducent $(ab)^2(al)^2$.} Da endlich $\varDelta$
eine reducible Form, so ist es auch $(f,\,\varDelta)^2$, und zwar findet
man ihren Werth auf folgende Weise. Es ist:
%CONT>>


% ===== pdf idx 293, printed 279.0 =====
% !PRINTED 279
\weberpage{279}
%<<CONT
\weq{1}{(ak)^3 a_x^3 k_y = -\frac{3}{4}\,l(yx)}
\weq{2}{(k\,k_1)\,k_y^3\,k_{1x}^3 = \frac{3}{2}\,\varDelta_{y^2}(yx) + \frac{B}{4}\,(yx)^3.}
Ersetzt man in (1) $y$ durch $k$, in (2) aber durch $a$, so erhält man
durch Comparation der beiden entstehenden Ueberschiebungsresultate die
Relation:
\weq{III}{\frac{3}{4}\,l\cdot k = \frac{3}{2}\,(\varDelta,\,f)^2 + \frac{B}{4}\,f.}
Anderntheils ist aber
\[
(ab)^2(al)^2\,b_x^4\,a_x^2 = \big((f,l)^2,\,f\big)^2 =
\Big(2\varDelta+\frac{A}{3}\,k,\ f\Big)^2
\]
\[
= 2(\varDelta,\,f)^2 + \frac{A}{3}\,(k,\,f)^2,
\]
also wegen (III)
\weq{IV}{(ab)^2(al)^2\,b_x^4\,a_x^2 = l\cdot k - \frac{B}{3}\,f + \frac{A}{3}\,(f,\,k)^2.}
Demnach ist das symbolische Product links reducibel, und weil jedes aus
ihm durch Faltung entstehende Product wegen des Factors $(ab)^3$
gleichfalls reducibel ist, so ist $(ab)^2(al)^2$ Reducent.

258. \textit{Rückblick.} Wir haben sonach sechs Reducenten erhalten:
\[
(ab)^3,\ (ak)^3,\ (k\,k_1),\ (ab)^2(ak)^2,
\]
\[
(ak)^2(al)^2,\ (ab)^2(al)^2.
\]
Jedes symbolische Product, das eine dieser sechs Grössen zum Factor
besitzt, ist auf andere reducibel, welche entweder eine der Invarianten
$A$ resp. $B$ zum Factor haben, oder welche durch Einführung der
Symbole $k$ resp. $l$ als einfachere Ueberschiebungen dieser Formen
über andere Formen des Systemes sich darstellen lassen. Diese
letzteren Ueberschiebungen werden wir, insoweit nicht selbst wieder
lineare Relationen unter ihnen bestehen, in das System $(A^{(2)})$
aufzunehmen haben.

Wir erhielten ferner folgende Beziehungen:
\weq{I}{2\varDelta = 2(k,\,k)^2 = (f,\,l)^2 - \frac{1}{3}\,A\cdot k}
\weq{II}{3\,l\cdot k = 6(\varDelta,\,f)^2 + B\cdot f}
\weq{III}{\big((f,l)^2,\,f\big)^2 = l\cdot k - \frac{B}{3}\,f + \frac{A}{3}\,(f,\,k)^2}
\[
(f,\,k)^3 = 0.
\]
Vermöge der Gleichung (I) erhält man hiezu noch die beiden
Invariantenrelationen:
\weq{IV}{A_u = (l,\,l)^2 = 2C + \frac{A}{3}\,B}
\weq{V}{(k,\,l^2)^4 = \frac{2}{3}\,(B^2+AC).}


% ===== pdf idx 294, printed 280.0 =====
% !PRINTED 280
\weberpage{280}

Die erste derselben entsteht aus (I), indem man diese Gleichung viermal
über $k$ schiebt:
\[
2(\varDelta,\,k)^4 = \big((f,l)^2,\,k\big)^4 - \frac{1}{3}\,AB
\]
oder:
\[
2j_k = \big((f,k)^4,\,l\big)^2 - \frac{1}{3}\,AB
\]
oder:
\[
2C = (l,\,l)^2 - \frac{1}{3}\,AB.
\]
Die zweite wird durch die viermalige Ueberschiebung von (I) über sich
selbst geliefert:
\[
4(\varDelta,\,\varDelta)^4 = \big((f,l)^2,(f,l)^2\big)^4 -
\frac{2}{3}\,A\big((f,l)^2,\,k\big)^4 + \frac{1}{9}\,A^2(k,\,k)^4
\]
oder
\weq{$\alpha$}{\frac{2}{3}\,B^2 = (ab)^4(al)^2(bl_1)^2 - \frac{2}{3}\,A\Big(2C+\frac{AB}{3}\Big) + \frac{1}{9}\,A^2 B.}
Der erste Term rechts ist aber ein Glied der Ueberschiebung
\[
(k,\,l^2)^4 = \frac{1}{3}\,(ab)^4\big\{(al)^2(bl_1)^2 + 2(al)(al_1)(bl)(bl_1)\big\},
\]
welche unter Anwendung des Productsatzes auf das zweite Glied rechts in
\weq{$\beta$}{(k,\,l^2)^4 = (ab)^4(al)^2(bl_1)^2 - \frac{1}{3}\,A\cdot A_u}
übergeht. Durch Substitution des hieraus sich ergebenden Werthes von
$(ab)^4(al)^2(bl_1)^2$ in Gleichung ($\alpha$) ergiebt sich die
Relation (V).

259. \textit{Aufstellung des Systemes $S\equiv 0$ mod. $(ak)^4$.} Nach
diesen Vorbereitungen sind wir nun im Stande, durch Ueberschiebung des
Systemes $(A^{(1)})$ über das System $(B^{(1)})$ in einfacher Weise das
neue System $S\equiv 0\ \mathrm{mod}\ (ak)^4$ herzustellen. Es ergiebt
sich nämlich wegen der Reducenten und der letzten Relationen direct,
dass hiebei überhaupt nur die Ueberschiebungen der Form $k$ über die
Formen des Systemes $(A^{(1)})$ in Betracht kommen und selbst hier nur
die Ueberschiebung der ersten Potenz von $k$.

Denn die Form $\varDelta$ ist nach Relation (I) Nr. 258 auf $(f,l)^2$,
also auf eine Form mit dem Modul $(ak)^4$, und auf $k$ reducibel, und
demnach reduciren sich die Ueberschiebungen zwischen ihr und den Formen
des Systemes $(A^{(1)})$ auf die Ueberschiebungen von $k$.

Die Form $(\varDelta,\,k)$ ist wegen der nämlichen Relation gleichfalls
auf Formen mit dem Modul $(ak)^4$ reducibel, sei es a priori oder sei es
mittelst des Reducenten $(k\,k_1)$, also ebenfalls überflüssig.

Die Formen $B=i_k$ und $C=j_k$ sind aber Invarianten. Also bleibt in der
That vom ganzen Systeme $(B^{(1)})$ nur die Form $k$ für die
%CONT>>


% ===== pdf idx 295, printed 281.0 =====
% !PRINTED 281
\weberpage{281}
%<<CONT
zu untersuchenden Ueberschiebungen mit $f$, $(f, f)^2$, $((f, f)^2, f)$ übrig.
Nun ist bereits $(ak)^3$ Reducent, und zwar ist $(ak)^3 Q \equiv 0 \bmod (ak)^4$,
wo $Q$ irgend ein symbolisches Product. Folglich sind höhere als
zweite Ueberschiebungen überflüssig, und wir können uns demnach
auf die erste Potenz von $k$ beschränken.

Durch ein- und zweimalige Ueberschiebung von $k$ über die
Formen des Systemes $(A^{(1)})$ entstehen aber die Covarianten:

\[
\begin{aligned}
&(f,\, k), \quad (f,\, k)^2, \quad ((f,\, f)^2,\, k); \quad ((f,\, f)^2,\, k)^2, \quad \bigl(((f,\, f)^2,\, f),\, k\bigr),\\
&\qquad \bigl(((f,\, f)^2,\, f),\, k\bigr)^2.
\end{aligned}
\]

Die drei letzten Formen sind überflüssig; denn $((f, f)^2, k)^2$ ist schon
nach Nr. 255 (II) als eine reducible Form erkannt worden. Die Form
$\bigl(((f, f)^2, f), k\bigr)$ ist aber reducibel als Functionaldeterminante von einer
Functionaldeterminante und im gleichen Sinne ist auch $\psi = \bigl(((f,f)^2, f), k\bigr)^2$
eine reducible Form, in so ferne die Glieder dieser Ueberschiebung
durch den Productsatz in Aggregate umgestaltet werden können, deren
einzelne Terme theils reducible, theils bereits in das System aufgenommene
 Formen zu Factoren haben. Hievon überzeugt man sich
unmittelbar auch auf folgende Weise. Gemäss den Ueberschiebungsgesetzen
 Nr. 44 ist, da $(f, f)^5 = k^{(1)} = 0$,

\[
\bigl(((f, f)^2, f), k\bigr)^2 = G_0 + \sum{}'(T^{(1)}, k)^1 + \sum{}'(T^{(2)}, k)^0 + \sum{}'(T, k^{(2)})^0,
\]

wobei das Anfangsglied $G_0 = (ab)^2 (ak)^2 (ac) a_x b_x^4 c_x^5 k_x^2$ und $T^{(1)}, T^{(2)}$;
$k^{(1)}, k^{(2)}$ Formen sind, die aus $T = ((f, f)^2, f)$ resp. $k = (f, f)^4$ durch
einmalige, resp. zweimalige Faltung entstehen. Die letzte Summe
rechts hat nur den einen Summanden $T \cdot A$; alle Summanden der
zweiten Summe zerfallen in Producte, deren einer Factor die Form $k$.
Die Formen $T^{(1)}$ in der ersten Summe sind

\[
\begin{aligned}
(ab)^3 (ac) a_x^2 b_x^3 c_x^5 &= \tfrac{1}{2} (ab)^3 a_x^2 b_x^2 c_x^5 \{(ac) b_x - (bc) a_x\} = \tfrac{1}{2} k \cdot f\\[4pt]
(ab)^2 (ac)^2 a_x^2 b_x^4 c_x^4 &= \tfrac{1}{4} \{(ab)^2 c_x^2 + (ac)^2 b_x^2 - (bc)^2 a_x^2\}^2 a_x^2 b_x^2 c_x^2 = \tfrac{1}{2} k \cdot f\\[4pt]
(ab)^2 (ac) (bc) a_x^2 b_x^3 c_x^4 &= \tfrac{1}{2} \{(ac)^2 b_x^2 + (bc)^2 a_x^2 - (ab)^2 c_x^2\} (ab)^2 a_x^2 b_x^2 c_x^2 = 0,
\end{aligned}
\]

so dass also auch die Formen $(T^{(1)}, k)^1$, wenn sie nicht verschwinden,
eine der Formen $k$ oder $f$ zu Factoren haben. $G_0$ ist aber ohnehin
reducibel nach Nr. 255 und damit die ganze rechte Seite.


% ===== pdf idx 296, printed 282.0 =====
% !PRINTED 282
\weberpage{282}
Das ganze System $S \equiv 0 \bmod (ab)^4$ besteht demnach aus den
zehn Formen:

\[
\left.
\begin{array}{l}
f, \quad (f, f)^2, \quad ((f, f)^2, f)\\[2pt]
k, \quad (f, k), \quad (f, k)^2, \quad ((f, f)^2, k)\\[2pt]
A, \quad B, \quad C.
\end{array}
\right\}\ S
\]

260. \textit{Ueberschiebung dieses Systemes $S$ über $l = (ak)^4 a_x^2$.} Unsere
letzte Aufgabe ist somit, die Ueberschiebungen der sieben Covarianten
dieses Systemes über die Covariante

\[
l = (ak)^4 a_x^2
\]

zu studiren. Dabei ist es überhaupt nicht nöthig, Producte dieser
Formen über $l$ zu schieben. Denn alle sieben Covarianten sind geraden
 Grades, und demnach würden in jeder Ueberschiebung eines Productes
 dieser Formen über Potenzen von $l$ zerfallende Glieder gebildet
 werden können.

Wir haben also nur die erste Potenz jeder der sieben Covarianten
über Potenzen von $l$ zu schieben.
Die Form $f$ liefert dadurch die sechs neuen Formen

\[
(f, l), \quad (f, l)^2, \quad (f, l^2)^3, \quad (f, l^2)^4, \quad (f, l^3)^5, \quad (f, l^3)^6.
\]

Die Form $(f, f)^2$ giebt nur zu einer neuen Form

\[
((f, f)^2, l)
\]

Veranlassung; denn bereits bei der zweiten Ueberschiebung tritt der
Reducent $(ab)^2 (al)^2$ auf. [Vgl. Nr. 257 (IV)].

Aus den Ueberschiebungen von $((f, f)^2, f)$ mit $l$ entstehen überhaupt
 keine neuen Formen; denn die erste ist als Functionaldeterminante
 von einer Functionaldeterminante reducibel; alle weiteren besitzen
 den Reducenten $(al)^2 (ab)^2$.

Als weitere Formen des Systemes ergeben sich ferner die Ueberschiebungen


\[
(k, l), \quad (k, l)^2, \quad (k, l^2)^3,
\]

während die Invariante $(k, l^2)^4$ nach (V) Nr. 258 reducibel ist.

Von den Ueberschiebungen der Form $(f, k)$ über Potenzen von $l$
liefern die ungeraden überflüssige Formen. Denn sie enthalten resp.
Glieder wie:

\[
\begin{aligned}
&(ak) (al) l_x a_x^4 k_x^3\\
&(ak) (al)^2 (al_1) l_x a_x^2 k_x^3\\
&(ak) (al)^2 (al_1)^2 (al_2) l_{2x} k_x^3\\
&(ak) (al)^2 (al_1)^2 (al_2) (kl_3)^2 l_{2x} k_x,
\end{aligned}
\]
%CONT>>


% ===== pdf idx 297, printed 283.0 =====
% !PRINTED 283
\weberpage{283}
%<<CONT
welche stets wegen des Productsatzes auf andere Formen des Systemes
reducibel sind. So ist z. B.

\[
\begin{aligned}
(ak) (al) k_x l_x \cdot a_x^4 k_x^2 &= a_x^4 k_x^2 \{(ak)^2 l_x^2 + (al)^2 k_x^2 - (kl)^2 a_x^2\}\\
&= (f, k)^2 \cdot l + (f, l)^2 \cdot k - (f, l)^2 \cdot (k, l)^2.
\end{aligned}
\]

Dagegen liefern die geraden Ueberschiebungen vier neue Formen:

\[
((f, k), l)^2, \quad ((f, k), l^2)^4, \quad ((f, k), l^3)^6, \quad ((f, k), l^4)^8.
\]

Hiezu tritt als weitere Covariante die erste Ueberschiebung der sechsten
Covariante $(f, k)^2$ über $l$, also die Form

\[
((f, k)^2, l),
\]

während bereits ihre zweite Ueberschiebung wegen des Reducenten
$(ak)^2 (al)^2$ reducibel ist.

Die siebente Covariante $((f, f)^2, k)$ endlich liefert keine neue Form
mehr. Denn ihre erste Ueberschiebung mit $l$ ist nach dem nun wiederholt
 citirten Satze als Functionaldeterminante reducibel, und jede
weitere hat den Reducenten $(ab)^2 (al)^2$.

261. \textit{Tabelle der Formen des Systemes $(A^{(2)})$.} Die Gesammtheit
aller Formen des vollen Systemes von $f = a_x^6$ ist also durch folgende
26 Formen gegeben.

1) vier gerade Invarianten:

\[
A, \ B, \ C, \ (f, l^3)^6 \text{ und eine schiefe } ((f, k), l^4)^8;
\]

2) sechs quadratische Covarianten:

\[
l, \ (k, l)^2, \ (f, l^2)^4, \ (k, l^2)^3 \ (f, l^3)^5, \ ((f, k), l^3)^6;
\]

3) fünf biquadratische Formen:

\[
k, \ (f, l)^2, \ (k, l), \ (f, l^2)^3, \ ((f, k), l^2)^4;
\]

4) fünf Covarianten sechsten Grades in $x$:

\[
f, \ (f, k)^2, {}^{*)} \ (f, l), \ ((f, k)^2, l), \ ((f, k), l)^2;
\]

5) drei Covarianten achten Grades:

\[
(f, f)^2, \ (f, k), \ ((f, f)^2, l);
\]

6) eine Covariante zehnten und eine zwölften Grades, nämlich:

\[
((f, f)^2, k), \ \text{resp.} \ ((f, f)^2, f).
\]

\bigskip
\noindent\textbf{§ 27. Relationen zwischen den Formen des Systemes von $f = a_x^6$.}
\bigskip

262. \textit{Einführung von Ueberschiebungsgliedern in das System.} Nachdem
 wir durch die vorausgegangenen Untersuchungen in den Besitz
sämmtlicher Formen des vollen Systemes von $f = a_x^6$ gelangt sind,
%CONT>>

\bigskip
\noindent\rule{4cm}{0.4pt}

\noindent{*)} $(f, k)^2 = p$ bei \gesperrt{Clebsch} \glqq Binäre Formen\grqq.


% ===== pdf idx 298, printed 284.0 =====
% !PRINTED 284
\weberpage{284}
%<<CONT
tritt an uns die Aufgabe heran, diese Formen, wenn nöthig, durch
eingliedrige symbolische Producte zu ersetzen, die für die symbolische
Rechnung als besonders tauglich erscheinen. Dies gilt insbesondere
für die fünf Invarianten des Systemes, sowie für die sechs einfachsten
Covarianten, d. i. für die sechs quadratischen Formen.

Von den Invarianten sind aber die drei ersten ohnehin einfache
symbolische Producte, nämlich

\[
\begin{aligned}
A &= (ab)^6, \quad B = (kk_1)^4,\\
C &= (kk_1)^2 (kk_2)^2 (k_1k_2)^2;
\end{aligned}
\]

diese werden wir beibehalten. Dagegen werden wir an Stelle der
Ueberschiebungen $(f, l^3)^6$, $((f, k), l^4)^8$ symbolische Producte einführen,
die sich aus dem simultanen System der sechs quadratischen Covarianten
 ergeben.

Von den quadratischen Covarianten sind die beiden Formen $l$
und $(k, l)^2$ durch je ein einfaches symbolisches Product dargestellt;
alle übrigen ersetzen wir durch eines ihrer Glieder.

263. \textit{Die quadratischen Formen.} Wie bei der Form fünften
Grades die linearen Covarianten, so spielen hier die quadratischen eine
wichtige Rolle. Auch hier ist das simultane System von vornherein
bekannt, und die erste Aufgabe ist, dessen Stellung zum vollen System
von $f = a_x^6$ kennen zu lernen. Diese Aufgabe aber deckt sich mit der
vorhin gestellten: die complicirteren quadratischen Covarianten durch
einfache symbolische Producte zu ersetzen. Setzen wir nämlich neben

\weq{1}{(ak)^4 a_x^2 = l, \quad (k, l)^2 = (kl)^2 k_x^2 = m,}

die wir a priori beibehalten, auch noch

\[
(k, m)^2 = n,
\]

dann zeigt sich, dass die in dem Systeme dieser drei Formen $l, m, n$
auftretenden Functionaldeterminanten

\weq{2}{\lambda = (m, n), \quad \mu = (n, l), \quad \nu = (l, m)}

Glieder von $(k, l^2)^3$ resp. $(f, l^3)^5$, $((f, k), l^3)^6$ sind, während $n$ an
Stelle der Ueberschiebung $(f, l^2)^4$ treten kann. Denn man hat wegen
der Relation [Nr. 258 (I)]

\[
(f, l)^2 = 2\varDelta + \tfrac{1}{3} A \cdot k
\]
% ?? uncertain: exact glyph of the script/italic Delta symbol, rendered here as \varDelta

durch zweimalige Ueberschiebung mit $l$ die Beziehung

\[
(f, l^2)^4 = 2(\varDelta, l)^2 + \tfrac{1}{3} A (k, l)^2.
\]

Hierin ist nach unseren Bezeichnungen (1) $(k, l)^2 = m$; das erste
%CONT>>


% ===== pdf idx 299, printed 285.0 =====
% !PRINTED 285
\weberpage{285}

Glied rechts ergiebt sich aber als zweite Ueberschiebung von $\Delta=(k,\ k)^2$
über $l$ aus der Reihenentwicklung

\[
(kk_1)^2 k_x^2 k_{1y}^2 = \Delta_{y^2} + \frac{1}{3} B (yx)^2 ,
\]

indem wir darin $y$ durch $l$ ersetzen; man erhält

\[
(kk_1)^2 (k_1 l)^2 k_x^2 = (\Delta,\ l)^2 + \frac{1}{3} B \cdot l_x^2
\]

oder

\[
\big((k,\ l)^2,\ k\big)^2 = (\Delta,\ l)^2 + \frac{1}{3} B \cdot l
\]

oder

\weq{3}{(m,\ k)^2 = n = (\Delta,\ l)^2 + \frac{1}{3} B \cdot l\,.}
Demnach wird:

\weq{4}{(f,\ l^2)^4 = 2n - \frac{2B}{3}\cdot l + \frac{A}{3} m\,,}
und folglich kann $(f,\ l^2)^4$ durch $n$ ersetzt werden.

Es ist ferner:

\[
\begin{aligned}
(k,\ l^2)^3 &= \frac{2}{2} (kl)^2(kl_1)k_xl_x = \big((k,\ l)^2,\ l\big) = (m,\ l) = -\nu\\[4pt]
(f,\ l^3)^5 &= \frac{3}{3} (al)^2(al_1)^2(al_2)a_xl_x = \big((f,\ l^2)^4,\ l\big)\\
&= \Big(2n-\frac{2}{3}Bl+\frac{1}{3}Am,\ l\Big) = 2(n,\ l)+\frac{1}{3}A(m,\ l)\\
&= 2\mu-\frac{A}{3}\nu\,,
\end{aligned}
\]
und endlich

\[
\begin{aligned}
-(ak)(al)^2(al_1)^2(kl_2)^2a_xk_x &= \big((kl_2)^2k_x^2,\ (al)^2(al_1)^2a_x^2\big)\\
&= \big(m,\ (f,\ l^2)^4\big)\\
&= \Big(m,\ 2n-\frac{2B}{3}l+\frac{A}{3}m\Big) = 2(m,\ n)+\frac{2}{3}B(l,\ m)\\
&= 2\lambda+\frac{2}{3}B\nu\,.
\end{aligned}
\]

Das symbolische Product links ist aber ein Glied der Ueberschiebung
$\big((f,\ k),\ l^3\big)^6$, so dass damit der lineare Zusammenhang von $\lambda$
mit dieser Form des Systemes hergestellt ist.

Die sechs quadratischen Covarianten des Systemes werden also der Reihe nach
durch die Formen:

\[
\begin{aligned}
&l = (ak)^4a_x^2, \quad m = (kl)^2k_x^2, \quad n = (km)^2k_x^2 = (kl)^2(kk_1)^2k_{1x}^2\\
&\lambda = (m,\ n), \quad \mu = (n,\ l), \quad \nu = (l,\ m)
\end{aligned}
\]
ersetzt.


% ===== pdf idx 300, printed 286.0 =====
% !PRINTED 286
\weberpage{286}

264. \textit{Die fünf Invarianten.} Von den Formen

\[
(f,\ f)^6,\quad (k,\ k)^4,\quad \big((k,\ k)^2,\ k\big)^4,\quad (f,\ l^3)^6,\quad \big((f,\ k),\ l^4\big)^8
\]
werden, wie schon Eingangs erwähnt, die drei ersten beibehalten. Die vierte
ersetzen wir durch die Invariante $(l,\ n)^2 = A_{ln} = D$, die fünfte durch
die Combinante $-(lm)(mn)(nl) = R$. Was nämlich die Form $(f,\ l^3)^6$
betrifft, so hat man:

\[
\begin{aligned}
(f,\ l^3)^6 &= (al)^2(al_1)^2(al_2)^2 = \big((al)^2(al_1)^2a_x^2,\ l_{2x}^2\big)^2\\
&= \big((f,\ l^2)^4,\ l\big)^2 = \Big(2(\Delta,\ l)^2+\frac{A}{3}m,\ l\Big)^2\\
&= \Big(2n-\frac{2}{3}Bl+\frac{A}{3}m,\ l\Big)^2\\
&= 2(n,\ l)^2-\frac{2}{3}B(l,\ l)^2+\frac{1}{3}A(m,\ l)^2\,.
\end{aligned}
\]
Diese Gleichung lehrt, dass in der That an Stelle von $(f,\ l^3)^6$ die Form
$(n,\ l)^2 = A_{ln}$ treten darf.

In Bezug auf die Invariante $\big((f,\ k),\ l^4\big)^8$ ergiebt sich aber, dass
sie sich überhaupt nur um einen numerischen Factor von der Form
$R=-(lm)(mn)(nl)$ unterscheiden kann. Denn im Systeme von $f=a_x^6$ tritt nur
diese eine fundamentale schiefe Invariante auf, und sie ist vom Grade 15 in
den Coefficienten von $f$. Die Form $R$ ist aber gleichfalls schiefe Invariante
und vom Grade 15 in den Coefficienten; sie kann also nur eine lineare Function
der ersteren sein, d. h.

\[
R = \varrho\cdot\big((f,\ k),\ l^4\big)^8\,.
\]
Die fünf Invarianten des Systemes können demnach dargestellt werden durch die
symbolischen Producte

\[
\begin{aligned}
&A=(ab)^6,\quad B=(kk_1)^4\\
&C=(kk_1)^2(k_1k_2)^2(k_2k)^2,\quad D=A_{ln}\\
&R=-(lm)(mn)(nl)\,.
\end{aligned}
\]

265. \textit{Die zweiten Ueberschiebungen von $\Delta$ und $k$ über die
quadratischen Formen $l$, $m$, $n$.} Neben den fundamentalen quadratischen
Formen interessiren uns noch eine Reihe anderer Covarianten zweiten Grades,
die wir im Folgenden berechnen wollen. Es sind das in erster Linie die
zweiten Ueberschiebungen von $\Delta$ und $k$ über $l$, $m$, $n$, so weit sie
nicht bisher schon in Betracht gezogen wurden.

Was die Ueberschiebungen von $k$ betrifft, so war

\weq{1}{\begin{aligned}
&(k,\ l)^2 = m\,, \qquad (k,\ m)^2 = n\\[2pt]
&\text{Hiezu ergiebt sich}\qquad (k,\ n)^2 = \frac{Bm}{2}+\frac{Cl}{3}\,.
\end{aligned}}


% ===== pdf idx 301, printed 287.0 =====
% !PRINTED 287
\weberpage{287}

Denn es ist

\[
(k,\ n)^2 = (kn)^2k_x^2 = (kk_1)^2(k_1m)^2k_x^2 = (kk_1)^2(k_1k_2)^2(k_2l)^2k_x^2\,.
\]
Nach den Gesetzen der Reihenentwicklung hat man aber

\[
(kk_1)^2(k_1k_2)^2k_{2y}^2k_x^2 = \frac{1}{2}Bk_{y^2}+\frac{1}{3}C(yx)^2
\]
und daraus folgt unmittelbar für $y=l$ die obige Relation. Von den zweiten
Ueberschiebungen der Form $\Delta$ über $l$, $m$, $n$ ist $(\Delta,\ l)^2$
bereits berechnet [Nr. 263 (3)]. Es war

\[
(\Delta,\ l)^2 = n-\frac{B}{3}l\,.
\]
Hiezu findet sich:

\weq{2}{\begin{aligned}
(\Delta,\ m)^2 &= \frac{B}{6}m+\frac{C}{3}l\\
(\Delta,\ n)^2 &= \frac{B}{6}n+\frac{C}{3}m\,.
\end{aligned}}
Die zweite dieser Relationen erhalten wir, indem wir in der
Reihenentwicklung

\weq{$\alpha$}{(k\Delta)^2k_y^2\Delta_x^2 = \frac{B}{6}k_{y^2}+\frac{C}{3}(yx)^2}
$y$ durch $l$ und alsdann $(k\Delta)^2(kl)^2\Delta_x^2 = \big(\Delta,\
(kl)^2k_x^2\big)^2$ durch $(\Delta,\ m)^2$ ersetzen, während die dritte aus
derselben Entwicklung $(\alpha)$ für $y=m$ hervorgeht.

266. \textit{Die Ueberschiebungen von $f$ über Producte von $l$, $m$, $n$.}
Eine weitere Reihe quadratischer Covarianten liefern die vierten
Ueberschiebungen von $f$ über $l^2$, $lm$, $ln$, $m^2$, $mn$, $n^2$. Die drei
ersten ergeben sich unmittelbar durch zweimalige Ueberschiebung der schon
früher gegebenen Relation:

\[
(f,\ l)^2 = 2\Delta + \frac{A}{3}k
\]
über $l$, $m$ und $n$ unter Benutzung der eben gefundenen Relationen (1) und
(2). Das Resultat ist:

\weq{3}{\left.\begin{aligned}
(f,\ l^2)^4 &= -\frac{2}{3}Bl+\frac{A}{3}m+2n\\
(f,\ lm)^4 &= +\frac{2}{3}Cl+\frac{B}{3}m+\frac{A}{3}n\\
(f,\ ln)^4 &= +\frac{1}{9}ACl+\Big(\frac{2}{3}C+\frac{1}{6}AB\Big)m+\frac{B}{3}n\,.
\end{aligned}\right\}}
Hiezu treten durch nachfolgende Untersuchungen die drei weiteren Relationen:
%CONT>>


% ===== pdf idx 302, printed 288.0 =====
% !PRINTED 288
\weberpage{288}
%<<CONT

\weq{4}{\left.\begin{aligned}
(f,\ m^2)^4 &= \Big\{\frac{1}{2}A_{lm}+\frac{AC}{9}\Big\}l+\frac{1}{6}\Big\{\frac{AB}{3}-A_{ll}\Big\}m+\frac{B}{3}n\\
(f,\ mn)^4 &= \Big\{\frac{1}{2}A_{ln}+\frac{BC}{9}\Big\}l+\frac{1}{2}\Big\{\frac{B^2}{2}+\frac{AC}{3}\Big\}m+\frac{1}{3}\Big\{\frac{AB}{6}-A_{ll}\Big\}n\\
(f,\ n^2)^4 &= \Big\{\frac{1}{2}A_{mn}-\frac{C^2}{9}\Big\}l+\Big\{A_{ln}-\frac{CB}{18}\Big\}m+\frac{1}{3}\Big\{\frac{B^2}{6}-2A_{lm}\Big\}n
\end{aligned}\right\}}
wenn wir die Invarianten $(l,\ l)^2$, $(l,\ m)^2$, $(l,\ n)^2$ etc. resp. mit
$A_{ll}$, $A_{lm}$, $A_{ln}$ etc. bezeichnen.

Zur Ermittelung dieser drei letzten Ueberschiebungen ist es nöthig, zuvor
$(f,\ m)^2$ und $(f,\ n)^2$ zu berechnen. Wir gehen aus von der schon in
Nr. 255 benutzten Reihenentwicklung

\weq{$\beta$}{(ak)^3a_x^2k_xa_y = \frac{l}{4}(yx)\,.}
Ersetzen wir in ihr $y$ durch $l$, so erhält man:

\[
2(ak)^3a_x^2k_x(al)l_x = \frac{l^2}{2}\,.
\]
Entwickelt man die linke Seite nach dem Productsatz, so kommt:

\[
\begin{aligned}
\frac{l^2}{2} &= a_x^2(ak)^2\big\{(ak)^2l+(al)^2k_x^2-(kl)^2a_x^2\big\}\\
&= l^2+(ak)^2(al)^2a_x^2k_x^2-(ak)^2(kl)^2a_x^4\\
&= l^2+\big((f,\ l)^2,\ k\big)^2-\big(f,\ (kl)^2k_x^2\big)^2\\
&= l^2+\Big(2\Delta+\frac{A}{3}k,\ k\Big)^2-(f,\ m)^2\\
&= l^2+\frac{B}{3}\cdot k+\frac{A}{3}\Delta-(f,\ m)^2\,,
\end{aligned}
\]
daher

\weq{5}{(f,\ m)^2 = \frac{l^2}{2}+\frac{B}{3}\cdot k+\frac{A}{3}\Delta\,.}
Ersetzt man dagegen in $(\beta)$ $y$ durch $m$, so erhalten wir

\[
2(ak)^3(am)a_x^2k_xm_x = \frac{l\cdot m}{2}\,.
\]
Wendet man wiederum auf die linke Seite den Productsatz an, dann erhält man:

\[
\begin{aligned}
\frac{lm}{2} &= (ak)^2a_x^2\big\{(ak)^2m_x^2+(am)^2k_x^2-(km)^2a_x^2\big\}\\
&= (f,\ k)^4\cdot m + \big((f,\ m)^2,\ k\big)^2-\big(f,\ (km)^2k_x^2\big)^2\\
&= lm+\Big(\frac{l^2}{2}+\frac{B}{3}k+\frac{A}{3}\Delta,\ k\Big)^2-(f,\ n)^2\,.
\end{aligned}
\]
Im mittleren Gliede rechts ist die Ueberschiebung $(l^2,\ k)^2$ zu berechnen.


% ===== pdf idx 303, printed 289.0 =====
% !PRINTED 289
\weberpage{289}

Wir benutzen hiezu die Reihe:
\weq{II}{l_x l_y \cdot l_x l_y = (l^2)_{y^2} + \frac{1}{3} A_{ll} (yx)^2.}

Ersetzt man hier $y$ durch $k$, so kommt

\[
(lk)^2 k_x^2 \cdot l_x^2 = (l^2,\, k)^2 + \frac{1}{3}\left(2C+\frac{AB}{3}\right)k_x^4,
\]

oder

\[
(l^2,\, k)^2 = lm - \frac{1}{3}\left(2C+\frac{AB}{3}\right)k.
\]

Trägt man diesen Werth in (6) ein, so erhält man:
\[
\frac{lm}{2} = lm + \left\{\frac{1}{2}\left(lm-\frac{2}{3}Ck-\frac{1}{9}ABk\right)+\frac{B}{3}\Delta+\frac{AB}{18}k\right\}-(f,\,n)^2.
\]
Daraus folgt endlich:
\weq{7}{(f,\,n)^2 = lm - \frac{1}{3}Ck + \frac{B}{3}\Delta.}

Mit Hilfe der Relationen (5) und (7) lassen sich nun in einfacher Weise die Ueberschiebungen $(f,\,m^2)^4$, $(f,\,mn)^4$, $(f,\,n^2)^4$ berechnen.

Es ist:
\[
(f,\,m^2)^4 = \frac{1}{2}(l^2,\,m)^2+\frac{B}{3}(k,\,m)^2+\frac{A}{3}(\Delta,\,m)^2.
\]

Hierin ist, wenn wir in der Reihenentwicklung (II) $y$ durch $m$ ersetzen:
\[
(l^2,\,m)^2=(l,\,m)^2\cdot l - \frac{1}{3}A_{ll}m.
\]

Ferner hatten wir in (2) Nr. 265
\[
(\Delta,\,m)^2=\frac{B}{6}m+\frac{C}{3}l,
\]
und da ausserdem $(k,\,m)^2=n$, $(l,\,m)^2=A_{lm}$, so ist endlich:
\[
(f,\,m^2)^4=\frac{1}{2}\left\{A_{lm}\cdot l - \frac{1}{3}A_{ll}\cdot m\right\}+\frac{B}{3}n+\frac{A}{3}\left\{\frac{B}{6}m+\frac{C}{3}l\right\}.
\]

Schieben wir aber die Relation (5) zweimal über $n$, so kommt:
\[
(f,\,m\cdot n)^4=\frac{1}{2}(l^2,\,n)+\frac{B}{3}(k,\,n)^2+\frac{A}{3}(\Delta,\,n)^2.
\]

Aus der Reihe (II) folgt wiederum für $y=n$
\[
(l^2,\,n)=(l,\,n)^2\cdot l - \frac{1}{3}A_{ll}n.
\]

Die Ueberschiebung wird also unter Berücksichtigung der Relationen (1) und (2) Nr. 265
\[
(f,\,mn)^4=\frac{1}{2}\left\{A_{ln}l-\frac{1}{3}A_{ll}n\right\}+\frac{B}{3}\left(\frac{B}{2}m+\frac{C}{3}l\right)+\frac{A}{3}\left(\frac{B}{6}n+\frac{C}{3}m\right)\cdot
\]

Endlich ergiebt sich aus der Gleichung (7):
\[
(f,\,n^2)^2=(lm,\,n)^2-\frac{1}{3}C(k,\,n)^2+\frac{B}{3}(\Delta,\,n)^2.
\]


% ===== pdf idx 304, printed 290.0 =====
% !PRINTED 290
\weberpage{290}

Zur Berechnung des ersten Gliedes der rechten Seite benutzen wir die Reihe
\[
lm_y^2+ml_y^2=2(lm)_{y^2}+\frac{2}{3}A_{lm}(yx)^2.
\]

Ersetzt man hier $y$ durch $n$, so kommt:
\[
2(lm,\,n)^2=lA_{mn}+mA_{ln}-\frac{2}{3}A_{lm}n.
\]

Demnach wird
\[
(f,\,n^2)^2=\frac{1}{2}\left\{A_{mn}l+A_{ln}m-\frac{2}{3}A_{lm}n\right\}-\frac{C}{3}\left\{\frac{C}{3}l+\frac{B}{2}m\right\}+\frac{B}{3}\left\{\frac{B}{6}n+\frac{C}{3}m\right\}.
\]

Damit wollen wir die Berechnung von quadratischen Covarianten der Form sechster Ordnung abschliessen, und zur Berechnung der Invarianten der quadratischen Fundamentalcovarianten übergehen.

267. \textit{Berechnung der Invarianten des simultanen Systemes $l$, $m$, $n$.} Die drei quadratischen Covarianten besitzen ausser der schiefen Invariante $R=(lm)(mn)(ln)$ noch sechs simultane Invarianten (vgl. Nr. 131)
\[
A_{ll},\quad A_{lm},\quad A_{ln},\quad A_{mm},\quad A_{mn},\quad A_{nn},
\]
die wir nun im Folgenden durch die Fundamentalinvarianten des Systemes von $f=a_x^6$ darstellen wollen.

Wir hatten bereits früher erhalten (vgl. Nr. 258 und Nr. 264)
\weq{1}{\left.\begin{aligned} A_{ll}&=2C+\frac{1}{3}AB\\ A_{lm}&=(k,\,l^2)^4=\frac{2}{3}(B^2+AC)\\ A_{ln}&=D. \end{aligned}\right\}}

Hiezu ergiebt sich:
\weq{2}{\left.\begin{aligned} A_{mm}&=D\\ A_{mn}&=\frac{1}{3}\left\{B^3+2C^2+\frac{4}{3}ABC\right\}\\ A_{nn}&=\frac{1}{2}BD+\frac{2}{9}C(B^2+AC) \end{aligned}\right\}}

Denn man hat:
\[
A_{mm}=((k,\,l)^2,\,m)^2=(k,\,ml)^4=((k,\,m)^2,\,l)^2=(n,\,l)^2
\]
\[
A_{mn}=((k,\,l)^2,\,n)^2=(k,\,nl)^4=((k,\,n)^2,\,l)^2=\left(\frac{B}{2}m+\frac{C}{3}l,\ l\right)^2
\]
\[
A_{nn}=((k,\,m)^2,\,n)^2=(k,\,mn)^4=((k,\,n)^2,\,m)^2=\left(\frac{B}{2}m+\frac{C}{3}l,\ l\right)^2.
\]


% ===== pdf idx 305, printed 291.0 =====
% !PRINTED 291
\weberpage{291}

268. \textit{Die Relationen für die schiefe Invariante.} Aus der Theorie der simultanen quadratischen Formen ergiebt sich unmittelbar, dass (vgl. Nr. 133) $R$ auch dargestellt werden kann durch
\[
R=A_{l\lambda}=A_{m\mu}=A_{n\nu},
\]
und dass ferner zwischen $R^2$ und den übrigen fundamentalen Covarianten eine Relation bestehen muss, welche aus
\[
2R^2=\begin{vmatrix} A_{ll}, & A_{lm}, & A_{ln} \\ A_{ml}, & A_{mm}, & A_{mn} \\ A_{nl}, & A_{nm}, & A_{nn} \end{vmatrix}
\]
durch Substitution der in Nr. 267 berechneten Werthe der Determinantenelemente hervorgeht.

Weitere Relationen will ich nicht mehr anführen. Die schon früher in der Theorie der quadratischen Formen gewonnenen zahlreichen Beziehungen lassen sich hier direct übertragen. In ähnlicher Weise würde man auch die Theorie der biquadratischen Formen benutzen können, um Relationen zwischen den Covarianten vierten Grades des Systemes aufzustellen.

\S\,28. \textbf{Die Discriminante der Form sechsten Grades und irrationale Relationen.}

269. \textit{Vorbereitende Untersuchungen.} Die erste Berechnung der Discriminante einer Form sechsten Grades haben wir von Brioschi. Er benutzte hiebei die allgemeinen Differentialgleichungen, denen Invarianten überhaupt genügen, und charakterisirte dann die Discriminante, indem er ihnen specielle hinzufügte, welche nur durch Discriminanten befriedigt werden (Crelle's Journal Bd. 53. Vgl. auch Faà di Bruno, \glqq Theorie der binären Formen\grqq{}, deutsche Uebersetzung von Walter, S. 275).

Die genaue Kenntniss des vollen Systemes von $f$ gestattet uns indess hier wiederum den nun schon mehrmals eingeschlagenen directen Weg zu verfolgen (vgl. Discrim. der Form vierten und fünften Grades), indem wir uns fragen:

\begin{quote}
\glqq Welche Relation besteht zwischen den Invarianten $A$, $B$, $C$, $D$, wenn $f=\alpha_x^2\cdot\varphi$, wo $\alpha_x$ eine lineare und $\varphi$ eine biquadratische Form ist?\grqq{}
\end{quote}

Wie bei der Form fünften Grades, und wie überhaupt allgemein (vgl. Nr. 99), gehen aus der Voraussetzung $f=\alpha_x^2\cdot\varphi$ folgende vier wichtige Beziehungen hervor:
%CONT>>


% ===== pdf idx 306, printed 292.0 =====
% !PRINTED 292
\weberpage{292}
%<<CONT

\weq{1}{(a\alpha)^6=0,\quad (a\alpha)^5a_x=0,}
\weq{2}{(a\alpha)^4a_x^2=\lambda\cdot\alpha_x^2,}
\weq{3}{\frac{5}{2}(ab)^2(a\alpha)^4b_x^4 - \frac{5}{2}(k\alpha)^2k_x^2\cdot\alpha_x^2=\frac{1}{3}A\cdot\alpha_x^4,}

wobei Gleichung (3) aus $(a\alpha)^4a_x$ durch die in Nr. 99 angegebene Reihenentwicklung gewonnen wird. Der Werth des Parameters $\lambda$ berechnet sich gleichfalls wie früher. Wir ersetzen in (2) $x$ durch $b$ und multipliciren mit $b_x^4$, dann erhalten wir:
\[
(a\alpha)^4(ab)^2b_x^4=\lambda\cdot(f,\,\alpha^2)^2
\]
und demnach wegen (3)
\weq{4}{\lambda\cdot(f,\,\alpha^2)^2=(k\alpha)^2k_x^2\cdot\alpha_x^2+\frac{2}{15}A\cdot\alpha_x^4.}

Ueberschiebt man aber diese Relation zweimal über $\alpha_x^2$, so kommt in Verbindung mit Gleichung (2) direct:
\weq{5}{\lambda^2=\frac{1}{6}(k\alpha)^4.}

Die nächste Aufgabe wäre nun, wiederum auch eine Relation für die erste Ueberschiebung $(f,\,\alpha)$ zu ermitteln und auf Grund der so berechneten Ueberschiebungen von $f$ und $\alpha$ das Formensystem der simultanen Invarianten von $f$ und $\alpha$ mit den Invarianten $A$, $B$, $C$, $D$ in Beziehung zu setzen. Hiebei zeigt sich aber, dass insbesondere die höheren Ueberschiebungen von $k$, $\Delta$ und $l$ über $\alpha$ die hervorragendste Rolle spielen, und wir suchen demnach hier von vornherein sie mit dem allgemeinen Formensystem in Beziehung zu setzen.

270. \textit{Berechnung von $(l\alpha)^2$ und $(\Delta\alpha)^2$.} Zur Berechnung der Ueberschiebungen von $l$ über $\alpha$ benutzen wir am geschicktesten die Relation (4), indem wir sie dreimal über $k$ schieben. Dann kommt:
\weq{1}{\begin{aligned} \lambda(a\alpha)^2(ak)^3a_xk_x&=\frac{1}{2}\left\{(kk_1)^2(\alpha k_1)\alpha_x+(kk_1)(\alpha k_1)^2k_x\right\}(k\alpha)^2k_{1x}-\frac{2}{15}A(k,\,\alpha^3)^3\alpha_x\\ &=-\frac{1}{2}(\Delta,\,\alpha^3)^3\alpha_x-\frac{2}{15}A(k,\,\alpha^3)^3\alpha_x,\end{aligned}}

weil das zweite Glied der Klammer rechts identisch verschwindet und $\Delta=(kk_1)^2k_x^2k_{1x}^2$ ist. Anderntheils ist aber, da nach Nr. 253, $(f,\,k)^3=0$,
\[
(ak)^3a_y^2k_xa_x=-\frac{1}{2}l_xl_y\cdot(xy),
\]
und demnach
\weq{2}{(a\alpha)^2(ak)^3a_xk_x=\frac{1}{2}(l,\,\alpha)\cdot\alpha_x.}

Führt man diesen Werth in (1) ein, so kommt nach Division mit $\frac{1}{2}\alpha_x$:
\weq{3}{\lambda\cdot(l,\,\alpha)=-\left\{(\Delta,\,\alpha^3)^3+\frac{4}{15}A\cdot(k,\,\alpha^3)^3\right\}}


% ===== pdf idx 307, printed 293.0 =====
% !PRINTED 293
\weberpage{293}

und hieraus durch abermalige Ueberschiebung mit $\alpha$:
\weq{4}{\lambda\cdot(l,\ \alpha^2)^2=-\left\{(\varDelta,\ \alpha^4)^4+\frac{4}{15}A\cdot(k,\ \alpha^4)^4\right\}\cdot}
Eine zweite Gleichung zwischen $(l\alpha)^2$ und $(\varDelta\alpha)^4$ erhalten wir durch
viermalige Ueberschiebung von Relation (4) mit der Form $k$, nämlich
\weq{5}{\lambda\cdot(l\alpha)^2=(\varDelta,\ \alpha^4)^4+\frac{2}{15}A\cdot(k,\ \alpha^4)^4.}
Indem wir beide Gleichungen addiren, erhalten wir wegen $(k\alpha)^4=6\lambda^2$
\weq{6}{(l\alpha)^2=-\frac{2}{5}A\cdot\lambda,}
und indem wir sie subtrahiren
\[
(\varDelta\alpha)^4=-\frac{6}{5}A\cdot\lambda^2.
\]
Nun lässt sich aber Gleichung (3) auch in der Form darstellen
\[
\left(\lambda\cdot l+(\varDelta,\ \alpha^2)^2+\frac{4}{15}A\cdot(k,\ \alpha^2)^2,\ \alpha\right)=0,
\]
d.\,h. (vgl. Nr. 99) die quadratische Form
\[
\lambda\cdot l_x^2+(\varDelta\alpha)^2\varDelta_x^2+\frac{4A}{15}(k\alpha)^2k_x^2
\]
ist eine Potenz von $\alpha$, so dass also, wenn $\varrho$ eine Proportionalitätsconstante
, die Relation besteht
\weq{7}{\varrho\cdot\alpha_x^2-\lambda\cdot l_x^2=(\varDelta\alpha)^2\varDelta_x^2+\frac{4}{15}A\cdot(k\alpha)^2k_x^2.}
Die Constante $\varrho$ finden wir sofort, wenn wir diese Relation (7) zweimal
über sich selbst schieben. Hiebei ergiebt sich links, da
\[
(l,\ l)^2=2C+\frac{1}{3}AB,
\]
der Werth:
\weq{8}{-2\lambda\cdot\varrho(l\alpha)^2+\lambda^2\left(2C+\frac{1}{3}AB\right);}
rechts dagegen erhalten wir:
\[
(\varDelta\alpha)^2(\varDelta_1\alpha)^2(\varDelta\varDelta_1)^2+\frac{8}{15}A\cdot(\varDelta\alpha)^2(k\alpha)^2(\varDelta k)^2+\frac{16}{225}A^2\cdot(k\alpha)^2(k_1\alpha)^2(kk_1)^2,
\]
oder, weil nach der Theorie der biquadratischen Formen $(\varDelta,\ k)^2=\frac{B}{6}k$
und $(\varDelta,\ \varDelta)^2=\frac{C}{3}k-\frac{B}{6}\varDelta$:
\weq{9}{\frac{C}{3}(k\alpha)^4-\frac{B}{6}(\varDelta\alpha)^4+\frac{8}{90}AB\cdot(k\alpha)^4+\frac{16}{225}A^2\cdot(\varDelta\alpha)^4.}
Tragen wir nun in (8) und (9) die Werthe von $(k\alpha)^4$, $(\varDelta\alpha)^4$, $(l\alpha)^2$
ein und setzen beide Ausdrücke wieder einander gleich, so kommt
nach einfacher Reduction:
\weq{10}{\varrho=\frac{1}{2}B-\frac{8}{75}A^2.}


% ===== pdf idx 308, printed 294.0 =====
% !PRINTED 294
\weberpage{294}

271. \textit{Die Discriminante.} Die Relation (7) gestattet nun aber
auch die Invariante $D=(mm_1)^2$ mit den simultanen Invarianten
$(\varDelta\alpha)^4$ und $(k\alpha)^4$ in Beziehung zu setzen und damit erreichen wir zugleich
 das Ziel unserer Aufgabe. Ueberschiebt man nämlich die erwähnte
 Relation (7) zweimal über $k$, so erhält man wegen $(k,\ l)^2=m$
zunächst:
\weq{1}{\varrho\cdot(k,\ \alpha^2)^2-\lambda\cdot m=(\varDelta\alpha)^2(\varDelta k)^2k_x^2+\frac{4}{15}A\cdot(k\alpha)^2(k_1k)^2k_x^2.}
Die Werthe der beiden Terme rechts ergeben sich aus den entsprechenden
 Reihenentwicklungen:
\[
\begin{aligned}
(\varDelta k)^2\varDelta_y^2k_x^2 &= \{(\varDelta,\ k)^2\}_{y^2}+\frac{1}{3}(\varDelta,\ k)^4(yx)^2=\frac{B}{6}k_x^2k_y^2+\frac{C}{3}(yx)^2\\
(kk_1)^2k_x^2k_{1y}^2 &= \{(k,\ k)^2\}_{y^2}+\frac{1}{3}(k,\ k)^4(xy)^2=\varDelta_x^2\varDelta_y^2+\frac{B}{3}(yx)^2,
\end{aligned}
\]
indem man darin $y$ durch $\alpha$ ersetzt. Führt man alsdann diese
Werthe in (1) ein und ersetzt $\varrho$ durch $\frac{1}{2}B-\frac{8}{75}A^2$, so kommt
nach Multiplication mit $-3$:
\[
3\lambda m=\left(B-\frac{8}{25}A^2\right)(k\alpha)^2k_x^2-\frac{4}{5}A\cdot(\varDelta\alpha)^2\varDelta_x^2-\left(C+\frac{4}{15}AB\right)\alpha_x^2.
\]
Bilden wir nun auf beiden Seiten dieser Gleichung
\[
3\lambda m=P\cdot(k\alpha)^2k_x^2+Q\cdot(\varDelta\alpha)^2\varDelta_x^2+R\cdot\alpha_x^2
\]
die Hesse'sche Form, so erhalten wir
\[
\begin{aligned}
9\lambda^2D &= P^2(\varDelta\alpha)^4+Q^2((\varDelta,\ \varDelta)^2,\ \alpha^4)^4\\
&\quad+2P\cdot Q\cdot((\varDelta,\ k)^2,\ \alpha^4)^4+2P\cdot R\cdot(k\alpha)^4+2QR(\varDelta\alpha)^4.
\end{aligned}
\]
Ersetzen wir hierin $(\varDelta,\ \varDelta)^2$ und $(\varDelta,\ k)^2$ durch ihre Werthe
$\frac{C}{3}k-\frac{B}{6}\varDelta$, resp. $\frac{B}{6}k$ und ordnen alsdann nach den Ueberschiebungen
$(\varDelta\alpha)^4$ und $(k\alpha)^4$, so geht diese Gleichung über in:
\[
9\lambda^2D=(\varDelta\alpha)^4\left\{P^2-\frac{B}{6}Q^2+2QR\right\}+(k\alpha)^4\left\{\frac{C}{3}\cdot Q^2+\frac{B}{3}\cdot PQ+2PR\right\}
\]
oder wenn wir nun $(\varDelta\alpha)^2$, $(k\alpha)^2$, $P$, $Q$, $R$ durch ihre Werthe ersetzen
und mit $\lambda^2$ dividiren: % ?? uncertain: faint mark over the middle equals sign a few lines above may be a printing defect rather than a notational symbol; reproduced here as a plain \textquotedbl{}=\textquotedbl{}
\[
\begin{aligned}
9D &= -\frac{6}{5}A\left\{B^2-\frac{8}{25}\cdot A^2B+\frac{64}{625}A^4+\frac{8}{5}AC\right\}\\
&\quad+6\left\{\frac{4}{3}\cdot\frac{16}{25}A^2C-\frac{4}{5}AB^2+\frac{32}{125}A^3B-2BC\right\}.
\end{aligned}
\]
Das ist die Invariantenrelation, welche die Voraussetzung $f=\alpha_x^2\cdot\varphi$
zur Folge hat. Wir können sie auf die einfachere Form bringen,
indem wir $\frac{2}{5}A=A_1$ setzen:
%CONT>>


% ===== pdf idx 309, printed 295.0 =====
% !PRINTED 295
\weberpage{295}
%<<CONT

\[
12A_1^5-30A_1^3B-20A_1^2C+15A_1B^2+12BC+9D=0,
\]
und die linke Seite dieser Gleichung müssen wir als Discriminante
von $f$ bezeichnen, mit der sie auch in Bezug auf den Grad in den
Coefficienten übereinstimmt.

272. \textit{Irrationale Relationen bei Combinanten.} Ehe wir die Untersuchungen
 über Co- und Invariantenrelationen verlassen, will ich
noch auf eine Gattung von irrationalen Beziehungen zwischen Combinanten
 hinweisen, auf welche bereits \gesperrt{Brill}, \gesperrt{Stephanos}, \gesperrt{Meyer}
und \gesperrt{Stroh} aufmerksam gemacht haben.

Sind nämlich $f=r_x^n$ und $\varphi=s_x^n$ zwei binäre Formen $n^{\mathrm{ten}}$ Grades,
so ist nach einem Satze von \gesperrt{Gordan} (vgl. Nr. 68 und Math. Ann.
Bd. V) jede Combinante der beiden Formen Covariante des Cayley'schen
Ausdruckes
\[
\varPsi=f(x)\varphi(y)-\varphi(x)f(y),
\]
also seiner Elementarcovarianten. Diese Elementarcovarianten sind
in unserem Falle die ungeraden Ueberschiebungen von $f$ über $\varphi$, da
die geraden sich gegenseitig aufheben, also die Formen:
\[
(f,\ \varphi),\quad (f,\ \varphi)^3,\quad (f,\ \varphi)^5\ \ldots\ \text{etc.}
\]
Aber diese Covarianten können nicht von einander unabhängig sein.
Denn wenn $\varphi$ proportional $f$ ist, so verschwindet $\varPsi$ und demnach alle
Elementarcovarianten, und umgekehrt: Ist $\varPsi$ identisch null, so ist
$f=\varrho\cdot\varphi$, wo $\varrho$ eine Constante. Damit aber $f$ und $\varphi$ einander proportional
 sind, ist es bekanntlich schon hinreichend, dass ihre Functionaldeterminante
 $(f,\ \varphi)$ identisch verschwindet. Wenn also schon das
Verschwinden der ersten Elementarcovariante auch das Verschwinden
von $\varPsi$ und damit aller übrigen Elementarcovarianten nach sich zieht,
so können diese letzteren nicht mehr unabhängig sein von $(f,\ \varphi)$.
Es müssen vielmehr die Coefficienten der späteren Ueberschiebungen
homogene und irrationale Functionen der Coefficienten von $(f,\ \varphi)$ sein.

Man kann daher sich die Fragen stellen: Welches sind diese
Coefficientenrelationen und wie gelangt man zu ihnen? Was die zweite
Frage betrifft, so kann man zu ihrer Beantwortung einen schon mehrfach
 benutzten Weg einschlagen. Wir bilden eine Covariante von
$(f,\ \varphi)$ durch Ueberschiebung dieser Form über sich selbst, und sehen
nach, ob sich diese Form nicht auch simultan durch eine Ueberschiebung
 von $F=(f,\ \varphi)$ über eine der übrigen Elementarcovarianten
darstellen lässt. Dies ist sicher der Fall, wenn z.\,B. ein Glied der
Ueberschiebung $(F,\ F)^\nu$ verschwindet, da man weiss, dass jedes Glied
einer Ueberschiebung nach dieser und niedrigeren Ueberschiebungen
von $F$ über andere Formen entwickelt werden kann.


% ===== pdf idx 310, printed 296.0 =====
% !PRINTED 296
\weberpage{296}

Nun tritt bereits für $\nu=4$, also in $(F,\ F)^4$ ein solches Glied
auf. Es ist, wenn $F=(f,\ \varphi)=(rs)r_x^{n-1}s_x^{n-1}$, das Glied
\[
(rs)(r_1s_1)(rr_1)(sr_1)(rs_1)(ss_1)r_x^{n-3}s_x^{n-3}r_{1x}^{n-3}s_{1x}^{n-3},
\]
welches verschwindet, weil es durch Vertauschung von $r$ mit $r_1$ nur
sein Zeichen ändert. Es muss also eine Relation bestehen von der
Form
\weq{I}{0=(F,\ F)^4+\sum_{\nu=0}^{\nu=3}\varrho_\nu(F,\ \varPhi_\mu)^\nu,}
wobei $\varPhi_\mu$ irgend welche spätere Elementarcovarianten niedrigerern % [sic]
Grades als $F$ sind. Durch Coefficientenvergleichung findet man alsdann
 jene Relationen, welche die Coefficienten der Formen $\varPhi_\mu$ mit
denen von $F$ in Beziehung setzen.

273. \textit{Beispiel.} Wir wollen hier an einem speciellen Falle, in
welchem die Functionaldeterminante $F=(f,\ \varphi)$ eine Form sechsten
Grades ist, diese allgemeinen Betrachtungen erläutern. Die Formen
$f$ und $\varphi$ sind alsdann zwei biquadratische Formen
\[
f=r_x^4,\quad \varphi=s_x^4,
\]
und ihre Combinanten entstehen durch Ueberschiebung der beiden
einzig vorhandenen Elementarcombinanten
\[
\begin{aligned}
F &= (rs)r_x^3s_x^3=a_x^6\\
\varPhi &= (rs)^3r_xs_x.
\end{aligned}
\]
Nach dem Vorausgegangenen sind die Coefficienten der Form $\varPhi$
irrationale Functionen der Coefficienten von $F$, und unsere Aufgabe
ist, die Beziehungen zwischen diesen Coefficienten herzustellen. Da
wir es hier ausser der Combinante $F$ nur noch mit einer Combinante
$\varPhi$ zu thun haben, so wird die oben aufgestellte allgemeine Relation
(I) in einfacher Weise zum Ziele führen. Wir bilden demnach $(F,\ F)^4$,
also die Covariante $k$ von $F=a_x^6$. Da diese Form $k=(ab)^4a_x^2b_x^2$
aus der gemischten Polare $a_x^2a_y^2a_z^2$ für $y=z=b$ hervorgeht, so wollen
wir zunächst diese gemischte Polare in Symbolen $r$, $s$ darstellen.

Wir substituiren zu dem Zwecke in der Identität:
\[
a_x^6=(rs)r_x^3s_x^3
\]
für $x_i$ den Werth $\lambda_1x_i+\lambda_2y_i+\lambda_3z_i$, dann kommt:
\[
(\lambda_1a_x+\lambda_2a_y+\lambda_3a_z)^6=(rs)(\lambda_1r_x+\lambda_2r_y+\lambda_3r_z)^3(\lambda_1s_x+\lambda_2s_y+\lambda_3s_z)^3.
\]
Durch Comparation der Glieder mit dem Factor $\lambda_1^2\lambda_2^2\lambda_3^2$ erhält man:
\[
90a_x^2a_y^2a_z^2=36(rs)r_xr_yr_zs_xs_ys_z+9(rs)\sum r_x^2r_ys_ys_z^2,
\]
%CONT>>


% ===== pdf idx 311, printed 297.0 =====
% !PRINTED 297
\weberpage{297}
%<<CONT
wobei sich die Summe über die sechs möglichen Producte von der Gestalt $r_x^2 r_y s_y s_z^2$ erstreckt. Addiren wir rechts sechsmal das Product $9\,(rs)\,r_x r_y r_z s_x s_y s_z$, indem wir gleichzeitig jeden solchen Summanden von je einem Gliede unter dem Summenzeichen subtrahiren, so findet man nach zweimaliger Anwendung der Identitäten

\[
r_x s_y - s_x r_y = (rs)(xy), \quad r_x s_z - s_x r_z = (rs)(xz), \quad r_y s_z - r_z s_y = (rs)(yz)
\]
die Relation:
\[
90\,a_x^6\,a_y^2\,a_z^2 = 90\,(rs)\,r_x r_y r_z s_x s_y s_z + 9\,(rs)^3 \sum (xy)^2 r_z s_z.
\]
Dividiren wir dieselbe durch 45 und setzen nun
\[
\frac{2}{5}\,\Phi = \frac{2}{5}\,(rs)^3\,r_x s_x = p_x^2\text{, etc.},
\]
so kommt:
\[
2\,a_x^2 a_y^2 a_z^2 = 2\,(rs)\,r_x r_y r_z s_x s_y s_z + \frac{1}{2}\left\{ p_x^2 (yz)^2 + p_y^2 (xz)^2 + p_z^2 (xy)^2 \right\},
\]
oder, indem wir den Productsatz
\[
p_y p_z (xy)(xz) + \frac{1}{2} p_x^2 (yz)^2 = \frac{1}{2}\left\{ p_y^2(xz)^2 + p_z^2(xy)^2 \right\}
\]
auf den zweiten Term rechts anwenden:
\weq{II}{
2\,a_x^2 a_y^2 a_z^2 = 2\,(rs)\,r_x r_y r_z s_x s_y s_z + p_x^2 (yz)^2 + p_y p_z (xy)(xz).
}
Ersetzen wir hierin $y$ und $z$ durch $b$ und multipliciren mit $b_x^2$, so kommt:
\weq{III}{
2\,(ab)^4 a_x^2 b_x^2 = 2\,(rs)(rb)^2(sb)^2 r_x s_x b_x^2 + (bp)^2 b_x^4.
}
Der zweite Term rechts ist nichts anderes als $\frac{2}{5}(F,\Phi)^2$. Der Werth des ersten Termes ergiebt sich aus derselben Gleichung (II), wenn wir darin $y$ durch $r$, $z$ durch $s$ ersetzen und mit $(rs)\,r_x s_x$ multipliciren; man erhält alsdann:
\weq{1}{
2\,(rs)(ar)^2(as)^2 a_x^2 r_x s_x = p_x^2\cdot(rs)^3 r_x s_x + (pr)(ps)(rs)\,r_x^2 s_x^2.
}
Hierin ist der zweite Term rechts ein Glied der Ueberschiebung
\[
(f,p)^2 = (rs)(r_x^3 s_x^3, p_x^2)^3 = \frac{(rs)}{5}\left\{ (pr)^2 r_x s_x^3 + (ps)^2 s_x r_x^3 + 3(pr)(ps) r_x^2 s_x^2 \right\}.
\]
Ersetzt man in der Klammer rechts die Grösse $(pr)^2 s_x^2 + (ps)^2 r_x^2$ nach dem Productsatz durch $2(pr)(ps) r_x s_x + (rs)^2 p_x^2$, so kommt:
\weq{2}{
(f,p)^2 = (ps)(pr)(rs)\,r_x^2 s_x^2 + \frac{1}{5}(rs)^3\,r_x s_x\cdot p_x^2.
}
Führt man also in die Gleichung (1) den aus dieser Gleichung (2) sich ergebenden Werth von $(ps)(pr)(rs)\,s_x^2 r_x^2$ ein, so erhalten wir
\[
2\,(rs)(ar)^2(as)^2\,a_x^2 r_x s_x = \frac{4}{5}\,p_x^2\cdot(rs)^3\,r_x s_x + (f,p)^2,
\]
%CONT>>


% ===== pdf idx 312, printed 298.0 =====
% !PRINTED 298
\weberpage{298}
%<<CONT
oder, weil $p_x^2 = p = \frac{2}{5}(rs)^3 r_x s_x$,
\[
2\,(rs)(ar)^2(as)^2\,a_x^2 r_x s_x = 2p^2 + (f,p)^2 = \frac{8}{25}\,\Phi^2 + \frac{2}{5}(F,\Phi)^2.
\]
Durch Substitution dieses Werthes in (III) ergiebt sich endlich nach Division mit dem Factor 2:
\weq{IV}{
(F,F)^4 = \frac{4}{25}\,\Phi^2 + \frac{2}{5}(F,\Phi)^2,
}
und dies ist die gesuchte Fundamentalrelation zwischen den Coefficienten von $F$ und $\Phi$.

274. Führen wir in dieselbe die Symbole $a$, $k$ und $p$ von $F$ resp. $(F,F)^4$ und $\frac{2}{5}\Phi$ ein, so nimmt sie die Gestalt an
\weq{3}{
k = p^2 + (ap)^2\,a_x^4.
}
Durch Vergleichung der Coefficienten von $x$ auf beiden Seiten dieser Identität erhält man fünf Gleichungen zur Berechnung der drei Grössen $p_0\,p_1\,p_2$. Doch lässt sich dieses System in einfacher Weise auf eine einzige Gleichung für die Discriminante $P=(pp_1)^2$ der quadratischen Form $p$ reduciren. Ich überschiebe diese Identität zunächst viermal über $f$, und erhalte, da $(k,f)^4=(ak)^4 a_x^2=l$,
\[
l = (f,p^2)^4 + (f,(f,p)^2)^4 = (ap)^2(ap_1)^2 a_x^2 + (ab)^4(bp)^2 a_x^2.
\]
Der Werth des ersten Termes rechts folgt aus der Fundamentalrelation (3), indem man sie zweimal über $p$ schiebt:
\[
(f,p^2)^4 = (k,p)^2 - \frac{2}{3}\,P\cdot p.
\]
Um das zweite Glied zu berechnen, bilde ich die Ueberschiebung
\[
\begin{aligned}
(k,p)^2 &= (ab)^4(a_x^2 b_x^2,\, p_x^2) = \frac{(ab)^4}{3}\left\{ (bp)^2 a_x^2 + 2(ap)(bp)\,a_x b_x \right\} \\
&= \frac{(ab)^4}{3}\left\{ (bp)^2 a_x^2 + (ap)^2 b_x^2 + (bp)^2 a_x^2 - (ab)^2 p \right\}.
\end{aligned}
\]
Hieraus folgt direct:
\[
(ab)^4(bp)^2 a_x^2 = (f,(f,p)^2)^4 = (k,p)^2 + \frac{A}{3}\,p.
\]
Demnach wird
\weq{V}{
l = 2(k,p)^2 + \frac{1}{3}\,p\,(A - 2P).
}
Bildet man nun die zweite Ueberschiebung dieser Relation über $p$, so findet man:
\weq{4}{
(lp)^2 = 2(k,p^2)^4 + \frac{P}{3}(A-2P).
}
Die Fundamentalrelation (3) liefert aber für $(k,p^2)^4$ unmittelbar durch viermalige Ueberschiebung über $k$:
\[
(k,p^2)^4 = B - (lp)^2.
\]


% ===== pdf idx 313, printed 299.0 =====
% !PRINTED 299
\weberpage{299}
Trägt man diesen Werth in die letzte Gleichung (4) ein, so kommt:
\weq{VI}{
3(lp)^2 = 2B + \frac{P}{3}(A-2P).
}
Gelingt es uns nun, die Ueberschiebung $(lp)^2$ noch auf eine zweite Art darzustellen, so erhalten wir in Verbindung mit Gleichung (VI) eine Relation für $P$.

275. Nun erinnern wir uns aber daran, dass die zweimalige Ueberschiebung von $\Delta$ über $l$ dargestellt ist durch (vergl. Nr. 265, 2)
\[
(\Delta,l)^2 = n - \frac{B}{3}\,l,
\]
dass also rechts wiederum die Covariante $l$ auftritt.

Wenn wir daher die Relation (V) zweimal über $\Delta$ schieben, so erhalten wir eine neue Relation für $l$, nämlich:
\weq{5}{
n - \frac{1}{3}\,Bl = 2\big((k,p)^2,\Delta\big)^2 + \frac{1}{3}(\Delta,p)^2(A-2P).
}
Es ist aber:
\[
\big((k,p)^2,\Delta_x^4\big)^2 = (kp)^2(k\Delta)^2\Delta_x^2.
\]
Der Term rechts ist ein Glied der Ueberschiebung
\[
\big((k,\Delta)^2,p\big)^2 = \left(\frac{B}{6}\,k,\,p\right)^2.
\]
Denn man hat:
\[
\big((k,\Delta)^2,p\big)^2 = \frac{(k\Delta)^2}{6}\left\{ (kp)^2\Delta_x^2 + (\Delta p)^2 k_x^2 + 4(kp)(\Delta p)\,k_x\Delta_x \right\},
\]
oder, da die beiden ersten Glieder rechts die nämliche Covariante darstellen, wie man unmittelbar erkennt, wenn man $\Delta$ durch $(k_1 k_2)^2 k_{1x}^2 k_{2x}^2$ ersetzt, und da der dritte Term der Klammer nach dem Productsatz sich durch $2(\Delta p)^2 k_x^2 + 2(kp)^2\Delta_x^2 - 2(k\Delta)^2 p$ ersetzen lässt, so wird
\[
\big((k,\Delta)^2,p\big)^2 = \frac{(k\Delta)^2}{6}\left\{ 6(kp)^2\Delta_x^2 - 2(k\Delta)^2 p \right\},
\]
oder
\[
\frac{B}{6}(k,p)^2 = (k\Delta)^2(kp)^2\Delta_x^2 - \frac{1}{3}\,C\,p,
\]
und folglich
\[
2(kp)^2(k\Delta)^2\Delta_x^2 = \frac{1}{3}\,B(k,p)^2 + \frac{2}{3}\,Cp.
\]
Durch Substitution dieses Werthes in (5) erhält man:
\weq{6}{
n - \frac{1}{3}\,Bl = \frac{1}{3}\,B(k,p)^2 + \frac{1}{3}(\Delta,p)^2(A-2P) + \frac{2}{3}\,Cp.
}
Hierin ist $(k,p)^2$ wegen Gleichung (V) bekannt. $(\Delta,p)^2$ aber geht aus der nämlichen Relation (V) durch zweimalige Ueberschiebung über $k$ hervor; dies liefert nämlich:
%CONT>>


% ===== pdf idx 314, printed 300.0 =====
% !PRINTED 300
\weberpage{300}
%<<CONT
\[
(k,l)^2 = m = 2\big(k,(k,p)^2\big)^2 + \frac{1}{3}(k,p)^2(A-2P),
\]
oder:
\[
m = 2(kk_1)^2(k_1p)^2 k_x^2 + \frac{1}{3}(k,p)^2(A-2P).
\]
Der erste Term rechts ist ein Glied der Ueberschiebung $(\Delta,p)^2$; denn man hat:
\[
\begin{aligned}
(\Delta,p)^2 &= (kk_1)^2(k_x^2 k_{1x}^2,\,p)^2 = \frac{(kk_1)^2}{3}\left\{ (k_1p)^2 k_x^2 + 2(kp)(k_1p)\,k_{1x} k_x \right\} \\
&= (kk_1)^2(k_1p)^2 k_x^2 - \frac{1}{3}\,Bp,
\end{aligned}
\]
oder:
\[
2(kk_1)^2(k_1p)^2 k_x^2 = 2(\Delta,p)^2 + \frac{2B}{3}\,p.
\]
Demnach wird:
\weq{7}{
m = 2(\Delta,p)^2 + \frac{2}{3}\,Bp + \frac{A-2P}{3}(k,p)^2.
}
Wir haben also, wenn wir der Einfachheit halber
\[
A - 2P = 6\varrho
\]
setzen, folgende drei Relationen:
\[
\begin{aligned}
\frac{l}{2} &= (k,p)^2 + p\cdot\varrho \quad\text{[vgl.\ (V)]} \\
\frac{m}{2} &= (k,p)^2\cdot\varrho + (\Delta,p)^2 + p\cdot\frac{B}{3} \quad\text{[vgl.\ (7)]} \\
\left(\frac{n}{2}-\frac{B}{6}l\right) &= (k,p)^2\cdot\frac{B}{6} + (\Delta,p)^2\cdot\varrho + p\cdot\frac{C}{3} \quad\text{[vgl.\ (6)].}
\end{aligned}
\]
276. Multiplicirt man die erste Gleichung mit $\varrho^2-\frac{B}{6}$, die zweite mit $-\varrho$, die dritte mit $1$ und addirt, so kommt:
\[
\frac{l}{2}\left(\varrho^2-\frac{B}{6}\right) - \frac{m}{2}\,\varrho + \left(\frac{n}{2}-\frac{B}{6}l\right) = p\left\{ \varrho^3 - \frac{B}{6}\varrho - \frac{B}{3}\varrho + \frac{C}{3} \right\},
\]
oder:
\[
\left(\frac{\varrho^2}{2}-\frac{B}{4}\right) l - \frac{\varrho}{2}\,m + \frac{n}{2} = p\left\{ \varrho^3 - \frac{B}{2}\varrho + \frac{C}{3} \right\}.
\]
Ueberschiebt man endlich diese Relation zweimal über $l$ und ersetzt $A_{ll}$, $A_{ml}$, $A_{nl}$ durch ihre Werthe $2C+\frac{AB}{3}$, resp. $\frac{2}{3}(B^2+AC)$ und $D$, so kommt als zweite Relation für $(lp)^2$:
\[
\left(\frac{\varrho^2}{2}-\frac{B}{4}\right)\left(2C+\frac{AB}{3}\right) - \frac{\varrho}{3}(B^2+AC) + \frac{D}{2} = (lp)^2\left\{ \varrho^3 - \frac{B}{2}\varrho + \frac{C}{3} \right\}.
\]
Ersetzen wir hierin gemäss Gleichung (VI) $(lp)^2$ durch
\[
\frac{2}{3}\,B + \frac{P}{9}(A-2P) = -2\varrho^2 + \frac{A}{3}\varrho + \frac{2}{3}B,
\]
so kommt nach gehöriger Vereinfachung:
%CONT>>


% ===== pdf idx 315, printed 301.0 =====
% !PRINTED 301
\weberpage{301}

\[
2\varrho^5 - \frac{A}{3}\varrho^4 - \frac{5B}{3}\varrho^3 + \frac{AB+5C}{3}\varrho^2 - \frac{4AC}{9}\varrho - \left(\frac{AB^2}{12} + \frac{13}{18}BC - \frac{D}{2}\right) = 0.
\]

Die Discriminante $P$ der quadratischen Form $p$ ist also von einer Gleichung fünften Grades abhängig, deren Coefficienten die Invarianten der ersten Elementarcovariante sind. In ähnlicher Weise würde man finden, dass für zwei Formen dritten Grades $f$ und $\varphi$ die zweite Elementarcovariante, welche in diesem Falle bereits Invariante ist, durch eine quadratische Gleichung sich aus den Coefficienten der ersten Elementarcovariante berechnen lässt.

277. \textit{Schlussbemerkungen.} Ich will die letzten Betrachtungen mit dem Hinweis auf deren Verallgemeinerung schliessen. Zunächst hatte Brill (Math. Annalen Bd. XX) die Bedeutung dieser ersten Elementarcombinante für alle übrigen Combinanten bemerkt, auch für den Fall, dass das zu Grunde gelegte simultane System aus mehr als zwei Formen $f, \varphi$ besteht. Unter anderm zeigte er, dass sie eine allgemeine Form ihres Grades ist und dass ihre Discriminante in zwei rationale Factoren zerfällt. Stephanos, der in den Comptes Rendus 1882 gleichzeitig auf diese Beziehungen aufmerksam machte, hat alsdann in seinen Abhandlungen: \glqq Mémoire sur les faisceaux de formes binaires ayant une même Jacobienne\grqq{} vom Jahre 1883, und \glqq Mémoire sur la Théorie des Formes binaires et sur l'élimination\grqq{} vom Jahre 1884 dieselben noch eingehender untersucht. Er sowohl als Franz Meyer (Apolarität und rationale Curven, Seite 391) zeigen, dass die in unserm speciellen Falle aufgestellte Endgleichung fünften Grades für die Coefficienten der zweiten Elementarcovariante im allgemeinen Falle, wenn $f$ und $\varphi$ vom $n^{\text{ten}}$ Grade sind, den Grad $N = \dfrac{(2(n-1))!}{(n-1)!\,n!}$ besitzt. Bezüglich weiterer interessanter Resultate müssen wir auf die citirten Quellen verweisen. Hier sei nur noch auf die Umkehrung der im Vorigen behandelten Aufgabe aufmerksam gemacht. Man kann sich nämlich fragen: welche Formen $f$ und $\varphi$ gehören zu einer gegebenen ersten Elementarcovariante $F$ von $f$ und $\varphi$? Für unser Beispiel lässt sich die Frage einfach beantworten. Es war $f = r_x^4$, $\varphi = s_x^4$, also:

\[
\psi = r_x^4 s_y^4 - s_x^4 r_y^4 = \mathrm{Const.}\, F_{y^3}(yx) + \mathrm{Const.}\, \Phi_y(yx).
\]

Da $F$ gegeben ist, $\Phi$ durch eine Gleichung fünften Grades aus $F$ sich berechnen lässt, $y$ aber willkürlich ist, so erkennt man, dass zu einer gegebenen Form $F$ stets fünf Schaaren oder lineare Systeme

% ?? uncertain: printed capital glyph appears stylized (script-like), transcribed as \Psi
\[
\Psi = \mu\, r_x^4 + \lambda\, s_x^4
\]

gehören. Einer solchen Schaar gehören insbesondere zwei specielle Formen an, für welche die Invariante $i = 0$ ist, also die beiden
%CONT>>


% ===== pdf idx 316, printed 302.0 =====
% !PRINTED 302
\weberpage{302}
%<<CONT
Tetraeder der Schaar, welche durch die aus $i = 0$ sich ergebende quadratische Gleichung für $\dfrac{\lambda}{\mu}$ bestimmt sind. In einer Mittheilung an die Société mathématique de France (1880) hat Stephanos gezeigt, dass, wenn $f$ und $\varphi$ Formen $n^{\text{ten}}$ Grades sind, die Zahl der linearen Systeme, welche zu einer gegebenen Functionaldeterminante zweier Formen $f$ und $\varphi$ gehören, gleich der oben angeführten Zahl $N = \dfrac{(2(n-1))!}{(n-1)!\,n!}$ ist.

\begin{center}
\S.\,29.\\
\textbf{Typische Darstellung der Form sechsten Grades.}
\end{center}

278. \textit{Allgemeiner Fall.} In den Systemen von Formen gerader Ordnung treten im allgemeinen Falle lineare Covarianten nicht auf. Dagegen besitzen dieselben im Allgemeinen, wenn $n > 4$, stets mindestens zwei quadratische Covarianten $\varphi$ und $\psi$, deren Resultante $R_{\varphi,\psi}$ nicht verschwindet. (Vergl. Clebsch, \glqq Binäre Formen\grqq, § 102, Seite 140.) Diese zwei quadratischen Formen können nun unter anderm als Grundlage für die typische Darstellung von Formen gerader Ordnung benutzt werden, wie Clebsch und Gordan, Annali di Mat., Ser. II., t. 1, gezeigt haben. Denn zwischen vier quadratischen Formen $a_x^2, l_x^2, m_x^2, n_x^2$ besteht stets die Relation (vergl. Nr. 135, I, worin $y = a$ zu setzen ist)

\[
a_x^2 \cdot R_{lmn} = l(a\lambda)^2 + m(a\mu)^2 + n(a\nu)^2,
\]

wobei $\lambda, \mu, \nu$ die Functionaldeterminanten $(mn)\,m_xn_x$, respective $(nl)\,n_xl_x$, $(lm)\,l_xm_x$ repräsentiren. Und diese Relation findet hier dieselbe Verwendung, wie bei den Formen ungeraden Grades die Identität

\[
a_x(\alpha\beta) = \alpha_x(a\beta) - \beta_x(a\alpha).
\]

Wir erheben sie auf die $\dfrac{n}{2}^{\text{te}}$ Potenz und erhalten die Beziehung

\weq{I}{R_{lmn}^{\frac{n}{2}} \cdot f = \left\{ l(a\lambda)^2 + m(a\mu)^2 + n(a\nu)^2 \right\}^{\frac{n}{2}},}

aus welcher die typische Darstellung von $f$ unmittelbar hervorgeht. Es treten dabei allerdings an Stelle der zwei alten Variabeln $x_1x_2$ drei neue $l, m, n$ auf; doch sind diese nicht von einander unabhängig. Denn zwischen ihnen besteht die Relation [vergl. Nr. 134, 3, $(\gamma)$]

\[
0 = \begin{vmatrix}
A_{ll}, & A_{lm}, & A_{ln}, & l \\
A_{lm}, & A_{mm}, & A_{mn}, & m \\
A_{ln}, & A_{mn}, & A_{nn}, & n \\
l, & m, & n, & 0
\end{vmatrix}
\]

und somit ist die typische Darstellung (I) von $f$ in der That eine binäre.


% ===== pdf idx 317, printed 303.0 =====
% !PRINTED 303
\weberpage{303}

279. Was nun die Form $f = a_x^6$ betrifft, so könnten wir deren typische Darstellung unmittelbar auf Gleichung (I) basiren. Indessen bietet in diesem Falle die in Nr. 135, (VII) gegebene Relation ein noch bequemeres Mittel, zum nämlichen Resultate zu gelangen, so lange die Invariante $R$ von $0$ verschieden ist.

Ersetzen wir in derselben nämlich $f, \varphi, \psi$ durch die drei quadratischen Formen $l, m, n$ des Systemes von $f = a_x^6$, so geht dieselbe über in:

% ?? uncertain: a small stray mark (ink speck or accent) appears above A_{mm} in this row in the print
\[
\begin{vmatrix}
A_{ll}, & A_{lm}, & A_{ln}, & l_x^2, & l_y^2 \\
A_{lm}, & A_{mm}, & A_{mn}, & m_x^2, & m_y^2 \\
A_{ln}, & A_{mn}, & A_{nn}, & n_x^2, & n_y^2 \\
l_x^2, & m_x^2, & n_x^2, & 0, & (xy)^2 \\
l_y^2, & m_y^2, & n_y^2, & (xy)^2, & 0
\end{vmatrix} = 0.
\]

Entwickeln wir diese Determinante mit Hilfe des Additionstheorems für Determinanten (vergl. Bd. I Nr. 35 u. 36) nach Potenzen von $(xy)$, so erhalten wir:

\[
0 = -2R^2\cdot(yx)^4 - 2
\begin{vmatrix}
A_{ll}, & A_{lm}, & A_{ln}, & l_x^2 \\
A_{lm}, & A_{mm}, & A_{mn}, & m_x^2 \\
A_{ln}, & A_{mn}, & A_{nn}, & n_x^2 \\
l_y^2, & m_y^2, & n_y^2, & 0
\end{vmatrix}(yx)^2 +
\begin{vmatrix}
A_{ll}, & A_{lm}, & A_{ln}, & l_x^2, & l_y^2 \\
A_{lm}, & A_{mm}, & A_{mn}, & m_x^2, & n_y^2 \\
A_{ln}, & A_{mn}, & A_{nm}, & n_x^2, & n_y^2 \\
l_x^2, & m_x^2, & n_x^2, & 0, & 0 \\
l_y^2, & m_y^2, & n_y^2, & 0, & 0
\end{vmatrix}.
\]
% [sic] third row, third entry printed "A_{nm}" where the pattern elsewhere in this row/column uses "A_{mn}"

Hierin besitzt der Coefficient von $2(yx)^2$ gemäss der Relation (VI) Nr. 135 den Werth $-2R^2(yx)^2$, wie sich aus derselben wiederum nach dem Additionsgesetze für Determinanten ergiebt. Demnach geht die letzte Gleichung, wenn wir überdies $y$ durch $a$ ersetzen und mit $a_x^2$ multipliciren, über in:

\[
-2R^2\cdot f = a_x^2
\begin{vmatrix}
A_{ll}, & A_{lm}, & A_{ln}, & l, & (al)^2 \\
A_{lm}, & A_{mm}, & A_{mn}, & m, & (am)^2 \\
A_{ln}, & A_{mn}, & A_{nn}, & n, & (an)^2 \\
l, & m, & n, & 0, & 0 \\
(al)^2, & (am)^2, & (an)^2, & 0, & 0
\end{vmatrix}.
\]

280. Entwickelt man diese Determinante nach den Minoren der letzten zwei Zeilen und Colonnen (vergl. insbesondere Bd. I Seite 72, Beispiel 2), so erhält man:

\weq{I}{\begin{aligned}
-2R^2\cdot f ={}& \sum A_{nn}\left\{ l^2(f,m^2)^4 - 2lm(f,ml)^4 + m^2(f,l^2)^4 \right\} \\
&{}- 2\sum A_{mn}\left\{ l^2(f,mn)^4 - lm(f,ln)^4 + mn(f,l^2)^4 - nl(f,ml)^4 \right\},
\end{aligned}}
% [sic] this display is printed with the margin label "(I)", reusing the same numeral already used for the equation on p. 302
%CONT>>


% ===== pdf idx 318, printed 304.0 =====
% !PRINTED 304
\weberpage{304}
%<<CONT
wobei jede Summe sich über drei Glieder erstreckt, deren zweites und drittes aus den angeschriebenen durch cyclische Vertauschung von $l\, m\, n$ hervorgehen.

Die in der Entwicklung rechts auftretenden quadratischen Formen

\[
(f, l^2)^4,\quad (f, m^2)^4,\quad (f, n^2)^4,\quad (f, lm)^4,\quad (f, ln)^4,\quad (f, mn)^4
\]

sind aber, wie wir bereits gezeigt haben, lineare Functionen der Covarianten $l, m, n$. Diese Entwicklung rechts ist also eine Function dritten Grades in den Variabeln $l, m, n$, und die darin auftretenden Coefficienten sind Invarianten. Die typische Darstellung von $f$ ist also in der That durch die Gleichung (I) geleistet. Betrachtet man in ihr die Covarianten $l, m, n$ als drei Coordinaten $x_1, x_2, x_3$, so stellt die rechte Seite dieser Gleichung eine Curve dritter Ordnung dar. Da aber nach der Schlussbemerkung in Nr. 278 dieselben drei Coordinaten den Kegelschnitt

% ?? uncertain: third row, first two entries print as lowercase "l_{ln}, l_{mn}" where the pattern elsewhere suggests "A_{ln}, A_{mn}"; reproduced as printed
\weq{II}{
\begin{vmatrix}
A_{ll} & A_{lm} & A_{ln} & x_1 \\
A_{lm} & A_{mm} & A_{mn} & x_2 \\
l_{ln} & l_{mn} & A_{nn} & x_3 \\
x_1 & x_2 & x_3 & 0
\end{vmatrix} = 0
}

befriedigen, so erkennt man: Die Wurzeln der Gleichung sechsten Grades $f = 0$ werden durch die typische Darstellung von $f$ als Schnittpunkte einer Curve dritter Ordnung und eines Kegelschnittes erhalten.

281. \textit{Der specielle Fall $R = 0$.} Wir hatten bisher vorausgesetzt, dass $R$ von null verschieden sei. Nehmen wir nun an, dass $R$ identisch verschwinde, so dürfen, wenn die typische Darstellung gleichwohl möglich sein soll, die Covarianten $l, m, n$ keinen gemeinsamen Factor haben. Zu dem Zwecke genügt es, dass $A_{ll}A_{mm} - A_{ml}^2$ von null verschieden sei.

Wir benutzen alsdann behufs typischer Darstellung die in der Theorie der quadratischen Formen gegebene Identität [vgl. Nr. 135, (V)]

\[
2v_x^2 \cdot v_y^2 =
\begin{vmatrix}
A_{ll} & A_{lm} & l \\
A_{lm} & A_{mm} & m \\
l_y^2 & m_y^2 & (xy)^2
\end{vmatrix},
\]

welche, nach $(xy)^2$ entwickelt, die Form annimmt:

\[
(xy)^2 A_{vv} = 2v_x^2v_y^2 -
\begin{vmatrix}
A_{ll} & l_{lm} & l \\
A_{lm} & A_{mm} & m \\
l_y^2 & m_y^2 & 0
\end{vmatrix}.
\]
% [sic] second row/first row printed "l_{lm}" where the pattern elsewhere suggests "A_{lm}"

Erheben wir dieselbe auf die dritte Potenz und ersetzen alsdann $y$ durch $a$, so kommt:
%CONT>>


% ===== pdf idx 319, printed 305.0 =====
% !PRINTED 305
\weberpage{305}
\[
\begin{aligned}
A_{\nu\nu}^3\cdot f&=\prod_{i=1}^{i=3}\left\{2\,\nu\,(\nu_ia)^2-\begin{vmatrix}A_{ll}, & A_{lm}, & l\\ A_{lm}, & A_{mm}, & m\\ (l_ia)^2, & (m_ia)^2, & 0\end{vmatrix}\right\}\\
&=8\,\nu^3\cdot(f,\,\nu^3)^6-4\,\nu^2\sum{}'(\nu_1a)^2(\nu_2a)^2\begin{vmatrix}A_{ll}, & A_{lm}, & l\\ A_{lm}, & A_{mm}, & m\\ (l_3a)^2, & (m_3a)^2, & 0\end{vmatrix}\\
&\quad+2\,\nu\sum{}'(\nu_1a)^2\begin{vmatrix}A_{ll}, & A_{lm}, & l\\ A_{lm}, & A_{mm}, & m\\ (l_2a)^2, & (m_2a)^2, & 0\end{vmatrix}\cdot\begin{vmatrix}A_{ll}, & A_{lm}, & l\\ A_{lm}, & A_{mm}, & m\\ (l_3a)^2, & (m_3a)^2, & 0\end{vmatrix}\\
&\quad-\prod_{i=1}^{i=3}\begin{vmatrix}A_{ll}, & A_{lm}, & l\\ A_{lm}, & A_{mm}, & m\\ (l_ia)^2, & (l_ia)^2, & 0\end{vmatrix},% [sic] printed (l_ia)^2 twice in this determinant, not (l_ia)^2,(m_ia)^2
\end{aligned}
\]
wobei sich die beiden Summen $\sum{}'$ über je drei Glieder erstrecken.

In der Entwicklung rechts verschwindet das erste und dritte Glied. Denn die Invarianten
\[
(f,\,\nu^3)^6,\quad (f,\,\nu l^2)^6,\quad (f,\,\nu m^2)^6,\quad (f,\,\nu lm)^6,
\]
welche in diesen zwei Gliedern als Coefficienten auftreten, sind schiefe Invarianten von den Graden 25, bezw. 15, 19, 17 in den Coefficienten von $f$, sie müssen daher, da nur die einzige schiefe Fundamentalinvariante $R$ im Systeme existirt, diese zum Factor haben, und sonach gleichzeitig mit $R$ verschwinden.

Es reducirt sich also die typische Darstellung auf:
\weq{III}{A_{\nu\nu}^3 f\equiv -\prod\begin{vmatrix}A_{ll}, & A_{lm}, & l\\ A_{lm}, & A_{mm}, & m\\ (l_ia)^2, & (l_ia)^2, & 0\end{vmatrix}% [sic] printed (l_ia)^2 twice, as above
-4\,\nu^2\sum{}'(\nu_1a)^2(\nu_2a)^2\begin{vmatrix}A_{ll}, & A_{lm}, & l\\ A_{lm}, & A_{mm}, & m\\ (l_3a)^2, & (m_3a)^2, & 0\end{vmatrix}.}

Da $\nu$ als Functionaldeterminante von $l$ und $m$ der Relation genügt
\weq{IV}{-2\,\nu^2=\begin{vmatrix}A_{ll}, & A_{lm}, & l\\ A_{lm}, & A_{mm}, & m\\ l, & m, & 0\end{vmatrix},}
so erhält man nach Substitution dieses Werthes von $\nu^2$ in die drei letzten Glieder der Gleichung (III) auf der rechten Seite einen in den Formen $l$ und $m$ homogenen cubischen Ausdruck von der Gestalt:
\[
A_{\nu\nu}^3 f=C_1l^3+C_2l^2m+C_3lm^2+C_4m^3.
\]


% ===== pdf idx 320, printed 306.0 =====
% !PRINTED 306
\weberpage{306}
Die Coefficienten $C_i$ dieser Darstellung sind Functionen der Invarianten $A_{ll}$, $A_{lm}$, $A_{mm}$ und ausserdem der Invarianten
\[
(f,\,l^3)^6,\quad (f,\,m^3)^6,\quad (f,\,l^2m)^6,\quad (f,\,lm^2)^6;
\]
denn die zwei Invarianten $(f,\,\nu^2l)^6$, $(f,\,\nu^2m)^6$ lassen sich sofort wegen (IV) auf dieselben reduciren.

Von diesen wurde die erste bereits früher berechnet (Nr. 264):
\[
(f,\,l^3)^6=2A_{ln}-\frac{2}{3}BA_{ll}+\frac{1}{3}A_{lm}.% [sic] index \textquotedbl{}ln\textquotedbl{} as printed
\]
Die drei andern ergeben sich direct aus den beiden ersten Relationen in (3) Nr. 266, indem man dieselben zweimal über $m$, resp. $l$ schiebt. Die Coefficienten der typischen Darstellung
\weq{V}{A_{\nu\nu}^3\cdot f=C_1l^3+C_2l^2m+C_3lm^2+C_4m^3}
sind somit vollständig bekannte Grössen.

Weitere specielle Fälle sind in Bezug auf typische Darstellung noch wenig untersucht worden. In den \glqq Binären Formen\grqq{} Clebsch's findet sich noch neben dem eben besprochenen Fall die typische Darstellung von $f$, wenn $A=0$, $B=0$, $C=0$, $D$ und $R$ von Null verschieden. Es ergiebt sich das Resultat:
\[
D^2f=2n^3+\frac{3D}{2}l^2m,
\]
wobei $n$ und $l$ reine Quadrate sind. (Vergl. a. a. O. Seite 467.)

\begin{center}
\S\ 30. \textbf{Auflösung specieller Gleichungen sechsten Grades.}
\end{center}

282. \textit{Ueberblick.} Von den Gleichungen sechsten Grades, welche vermöge gewisser Invarianten- oder Covariantenbedingungen den Charakter von speciellen Gleichungen tragen, sind einmal jene hervorzuheben, deren Invariantenrelationen ihre Reduction auf Gleichungen niedrigerern Grades involviren, dann aber auch jene, durch deren Invariantenbeziehungen zwar nicht ihr Grad reducirt wird, welche aber durch andere bemerkenswerthe Eigenschaften unser Interesse erregen, insofern sie mit Modulargleichungen oder Multiplicatorgleichungen identificirt werden können. Wir wollen einige wichtige Fälle dieser speciellen Gleichungen sechsten Grades im Folgenden eingehender untersuchen, insbesondere um auf die Art und Weise hinzudeuten, wie die hier einschlägigen Fragen mit unsern Mitteln behandelt werden können. Der uns zunächst interessirende Fall ist der Fall, in welchem $R=0$, und die Resultante $R_{lm}=-2A_{\nu\nu}$ von $l$ und $m$ von Null verschieden ist. Hieran schliessen sich die Untersuchungen über den Fall $R=0$, $A_{\nu\nu}=0$, die zu drei Unterabtheilungen von Bedingungen Veranlassung geben, jenachdem das Verschwinden von $A_{\nu\nu}$ eine Folge davon, dass $l$ und
%CONT>>


% ===== pdf idx 321, printed 307.0 =====
% !PRINTED 307
\weberpage{307}
%<<CONT
$m$ einen linearen Factor gemeinschaftlich besitzen, oder dass $l$ proportional $m$ ist, oder endlich, dass $l$ oder $m$ selbst identisch null wird. Im letzten Falle wird sich zeigen, dass, wenn eine der drei quadratischen Formen verschwindet, nothwendig auch die zwei andern null werden, und die Form $(f,f)^4$ entweder das Quadrat oder die vierte Potenz einer Form zweiten resp. ersten Grades ist, so dass also nur noch diese beiden Möglichkeiten zu besprechen sind, um den Fall $R=0$, $A_{\nu\nu}=0$ völlig zu erschöpfen.

Zum Schlusse weisen wir noch auf die beiden Fälle: $A=0$, $C=0$ und $B=\text{Const.}\times A^2$, $C=\text{Const.}\times A^3$ hin, unter welchen Bedingungen sich die allgemeine Gleichung sechsten Grades durch lineare Transformation auf Gleichungen mit einem einzigen Parameter reduciren lässt.

283. \textit{Erster Fall: $R=0$, $A_{\nu\nu}\gtrless 0$.} Die typische Darstellung der Form $f=a_x^6$ für $R=0$ am Schlusse des letzten Paragraphen lehrt:

\glqq Verschwindet die Invariante $R$, so ist die Gleichung sechsten Grades $f=0$ algebraisch lösbar, und ihre Wurzeln sind drei Elementenpaare einer Involution, deren Doppelpunkte durch die linearen Factoren von $\nu=0$ bestimmt sind.\grqq{}

Denn durch die Substitution
\[
y=\frac{l}{m}
\]
geht die unter Nr. 281 gegebene Gleichung (V) über in
\[
C_1y^3+C_2y^2+C_3y+C_4=0.
\]
Die Wurzeln dieser Gleichung seien $y_1$, $y_2$, $y_3$; die Wurzeln von $f=0$ werden alsdann durch die drei quadratischen Gleichungen:
\weq{1}{l-y_1m=0,\quad l-y_2m=0,\quad l-y_3m=0}
geliefert. Da nach Voraussetzung $A_{\nu\nu}\gtrless 0$, also $l$ und $m$ keinen gemeinschaftlichen Factor haben, so kann man diese beiden quadratischen Covarianten auf die canonische Form:
\[
\begin{aligned}
l&=\alpha_1\xi^2+\beta_1\eta^2\\
m&=\alpha_2\xi^2+\beta_2\eta^2
\end{aligned}
\]
durch eine gemeinschaftliche lineare Transformation bringen (vergl. Nr. 130), wobei $\xi$ und $\eta$ die linearen Factoren von $\nu=0$ sind. Alsdann gehen die drei quadratischen Formen (1) über in:
\[
c_1\xi^2-c_2\eta^2=0,\quad c_1'\xi^2-c_2'\eta^2=0,\quad c_1''\xi^2-c_2''\eta^2=0
\]
und diese drei Gleichungen repräsentiren drei Punktepaare, welche zu den Punkten $\xi=0$, $\eta=0$ harmonisch liegen, was zu beweisen war.


% ===== pdf idx 322, printed 308.0 =====
% !PRINTED 308
\weberpage{308}
284. \textit{Zweiter Fall; erste Unterabtheilung: $R=0$, $A_{\nu\nu}=0$.} Wir nehmen an, das Verschwinden von $A_{\nu\nu}=A_{ll}A_{mm}-A_{lm}^2$ sei dadurch bedingt, dass $l$ und $m$ einen gemeinschaftlichen linearen Factor $\alpha$ haben. Da aber die Bedingung $R=0$ eine lineare Relation:
\[
2\,\nu\cdot R=\begin{vmatrix}A_{ll}, & A_{lm}, & A_{ln}\\ A_{ml}, & A_{mm}, & A_{mn}\\ l, & m, & n\end{vmatrix}=c_1l+c_2m+c_2n=0% [sic] third coefficient printed c_2 twice, not c_3
\quad\text{(vgl. Nr. 135, IV$_{\text{b}}$)}
\]
zwischen den drei Covarianten $l$, $m$, $n$ involvirt, so muss auch $n$ mit $l$ und $m$ denselben linearen Factor $\alpha$ gemeinschaftlich haben.

Die Bedingungen für unsern Fall lassen sich demnach auch in der Form schreiben
\weq{1}{l=\alpha\cdot r,\quad m=\alpha\cdot s,\quad n=\alpha\cdot t,}
wobei die linearen Factoren $r$, $s$, $t$ von einander verschieden sein sollen. Aus diesen drei Bedingungen (1) gehen direct die Relationen
\weq{2}{(l\alpha)^2=0,\quad (m\alpha)^2=0,\quad (n\alpha)^2=0}
hervor, während nach der Theorie der quadratischen Formen
\[
\alpha^2=(l,\,m)=\nu\quad\text{ist.}
\]
Nun ist aber die Covariante $l$ nichts anderes als $(f,k)^4$ und es liegt die Frage nahe, ob eine dieser beiden Formen $f$ und $k$ den rationalen linearen Factor $\alpha_x$ besitzt. Um dies zu untersuchen, hat man nur die Ueberschiebungen von $k$, resp. $f$ über $\alpha$ zu bilden und nachzusehen, welche derselben identisch verschwindet. Man findet alsdann:
\weq{3}{(k\alpha)^3k_x=0}
und
\weq{4}{(f\alpha)^4f_x^2=0.}

Die erste Gleichung lehrt: $k$ besitzt den Factor $\alpha_x$ quadratisch, die zweite: $f$ enthält ihn cubisch*).

Denn multipliciren wir, um die Bedingungen (2) zum Beweise einführen zu können, $(k\alpha)^3k_x$ mit $(rs)\gtrless 0$, so erhalten wir unter Anwendung des Identitätssatzes:

\noindent\rule{4cm}{0.4pt}

*) Dass die Deutung dieser beiden Gleichungen richtig ist, kann man sofort erkennen, wenn man eine Form $f=a_x^n$ mit Hilfe zweier linearer Formen $\alpha_x$ und $\beta_x$ vermöge der Identität
\[
\alpha_x(\alpha\beta)=(\alpha\beta)\alpha_x-(a\alpha)\beta_x
\]
darzustellen versucht. Erhebt man dieselbe auf die $n^{\text{te}}$ Potenz und entwickelt rechts nach dem binomischen Lehrsatz, so ergiebt sich unmittelbar unter der Voraussetzung $(f,\alpha)^{n-k}=0$, dass $f$ den Factor $\alpha^{k+1}$ besitzt.
%CONT>>


% ===== pdf idx 323, printed 309.0 =====
% !PRINTED 309
\weberpage{309}

\[
\begin{aligned}
(k\alpha)^3k_x(rs) &= (k\alpha)^3\{(ks)r_x-(kr)s_x\} = (k,\alpha^3s)^4-(k,\alpha^3r)^4\\
&= ((k,\alpha s)^2,\alpha^2)^2-((k,\alpha r)^2,\alpha^2)^2=((k,m)^2,\alpha^2)^2-((k,l)^2,\alpha^2)^2\\
&= (n\alpha)^2-(m\alpha)^2=0.
\end{aligned}
\]
Da aber $(rs)\gtrless 0$ ist, so muss $(k\alpha)^3k_x=0$ sein. Es enthält demnach in der That $k$ den Factor $\alpha_x$ zweifach, und folglich nach der Theorie der biquadratischen Formen auch die Hesse'sche Form $\varDelta$ von $k$. Weil sich nun aber alle zweiten Ueberschiebungen von $f$ über $l$, $m$, $n$ als lineare Functionen von $k$, $\varDelta$, $l^2$, $lm$ darstellen lassen, so muss auch jede dieser Formen $(f,l)^2$, $(f,m)^2$, $(f,n)^2$ den linearen Factor $\alpha$ quadratisch enthalten, d.\,h. es muss sein:
\weq{5}{(f,l)^2=a_x^4(a\alpha)(ar)=\alpha_x^2\cdot\varphi}
\weq{6}{(f,m)^2=a_x^4(a\alpha)(as)=\alpha_x^2\cdot\psi,}
wo $\varphi$ und $\psi$ quadratische Formen sind, die für $x=\alpha$ nicht verschwinden. Ueberschiebt man diese beiden Formen dreimal über $\alpha^3$, so müssen diese Ueberschiebungen wegen $(\alpha,\alpha)=0$ verschwinden, d.\,h. man hat:
\[
\begin{aligned}
(\alpha_x^2\cdot\varphi,\alpha_x^3)^3&=a_x(a\alpha)^4(ar)=0\\
(\alpha_x^2\cdot\psi,\alpha_x^3)^3&=a_x(a\alpha)^4(as)=0.
\end{aligned}
\]

Multiplicirt man die erste dieser Gleichungen mit $-s_x$, die zweite mit $r_x$ und addirt, so kommt
\[
a_x(a\alpha)^4\big((as)r_x-(ar)s_x\big)=a_x^2(a\alpha)^4(rs)=0,
\]
oder:
\[
a_x^2(a\alpha)^4=0,
\]
d.\,h. $f$ besitzt den Factor $\alpha_x$ cubisch. Man hat also den Satz:
\glqq Haben die drei quadratischen Formen $l$, $m$, $n$ einen gemeinsamen Factor $\alpha_x$, so besitzt $f$ denselben dreifach.\grqq{}

285. \textit{Umkehrung.} Diese Bedingungen sind nöthig und hinreichend.
Denn besitzt $f$ den Factor $\alpha_x$ dreifach, so ist zunächst
\[
a_x^2(a\alpha)^4=0.
\]
Multipliciren wir anderntheils $(k\alpha)^3k_x$ mit $\alpha_x^4$, so erhalten wir:
\[
(k\alpha)^3k_x\cdot\alpha_x^4=\frac{(ab)^4}{2}\{(a\alpha)^2(b\alpha)b_x+(b\alpha)^2(a\alpha)a_x\}\cdot\alpha_x^4,
\]
oder, weil beide Terme rechts gleich sind:
\[
\begin{aligned}
(k\alpha)^3k_x\cdot\alpha_x^4&=(ab)^4(a\alpha)^2(b\alpha)b_x\cdot\alpha_x^4\\
&=(a\alpha)^2(b\alpha)b_x\cdot\{b_x(a\alpha)-a_x(b\alpha)\}^4.
\end{aligned}
\]
Entwickelt man rechts nach dem binomischen Lehrsatz und führt die Multiplication aus, so besitzt jedes Glied den Factor
\[
(a\alpha)^4a_x^2=(b\alpha)^4b_x^2=0.
\]
%CONT>>


% ===== pdf idx 324, printed 310.0 =====
% !PRINTED 310
\weberpage{310}

%<<CONT
Es verschwindet demnach auch die linke Seite; folglich besitzt $k$ und mit ihm $\varDelta$ den Factor $\alpha_x$ quadratisch. Daraus ergiebt sich endlich unmittelbar, dass
\[
(l\alpha)^2=(m\alpha)^2=(n\alpha)^2=0.
\]
Denn es ist:
\[
(l,\alpha^2)^2=(ak)^4(a\alpha)^2,
\]
also
\[
\alpha^4(l\alpha)^2=(a\alpha)^2\{k_x(a\alpha)-a_x(k\alpha)\}^4=0,
\]
weil
\[
(k\alpha)^3k_x=0,\quad (a\alpha)^4a_x^2=0.
\]
Ferner ist:
\[
(m,\alpha^2)^2=(kl)^2(k\alpha)^2,
\]
also
\[
\alpha^2(m\alpha)^2=(k\alpha)^2\{k_x(l\alpha)-l_x(k\alpha)\}^2=0,
\]
weil
\[
(k\alpha)^3k_x=0,\quad (l\alpha)^2=0.
\]
Endlich ist:
\[
(n\alpha)^2=(km)^2(k\alpha)^2,
\]
also
\[
\alpha^2(n\alpha)^2=(k\alpha)^2\{k_x(m\alpha)-m_x(k\alpha)\}^2=0,
\]
weil
\[
(k\alpha)^3k_x=0,\quad (m\alpha)^2=0.
\]
Die drei quadratischen Formen $l$, $m$, $n$ besitzen also in der That einen gemeinsamen Factor $\alpha$ und es ist demnach $R=0$, $A_{\nu\nu}=0$.

286. \textit{Zweiter Fall; zweite Unterabtheilung.} Die Resultante von $l$ und $m$ verschwindet aber auch, wenn $l$ und $m$ zwei lineare Factoren gemeinschaftlich haben, d.\,h. wenn
\weq{1}{m=\varrho\cdot l,}
wo $\varrho$ eine Proportionalitätsconstante. In diesem Falle ist aber auch $n$ nur bis auf eine Constante von $l$ und $m$ verschieden. Denn
\[
n=(k,m)^2=\varrho(k,l)^2,
\]
und da $(k,l)^2=m$, so hat man:
\[
n=\varrho\cdot m=\varrho^2\cdot l.
\]
Unter diesen Voraussetzungen aber gilt zunächst der Satz:
\glqq Sind die drei quadratischen Covarianten $l$, $m$, $n$ einander proportional und von null verschieden, so kann $A_u$ nicht verschwinden.\grqq{}
Wir beweisen diesen Satz indirect. Wäre nämlich $A_u=0$, also $l$ und damit auch $n$ und $m$ ein volles Quadrat, so hätte man zunächst
\weq{2}{A_u=2C+\frac{AB}{3}=0\ (\text{vgl.\ Nr.\ 267}).}
%CONT>>


% ===== pdf idx 325, printed 311.0 =====
% !PRINTED 311
\weberpage{311}

%<<CONT
Anderntheils war (vgl.\ Nr.\ 266):
\weq{3}{
\left.\begin{aligned}
(f,n)^2&=-\frac{C}{3}k+\frac{B}{3}\varDelta+lm\\
(f,l)^2&=\phantom{-}\frac{A}{3}k+2\varDelta
\end{aligned}\right\}
}
Demnach ist
\[
\begin{aligned}
\Big(f,n-\frac{B}{6}l\Big)^2&=-\frac{C}{3}k+\frac{B}{3}\varDelta+lm-\frac{B}{6}\Big(\frac{A}{3}k+2\varDelta\Big)\\
&=lm-\frac{1}{6}\Big(2C+\frac{AB}{3}\Big)k,
\end{aligned}
\]
oder wegen (2)
\weq{4}{\Big(f,n-\frac{B}{6}l\Big)^2=lm=\varrho l^2.}
Weil aber $n=\varrho^2l$, so ist auch
\weq{5}{\Big(f,n-\frac{B}{6}l\Big)^2=\Big(\varrho^2-\frac{B}{6}\Big)(f,l)^2}
und also durch Comparation
\[
(f,l)^2=\frac{6\varrho}{6\varrho^2-B}l^2.
\]
Nun soll aber $l$ das Quadrat einer linearen Form $\alpha_x$ sein; demnach können wir die letzte Gleichung auch schreiben:
\[
(f,\alpha^2)^2=(a\alpha)^2a_x^4=\frac{6\varrho}{6\varrho^2-B}\alpha_x^4=\text{Const.}\times\alpha_x^4.
\]
Durch einmalige Ueberschiebung dieser Gleichung über $\alpha_x$ folgt:
\[
(a\alpha)^3a_x^3=0,
\]
d.\,h. $f$ besitzt den Factor $\alpha_x$ vierfach.
Dann ist aber auch $k$ das Biquadrat von $\alpha$. Denn in der Gleichung
\[
\alpha^4(k,\alpha)=(ab)^4\alpha^4\cdot a_xb_x^2(a\alpha)=(a\alpha)a_xb_x^2\{(a\alpha)b_x-(b\alpha)a_x\}^4
\]
verschwindet wegen $(a\alpha)^3a_x^3=(b\alpha)^3b_x^3=0$ die ganze rechte Seite und es ist sonach
\[
(k,\alpha)=0,\quad\text{also}\quad k=\text{Const.}\times\alpha^4.
\]
Bilden wir daher die Covariante $l$, also die vierte Ueberschiebung von $f$ über $k$, so ist diese vierte Ueberschiebung null, weil sowohl $f$ als $k$ den linearen Factor $\alpha_x$ vierfach besitzen. Es wäre demnach unter der Annahme $A_u=0$
\[
(f,k)^4=l=0,
\]
das ist aber gegen die Voraussetzung, und sonach muss $A_u\gtrless 0$ sein.

287. \textit{Schärfere Formulirung unserer Voraussetzungen.} Wir können nun unsere Bedingungen schärfer abgrenzen. Sind die drei quadratischen Formen einander proportional, so muss $l$, wegen $A_u\gtrless 0$, zwei von einander verschiedene lineare Factoren besitzen, d.\,h. es muss sein:
%CONT>>


% ===== pdf idx 326, printed 312.0 =====
% !PRINTED 312
\weberpage{312}

%<<CONT
\[
m=\varrho\cdot l=\varrho\cdot\alpha\cdot\beta\quad\text{und}\quad n=\varrho^2\cdot\alpha\cdot\beta,
\]
wo $\alpha$ und $\beta$ die beiden linearen Factoren sind. Es ist demnach
\weq{1}{
\left.\begin{aligned}
(l\alpha)^2&=0,\quad (m\alpha)^2=0,\quad (n\alpha)^2=0\\
(l\beta)^2&=0,\quad (m\beta)^2=0,\quad (n\beta)^2=0.
\end{aligned}\right\}
}
Da aber:
\[
m=\varrho\cdot\alpha\cdot\beta=(k,l)^2=\varrho(k,\alpha\beta)^2
\]
ist, so erhalten wir hieraus:
\weq{$2_{\mathrm{a}}$}{
\left.\begin{aligned}
((k,\alpha\beta)^2,\alpha^2)^2&=(k,\alpha^3\beta)^4=0\\
(k,\alpha\beta^3)^4&=0
\end{aligned}\right\}
}
Ebenso folgen aus diesen Gleichungen (1) für die Form $\varDelta$ die beiden Relationen:
\weq{$2_{\mathrm{b}}$}{
\left.\begin{aligned}
(\varDelta,\alpha^3\beta)&=0\\
(\varDelta,\alpha\beta^3)&=0,
\end{aligned}\right\}
}
wie man erkennt, wenn man diese Ueberschiebungen mit $(\alpha\beta)^2$ multiplicirt, und alsdann $(\alpha\beta)^2(kk_1)^2$ nach dem Identitätssatze entwickelt.
Nun sind die zweiten Ueberschiebungen von $f$ über $l$ und $m$ lineare Functionen von $k$, $\varDelta$ und $l^2$. Wenn wir daher $(f,l)^2$ und $(f,m)^2$ noch viermal über $\alpha^3\cdot\beta$ schieben, so verschwinden die so entstehenden Invarianten, d.\,h. es ist:
\weq{3}{
\left.\begin{aligned}
(f,\alpha^4\beta^2)^6&=0\\
(f,\alpha^2\beta^4)^6&=0
\end{aligned}\right\}
}
In die Relationen $(2_{\mathrm{a}})$ können wir noch die Symbole von $f$ einführen. Es ist:
\[
(k,\alpha^3\beta)^4=(ab)^4(a_x^2b_x^2,\alpha_x^3\beta_x)^4=(ab)^4(a\alpha)^2(b\alpha)(b\beta)=0.
\]
Multipliciren wir diesen Ausdruck mit $(\alpha\beta)^4$, und wenden auf $(ab)^4(\alpha\beta)^4$ den Identitätssatz an, so kommt:
\[
0=(ab)^4(\alpha\beta)^4(a\alpha)^2(b\alpha)(b\beta)\doteq(a\alpha)^2(b\alpha)(b\beta)\{(a\alpha)(b\beta)-(b\alpha)(b\beta)\}^4.
% ?? uncertain: second bracketed term on the right of the doteq -- printed factor may read (b\alpha)(a\beta) rather than (b\alpha)(b\beta); reproduced as best read from the scan
\]
Entwickelt man die rechte Seite nach dem binomischen Lehrsatz, so findet man, dass alle Glieder bis auf das erste wegen der Relationen (3) für sich allein verschwinden. Es muss sonach das erste Glied selbst null sein, d.\,h. wir haben statt der Relationen (1) die beiden folgenden:
\weq{4}{
\left.\begin{aligned}
(a\alpha)^6\cdot(a\alpha)(a\beta)^5&=0\\
(a\beta)^6\cdot(a\beta)(a\alpha)^5&=0,
\end{aligned}\right\}
}
wobei die zweite aus der ersten durch Vertauschung von $\alpha$ mit $\beta$ gewonnen wurde. Hiezu treten die beiden Relationen (3), welche wir schreiben können:
\weq{5}{
\left.\begin{aligned}
(a\alpha)^4(a\beta)^2&=0\\
(a\beta)^4(a\alpha)^2&=0
\end{aligned}\right\}
}


% ===== pdf idx 327, printed 313.0 =====
% !PRINTED 313
\weberpage{313}

Die vier Gleichungen (4) und (5) müssen also zusammen bestehen, wenn $m = \varrho \cdot l = \varrho \cdot \alpha \cdot \beta$ und $n = \varrho^2 \cdot \alpha \cdot \beta$ ist.

Nun kann das erste Gleichungspaar (4) aus drei verschiedenen Ursachen verschwinden. Es kann sein

\weq{1}{(a\alpha)^6=0,\quad (a\beta)(a\alpha)^5=0\ *)}

\weq{2}{(a\alpha)^6=0,\quad (a\beta)^6=0}

\weq{3}{(a\alpha)(a\beta)^5=0,\quad (a\beta)(a\alpha)^5=0.}

Demnach spaltet sich der Fall

\[ n = \varrho\, m = \varrho^2\, l \]

in drei Gruppen, deren jede wir nun einzeln untersuchen wollen.

288. \textit{Erste Gruppe von Bedingungen.} Es sei:

\weq{1}{(a\alpha)^6=0,}

\weq{2}{(a\alpha)^5(a\beta)=0,}

und hiezu

\weq{3}{(a\alpha)^4(a\beta)^2=0,}

\weq{4}{(a\beta)^4(a\alpha)^2=0.}

Sondert man von den ersten drei Bedingungen einen Moment den Factor $(a\alpha)^4$ ab, so sind die übrigbleibenden Factoren

\[ (a\alpha)^2,\quad (a\alpha)(a\beta),\quad (a\beta)^2 \]

die Coefficienten von $\alpha_x$ und $\beta_x$ in der Entwicklung

\[ (\alpha\beta)^2 a_x^2 = \{(a\beta)\alpha_x - (a\alpha)\beta_x\}^2. \]

Multiplicirt man demnach diese Relation wieder mit $(a\alpha)^4$, so kommt:

\[ (a\alpha)^4 a_x^2 (\alpha\beta)^2 = (a\alpha)^4(a\beta)^2\alpha_x^2 - 2(a\alpha)^5(a\alpha) a_x\beta_x + (a\alpha)^6\beta_x^2 = 0. \]

Es ist daher:

\[ (a\alpha)^4 a_x^2 = 0,\ \text{d. h.} \]

\glqq der Factor $\alpha_x$ ist ein cubischer Factor von $f$\grqq{}.

Dieser Fall führt also zu dem nämlichen Resultaten wie der bereits vorhin behandelte: % [sic]

\[ l = \alpha r,\quad m = \alpha s,\quad n = \alpha t. \]

289. \textit{Zweite Gruppe von Bedingungen.} Es möge ferner stattfinden:

\weq{II}{\begin{aligned} (a\alpha)^6&=0, & (a\beta)^6&=0\\ (a\alpha)^4(a\beta)^2&=0, & (a\beta)^4(a\alpha)^2&=0 \end{aligned}}

Versucht man die Form $f$ mit Hilfe der linearen Factoren $\alpha_x$ und $\beta_x$ darzustellen, so erhält man aus der binomischen Entwicklung

\noindent\rule{4cm}{0.4pt}

\noindent *) Der Fall $(a\beta)^6=0,\ (a\alpha)(a\beta)^5=0$ ist nicht verschieden von (1); denn er entsteht durch Vertauschung von $\alpha$ mit $\beta$.
%CONT>>


% ===== pdf idx 328, printed 314.0 =====
% !PRINTED 314
\weberpage{314}

%<<CONT
\[ \{(\alpha\beta) a_x\}^6 = \{(a\beta)\alpha_x - (a\alpha)\beta_x\}^6 \]

unter Anwendung der Relationen (II) das Resultat:

\[ (\alpha\beta)^6 a_x^6 = -\,6(a\beta)(a\alpha)\alpha_x\beta_x \left\{ (a\beta)^4\alpha_x^4 + \frac{10}{3}(a\beta)^2(a\alpha)^2\alpha_x^2\beta_x^2 + (a\alpha)^4\beta_x^4 \right\}, \]

d. h. die Form sechsten Grades zerfällt in das Product einer biquadratischen Form $\Psi$, multiplicirt in die Covariante $l=\alpha\beta$. Die Gleichung sechsten Grades kann also in diesem Falle durch Division mit der bekannten Covariante $l$ auf eine biquadratische reducirt werden. Der biquadratische Factor steht aber noch in engerem Zusammenhange mit $l$. Untersucht man nämlich wie bisher auch die Ueberschiebungen von $k$ über Potenzen von $\alpha$ und $\beta$, so zeigt sich, dass die vierte Ueberschiebung von $k$ über $\alpha^3\beta$ und über $\beta^3\alpha$ verschwindet. Denn man erhält, wenn man um auf die Bedingungen zu gelangen,

\[ (k\alpha)^3(k\beta) = 2(ab)^4(a\alpha)^2(b\alpha)(b\beta) \]

mit $(\alpha\beta)^4$ multiplicirt:

\begin{align*}
(ab)^4(\alpha\beta)^4(a\alpha)^2(b\alpha)(b\beta) &= (a\alpha)^2(b\alpha)(b\beta)\{(b\beta)(a\alpha) - (a\beta)(b\alpha)\}^4\\
&= (a\alpha)^6(b\alpha)(b\beta)^5 - 4(b\beta)^4(b\alpha)^2(a\alpha)^5(a\beta) + 6(a\alpha)^4(a\beta)^2(b\alpha)^3(b\beta)^3\\
&\quad - 4(b\alpha)^4(b\beta)^2(a\alpha)^3(a\beta)^3 + (a\beta)^4(a\alpha)^2(b\alpha)^5(b\beta) = 0.
\end{align*}

Es ist also in der That

\[ (k,\ \alpha^3\beta)^4 = 0 \quad \text{und ebenso} \quad (k,\ \beta^3\alpha)^4=0, \]

und deshalb auch [vgl. Nr. 287 $(2_a)$ und $(2_b)$]

\[ (\varDelta,\ \alpha^3\beta)^4 = 0 \quad \text{und} \quad (\varDelta,\ \beta^3\alpha)^4=0. \]

Will man also $k$ oder $\varDelta$ durch $\alpha$ und $\beta$ darstellen, so verschwinden aus diesem Grunde die Coefficienten der ungeraden Potenzen von $\alpha_x$ und $\beta_x$, d. h. $k$ sowohl als $\varDelta$ werden dadurch in die Normalform transformirt. Demnach sind $\alpha$ und $\beta$ gemäss der Theorie der biquadratischen Formen (vgl. Nr. 179) nichts anderes, als die linearen Factoren einer der drei irrationalen quadratischen Covarianten, in welche die Form $t$ von $k$ sich spalten lässt, d. h.

\glqq\textit{Die Form $f$ ist das Product einer biquadratischen Form in einer ihrer irrationalen quadratischen Covarianten.}\grqq{}

290. \textit{Dritte Gruppe von Bedingungen.} Ist endlich:

\[ (a\alpha)\,(a\beta)^5=0,\quad (a\beta)(a\alpha)^5=0, \]

\[ (a\alpha)^4(a\beta)^2=0,\quad (a\beta)^4(a\alpha)^2=0, \]
%CONT>>


% ===== pdf idx 329, printed 315.0 =====
% !PRINTED 315
\weberpage{315}

%<<CONT
so reducirt sich

\[ \{(\alpha\beta) a_x\}^6 = \{(a\beta)\alpha_x - (a\alpha)\beta_x\}^6 \]

auf die drei Glieder:

\[ (\alpha\beta)^6 a_x^6 = (\alpha\beta)^6\alpha_x^6 - 20(\alpha\beta)^3(a\alpha)^3\alpha_x^3\beta_x^3 + (a\alpha)^6\beta_x^6, \]

d. h. $f=0$ reducirt sich durch die Substitution $\alpha_x^3=x_1$, $\beta_x^3=x_2$ auf eine quadratische Gleichung. Untersucht man überdies die Ueberschiebungen von $k$ über $\alpha$ und $\beta$, so findet sich, dass nicht nur

\[ (k,\ \alpha^3\beta)^4=0 \quad \text{und} \quad (k,\ \alpha\beta^3)^4=0, \]

sondern auch, dass

\[ (k\alpha)^4=0 \quad \text{und} \quad (k\beta)^4=0, \]

wovon man sich leicht auf dem nun schon öfters eingeschlagenen Wege überzeugen kann. Daraus geht aber hervor, dass $k$ das Quadrat von $l$ ist; denn in der binomischen Entwicklung

\[ k_x^4(\alpha\beta)^4 = \{(k\beta)\alpha_x - (k\alpha)\beta_x\}^4 \]

verschwinden alle Glieder bis auf das Glied $(k\alpha)^2(k\beta)^2\alpha_x^2\beta_x^2$.

291. \textit{Zusammenfassung der Resultate.} Wir haben demnach folgende Sätze: Ist

\[ m = \varrho \cdot l = \varrho \cdot \alpha \cdot \beta, \]

so ist

\begin{enumerate}
\item[1)] entweder $\alpha_x$ ein cubischer Factor von $f$, und $\alpha_x$ ergiebt sich rational nach der Methode für Aufsuchung des grössten gemeinschaftlichen Theilers von $f$ und $l$,
\item[2)] oder $f$ zerfällt in eine biquadratische Form und die Covariante $l$, welche zugleich Factor der Form $t$ von jener Form $k$ ist,
\item[3)] oder endlich $f$ lässt sich durch die Substitution
\end{enumerate}

\[ x_1 = \alpha_x^3,\quad x_2 = \beta_x^3 \]

auf eine quadratische Gleichung reduciren.

292. \textit{Zweiter Fall; dritte Unterabtheilung.} Die Resultante $A_{\nu\nu}$ verschwindet endlich auch, wenn eine der quadratischen Formen $l$ oder $m$ null ist. Verschwindet $l$, so ist $m$ a priori null, da ja $m=(k,l)^2$ ist. Man kann aber auch beweisen, dass, wenn $m=0$, nothwendig auch $l$ verschwinden muss. Nehmen wir nämlich an, $l$ sei $\gtrless 0$, wenn $m=0$, dann ist wegen

\[ 0 = (\varDelta,\ m)^2 = \frac{B}{6}\, m + \frac{C}{3}\, l \quad \text{(vgl. Nr. 265, 2)}, \]

\[ \frac{C}{3}\, l = 0,\quad \text{also}\quad C=0, \]

und demnach, da ja $A_{lm}=(l,m)^2=0=\dfrac{2}{3}(B^2+AC)$ auch $B=0$.


% ===== pdf idx 330, printed 316.0 =====
% !PRINTED 316
\weberpage{316}

Es ist aber

\[ 2C + \frac{AB}{3} = A_{ll} \]

und somit verschwindet auch $A_{ll}$.

Nun ist (vgl. Nr. 266, 5)

\[ 0 = (f,\ m)^2 = \frac{l^2}{2} + \frac{B}{3}\, k + \frac{A}{3}\, \varDelta, \]

also

\[ 0 = \frac{l^2}{2} + \frac{A}{3}\, \varDelta. \]

Schieben wir andererntheils $(f,\ m)^2$ zweimal über $f$, so erhalten wir % ?? uncertain: reading of \textquotedbl{}andererntheils\textquotedbl{}

\begin{align*}
\big((f,\ m)^2,\ f\big)^2 &= \frac{1}{2}(l^2,\ f)^2 + \frac{A}{3}(\varDelta,\ f)^2\\
&= \frac{1}{2}\left\{\Big(2\varDelta + \frac{A}{3}\, k\Big) l - \frac{1}{3} A_{ll} f\right\} + \frac{A}{3}\left\{\frac{lk}{2} + \frac{B}{6} f\right\}\\
&= l\varDelta + \frac{lkA}{3}.
\end{align*}

Dieser Ausdruck ist null, da $(f,\ m)^2$ verschwindet; da aber $l$ nach Annahme von null verschieden sein soll, so ist

\[ \varDelta + \frac{A}{3}\, k = 0. \]

Schieben wir diesen Ausdruck zweimal über $k$, so kommt

\[ (\varDelta,\ k)^2 + \frac{A}{3}\, \varDelta = 0 = \frac{B}{6}\, f + \frac{A}{3}\, \varDelta. \]

Daher ist:

\[ \frac{A}{3}\, \varDelta = 0. \]

Es müsste also auch wegen

\[ \frac{l^2}{2} + \frac{A}{3}\, \varDelta = 0 \]

$l^2=0$ sein, was gegen unsere Annahme ist.

Wenn also eine der drei quadratischen Formen verschwindet, so verschwinden auch die beiden andern.

Nehmen wir nun an

\[ (1)\ l=0,\quad (2)\ k \lessgtr 0. \]

Dann ist, wegen $(f,\ l)^2 = 2\varDelta + \dfrac{A}{3} k = 0$

\[ \varDelta = \varrho \cdot k, \]

d. h. nach der Theorie der biquadratischen Formen: $k$ ist entweder ein volles Quadrat, etwa

\[ k = \alpha^2 \cdot \beta^2 \]
%CONT>>


% ===== pdf idx 331, printed 317.0 =====
% !PRINTED 317
\weberpage{317}

oder auch $k$ ist ein Biquadrat, etwa

\[
k = \alpha^4.
\]

Im ersten Falle, in welchem $(\alpha\beta)\genfrac{}{}{0pt}{}{\leqq}{\geqq}0$, ist also nach Voraussetzung

\[
l = (f, k)^4 = 0.
\]

Weil aber nach Nr. 253 (1) in jedem Falle $(f, k)^3$ gleich null ist, so hat man auch:

\[
(ak)^3 a_x^3 k_y = 0,
\]

da die Elementarcovarianten rechts verschwinden. Ersetzen wir links $k$ durch $\alpha^2\beta^2$, so erhalten wir

\[
a_x^3 (a\alpha)(a\beta)\{\alpha_y(a\beta) + \beta_y(a\alpha)\} = 0.
\]

Nun sind $\alpha$ und $\beta$ von einander verschieden und da diese Relation für jeden Werth von $y$ gilt, so müssen die Coefficienten von $\alpha_y$ und $\beta_y$ einzeln verschwinden, d. h. es muss sein:

\[
\begin{aligned}
a_x^3(a\alpha)(a\beta)^2 &= 0\\
a_x^3(a\alpha)^2(a\beta) &= 0.
\end{aligned}
\]

Multipliciren wir die erste Gleichung mit $\alpha_x$, die zweite mit $-\beta_x$ und addiren, so kommt:

\[
a_x^3(a\alpha)(a\beta)\{(a\beta)\alpha_x - (\alpha a)\beta_x\} = a_x^4(a\alpha)(a\beta)(\alpha\beta) = 0,
\]

also:

\[
a_x^4(a\alpha)(a\beta) = 0.
\]

Versucht man daher $a_x^6$ vermittelst der Identität

\[
a_x^6(\alpha\beta)^6 = \{(\alpha\beta)\alpha_x - (a\alpha)\beta_x\}^6
\]

darzustellen, so reducirt sich die binomische Entwicklung in diesem Falle auf

\[
a_x^6(a\beta)^6 = (a\beta)^6\alpha_x^6 - (a\alpha)^6\beta_x^6,
\]

d. h. \glqq Die Form $f$ reducirt sich auf eine binomische Gleichung, wenn $R=0$, $l=0$, $k=\alpha^2\beta^2$.\grqq{}

Ist aber neben $l=0$ die Form $k$ ein Biquadrat, also

\weq{1}{k = \alpha^4,}

dann ist immer noch:

\[
(f, k)^3 = 0,
\]

also wegen (1)

\[
a_x^3(a\alpha)^3 = 0,
\]

d. h. \glqq Die Form $f$ besitzt den linearen Factor $\alpha$ vierfach.\grqq{}

Umgekehrt: Besitzt $f$ einen vierfachen linearen Factor $\alpha$, so ist $k$ ein Biquadrat und $l$ identisch null. Denn es ist dann:
%CONT>>


% ===== pdf idx 332, printed 318.0 =====
% !PRINTED 318
\weberpage{318}
%<<CONT

\[
\begin{aligned}
\alpha^4(k\alpha) &= (ab)^4\alpha^4\cdot(a\alpha)\alpha_x\beta_x^2\\
&= (a\alpha)a_xb_x^2\{(a\alpha)b_x - (b\alpha)a_x\}^4\\
&= 0, \quad \text{wegen } a_x^3(a\alpha)^3 = 0.
\end{aligned}
\]

Also: $(k, \alpha) = 0$, d. h. $k = \alpha_x^4$.

Deshalb ist aber auch $(f, k)^4 = 0$, d. h. $l = 0$, wie behauptet wurde.

\textit{Anmerkung.} Ist $k$ identisch null, dann ist $f$ entweder eine Octaederform und kann nach der in Nr. 195 gegebenen Methode aufgelöst werden, oder $f$ besitzt einen linearen Factor $\alpha_x$ mindestens fünffach [vgl. § 19, über die Formen, für welche $(f,f)^4=0$].

293. \textit{Dritter Fall}: $A=0$, $C=0$. Wir wollen endlich noch zwei Fälle erwähnen, in welchen die Invariantenrelationen zwar nicht auf algebraisch, so doch mit Hilfe elliptischer Functionen lösbare Gleichungen herführen. (Vgl. Clebsch § 114 der \glqq Binären Formen\grqq{}.) Bildet man nämlich das Formensystem der Modulargleichung (Jacobi, \glqq Fundamenta\grqq{}, S. 27) mit dem einen Parameter $u$:

\weq{1}{v^6 + 4u^5v^5 + 5u^2v^4 - 5u^4v^2 - 4uv - u^6 = 0}

für die Transformation fünfter Ordnung*), so zeigt sich, dass die Invarianten $A$ und $C$ identisch verschwinden. Umgekehrt hat dann Clebsch (a. a. O.) gezeigt, dass $f = a_x^6$ sich durch eine lineare Transformation in die Form

\weq{2}{\varphi = \xi^6 + 5\xi^4\eta^2 + 15\xi^2\eta^4 - 4Z\xi\eta^5 - 5\eta^6}

mit dem einen Parameter $Z$, welche aus (1) durch die Substitution

\[
v = u\,\frac{\dfrac{\xi}{\eta} + u^4}{1 - u^4\dfrac{\xi}{\eta}}
\]

hervorgeht, stets überführen lässt, sobald ihre Invarianten $A$ und $C$ verschwinden. In diesem Falle besitzen die beiden quadratischen Covarianten $l$ und $\mu$ einen gemeinsamen Factor und

\[
\frac{\eta}{\xi} = \frac{l\sqrt{-BD}}{2\mu}
\]

ist die erwähnte lineare Transformation der Form $f$ in $\varphi$.

\footnote{*) Zur Erklärung möge hier beigefügt werden: Ist $k$ der Modul des elliptischen Integrals $\displaystyle\int\frac{dx}{\sqrt{(1-x^2)(1-k^2x^2)}}$, $\lambda$ der Modul des neuen Integrals, das aus ihm durch eine Transformation fünfter Ordnung hervorgeht, dann besteht zwischen $\sqrt[4]{k}=u$ und $\sqrt[4]{\lambda}=v$ die Gleichung (1).}


% ===== pdf idx 333, printed 319.0 =====
% !PRINTED 319
\weberpage{319}

294. \textit{Vierter Fall}: $B=\dfrac{7}{50}A^2$, $C=-\dfrac{9}{500}A^3$. Eine zweite Gleichungsgattung mit einem einzigen Parameter repräsentirt die Multiplicatorgleichung von Brioschi:

\weq{3}{z^6 - 4z^5 + 256k^2k'^2(z+1) = 0, \quad (k'^2 = 1 - k^2),}

auf welche die Frage nach dem Multiplicator $M$ führt, vermöge welchem die Beziehung besteht:

\[
\frac{dy}{\sqrt{(1-y^2)(1-\lambda^2y^2)}} = \frac{1}{M}\cdot\frac{dx}{\sqrt{(1-x^2)(1-\varkappa^2x^2)}}.
\]

Dieser Multiplicator $M$ hängt mit $v$ in (1) durch die Gleichung

\[
M = \frac{v(1-uv^3)}{v-u^5}
\]

und mit $z$ in (3) durch

\[
M = \frac{1}{z+1}
\]

zusammen. Setzt man in (3) $z=\lambda\cdot\dfrac{\xi}{\eta}$, so kann man sie nach Clebsch in der Form schreiben:

\weq{4}{a_0\xi^6 + 6a_1\xi^5\eta + 6a_5\xi\eta^5 + a_6\eta^6 = 0,}

wobei

\[
a_0 = \varrho\cdot\lambda^6, \quad a_1 = -\frac{2}{3}\varrho\lambda^5,
\]

\[
a_5 = \frac{128}{3}\varrho\cdot\lambda k^2k'^2, \quad a_6 = 256\varrho\cdot k^2k'^2
\]

und

\[
a_0a_6 + 9a_1a_5 = 0.
\]

Für diese Gleichung (4) findet man nun durch directe Rechnung die beiden sie charakterisirenden Invarianteneigenschaften:

\[
B = \frac{7}{50}A^2, \quad C = -\frac{9}{500}A^3,
\]

und umgekehrt sind dieselben genügend, um die allgemeine Form $f$ durch die aus

\[
-8(a_1^3a_6^2 + a_0^2a_5^3)\xi\cdot\eta = m - \frac{A}{10}l
\]

sich ergebende Transformation in die Form (4) überzuführen.

\textit{Anmerkung 1.} Die Modulargleichung (1) wurde von Hermite zur Auflösung der Gleichung fünften Grades benutzt. Er fand, dass diese Gleichung sechsten Grades eine Resolvente fünften Grades besitzt, für welche $C=0$ ist, die sich also in die Form bringen lässt:

\[
t^5 - t - \frac{1}{5\sqrt[4]{z}} = 0.
\]


% ===== pdf idx 334, printed 320.0 =====
% !PRINTED 320
\weberpage{320}

Brioschi dagegen ging von der Multiplicatorgleichung (3) aus und kam von dieser auf dieselbe Resolvente wie Hermite. (Vgl. die diesbezüglichen interessanten historischen Entwicklungen in: Klein, \glqq Vorlesungen über das Ikosaeder\grqq{}, insbesondere § 3, S. 144.)

\textit{Anmerkung 2.} Ich erwähne hier noch der Vollständigkeit halber die beiden Brill'schen Normalformen der Gleichung sechsten Grades

\[
\begin{aligned}
&x^6 + 2px^5 + 3qx^4 + 4rx^3 + 3x^2 + 2px + q = 0\\
&x^6 + ax^4 + bx^3 + cx^2 + 1 = 0,
\end{aligned}
\]

in welche sich jede Gleichung sechsten Grades nach Adjunction einer Wurzel einer Gleichung fünften Grades (vgl. Nr. 273--276) linear transformiren lässt. Die Discriminante dieser Gleichungen zerfällt in zwei rationale Factoren (vgl. Math. Annalen, Bd. XX, S. 330).

\begin{center}
\S.\,31. \gesperrt{Simultanes System einer quadratischen und cubischen Form.}
\end{center}

295. \textit{Allgemeine Betrachtungen.} Die simultanen Systeme mehrerer binärer Formen spielen in vielen Problemen der Mathematik eine hervorragende Rolle. Ich erinnere nur an die Lösung der Gleichung neunten Grades, welcher die neun Wendepunkte einer Curve dritter Ordnung genügen, an die Dreitheilung elliptischer und hyperelliptischer Functionen, an die Reduction elliptischer Integrale erster Gattung auf die Normalform, an die Theorie rationaler Curven, an die Untersuchungen über apolare und involutorische Systeme etc. Es erscheint daher angezeigt, an Beispielen zu zeigen, in welcher Weise derartige Systeme aufgestellt werden. Die einfachsten und sofort allgemein zu behandelnden derartigen Systeme sind simultane Systeme von $n$ linearen Formen, die aus den $n$ Formen selbst und den $\dfrac{n(n-1)}{2}$ Determinanten $(ab)$ bestehen; ferner das simultane System von $n$ quadratischen Formen, das wir bereits früher aufgestellt haben. Es bietet auch wenig Schwierigkeiten, die Formen des vollständigen Systemes zu ermitteln, das entsteht, wenn ein bereits bekanntes volles System mit dem vollen System von $n$ linearen oder $n$ quadratischen Formen überschoben werden soll. Dagegen werden die Untersuchungen bereits complicirter, sobald es sich auch nur um das simultane System von zwei Formen handelt, von denen keine quadratisch oder linear ist. Wir werden demnach im Folgenden zwei Beispiele ins Auge fassen; in dem ersten wird eine der Grundformen eine quadratische Form sein, die andere cubisch; im zweiten werden beide Formen den dritten Grad besitzen. Wir wenden uns nun dazu, das simultane System einer quadratischen und cubischen Form zu studiren.


% ===== pdf idx 335, printed 321.0 =====
% !PRINTED 321
\weberpage{321}

296. \textit{Disposition der Aufgabe.} Die beiden gegebenen Formen seien
\[
f = a_x^2, \quad \varphi = \alpha_x^3;
\]
das volle System von $f$ besteht aus den zwei Formen
\[
f, \quad A_{ff},
\]
das von $\varphi$ aus den vier Formen
\weq{2}{\varphi, \quad \Delta, \quad Q, \quad R = A_{\Delta\Delta}.}
Durch Ueberschiebung vollständiger Systeme entstehen wieder vollständige Systeme. Das volle simultane System von $f$ und $\varphi$ erhalten wir also, wenn wir Producte
\[
U = f^\varrho \cdot A_{ff}^\sigma
\]
mit Producten
\[
V = \varphi^\varkappa \cdot \Delta^\lambda \cdot Q^\mu \cdot A_{\Delta\Delta}^\nu
\]
überschieben, d. h. mit einander multipliciren und auf alle möglichen Arten zur Faltung bringen. Da bei der Faltung Invarianten nicht betheiligt sind, und da ferner $Q$ Functionaldeterminante von $\varphi$ und $\Delta$ ist, so können wir in diesen Producten von Factoren $A_{ff}$ und $A_{\Delta\Delta}$ absehen und höhere Potenzen von $Q$ nach der Relation
\[
Q^2 = -\frac{1}{2}\{\Delta^3 + A_{\Delta\Delta}\varphi^2\}
\]
durch $\varphi$ und $\Delta$ ersetzt denken. Unsere Aufgabe reducirt sich daher darauf:

\glqq Alle jene nicht zerfallenden Formen zu finden, die aus den Faltungen
\[
(U, \; V)^\alpha = (f^\varrho, \; \varphi^\varkappa \cdot \Delta^\lambda \cdot Q)^\alpha
\]
hervorgehen, wenn $\varrho, \varkappa, \lambda, \alpha$ alle möglichen Werthe annehmen.\grqq{}

Wir gehen hiebei, wie früher, schrittweise vor, indem wir zunächst solche Producte $V$ mit $f^\varrho$ falten, die nur eine Form $\varphi^\mu$, oder $\Delta^\mu$, oder $Q$ enthalten. Erst in zweiter Linie werden Faltungen mit Producten $V$ untersucht, die mehr als eine der drei Formen zu Factoren haben.

297. \textit{Erste Gruppe von Ueberschiebungen.} Wir haben demnach zunächst zu untersuchen: Zu welchen Formen geben die Ueberschiebungen
\[
(f^\varrho, \; \Delta^\mu)^\alpha, \quad (f^\varrho, \; \varphi^\mu)^\alpha, \quad (f^\varrho, \; Q)^\alpha
\]
Veranlassung? Aus dem ersten Operationsymbol ergeben sich nur zwei Formen:
\[
(f, \; \Delta) \text{ und } (f, \; \Delta)^2 = A_{f\Delta},
\]
%CONT>>


% ===== pdf idx 336, printed 322.0 =====
% !PRINTED 322
\weberpage{322}
%<<CONT
wie aus der Theorie zweier quadratischer Formen bekannt ist. Das zweite Operationssymbol führt zunächst zu folgenden sechs Formen:
\[
(f, \; \varphi), \quad (f, \; \varphi)^2, \quad (f^2, \; \varphi)^3, \quad (f^2, \; \varphi^2)^4, \quad (f^3, \; \varphi^2)^5, \quad (f^3, \; \varphi^2)^6.
\]
Von diesen sind aber die vierte und fünfte Form reducibel. Denn
\[
(f^2, \; \varphi^2)^4 \text{ enthält das zerfallende Glied } (f, \; \varphi)^2 \cdot (f, \; \varphi)^2
\]
und
\[
(f^3, \; \varphi^2)^5 \quad \text{\textquotedbl} \quad \text{\textquotedbl} \quad \text{\textquotedbl} \quad \text{\textquotedbl} \quad (f, \; \varphi)^2 \cdot (f^2, \; \varphi)^3.
\]
Die vier anderen Formen weisen dagegen keine zerfallenden Glieder auf. Höhere Potenzen von $f$ und $\varphi$ brauchen nicht mehr überschoben zu werden. Denn jede solche Ueberschiebung $(f^\varrho, \; \varphi^\mu)^\alpha$ würde a priori ein Glied besitzen, welches zerfällt, nämlich das Glied
\[
(f^3, \; \varphi^2)^6 \cdot (f^{\varrho-3}, \; \varphi^{\mu-2})^{\alpha-6}.
\]
Das dritte Operationssymbol endlich liefert überhaupt nur drei Formen
\[
(Q, \; f), \quad (Q, \; f)^2, \quad (Q, \; f^2)^3,
\]
von denen die erste als Functionaldeterminante einer Functionaldeterminante reducibel ist, während die beiden anderen Formen in das System aufgenommen werden müssen.

298. \textit{Zweite Gruppe von Ueberschiebungen.} Wir überschieben nun Producte $V$ über $f^\varrho$ von der Form
\[
V = \Delta\varphi, \quad \Delta Q, \quad \Delta\varphi Q, \quad \varphi Q.
\]
Allein da die drei ersten Producte den Factor $\Delta$ besitzen, so wird jede der Ueberschiebungen
\[
(f^\varrho, \; \Delta\varphi)^\alpha, \quad (f^\varrho, \; \Delta Q)^\alpha, \quad (f^\varrho, \; \Delta\varphi Q)^\alpha
\]
ein Glied enthalten, welches eine Potenz von $A_{f\Delta}$ zum Factor hat, also zerfällt. Es erübrigt daher nur noch, die Ueberschiebungen
\[
(f^\varrho, \; \varphi Q)^\alpha
\]
zu untersuchen. Hiebei muss $\alpha$ a priori grösser als 3 und demnach $\varrho$ mindestens gleich 2 sein. Wir haben somit zunächst die drei Formen:
\[
(f^2, \; \varphi Q)^4, \quad (f^3, \; \varphi Q)^5, \quad (f^3, \; \varphi Q)^6.
\]
Nur die letzte derselben, eine Invariante, enthält kein zerfallendes Glied; die erste dagegen enthält das Glied $(f, \; \varphi)^2 \cdot (f, \; Q)^2$, die zweite das Glied $(f, \; \varphi)^2 \cdot (f^2, \; Q)^3$. Weitere Ueberschiebungen zu untersuchen ist überflüssig; sie besitzen stets Glieder, welche die Invariante $(f^3, \; \varphi Q)^6$ zum Factor haben.

299. \textit{Zusammenfassung der Resultate.} Das simultane vollständige System von $f$ und $\varphi$ besteht demnach aus folgenden 15 Formen:
%CONT>>


% ===== pdf idx 337, printed 323.0 =====
% !PRINTED 323
\weberpage{323}
%<<CONT
\begin{tabular}{@{}ll@{}}
Drei cubische Covarianten: & $\varphi, \; Q, \; (f, \; \varphi) = \vartheta$ \\
Drei quadratische \quad ,, & $f, \; \Delta, \; (f, \; \Delta) = \Theta$ \\
Vier lineare \quad ,, & $(f, \; \varphi)^2, \; (f^2, \; \varphi)^3, \; (f, \; Q)^2, \; (f^2, \; Q)^3$ \\
Fünf Invarianten: & $A_{ff}, \; A_{\Delta\Delta}, \; A_{f\Delta}, \; (f^3, \; \varphi^2)^6, \; (f^3, \; \varphi \cdot Q)^6.$ \\
\end{tabular}

Zwischen den fünf Invarianten muss eine Relation bestehen. Denn die beiden Formen $f$ und $\varphi$ besitzen nur sieben Constante, von denen drei durch lineare Transformation zum Verschwinden gebracht werden können: demnach können nur vier von einander unabhängig sein.

300. \textit{Einführung einfacherer Symbole für die Formen des Systemes.} Für die symbolische Rechnung ist es wiederum vortheilhaft, jene Formen des Systemes, die noch nicht durch ein einfaches symbolisches Product dargestellt sind, durch eines ihrer Glieder zu ersetzen. Solche Formen sind die beiden linearen Covarianten $(f, \; Q)^2$ und $(f^2, \; Q)^3$ und die beiden Invarianten $(f^3, \; \varphi^2)^6$, $(f^3, \; \varphi \cdot Q)^6$. Setzen wir nämlich
\[
(f, \; \varphi)^2 = (a\alpha)^2 \alpha_x = p
\]
\[
(f^2, \; \varphi)^3 = (a\alpha)^2 (b\alpha) b_x = q = (f, \; p),
\]
so lassen sich an Stelle der beiden andern linearen Covarianten $(f, \; Q)^2$ und $(f^2, \; Q)^3$ die Formen:
\weq{1}{(p, \; \Delta) = (p\Delta)\Delta_x = r}
resp.:
\weq{2}{(\Theta, \; p) = (\Theta p)\Theta_x = s}
und an Stelle der beiden erwähnten Invarianten die Formen
\weq{3}{(f, \; p^2)^2 = (ap)^2 = F}
resp.
\weq{4}{(\Theta, \; p)^2 = (sp) = -M}
einführen.

Denn es ist:
\[
(f, \; Q)^2 = (a_x^2, \; (\alpha\Delta)\Delta_x \alpha_x^2)^2
\]
und $(\alpha\Delta)(a\alpha)^2\Delta_x = (p, \; \Delta) = r$ ist ein Glied dieser Ueberschiebung. Es ist ferner:
\[
(f^2, \; Q)^3 = (a_x^2 b_x^2, \; (\alpha\Delta)\Delta_x \alpha_x^2)^3.
\]
Greifen wir auch hier das Glied heraus, das durch Faltung der Factoren $a_x^2$ mit $\alpha_x^2$, $b_x$ mit $\Delta_x$ entsteht, so können wir dasselbe durch:
\[
(a\alpha)^2(\alpha\Delta)(b\Delta)b_x = (p\Delta)(b\Delta)b_x
\]
darstellen. Das symbolische Product $(p\Delta)(b\Delta)b_x$ ist aber ein Glied der Ueberschiebung $(p, \; (b\Delta)b_x\Delta_x) = (p, \; \Theta)$, und demnach lässt sich in der That $(f^2, \; \Theta)^3$ durch $(\Theta, \; p) = -(p, \; \Theta)$ ersetzen.


% ===== pdf idx 338, printed 324.0 =====
% !PRINTED 324
\weberpage{324}
Des weitern besitzt die Invariante $(f^3, \; \varphi^2)^6$ das Glied
\[
(a\alpha)^2(b\beta)^2(c\alpha)(c\beta) = (f, \; p^2)^2 = (ap)^2 = F
\]
und endlich die Invariante $(f^3, \; \varphi \cdot Q)^6 = (a_x^2 b_x^2 c_x^2, \; \alpha_x^3 \cdot (\beta\Delta)\Delta_x \beta_x^2)^6$ das Glied:
\[
(\beta\Delta) \cdot (a\alpha)^2(b\alpha)(b\Delta)(c\beta)^2 = -((b\Delta)b_x\Delta_x, \; p^2)^2 = -(\Theta, \; p)^2.
\]
Die Tabelle des simultanen Formensystemes wird demnach nach Einführung dieser neuen Symbole:
\[
1) \quad \varphi = a_x^3, \qquad Q = (\alpha\Delta)a_x^2\Delta_x, \qquad \vartheta = (a\alpha)a_x\alpha_x^2;
\]
\[
2) \quad f = a_x^2, \qquad \Delta = (\alpha\beta)^2\alpha_x\beta_x, \qquad \Theta = (a\Delta)a_x\Delta_x;
\]
\[
3) \quad p = (a\alpha)^2\alpha_x, \qquad q = (ap)a_x, \qquad r = (p\Delta)\Delta_x,
\]
\[
\qquad s = (\Theta p)\Theta_x;
\]
\[
4) \quad A_{ff} = (ab)^2, \qquad A_{\Delta\Delta} = (\Delta\Delta_1)^2, \qquad A_{f\Delta} = (a\Delta)^2,
\]
\[
\qquad F = (ap)^2, \qquad M = -(\Theta p)^2.
\]
Die Ueberschiebungen dieser Formen geben zu zahlreichen Relationen Veranlassung, welche bei den verschiedenen Aufgaben, in denen gerade dieses simultane System eine Rolle spielt, von Wichtigkeit werden. Wir gelangen zu diesen Relationen einmal, indem wir frühere allgemeine Theorien über simultane Systeme oder Functionaldeterminanten benutzen, dann durch directes Ueberschieben, insbesondere der quadratischen über die linearen Formen, und der linearen unter sich, endlich aber auch, indem wir versuchen, ein und dieselbe Form auf mehrere Arten darzustellen. Jedem dieser drei Wege entsprechend, wollen wir nun eine Reihe von Relationen ermitteln.

301. \textit{Erste Reihe von Relationen.} Das simultane System von $f$ und $\Delta$, welches aus den sechs Formen
\[
f, \quad \Delta, \quad \Theta, \quad A_{ff}, \quad A_{\Delta\Delta}, \quad A_{f\Delta}
\]
besteht, liefert sofort eine Reihe von Relationen, die wir schon in der Theorie simultaner quadratischer Formen aufgestellt haben. Wir erhielten damals (vgl. Nr. 127):
\weq{I}{\left.\begin{aligned} 2A_{\Theta\Theta} &= A_{ff}A_{\Delta\Delta} - A_{f\Delta}^2 \\ A_{f\Theta} &= 0, \quad A_{\Delta\Theta} = 0 \end{aligned}\right\}}
\weq{II}{\left.\begin{aligned} (1) &\quad 2(\Theta, \; f) = A_{ff}\Delta - A_{f\Delta}f \\ (2) &\quad 2(\Delta, \; \Theta) = A_{\Delta\Delta}f - A_{f\Delta}\Delta \\ (3) &\quad -2\Theta^2 = A_{ff}\Delta^2 - 2A_{f\Delta}f\Delta + A_{\Delta\Delta}f^2 \end{aligned}\right\}}
Die letzte dieser Relationen liefert sofort die in Nr. 299 erwähnte Invarianten-Beziehung; denn ersetzen wir in ihr $x$ durch $p$ und setzen einstweilen
\[
(\Delta p)^2 = L,
\]
%CONT>>


% ===== pdf idx 339, printed 325.0 =====
% !PRINTED 325
\weberpage{325}
so kommt:

\weq{III}{-\,2\,M^2 = A_{ff}L^2 - 2\,A_{f\Delta}F\cdot L + A_{\Delta\Delta}F^2\,.}
Wie wir hier das Quadrat der beiden schiefen Formen $\Theta$ und $M$ durch
gerade Formen dargestellt haben, so können wir auch die Quadrate der drei
anderen schiefen Formen, $q$, $r$, $s$ des Systemes durch gerade Formen
darstellen. Man findet nämlich in einfacher Weise

\weq{IV}{\left.\begin{aligned}
(1)&\quad q^2 = f\cdot F - \tfrac{1}{2}\,p^2\,A_{ff}\\[6pt]
(2)&\quad r^2 = \Delta\cdot L - \tfrac{1}{2}\,p^2\,A_{\Delta\Delta}\\[6pt]
(3)&\quad -\,s^2 = \Theta\cdot M + \tfrac{1}{2}\,p^2\,A_{\Theta\Theta}
\end{aligned}\right\}\,,}
indem man entweder die früher aufgestellten Sätze über
Functionaldeterminanten zu Grunde legt, oder noch kürzer, indem man
Identitäten wie $p_x\,(ab) = a_x\,(bp) - b_x\,(ap)$ und die beiden analogen
quadrirt.

302. \textit{Zweite Reihe von Relationen.} Wir können ferner
Formenbeziehungen aufstellen durch directe Ueberschiebung der Formen $p$,
$q$, $r$, $s$ unter sich oder über die Formen $f$, $\Delta$, $\Theta$. Sie
sind durchwegs leicht zu berechnen, und wir können uns daher auf solche
Beispiele beschränken, die wir später bei der typischen Darstellung von $f$
und $\varphi$ verwenden werden. Das sind zunächst die Ueberschiebungen von
$f$ über die linearen Covarianten. Man findet:

\weq{V}{\left.\begin{aligned}
(f,\ p) &= \phantom{-}q \quad \text{(vgl. Nr. 300)}\\[6pt]
(f,\ q) &= -\tfrac{1}{2}\,A_{ff}\,p\\[6pt]
(f,\ r) &= -\,s + \tfrac{1}{2}\,A_{f\Delta}\,p\,^{*)}\\[6pt]
(f,\ s) &= +\tfrac{1}{2}\left\{A_{f\Delta}\,q + A_{ff}\,r\right\}
\end{aligned}\right\}.}
Zur Verification der drei letzten Relationen dieser Gruppe mögen die
Entwicklungen dienen:

\[
\begin{aligned}
(f,\ q) &= (f,\ (f,\ p)) = (ab)(bp)a_x = \tfrac{1}{2}\,(ab)\left\{(bp)a_x -
(ap)b_x\right\}\\
&= -\tfrac{1}{2}\,(ab)^2\cdot p\\
(f,\ r) &= (f,\ (p,\ \Delta)) = (p\Delta)(a\Delta)a_x\,;\quad
(a\Delta)a_x\Delta_y = \Theta_x\Theta_y + \tfrac{1}{2}\,A_{f\Delta}(xy)\\
&\quad\text{[vgl. Nr. 125 (2)]}\\
(f,\ s) &= ((f,\ (\Theta,\ p))) = (\Theta p)(a\Theta)a_x = -((\Theta,f),p) =
-\tfrac{1}{2}(A_{ff}\Delta - A_{f\Delta}f,\ p)\,.
\end{aligned}
\]

\par\noindent\rule{0.3\linewidth}{0.4pt}\par
{\footnotesize $^{*)}$ In Clebsch, \glqq Bin\"are Formen\grqq{} ist $(f,\
r) = +\,s$.}


% ===== pdf idx 340, printed 326.0 =====
% !PRINTED 326
\weberpage{326}
Wir fügen hinzu:

\weq{VI}{\left.\begin{aligned}
(qp) &= ((f,\ p),\ p) = (ap)^2 = F\\
(rp) &= -\,((\Delta,\ p),\ p) = -\,(\Delta p)^2 = -\,L\\
(sp) &= ((\Theta,\ p),\ p) = (\Theta p)^2 = -\,M \quad\text{\% ?? uncertain: faint mark after ``='' before ``$-\,M$''}\\
(qr) &= ((f,\ p),\ (p,\ \Delta)) = (-(f,\ \Delta),\ p^2)^2 = -\,(\Theta p)^2 = +\,M\\
(qs) &= ((f,\ p),\ (\Theta,\ p)) = ((f,\ \Theta),\ p^2)^2 = -\tfrac{1}{2}\left\{A_{ff}L - A_{f\Delta}F\right\}\\
(rs) &= ((\Delta,\ p),\ (\Theta,\ p)) = ((\Delta,\ \Theta),\ p^2)^2 = -\tfrac{1}{2}\left\{A_{f\Delta}L - A_{\Delta\Delta}F\right\}
\end{aligned}\right\}.}
Mit Hilfe dieser beiden Gruppen (V) und (VI) von Relationen lassen sich nun
auch sofort die Werthe jener Invarianten anschreiben, welche entstehen,
wenn wir die Form $f$ zweimal über irgend ein Product zweier linearer
Covarianten schieben. Man findet unmittelbar:

\weq{VII}{\left.\begin{aligned}
(f,\ pq)^2 &= (q,\ q) = 0\\
(f,\ pr)^2 &= (q,\ r) = +\,M\\
(f,\ ps)^2 &= (q,\ s) = -\tfrac{1}{2}\left\{A_{ff}L - A_{f\Delta}F\right\}\\
(f,\ qr)^2 &= -\tfrac{1}{2}\,A_{ff}(p,\ r) = -\tfrac{1}{2}\,A_{ff}L\\
(f,\ qs)^2 &= -\tfrac{1}{2}\,A_{ff}(p,\ s) = -\tfrac{1}{2}\,A_{ff}M\\
(f,\ rs)^2 &= \left(-\,s + \tfrac{1}{2}\,A_{f\Delta}p,\ s\right) = \tfrac{1}{2}\,A_{f\Delta}M
\end{aligned}\right\}.}
303. \textit{Dritte Reihe von Relationen.} Die dritte Reihe von Relationen
hat ihren Ursprung in der Fundamentalrelation dieses simultanen Systemes

\weq{VIII}{\Theta = (f,\ \Delta) = (\varphi,\ p)\,.}
Wir beweisen dieselbe, indem wir die Form

\[
\Theta = (a\Delta)\,a_x\,\Delta_x = (\alpha\beta)^2(a\alpha)\,a_x\,\beta_x
\]
mit Hilfe des Identitätssatzes umformen. Man hat nämlich

\[
\begin{aligned}
\Theta &= (\alpha\beta)^2(a\alpha)\,a_x\,\beta_x = (a\alpha)(\alpha\beta)\,\beta_x\left\{\alpha_x(a\beta) - \beta_x(a\alpha)\right\}\\
&= (a\alpha)(a\beta)(\alpha\beta)\,\alpha_x\,\beta_x - (a\alpha)^2(\alpha\beta)\,\beta_x^2\\
&= -\,(a\alpha)^2(\alpha\beta)\,\beta_x^2\,,
\end{aligned}
\]
da das erste Glied rechts identisch verschwindet. Führen wir aber in
$(a\alpha)^2(\alpha\beta)\,\beta_x$ das Symbol $p_x = (a\alpha)^2\alpha_x$
ein, so kommt:

\[
\Theta = -\,(p\beta)\,\beta_x^2 = (\varphi,\ p)\,.
\]
Aus dieser fundamentalen Relation (VIII) gehen insbesondere folgende
wichtige Beziehungen hervor:
%CONT>>


% ===== pdf idx 341, printed 327.0 =====
% !PRINTED 327
\weberpage{327}
%<<CONT
\weq{IX}{\left.\begin{aligned}
s &= (\Theta,\ p) = ((\varphi,\ p),\ p) = (\varphi,\ p)^2 = (\alpha p)^2\alpha_x\\
M &= -\,(s,\ p) = -\,((\varphi,\ p)^2,\ p) = -\,(\alpha p)^3\\
-\,r &= (\Delta p)\,\Delta_x = (\alpha\beta)^2(\beta p)\,\alpha_x^2 = (\alpha_x^3,\ (\beta p)\,\beta_x)^2 = (\varphi,\ \Theta)^2\\
L &= (\Delta p)^2 = (\alpha\beta)^2(\alpha p)(\beta p) = ((\alpha p)\,\alpha_x^2,\ (\beta p)\,\beta_x^2) = A_{\Theta\Theta}\\
&= \tfrac{1}{2}\left\{A_{ff}A_{\Delta\Delta} - A_{f\Delta}^2\right\}\\
(\Theta,f) &= p^2 + (\varphi,\ q)
\end{aligned}\right\}.}
Was die letzte dieser Relationen betrifft, so ergiebt sich ihre
Richtigkeit aus folgender Entwicklung. Es ist

\[
\begin{aligned}
(\Theta,\ f) &= ((\varphi,\ p),\ f) = ((\alpha p)\,\alpha_x,\ a_x^2) = (\alpha p)(\alpha a)\,\alpha_x\,a_x\\
&= (\alpha a)\,\alpha_x\left\{\alpha_x(ap) + p_x(\alpha a)\right\}\\
&= (\alpha q)\,\alpha_x^2 + p_x^2 = (\varphi,\ q) + p^2\,.
\end{aligned}
\]
Wir wollen damit die Untersuchungen über die Covariantenbeziehungen des
simultanen Systemes von $f$ und $\varphi$ beschliessen und nun zu einigen
Anwendungen übergehen.

304. \textit{Typische Darstellung.} Die beiden Formen $f$ und $\varphi$
lassen sich mit Hilfe zweier ihrer vier simultanen linearen Invarianten
typisch darstellen. Gehen wir zu dem Zwecke wiederum von Identitäten wie

\[
\begin{aligned}
(pq)^2\,a_x^2 &= \left\{(aq)\,p_x - (ap)\,q_x\right\}^2\\
(pq)^3\,\alpha_x^3 &= \left\{(\alpha q)\,p_x - (\alpha p)\,q_x\right\}^3
\end{aligned}
\]
aus, so ergeben sich für die erste derselben die Werthe der typischen
Coefficienten unmittelbar aus den Relationen in Nr. 302. Man hat
beispielsweise:

\[
\begin{aligned}
F^2\cdot f &= (aq)^2\cdot p_x^2 - 2\,(ap)(aq)\,p_x\,q_x + (ap)^2\,q_x^2\\
&= (f,\ q)^2\cdot p_x^2 - 2\,(f,\ pq)\,p_x\,q_x + (f,\ p)^2\cdot q_x^2\\
&= \tfrac{1}{2}\,A_{ff}F\cdot p_x^2 + F\cdot q_x^2
\end{aligned}
\]
oder

\[
F\cdot f = \tfrac{1}{2}\,A_{ff}\,p_x + q_x^2\,,
\]
ein Resultat, das wir auch schon an anderer Stelle (Nr. 301, IV) ermittelt
haben.

Die typischen Coefficienten für die zweite Identität

\[
-\,F^3\cdot\varphi = (\alpha q)^3\,p_x^3 - 3\,(\alpha q)^2(\alpha p)\,p_x^2\,q_x + 3\,(\alpha q)(\alpha p)^2\,p_x\,q_x^2 - (\alpha p)^3\,q_x^3
\]
berechnet man mit Hilfe der in Nr. 303 entwickelten Relationen:

\[
\begin{aligned}
\Theta &= (\varphi,\ p)\ = (\alpha p)\,\alpha_x^2\\
s &= (\varphi,\ p)^2 = (\alpha p)^2\,\alpha_x\,.
\end{aligned}
\]


% ===== pdf idx 342, printed 328.0 =====
% !PRINTED 328
\weberpage{328}
Aus diesen ergiebt sich sofort:

\[
\begin{aligned}
(\alpha p)^3 &= (sp) = -\,M\\
(\alpha p)^2(\alpha q) &= (sq) = \tfrac{1}{2}\left\{A_{ff}L - A_{f\Delta}F\right\}\\
(\alpha p)(\alpha q)^2 &= (\Theta,\ q^2)^2 = \left(\Theta,\ F\cdot f - \tfrac{1}{2}\,A_{ff}p^2\right)^2 = -\tfrac{1}{2}\,A_{ff}(\Theta p)^2 = \tfrac{M}{2}\cdot A_{ff}\\
(\alpha q)^3 &= \left((\alpha q)\,\alpha_x^2,\ q^2\right)^2 = \left((\alpha q)\,\alpha_x^2,\ Ff - \tfrac{1}{2}\,A_{ff}p^2\right)^2\\
&= F\cdot(\alpha q)(\alpha a)^2 - \tfrac{1}{2}\,A_{ff}\cdot(\alpha q)(\alpha p)^2 = F\cdot(pq) - \tfrac{1}{2}\,A_{ff}(sq)\,.
\end{aligned}
\]
In ebenso einfacher Weise würden sich die typischen Coefficienten
berechnen lassen, wenn wir $f$ und $\varphi$ durch irgend zwei andere
lineare Covarianten hätten darstellen wollen.

305. \textit{Ausnahmefälle für typische Darstellung.} Die typische
Darstellung wird unmöglich, sobald die beiden Invarianten

\[
\begin{aligned}
F &= (ap)^2\\
L &= (\Delta p)^2
\end{aligned}
\]
gleichzeitig identisch verschwinden; denn in diesem Falle verschwinden
auch alle Invarianten

\[
(pq),\ (pr),\ (ps),\ (qr),\ (qs),\ (rs)\,,
\]
wie ein Blick auf die Relationen (VI) Nr. 302 lehrt, und demnach sind die
vier linearen Covarianten nicht mehr von einander unabhängig.

Das Verschwinden von $F$ und $L$ sagt aber nichts anderes aus, als dass
$f$ und $\Delta$ den gemeinsamen Factor $p$ besitzen. In diesem Falle muss
aber $\varphi$ den linearen Factor $p$ quadratisch enthalten. Denn
multipliciren wir

\weq{1}{(ap)\,\Delta_x - (\Delta p)\,a_x = (a\Delta)\,p_x}
mit

\weq{2}{a_y\,\Delta_x + \Delta_y\,a_x = 2\,(a_x\,\Delta_x)_y\,,}
so kommt:

\[
\bigl((ap)\,\Delta_x - (Ap)\,a_x\bigr)\,\bigl(a_y\,\Delta_x + \Delta_y\,a_x\bigr) = 2\,\bigl((a\Delta)\,a_x\,\Delta_x\bigr)_y\,p_x\,. % [sic] \textquotedbl{}(Ap)\textquotedbl{} printed with roman A, not Delta, as elsewhere in this line
\]
Ersetzen wir aber hierin $y$ durch $p$, so erhalten wir

\[
(ap)^2\,\Delta_x^2 - (\Delta p)^2\,a_x^2 = 2\,p_x\,(\Theta,\ p)\,.
\]
Die linke Seite ist null wegen $F = L = 0$; also ist auch

\[
(\Theta,\ p) = (\alpha p)^2\,\alpha_x = 0\,,
\]
d.\,h. $\varphi$ besitzt den Factor $p$ im zweiten Grade.

Ebenso findet man umgekehrt, dass, wenn $\varphi$ einen linearen Factor
von $f$ quadratisch besitzt, $F$ und $L$ verschwinden müssen und sonach
$p$, $q$, $r$, $s$ einander proportional sind.


% ===== pdf idx 343, printed 329.0 =====
% !PRINTED 329
\weberpage{329}

$F$ und $L$ verschwinden auch, wenn $\Delta$ proportional $f$ ist, also
wenn $\Theta = 0$; dann kann aber überhaupt nicht mehr von einem simultanen
System von $f$ und $\varphi$ die Rede sein.

306. \textit{Resultante von f und $\varphi$.} Um die Resultante von $f$ und
$\varphi$ aufzustellen, benutzen wir die nämlichen Principien, wie wir sie
beim Aufsuchen der Discriminanten der Form fünfter und sechster Ordnung
verwendeten. Es möge $\lambda$ der gemeinsame Factor von $f$ und $\varphi$
sein. Diesen Factor besitzt alsdann insbesondere auch die erste
Ueberschiebung von $\varphi$ über $\lambda$. Sei $\mu$ der zweite lineare
Factor von $f$, $\nu$ der zweite von $(\varphi, \lambda)$, so können wir
unsere Bedingungen formuliren:
\[
a_x^2 = \lambda\mu, \quad \alpha_x^2(\alpha\lambda) = \lambda\nu.
\]
Versuchen wir auf Grund dieser Annahmen die Invarianten des Systemes zu
berechnen, so wird sich zwischen denselben eine Relation ergeben, die
gemäss der Voraussetzung und nach dem Grade in den Coefficienten nur die
Resultante von $f$ und $\varphi$ sein kann.

Wir berechnen zunächst die Invariante $F = (ap)^2$. Die Form $p$ war
dargestellt durch:
\[
p = (\alpha a)^2 \alpha_x, \quad \text{oder weil } a_x^2 = \lambda\mu,
\]
\[
p = (\alpha\lambda)(\alpha\mu)\alpha_x = ((\alpha\lambda)a_x^2, \mu) = (\lambda\nu, \mu).
\]
Nun ist:
\[
(\lambda\nu, \mu) = \frac{1}{2}\nu(\lambda\mu) + \frac{1}{2}\lambda(\nu\mu).
\]
Da aber
\[
\nu(\lambda\mu) = (\nu\mu)\lambda + (\lambda\nu)\mu,
\]
so kommt:
\[
(\lambda\nu, \mu) = \lambda(\nu\mu) + \frac{1}{2}\mu(\lambda\nu) = p.
\]
Bilden wir somit die zweite Ueberschiebung $(ap)^2 = F$, so erhalten wir
\[
F = \left(a_x^2, \ \lambda^2(\nu\mu)^2 + \lambda\cdot\mu(\nu\mu)(\lambda\nu) + \frac{1}{4}\mu^2(\lambda\nu^2)\right)^2,
\]
oder weil $a_x^2 = \lambda\cdot\mu$, und demnach das erste und dritte
Ueberschiebungsglied gleich null, so wird:
\weq{1}{F = (a_x^2, \ b_x^2(\nu\mu)(\lambda\nu))^2 = (\nu\mu)(\lambda\nu)\cdot A_{ff}.}
Berechnen wir sodann die simultane Invariante $A_{f\Delta}$, so findet man:
\[
A_{f\Delta} = (\Delta\lambda)(\Delta\mu) = (\alpha\beta)^2(\beta\lambda)(\alpha\mu).
\]
Der Ausdruck rechts ist nichts anderes als die zweite Ueberschiebung
\[
((\beta\lambda)\beta_x^2, \ (\alpha\mu)\alpha_x^2)^2 = (\lambda\nu, \ (\alpha\mu)\alpha_x^2)^2 = (\alpha\lambda)(\alpha\nu)(\alpha\mu).
\]


% ===== pdf idx 344, printed 330.0 =====
% !PRINTED 330
\weberpage{330}

Nun ist aber:
\[
(\alpha\lambda)(\alpha\nu)(\alpha\mu) = ((\alpha\lambda)a_x^2, \ \mu\nu)^2 = (\lambda\nu, \ \mu\nu) = \frac{1}{2}(\nu\mu)(\lambda\nu),
\]
somit:
\weq{2}{2A_{f\Delta} = (\nu\mu)(\lambda\nu).}
Setzen wir den Werth von $(\nu\mu)(\lambda\nu)$ in die Relation (1) ein, so
kommt:
\[
F - 2A_{f\Delta}\cdot A_{ff} = 0.
\]
Diese Relation besteht also unter der Voraussetzung, dass $f$ und $\varphi$
einen gemeinsamen Factor haben. Sie ist überdies vom Grade $2 + 3 = 5$ in
den Coefficienten, und ihre linke Seite darf somit als Resultante von $f$
und $\varphi$ betrachtet werden.

\begin{center}
\S.\,32. \\
\textbf{Das simultane System zweier cubischer Formen.}
\end{center}

307. \textit{Erste Formulirung der Aufgabe.} Das vollständige System einer
cubischen Form besteht aus den vier Formen:
\[
a_x^3 = f, \quad (f,f)^2 = \Delta, \quad (f,\Delta) = Q, \quad (\Delta,\Delta)^2 = A_{\Delta\Delta}.
\]
Es sei $\varphi$ eine zweite cubische Form; die Formen ihres vollen
Systemes wollen wir bezeichnen mit
% [note: printed capital letter here is a distinct script/Gothic form, not the roman letter V; rendered as \Gamma for lack of a closer available symbol]
\[
\alpha_x^3 = \varphi, \quad (\varphi,\varphi)^2 = \Gamma, \quad (\varphi,\Gamma) = K, \quad (\Gamma,\Gamma)^2 = A_{\Gamma\Gamma}.
\]
Wir stellen uns die Aufgabe, das simultane System der beiden Formen $f$ und
$\varphi$ zu ermitteln. Nach den allgemeinen Sätzen, § 38, erhalten wir
dasselbe, wenn wir die beiden Producte
\[
U = f^{\lambda_1}\cdot Q^{\lambda_2}\cdot \Delta^{\lambda_3}
\]
\[
V = \varphi^{\mu_1}\cdot K^{\mu_2}\cdot \Gamma^{\mu_3}
\]
über einander schieben, d.\,h. also mit einander multipliciren und das
Product $U\cdot V$ auf alle möglichen Arten in bekannter Weise falten.
Enthält eine solche Ueberschiebung
\[
F = (U,V)^{\varrho} = (f^{\lambda_1}\cdot Q^{\lambda_2}\cdot \Delta^{\lambda_3}, \ \varphi^{\mu_1}\cdot K^{\mu_2}\cdot \Gamma^{\mu_3})^{\varrho}
\]
kein zerfallendes Glied, so liefert sie eine Form des Systemes.

308. \textit{Reduction der Aufgabe.} Die Zahl der zu faltenden Producte
$U\cdot V$ lässt sich von vornherein reduciren, insofern sich bis auf
wenige Ausnahmen leicht erkennen lässt, welche derselben auf
Ueberschiebungen mit zerfallenden Gliedern führen.

Erstens: Producte, welche höhere als die zweite Potenz von $\Delta$ resp.
$\Gamma$ enthalten, bleiben unberücksichtigt, wegen der bekannten Relation:
\[
\Delta^3 + 2Q^2 + A_{\Delta\Delta}f^2 = 0.
\]
Zweitens: Enthält $U$ eine Form $\Delta$ und gleichzeitig $V$ eine Form
$\Gamma$, so sind höhere Faltungen als die zweite ausgeschlossen, da stets
Glieder mit dem Factor $A_{\Delta\Gamma} = (\Delta,\Gamma)^2$ gebildet
werden können.


% ===== pdf idx 345, printed 331.0 =====
% !PRINTED 331
\weberpage{331}

Drittens: Bei Ueberschiebungen, welche den Grad 3 übersteigen, kann $U$ nur
Einen Factor $f$ oder $Q$, und $V$ nur Einen Factor $\varphi$ oder $K$
enthalten, da diese Formen von gleichem Grade sind und somit stets Glieder
auftreten würden, welche die dritten Ueberschiebungen $(f,\varphi)^3$,
$(f,K)^3$, $(\varphi,Q)^3$ zum Factor haben. Bei niedrigeren
Ueberschiebungen ist aber ohnehin der Fall ausgeschlossen, dass $U$ oder
$V$ aus mehr als aus einem Factor besteht, wenn dieser dritten Grades ist.

Viertens: Ueberschiebungen einer Potenz von $f$, oder $Q$, über eine Potenz
von $\Gamma$, bleiben unberücksichtigt, da schon $(f^2,\Gamma^2)^4$, oder
$(Q^2,\Gamma^2)^4$ Glieder enthalten, welche Producte früherer
Ueberschiebungen sind, nämlich: $(f,\Gamma)^2\cdot(f,\Gamma)^2$ oder
$(Q,\Gamma)^2\cdot(Q,\Gamma)^2$.

Von allen Producten $U\cdot V$ sind also nur jene zu falten, in welchen\\
1) sowohl $U$ als $V$ nur \gesperrt{eine} Form enthält,\\
2) $U$ (oder $V$) eine Form, $V$ (oder $U$) eine zweite Potenz von
$\Gamma$ (resp. $\Delta$).

309. \textit{Aufstellung der einzelnen Formen.} Wir bilden nun die
einzelnen Ueberschiebungen:
\[
1) \quad U = f, \quad V = \varphi.
\]
Ihre Ueberschiebung liefert drei Formen:
\[
(f,\varphi) = \vartheta, \quad (f,\varphi)^2 = \Theta, \quad (f,\varphi)^3 = J.
\]
\[
2) \quad U = f, \quad V = \Gamma, \quad \text{oder} \quad U = \Delta, \quad V = \varphi.
\]
Wir erhalten sechs Formen:
\[
\begin{aligned}
&(f,\Gamma), \ (f,\Gamma)^2, \ (f,\Gamma^2)^3\\
&(\varphi,\Delta), \ (\varphi,\Delta)^2, \ (\varphi,\Gamma^2)^3.
\end{aligned}
\]
\[
3) \quad U = \Delta, \quad V = \Gamma.
\]
Es entstehen die zwei Formen:
\[
(\Delta,\Gamma) \quad \text{und} \quad (\Delta,\Gamma)^2.
\]
\[
4) \quad U = Q \quad \text{und} \quad V
\]
sei irgend eine Form des zweiten Systems. Als erste Ueberschiebungen
treten die Formen auf:
\[
(Q,\varphi), \ (Q,\Gamma); \ (Q,K),
\]
und entsprechend
\[
(K,f), \ (Q,\Delta);
\]
sie sind als Functionaldeterminanten von Functionaldeterminanten
reducibel.

Als zweite Ueberschiebungen erhält man:
\[
\begin{aligned}
&(Q,\varphi)^2, \ (Q,\Gamma)^2, \ (Q,K)^2\\
&(K,f)^2, \ (Q,\Delta)^2.
\end{aligned}
\]


% ===== pdf idx 346, printed 332.0 =====
% !PRINTED 332
\weberpage{332}

Von diesen ist $(Q,K)^2$ reducibel; denn falten wir das Product
\[
a_x^2\Delta_x(a\Delta)\cdot \alpha_x^2\Gamma_x(\alpha\Gamma)
\]
zweimal, und bilden insbesondere das Glied
\[
G = a_x\alpha_x(a\Gamma)(\alpha\Gamma)(a\Delta)(\alpha\Delta),
\]
% ?? uncertain: a small mark above "a" in "(a\Gamma)" in this line of the print may be ink noise rather than a diacritic; transcribed here without accent
so lässt sich auf dasselbe der Productsatz unmittelbar anwenden und man
findet
\[
G = \alpha_x\alpha_x(a\Gamma)(\alpha\Gamma)(a\Delta)(\alpha\Delta) % [sic]
\]
\[
= \frac{1}{2}a_x\alpha_x\{(\Gamma\alpha)^2(\Delta a)^2 + (\Gamma a)^2(\Delta\alpha)^2 - (\Gamma\Delta)^2(a\alpha)^2\}.
\]
Das erste Glied rechts ist wegen $(\Delta a)^2$ oder $(\Gamma\alpha)^2$
null. [Vergl. Nr. 145, (I).] Das zweite ist $(f,\Gamma)^2\cdot(\varphi,\Delta)^2$,
das dritte $(\Gamma,\Delta)^2\cdot(f,\varphi)^2$. Es ist somit das Glied
$G$ reducibel, folglich $(Q,K)^2$ selbst.

Als dritte Ueberschiebungen erhält man endlich:
\[
\begin{aligned}
&(Q,\varphi)^3, \ (Q,f)^3, \ (Q,K)^3\\
&(Q,\Gamma^2)^3, \ (K,\Gamma^2)^3.
\end{aligned}
\]
Von diesen Ueberschiebungen sind die beiden letzten reducibel. Es ist nur
nöthig, dies von $(K,\Delta^2)^3$ zu zeigen; für die Form $(Q,\Gamma^2)^3$
folgt die Reducibilität aus der entsprechenden Bedeutung der Formen $Q$
und $\Gamma$. Die dritte Polare von $K$ ist:
\[
K_{y^3} = (\alpha\Gamma)\alpha_y^2\Gamma_y,
\]
ferner die dritte Polare von $\Delta_x^4$
\[
\Delta_{y^3}^2 = \Delta_{y^2}\Delta_{1y}\Delta_{1x}.
\]
Indem wir das Product
\[
(\alpha\Gamma)\alpha_y^2\Gamma_y\cdot \Delta_{y^2}\Delta_{1y}\Delta_{1x}
\]
auf alle möglichen Arten dreimal falten, erhalten wir die Ueberschiebung
\[
(K,\Delta^2)^3 = \frac{1}{3}(\alpha\Gamma)(\alpha\Delta)^2(\Gamma\Delta_1)\Delta_{1x} + \frac{2}{3}(\alpha\Gamma)(\alpha\Delta_1)(\Gamma\Delta)(\alpha\Delta).
\]
Die Differenz dieser beiden Glieder ist aber null, da sie den Factor
$(\alpha\Gamma)^2$ enthält (vergl. Nr. 145), folglich sind beide einander
gleich, und wir erhalten:
\[
(K,\Delta^2)^3 = (\alpha\Gamma)(\alpha\Delta)^2(\Gamma\Delta_1)\Delta_{1x}.
\]
Ersetzen wir hier $\Delta_1$ durch $(ab)^2a_xb_x$, so kommt:
\[
(K,\Delta^2)^3 = (ab)^2(\alpha\Delta)^2(\alpha\Gamma)(\Gamma a)b_x
\]
\[
= (\alpha\Gamma)(\Gamma a)b_x\{(b\Delta)(a\alpha) - (a\Delta)(b\alpha)\}^2.
\]
Entwickelt man rechts nach dem binomischen Lehrsatz, so ist das erste und
letzte Glied null wegen des Factors $(b\Delta)^2$. Es reducirt sich also:
%CONT>>


% ===== pdf idx 347, printed 333.0 =====
% !PRINTED 333
\weberpage{333}

\[
(K, \Delta^2)^3 = -2(\alpha\Gamma)(\Gamma a)(a\alpha)(b\alpha)(a\Delta)(b\Delta)b_x
\]
\[
= -2(\alpha\Gamma)(a\Delta)(b\Delta)(a\alpha)b_x\{(\Gamma b)(a\alpha)-(ab)(\Gamma\alpha)\}
\]
\[
= -2(a\alpha)^2(\alpha\Gamma)(\Gamma b)(a\Delta)(b\Delta)b_x.
\]
% ?? uncertain: the symbol here transcribed as \Gamma is a stylised capital letter paired
% throughout this page with \Delta; the print quality does not let me distinguish it with
% certainty from a script V, but \Gamma is consistent with its parallel role to \Delta and
% with later occurrences on this page and the following pages.

Da rechts der Klammerfactor $(a\alpha)^2$ auftritt, so suchen wir diesen
Ausdruck mittelst der Reihenentwicklung
\[
(a\alpha)^2 a_y\alpha_x = \Theta_x\Theta_y + \frac{1}{2}J(yx)
\]
umzuformen. Ersetzt man nämlich in ihr $y$ durch $\Delta$, $x$ durch $\Gamma$,
so kommt:
\[
(a\alpha)^2(a\Delta)(\alpha\Gamma) = (\Theta\Gamma)(\Theta\Delta) + \frac{1}{2}J(\Delta\Gamma).
\]

Tragen wir dies in $(K, \Delta^2)^3$ ein, so erhalten wir:
\[
(K, \Delta^2)^3 = -2b_x(b\Delta)(\Gamma b)\{(\Theta\Gamma)(\Theta\Delta)+\frac{1}{2}J(\Delta\Gamma)\}
\]
\[
= +2b_x(b\Delta)(\Theta\Delta)(b\Gamma)(\Theta\Gamma) + J(\Delta\Gamma)(b\Delta)(b\Gamma)b_x.
\]
Für das erste Glied liefert der Productsatz:
\[
(b\Delta)(\Theta\Delta)(b\Gamma)(\Theta\Gamma) = \frac{1}{2}\{(b\Delta)^2(\Theta\Gamma)^2+(\Theta\Delta)^2(b\Gamma)^2-(\Delta\Gamma)^2(b\Theta)^2\}.
\]

Das erste Glied in dieser Summe ist null wegen $(b\Delta)^2$. Man erhält
also:
\[
(K, \Delta^2)^3 = (\Theta, \Delta)^2\cdot(f, \Gamma)^2 - (\Delta, \Gamma)^2\cdot(f, (f, \varphi))^2 + J\cdot(f, (\Delta, \Gamma))^2.
\]

Die beiden Covarianten $(K\Delta^2)^3$ und $(Q\Gamma^2)^3$ sind somit
reducibel und das ganze Formensystem der zwei cubischen Formen $f$ und
$\varphi$ besteht also aus 26 Formen.

Eine biquadratische:
\[
(f, \varphi) = (a\alpha)a_x^2\alpha_x^2 = \vartheta.
\]

Sechs cubische:
\[
f,\ \varphi,\ Q,\ K,\ (f, \Gamma),\ (\varphi, \Delta).
\]

Sechs quadratische:
\[
\Delta,\ \Gamma,\ (f, \varphi)^2 = \Theta,\ (\Delta, \Gamma),\ (Q, \varphi)^2,\ (K, f)^2.
\]

Sechs lineare:
\[
(f, \Gamma)^2,\ (\varphi, \Delta)^2,\ (f, \Gamma^2)^3,\ (\varphi, \Delta^2)^3,\ (Q, \Gamma)^2,\ (K, \Delta)^2.
\]

Sieben Invarianten:
\[
A_{\Delta\Delta},\ A_{\Gamma\Gamma},\ A_{\Delta\Gamma},\ (f, \varphi)^3 = J,\ (Q, K)^3,\ (Q, \varphi)^3,\ (Q, f)^3.
\]
Von den linearen Formen erhalten zwei eine besondere Bezeichnung.
Wir setzen
\[
(f, \Gamma)^2 = (a\Gamma)^2a_x = \pi_x = \pi
\]
\[
(\varphi, \Delta)^2 = (\alpha\Delta)^2\alpha_x = p_x = p.
\]


% ===== pdf idx 348, printed 334.0 =====
% !PRINTED 334
\weberpage{334}

310. \textit{Exposition der folgenden Aufgaben.} Nachdem wir im
Vorhergehenden die Formen ermittelt haben, aus denen das vollständige
System von $f$ und $\varphi$ bestehen muss, wird unser nächstes Ziel nach
den Beziehungen und Relationen gerichtet sein, welche, sei es zwischen
den Formen des Systemes selbst, sei es zwischen deren Ueberschiebungen
bestehen. Die allgemeinen Wege, welche dazu führen, sowie die hiebei
sich einstellende Nothwendigkeit der Einführung neuer gleich
berechtigter Formen an Stelle alter Formen, haben wir wiederholt an
den geeigneten Stellen klar gelegt. Wir greifen hier indes der
systematischen Methode vor, indem wir uns zunächst, in Folge der
Wichtigkeit der sich hiebei einstellenden Resultate, mit den linearen
Covarianten $p$ und $\pi$ beschäftigen, deren erste Ueberschiebung über
$f$ und $\varphi$ sowohl, als deren Quadrate wir vor allem berechnen
wollen. Hiebei ist es, wie überhaupt bei allen kommenden
Untersuchungen, stets hinreichend, die Berechnungen für eine je zweier
gleich berechtigter Formen $(f, p)$ und $(\varphi, \pi)$, oder $(f, \pi)$
und $(\varphi, p)$, oder $p^2$ und $\pi^2$ durchzuführen, da die
Resultate für die andere der beiden Formen einfach durch Vertauschung
von $a$ mit $\alpha$ gewonnen werden.

311. \textit{Die ersten Ueberschiebungen von $f$ und $\varphi$ über $p$
und $\pi$.} Die erste Ueberschiebung von $f$ über $p$ ist dargestellt
durch:
\[
(f, p) = a_x^2(a\alpha)(\alpha\Delta)^2.
\]
Subtrahiren wir auf der rechten Seite die Grösse
\[
0 = (a\alpha)(a\Delta)^2\alpha_x^2,
\]
so erhalten wir
\[
(f, p) = (a\alpha)\{a_x^2(\alpha\Delta)^2 - \alpha_x^2(a\Delta)^2\}
\]
\[
= (a\alpha)\{a_x(\alpha\Delta) + \alpha_x(a\Delta)\}\{a_x(\alpha\Delta) - \alpha_x(a\Delta)\}.
\]
Der letzte Factor rechts ist aber gleich $-(a\alpha)\Delta_x$ gemäss dem
Identitätssatze. Folglich ist:
\[
(f, p) = -(a\alpha)^2\Delta_x\{a_x(\alpha\Delta) + \alpha_x(a\Delta)\},
\]
oder:
\weq{1}{(f, p) = -2\big((a\alpha)^2a_x\alpha_x, \Delta\big) = -2(\Theta, \Delta).}
% ?? uncertain: the printed bracketing around ((a\alpha))^2 a_x\alpha_x, \Delta) is not
% fully legible; reproduced here in the sensible reading consistent with the equation
% being used as -2(\Theta,\Delta) two lines below and again in \S 311.
Also ebenso:
\weq{2}{(\varphi, \pi) = -2(\Theta, \Gamma).}
Es ist ferner:
\[
(f, \pi) = a_x^2(ab)(b\Gamma)^2.
\]
Vertauscht man wegen des Factors $(ab)$ das Symbol $a$ mit $b$ und
nimmt die halbe Summe, so kommt:
\[
(f, \pi) = \frac{1}{2}(ab)\{a_x^2(b\Gamma)^2 - b_x^2(a\Gamma)^2\}
\]
\[
= -\frac{1}{2}(ab)^2\Gamma_x\{a_x(b\Gamma) + b_x(a\Gamma)\},
\]
%CONT>>


% ===== pdf idx 349, printed 335.0 =====
% !PRINTED 335
\weberpage{335}
%<<CONT
oder:
\weq{3}{(f, \pi) = -\big((ab)^2a_xb_x, \Gamma\big) = -(\Delta, \Gamma).}
Also auch:
\weq{4}{(\varphi, p) = -(\Gamma, \Delta).}
Diese beiden Ueberschiebungen sind also bis aufs Vorzeichen einander
gleich, und man gewinnt daher aus (3) und (4) die
\gesperrt{Fundamentalrelation:}
\weq{5}{(f, \pi) + (\varphi, p) = 0.}

312. \textit{Eine Untersuchung über die Fundamentalrelation
$(f, \pi)+(\varphi, p)=0$.} Ehe wir dazu übergehen, die Quadrate von $p$
und $\pi$ zu berechnen, wollen wir uns die Frage vorlegen, ob nicht die
Formen $p$ und $\pi$ bereits durch die Relation
\weq{1}{(f, \pi) + (\varphi, p) = 0}
vollständig definirt sind. Die zu stellende Frage können wir in der
Form aussprechen: Für welche Werthe von $y$ und $z$ wird die Summe der
beiden Polaren
\[
f_y + \varphi_z
\]
identisch null? So lange $f$ und $\varphi$ Formen beliebiger Ordnung
sind, wird es nicht möglich sein, wenn nicht anderweitige Bedingungen
stattfinden, solche Werthe zu bestimmen. Nur in dem Falle, wo $f$ und
$\varphi$ quadratische Formen, giebt es unendlich viele Werthe. Bei
Formen dritter Ordnung findet man nur ein einziges Werthsystem, bei
biquadratischen überhaupt keines, wenn nicht eine bestimmte simultane
Invariante verschwindet. Dies lehrt schon eine einfache
Constantenabzählung.

Wenn wir nun hier für cubische Formen die Frage beantworten wollen, so
verlangt zunächst die Identität:
\[
f_y + \varphi_z = 0,
\]
da sie unabhängig von $x$ erfüllt sein soll, dass die Coefficienten
von $x_1^2, 2x_1x_2, x_2^2$ für sich allein verschwinden. Das liefert
drei Gleichungen ersten Grades für die Verhältnisse der vier Grössen
$y_1\ y_2\ z_1\ z_2$, nämlich:
\[
\bar a_0y_1 + \bar a_1y_2 + \bar\alpha_0z_1 + \bar\alpha_1z_2 = 0
\]
\[
\bar a_1y_1 + \bar a_2y_2 + \bar\alpha_1z_1 + \bar\alpha_2z_2 = 0
\]
\[
\bar a_2y_1 + \bar a_3y_2 + \bar\alpha_2z_1 + \bar\alpha_3z_2 = 0.
\]
Gemäss der Theorie der Auflösung linearer Gleichungen (vergl. Bd. I,
Nr. 111) sind also die Werthe von $y_1\ y_2\ z_1\ z_2$ proportional den
einzelnen Determinanten der Matrix
\[
V = \begin{vmatrix} \bar a_0 & \bar a_1 & \bar\alpha_0 & \bar\alpha_1 \\ \bar a_1 & \bar a_2 & \bar\alpha_1 & \bar\alpha_2 \\ \bar a_2 & \bar a_3 & \bar\alpha_2 & \bar\alpha_3 \end{vmatrix}.
\]


% ===== pdf idx 350, printed 336.0 =====
% !PRINTED 336
\weberpage{336}

Es ist:
\[
\varrho y_1 = \begin{vmatrix} \bar a_1 & \bar\alpha_0 & \bar\alpha_1 \\ \bar a_2 & \bar\alpha_1 & \bar\alpha_2 \\ \bar a_3 & \bar\alpha_2 & \bar\alpha_3 \end{vmatrix} = \begin{vmatrix} a_1^2a_2 & \alpha_1^3 & \beta_1^2\beta_2 \\ a_1a_2^2 & \alpha_1^2\alpha_2 & \beta_1\beta_2^2 \\ a_2^3 & \alpha_1\alpha_2^2 & \beta_2^3 \end{vmatrix} = a_2\alpha_1\beta_2(a\alpha)(a\beta)(\alpha\beta).
\]
Vertauscht man rechts $\alpha$ mit $\beta$ und nimmt die halbe Summe,
so kommt:
\[
\varrho y_1 = \frac{1}{2}a_2(\alpha\beta)^2(a\alpha)(a\beta)
\]
oder:
\[
\varrho y_1 = \frac{1}{2}a_2(a\Gamma)^2 = \frac{1}{2}\bar\pi_2.
\]
Ebenso:
\[
\varrho y_2 = -\frac{1}{2}a_1(a\Gamma)^2 = -\frac{1}{2}\bar\pi_1
\]
\[
\varrho z_1 = \frac{1}{2}\alpha_2(\alpha\Delta)^2 = \frac{1}{2}\bar p_2
\]
\[
\varrho z_2 = -\frac{1}{2}\alpha_1(\alpha\Delta)^2 = -\frac{1}{2}\bar p_1.
\]
% ?? uncertain: the overbars on \pi_1, \pi_2, p_1, p_2 in this determinant expansion are
% faint in the scan; reproduced as printed to the best of my reading.

Man erkennt, $p$ und $\pi$ sind die einzigen Werthe, für welche die
Identität (1) besteht, und sie sind demnach durch sie vollständig
definirt.

Bei quadratischen Formen wäre die Matrix $V$ nur zweizeilig gewesen;
zwei lineare homogene Gleichungen mit vier Unbekannten werden aber
durch unendlich viele Werthe befriedigt; bei biquadratischen
Gleichungen aber wäre die Matrix $V$ vierzeilig gewesen, und die
Bedingung, dass vier lineare homogene Gleichungen mit vier Unbekannten
durch ein bestimmtes Werthsystem derselben befriedigt werden können,
ist gerade durch das Verschwinden dieser Determinante $V$ ausgedrückt.

313. \textit{Berechnung von $p^2$ und $\pi^2$.} Unsere nächste Aufgabe
ist, $p^2$ als lineare Function quadratischer Covarianten
darzustellen. Es war
\[
p = (\beta\Delta)^2\beta_x = (\alpha\Delta)^2\alpha_x.
\]
Führen wir für $\Delta$ seinen Werth $(ab)^2a_xb_x$ ein, so wird:
\[
p = (ab)^2(a\beta)(b\beta)\beta_x.
\]
Multipliciren wir die beiden so erhaltenen Werthe für $p$, so kommt:
\[
p^2 = (\alpha\Delta)^2\alpha_x(ab)^2(a\beta)(b\beta)\beta_x.
\]
Hierin fordern die beiden quadratischen Klammerfactoren zur Benutzung
der Identität auf:
\[
(\alpha\Delta)^2(ab)^2 = \{(a\alpha)(b\Delta) - (b\alpha)(a\Delta)\}^2.
\]
Wegen der rechts auftretenden Factoren $(b\Delta)^2$, $(a\Delta)^2$
reducirt sich also der Werth von $p^2$ nach Einführung dieses Werthes
von
%CONT>>


% ===== pdf idx 351, printed 337.0 =====
% !PRINTED 337
\weberpage{337}
%<<CONT
$(a\Delta)^2(ab)^2$ auf:
\[
p^2 = -2(a\alpha)(b\alpha)(a\beta)(b\beta)(a\Delta)(b\Delta)\alpha_x\beta_x.
\]
Diesen Ausdruck wird man wegen der vier ersten Factoren mit Hilfe des Productsatzes umformen. Man erhält:
\[
p^2 = -(a\Delta)(b\Delta)\alpha_x\beta_x\left\{(a\alpha)^2(b\beta)^2+(a\beta)^2(b\alpha)^2-(ab)^2(\alpha\beta)^2\right\}.
\]
Die beiden ersten Glieder rechts gehen durch Vertauschung von $a$ mit $b$ in einander über, sind also einander gleich; daher:
\weq{6}{p^2 = (ab)^2(\alpha\beta)^2(a\Delta)(b\Delta)\alpha_x\beta_x - 2(a\alpha)^2(b\beta)^2(a\Delta)(b\Delta)\alpha_x\beta_x.}
Der erste Term rechts ist nichts anderes als das Product
\weq{7}{(ab)^2(a\Delta)(b\Delta)(\alpha\beta)^2\alpha_x\beta_x = A_{\Delta\Delta}\cdot \nabla.}
Um den zweiten Term umzuformen, benutzen wir wegen des Factors $(a\alpha)^2$ die Reihenentwicklung:
\weq{8}{(a\alpha)^2 a_y \alpha_x = \Theta_x\Theta_y + \frac{1}{2}J(yx).}
Denken wir uns in ihr $y$ durch $\Delta$ ersetzt und alsdann den Werth von $(a\alpha)^2(a\Delta)\alpha_x$ in den zweiten Term (6) eingetragen, so wird derselbe:
\[
-2(a\alpha)^2(a\Delta)\alpha_x\cdot(b\beta)^2(b\Delta)\beta_x = -2(b\beta)^2(b\Delta)\beta_x\left\{\Theta_x(\Theta\Delta)+\frac{1}{2}J\Delta_x\right\}.
\]
Transformiren wir die Factoren vor der Klammer nochmals durch dieselbe Reihenentwicklung (8), so kommt:
\weq{9}{\begin{aligned}
&-2(a\alpha)^2(a\Delta)\alpha_x\cdot(b\beta)^2(b\Delta)\beta_x\\
&= -2\Theta_x(\Theta\Delta)\cdot\left\{\Theta'_x(\Theta'\Delta)+\frac{1}{2}J\Delta_x\right\} - J\Delta_x\left\{\Theta_x(\Theta\Delta)+\frac{1}{2}J\Delta_x\right\}\\
&= -2(\Theta\Delta)\Theta_x(\Theta'\Delta)\Theta'_x - 2J\Delta_x(\Theta\Delta)\Theta_x - \frac{J^2}{2}\Delta\\
&= -2(\Theta\Delta)(\Theta'\Delta)\Theta_x\Theta'_x - 2J(\Theta,\Delta) - \frac{J^2}{2}\Delta.
\end{aligned}}
Das erste Glied rechts ist aber nach dem Productsatze:
\weq{10}{\begin{aligned}
2(\Theta\Delta)(\Theta'\Delta)\Theta_x\Theta'_x &= \left\{(\Theta\Delta)^2\Theta' + (\Theta'\Delta)^2\Theta - (\Theta\Theta')^2\Delta\right\}\\
&= 2A_{\Theta\Delta}\cdot\Theta - A_{\Theta\Theta}\cdot\Delta.
\end{aligned}}
Es wird also, wenn wir die in (7), (9) und (10) berechneten Werthe in (6) eintragen und überdies $-2(\Theta,\Delta)$ durch $(f,p)$ ersetzen [vergl. Nr. 310, (1)]
\weq{11}{p^2 = A_{\Delta\Delta}\nabla - 2A_{\Theta\Delta}\Theta + \left(A_{\Theta\Theta}-\frac{J^2}{2}\right)\Delta + J(f,p),}
und demgemäss, da $J$ bei Vertauschung von $a$ mit $\alpha$ das Zeichen wechselt:
\weq{12}{\pi^2 = A_{\nabla\nabla}\Delta - 2A_{\Theta\nabla}\Theta + \left(A_{\Theta\Theta}-\frac{J^2}{2}\right)\nabla - J(\varphi,\pi).}


% ===== pdf idx 352, printed 338.0 =====
% !PRINTED 338
\weberpage{338}
314. \textit{Exposition der folgenden Aufgaben.} Die vorausgegangenen Abschnitte dienen dazu, uns über einige Eigenschaften jener beiden linearen Covarianten $p$ und $\pi$ zu instruiren, die im ganzen System eine so hervorragende Rolle spielen. Wir können nun wiederum zu einer systematischen Methode zurückkehren, die uns über die Beziehungen der Formen des Systemes Aufschluss verschaffen soll. Indem wir diese Relationen zu ermitteln streben, werden wir uns gleichzeitig veranlasst sehen, gewisse Ueberschiebungen des Systemes durch ein geeignetes ihrer Glieder zu ersetzen. Wir beginnen damit, die Ueberschiebungen der einzelnen Covariantengruppen des Systemes der Reihe nach zu studiren, so weit sie für uns Interesse haben. Daran schliessen sich die Untersuchungen über die Relationen zwischen den Invarianten des Systemes, deren nothwendig zwei existiren müssen, da zwei cubische Formen nur $8-3=5$ unabhängige Constante besitzen, während doch 7 Invarianten im Systeme auftraten.

315. \textit{Die biquadratische Covariante $\vartheta$.} Die einzige Form vierter Ordnung, die im simultanen System von $f$ und $\varphi$ auftritt, ist die schiefe Covariante: $\vartheta = (f,\varphi) = (a\alpha)a_x^2\alpha_x^2$.

Von den Ueberschiebungen dieser Form über andere besprechen wir hier vier, die ein allgemeineres Interesse haben, nämlich die dritte Ueberschiebung über $f$ und $\varphi$, die erste über $p$ und $\pi$. Es ist:
\[
(\vartheta,f)^3 = \left((a\alpha)a_x^2\alpha_x^2,\, b_x^3\right)^3 = \frac{1}{2}(a\alpha)\left\{(ab)^2(\alpha b)\alpha_x + (\alpha b)^2(ab)a_x\right\}
\]
\[
= \frac{1}{2}(a\alpha)(ab)(\alpha b)\left\{(ab)\alpha_x + (\alpha b)a_x\right\}.
\]
Das letzte Glied $(\alpha b)a_x$ in der Klammer kann durch $\frac{1}{2}(ab)\alpha_x$ ersetzt werden, da die Identität besteht:
\[
(\alpha b)a_x + (a\alpha)b_x = (ab)\alpha_x
\]
und $(\alpha b)a_x$ hier mit $(a\alpha)b_x$ gleichbedeutend ist. Man erhält also:
\[
(\vartheta,f)^3 = \frac{3}{4}(a\alpha)(ab)^2(\alpha b)\alpha_x,
\]
oder, da $\Delta = (ab)^2a_xb_x$:
\weq{1}{(\vartheta,f)^3 = \frac{3}{4}(\Delta\alpha)(\alpha\Delta)\alpha_x = -\frac{3}{4}(\varphi,\Delta)^2 = -\frac{3}{4}p.}
Folglich auch, da $\vartheta$ durch Vertauschung von $a$ mit $\alpha$ sein Zeichen ändert:
\[
(\vartheta,\varphi)^3 = +\frac{3}{4}\pi.
\]
Zur Berechnung der ersten Ueberschiebung von $\vartheta$ über $p$ und $\pi$ gehen wir aus von der Fundamentalrelation Nr. 310, (5)
\[
a_x^2(a\pi) + \alpha_x^2(\alpha p) = 0.
\]


% ===== pdf idx 353, printed 339.0 =====
% !PRINTED 339
\weberpage{339}
Schieben wir diese Relation einmal über $f$, so kommt:
\weq{2}{0 = (ab)(a\pi)a_xb_x^2 + (\alpha b)(\alpha p)\alpha_xb_x^2.}
Das zweite Glied geht aus der Reihe
\[
(a\alpha)a_x^2\alpha_x\alpha_y = \vartheta_y - \frac{1}{2}\Theta(yx)
\]
hervor, wenn wir in ihr $a$ durch $b$ und $y$ durch $p$ ersetzen; es folgt dann
\weq{3}{(\alpha b)(\alpha p)a_xb_x^2 = -(\vartheta,p) + \frac{1}{2}\Theta p.}
Vertauschen wir im ersten Glied der Relation (2) $a$ mit $b$ und nehmen die halbe Summe, so erhalten wir:
\weq{4}{(ab)(a\pi)a_xb_x^2 = \frac{1}{2}(ab)a_xb_x\left\{(a\pi)b_x - (b\pi)a_x\right\} = \frac{1}{2}(ab)^2a_xb_x\cdot\pi.}
Tragen wir diese Werthe der beiden Glieder aus (3) und (4) in (2) ein, so erhalten wir:
\[
0 = \frac{\pi\Delta}{2} - (\vartheta,p) + \frac{1}{2}\Theta p,
\]
oder:
\weq{5}{(\vartheta,p) = \frac{1}{2}\left\{\Delta\pi + \Theta p\right\},}
und demgemäss:
\weq{6}{(\vartheta,\pi) = -\frac{1}{2}\left\{\nabla p + \Theta \pi\right\}.}
316. \textit{Die cubischen und quadratischen Covarianten.} Die Formen dritter Ordnung, welche das simultane System aufweist, wollen wir genau so beibehalten, wie sie der Faltungsprocess bei Aufstellung des Systemes geliefert hat. Es sind sechs an der Zahl:
\[
f,\ \varphi,\ Q,\ K,\ (f,\nabla),\ (\varphi,\Delta).
\]
Zu weitern Untersuchungen von Ueberschiebungen derselben über einander oder über andere Formen haben wir hier keine Veranlassung.

Von den sechs quadratischen Formen, die im Systeme auftreten, wollen wir drei durch andere ersetzen, nämlich die drei Formen $(\Delta,\nabla)$, $(Q,\varphi)^2$ und $(K,f)^2$. Die Gründe, die uns hiezu bewegen, sind rein praktischer Natur. \gesperrt{Clebsch} hatte dieselben ersetzt durch die drei Functionaldeterminanten der drei andern quadratischen Formen $\Delta$, $\nabla$, $\Theta$. Wir wollen an ihrer Stelle einführen: $(f,\pi)$, $(f,p)$, $(\varphi,\pi)$, also Ueberschiebungen über lineare Formen. Dass wir $(\Delta,\nabla)$ durch $(f,\pi)$ ersetzen können, geht aus Nr. 311, (3) hervor.

Berechnet man ferner die zweite Ueberschiebung von $Q$ über $\varphi$, so findet man Folgendes. Die zweite Polare von $Q$ ist [vgl. Nr. 146, (2)]:
\[
Q_xQ_y^2 = (a\Delta)a_y^2\Delta_x,
\]


% ===== pdf idx 354, printed 340.0 =====
% !PRINTED 340
\weberpage{340}
Ersetzen wir darin $y$ durch $\alpha$, und multipliciren mit $\alpha_x$, so kommt:
\[
(Q,\varphi)^2 = (a\Delta)(a\alpha)^2\Delta_x\alpha_x.
\]
Der Ausdruck rechts entsteht aber aus $(a\alpha)^2a_y\alpha_x$, wenn wir darin $y$ durch $\Delta$ ersetzen und mit $\Delta_x$ multipliciren. Wir erhalten [vergl. Nr. 313, (8)] daher:
\[
(Q,\varphi)^2 = (\Theta,\Delta) + \frac{J}{2}\Delta,
\]
also wegen Nr. 311, (1)
\weq{1}{(Q,\varphi)^2 = -\frac{1}{2}(f,p) + \frac{1}{2}J\Delta.}
Man darf also in der That $(Q,\varphi)^2$ durch $(f,p)$ ersetzen und demgemäss auch $(K,f)^2$ durch $(\varphi,\pi)$. Die sechs quadratischen Formen sind also:
\[
\Delta,\ \Theta,\ \nabla,\ (f,p),\ (\varphi,\pi),\ (f,\pi) = -(\varphi,p).
\]
317. \textit{Die sechs linearen Formen.} Von den sechs linearen Formen des Systemes behalten wir die beiden Formen $p$ und $\pi$ im Folgenden bei. Dagegen ersetzen wir die vier übrigen durch
\[
(\Delta,p),\ (\Delta,\pi),\ (\nabla,p),\ (\nabla,\pi).
\]
Dass dies gestattet ist, kann man leicht zeigen. Denn da die zweite Polare von $Q$, wie schon erwähnt, sich auf $(a\Delta)a_y^2\Delta_x$ reducirt, so wird
\[
(Q,\nabla)^2 = (a\Delta)(a\nabla)^2\Delta_x.
\]
Nun ist aber die lineare Form
\[
\pi_x = (a\nabla)^2a_x,
\]
also
\weq{1}{(Q,\nabla)^2 = \left((a\nabla)^2a_x,\, \Delta\right) = (\pi,\Delta)}
und ebenso
\weq{2}{(K,\Delta)^2 = \left((\alpha\Delta)^2\alpha_x,\, \nabla\right) = (p,\nabla).}
Man erhält ferner:
\[
(\varphi,\Delta^2)^3 = (\alpha\Delta)^2(\alpha\Delta_1)\Delta_{1x},
\]
oder, weil:
\[
(\alpha\Delta)^2\alpha_x = p,
\]
\weq{3}{(\varphi,\Delta^2)^3 = \left((\alpha\Delta)^2\alpha_x,\, \Delta_1\right) = (p,\Delta),}
und ebenso
\weq{4}{(f,\nabla^2)^3 = \left((a\nabla)^2a_x,\, \nabla_1\right) = (\pi,\nabla).}
Die sechs linearen Covarianten des Systemes sind also:
\[
p,\ \pi,\ (\Delta,p),\ (\Delta,\pi),\ (\nabla,p),\ (\nabla,\pi).
\]
318. \textit{Weitere lineare Covarianten.} Den eben besprochenen sechs Formen, insbesondere den beiden Ueberschiebungen von $f$ und $\varphi$ über
%CONT>>


% ===== pdf idx 355, printed 341.0 =====
% !PRINTED 341
\weberpage{341}
$V$, resp. $\Delta$, sehr nahe liegend sind auch die beiden Ueberschiebungen
von $f$ und $\varphi$ über die andere quadratische Form $\Theta$, nämlich
$(f, \Theta)^2$ und $(\varphi, \Theta)^2$. Wir erhalten:
\[
\begin{aligned}
(\Theta, f)^2 &= (a\alpha)^2(ab)(\alpha b)\,b_x\\
&= \frac{1}{2}\{(a\alpha)^2(ab)(\alpha b)\,b_x + (b\alpha)^2(ba)(\alpha a)\,a_x\}\\
&= \frac{1}{2}(a\alpha)(ab)(\alpha b)\{(a\alpha)\,b_x + (\alpha b)\,a_x\}\\
&= \frac{1}{2}(a\alpha)(ab)^2(\alpha b)\,\alpha_x = -\frac{1}{2}(a\alpha)(b\alpha)(ab)^2\alpha_x,
\end{aligned}
\]
oder, weil $(ab)^2a_xb_x = \Delta$ und nach Nr.\ 309, $(\varphi, \Delta)^2 = p$,
\weq{1}{(f, \Theta)^2 = -\frac{1}{2}(\Delta, \varphi)^2 = -\frac{1}{2}\,p,}
also auch
\weq{2}{(\varphi, \Theta)^2 = -\frac{1}{2}(V, f)^2 = -\frac{1}{2}\,\pi.}

Wir haben damit auch für die beiden wichtigen linearen Covarianten $p$ und
$\pi$ einen neuen Ausdruck gefunden.

Eine zweite Reihe einfach zu berechnender verwandter linearer Covarianten
ergiebt sich aus den Ueberschiebungen von $f$ und $\varphi$ über Producte
und Potenzen von $p$ und $\pi$, also die Covarianten:
\[
(f, p^2)^2, \ (f, p\pi)^2, \ (f, \pi^2)^2; \quad (\varphi, p^2), \ (\varphi, p\pi)^2, \ (\varphi, \pi^2)^2. % [sic] first term lacks outer exponent 2 as printed
\]

Da wir die Werthe von $p^2$ und $\pi^2$ bereits als lineare Functionen
quadratischer Formen ermittelt haben, so hat die Berechnung dieser
Covarianten keine weitere Schwierigkeit. Wir schieben die Relation
\[
p^2 = A_{\Delta\Delta}\,V - 2A_{\Theta\Delta}\,\Theta + \left(A_{\Theta\Theta} - \frac{J^2}{2}\right)\Delta + J(f, p)
\]
zweimal über $f$ und erhalten:
\[
\begin{aligned}
(f, p^2)^2 = (ap)^2a_x &= A_{\Delta\Delta}(f,V)^2 - 2A_{\Theta\Delta}(f,\Theta)^2 + \left(A_{\Theta\Theta} - \frac{J^2}{2}\right)(f,\Delta)^2\\
&\quad + J(f, (ap)a_x^2)^2\\
&= A_{\Delta\Delta}\pi + A_{\Theta\Delta}p + J(ab)^2(ap)b_x,
\end{aligned}
\]
weil ja die Ueberschiebung $(f, \Delta)^2 = 0$ ist. Es ist aber
\[
(ab)^2(ap)b_x = (\Delta, p),
\]
somit
\weq{3}{(f, p^2)^2 = A_{\Delta\Delta}\pi + A_{\Theta\Delta}p + J(\Delta, p),}
und ebenso
\weq{4}{(\varphi, \pi^2)^2 = A_{VV}p + A_{\Theta V}\pi - J(V, \pi).}

Ferner ergiebt sich:
\[
\begin{aligned}
(\varphi, p^2)^2 = (\alpha p)^2\alpha_x &= A_{\Delta\Delta}(\varphi,V)^2 - 2A_{\Theta\Delta}(\varphi,\Theta)^2 + \left(A_{\Theta\Theta} - \frac{J^2}{2}\right)(\varphi,\Delta)^2\\
&\quad + J(\varphi, (ap)a_x^2)^2.
\end{aligned}
\]


% ===== pdf idx 356, printed 342.0 =====
% !PRINTED 342
\weberpage{342}
In dieser Relation ist einestheils $(\varphi, V)^2 = 0$, und anderntheils:
\[
\begin{aligned}
(\varphi, (ap)a_x^2)^2 = (ap)(\alpha a)^2\alpha_x &= (a\alpha)^2\{a_x(p\alpha) + p_x(a\alpha)\}\\
&= (a\alpha)^3p + (a\alpha)^2(p\alpha)a_x.
\end{aligned}
\]
Da aber wegen der Fundamentalrelation $(f, \pi) + (\varphi, p) = 0$:
\[
(\alpha a)^2(p\alpha)a_x = ((p\alpha)\alpha_x^2, a_x^3)^2 = ((b\pi)b_x^2, a_x^3)^2 = (ab)^2(b\pi)a_x = (\Delta,\pi),
\]
so erhalten wir:
\[
(\varphi, (ap)a_x^2)^2 = Jp + (\Delta,\pi).
\]

Es wird sonach, unter Berücksichtigung dieser Werthe der einzelnen Glieder,
und wegen Nr.\ 318, (2):
\weq{5}{(\varphi, p^2)^2 = A_{\Theta\Delta}\pi + \left(A_{\Theta\Theta} + \frac{J^2}{2}\right)p + J(\Delta,\pi),}
also auch
\weq{6}{(f, \pi^2)^2 = A_{\Theta V}p + \left(A_{\Theta\Theta} + \frac{J^2}{2}\right)\pi - J(V,p).}

Nun kann man die Ueberschiebung $(\varphi, p^2)^2$ auch schreiben:
\[
((\varphi, p), p);
\]
weil aber $(\varphi, p) = -(f, \pi)$, so ist auch
\[
(\varphi, p^2)^2 = -((f, \pi), p)^2 = -(f, \pi\cdot p)^2,
\]
und ebenso
\[
(f, \pi^2)^2 = -((\varphi, p), \pi)^2 = -(\varphi, p\cdot\pi)^2.
\]

Diese beiden letzten linearen Covarianten haben sonach dieselben Werthe,
die wir bereits in (5) und (6) berechnet haben.

319. \textit{Die Invarianten des simultanen Systems.} Von den sieben
Invarianten, auf welche wir durch Aufstellung des Systemes geführt wurden,
behalten wir folgende vier bei: $J$, $A_{\Delta\Delta}$, $A_{VV}$,
$A_{\Delta V}$. Die drei übrigen: $(f, K)^3$, $(\varphi, Q)^3$, $(Q, K)^3$
ersetzen wir durch andere, indem wir die übrigen Invarianten des
simultanen Systemes $\Delta$, $V$, $\Theta$ zu Hilfe nehmen. Es ist
nämlich:
\weq{1}{(f, K)^3 = (\alpha V)(a\alpha)^2(aV) = ((a\alpha)^2a_x\alpha_x, V)^2 = (\Theta, V)^2 = A_{\Theta V},}
und ebenso
\[
(\varphi, Q)^3 = A_{\Theta\Delta}.
\]
Ferner ist ein Glied der Ueberschiebung $(Q, K)^3$ dargestellt durch:
\[
G = (a\Delta)(\alpha V)(a\alpha)^2(\Delta V).
\]
Die rechte Seite geht aus der schon mehrmals benutzten
Reihenentwicklung hervor:
\[
(a\alpha)^2a_y\alpha_x = \Theta_x\Theta_y + \frac{1}{2}J(yx),
\]
wenn wir in ihr $y$ durch $\Delta$, $x$ durch $V$ ersetzen und mit
$(\Delta V)$ multipliciren. Wir erhalten:
%CONT>>


% ===== pdf idx 357, printed 343.0 =====
% !PRINTED 343
\weberpage{343}
%<<CONT
\[
\begin{aligned}
G = (\Delta V)\left\{(\Theta V)(\Theta\Delta) + \frac{1}{2}J(\Delta V)\right\}
&= (\Delta V)(\Theta V)(\Theta\Delta) + \frac{1}{2}J A_{\Delta V}.
\end{aligned}
\]
Man kann somit $(Q, K)^3$ ersetzen durch das erste Glied rechts
\weq{2}{\Omega = (\Delta V)(\Theta V)(\Theta\Delta),}
welches nichts anderes ist, als die simultane Invariante der drei
quadratischen Formen $\Delta$, $V$ und $\Theta$.

Diese Form $\Omega$ lässt sich auch noch auf eine andere Weise
darstellen. Wir hatten gefunden [vergl. Nr.\ 311, (1)]
\[
(f, p) = (ap)a_x^2 = 2(\Delta, \Theta) = 2(\Delta\Theta)\Delta_x\Theta_x.
\]
Aus dem Ausdrucke rechts geht aber $\Omega$ hervor, wenn man diese
Relation zweimal über $V$ schiebt; man erhält also:
\[
((ap)a_x^2, V)^2 = 2((\Delta\Theta)\Delta_x\Theta_x, V)^2,
\]
oder
\[
(aV)^2(ap) = -2\Omega.
\]
Es ist aber
\[
(aV)^2(ap) = ((aV)^2a_x, p) = (\pi, p);
\]
also ist
\[
2\Omega = (p, \pi).
\]

Es erübrigt noch, den Werth der nicht zum System gehörigen Invariante
$A_{\Theta\Theta}$ zu bestimmen. Zu dem Zwecke ersetzen wir in der
Reihenentwicklung
\[
(a\alpha)^2a_y\alpha_x = \Theta_x\Theta_y + \frac{1}{2}J(yx)
\]
$y$ durch $b$ und $x$ durch $\beta$ und multipliciren die Entwicklung
mit $(b\beta)^2$; dann erhalten wir:
\weq{3}{
\begin{aligned}
(a\alpha)^2(ab)(\alpha\beta)(b\beta)^2 &= (b\beta)^2\left\{(\Theta\beta)(\Theta b) + \frac{J}{2}(b\beta)\right\}\\
&= \left(\Theta, (b\beta)^2b_x\beta_x\right) + \frac{J^2}{2}\\
&= A_{\Theta\Theta} + \frac{J^2}{2}.
\end{aligned}
}
Um den Ausdruck links umzuformen beachte man, dass wegen des Factors
$(ab)$ es vortheilhaft ist, $a$ mit $b$ zu vertauschen und die halbe
Summe zu nehmen; dann wird:
\[
\begin{aligned}
(a\alpha)^2(b\beta)^2(ab)(\alpha\beta) &= \frac{1}{2}(ab)(\alpha\beta)\{(a\alpha)^2(b\beta)^2 - (b\alpha)^2(a\beta)^2\}\\
&= \frac{1}{2}(ab)(\alpha\beta)\{(a\alpha)(b\beta) + (b\alpha)(a\beta)\}\\
&\quad \{(a\alpha)(b\beta) - (b\alpha)(a\beta)\}\\
&= \frac{1}{2}(ab)^2(\alpha\beta)^2\{(a\alpha)(b\beta) + (b\alpha)(a\beta)\}.
\end{aligned}
\]


% ===== pdf idx 358, printed 344.0 =====
% !PRINTED 344
\weberpage{344}
Die beiden Terme rechts stimmen überein; also ist:
\weq{4}{(a\alpha)^2(b\beta)^2(ab)(\alpha\beta) = (ab)^2(\alpha\beta)^2(a\alpha)(b\beta) = (\Delta V)^2 = A_{\Delta V}.}
Substituirt man daher (4) in (3), so erhält man:
\weq{5}{A_{\Theta\Theta} = A_{\Delta V} - \frac{J^2}{2}.}

Fassen wir wiederum die sieben Invarianten des Systemes zusammen, so
sind dieselben repräsentirt durch:
\[
A_{\Delta\Delta}, \ A_{VV}, \ J, \ A_{\Delta V}, \ A_{\Delta\Theta}, \ A_{V\Theta}, \ \Omega = R_{\Delta V\Theta}.
\]
Zwischen diesen sieben Invarianten müssen, wie schon früher erwähnt,
zwei Relationen bestehen. Wir wollen sie im Folgenden berechnen.

320. \textit{Invariantenrelation für $\Omega^2$.} Zunächst weiss man aus
der Theorie der quadratischen Formen, dass das Quadrat der schiefen
Invariante $\Omega$ sich darstellen lässt durch [vgl. Nr.\ 134, (2)
$\beta$)] die Determinante:
\[
2\Omega^2 = \begin{vmatrix}
A_{\Delta\Delta}, & A_{\Theta\Delta}, & A_{V\Delta} \\
A_{\Delta\Theta}, & A_{\Theta\Theta}, & A_{V\Theta} \\
A_{\Delta V}, & A_{\Theta V}, & A_{VV}
\end{vmatrix}.
\]

Wir bezeichnen diese Determinante mit $U$ und ihre Unterdeterminanten mit
$U_{ik}$. Diese Unterdeterminanten $U_{ik}$ sind nach der Theorie der
quadratischen Formen nichts anderes als die zweiten Ueberschiebungen von
Functionaldeterminanten der drei Formen $\Delta$, $V$, $\Theta$. Denn
ersetzen wir die dort [vgl. Nr.\ 134, (1) $\beta$) und $\gamma$)] mit
$f$, $\varphi$, $\psi$ bezeichneten quadratischen Formen durch $\Delta$,
resp. $\Theta$, $V$ und demnach $\vartheta_1 = (\varphi,\psi)$,
$\vartheta_2 = (\psi,f)$, $\vartheta_3 = (f,\varphi)$ durch $(\Theta,V)$,
resp. $(V,\Delta)$, $(\Delta\Theta)$, so erhalten wir unmittelbar die
folgenden sechs Relationen:
\[
\begin{aligned}
2A_{\vartheta_1\vartheta_1} &= 2((\Theta,V),(\Theta,V))^2 = A_{\Theta\Theta}A_{VV} - A_{V\Theta}^2 = +U_{11}\\
2A_{\vartheta_2\vartheta_2} &= 2((V,\Delta),(V,\Delta))^2 = A_{VV}A_{\Delta\Delta} - A_{V\Delta}^2 = +U_{22}\\
2A_{\vartheta_3\vartheta_3} &= 2((\Delta,\Theta),(\Delta,\Theta))^2 = A_{\Delta\Delta}A_{\Theta\Theta} - A_{\Delta\Theta}^2 = +U_{33}\\
2A_{\vartheta_1\vartheta_2} &= 2((\Theta,V),(V,\Delta))^2 = A_{V\Delta}A_{V\Theta} - A_{VV}A_{\Delta\Theta} = -U_{12}\\
2A_{\vartheta_1\vartheta_3} &= 2((\Theta,V),(\Delta,\Theta))^2 = A_{\Theta\Delta}A_{V\Theta} - A_{\Theta\Theta}A_{\Delta V} = +U_{13}\\
2A_{\vartheta_2\vartheta_3} &= 2((V,\Delta),(\Delta,\Theta))^2 = A_{\Delta\Theta}A_{\Delta V} - A_{\Delta\Delta}A_{\Theta V} = -U_{23}.
\end{aligned}
\]

Diese Unterdeterminanten spielen bei der typischen Darstellung von
$\Delta$, $V$ und $\Theta$ eine besondere Rolle. Denn sie sind es, welche
die typischen Coefficienten liefern. Sie dienen ferner dazu, die
Relation, welche das Product $\Omega J$ mit den andern fundamentalen
Invarianten bildet, in einfacher Weise darzustellen.


% ===== pdf idx 359, printed 345.0 =====
% !PRINTED 345
\weberpage{345}

321. \textit{Invariantenrelation für $\Omega\cdot J$.} Man hat für das
Product dieser beiden Invarianten den symbolischen Ausdruck:
\[
2J\cdot\Omega=(a\alpha)^3\cdot(p\pi),
\]
oder, indem wir den Identitätssatz benutzen:
\[
\begin{aligned}
2J\cdot\Omega&=(a\alpha)^2\{(ap)(\alpha\pi)-(a\pi)(\alpha p)\}\\
&=(a\alpha)^2(ap)(\alpha\pi)-(a\alpha)^2(a\pi)(\alpha p)\\
&=\mathbf{(}(ap)a_x^2,(\alpha\pi)\alpha_x^2\mathbf{)}^2-\mathbf{(}(a\pi)a_x^2,(\alpha p)\alpha_x^2\mathbf{)}^2\\
&=((f,p),(\varphi,\pi))^2-((f,\pi),(\varphi,p))^2.
\end{aligned}
\]

Nun bestehen aber nach dem Früheren die Beziehungen:
\[
\left.\begin{array}{l}
(f,\pi)=-(\Delta,V)\\
(\varphi,p)=-(V,\Delta)\\
(f,p)=-2(\Theta,\Delta)\\
(\varphi,\pi)=-2(\Theta,V)
\end{array}\right\}\qquad\text{(I)}
\]
Demnach wird die Relation für $2J\cdot\Omega$:
\[
\begin{aligned}
2J\cdot\Omega&=4((\Theta,\Delta),(\Theta,V))^2-((\Delta,V),(V,\Delta))^2\\
&=((V,\Delta),(V,\Delta))^2-4((\Theta,V),(\Delta,\Theta))^2,
\end{aligned}
\]
oder:
\[
4J\cdot\Omega=U_{22}-4U_{13}.
\]

Hiemit seien die Untersuchungen über die Relationen der Formen des
Systems von $f$ und $\varphi$ beschlossen.

322. \textit{Zusammenfassung.} Ehe wir zu Anwendungen der bisherigen
theoretischen Untersuchungen übergehen, dürfte es angezeigt sein, eine
tabellarische Uebersicht über die gewonnenen Resultate zu entwerfen.

Die 26 Formen des simultanen Systemes von $f$ und $\varphi$ sind durch
folgende Symbole dargestellt:
\[
\vartheta=(a\alpha)a_x^2\alpha_x^2,
\]
\[
\begin{array}{llll}
f=a_x^3, & \Delta=(ab)^2a_xb_x, & \pi=(aV)^2a_x, & A_{\Delta\Delta}\\
\varphi=\alpha_x^3, & V=(\alpha\beta)^2\alpha_x\beta_x, & p=(\alpha\Delta)^2\alpha_x, & A_{VV}\\
Q=(a\Delta)a_x^2\Delta_x, & \Theta=(a\alpha)^2a_x\alpha_x, & (\Delta,p), & A_{\Delta V}\\
K=(\alpha V)\alpha_x^2V_x, & (f,p), & (\Delta,\pi), & A_{\Delta\Theta}\\
(f,V), & (\varphi,\pi), & (V,p), & A_{V\Theta}\\
(\varphi,\Delta), & (f,\pi)=-(\varphi,p), & (V,\pi), & J=(a\alpha)^3
\end{array}
\]
\[
\Omega=(\Delta V)(V\Theta)(\Delta\Theta).
\]

Zwischen den Ueberschiebungen dieser Formen bestehen insbesondere
folgende Relationen:
%CONT>>


% ===== pdf idx 360, printed 346.0 =====
% !PRINTED 346
\weberpage{346}
%<<CONT

\[
\left.\begin{array}{l}
(f,\Theta)^2=-\dfrac{1}{2}p,\qquad (\vartheta,f)^3=-\dfrac{3}{4}p\\[6pt]
(\varphi,\Theta)^2=-\dfrac{1}{2}\pi,\qquad (\vartheta,\varphi)^3=+\dfrac{3}{4}\pi
\end{array}\right\}\qquad\text{(I)}
\]
\[
\left.\begin{array}{l}
(f,p)=-2(\Theta,\Delta),\qquad (\vartheta,p)=\dfrac{1}{2}\{\Delta\pi+\Theta p\}\\[6pt]
(\varphi,\pi)=-2(\Theta,V),\qquad (\vartheta,\pi)=-\dfrac{1}{2}\{Vp+\Theta\pi\}
\end{array}\right\}\qquad\text{(II)}
\]
\[
\left.\begin{array}{l}
(f,\pi)=-(\Delta,V)\\
(\varphi,p)=-(V,\Delta)
\end{array}\right\},\quad\text{daher: }(f,\pi)+(\varphi,p)=0\qquad\text{(III)}
\]
\[
\left.\begin{array}{l}
p^2=A_{\Delta\Delta}V-2A_{\Theta\Delta}\Theta+\left(A_{\Theta\Theta}-\dfrac{J^2}{2}\right)\Delta+J(f,p)\\[6pt]
\pi^2=A_{V\Delta}\Delta-2A_{\Theta V}\Theta+\left(A_{\Theta\Theta}-\dfrac{J^2}{2}\right)V-J(\varphi,\pi)
\end{array}\right\}\qquad\text{(IV)}
\]
\[
\left.\begin{array}{l}
(f,p^2)^2=A_{\Delta\Delta}\pi+A_{\Theta\Delta}p+J(\Delta,p)\\
(\varphi,\pi^2)^2=A_{V\Delta}p+A_{\Theta V}\pi-J(V,\pi)\\
(\varphi,p^2)^2=A_{\Theta\Delta}\pi+\left(A_{\Theta\Theta}+\dfrac{J^2}{2}\right)p+J(\Delta,\pi)=-(f,\pi p)^2\\
(f,\pi^2)^2=A_{\Theta V}p+\left(A_{\Theta\Theta}+\dfrac{J^2}{2}\right)\pi-J(V,p)=-(\varphi,\pi p)^2
\end{array}\right\}\qquad\text{(V)}
\]
\[
A_{\Theta\Theta}=A_{V\Delta}-\dfrac{J^2}{2}\qquad\text{(VI)}
\]
\[
2\Omega=(p\pi),\qquad 2\Omega^2=(A_{\Delta\Delta}A_{\Theta\Theta}A_{VV}).\qquad\text{(VII)}
\]

Die Unterdeterminanten $U_{ik}$ dieser Determinante
$(A_{\Delta\Delta}A_{\Theta\Theta}A_{VV})$ sind den sechs Invarianten,
die zum simultanen System der drei Functionaldeterminanten $(\Theta,V)$,
$(V,\Delta)$, $(\Delta,\Theta)$ gehören, proportional. (Vergl. Nr. 320.)
Endlich ist:
\[
4\Omega J=U_{22}-4U_{13}.\qquad\text{(VIII)}
\]

\begin{center}
\S\ 33.\quad\textbf{Anwendungen des simultanen Formensystems\\
zweier cubischer Formen.}
\end{center}

323. \textit{Typische Darstellung der quadratischen Covarianten
$\Delta$, $\Theta$, $V$.} Die ersten Anwendungen, die wir von den im
vorhergegangenen Paragraphen gewonnenen Resultaten machen, beziehen
sich auf typische Darstellungen von Formen des Systemes unter
Zugrundelegung der beiden linearen Covarianten $p$ und $\pi$. Um eine
der drei quadratischen Covarianten $\Delta$, $V$ oder $\Theta$, etwa
die Covariante $\Delta$, typisch darzustellen, gehen wir wiederum aus
von der Identität
\[
(p\pi)\Delta_x=(\Delta\pi)p_x-(\Delta p)\pi_x.
\]


% ===== pdf idx 361, printed 347.0 =====
% !PRINTED 347
\weberpage{347}

Wir erheben sie auf die zweite Potenz und erhalten:
\[
4\Omega^2\Delta=(\Delta\pi)^2p^2-2(\Delta p)(\Delta\pi)p\pi+(\Delta p)^2\pi.
\]
Für die typischen Coefficienten $(\Delta\pi)^2$, $(\Delta p)(\Delta\pi)$,
$(\Delta p)^2$ ergeben sich aber unmittelbar die Werthe:
\[
\begin{aligned}
(\Delta p)^2&=2U_{33}=2\{A_{\Delta\Delta}A_{\Theta\Theta}-A_{\Delta\Theta}^2\}\\
(\Delta p)(\Delta\pi)&=U_{23}=A_{\Theta V}A_{\Delta\Delta}-A_{\Theta\Delta}A_{V\Delta}\\
(\Delta\pi)^2&=\dfrac{1}{2}U_{22}=\dfrac{1}{2}\{A_{\Delta\Delta}A_{VV}-A_{V\Delta}^2\}.
\end{aligned}
\]
Denn es ist (vergl. Nr. 320 und 322):
\[
\begin{aligned}
2U_{33}&=4((\Delta,\Theta),(\Delta,\Theta))^2=((f,p),(f,p))^2=(ab)^2(ap)(bp)=(\Delta p)^2\\
U_{23}&=2((V,\Delta),(\Theta,\Delta))^2=((f,p),(f,\pi))^2=(ab)^2(ap)(b\pi)=(\Delta p)(\Delta\pi)\\
\dfrac{1}{2}U_{22}&=((\Delta,V),(\Delta,V))^2=((f,\pi),(f,\pi))^2=(ab)^2(a\pi)(b\pi)=(\Delta\pi)^2.
\end{aligned}
\]

In derselben Weise gelangt man zu den typischen Coefficienten von $V$
und $\Theta$. Man braucht nur die Ueberschiebungen:
\[
((\varphi,\pi),(\varphi,\pi))^2,\quad ((\varphi,\pi),(\varphi,p))^2,\quad ((\varphi,p),(\varphi,p))^2
\]
\[
((f,p),(\varphi,p))^2,\quad ((f,\pi),(\varphi,p))^2,\quad ((f,\pi),(\varphi,\pi))^2
\]
zu berechnen.

324. \textit{Die typischen Coefficienten von $f$ und $\varphi$.} Ebenso
bietet, nachdem wir die Ueberschiebungen
\[
(f,p^2)^2,\quad (f,p\pi)^2,\quad (f,\pi^2)^2
\]
\[
(\varphi,p^2)^2,\quad (\varphi,p\pi)^2,\quad (\varphi,\pi^2)^2
\]
bereits berechnet haben, die typische Darstellung von $f$ und $\varphi$
mit Hilfe der linearen Covarianten $p$ und $\pi$ keine Schwierigkeiten
mehr. Wir haben die für diese Ueberschiebungen gegebenen Relationen
(vgl. Nr. 322) nur noch einmal über die Formen $p$ und $\pi$ zu
schieben und hiebei die sich aus den Betrachtungen in Nr. 323
ergebenden Werthe für $(\Delta p)^2$, $(\Delta\pi)^2$ etc. zu benutzen.

Sind auf diese Weise die typischen Coefficienten berechnet, so ist die
typische Darstellung von $f$ und $\varphi$ wieder durch die
Identitäten:
\[
(p\pi)^3a_x^3=\{(a\pi)p_x-(ap)\pi_x\}^3
\]
\[
(p\pi)^3\alpha_x^3=\{(\alpha\pi)p_x-(\alpha p)\pi_x\}^3
\]
unmittelbar gegeben.

Die leicht auszuführende Rechnung ergiebt, dass $f$ und $\varphi$ in
der typischen Darstellung bis auf einen Invariantenfactor als die
beiden
%CONT>>


% ===== pdf idx 362, printed 348.0 =====
% !PRINTED 348
\weberpage{348}

%<<CONT
ersten Differentialquotienten einer gewissen Form vierten Grades
erscheinen, nämlich der Form:
\[
\begin{aligned}
F&=-8\Omega^3\{f\pi+\varphi p\}\\
&=Ap^4+4Bp^3\pi+3Cp^2\pi^2+4Dp\pi^3+E\pi^4.
\end{aligned}
\]
Diese Eigenschaft hat \gesperrt{Clebsch} für das Problem der
Transformation dritter Ordnung in der Theorie der elliptischen
Functionen ausgebeutet. Da sie auch noch anderweitig interessante
Beziehungen bietet, so wollen wir bei Betrachtung dieser Form $F$ noch
ein wenig verweilen.

325. \textit{Die ersten Polaren der Form $F$ nach $p$ und $\pi$.} Wir
werden zunächst noch in anderer Weise zeigen, dass, wenn
\[
F=f\pi+\varphi p
\]
ist, die ersten Differentialquotienten von $F$ nach $p$ und $\pi$,
oder was dasselbe ist, die ersten Polaren nach $p$ und $\pi$
proportional sind den Formen $f$ resp. $\varphi$.

Setzen wir nämlich etwa
\[
F=r_x^4,
\]
so ist:
\[
4(F,\pi)=4r_x^3(r\pi);
\]
und wegen der Voraussetzung auch
\[
\begin{aligned}
4(F,\pi)&=3(f,\pi)\pi+(\pi,\pi)f+3(\varphi,\pi)p+\varphi(p,\pi)\\
&=3(a\pi)a_x^2\pi_x+3(\alpha\pi)\alpha_x^2p_x+\varphi(p,\pi).
\end{aligned}
\]
Berücksichtigen wir, dass $(f,\pi)=-(\varphi,p)$, so wird
\[
\begin{aligned}
4r_x^3(r\pi)&=-3(\alpha p)\alpha_x^2\pi_x+3(\alpha\pi)\alpha_x^2p_x+\varphi(p\pi)\\
&=3\alpha_x^2\{(\alpha\pi)p_x-(\alpha p)\pi_x\}+\varphi(p\pi)\\
&=4\varphi(p\pi)=8\Omega\varphi.
\end{aligned}
\]
Demnach ist:
\[
r_x^3(r\pi)=2\Omega\varphi,\qquad\text{(1)}
\]
und ebenso findet sich:
\[
r_x^3(rp)=-2\Omega f.\qquad\text{(2)}
\]

326. \textit{Das System der Form $F$.} Dasselbe besteht bekanntlich aus
den Formen:
\[
\Delta_F=(F,F)^2,\quad t_F=(F,\Delta_F),\quad i_F=(F,F)^4,\quad j_F=(F,\Delta_F)^4.
\]
Mit Hilfe der Relationen (1) und (2) können wir dasselbe leicht durch
die simultanen Formen von $f$ und $\varphi$ ausdrücken.

Ueberschiebt man nämlich diese beiden Relationen dreimal übereinander,
so kommt:
\[
(rr_1)^3(r\pi)(r_1p)=4\Omega^2J.
\]


% ===== pdf idx 363, printed 349.0 =====
% !PRINTED 349
\weberpage{349}
Vertauscht man links $r$ mit $r_1$ und nimmt die halbe Summe, so kommt:
\[
-\frac{1}{2}(rr_1)^4(p\pi) = 4\Omega^2J,
\]
oder, da $(p\pi)=2\Omega$:
\weq{1}{i_F = -4\Omega J.}
Ebenso findet man durch einmalige Ueberschiebung von (1) und (2):
\[
(rr_1)(r\pi)(r_1p)r_x^2r_{1x}^2 = -4\Omega^2(\varphi,f) = 4\Omega^2\vartheta.
\]
Vertauscht man wieder $r$ mit $r_1$ und nimmt die halbe Summe, so kommt:
\[
-\Omega\Delta_F = 4\Omega^2\vartheta,
\]
oder
\weq{2a}{\Delta_F = -4\Omega\vartheta.}
Zur Berechnung von $t_F$ und $j_F$ benutzen wir die in Nr. 322 aufgestellten vier Relationen:
\weq{2b}{\left.
\begin{aligned}
(\vartheta,f)^3 &= -\frac{3}{4}p, & (\vartheta,\varphi)^3 &= +\frac{3}{4}\pi\\
(\vartheta,p) &= \frac{1}{2}(\Delta\pi+\Theta p), & (\vartheta,\pi) &= -\frac{1}{2}(Vp+\Theta\pi)
\end{aligned}
\right\}}
Es ist dann:
\[
\begin{aligned}
j_F &= (F,-4\Omega\vartheta)^4 = -4\Omega(f\pi+\varphi p,\vartheta)^4\\
&= -4\Omega\{(\vartheta,f\pi)^4+(\vartheta,fp)^4\}\\
&= -4\Omega\{((\vartheta,f)^3,\pi)+((\vartheta,\varphi),p)\}\\
&= -4\Omega\left\{-\frac{3}{4}(p\pi)+\frac{3}{4}(\pi p)\right\},
\end{aligned}
\]
also:
\weq{3}{j_F = +12\Omega^2.}
Ferner geht mit R\"ucksicht auf die Formeln (2a) und (2b) und (1) und (2) Nr. 325 aus der Identit\"at:
\[
(F,\Delta_F)(p,\pi) = (F,p)(\Delta_F,\pi) - (F,\pi)(\Delta_F,p)
\]
die Relation hervor:
\[
\begin{aligned}
t_F\cdot 2\Omega &= (-2\Omega f)\left(-4\Omega\cdot-\frac{1}{2}\{Vp+\Theta\pi\}\right)\\
&\quad -(+2\Omega\varphi)\left(-4\Omega\cdot\frac{1}{2}(\Delta\pi+\Theta p)\right),
\end{aligned}
\]
oder:
\weq{4}{t_F = 2\Omega\cdot\{\varphi(\Delta\pi+\Theta p)-f(Vp+\Theta\pi)\}.}
Die eben angestellten Berechnungen zeigen einen merkw\"urdigen Zusammenhang zwischen den drei Formen $F$, $p$, $\pi$ einerseits und den beiden Formen $f$ und $\varphi$ andererseits. Sind die Co\"efficienten einer biquadratischen Form so beschaffen, dass die Invariante $j$ proportional der Invariante $(p\pi)$, dass also die Relation (3)
\[
12\Omega^2 = j_F
\]
%CONT>>


% ===== pdf idx 364, printed 350.0 =====
% !PRINTED 350
\weberpage{350}
%<<CONT
besteht, so stimmt das simultane System von $F$, $p$, $\pi$ \"uberein mit dem simultanen System zweier cubischen Formen $f$ und $\varphi$, welche $p$ und $\pi$ zu linearen Covarianten haben. In der That reducirt sich auch durch die Relation (3) die Zahl 9 der Constanten von $F$, $p$ und $\pi$ auf 8 willk\"urliche, was mit der Zahl der unabh\"angigen Coefficienten von $f$ und $\varphi$ \"ubereinstimmt.

327. \gesperrt{Resultante von $f$ und $\varphi$.} Die Resultante zweier Formen $F_1$ und $F_2$, also $R_{F_1,F_2}$, ist bis auf den Factor $(p\pi)^3$ gleich der Resultante der beiden Formen $\pi_2F_1-\pi_1F_2$ und $p_2F_1-p_1F_2$, d.\,h. es ist:
\[
(p\pi)^3R_{F_1,F_2} = R_{\pi_2F_1-\pi_1F_2,\,p_2F_1-p_1F_2}, \quad (\text{vgl. Bd. I Nr. 256}),
\]
wenn $F_1$ und $F_2$ Formen dritten Grades in $x$ sind. Sind nun $F_1$ und $F_2$ die ersten Differentialquotienten von $F$, so ist einestheils $R_{F_1,F_2}$ gleich der Discriminante $R$ von $F$ und andererntheils ist wegen der Relationen (1) und (2) Nr. 325
\[
R_{\pi_2F_1-\pi_1F_2,\,p_2F_1-p_1F_2} = R_{2\Omega\varphi,\,-2\Omega f},
\]
also bis auf einen Zahlenfactor gleich
\[
\Omega^6R_{f,\varphi}.
\]
Man hat also die Relation (von einem Zahlenfactor abgesehen):
\[
\Omega^3R = \Omega^6R_{f,\varphi},
\]
oder weil, wiederum von numerischen Factoren abgesehen,
\[
R = 6j_F^2 - i_F^3 = 6\cdot12^2\Omega^4 - 64\Omega^3\cdot J^3,
\]
so ist:
\[
R_{f,\varphi} = 32\{27\Omega - 2J^3\}.
\]

328. \gesperrt{F\"alle, in denen die typische Darstellung unm\"oglich ist.} Sind alle sechs linearen Formen
\[
\begin{gathered}
p,\quad (\Delta,p),\quad (V,p)\\
\pi,\quad (\Delta,\pi),\quad (V,\pi)
\end{gathered}
\]
einander proportional, so kann eine typische Darstellung durch lineare Formen nicht mehr geleistet werden. In diesem Falle verschwindet zun\"achst jedenfalls $\Omega=(p\pi)$. Dies kann eine Folge des Verschwindens aller Unterdeterminanten $U_{ik}$ sein oder nicht, und je nachdem werden entweder alle linearen Covarianten identisch verschwinden, oder alle einander proportional sein, oder bei nicht verschwindenden $U_{ik}$ zwei von einander verschiedene lineare Formen existiren. Wir wollen daher der Reihe nach folgende zwei Fragen untersuchen:

1) Welche Eigenschaften haben $f$ und $\varphi$, wenn $\Omega=0$, $U_{ik}\lessgtr0$?

2) Welche Eigenschaften haben $f$ und $\varphi$, wenn $U_{ik}=0$, und also auch $\Omega=0$?

329. \gesperrt{Beantwortung der ersten Frage.} Wir setzen voraus:
\[
\Omega=0,\quad U_{ik}\gtrless0.
\]
%CONT>>


% ===== pdf idx 365, printed 351.0 =====
% !PRINTED 351
\weberpage{351}
Wenn $\Omega$ verschwindet, also immer $p$ f\"ur einen bestimmten Werth von $\lambda$ mit $\lambda\cdot\pi$ \"ubereinstimmt, so muss wegen der Relation:
\[
(f,\pi)+(\varphi,p)=0
\]
die Ueberschiebung:
\[
(f,\pi)+\lambda\cdot(\varphi,p)=(f+\lambda\varphi,\pi)
\]
identisch verschwinden. Das Verschwinden der Functionaldeterminante
\[
(f+\lambda\varphi,\pi)=0
\]
lehrt aber, dass $f+\lambda\varphi$ eine Potenz von $\pi$ sein muss. Wir haben also in diesem Falle:
\[
f+\lambda\varphi = c\cdot\pi^3,
\]
d.\,h. es giebt eine Combination der beiden Formen $f$ und $\varphi$, welche ein vollst\"andiger Cubus ist. Umgekehrt: Ist $f+\lambda\varphi$ ein vollst\"andiger Cubus, also
\[
f+\lambda\varphi = r_x^3,
\]
dann ist die zweite Ueberschiebung von $f+\lambda\varphi$ \"uber $V$, wegen $(\varphi,V)^2=0$:
\weq{1}{p = r(Vr)^2}
und ferner ist die zweite Ueberschiebung \"uber $\Delta$, wegen $(f,\Delta)^2=0$:
\weq{2}{\lambda\pi = r(\Delta r)^2}
Aus der Comparation von (1) mit (2) folgt:
\[
p = \lambda\cdot\frac{(Vr)^2}{(\Delta r)^2}\cdot\pi = \varrho\cdot\pi
\]
und demnach:
\[
\Omega = \frac{(p\pi)}{2} = 0.
\]

330. \gesperrt{Beantwortung der zweiten Frage.} Wir setzen voraus:
\[
U_{ik}=0,\quad\text{und demnach}\quad \Omega=0,\quad i\text{ und }k=1,2,3.
\]
In diesem Falle sind aber auch $p$ und $\pi$ identisch null. Um diesen Lehrsatz zu beweisen, haben wir nur zu zeigen, dass $\Delta$, resp. $V$ identisch verschwindet; dann ist auch
\[
(\varphi,\Delta)^2 = p = 0,\quad\text{und}\quad (f,V)^2 = \pi = 0.
\]
Nun ist (vgl. Nr. 323)
\[
U_{33} = \frac{1}{2}(\Delta p)^2.
\]
Da aber $U_{33}=0$, so folgt
\[
(\Delta p)^2 = 0
\]
und ebenso
\[
(Vp)^2 = 0,\quad (\Theta p)^2 = 0\ \text{etc.}
\]
Diese Relationen lehren, dass $\Delta$, $V$, $\Theta$ den Factor $p$ im ersten Grade besitzen. Nun ist aber:
%CONT>>


% ===== pdf idx 366, printed 352.0 =====
% !PRINTED 352
\weberpage{352}
%<<CONT
\[
\begin{aligned}
\Theta_x^2(\Delta p)^2 - \Delta_x^2(\Theta p)^2 &= (\Theta_x(\Delta p)+\Delta_x(\Theta p))(\Theta_x(\Delta p)-\Delta_x(\Theta p))\\
&= (\Theta_x(\Delta p)+\Delta_x(\Theta p))(\Delta\Theta)p_x\\
&= ((\Delta\Theta)\Delta_x\Theta_x,p)p\\
&= ((\Delta,\Theta),p)p.
\end{aligned}
\]
Die linke Seite ist null; folglich auch, wenn $p\gtrless0$ ist:
\[
((\Delta,\Theta),p) = 0.
\]
Es ist aber (vgl. Nr. 322):
\[
(\Delta,\Theta) = \frac{1}{2}(f,p) = \frac{1}{2}(ap)a_x^2,
\]
also
\[
(ap)^2a_x = 0
\]
und ebenso
\[
(\alpha p)^2\alpha_x = 0,
\]
d.\,h. die cubischen Formen $f$ und $\varphi$ besitzen $p$ als quadratischen Factor. Daraus folgt nach der Theorie der cubischen Formen, dass auch $\Delta$ und $V$ ihn quadratisch enthalten. Ebenso aber auch $\Theta$. Denn wenn $f$ und $\varphi$ den Factor $p$ quadratisch enthalten, so besitzt ihn auch $f+\lambda\varphi$ und demnach auch die Hesse'sche $\Delta_{f+\lambda\varphi}$ dieser Form. Nun ist
\[
\Delta_{f+\lambda\varphi} = \Delta + 2\lambda\Theta + \lambda^2V.
\]
Hierin besitzen $\Delta_{f+\lambda\varphi}$, $\Delta$, $V$ den Factor $p$ quadratisch, also auch $\Theta$. Wenn aber $\Delta$, $V$, $\Theta$ diesen Factor quadratisch enthalten, so ist
\[
(\Delta p)\Delta_x = 0,\quad (Vp)V_x = 0,\quad (\Theta p)\Theta_x = 0.
\]
Bilden wir also:
\[
\begin{aligned}
\Theta_x^2(\Delta p)\Delta_x - \Delta_x^2(\Theta p)\Theta_x &= \Theta_x\Delta_x\{(\Delta p)\Theta_x-(\Theta p)\Delta_x\}\\
&= p\cdot(\Delta\Theta)\Theta_x\Delta_x = p\cdot(\Delta,\Theta),
\end{aligned}
\]
so ist, da die linke Seite verschwindet und $p$ von Null verschieden sein soll:
\[
(\Delta,\Theta) = (f,p) = a_x^2(ap) = 0,
\]
d.\,h. $f$ besitzt den Factor $p$ im dritten Grade, und also ebenso auch $\varphi$. Wenn aber $f$ und $\varphi$ reine Cuben sind, dann verschwinden ihre Hesse'schen Formen und folglich ist, wie behauptet wurde:
\[
(\varphi,\Delta)^2 = p = 0,\quad (f,V)^2 = \pi = 0.
\]

331. \gesperrt{Folgerungen f\"ur $f$ und $\varphi$.} Wir setzen also voraus
\[
p=0,\quad \pi=0.
\]
Wenn $p$ und $\pi$ identisch verschwinden, so m\"ussen ihre Co\"efficienten $p_1$ $p_2$ $\pi_1$ $\pi_2$ einzeln null sein, d.\,h. (vgl. Nr. 312) es m\"ussen die Determinanten der Matrix
%CONT>>


% ===== pdf idx 367, printed 353.0 =====
% !PRINTED 353
\weberpage{353}
Systeme einzelner und simultaner Formen höheren Grades.

\[
V = \begin{vmatrix}
\overline{a_0} & \overline{a_1} & \overline{\alpha_0} & \overline{\alpha_1} \\
\overline{a_1} & \overline{a_2} & \overline{\alpha_1} & \overline{\alpha_2} \\
\overline{a_2} & \overline{a_3} & \overline{\alpha_2} & \overline{\alpha_3}
\end{vmatrix}
\]
verschwinden. Es mögen hiebei zunächst wenigstens die Minoren dieser Determinanten von null verschieden sein. Wir multipliciren alsdann mit dieser Matrix Zeile für Zeile die Determinante
\[
\begin{vmatrix}
x_1 & x_2 & 0 & 0 \\
\Delta_1 & -\Delta_2 & 0 & 0 \\
0 & 0 & x_1 & x_2 \\
0 & 0 & V_2 & -V_1
\end{vmatrix}.
\]
Das Product ist einerseits Null, da alle Determinanten von $V$ verschwinden. Andererseits ist es dargestellt durch die Matrix:
\[
\begin{vmatrix}
f_{11} & Q_{11} & \varphi_{11} & K_{11} \\
f_{12} & Q_{12} & \varphi_{12} & K_{12} \\
f_{22} & Q_{22} & \varphi_{22} & K_{22}
\end{vmatrix}.
\]
Sie verschwindet, d. h. ihre einzelnen Determinanten sind null, nicht aber deren Minoren. Greifen wir nun etwa die erste Determinante heraus
\weq{1}{
\begin{vmatrix}
f_{11} & Q_{11} & \varphi_{11} \\
f_{12} & Q_{12} & \varphi_{12} \\
f_{22} & Q_{22} & \varphi_{22}
\end{vmatrix} = 0.
}
Da die Minoren nicht verschwinden, so muss $\varphi$ eine lineare Function von $f$ und $Q$ sein. Denn die Determinante (1) lässt sich schreiben:
\[
\begin{vmatrix}
\overline{a_0}\,x_1 + \overline{a_1}\,x_2, & \overline{Q_0}\,x_1 + \overline{Q_1}\,x_2, & \overline{\alpha_0}\,x_1 + \overline{\alpha_1}\,x_2 \\
\overline{a_1}\,x_1 + \overline{a_2}\,x_2, & \overline{Q_1}\,x_1 + \overline{Q_2}\,x_2, & \overline{\alpha_1}\,x_1 + \overline{\alpha_2}\,x_2 \\
\overline{a_2}\,x_1 + \overline{a_3}\,x_2, & \overline{Q_2}\,x_1 + \overline{Q_3}\,x_2, & \overline{\alpha_2}\,x_1 + \overline{\alpha_3}\,x_2
\end{vmatrix} = 0.
\]
Zerlegt man diese Determinante in ihre 8 Summanden und fasst gleiche Potenzen von $x$ zusammen, so müssen die Coefficienten von $x^\lambda$ einzeln verschwinden. Nun ist z. B. der Coefficient von $x_1{}^3$
\[
\begin{vmatrix}
\overline{a_0} & \overline{Q_0} & \overline{\alpha_0} \\
\overline{a_1} & \overline{Q_1} & \overline{\alpha_1} \\
\overline{a_2} & \overline{Q_2} & \overline{\alpha_2}
\end{vmatrix} = 0,
\]
d. h. es müssen die Gleichungen bestehen:
\[
\begin{aligned}
k\overline{a_0} + \lambda \overline{Q_0} + \varrho\overline{\alpha_0} &= 0 \\
k\overline{a_1} + \lambda \overline{Q_1} + \varrho\overline{\alpha_1} &= 0 \\
k\overline{a_2} + \lambda \overline{Q_2} + \varrho\overline{\alpha_2} &= 0.
\end{aligned}
\]
%CONT>>


% ===== pdf idx 368, printed 354.0 =====
% !PRINTED 354
\weberpage{354}
Dritter Theil.

%<<CONT
Daraus geht aber die Richtigkeit der Behauptung direct hervor, dass $\varphi = a_x^3$ eine lineare Function $f = a_x^3$ und $Q = Q_x^3$ ist.

Wenn also $p$ und $\pi$ identisch null sind, so ist $\varphi$ keine selbstständige Form, sondern eine Covariante $f + \mu Q$. Dies gilt, so lange die Minoren der Determinanten $\gtrless 0$ sind. Verschwinden auch diese identisch, die, in so weit sie aus den beiden ersten resp. aus den beiden letzten Colonnen von $V$ entnommen sind, nichts anderes als die Coefficienten von $\Delta$ und $V$ darstellen, so sind $f$ und $\varphi$ vollständige Cuben, da $\Delta$ und $V$ in diesem Falle null sind.

Wir haben demnach folgende Resultate:

\glqq Ist $\Omega = 0$, aber $U_{ik}$ im Allgemeinen von null verschieden, so existirt eine lineare Combination $f + \lambda\varphi$, welche dem Cubus der linearen Form $\pi$ proportional ist.\grqq{}

\glqq Ist $U_{ik} = 0$, $i$ und $k = 1, 2, 3$ und daher $\Omega = 0$, so ist entweder $\varphi = \mathrm{Const.}\,(f + \mu Q)$, oder $f$ und $\varphi$ sind vollständige Cuben; in diesen beiden Fällen ist eine typische Darstellung durch $p$ und $\pi$ unmöglich.\grqq{}

\bigskip
\centerline{\S\ 34. \textbf{Die Schwesterformen.}}
\bigskip

332. \textit{Einleitende Bemerkungen.} Wir haben uns bisher mit der Aufstellung des Systemes jener Grundformen beschäftigt, welche die Eigenschaft haben, dass alle Co- und Invarianten, die eine bestimmte Form besitzt, rationale und ganze Functionen derselben sind.

Einem solchen System gegenüber steht nun eine Reihe anderer Systeme von Formen von der Eigenschaft, dass sämmtliche Co- und Invarianten von $f = a_x^n$ zwar noch rationale Functionen derselben sind, aber nicht mehr ganze. Dies ist von vornherein einzusehen. Denn die Form $f$ besitzt zunächst $(n+1)$ Constante. Vier derselben können durch lineare Transformationen bestimmte Zahlenwerthe annehmen; die ursprüngliche Form hängt alsdann von den $n - 3$ transformirten Constanten und der Transformationsdeterminante ab, d. i. von $n - 2$ Parametern. Es muss also bereits zwischen $n - 1$ Invarianten der Form $f$ eine Relation bestehen. Nehmen wir die beiden Variabeln $x_1$ und $x_2$ als zwei weitere Parameter hinzu, so hängt jede Covariante der Form $f$ von $n$ willkürlichen Grössen ab; demnach sind auch stets $n + 1$ Covarianten von $f$ durch eine Relation verknüpft, deren Coefficienten nur numerische Constante sind. Wenn nun das vollständige System von $f$ aus $s$ Formen besteht, so müssen zwischen denselben $s - n$ Relationen existiren. Jede Covariante lässt sich aber durch die $s$ Formen rational und ganz darstellen. Eliminirt man ver-
%CONT>>


% ===== pdf idx 369, printed 355.0 =====
% !PRINTED 355
\weberpage{355}
Systeme einzelner und simultaner Formen höheren Grades.

%<<CONT
möge der $s - n$ Relationen aus jeder solchen Darstellung $s - n$ bestimmte Formen, ein Process, der immer rational ausführbar ist, aber im Allgemeinen zu einer gebrochenen Function führt, so ist jede Covariante durch eine rationale und gebrochene Function von $n$ bestimmten Covarianten dargestellt. Diese $n$ ausgewählten Covarianten bilden in diesem Sinne ein System, und die Aufgabe, ein solches System von Formen herzustellen, ist auf unendlich mannigfaltige Weisen lösbar, da ja auch die $s$ Formen des vollständigen Systemes nicht a priori fest sind, sondern stets durch andere ersetzt werden können. Der einfachste Fall eines solchen Systemes von $n$ Formen ist jener, in welchem der Nenner dieser rationalen Functionen eine Potenz von $f$ ist, während die Zähler Aggregate von Producten ganz bestimmter Covarianten sind.

Wir werden uns hier nur mit diesem einfachsten Falle beschäftigen und verweisen wegen der allgemeinen Untersuchungen auf Clebsch: \glqq Binäre Formen\grqq{}, siebenter Abschnitt. Hier sei nur noch erwähnt, dass die Aufgabe, jene Formen zu ermitteln, durch welche sich alle Formen des Systems rational ausdrücken lassen, im innigsten Zusammenhange steht mit der typischen Darstellung von Formen, sowie mit der Aufgabe, Covariantenrelationen herzustellen, wie auch aus diesen einleitenden Betrachtungen hervorgehen mag (vgl. auch Nr. 335).

333. \textit{Jedes symbolische Product ist gleich einer rationalen Function von Covarianten, deren Nenner eine Potenz von $f$ ist.} Um ein beliebiges symbolisches Product $P$ durch eine rationale Function von Formen darzustellen, deren Nenner eine Potenz von $f$ ist, können wir auf folgende Weise verfahren. Wir denken uns zunächst dasselbe in Symbolen von $f = a_x^n = b_x^n = c_x^n$ etc. allein dargestellt und erinnern uns, dass wir früher gesetzt hatten:
\[
f_1 = a_x^{n-1}a_1 = \frac{1}{n}\frac{\partial f}{\partial x_1}
\]
\[
f_2 = a_x^{n-2}a_2 = \frac{1}{n}\frac{\partial f}{\partial x_2}.
\]
Es ist dann in Analogie mit $a_x = a_1x_1 + a_2x_2$ auch $f_x = f_1x_1 + f_2x_2$, und anderntheils ist $f_1x_1 + f_2x_2$ nach dem Euler'schen Satze auch identisch mit $f = a_x^n$. Jeder Klammerfactor $(ab)$, $(cd)$ etc. des symbolischen Products lässt sich aber nach dem Identitätssatze ersetzen durch
\weq{1}{
(ab) = \frac{b_x(af) - a_x(bf)}{f_x}.
}
Enthält das symbolische Product $P$ $\nu$ solcher Klammerfactoren, so geht
%CONT>>


% ===== pdf idx 370, printed 356.0 =====
% !PRINTED 356
\weberpage{356}
Dritter Theil.

%<<CONT
dasselbe durch Substitution ihrer Werthe (1) in einen Quotienten über, dessen Nenner eine Potenz von $f$ ist, und dessen Zähler ein Aggregat von Producten ist, die sich aus Factoren $a_x, b_x, c_x \ldots$ erster Art, und Klammerfactoren $(af), (bf), (cf) \ldots$ etc. zusammensetzen, während Klammerfactoren vom Typus $(ab), (ac), (bc) \ldots$ etc. nicht mehr vorkommen. Das Symbol $a$ kann also beispielsweise nur mehr in einer Form vorkommen wie:
\weq{2}{
\psi_k = a_x^{n-k}(af)^k,
}
deren unsymbolische Bedeutung die folgende Untersuchung klar legen soll.

334. \textit{Alle Formen $\psi$ haben die Form $f$ als Factor.} Zunächst wollen wir zeigen, dass jede Form $\psi$ in die Factoren $f \cdot u$ zerfällt, wo $u$ ein Aggregat von symbolischen Producten ist. Dies ist für $k = 0$ und $k = 1$ von vornherein ersichtlich. Denn für $k = 0$ hat man
\[
\psi_0 = a_x^{n-0}(af)^0 = a_x^n = f \cdot u, \text{ wo } u_0 = 1;
\]
für $k = 1$ folgt
\[
\begin{aligned}
\psi_1 = a_x^{n-1}(af)^1 &= a_x^{n-1}\{a_1b_2b_x^{n-1} - a_2b_1b_x^{n-1}\} \\
&= (ab)\,a_x^{n-1}b_x^{n-1} = 0 = f \cdot u_1, \text{ wo } u_1 = 0.
\end{aligned}
\]
Wir setzen nun voraus, der Satz gelte bereits für $k = \varrho$, und zeigen, dass er alsdann auch für $k = \varrho + 1$ giltig ist. Es sei demnach:
\[
\psi_\varrho = a_x^{n-\varrho}(af)^\varrho = f \cdot U_\varrho.
\]
Die erste Polare von $\psi_\varrho$ wird:
\[
\begin{aligned}
&(n - \varrho)\,a_x^{n-\varrho-1}a_y(af)^\varrho + \varrho\,a_x^{n-\varrho}(af)^{\varrho-1}\cdot(n-1)b_x^{n-2}b_y(ab) \\
&= nf_y + \mathrm{Const.}\times f\cdot U_{\varrho y}.
\end{aligned}
\]
Ersetzen wir hierin $y_1$ durch $f_2$ und $y_2$ durch $-f_1$, so erhalten wir, wenn wir jene Grösse, die dadurch aus $U_{\varrho y}$ wird, mit $M$ bezeichnen:
\weq{3}{
\begin{aligned}
&(n - \varrho)\,a_x^{n-\varrho-1}(af)^{\varrho+1} + \varrho(n-1)\,a_x^{n-\varrho}(ab)b_x^{n-2}(af)^{\varrho-1}(bf) \\
&= 0 + \mathrm{Const.}\times f \cdot M.
\end{aligned}
}
In dieser Identität ist das erste Glied links eine Form $\psi$ und zwar gerade die nächstfolgende Form $\psi_{\varrho+1}$ (vom Factor $n - \varrho$ abgesehen). Das zweite Glied rechts besitzt den Factor $f$; können wir also zeigen, dass der zweite Term links ebenfalls diesen Factor enthält, so ist der zu führende Beweis erbracht. Nun ist aber:
\[
\begin{aligned}
a_x^{n-\varrho}b_x^{n-2}(ab)(bf)(af)^{\varrho-1} &= a_x^{n-\varrho}b_x^{n-\varrho}\cdot(ab)(bf)(af)^{\varrho-1}b_x^{\varrho-2} \\
&= a_x^{n-\varrho}b_x^{n-\varrho}(ab)(bf)(af)\cdot b_x^{\varrho-2}(af)^{\varrho-2} \\
&= \frac{1}{2}a_x^{n-\varrho}b_x^{n-\varrho}(ab)(bf)(af)\cdot\{(af)^{\varrho-2}b_x^{\varrho-2} - (bf)^{\varrho-2}a_x^{\varrho-2}\}.
\end{aligned}
\]
%CONT>>


% ===== pdf idx 371, printed 357.0 =====
% !PRINTED 357
\weberpage{357}
Die in der Klammer stehende Differenz besitzt den Factor
\[
(af)\,b_x - (bf)\,a_x = f_x\,(ab) = f\cdot (ab).
\]
Es besitzt daher in der That der zweite Term links in (3) den Factor
$f$ und damit auch der erste Term, was zu beweisen war.

335. \textit{Definition der Schwesterformen.} Jede Form $\psi_k$ zerfällt also
in das Product $u_k\cdot f$; die Formen
\[
u_k = \frac{a_x^{n-k}(af)^k}{f}
\]
nennen wir \glqq Schwesterformen\grqq{} (formes associées) der Form $f$. Die
Anzahl der Schwesterformen einer Form $f$ ist $n$, oder wenn wir $f$ mit
einschliessen gleich $n+1$, denn $k$ kann nur die Werthe $0, 1, 2, \dots n$
annehmen. Indess ist, wie wir schon oben gesehen haben, $u_0=1$,
$u_1=0$.

Es ist ferner:
\[
u_2 = \frac{1}{2}(f,f)^2 = \frac{1}{2}H.
\]
Denn man hat:
\[
\begin{aligned}
a_x^{n-2}(af)^2 &= a_x^{n-2}(af)\cdot(ab)\,b_x^{n-1}\\
&= \frac{1}{2}\,a_x^{n-2}b_x^{n-2}(ab)\,\{(af)\,b_x - a_x(bf)\}\\
\text{oder}\qquad (af)^2 a_x^{n-2} &= \frac{1}{2}\,a_x^{n-2}b_x^{n-2}(ab)^2\cdot f = \frac{1}{2}(f,f)^2\cdot f.
\end{aligned}
\]
Hieraus erhält man $u_3$, wenn man von dieser Gleichung die erste
Polare bildet und wie in Nr. 334 (3) $y$ durch $f_2$, $y_2$ durch
$-f_1$ ersetzt. Dann kommt:
\[
(n-2)a_x^{n-3}(af)^3+2a_x^{n-2}(af)\cdot(n-1)(bf)b_x^{n-2}(ab)=(n-2)f\cdot((f,f)^2,f).
\]
Der zweite Term links verschwindet, da er durch Vertauschung von
$a$ mit $b$ nur sein Zeichen ändert; demnach wird:
\[
a_x^{n-3}(af)^3 = f\cdot((f,f)^2,f) = f\cdot T
\]
\[
u_3 = T = ((f,f)^2,f).
\]
Auf analoge Weise findet man der Reihe nach:
\[
u_4 = \frac{1}{2}(f,f)^4\cdot f^2 - \frac{3}{4}\left\{(f,f)^2\right\}^2
\]
\[
u_5 = (f,f)^4\cdot f^3 - (f,f)^2\cdot((f,f),f)
\]
\[
u_6 = \frac{1}{2}(f,f)^6\cdot f^4 - \frac{15}{4}(f,f)^4\cdot(f,f)^2\cdot f^2 + \frac{45}{8}\left\{(f,f)^2\right\}^3 + 10\left\{((f,f)^2,f)\right\}^2 \qquad \text{u.\,s.\,w.}
\]
Diese Formen $u_\varrho$ sind es also, durch welche sich jede Covariante von
$f$ rational darstellen lässt. Insbesondere bilden sie die typischen
%CONT>>


% ===== pdf idx 372, printed 358.0 =====
% !PRINTED 358
\weberpage{358}
%<<CONT
Coefficienten, wenn die Form $f$ selbst durch sie ausgedrückt werden
soll. Zu dem Zwecke hat man nur $f(y)$ durch die lineare Substitution
\[
\begin{aligned}
\xi &= f_1y_1+f_2y_2 = f_y\\
\eta &= -x_2y_1+x_1y_2 = (xy)
\end{aligned}
\]
vom Modul $\Delta = f_1x_1+f_2x_2=f$ zu transformiren und erhält dann
\[
f^{n-1}\cdot f(y) = u_0\xi^n+\binom{n}{1}u_1\xi^{n-1}\eta+\binom{n}{2}u_2\xi^{n-2}\eta^2+\cdots u_n\eta^n.
\]
So erhält man beispielsweise für $n=4$
\[
\begin{aligned}
f^3\cdot f(y) &= u_0\xi^4+4u_1\xi^3\eta+6u_2\xi^2\eta^2+4u_3\xi\eta^3+u_4\eta^4\\
&= \xi^4+3H\xi^2\eta^2+4T\xi\eta^3+\left(\frac{i}{2}f^2-\frac{3}{4}H^2\right)\eta^4.
\end{aligned}
\]
Will man nun irgend eine Covariante $P$ von $f$ durch die
Schwesterformen $u_0u_1\dots u_n$ darstellen, so genügt bereits deren
Anfangsglied, das im Falle einer Invariante mit dieser zusammenfällt.
Man ersetzt alsdann im Anfangsglied die ursprünglichen Coefficienten
$\overline{a}_0\,\overline{a}_1\dots\overline{a}_n$ durch die typischen $u_0=1$, $u_1=0$,
$u_2=\frac{H}{2}$ etc. und dividirt mit einer passenden Potenz von $f$.
So ist in unserm Beispiele $n=4$ die Form $j$ keine Schwesterform.
Wenn wir daher in der Determinante
\[
j = 6\begin{vmatrix} \overline{a}_0 & \overline{a}_1 & \overline{a}_2 \\ \overline{a}_1 & \overline{a}_2 & \overline{a}_3 \\ \overline{a}_2 & \overline{a}_3 & \overline{a}_4 \end{vmatrix}
\]
die Coefficienten $\overline{a}_i$ durch die Schwesterformen $u_i$ ersetzen und
mit $f^3$ dividiren, so kommt:
\[
j = \frac{1}{f^3}\left\{\frac{3}{2}\,iHf^2 - 3H^3 - 6T^2\right\}\cdot
\]
Man erkennt zugleich in dieser Gleichung jene einzige Relation, welche
zwischen den fünf Covarianten von $f=a_x^4$ bestehen muss. (Vgl. auch
Nr. 162 und Clebsch, a.\,a.\,O., 7. Abschnitt.)

336. \textit{Stellung der Schwesterformen zum vollständigen System von $f$.}
Betrachtet man die Nr. 335 gegebenen Darstellungen der Schwesterformen
$u_0, u_1, u_2, u_3, u_4, u_5, u_6$ genauer, so sieht man, dass diese sieben
Formen ganze und rationale Functionen von $f$, von Ueberschiebungen
zweiten Grades $(f,f)^{2k}$ in den Coefficienten von $f$ und von
Ueberschiebungen dritten Grades $((f,f)^{2k},f)$ sind. Nachdem nun bereits
Hermite im Jahre 1854, Crelle's Journ. Bd.\,52 seine grundlegenden
Gedanken über associirte Formen und typische Darstellung
veröffentlicht hatte, unterwarf später Clebsch diese Begriffe einer
allgemeinen
%CONT>>


% ===== pdf idx 373, printed 359.0 =====
% !PRINTED 359
\weberpage{359}
%<<CONT
Untersuchung und fragte sich insbesondere auch, welche Stellung das
einfachste System von associirten Formen $u_0u_1\dots u_n$ zum vollen
System von $f$ einnehme. Dabei zeigte sich, dass, was eben die ersten
sieben Formen schon vermuthen lassen, jede Form $u_i$ eine ganze und
rationale Function\par
1) der Form $f$ selbst,\\
2) der Formen $(f,f)^{2k}=G_{2k}$,\\
3) der Functionaldeterminanten $((f,f)^{2k},f)=H_{2k}$\\
ist. Die Anzahl dieser Formen $f$, $G_{2k}$, $H_{2k}$, $k=1,2\dots n$ ist in
der That gleich der Anzahl der Schwesterformen $u_0u_1\dots u_n$ und es
handelt sich nur, Recursionsformeln herzustellen, welche die Formen
$f$, $G_{2k}$, $H_{2k}$ einerseits mit den Formen $u_\varrho$ andererseits in
Beziehung setzen. Bestätigt sich die eben aufgestellte Behauptung, so
können wir folgenden Satz aussprechen:

\glqq Jede Covariante und Invariante von $f$ ist eine rationale
gebrochene Function der $(n+1)$ Formen $f$, $G_{2k}$, $H_{2k}$, deren
Nenner eine Potenz von $f$ ist.\grqq{}

337. \textit{Recursionsformeln für $G_{2k}$ und $H_{2k}$.} Um die gesuchten
Beziehungen aufzustellen, verfahren wir ganz nach der schon im
Eingange dieses Paragraphen gegebenen allgemeinen Methode. Wir
ersetzen in den Formen
\[
\begin{aligned}
G_{2k} &= (ab)^{2k}a_x^{n-2k}b_x^{n-2k}\\
H_{2k} &= (ab)^{2k}(a_x^{n-2k}b_x^{n-2k},\,c_x^n)\\
&= \frac{1}{2}(ab)^{2k}\big(b_x^{n-2k}a_x^{n-2k-1}c_x^{n-1}(ac)+a_x^{n-2k}b_x^{n-2k-1}(bc)c_x^{n-1}\\
&= (ab)^{2k}(ac)a_x^{n-2k-1}b_x^{n-2k}c_x^{n-1}
\end{aligned}
\]
den Klammerfactor $(ab)$ durch $\dfrac{b_x(af)-a_x(bf)}{f}$, und die Grösse
$(ac)c_x^{n-1}$ durch $(af)$. Dann erhalten wir:
\[
\begin{aligned}
G_{2k}\cdot f^{2k} &= a_x^{n-2k}b_x^{n-2k}\big((af)\,b_x-(bf)\,a_x\big)^{2k}\\
H_{2k}\cdot f^{2k} &= a_x^{n-2k-1}b_x^{n-2k}(af)\big((af)\,b_x-(bf)\,a_x\big)^{2k}.
\end{aligned}
\]
Entwickelt man rechts nach dem binomischen Lehrsatze und ersetzt
sofort $(af)^\varrho a_x^{n-\varrho}$ durch $f\cdot u_\varrho$, so kommt:
{\scriptsize
\[
G_{2k}\cdot f^{2k} = u_{2k}\cdot f^2 - \binom{2k}{1}u_{2k-1}\cdot u_1f^2+\binom{2k}{2}u_{2k-2}\cdot u_2\cdot f^2
-\cdots(-1)^k\binom{2k}{k}u_k^2\cdot f^2\cdots-\binom{2k}{1}u_1\cdot u_{2k-1}\cdot f^2+u_{2k}\cdot f^2;
\]
}
%CONT>>


% ===== pdf idx 374, printed 360.0 =====
% !PRINTED 360
\weberpage{360}
%<<CONT
{\scriptsize
\[
H_{2k}\cdot f^{2k} = u_{2k+1}\cdot f^2 - \binom{2k}{1}u_{2k}\cdot u_1\cdot f^2+\binom{2k}{2}u_{2k-1}u_2\cdot f^2
-\cdots(-1)^k\binom{2k}{k}u_{k+1}\cdot u_k\cdot f^2\cdots-\binom{2k}{1}u_{2k-1}\cdot u_2\cdot f^2+u_{2k}u_1\cdot f^2.
\]
}
Dividirt man die beiden Gleichungen mit $f^2$, führt in sie sodann die
Werthe $u_0=1$, $u_1=0$ ein, und zieht endlich je zwei gleichwerthige
Glieder zusammen, so kommt:
\[
G_{2k}\cdot f^{2k-2} = 2\left\{u_{2k}u_0+\binom{2k}{2}u_{2k-2}u_2-\binom{2k}{3}u_{2k-3}u_3-\cdots\frac{1}{2}(-1)^k\binom{2k}{k}u_k^2\right\}
\]
\[
\begin{aligned}
H_{2k}\cdot f^{2k-2} &= u_{2k+1}u_0 + \frac{2k(2k-3)}{2!}u_{2k-1}u_2\\
&\qquad -\frac{2k(2k-1)(2k-5)}{3!}u_{2k-2}u_3\cdots(-1)^k\binom{2k}{k}u_{k+1}u_k.
\end{aligned}
\]
Dies sind die beiden Recursionsformeln, aus denen sich successive jede
Schwesterform $u_\varrho$ als rationale und ganze Function der $G_\varrho$ und
$H_\varrho$ berechnen lässt. So erhält man z.\,B.
\[
\begin{aligned}
G_2 &= (f,f)^2 = 2u_2,\\
G_4\cdot f^2 &= (f,f)^4\cdot f^2 = 2\{u_4+3u_2^{\,2}\},\\
G_6\cdot f^4 &= (f,f)^6\cdot f^4 = 2\{u_6+15u_4u_2-10u_3^{\,2}\};\\
H_2 &= ((f,f)^2,f) = u_3,\\
H_4\cdot f^2 &= ((f,f)^2,f)\cdot f^2 = u_5+2u_2u_3,\\
H_6\cdot f^4 &= ((f,f)^6,f)\cdot f^4 = u_7+9u_2u_5-5u_3u_4,\quad \text{u.\,s.\,w.}
\end{aligned}
\]
Bestimmt man hieraus die Werthe von $u_2, u_3, u_4\dots u_7$, so erhält
man genau dieselben Functionen in $H$ und $G$, welche wir bereits
früher auf anderm Wege ermittelt haben. Eine independente
Darstellung der Schwesterform $u_\varrho$ ist bis jetzt noch nicht geliefert
worden.

\weberfigure{360}{Typographischer Schlussschmuck: eine waagrechte Zierlinie mit rautenförmigen Verzierungen in der Mitte, darunter ein kurzer versprengter Druckstrich; markiert das Ende von Nr.\ 337.}


% ===== pdf idx 375, printed 360.001 =====
% !PRINTED none
% [no running-head folio number visible on this page image; it is a publisher's
%  advertisement page (Verlagsanzeige) listing other Teubner/Vieweg mathematics texts]
%<<CONT
\textbf{Graefe, Dr. Fr.,} Professor, Auflösungen u. Beweise der Aufgaben und Lehrsätze
aus der analytischen Geometrie des Punktes, der geraden Linie, des Kreises und der
Kegelschnitte. Für Studierende an Universitäten und technischen Hochschulen.
[IV u. 259 S.] gr. 8. geh. n. M 4.80.

\textbf{Heinze, Dr. Karl,} weiland Professor in Cöthen, genetische Stereometrie,
bearbeitet von Franz Lucke, Gymnasiallehrer in Zerbst. Mit lithographierten Tafeln.
[XII u. 194 S.] gr. 8. geh. n. M 6.\,---

\textbf{Hochheim, Dr. Adolf,} Professor, Aufgaben aus der analytischen Geometrie der
Ebene. Heft III. Die Kegelschnitte. Abteilung II. A. Aufgaben. [67 S.] gr. 8. geh.
n. M 1.20.

\rule{1.5em}{0.4pt} B. Auflösungen. [94 S.] gr. 8. geh. n. M 1.60.

\textbf{Hofmann, Fritz,} die Constructionen doppelt berührender Kegelschnitte mit
imaginären Bestimmungsstücken. Eine Wanderung durch die Theorie der Kegelschnitte in
doppelter Berührung an der Hand anschaulicher Methoden. [Mit Figuren im Text.]
[IV u. 109 S.] gr. 8. geh. n. M 3.20.

\textbf{Holzmüller, Dr. G.,} Direktor der Gewerbeschule zu Hagen, Einführung in das
stereometrische Zeichnen. Mit Berücksichtigung der Krystallographie und Kartographie.
[VIII u. 102 S. mit 16 lithographierten Tafeln.] gr. 8. kart. n. M 4.40.

\textbf{Januschke, Hans,} k. k. Professor an der Staats-Oberrealschule in Troppau,
das Princip der Erhaltung der Energie in der elementaren Elektricitätslehre.
[VIII u. 186 S.] gr. 8. geh. n. M 4.\,---

\textbf{Kohlrausch, Dr. F.,} ordentl. Professor an der Universität Würzburg, Leitfaden
der praktischen Physik mit einem Anhange: Das absolute Mafs-System. Mit in den Text
gedruckten Figuren. Sechste vermehrte Auflage. [XXIII u. 364 S.] gr. 8. geh. n. M
5.60, geb. M 6.10.

In leichtem Leinwandband [nach Art von Baedekers Reisebüchern] elegant gebundene
Exemplare n. M 6.10.

\textbf{Krause, Martin,} Professor der Mathematik an der Universität zu Rostock, die
Transformation der hyperelliptischen Functionen erster Ordnung. Nebst Anwendungen.
[VII u. 276 S.] gr. 8. geh. n. M 10.\,---

\textbf{Legendre, Adrien-Marie,} Zahlentheorie. Nach der dritten Auflage ins Deutsche
übertragen von H. Maser. Zwei Bände. [I. Bd. XVIII u. 442, II. Bd. XII u. 453 S.]
gr. 8. geh. n. M 23.20.

\textbf{Mansion, Dr. P.,} Professor an der Universität zu Gent, Elemente der Theorie
der Determinanten mit vielen Übungsaufgaben. Zweite vermehrte Aufl. [XXIV u. 56 S.]
gr. 8. geh. n. M 1.20.

\textbf{Reidt, Dr. Friedrich,} Professor am Gymnasium und dem Realprogymnasium in
Hamm, Sammlung von Aufgaben und Beispielen aus der Trigonometrie und Stereometrie.
Zweiter Teil: Stereometrie. Dritte Auflage. [VIII u. 190 S.] gr. 8. geh. M 3.\,---
%CONT>>


% ===== pdf idx 376, printed 360.002 =====
% !PRINTED none
% [no running-head folio number visible on this page image; it continues the
%  publisher's advertisement page (Verlagsanzeige) begun on the previous scan]
%<<CONT
\textbf{Reidt, Dr. Friedrich,} Professor am Gymnasium und dem Realprogymnasium in
Hamm, Resultate der Rechnungs-Aufgaben in der Sammlung von Aufgaben und Beispielen
aus der Trigonometrie und Stereometrie. I. Teil: Trigonometrie. Dritte Auflage.
[84 S.] gr. 8. geh. n. M 1.80.

\rule{1.5em}{0.4pt}\ \rule{1.5em}{0.4pt} II. Teil: Stereometrie. Dritte Auflage.
[48 S.] gr. 8. geh. M 1.\,---

\textbf{Schoenflies, Dr. Arthur,} Privatdocent der Mathematik an der Universität
Göttingen, Geometrie der Bewegung in synthetischer Darstellung. Mit Figuren im Text.
[VI u. 185 S.] gr. 8. geh. n. M 4.\,---

\textbf{Stolz, Dr. Otto,} ord. Professor an der Universität zu Innsbruck,
Vorlesungen über allgemeine Arithmetik. Nach den neueren Ansichten bearbeitet.
Erster Theil: Allgemeines und Arithmetik der reellen Zahlen. [VI u. 344 S.] gr. 8.
geh. n. M 8.\,---

\rule{1.5em}{0.4pt}\ \rule{1.5em}{0.4pt} Zweiter Theil: Arithmetik der complexen
Zahlen mit geometrischen Anwendungen. [VIII u. 326 S.] gr. 8. geh. n. M 8.\,---

\textbf{Wertheim, Gustav,} Elemente der Zahlentheorie. [IX u. 382 S.] gr. 8. geh.
n. M 8.40.

\textbf{Wiener, Dr. Christ.,} Geh. Hofrath und Professor an der Grossh. technischen
Hochschule zu Karlsruhe, Lehrbuch der darstellenden Geometrie. In zwei Bänden.
Zweiter Band. Krumme Linien (zweiter Theil) und krumme Flächen. Beobachtungslehre,
Perspektive. Mit Figuren im Text. [XXX u. 649 S.] gr. 8. geh. M 18.\,---

\textbf{Wüllner, Dr. Adolph,} Professor der Physik an der Königl. technischen
Hochschule zu Aachen, Lehrbuch der Experimentalphysik. Vierter Band: Die Lehre vom
Magnetismus und von der Electricität. Mit einer Einleitung \glqq Grundzüge der Lehre
vom Potential\grqq{}. Vierte vielfach umgearbeitete und verbesserte Auflage. (Mit
vielen in den Text gedruckten Holzschnitten.) [XII u. 1231 S.] gr. 8. geh. n.
M 16.80. --- Preis des ganzen Werkes M 48.80.

% [horizontal rule printed below this line, marking the end of the advertisement list]

\end{document}
